\documentclass[11pt]{article}
\usepackage[T1]{fontenc}
\usepackage{amsmath}
\usepackage{amssymb}
\usepackage{amsthm}
\usepackage[hidelinks]{hyperref}
\newtheorem{theorem}{Theorem}[section]
\newtheorem{lemma}[theorem]{Lemma}
\newtheorem{corollary}[theorem]{Corollary}
\theoremstyle{definition}
\newtheorem{definition}[theorem]{Definition}
\begin{document}
\begin{center}
{\LARGE Rule thirty: the two places}
\end{center}
\section{Local laws}
\label{ll:paper:local-laws}
\begin{lemma}
\label{ll:paper:white-interior}
For every truth value \(a\) and truth value \(b\) and truth value \(c\), if \(a = \operatorname{false}\), then if \(b = \operatorname{false}\), then if \(c = \operatorname{false}\), then \(\operatorname{rule}(a, b, c) = \operatorname{false}\).
\end{lemma}
\begin{proof}
Assume \(ha\), \(hb\), \(hc\).
Rewrite the goal using \(\rightarrow ha\), \(\rightarrow hb\), \(\rightarrow hc\).
The goal follows by reflexivity.
\end{proof}
\begin{lemma}
\label{ll:paper:white-run-interior}
For every natural number \(t\) and integer \(p\), if \(\operatorname{cell}(t, \operatorname{leftof}(p)) = \operatorname{false}\), then if \(\operatorname{cell}(t, p) = \operatorname{false}\), then if \(\operatorname{cell}(t, \operatorname{rightof}(p)) = \operatorname{false}\), then \(\operatorname{cell}(\operatorname{succ}(t), p) = \operatorname{false}\).
\end{lemma}
\begin{proof}
Assume \(hl\), \(hc\), \(hr\).
Rewrite the goal using \(\rightarrow \operatorname{cellsucc}(t, p)\), \(\rightarrow hl\), \(\rightarrow hc\), \(\rightarrow hr\).
The goal follows by reflexivity.
\end{proof}
\begin{lemma}
\label{ll:paper:black-run-interior}
For every natural number \(t\) and integer \(p\), if \(\operatorname{cell}(t, \operatorname{leftof}(p)) = \operatorname{true}\), then if \(\operatorname{cell}(t, p) = \operatorname{true}\), then \(\operatorname{cell}(\operatorname{succ}(t), p) = \operatorname{false}\).
\end{lemma}
\begin{proof}
Assume \(hl\), \(hc\).
Rewrite the goal using \(\rightarrow \operatorname{cellsucc}(t, p)\), \(\rightarrow hl\), \(\rightarrow hc\).
The goal follows from \(\operatorname{ruleblackblack}(\operatorname{cell}(t, \operatorname{rightof}(p)))\).
\end{proof}
\begin{lemma}
\label{ll:paper:white-pair-right}
For every natural number \(t\) and integer \(p\), if \(\operatorname{cell}(t, p) = \operatorname{false}\), then if \(\operatorname{cell}(t, \operatorname{rightof}(p)) = \operatorname{false}\), then \(\operatorname{cell}(\operatorname{succ}(t), p) = \operatorname{cell}(t, \operatorname{leftof}(p))\).
\end{lemma}
\begin{proof}
Assume \(hc\), \(hr\).
Rewrite the goal using \(\rightarrow \operatorname{cellsucc}(t, p)\), \(\rightarrow hc\), \(\rightarrow hr\), \(\rightarrow \operatorname{rulewhitewhite}(\operatorname{cell}(t, \operatorname{leftof}(p)))\).
\end{proof}
\begin{lemma}
\label{ll:paper:white-pair-left}
For every natural number \(t\) and integer \(p\), if \(\operatorname{cell}(t, \operatorname{leftof}(p)) = \operatorname{false}\), then if \(\operatorname{cell}(t, p) = \operatorname{false}\), then \(\operatorname{cell}(\operatorname{succ}(t), p) = \operatorname{cell}(t, \operatorname{rightof}(p))\).
\end{lemma}
\begin{proof}
Assume \(hl\), \(hc\).
Rewrite the goal using \(\rightarrow \operatorname{cellsucc}(t, p)\), \(\rightarrow hl\), \(\rightarrow hc\), \(\rightarrow \operatorname{ruleleftedge}(\operatorname{cell}(t, \operatorname{rightof}(p)))\).
\end{proof}
\begin{theorem}
\label{ll:paper:white-run-shrinks}
For every natural number \(t\) and integer \(a\) and integer \(b\), if for every integer \(q\), if \(a \le q\), then if \(q \le b\), then \(\operatorname{cell}(t, q) = \operatorname{false}\), then for every integer \(q\), if \(\operatorname{rightof}(a) \le q\), then if \(q \le \operatorname{leftof}(b)\), then \(\operatorname{cell}(\operatorname{succ}(t), q) = \operatorname{false}\).
\end{theorem}
\begin{proof}
Assume \(h\), \(q\), \(h1\), \(h2\).
Apply \(\texttt{Paper::white-run-interior}(t, q)\).
Premise \(1\):
\begin{quote}
The goal follows from \(h(\operatorname{leftof}(q), \operatorname{leleftof}(h1), \operatorname{leofrightof}(\operatorname{leftofle}(\operatorname{leofrightof}(\operatorname{rightofle}(h2)))))\).
\end{quote}
Premise \(2\):
\begin{quote}
The goal follows from \(h(q, \operatorname{leofrightof}(h1), \operatorname{leofrightof}(\operatorname{rightofle}(h2)))\).
\end{quote}
Premise \(3\):
\begin{quote}
The goal follows from \(h(\operatorname{rightof}(q), \operatorname{lerightof}(\operatorname{leofrightof}(h1)), \operatorname{rightofle}(h2))\).
\end{quote}
\end{proof}
\section{Diagonal laws}
\label{ll:paper:diagonal-laws}
\begin{lemma}
\label{ll:paper:reset}
For every natural number \(t\) and integer \(p\), if \(\operatorname{cell}(t, p) = \operatorname{true}\), then \(\operatorname{cell}(\operatorname{succ}(t), p) = \operatorname{not}(\operatorname{cell}(t, \operatorname{leftof}(p)))\).
\end{lemma}
\begin{proof}
Assume \(hc\).
Rewrite the goal using \(\rightarrow \operatorname{cellsucc}(t, p)\), \(\rightarrow hc\).
The goal follows from \(\operatorname{rulereset}(\operatorname{cell}(t, \operatorname{leftof}(p)), \operatorname{cell}(t, \operatorname{rightof}(p)))\).
\end{proof}
\begin{lemma}
\label{ll:paper:toggle}
For every natural number \(t\) and integer \(p\), if \(\operatorname{cell}(t, p) = \operatorname{false}\), then \(\operatorname{cell}(\operatorname{succ}(t), p) = \operatorname{bxor}(\operatorname{cell}(t, \operatorname{leftof}(p)), \operatorname{cell}(t, \operatorname{rightof}(p)))\).
\end{lemma}
\begin{proof}
Assume \(hc\).
Rewrite the goal using \(\rightarrow \operatorname{cellsucc}(t, p)\), \(\rightarrow hc\).
The goal follows from \(\operatorname{ruletoggle}(\operatorname{cell}(t, \operatorname{leftof}(p)), \operatorname{cell}(t, \operatorname{rightof}(p)))\).
\end{proof}
\section{The left system}
\label{ll:paper:left-system}
\begin{lemma}
\label{ll:paper:edge-black}
For every natural number \(t\), \(\operatorname{ldiag}(0, t) = \operatorname{true}\).
\end{lemma}
\begin{proof}
Consider the cases of \(t\).
Case zero:
\begin{quote}
The goal follows by reflexivity.
\end{quote}
Case \(\operatorname{succ}\) with \(s\):
\begin{quote}
The goal follows from \(\operatorname{ldiagzero}(s)\).
\end{quote}
\end{proof}
\begin{lemma}
\label{ll:paper:diag-reset}
For every natural number \(k\) and natural number \(s\), if \(\operatorname{ldiag}(\operatorname{succ}(k), s) = \operatorname{true}\), then \(\operatorname{ldiag}(\operatorname{succ}(\operatorname{succ}(k)), \operatorname{succ}(s)) = \operatorname{not}(\operatorname{ldiag}(k, s))\).
\end{lemma}
\begin{proof}
Assume \(h\).
Rewrite the goal using \(\rightarrow \operatorname{ldiagstep}(k, s)\), \(\rightarrow h\).
The goal follows from \(\operatorname{bxorreset}(\operatorname{ldiag}(k, s), \operatorname{ldiag}(\operatorname{succ}(\operatorname{succ}(k)), s))\).
\end{proof}
\begin{lemma}
\label{ll:paper:diag-toggle}
For every natural number \(k\) and natural number \(s\), if \(\operatorname{ldiag}(\operatorname{succ}(k), s) = \operatorname{false}\), then \(\operatorname{ldiag}(\operatorname{succ}(\operatorname{succ}(k)), \operatorname{succ}(s)) = \operatorname{bxor}(\operatorname{ldiag}(k, s), \operatorname{ldiag}(\operatorname{succ}(\operatorname{succ}(k)), s))\).
\end{lemma}
\begin{proof}
Assume \(h\).
Rewrite the goal using \(\rightarrow \operatorname{ldiagstep}(k, s)\), \(\rightarrow h\).
The goal follows from \(\operatorname{bxortoggle}(\operatorname{ldiag}(k, s), \operatorname{ldiag}(\operatorname{succ}(\operatorname{succ}(k)), s))\).
\end{proof}
\begin{theorem}
\label{ll:paper:leading-digits}
For every natural number \(j\), there exists a natural number \(P\) such that \(0 < P\) and there exists a natural number \(T\) such that for every natural number \(t\), if \(T \le t\), then \(\operatorname{ldiag}(j, t + P) = \operatorname{ldiag}(j, t)\).
\end{theorem}
\begin{proof}
The goal follows from \(\operatorname{ldiagperiodic}(j)\).
\end{proof}
\section{Intent}
\label{ll:paper:intent}
\begin{theorem}
\label{ll:paper:gate-theorem}
If \(\operatorname{frontobligation}\), then if \(\operatorname{zeroobligation}\), then \(\operatorname{problemone}\).
\end{theorem}
\begin{proof}
Assume \(hf\), \(hz\).
The goal follows from \(\operatorname{gate}(hf, hz)\).
\end{proof}
\begin{theorem}
\label{ll:paper:recompile-points}
For every natural number \(N\), there exists a natural number \(z\) such that \(N \le z\) and \(\operatorname{zerodiagonal}(z)\).
\end{theorem}
\begin{proof}
The goal follows from \(\operatorname{zerodiagonals}(N)\).
\end{proof}
\begin{theorem}
\label{ll:paper:unbounded-periods}
For every natural number \(P\), if \(0 < P\), then not \(\operatorname{commonperiod}(P)\).
\end{theorem}
\begin{proof}
Assume \(h\).
The goal follows from \(\operatorname{nocommonperiod}(P, h)\).
\end{proof}
\section{Compile points}
\label{ll:paper:compile-points}
\begin{theorem}
\label{ll:paper:black-after-zero}
For every natural number \(z\), if \(\operatorname{zerodiagonal}(z)\), then \(\operatorname{blackdiagonal}(z + 2)\).
\end{theorem}
\begin{proof}
Assume \(h\).
The goal follows from \(\operatorname{blackafterzero}(z, h)\).
\end{proof}
\begin{theorem}
\label{ll:paper:zero-before-black}
For every natural number \(b\), if \(\operatorname{blackdiagonal}(b + 2)\), then \(\operatorname{zerodiagonal}(b)\).
\end{theorem}
\begin{proof}
Assume \(h\).
The goal follows from \(\operatorname{zerobeforeblack}(b, h)\).
\end{proof}
\begin{theorem}
\label{ll:paper:isolated-zeros}
For every natural number \(z\), if \(\operatorname{zerodiagonal}(z)\), then not \(\operatorname{zerodiagonal}(z + 1)\).
\end{theorem}
\begin{proof}
Assume \(h\).
The goal follows from \(\operatorname{zerothennonzero}(z, h)\).
\end{proof}
\begin{theorem}
\label{ll:paper:doubling-points}
For every natural number \(P\), if \(0 < P\), then there exists a natural number \(z\) such that \(\operatorname{zerodiagonal}(z)\) and \(\operatorname{doublesat}(P, z)\).
\end{theorem}
\begin{proof}
Assume \(h\).
The goal follows from \(\operatorname{doublingpoints}(P, h)\).
\end{proof}
\begin{theorem}
\label{ll:paper:antiperiodic-doubling}
For every natural number \(P\) and natural number \(z\), if \(\operatorname{zerodiagonal}(z)\), then if \(\operatorname{doublesat}(P, z)\), then \(\operatorname{antiperiodicdiagonal}(P, z + 1)\).
\end{theorem}
\begin{proof}
Assume \(hz\), \(hd\).
The goal follows from \(\operatorname{antiperiodicafterdoubles}(P, z, hz, hd)\).
\end{proof}
\begin{theorem}
\label{ll:paper:compile-bit}
\(\operatorname{compilebitlaw}\).
\end{theorem}
\begin{proof}
The goal follows from \(\operatorname{compilebit}\).
\end{proof}
\section{The right system}
\label{ll:paper:right-system}
\begin{lemma}
\label{ll:paper:right-law}
For every natural number \(k\) and natural number \(s\), \(\operatorname{rdiag}(\operatorname{succ}(\operatorname{succ}(k)), \operatorname{succ}(s)) = \operatorname{bxor}(\operatorname{rdiag}(\operatorname{succ}(\operatorname{succ}(k)), s), \operatorname{rdiag}(\operatorname{succ}(k), s) \operatorname{or} \operatorname{rdiag}(k, s))\).
\end{lemma}
\begin{proof}
The goal follows from \(\operatorname{rdiagstep}(k, s)\).
\end{proof}
\begin{theorem}
\label{ll:paper:right-periods}
For every natural number \(j\) and natural number \(t\), \(\operatorname{rdiag}(j, t + \operatorname{rpow}(j)) = \operatorname{rdiag}(j, t)\).
\end{theorem}
\begin{proof}
The goal follows from \(\operatorname{rdiagperiodic}(j, t)\).
\end{proof}
\begin{theorem}
\label{ll:paper:mask-form}
If \(\operatorname{maskobligation}\), then for every natural number \(T\) and natural number \(P\), if \(0 < P\), then if \(\operatorname{readoutperiodic}(T, P)\), then \(\operatorname{maskedperiodic}(T, P)\).
\end{theorem}
\begin{proof}
Assume \(h\).
The goal follows from \(\operatorname{maskunfold}(h)\).
\end{proof}
\begin{theorem}
\label{ll:paper:mask-sufficient}
If \(\operatorname{maskobligation}\), then \(\operatorname{problemone}\).
\end{theorem}
\begin{proof}
Assume \(hm\).
The goal follows from \(\operatorname{problemoneofmaskpos}(hm)\).
\end{proof}
\begin{theorem}
\label{ll:paper:mask-necessary}
If \(\operatorname{problemone}\), then \(\operatorname{maskobligation}\).
\end{theorem}
\begin{proof}
Assume \(h\).
The goal follows from \(\operatorname{maskposofproblemone}(h)\).
\end{proof}
\begin{theorem}
\label{ll:paper:center-reproduces}
For every natural number \(t\), \(\operatorname{centercol}(t + 1) = \operatorname{phicenter}(t)\).
\end{theorem}
\begin{proof}
The goal follows from \(\operatorname{centerfixed}(t)\).
\end{proof}
\begin{theorem}
\label{ll:paper:fixed-sufficient}
If \(\operatorname{noperiodicfixed}\), then \(\operatorname{problemone}\).
\end{theorem}
\begin{proof}
Assume \(h\).
The goal follows from \(\operatorname{problemoneofnpf}(h)\).
\end{proof}
\begin{theorem}
\label{ll:paper:fixed-necessary}
If \(\operatorname{problemone}\), then \(\operatorname{noperiodicfixed}\).
\end{theorem}
\begin{proof}
Assume \(h\).
The goal follows from \(\operatorname{npfofproblemone}(h)\).
\end{proof}
\section{Projection onto known rules}
\label{ll:paper:known-rules}
\begin{theorem}
\label{ll:paper:split-onefifty}
For every truth value \(a\) and truth value \(b\) and truth value \(c\), \(\operatorname{rule}(a, b, c) = \operatorname{bxor}(\operatorname{ruleonefifty}(a, b, c), b \operatorname{and} c)\).
\end{theorem}
\begin{proof}
The goal follows from \(\operatorname{splitonefifty}(a, b, c)\).
\end{proof}
\begin{theorem}
\label{ll:paper:split-ninety}
For every truth value \(a\) and truth value \(b\) and truth value \(c\), \(\operatorname{rule}(a, b, c) = \operatorname{bxor}(\operatorname{ruleninety}(a, b, c), b \operatorname{and} \operatorname{not}(c))\).
\end{theorem}
\begin{proof}
The goal follows from \(\operatorname{splitninety}(a, b, c)\).
\end{proof}
\begin{theorem}
\label{ll:paper:split-sixty}
For every truth value \(a\) and truth value \(b\) and truth value \(c\), \(\operatorname{rule}(a, b, c) = \operatorname{bxor}(\operatorname{rulesixty}(a, b, c), \operatorname{not}(b) \operatorname{and} c)\).
\end{theorem}
\begin{proof}
The goal follows from \(\operatorname{splitsixty}(a, b, c)\).
\end{proof}
\begin{theorem}
\label{ll:paper:split-fifteen}
For every truth value \(a\) and truth value \(b\) and truth value \(c\), \(\operatorname{rule}(a, b, c) = \operatorname{bxor}(\operatorname{rulefifteen}(a, b, c), \operatorname{not}(b) \operatorname{and} \operatorname{not}(c))\).
\end{theorem}
\begin{proof}
The goal follows from \(\operatorname{splitfifteen}(a, b, c)\).
\end{proof}
\begin{theorem}
\label{ll:paper:deep-hole}
For every truth value \(u0\) and truth value \(u1\) and truth value \(u2\) and truth value \(u3\), \(2 \le \operatorname{distthirty}(u0, u1, u2, u3)\).
\end{theorem}
\begin{proof}
The goal follows from \(\operatorname{deephole}(u0, u1, u2, u3)\).
\end{proof}
\begin{theorem}
\label{ll:paper:center-response}
For every natural number \(t\), \(\operatorname{centercol}(t) = \operatorname{not}(\operatorname{respsixty}(t))\).
\end{theorem}
\begin{proof}
The goal follows from \(\operatorname{centerrespsixty}(t)\).
\end{proof}
\begin{theorem}
\label{ll:paper:orthogonal-responses}
For every natural number \(t\), if \(1 \le t\), then \(\operatorname{bxor}(\operatorname{respsixty}(t), \operatorname{respfifteen}(t)) = \operatorname{eventime}(t)\).
\end{theorem}
\begin{proof}
Assume \(h\).
The goal follows from \(\operatorname{resprelation}(t, h)\).
\end{proof}
\begin{theorem}
\label{ll:paper:projection-sixty-sufficient}
If \(\operatorname{respsixtyaperiodic}\), then \(\operatorname{problemone}\).
\end{theorem}
\begin{proof}
Assume \(h\).
The goal follows from \(\operatorname{poneofrsixty}(h)\).
\end{proof}
\begin{theorem}
\label{ll:paper:projection-sixty-necessary}
If \(\operatorname{problemone}\), then \(\operatorname{respsixtyaperiodic}\).
\end{theorem}
\begin{proof}
Assume \(h\).
The goal follows from \(\operatorname{rsixtyofpone}(h)\).
\end{proof}
\begin{theorem}
\label{ll:paper:projection-fifteen-sufficient}
If \(\operatorname{respfifteenaperiodic}\), then \(\operatorname{problemone}\).
\end{theorem}
\begin{proof}
Assume \(h\).
The goal follows from \(\operatorname{poneofrfifteen}(h)\).
\end{proof}
\begin{theorem}
\label{ll:paper:projection-fifteen-necessary}
If \(\operatorname{problemone}\), then \(\operatorname{respfifteenaperiodic}\).
\end{theorem}
\begin{proof}
Assume \(h\).
The goal follows from \(\operatorname{rfifteenofpone}(h)\).
\end{proof}
\begin{theorem}
\label{ll:paper:displacement-carry}
For every truth value \(a\) and truth value \(b\) and truth value \(c\), \(\operatorname{displacement}(a, b, c) = b \operatorname{and} c\).
\end{theorem}
\begin{proof}
The goal follows from \(\operatorname{displacementiscarry}(a, b, c)\).
\end{proof}
\begin{theorem}
\label{ll:paper:frame-direction}
For every truth value \(a\) and truth value \(b\) and truth value \(c\), \(\operatorname{displacement}(\operatorname{not}(a), b, c) = \operatorname{displacement}(a, b, c)\).
\end{theorem}
\begin{proof}
The goal follows from \(\operatorname{framedirection}(a, b, c)\).
\end{proof}
\begin{theorem}
\label{ll:paper:affine-left-linear}
For every truth value \(a\) and truth value \(b\) and truth value \(c\), \(\operatorname{ruleonefifty}(\operatorname{not}(a), b, c) = \operatorname{not}(\operatorname{ruleonefifty}(a, b, c))\).
\end{theorem}
\begin{proof}
The goal follows from \(\operatorname{affineleftlinear}(a, b, c)\).
\end{proof}
\begin{theorem}
\label{ll:paper:irrational-sufficient}
If \(\operatorname{irrationalthirty}\), then \(\operatorname{problemone}\).
\end{theorem}
\begin{proof}
Assume \(h\).
The goal follows from \(\operatorname{poneofirr}(h)\).
\end{proof}
\begin{theorem}
\label{ll:paper:irrational-necessary}
If \(\operatorname{problemone}\), then \(\operatorname{irrationalthirty}\).
\end{theorem}
\begin{proof}
Assume \(h\).
The goal follows from \(\operatorname{irrofpone}(h)\).
\end{proof}
\begin{theorem}
\label{ll:paper:frame-projection}
For every truth value \(a\) and truth value \(b\) and truth value \(c\), \(\operatorname{iadd}(\operatorname{sgn}(\operatorname{rule}(a, b, c)), \operatorname{sgn}(\operatorname{rule}(a, b, c))) = \operatorname{iadd}(\operatorname{iadd}(\operatorname{sgn}(\operatorname{ruleonefifty}(a, b, c)), \operatorname{sgn}(\operatorname{ruleninety}(a, b, c))), \operatorname{iadd}(\operatorname{sgn}(\operatorname{rulesixty}(a, b, c)), \operatorname{sgn}(\operatorname{rulefifteen}(a, b, c))))\).
\end{theorem}
\begin{proof}
The goal follows from \(\operatorname{frameprojection}(a, b, c)\).
\end{proof}
\section{The dyadic scale}
\label{ll:paper:dyadic-scale}
\begin{lemma}
\label{ll:paper:window-law}
For every natural number \(t0\) and natural number \(m\) and integer \(p\), \(\operatorname{window}(t0, \operatorname{succ}(m), p) = \operatorname{bxor}(\operatorname{bxor}(\operatorname{window}(t0, m, \operatorname{leftof}(p)), \operatorname{window}(t0, m, p)), \operatorname{carry}(t0 + m, p))\).
\end{lemma}
\begin{proof}
The goal follows from \(\operatorname{windowstep}(t0, m, p)\).
\end{proof}
\begin{lemma}
\label{ll:paper:scale-doubles}
For every natural number \(j\), \(\operatorname{twopow}(\operatorname{succ}(j)) = \operatorname{twopow}(j) + \operatorname{twopow}(j)\).
\end{lemma}
\begin{proof}
The goal follows from \(\operatorname{twopowsucc}(j)\).
\end{proof}
\begin{theorem}
\label{ll:paper:scale-identity}
For every natural number \(t\) and natural number \(k\), \(\operatorname{center}(t + \operatorname{twopow}(k)) = \operatorname{bxor}(\operatorname{bxor}(\operatorname{center}(t), \operatorname{cell}(t, \operatorname{negpow}(k))), \operatorname{response}(t, k))\).
\end{theorem}
\begin{proof}
The goal follows from \(\operatorname{scalecenter}(t, k)\).
\end{proof}
\begin{theorem}
\label{ll:paper:periodic-scale}
For every natural number \(T\) and natural number \(P\), if \(0 < P\), then if \(\operatorname{centerperiodic}(T, P)\), then for every natural number \(k\) and natural number \(t\), if \(T \le t\), then \(\operatorname{bxor}(\operatorname{cell}(t, \operatorname{negpow}(k)), \operatorname{response}(t, k)) = \operatorname{bxor}(\operatorname{center}(t + \operatorname{lag}(P, k)), \operatorname{center}(t))\).
\end{theorem}
\begin{proof}
Assume \(hP\), \(h\).
The goal follows from \(\operatorname{periodicscale}(T, P, hP, h)\).
\end{proof}
\begin{theorem}
\label{ll:paper:periodic-scale-cone}
For every natural number \(T\) and natural number \(P\), if \(0 < P\), then if \(\operatorname{centerperiodic}(T, P)\), then for every natural number \(k\) and natural number \(t\), if \(T \le t\), then if \(t < \operatorname{twopow}(k)\), then \(\operatorname{response}(t, k) = \operatorname{bxor}(\operatorname{center}(t + \operatorname{lag}(P, k)), \operatorname{center}(t))\).
\end{theorem}
\begin{proof}
Assume \(hP\), \(h\).
The goal follows from \(\operatorname{periodicscalecone}(T, P, hP, h)\).
\end{proof}
\begin{theorem}
\label{ll:paper:response-split}
For every natural number \(t\) and natural number \(k\), \(\operatorname{response}(t, k) = \operatorname{bxor}(\operatorname{leftresponse}(t, k), \operatorname{maskparity}(t, k))\).
\end{theorem}
\begin{proof}
The goal follows from \(\operatorname{split}(t, k)\).
\end{proof}
\begin{theorem}
\label{ll:paper:periodic-mask}
For every natural number \(T\) and natural number \(P\), if \(0 < P\), then if \(\operatorname{centerperiodic}(T, P)\), then for every natural number \(k\) and natural number \(t\), if \(T \le t\), then if \(t < \operatorname{twopow}(k)\), then \(\operatorname{maskparity}(t, k) = \operatorname{bxor}(\operatorname{leftresponse}(t, k), \operatorname{bxor}(\operatorname{center}(t + \operatorname{lag}(P, k)), \operatorname{center}(t)))\).
\end{theorem}
\begin{proof}
Assume \(hP\), \(h\).
The goal follows from \(\operatorname{periodicmask}(T, P, hP, h)\).
\end{proof}
\begin{theorem}
\label{ll:paper:noperiod-sufficient}
If \(\operatorname{noperiod}\), then \(\operatorname{problemone}\).
\end{theorem}
\begin{proof}
Assume \(h\).
The goal follows from \(\operatorname{poneofnoperiod}(h)\).
\end{proof}
\begin{theorem}
\label{ll:paper:noperiod-necessary}
If \(\operatorname{problemone}\), then \(\operatorname{noperiod}\).
\end{theorem}
\begin{proof}
Assume \(h\).
The goal follows from \(\operatorname{noperiodofpone}(h)\).
\end{proof}
\section{Witt forms and linear interfaces}
\label{ll:paper:witt-interfaces}
\begin{theorem}
\label{ll:paper:witt-split}
For every natural number \(y\), \(\operatorname{wittphi}(y) = \operatorname{wittxor}(\operatorname{wittlinear}(y), \operatorname{wittshl}(\operatorname{wittcarry}(y), 1))\).
\end{theorem}
\begin{proof}
The goal follows from \(\operatorname{phisplit}(y)\).
\end{proof}
\begin{theorem}
\label{ll:paper:parity-interface}
For every natural number \(y\), \(\operatorname{rowpar}(\operatorname{wittphi}(y)) = \operatorname{bxor}(\operatorname{rowpar}(y), \operatorname{rowpar}(\operatorname{wittcarry}(y)))\).
\end{theorem}
\begin{proof}
The goal follows from \(\operatorname{parphi}(y)\).
\end{proof}
\begin{theorem}
\label{ll:paper:cyclotomic-first-interface}
For every natural number \(y\), \(\operatorname{vala}(\operatorname{wittphi}(y)) = \operatorname{valb}(\operatorname{wittcarry}(y))\).
\end{theorem}
\begin{proof}
The goal follows from \(\operatorname{valaphi}(y)\).
\end{proof}
\begin{theorem}
\label{ll:paper:cyclotomic-second-interface}
For every natural number \(y\), \(\operatorname{valb}(\operatorname{wittphi}(y)) = \operatorname{bxor}(\operatorname{vala}(\operatorname{wittcarry}(y)), \operatorname{valb}(\operatorname{wittcarry}(y)))\).
\end{theorem}
\begin{proof}
The goal follows from \(\operatorname{valbphi}(y)\).
\end{proof}
\begin{theorem}
\label{ll:paper:witt-row-rdiag}
For every natural number \(t\) and natural number \(j\), \(\operatorname{wittbit}(\operatorname{wittrow}(t), j) = \operatorname{bton}(\operatorname{rdiag}(j, t))\).
\end{theorem}
\begin{proof}
The goal follows from \(\operatorname{bitrowrdiag}(t, j)\).
\end{proof}
\begin{theorem}
\label{ll:paper:witt-center-rcenter}
For every natural number \(t\), \(\operatorname{wittcenter}(t) = \operatorname{bton}(\operatorname{rcenter}(t))\).
\end{theorem}
\begin{proof}
The goal follows from \(\operatorname{rowcenterrcenter}(t)\).
\end{proof}
\begin{theorem}
\label{ll:paper:witt-center-cell}
For every natural number \(t\), \(\operatorname{wittcenter}(t) = \operatorname{bton}(\operatorname{center}(t))\).
\end{theorem}
\begin{proof}
The goal follows from \(\operatorname{rowcentercenter}(t)\).
\end{proof}
\begin{theorem}
\label{ll:paper:witt-noperiod-sufficient}
If \(\operatorname{norowcenterperiod}\), then \(\operatorname{problemone}\).
\end{theorem}
\begin{proof}
Assume \(h\).
The goal follows from \(\operatorname{poneofnorowcenterperiod}(h)\).
\end{proof}
\begin{theorem}
\label{ll:paper:witt-noperiod-necessary}
If \(\operatorname{problemone}\), then \(\operatorname{norowcenterperiod}\).
\end{theorem}
\begin{proof}
Assume \(h\).
The goal follows from \(\operatorname{norowcenterperiodofpone}(h)\).
\end{proof}
\begin{theorem}
\label{ll:paper:witt-freshman}
For every natural number \(k\) and natural number \(y\), \(\operatorname{wittlinpow}(\operatorname{twopow}(k), y) = \operatorname{wittscale}(k, y)\).
\end{theorem}
\begin{proof}
The goal follows from \(\operatorname{wittfreshman}(k, y)\).
\end{proof}
\begin{theorem}
\label{ll:paper:witt-duhamel}
For every natural number \(t0\) and natural number \(n\), \(\operatorname{wittrow}(t0 + n) = \operatorname{wittxor}(\operatorname{wittlinpow}(n, \operatorname{wittrow}(t0)), \operatorname{wittwin}(t0, n))\).
\end{theorem}
\begin{proof}
The goal follows from \(\operatorname{wittduhamel}(t0, n)\).
\end{proof}
\begin{theorem}
\label{ll:paper:witt-scale-row}
For every natural number \(t0\) and natural number \(k\), \(\operatorname{wittrow}(t0 + \operatorname{twopow}(k)) = \operatorname{wittxor}(\operatorname{wittscale}(k, \operatorname{wittrow}(t0)), \operatorname{wittwin}(t0, \operatorname{twopow}(k)))\).
\end{theorem}
\begin{proof}
The goal follows from \(\operatorname{wittrowscale}(t0, k)\).
\end{proof}
\begin{theorem}
\label{ll:paper:witt-scale-cone}
For every natural number \(t\) and natural number \(k\), if \(t < \operatorname{twopow}(k)\), then \(\operatorname{center}(t + \operatorname{twopow}(k)) = \operatorname{bxor}(\operatorname{center}(t), \operatorname{wittresp}(t, k))\).
\end{theorem}
\begin{proof}
The goal follows from \(\operatorname{wittconescale}(t, k)\).
\end{proof}
\begin{theorem}
\label{ll:paper:witt-cone-duality}
For every natural number \(t\) and natural number \(k\), if \(t < \operatorname{twopow}(k)\), then \(\operatorname{wittresp}(t, k) = \operatorname{response}(t, k)\).
\end{theorem}
\begin{proof}
The goal follows from \(\operatorname{wittconeduality}(t, k)\).
\end{proof}
\begin{theorem}
\label{ll:paper:witt-late-center}
For every natural number \(d\) and natural number \(k\), if \(d \le k\), then \(\operatorname{center}(d) = \operatorname{bxor}(\operatorname{wittresp}(d + \operatorname{twopow}(k), k), \operatorname{response}(d + \operatorname{twopow}(k), k))\).
\end{theorem}
\begin{proof}
The goal follows from \(\operatorname{wittlatecenter}(d, k)\).
\end{proof}
\begin{theorem}
\label{ll:paper:two-places-witt}
For every natural number \(d\) and natural number \(k\), if \(d \le k\), then \(\operatorname{bxor}(\operatorname{center}(d), \operatorname{ldiag}(d, d + \operatorname{twopow}(k))) = \operatorname{bxor}(\operatorname{bxor}(\operatorname{center}(d + \operatorname{twopow}(k) + \operatorname{twopow}(k)), \operatorname{center}(d + \operatorname{twopow}(k))), \operatorname{wittresp}(d + \operatorname{twopow}(k), k))\).
\end{theorem}
\begin{proof}
The goal follows from \(\operatorname{twoplaceswitt}(d, k)\).
\end{proof}
\begin{theorem}
\label{ll:paper:maskparity-coboundary}
For every natural number \(d\) and natural number \(k\), \(\operatorname{bxor}(\operatorname{bxor}(\operatorname{maskparity}(d, \operatorname{succ}(k)), \operatorname{maskparity}(d, k)), \operatorname{maskparity}(d + \operatorname{twopow}(k), k)) = \operatorname{false}\).
\end{theorem}
\begin{proof}
The goal follows from \(\operatorname{maskparitycoboundary}(d, k)\).
\end{proof}
\begin{theorem}
\label{ll:paper:ldiag-response-coboundary}
For every natural number \(d\) and natural number \(k\), if \(d \le k\), then \(\operatorname{ldiag}(d, d + \operatorname{twopow}(k)) = \operatorname{bxor}(\operatorname{bxor}(\operatorname{response}(d, \operatorname{succ}(k)), \operatorname{response}(d, k)), \operatorname{response}(d + \operatorname{twopow}(k), k))\).
\end{theorem}
\begin{proof}
The goal follows from \(\operatorname{ldiagresponsecoboundary}(d, k)\).
\end{proof}
\begin{theorem}
\label{ll:paper:ldiag-leftresponse-coboundary}
For every natural number \(d\) and natural number \(k\), if \(d \le k\), then \(\operatorname{ldiag}(d, d + \operatorname{twopow}(k)) = \operatorname{bxor}(\operatorname{bxor}(\operatorname{leftresponse}(d, \operatorname{succ}(k)), \operatorname{leftresponse}(d, k)), \operatorname{leftresponse}(d + \operatorname{twopow}(k), k))\).
\end{theorem}
\begin{proof}
The goal follows from \(\operatorname{ldiagleftresponsecoboundary}(d, k)\).
\end{proof}
\begin{theorem}
\label{ll:paper:two-places-witt-coboundary}
For every natural number \(d\) and natural number \(k\), if \(d \le k\), then \(\operatorname{bxor}(\operatorname{center}(d), \operatorname{ldiag}(d, d + \operatorname{twopow}(k))) = \operatorname{bxor}(\operatorname{bxor}(\operatorname{wittresp}(d, \operatorname{succ}(k)), \operatorname{wittresp}(d, k)), \operatorname{wittresp}(d + \operatorname{twopow}(k), k))\).
\end{theorem}
\begin{proof}
The goal follows from \(\operatorname{twoplaceswittcoboundary}(d, k)\).
\end{proof}
\begin{theorem}
\label{ll:paper:rdiag-scale}
For every natural number \(d\) and natural number \(s\) and natural number \(k\), if \(d \le k\), then \(\operatorname{rdiag}(d, s + \operatorname{twopow}(k)) = \operatorname{rdiag}(d, s)\).
\end{theorem}
\begin{proof}
The goal follows from \(\operatorname{rdiagscaleeq}(d, s, k)\).
\end{proof}
\begin{theorem}
\label{ll:paper:two-places-settled}
For every natural number \(d\) and natural number \(k\), if \(d \le k\), then if \(\operatorname{settledat}(d, k)\), then \(\operatorname{damageeq}(d, \operatorname{bxor}(\operatorname{bxor}(\operatorname{center}(d + \operatorname{twopow}(k) + \operatorname{twopow}(k)), \operatorname{center}(d + \operatorname{twopow}(k))), \operatorname{wittresp}(d + \operatorname{twopow}(k), k)))\).
\end{theorem}
\begin{proof}
The goal follows from \(\operatorname{twoplacessettled}(d, k)\).
\end{proof}
\begin{theorem}
\label{ll:paper:zcenter-leftresponse-coboundary}
For every natural number \(d\) and natural number \(k\), if \(d \le k\), then if \(\operatorname{settledat}(d, k)\), then \(\operatorname{zcentereq}(d, \operatorname{bxor}(\operatorname{bxor}(\operatorname{leftresponse}(d, \operatorname{succ}(k)), \operatorname{leftresponse}(d, k)), \operatorname{leftresponse}(d + \operatorname{twopow}(k), k)))\).
\end{theorem}
\begin{proof}
The goal follows from \(\operatorname{zcenterleftresponsecoboundary}(d, k)\).
\end{proof}
\begin{theorem}
\label{ll:paper:damage-witt-coboundary}
For every natural number \(d\) and natural number \(k\), if \(d \le k\), then if \(\operatorname{settledat}(d, k)\), then \(\operatorname{damageeq}(d, \operatorname{bxor}(\operatorname{bxor}(\operatorname{wittresp}(d, \operatorname{succ}(k)), \operatorname{wittresp}(d, k)), \operatorname{wittresp}(d + \operatorname{twopow}(k), k)))\).
\end{theorem}
\begin{proof}
The goal follows from \(\operatorname{damagewittcoboundary}(d, k)\).
\end{proof}
\begin{theorem}
\label{ll:paper:zcenter-zero}
For every natural number \(z\), if \(\operatorname{zerodiagonal}(z)\), then \(\operatorname{zcentereq}(z, \operatorname{false})\).
\end{theorem}
\begin{proof}
Assume \(hz\).
The goal follows from \(\operatorname{zcenterofzero}(z, hz)\).
\end{proof}
\begin{theorem}
\label{ll:paper:zcenter-after-zero}
For every natural number \(z\), if \(\operatorname{zerodiagonal}(z)\), then \(\operatorname{zcentereq}(z + 2, \operatorname{true})\).
\end{theorem}
\begin{proof}
Assume \(hz\).
The goal follows from \(\operatorname{zcenterafterzero}(z, hz)\).
\end{proof}
\begin{theorem}
\label{ll:paper:zcenter-oscillates}
For every natural number \(N\), there exists a natural number \(z\) such that \(N \le z\) and \(\operatorname{zcentereq}(z, \operatorname{false})\) and \(\operatorname{zcentereq}(z + 2, \operatorname{true})\).
\end{theorem}
\begin{proof}
The goal follows from \(\operatorname{zcenteroscillates}(N)\).
\end{proof}
\begin{theorem}
\label{ll:paper:leftresponse-coboundary-zero}
For every natural number \(z\) and natural number \(k\), if \(\operatorname{zerodiagonal}(z)\), then if \(z \le k\), then if \(\operatorname{settledat}(z, k)\), then \(\operatorname{bxor}(\operatorname{bxor}(\operatorname{leftresponse}(z, \operatorname{succ}(k)), \operatorname{leftresponse}(z, k)), \operatorname{leftresponse}(z + \operatorname{twopow}(k), k)) = \operatorname{false}\).
\end{theorem}
\begin{proof}
The goal follows from \(\operatorname{leftresponsecoboundaryzero}(z, k)\).
\end{proof}
\begin{theorem}
\label{ll:paper:leftresponse-coboundary-after-zero}
For every natural number \(z\) and natural number \(k\), if \(\operatorname{zerodiagonal}(z)\), then if \(z + 2 \le k\), then if \(\operatorname{settledat}(z + 2, k)\), then \(\operatorname{bxor}(\operatorname{bxor}(\operatorname{leftresponse}(z + 2, \operatorname{succ}(k)), \operatorname{leftresponse}(z + 2, k)), \operatorname{leftresponse}(z + 2 + \operatorname{twopow}(k), k)) = \operatorname{true}\).
\end{theorem}
\begin{proof}
The goal follows from \(\operatorname{leftresponsecoboundaryafterzero}(z, k)\).
\end{proof}
\begin{theorem}
\label{ll:paper:witt-coboundary-zero}
For every natural number \(z\) and natural number \(k\), if \(\operatorname{zerodiagonal}(z)\), then if \(z \le k\), then if \(\operatorname{settledat}(z, k)\), then \(\operatorname{bxor}(\operatorname{bxor}(\operatorname{wittresp}(z, \operatorname{succ}(k)), \operatorname{wittresp}(z, k)), \operatorname{wittresp}(z + \operatorname{twopow}(k), k)) = \operatorname{center}(z)\).
\end{theorem}
\begin{proof}
The goal follows from \(\operatorname{wittcoboundaryzero}(z, k)\).
\end{proof}
\begin{theorem}
\label{ll:paper:witt-coboundary-after-zero}
For every natural number \(z\) and natural number \(k\), if \(\operatorname{zerodiagonal}(z)\), then if \(z + 2 \le k\), then if \(\operatorname{settledat}(z + 2, k)\), then \(\operatorname{bxor}(\operatorname{bxor}(\operatorname{wittresp}(z + 2, \operatorname{succ}(k)), \operatorname{wittresp}(z + 2, k)), \operatorname{wittresp}(z + 2 + \operatorname{twopow}(k), k)) = \operatorname{not}(\operatorname{center}(z + 2))\).
\end{theorem}
\begin{proof}
The goal follows from \(\operatorname{wittcoboundaryafterzero}(z, k)\).
\end{proof}
\begin{theorem}
\label{ll:paper:periodic-response-lag}
For every natural number \(T\) and natural number \(P\), if \(0 < P\), then if \(\operatorname{centerperiodic}(T, P)\), then for every natural number \(k1\) and natural number \(k2\) and natural number \(t\), if \(\operatorname{lag}(P, k1) = \operatorname{lag}(P, k2)\), then if \(T \le t\), then if \(t < \operatorname{twopow}(k1)\), then if \(t < \operatorname{twopow}(k2)\), then \(\operatorname{response}(t, k1) = \operatorname{response}(t, k2)\).
\end{theorem}
\begin{proof}
Assume \(hP\), \(h\).
The goal follows from \(\operatorname{periodicresponselag}(T, P, hP, h)\).
\end{proof}
\begin{theorem}
\label{ll:paper:periodic-witt-lag}
For every natural number \(T\) and natural number \(P\), if \(0 < P\), then if \(\operatorname{centerperiodic}(T, P)\), then for every natural number \(k1\) and natural number \(k2\) and natural number \(t\), if \(\operatorname{lag}(P, k1) = \operatorname{lag}(P, k2)\), then if \(T \le t\), then if \(t < \operatorname{twopow}(k1)\), then if \(t < \operatorname{twopow}(k2)\), then \(\operatorname{wittresp}(t, k1) = \operatorname{wittresp}(t, k2)\).
\end{theorem}
\begin{proof}
Assume \(hP\), \(h\).
The goal follows from \(\operatorname{periodicwittlag}(T, P, hP, h)\).
\end{proof}
\begin{theorem}
\label{ll:paper:periodic-left-mask-lag}
For every natural number \(T\) and natural number \(P\), if \(0 < P\), then if \(\operatorname{centerperiodic}(T, P)\), then for every natural number \(k1\) and natural number \(k2\) and natural number \(t\), if \(\operatorname{lag}(P, k1) = \operatorname{lag}(P, k2)\), then if \(T \le t\), then if \(t < \operatorname{twopow}(k1)\), then if \(t < \operatorname{twopow}(k2)\), then \(\operatorname{bxor}(\operatorname{leftresponse}(t, k1), \operatorname{leftresponse}(t, k2)) = \operatorname{bxor}(\operatorname{maskparity}(t, k1), \operatorname{maskparity}(t, k2))\).
\end{theorem}
\begin{proof}
Assume \(hP\), \(h\).
The goal follows from \(\operatorname{periodicleftmasklag}(T, P, hP, h)\).
\end{proof}
\begin{theorem}
\label{ll:paper:periodic-late-response-lag}
For every natural number \(T\) and natural number \(P\), if \(0 < P\), then if \(\operatorname{centerperiodic}(T, P)\), then for every natural number \(d\) and natural number \(k1\) and natural number \(k2\), if \(\operatorname{lag}(P, k1) = \operatorname{lag}(P, k2)\), then if \(T \le d\), then if \(d \le k1\), then if \(d \le k2\), then if \(\operatorname{settledat}(d, k1)\), then if \(\operatorname{settledat}(d, k2)\), then \(\operatorname{response}(d + \operatorname{twopow}(k1), k1) = \operatorname{response}(d + \operatorname{twopow}(k2), k2)\).
\end{theorem}
\begin{proof}
Assume \(hP\), \(h\).
The goal follows from \(\operatorname{periodicUDlateresponselag}(T, P, hP, h)\).
\end{proof}
\begin{theorem}
\label{ll:paper:periodic-late-witt-lag}
For every natural number \(T\) and natural number \(P\), if \(0 < P\), then if \(\operatorname{centerperiodic}(T, P)\), then for every natural number \(d\) and natural number \(k1\) and natural number \(k2\), if \(\operatorname{lag}(P, k1) = \operatorname{lag}(P, k2)\), then if \(T \le d\), then if \(d \le k1\), then if \(d \le k2\), then if \(\operatorname{settledat}(d, k1)\), then if \(\operatorname{settledat}(d, k2)\), then \(\operatorname{wittresp}(d + \operatorname{twopow}(k1), k1) = \operatorname{wittresp}(d + \operatorname{twopow}(k2), k2)\).
\end{theorem}
\begin{proof}
Assume \(hP\), \(h\).
The goal follows from \(\operatorname{periodiclatescalelag}(T, P, hP, h)\).
\end{proof}
\begin{theorem}
\label{ll:paper:periodic-late-left-mask-lag}
For every natural number \(T\) and natural number \(P\), if \(0 < P\), then if \(\operatorname{centerperiodic}(T, P)\), then for every natural number \(d\) and natural number \(k1\) and natural number \(k2\), if \(\operatorname{lag}(P, k1) = \operatorname{lag}(P, k2)\), then if \(T \le d\), then if \(d \le k1\), then if \(d \le k2\), then if \(\operatorname{settledat}(d, k1)\), then if \(\operatorname{settledat}(d, k2)\), then \(\operatorname{bxor}(\operatorname{leftresponse}(d + \operatorname{twopow}(k1), k1), \operatorname{leftresponse}(d + \operatorname{twopow}(k2), k2)) = \operatorname{bxor}(\operatorname{maskparity}(d + \operatorname{twopow}(k1), k1), \operatorname{maskparity}(d + \operatorname{twopow}(k2), k2))\).
\end{theorem}
\begin{proof}
Assume \(hP\), \(h\).
The goal follows from \(\operatorname{periodiclateleftmasklag}(T, P, hP, h)\).
\end{proof}
\begin{theorem}
\label{ll:paper:winwitt-ray-zero}
For every natural number \(t0\) and natural number \(k\) and natural number \(j\), if \(j < \operatorname{twopow}(k)\), then \(\operatorname{wittbit}(\operatorname{wittwin}(t0, \operatorname{twopow}(k)), j) = \operatorname{bton}(\operatorname{bxor}(\operatorname{rdiag}(j, t0 + \operatorname{twopow}(k)), \operatorname{rdiag}(j, t0)))\).
\end{theorem}
\begin{proof}
The goal follows from \(\operatorname{bitwinwittrayzero}(t0, k, j)\).
\end{proof}
\begin{theorem}
\label{ll:paper:winwitt-low-zero}
For every natural number \(t0\) and natural number \(k\) and natural number \(j\), if \(j \le k\), then \(\operatorname{wittbit}(\operatorname{wittwin}(t0, \operatorname{twopow}(k)), j) = 0\).
\end{theorem}
\begin{proof}
The goal follows from \(\operatorname{bitwinwittlow}(t0, k, j)\).
\end{proof}
\begin{theorem}
\label{ll:paper:winwitt-ray-one}
For every natural number \(t0\) and natural number \(k\) and natural number \(i\), if \(i < \operatorname{twopow}(k)\), then if \(2 \cdot t0 < i + \operatorname{twopow}(k)\), then \(\operatorname{wittbit}(\operatorname{wittwin}(t0, \operatorname{twopow}(k)), i + \operatorname{twopow}(k)) = \operatorname{bton}(\operatorname{bxor}(\operatorname{rdiag}(i + \operatorname{twopow}(k), t0 + \operatorname{twopow}(k)), \operatorname{rdiag}(i, t0)))\).
\end{theorem}
\begin{proof}
The goal follows from \(\operatorname{bitwinwittrayone}(t0, k, i)\).
\end{proof}
\begin{theorem}
\label{ll:paper:winwitt-ray-two}
For every natural number \(d\) and natural number \(m\) and natural number \(k\), if \(d + m < \operatorname{twopow}(k)\), then \(\operatorname{wittbit}(\operatorname{wittwin}(d + m, \operatorname{twopow}(k)), d + 2 \cdot m + \operatorname{twopow}(\operatorname{succ}(k))) = \operatorname{bton}(\operatorname{bxor}(\operatorname{ldiag}(d, d + m), \operatorname{ldiag}(d, d + m + \operatorname{twopow}(k))))\).
\end{theorem}
\begin{proof}
The goal follows from \(\operatorname{bitwinwittraytwo}(d, m, k)\).
\end{proof}
\begin{theorem}
\label{ll:paper:winwitt-archimedean}
For every natural number \(d\) and natural number \(k\), if \(d \le k\), then \(\operatorname{wittbit}(\operatorname{wittwin}(d, \operatorname{twopow}(k)), d + \operatorname{twopow}(\operatorname{succ}(k))) = \operatorname{bton}(\operatorname{bxor}(\operatorname{center}(d), \operatorname{ldiag}(d, d + \operatorname{twopow}(k))))\).
\end{theorem}
\begin{proof}
The goal follows from \(\operatorname{bitwinwittarch}(d, k)\).
\end{proof}
\begin{theorem}
\label{ll:paper:winwitt-settled-zero}
For every natural number \(d\) and natural number \(m\) and natural number \(k\), if \(d \le k\), then if \(d + m < \operatorname{twopow}(k)\), then if \(\operatorname{settledby}(d, d + m)\), then \(\operatorname{wittbit}(\operatorname{wittwin}(d + m, \operatorname{twopow}(k)), d + 2 \cdot m + \operatorname{twopow}(\operatorname{succ}(k))) = 0\).
\end{theorem}
\begin{proof}
The goal follows from \(\operatorname{bitwinwittsettledzero}(d, m, k)\).
\end{proof}
\begin{theorem}
\label{ll:paper:winwitt-degree}
For every natural number \(t0\) and natural number \(k\) and natural number \(j\), if \(2 \cdot (t0 + \operatorname{twopow}(k)) < j\), then \(\operatorname{wittbit}(\operatorname{wittwin}(t0, \operatorname{twopow}(k)), j) = 0\).
\end{theorem}
\begin{proof}
The goal follows from \(\operatorname{bitwinwittdeg}(t0, k, j)\).
\end{proof}
\begin{theorem}
\label{ll:paper:right-threshold-exists}
For every natural number \(k\), there exists a natural number \(J\) such that \(k \le J\) and \(J + 2 \le \operatorname{twopow}(\operatorname{succ}(k))\) and \(\operatorname{rightthreshold}(k, J)\).
\end{theorem}
\begin{proof}
The goal follows from \(\operatorname{existsrightthreshold}(k)\).
\end{proof}
\begin{theorem}
\label{ll:paper:winwitt-threshold-gap}
For every natural number \(t0\) and natural number \(k\) and natural number \(J\) and natural number \(j\), if \(\operatorname{rightthreshold}(k, J)\), then if \(j \le J\), then if \(j < \operatorname{twopow}(k)\), then \(\operatorname{wittbit}(\operatorname{wittwin}(t0, \operatorname{twopow}(k)), j) = 0\).
\end{theorem}
\begin{proof}
The goal follows from \(\operatorname{bitwinwittthreshgap}(t0, k, J, j)\).
\end{proof}
\begin{theorem}
\label{ll:paper:winwitt-threshold-lead}
For every natural number \(t0\) and natural number \(k\) and natural number \(J\), if \(\operatorname{rightthreshold}(k, J)\), then if \(J + 1 < \operatorname{twopow}(k)\), then \(\operatorname{wittbit}(\operatorname{wittwin}(t0, \operatorname{twopow}(k)), J + 1) = 1\).
\end{theorem}
\begin{proof}
The goal follows from \(\operatorname{bitwinwittthreshlead}(t0, k, J)\).
\end{proof}
\begin{theorem}
\label{ll:paper:threshold-lt-center-black}
For every natural number \(t\) and natural number \(k\) and natural number \(J\), if \(\operatorname{rightthreshold}(k, J)\), then if \(\operatorname{twopow}(k) \le t\), then if \(\operatorname{center}(t) = \operatorname{true}\), then \(J < t\).
\end{theorem}
\begin{proof}
The goal follows from \(\operatorname{thresholdltcentertrue}(t, k, J)\).
\end{proof}
\begin{theorem}
\label{ll:paper:winwitt-cyclotomic-first}
For every natural number \(t0\) and natural number \(k\), \(\operatorname{vala}(\operatorname{wittwin}(t0, \operatorname{twopow}(k))) = \operatorname{vala}(\operatorname{wittrow}(t0 + \operatorname{twopow}(k)))\).
\end{theorem}
\begin{proof}
The goal follows from \(\operatorname{valawinwitt}(t0, k)\).
\end{proof}
\begin{theorem}
\label{ll:paper:winwitt-cyclotomic-second}
For every natural number \(t0\) and natural number \(k\), \(\operatorname{valb}(\operatorname{wittwin}(t0, \operatorname{twopow}(k))) = \operatorname{valb}(\operatorname{wittrow}(t0 + \operatorname{twopow}(k)))\).
\end{theorem}
\begin{proof}
The goal follows from \(\operatorname{valbwinwitt}(t0, k)\).
\end{proof}
\begin{theorem}
\label{ll:paper:winwitt-parity}
For every natural number \(t0\) and natural number \(k\), \(\operatorname{rowpar}(\operatorname{wittwin}(t0, \operatorname{twopow}(k))) = \operatorname{bxor}(\operatorname{rowpar}(\operatorname{wittrow}(t0 + \operatorname{twopow}(k))), \operatorname{rowpar}(\operatorname{wittrow}(t0)))\).
\end{theorem}
\begin{proof}
The goal follows from \(\operatorname{parwinwitt}(t0, k)\).
\end{proof}
\begin{theorem}
\label{ll:paper:witt-carry-xor}
For every natural number \(a\) and natural number \(b\), \(\operatorname{wittcarry}(\operatorname{wittxor}(a, b)) = \operatorname{wittxor}(\operatorname{wittxor}(\operatorname{wittcarry}(a), \operatorname{wittcarry}(b)), \operatorname{crosscarry}(a, b))\).
\end{theorem}
\begin{proof}
The goal follows from \(\operatorname{carryxor}(a, b)\).
\end{proof}
\begin{theorem}
\label{ll:paper:witt-phi-xor}
For every natural number \(a\) and natural number \(b\), \(\operatorname{wittphi}(\operatorname{wittxor}(a, b)) = \operatorname{wittxor}(\operatorname{wittxor}(\operatorname{wittphi}(a), \operatorname{wittphi}(b)), \operatorname{wittshl}(\operatorname{crosscarry}(a, b), 1))\).
\end{theorem}
\begin{proof}
The goal follows from \(\operatorname{phixor}(a, b)\).
\end{proof}
\begin{theorem}
\label{ll:paper:winwitt-base-step}
For every natural number \(t0\) and natural number \(k\), \(\operatorname{wittwin}(t0 + 1, \operatorname{twopow}(k)) = \operatorname{wittxor}(\operatorname{wittlinear}(\operatorname{wittwin}(t0, \operatorname{twopow}(k))), \operatorname{wittshl}(\operatorname{wittxor}(\operatorname{wittcarry}(\operatorname{wittrow}(t0 + \operatorname{twopow}(k))), \operatorname{wittscale}(k, \operatorname{wittcarry}(\operatorname{wittrow}(t0)))), 1))\).
\end{theorem}
\begin{proof}
The goal follows from \(\operatorname{winwittbasestep}(t0, k)\).
\end{proof}
\begin{theorem}
\label{ll:paper:winwitt-cone-step}
For every natural number \(t0\) and natural number \(k\), if \(2 \cdot t0 + 1 < \operatorname{twopow}(k)\), then \(\operatorname{wittwin}(t0 + 1, \operatorname{twopow}(k)) = \operatorname{wittxor}(\operatorname{wittphi}(\operatorname{wittwin}(t0, \operatorname{twopow}(k))), \operatorname{wittshl}(\operatorname{crosscarry}(\operatorname{wittscale}(k, \operatorname{wittrow}(t0)), \operatorname{wittwin}(t0, \operatorname{twopow}(k))), 1))\).
\end{theorem}
\begin{proof}
The goal follows from \(\operatorname{winwittconestep}(t0, k)\).
\end{proof}
\begin{theorem}
\label{ll:paper:winwitt-triality-outside}
For every natural number \(t0\) and natural number \(k\) and natural number \(j\), if \(2 \cdot t0 < \operatorname{twopow}(k)\), then if \(3 \cdot \operatorname{twopow}(k) \le j\), then \(\operatorname{wittbit}(\operatorname{wittwin}(t0, \operatorname{twopow}(k)), j) = 0\).
\end{theorem}
\begin{proof}
The goal follows from \(\operatorname{bitwinwitttrialityoutside}(t0, k, j)\).
\end{proof}
\begin{theorem}
\label{ll:paper:winwitt-sector-one-cone}
For every natural number \(t0\) and natural number \(k\) and natural number \(r\), if \(2 \cdot t0 < \operatorname{twopow}(k)\), then if \(r < \operatorname{twopow}(k)\), then \(\operatorname{wittbit}(\operatorname{wittwin}(t0, \operatorname{twopow}(k)), r + \operatorname{twopow}(k)) = \operatorname{bton}(\operatorname{bxor}(\operatorname{rdiag}(r + \operatorname{twopow}(k), t0 + \operatorname{twopow}(k)), \operatorname{rdiag}(r, t0)))\).
\end{theorem}
\begin{proof}
The goal follows from \(\operatorname{bitwinwittsectoronecone}(t0, k, r)\).
\end{proof}
\begin{theorem}
\label{ll:paper:winwitt-center-slot}
For every natural number \(t0\) and natural number \(k\), if \(t0 < \operatorname{twopow}(k)\), then \(\operatorname{wittbit}(\operatorname{wittwin}(t0, \operatorname{twopow}(k)), t0 + \operatorname{twopow}(k)) = \operatorname{bton}(\operatorname{bxor}(\operatorname{center}(t0 + \operatorname{twopow}(k)), \operatorname{center}(t0)))\).
\end{theorem}
\begin{proof}
The goal follows from \(\operatorname{bitwinwittcenterslot}(t0, k)\).
\end{proof}
\begin{theorem}
\label{ll:paper:winwitt-sector-one-below-thresh}
For every natural number \(t0\) and natural number \(k\) and natural number \(J\) and natural number \(r\), if \(\operatorname{rightthreshold}(k, J)\), then if \(r \le J\), then if \(r < \operatorname{twopow}(k)\), then if \(2 \cdot t0 < r + \operatorname{twopow}(k)\), then \(\operatorname{wittbit}(\operatorname{wittwin}(t0, \operatorname{twopow}(k)), r + \operatorname{twopow}(k)) = \operatorname{bton}(\operatorname{bxor}(\operatorname{rdiag}(r + \operatorname{twopow}(k), t0 + \operatorname{twopow}(k)), \operatorname{rdiag}(r, t0 + \operatorname{twopow}(k))))\).
\end{theorem}
\begin{proof}
The goal follows from \(\operatorname{bitwinwittsectoronebelowthresh}(t0, k, J, r)\).
\end{proof}
\begin{theorem}
\label{ll:paper:winwitt-sector-one-at-thresh}
For every natural number \(t0\) and natural number \(k\) and natural number \(J\), if \(\operatorname{rightthreshold}(k, J)\), then if \(J + 1 < \operatorname{twopow}(k)\), then if \(2 \cdot t0 < J + 1 + \operatorname{twopow}(k)\), then \(\operatorname{wittbit}(\operatorname{wittwin}(t0, \operatorname{twopow}(k)), J + 1 + \operatorname{twopow}(k)) = \operatorname{bton}(\operatorname{not}(\operatorname{bxor}(\operatorname{rdiag}(J + 1 + \operatorname{twopow}(k), t0 + \operatorname{twopow}(k)), \operatorname{rdiag}(J + 1, t0 + \operatorname{twopow}(k)))))\).
\end{theorem}
\begin{proof}
The goal follows from \(\operatorname{bitwinwittsectoroneatthresh}(t0, k, J)\).
\end{proof}
\begin{theorem}
\label{ll:paper:threshold-center-antiperiod}
For every natural number \(k\) and natural number \(J\), if \(\operatorname{rightthreshold}(k, J)\), then if \(J + 1 < \operatorname{twopow}(k)\), then \(\operatorname{rdiag}(J + 1, J + 1 + \operatorname{twopow}(k)) = \operatorname{not}(\operatorname{center}(J + 1))\).
\end{theorem}
\begin{proof}
The goal follows from \(\operatorname{threshcenterantiperiod}(k, J)\).
\end{proof}
\begin{theorem}
\label{ll:paper:threshold-center-witt-resp}
For every natural number \(k\) and natural number \(J\), if \(\operatorname{rightthreshold}(k, J)\), then if \(J + 1 < \operatorname{twopow}(k)\), then \(\operatorname{wittresp}(J + 1, k) = \operatorname{not}(\operatorname{bxor}(\operatorname{center}(J + 1 + \operatorname{twopow}(k)), \operatorname{rdiag}(J + 1, J + 1 + \operatorname{twopow}(k))))\).
\end{theorem}
\begin{proof}
The goal follows from \(\operatorname{threshcenterwittresp}(k, J)\).
\end{proof}
\begin{theorem}
\label{ll:paper:phiscale-row-cone-step}
For every natural number \(t0\) and natural number \(k\), if \(2 \cdot t0 + 1 < \operatorname{twopow}(k)\), then \(\operatorname{wittscale}(k, \operatorname{wittrow}(t0 + 1)) = \operatorname{wittphi}(\operatorname{wittscale}(k, \operatorname{wittrow}(t0)))\).
\end{theorem}
\begin{proof}
The goal follows from \(\operatorname{phiscalerowconestep}(t0, k)\).
\end{proof}
\begin{theorem}
\label{ll:paper:winwitt-unperturbed-sector-zero}
For every natural number \(t0\) and natural number \(k\) and natural number \(J\) and natural number \(j\), if \(2 \cdot t0 + 1 < \operatorname{twopow}(k)\), then if \(\operatorname{rightthreshold}(k, J)\), then if \(2 \cdot t0 < J\), then if \(j \le \operatorname{twopow}(k)\), then \(\operatorname{wittbit}(\operatorname{wittwin}(t0 + 1, \operatorname{twopow}(k)), j) = \operatorname{wittbit}(\operatorname{wittphi}(\operatorname{wittwin}(t0, \operatorname{twopow}(k))), j)\).
\end{theorem}
\begin{proof}
The goal follows from \(\operatorname{bitwinwittunpertsectzero}(t0, k, J, j)\).
\end{proof}
\begin{theorem}
\label{ll:paper:winwitt-unperturbed-sector-one-gap}
For every natural number \(t0\) and natural number \(k\) and natural number \(r\), if \(2 \cdot t0 + 1 < \operatorname{twopow}(k)\), then if \(2 \cdot t0 + 2 < r\), then if \(r \le \operatorname{twopow}(k)\), then \(\operatorname{wittbit}(\operatorname{wittwin}(t0 + 1, \operatorname{twopow}(k)), r + \operatorname{twopow}(k)) = \operatorname{wittbit}(\operatorname{wittphi}(\operatorname{wittwin}(t0, \operatorname{twopow}(k))), r + \operatorname{twopow}(k))\).
\end{theorem}
\begin{proof}
The goal follows from \(\operatorname{bitwinwittunpertsectonegap}(t0, k, r)\).
\end{proof}
\begin{theorem}
\label{ll:paper:winwitt-sector-one-gap-ldiag}
For every natural number \(t0\) and natural number \(k\) and natural number \(m\), if \(1 \le m\), then if \(2 \cdot t0 + m < \operatorname{twopow}(k)\), then \(\operatorname{wittbit}(\operatorname{wittwin}(t0, \operatorname{twopow}(k)), 2 \cdot t0 + m + \operatorname{twopow}(k)) = \operatorname{bton}(\operatorname{ldiag}(\operatorname{twopow}(k) - m, t0 + \operatorname{twopow}(k)))\).
\end{theorem}
\begin{proof}
The goal follows from \(\operatorname{bitwinwittsectonegapldiag}(t0, k, m)\).
\end{proof}
\begin{theorem}
\label{ll:paper:winwitt-sector-one-cone-edge}
For every natural number \(t0\) and natural number \(k\), if \(2 \cdot t0 < \operatorname{twopow}(k)\), then \(\operatorname{wittbit}(\operatorname{wittwin}(t0, \operatorname{twopow}(k)), 2 \cdot t0 + \operatorname{twopow}(k)) = \operatorname{bton}(\operatorname{not}(\operatorname{ldiag}(\operatorname{twopow}(k), t0 + \operatorname{twopow}(k))))\).
\end{theorem}
\begin{proof}
The goal follows from \(\operatorname{bitwinwittsectoneconeedge}(t0, k)\).
\end{proof}
\begin{theorem}
\label{ll:paper:ldiag-eq-rdiag}
For every natural number \(d\) and natural number \(s\), \(\operatorname{ldiag}(d, d + s) = \operatorname{rdiag}(d + 2 \cdot s, d + s)\).
\end{theorem}
\begin{proof}
The goal follows from \(\operatorname{ldiageqrdiag}(d, s)\).
\end{proof}
\begin{theorem}
\label{ll:paper:rdiag-threshold-sub-ldiag}
For every natural number \(d\) and natural number \(s\) and natural number \(k\) and natural number \(J\), if \(\operatorname{rightthreshold}(k, J)\), then if \(d + 2 \cdot s \le J\), then if \(d + 2 \cdot s < \operatorname{twopow}(k)\), then \(\operatorname{rdiag}(d + 2 \cdot s, d + s + \operatorname{twopow}(k)) = \operatorname{ldiag}(d, d + s)\).
\end{theorem}
\begin{proof}
The goal follows from \(\operatorname{rdiagthreshsubldiag}(d, s, k, J)\).
\end{proof}
\begin{theorem}
\label{ll:paper:rdiag-threshold-succ-not-ldiag}
For every natural number \(d\) and natural number \(s\) and natural number \(k\) and natural number \(J\), if \(\operatorname{rightthreshold}(k, J)\), then if \(d + 2 \cdot s = J + 1\), then if \(J + 1 < \operatorname{twopow}(k)\), then \(\operatorname{rdiag}(J + 1, d + s + \operatorname{twopow}(k)) = \operatorname{not}(\operatorname{ldiag}(d, d + s))\).
\end{theorem}
\begin{proof}
The goal follows from \(\operatorname{rdiagthreshsuccnotldiag}(d, s, k, J)\).
\end{proof}
\begin{theorem}
\label{ll:paper:winwitt-sector-one-ldiag}
For every natural number \(d\) and natural number \(s\) and natural number \(k\), if \(d + 2 \cdot s < \operatorname{twopow}(k)\), then \(\operatorname{wittbit}(\operatorname{wittwin}(d + s, \operatorname{twopow}(k)), d + 2 \cdot s + \operatorname{twopow}(k)) = \operatorname{bton}(\operatorname{bxor}(\operatorname{ldiag}(d + \operatorname{twopow}(k), d + s + \operatorname{twopow}(k)), \operatorname{ldiag}(d, d + s)))\).
\end{theorem}
\begin{proof}
The goal follows from \(\operatorname{bitwinwittsectoneldiag}(d, s, k)\).
\end{proof}
\begin{theorem}
\label{ll:paper:winwitt-top-edge-zero}
For every natural number \(t0\) and natural number \(k\), if \(t0 < \operatorname{twopow}(k)\), then \(\operatorname{wittbit}(\operatorname{wittwin}(t0, \operatorname{twopow}(k)), 2 \cdot t0 + \operatorname{twopow}(\operatorname{succ}(k))) = 0\).
\end{theorem}
\begin{proof}
The goal follows from \(\operatorname{bitwinwitttopedgezero}(t0, k)\).
\end{proof}
\begin{theorem}
\label{ll:paper:winwitt-second-edge-zero}
For every natural number \(m\) and natural number \(k\), if \(1 + m < \operatorname{twopow}(k)\), then \(\operatorname{wittbit}(\operatorname{wittwin}(1 + m, \operatorname{twopow}(k)), 1 + 2 \cdot m + \operatorname{twopow}(\operatorname{succ}(k))) = 0\).
\end{theorem}
\begin{proof}
The goal follows from \(\operatorname{bitwinwittsecondedgezero}(m, k)\).
\end{proof}
\begin{theorem}
\label{ll:paper:rdiag-threshold-late-false-window}
For every natural number \(t0\) and natural number \(k\) and natural number \(J\) and natural number \(j\), if \(\operatorname{rightthreshold}(k, J)\), then if \(2 \cdot t0 < j\), then if \(j \le J\), then \(\operatorname{rdiag}(j, t0 + \operatorname{twopow}(k)) = \operatorname{false}\).
\end{theorem}
\begin{proof}
The goal follows from \(\operatorname{rdiagthreshlatefalsewin}(t0, k, J, j)\).
\end{proof}
\begin{theorem}
\label{ll:paper:rdiag-threshold-late-true-plateau}
For every natural number \(t0\) and natural number \(k\) and natural number \(J\), if \(\operatorname{rightthreshold}(k, J)\), then if \(2 \cdot t0 \le J\), then \(\operatorname{rdiag}(J + 1, t0 + \operatorname{twopow}(k)) = \operatorname{true}\).
\end{theorem}
\begin{proof}
The goal follows from \(\operatorname{rdiagthreshlatetrueplat}(t0, k, J)\).
\end{proof}
\begin{theorem}
\label{ll:paper:rdiag-threshold-late-false-odd-hit}
For every natural number \(t0\) and natural number \(k\) and natural number \(J\), if \(\operatorname{rightthreshold}(k, J)\), then if \(2 \cdot t0 = J + 1\), then \(\operatorname{rdiag}(J + 1, t0 + \operatorname{twopow}(k)) = \operatorname{false}\).
\end{theorem}
\begin{proof}
The goal follows from \(\operatorname{rdiagthreshlatefalseodd}(t0, k, J)\).
\end{proof}
\begin{theorem}
\label{ll:paper:rdiag-threshold-late-false-even-hit}
For every natural number \(t0\) and natural number \(k\) and natural number \(J\), if \(\operatorname{rightthreshold}(k, J)\), then if \(2 \cdot t0 = J\), then \(\operatorname{rdiag}(J + 1, t0 + 1 + \operatorname{twopow}(k)) = \operatorname{false}\).
\end{theorem}
\begin{proof}
The goal follows from \(\operatorname{rdiagthreshlatefalseeven}(t0, k, J)\).
\end{proof}
\begin{theorem}
\label{ll:paper:rdiag-threshold-plus-two-step-neg}
For every natural number \(t0\) and natural number \(k\) and natural number \(J\), if \(\operatorname{rightthreshold}(k, J)\), then if \(2 \cdot t0 \le J\), then \(\operatorname{rdiag}(J + 2, t0 + 1 + \operatorname{twopow}(k)) = \operatorname{not}(\operatorname{rdiag}(J + 2, t0 + \operatorname{twopow}(k)))\).
\end{theorem}
\begin{proof}
The goal follows from \(\operatorname{rdiagthreshplustwostepneg}(t0, k, J)\).
\end{proof}
\begin{theorem}
\label{ll:paper:winwitt-thresh-plus-two-freeze}
For every natural number \(t0\) and natural number \(k\) and natural number \(J\), if \(\operatorname{rightthreshold}(k, J)\), then if \(2 \cdot t0 = J\), then if \(J + 2 < \operatorname{twopow}(k)\), then \(\operatorname{wittbit}(\operatorname{wittwin}(t0 + 1, \operatorname{twopow}(k)), J + 2) = \operatorname{wittbit}(\operatorname{wittwin}(t0, \operatorname{twopow}(k)), J + 2)\).
\end{theorem}
\begin{proof}
The goal follows from \(\operatorname{bitwinwittthreshplustwofreeze}(t0, k, J)\).
\end{proof}
\begin{theorem}
\label{ll:paper:ldiag-two-false}
For every natural number \(m\), \(\operatorname{ldiag}(2, 2 + m) = \operatorname{false}\).
\end{theorem}
\begin{proof}
The goal follows from \(\operatorname{ldiagtwofalse}(m)\).
\end{proof}
\begin{theorem}
\label{ll:paper:winwitt-third-edge-zero}
For every natural number \(m\) and natural number \(k\), if \(2 + m < \operatorname{twopow}(k)\), then \(\operatorname{wittbit}(\operatorname{wittwin}(2 + m, \operatorname{twopow}(k)), 2 + 2 \cdot m + \operatorname{twopow}(\operatorname{succ}(k))) = 0\).
\end{theorem}
\begin{proof}
The goal follows from \(\operatorname{bitwinwittthirdedgezero}(m, k)\).
\end{proof}
\begin{theorem}
\label{ll:paper:ldiag-three-step-neg}
For every natural number \(m\), \(\operatorname{ldiag}(3, 3 + m) = \operatorname{not}(\operatorname{ldiag}(3, 2 + m))\).
\end{theorem}
\begin{proof}
The goal follows from \(\operatorname{ldiagthreestepneg}(m)\).
\end{proof}
\begin{theorem}
\label{ll:paper:winwitt-sector-one-settled}
For every natural number \(d\) and natural number \(s\) and natural number \(k\), if \(d \le k\), then if \(d + 2 \cdot s < \operatorname{twopow}(k)\), then if \(\operatorname{settledby}(d, d + s)\), then \(\operatorname{wittbit}(\operatorname{wittwin}(d + s, \operatorname{twopow}(k)), d + 2 \cdot s + \operatorname{twopow}(k)) = \operatorname{bton}(\operatorname{bxor}(\operatorname{ldiag}(d + \operatorname{twopow}(k), d + s + \operatorname{twopow}(k)), \operatorname{ldiag}(d, d + s + \operatorname{twopow}(k))))\).
\end{theorem}
\begin{proof}
The goal follows from \(\operatorname{bitwinwittsectoronesettled}(d, s, k)\).
\end{proof}
\begin{theorem}
\label{ll:paper:rdiag-eq-ldiag-settled-below-thresh}
For every natural number \(d\) and natural number \(s\) and natural number \(k\) and natural number \(J\), if \(\operatorname{rightthreshold}(k, J)\), then if \(d \le k\), then if \(d + 2 \cdot s \le J\), then if \(d + 2 \cdot s < \operatorname{twopow}(k)\), then if \(\operatorname{settledby}(d, d + s)\), then \(\operatorname{rdiag}(d + 2 \cdot s, d + s + \operatorname{twopow}(k)) = \operatorname{ldiag}(d, d + s + \operatorname{twopow}(k))\).
\end{theorem}
\begin{proof}
The goal follows from \(\operatorname{rdiageqldiagsettledbelowthresh}(d, s, k, J)\).
\end{proof}
\begin{theorem}
\label{ll:paper:rdiag-not-ldiag-settled-at-thresh}
For every natural number \(d\) and natural number \(s\) and natural number \(k\) and natural number \(J\), if \(\operatorname{rightthreshold}(k, J)\), then if \(d \le k\), then if \(d + 2 \cdot s = J + 1\), then if \(J + 1 < \operatorname{twopow}(k)\), then if \(\operatorname{settledby}(d, d + s)\), then \(\operatorname{rdiag}(J + 1, d + s + \operatorname{twopow}(k)) = \operatorname{not}(\operatorname{ldiag}(d, d + s + \operatorname{twopow}(k)))\).
\end{theorem}
\begin{proof}
The goal follows from \(\operatorname{rdiagnotldiagsettledatthresh}(d, s, k, J)\).
\end{proof}
\begin{theorem}
\label{ll:paper:winwitt-sector-two-scale-eq}
For every natural number \(d\) and natural number \(s\) and natural number \(k1\) and natural number \(k2\), if \(d \le k1\), then if \(d \le k2\), then if \(d + s < \operatorname{twopow}(k1)\), then if \(d + s < \operatorname{twopow}(k2)\), then if \(\operatorname{settledby}(d, d + s + \operatorname{twopow}(k1))\), then if \(\operatorname{settledby}(d, d + s + \operatorname{twopow}(k2))\), then \(\operatorname{wittbit}(\operatorname{wittwin}(d + s, \operatorname{twopow}(k1)), d + 2 \cdot s + \operatorname{twopow}(\operatorname{succ}(k1))) = \operatorname{wittbit}(\operatorname{wittwin}(d + s, \operatorname{twopow}(k2)), d + 2 \cdot s + \operatorname{twopow}(\operatorname{succ}(k2)))\).
\end{theorem}
\begin{proof}
The goal follows from \(\operatorname{bitwinwittsectortwoscaleeq}(d, s, k1, k2)\).
\end{proof}
\begin{theorem}
\label{ll:paper:center-late-xor-ldiag-witt}
For every natural number \(d\) and natural number \(k\), if \(d \le k\), then \(\operatorname{bxor}(\operatorname{center}(d + \operatorname{twopow}(k)), \operatorname{ldiag}(d, d + \operatorname{twopow}(k))) = \operatorname{bxor}(\operatorname{wittresp}(d, \operatorname{succ}(k)), \operatorname{wittresp}(d + \operatorname{twopow}(k), k))\).
\end{theorem}
\begin{proof}
The goal follows from \(\operatorname{centerlatexorldiagwitt}(d, k)\).
\end{proof}
\begin{theorem}
\label{ll:paper:periodic-ldiag-xor-late-witt-lag}
For every natural number \(T\) and natural number \(P\), if \(0 < P\), then if \(\operatorname{centerperiodic}(T, P)\), then for every natural number \(d\) and natural number \(k1\) and natural number \(k2\), if \(\operatorname{lag}(P, k1) = \operatorname{lag}(P, k2)\), then if \(T \le d\), then if \(d \le k1\), then if \(d \le k2\), then \(\operatorname{bxor}(\operatorname{ldiag}(d, d + \operatorname{twopow}(k1)), \operatorname{ldiag}(d, d + \operatorname{twopow}(k2))) = \operatorname{bxor}(\operatorname{wittresp}(d + \operatorname{twopow}(k1), k1), \operatorname{wittresp}(d + \operatorname{twopow}(k2), k2))\).
\end{theorem}
\begin{proof}
Assume \(hP\), \(h\).
The goal follows from \(\operatorname{periodicldiagxorlatewittlag}(T, P, hP, h)\).
\end{proof}
\begin{theorem}
\label{ll:paper:periodic-response-shift-p}
For every natural number \(T\) and natural number \(P\), if \(0 < P\), then if \(\operatorname{centerperiodic}(T, P)\), then for every natural number \(t\) and natural number \(k\), if \(T \le t\), then if \(t + P < \operatorname{twopow}(k)\), then \(\operatorname{response}(t + P, k) = \operatorname{response}(t, k)\).
\end{theorem}
\begin{proof}
Assume \(hP\), \(h\).
The goal follows from \(\operatorname{periodicresponseshiftp}(T, P, hP, h)\).
\end{proof}
\begin{theorem}
\label{ll:paper:periodic-witt-shift-p}
For every natural number \(T\) and natural number \(P\), if \(0 < P\), then if \(\operatorname{centerperiodic}(T, P)\), then for every natural number \(t\) and natural number \(k\), if \(T \le t\), then if \(t + P < \operatorname{twopow}(k)\), then \(\operatorname{wittresp}(t + P, k) = \operatorname{wittresp}(t, k)\).
\end{theorem}
\begin{proof}
Assume \(hP\), \(h\).
The goal follows from \(\operatorname{periodicwittshiftp}(T, P, hP, h)\).
\end{proof}
\begin{theorem}
\label{ll:paper:periodic-winwitt-sector-one-shift-p}
For every natural number \(T\) and natural number \(P\), if \(0 < P\), then if \(\operatorname{centerperiodic}(T, P)\), then for every natural number \(d\) and natural number \(k\), if \(T \le d\), then if \(d + P < \operatorname{twopow}(k)\), then \(\operatorname{wittbit}(\operatorname{wittwin}(d + P, \operatorname{twopow}(k)), d + P + \operatorname{twopow}(k)) = \operatorname{wittbit}(\operatorname{wittwin}(d, \operatorname{twopow}(k)), d + \operatorname{twopow}(k))\).
\end{theorem}
\begin{proof}
Assume \(hP\), \(h\).
The goal follows from \(\operatorname{bitperiodicwinwittsectoneshiftp}(T, P, hP, h)\).
\end{proof}
\begin{theorem}
\label{ll:paper:periodic-late-response-unsettled-decomp}
For every natural number \(T\) and natural number \(P\), if \(0 < P\), then if \(\operatorname{centerperiodic}(T, P)\), then for every natural number \(d\) and natural number \(k\), if \(T \le d\), then if \(d \le k\), then \(\operatorname{response}(d + \operatorname{twopow}(k), k) = \operatorname{bxor}(\operatorname{bxor}(\operatorname{center}(d + \operatorname{lag}(P, \operatorname{succ}(k))), \operatorname{center}(d + \operatorname{lag}(P, k))), \operatorname{ldiag}(d, d + \operatorname{twopow}(k)))\).
\end{theorem}
\begin{proof}
Assume \(hP\), \(h\).
The goal follows from \(\operatorname{periodiclaterespunsettled}(T, P, hP, h)\).
\end{proof}
\begin{theorem}
\label{ll:paper:periodic-zcenter-from-lag-late-response}
For every natural number \(T\) and natural number \(P\), if \(0 < P\), then if \(\operatorname{centerperiodic}(T, P)\), then for every natural number \(d\) and natural number \(k\), if \(T \le d\), then if \(d \le k\), then if \(\operatorname{settledat}(d, k)\), then \(\operatorname{zcentereq}(d, \operatorname{bxor}(\operatorname{bxor}(\operatorname{center}(d + \operatorname{lag}(P, \operatorname{succ}(k))), \operatorname{center}(d + \operatorname{lag}(P, k))), \operatorname{response}(d + \operatorname{twopow}(k), k)))\).
\end{theorem}
\begin{proof}
Assume \(hP\), \(h\).
The goal follows from \(\operatorname{periodiczcenterfromlaglateresp}(T, P, hP, h)\).
\end{proof}
\begin{theorem}
\label{ll:paper:periodic-zcenter-shift-xor-late-response}
For every natural number \(T\) and natural number \(P\), if \(\operatorname{centerperiodic}(T, P)\), then for every natural number \(d\) and natural number \(k\), if \(T \le d\), then if \(d + P \le k\), then if \(\operatorname{settledat}(d, k)\), then if \(\operatorname{settledat}(d + P, k)\), then \(\operatorname{zcenterxoreq}(d + P, d, \operatorname{bxor}(\operatorname{response}(d + P + \operatorname{twopow}(k), k), \operatorname{response}(d + \operatorname{twopow}(k), k)))\).
\end{theorem}
\begin{proof}
Assume \(h\).
The goal follows from \(\operatorname{periodiczcentershiftxorlateresp}(T, P, h)\).
\end{proof}
\begin{theorem}
\label{ll:paper:periodic-zcenter-shift-xor-late-witt}
For every natural number \(T\) and natural number \(P\), if \(\operatorname{centerperiodic}(T, P)\), then for every natural number \(d\) and natural number \(k\), if \(T \le d\), then if \(d + P \le k\), then if \(\operatorname{settledat}(d, k)\), then if \(\operatorname{settledat}(d + P, k)\), then \(\operatorname{zcenterxoreq}(d + P, d, \operatorname{bxor}(\operatorname{wittresp}(d + P + \operatorname{twopow}(k), k), \operatorname{wittresp}(d + \operatorname{twopow}(k), k)))\).
\end{theorem}
\begin{proof}
Assume \(h\).
The goal follows from \(\operatorname{periodiczcentershiftxorlatewitt}(T, P, h)\).
\end{proof}
\begin{theorem}
\label{ll:paper:exists-settled-scale-pair}
For every natural number \(d1\) and natural number \(d2\), there exists a natural number \(k\) such that \(d1 \le k\) and \(d2 \le k\) and \(\operatorname{settledat}(d1, k)\) and \(\operatorname{settledat}(d2, k)\).
\end{theorem}
\begin{proof}
The goal follows from \(\operatorname{existssettledscalepair}(d1, d2)\).
\end{proof}
\begin{theorem}
\label{ll:paper:zcenter-not-periodic-one}
For every natural number \(T\), not \(\operatorname{zcenterperiodic}(T, 1)\).
\end{theorem}
\begin{proof}
The goal follows from \(\operatorname{zcenternotperiodicone}(T)\).
\end{proof}
\begin{theorem}
\label{ll:paper:zcenter-not-periodic-two}
For every natural number \(T\), not \(\operatorname{zcenterperiodic}(T, 2)\).
\end{theorem}
\begin{proof}
The goal follows from \(\operatorname{zcenternotperiodictwo}(T)\).
\end{proof}
\begin{theorem}
\label{ll:paper:zcenter-periodic-not-zero-diag-shift}
For every natural number \(T\) and natural number \(Q\), if \(\operatorname{zcenterperiodic}(T, Q)\), then for every natural number \(z\) and natural number \(m\), if \(\operatorname{zerodiagonal}(z)\), then if \(T \le z\), then not \(\operatorname{zerodiagonal}(z + 2 + m \cdot Q)\).
\end{theorem}
\begin{proof}
Assume \(h\).
The goal follows from \(\operatorname{zcenterperiodicnotzerodiagshift}(T, Q, h)\).
\end{proof}
\begin{theorem}
\label{ll:paper:winwitt-scale-cocycle}
For every natural number \(t0\) and natural number \(k\), \(\operatorname{wittwin}(t0 + \operatorname{twopow}(k), \operatorname{twopow}(k)) = \operatorname{wittxor}(\operatorname{wittwin}(t0, \operatorname{twopow}(\operatorname{succ}(k))), \operatorname{wittscale}(k, \operatorname{wittwin}(t0, \operatorname{twopow}(k))))\).
\end{theorem}
\begin{proof}
The goal follows from \(\operatorname{winwittscalecocycle}(t0, k)\).
\end{proof}
\begin{theorem}
\label{ll:paper:witt-resp-late-from-early-triality}
For every natural number \(d\) and natural number \(k\), if \(d \le k\), then \(\operatorname{bitwinwitteq}(d, \operatorname{twopow}(k), d + \operatorname{twopow}(\operatorname{succ}(k)), \operatorname{bxor}(\operatorname{wittresp}(d + \operatorname{twopow}(k), k), \operatorname{bxor}(\operatorname{wittresp}(d, \operatorname{succ}(k)), \operatorname{wittresp}(d, k))))\).
\end{theorem}
\begin{proof}
The goal follows from \(\operatorname{wittresplatefromearlytriality}(d, k)\).
\end{proof}
\begin{theorem}
\label{ll:paper:periodic-late-witt-from-lag-sector-two}
For every natural number \(T\) and natural number \(P\), if \(0 < P\), then if \(\operatorname{centerperiodic}(T, P)\), then for every natural number \(d\) and natural number \(k\), if \(T \le d\), then if \(d \le k\), then \(\operatorname{bitwinwitteq}(d, \operatorname{twopow}(k), d + \operatorname{twopow}(\operatorname{succ}(k)), \operatorname{bxor}(\operatorname{wittresp}(d + \operatorname{twopow}(k), k), \operatorname{bxor}(\operatorname{center}(d + \operatorname{lag}(P, \operatorname{succ}(k))), \operatorname{center}(d + \operatorname{lag}(P, k)))))\).
\end{theorem}
\begin{proof}
Assume \(hP\), \(h\).
The goal follows from \(\operatorname{periodiclatewittfromlagsectortwo}(T, P, hP, h)\).
\end{proof}
\begin{theorem}
\label{ll:paper:periodic-late-maskparity-decomp}
For every natural number \(T\) and natural number \(P\), if \(0 < P\), then if \(\operatorname{centerperiodic}(T, P)\), then for every natural number \(d\) and natural number \(k\), if \(T \le d\), then if \(d \le k\), then \(\operatorname{maskparity}(d + \operatorname{twopow}(k), k) = \operatorname{bxor}(\operatorname{bxor}(\operatorname{leftresponse}(d, \operatorname{succ}(k)), \operatorname{leftresponse}(d, k)), \operatorname{bxor}(\operatorname{center}(d + \operatorname{lag}(P, \operatorname{succ}(k))), \operatorname{center}(d + \operatorname{lag}(P, k))))\).
\end{theorem}
\begin{proof}
Assume \(hP\), \(h\).
The goal follows from \(\operatorname{periodiclatesmaskparitydecomp}(T, P, hP, h)\).
\end{proof}
\begin{theorem}
\label{ll:paper:leftresponse-zero}
For every natural number \(t\), \(\operatorname{leftresponse}(t, 0) = \operatorname{false}\).
\end{theorem}
\begin{proof}
The goal follows from \(\operatorname{leftresponsezero}(t)\).
\end{proof}
\begin{theorem}
\label{ll:paper:response-zero-maskparity}
For every natural number \(t\), \(\operatorname{response}(t, 0) = \operatorname{maskparity}(t, 0)\).
\end{theorem}
\begin{proof}
The goal follows from \(\operatorname{responsezeromaskparity}(t)\).
\end{proof}
\begin{theorem}
\label{ll:paper:maskparity-step-diff}
For every natural number \(t\) and natural number \(k\), \(\operatorname{bxor}(\operatorname{maskparity}(t + 1, k), \operatorname{maskparity}(t, k)) = \operatorname{bxor}(\operatorname{maskparity}(t + \operatorname{twopow}(k), 0), \operatorname{maskparity}(t, 0))\).
\end{theorem}
\begin{proof}
The goal follows from \(\operatorname{maskparitystepdiff}(t, k)\).
\end{proof}
\begin{theorem}
\label{ll:paper:maskparity-succ-zero-rdiag}
For every natural number \(s\), \(\operatorname{maskparity}(s + 1, 0) = \operatorname{not}(\operatorname{rdiag}(s + 1, s + 1)) \operatorname{and} \operatorname{rdiag}(s, s + 1)\).
\end{theorem}
\begin{proof}
The goal follows from \(\operatorname{maskparitysucczerordiag}(s)\).
\end{proof}
\begin{theorem}
\label{ll:paper:maskparity-zero-periodic-of-maskedper}
For every natural number \(T\) and natural number \(P\), if \(\operatorname{centerperiodic}(T, P)\), then if \(\operatorname{maskedperiodic}(T, P)\), then \(\operatorname{maskparityzeroperiodic}(T, P)\).
\end{theorem}
\begin{proof}
Assume \(h\), \(hm\).
The goal follows from \(\operatorname{maskparityzeroperiodicofmaskedper}(T, P, h, hm)\).
\end{proof}
\begin{theorem}
\label{ll:paper:maskedper-of-maskparity-zero-periodic}
For every natural number \(T\) and natural number \(P\), if \(\operatorname{centerperiodic}(T, P)\), then if \(\operatorname{maskparityzeroperiodic}(T, P)\), then \(\operatorname{maskedperiodic}(T, P)\).
\end{theorem}
\begin{proof}
Assume \(h\), \(hmp\).
The goal follows from \(\operatorname{maskedperofmaskparityzeroperiodic}(T, P, h, hmp)\).
\end{proof}
\begin{theorem}
\label{ll:paper:maskparity-periodic-of-maskedper}
For every natural number \(T\) and natural number \(P\), if \(\operatorname{centerperiodic}(T, P)\), then if \(\operatorname{maskedperiodic}(T, P)\), then for every natural number \(t\) and natural number \(k\), if \(T + 1 \le t\), then \(\operatorname{maskparity}(t + P, k) = \operatorname{maskparity}(t, k)\).
\end{theorem}
\begin{proof}
Assume \(h\), \(hm\).
The goal follows from \(\operatorname{maskparityperiodicofmaskedper}(T, P, h, hm)\).
\end{proof}
\begin{theorem}
\label{ll:paper:periodic-maskparity-shift-p-cone}
For every natural number \(T\) and natural number \(P\), if \(0 < P\), then if \(\operatorname{centerperiodic}(T, P)\), then for every natural number \(t\) and natural number \(k\), if \(T \le t\), then if \(t + P < \operatorname{twopow}(k)\), then \(\operatorname{bxor}(\operatorname{maskparity}(t + P, k), \operatorname{maskparity}(t, k)) = \operatorname{bxor}(\operatorname{leftresponse}(t + P, k), \operatorname{leftresponse}(t, k))\).
\end{theorem}
\begin{proof}
Assume \(hP\), \(h\).
The goal follows from \(\operatorname{periodicmaskparityshiftpcone}(T, P, hP, h)\).
\end{proof}
\begin{theorem}
\label{ll:paper:periodic-maskparity-zero-dyadic-jump-cone}
For every natural number \(T\) and natural number \(P\), if \(0 < P\), then if \(\operatorname{centerperiodic}(T, P)\), then for every natural number \(t\) and natural number \(k\), if \(T \le t\), then if \(t + 1 < \operatorname{twopow}(k)\), then \(\operatorname{bxor}(\operatorname{maskparity}(t + \operatorname{twopow}(k), 0), \operatorname{maskparity}(t, 0)) = \operatorname{bxor}(\operatorname{bxor}(\operatorname{leftresponse}(t + 1, k), \operatorname{leftresponse}(t, k)), \operatorname{bxor}(\operatorname{bxor}(\operatorname{center}(t + 1 + \operatorname{lag}(P, k)), \operatorname{center}(t + \operatorname{lag}(P, k))), \operatorname{bxor}(\operatorname{center}(t + 1), \operatorname{center}(t))))\).
\end{theorem}
\begin{proof}
Assume \(hP\), \(h\).
The goal follows from \(\operatorname{periodicmaskparityzerodyadicjumpcone}(T, P, hP, h)\).
\end{proof}
\begin{theorem}
\label{ll:paper:periodic-maskparity-zero-defect-dyadic-jump-cone}
For every natural number \(T\) and natural number \(P\), if \(0 < P\), then if \(\operatorname{centerperiodic}(T, P)\), then for every natural number \(t\) and natural number \(k\), if \(T \le t\), then if \(t + P + 1 < \operatorname{twopow}(k)\), then \(\operatorname{bxor}(\operatorname{bxor}(\operatorname{maskparity}(t + P + \operatorname{twopow}(k), 0), \operatorname{maskparity}(t + \operatorname{twopow}(k), 0)), \operatorname{bxor}(\operatorname{maskparity}(t + P, 0), \operatorname{maskparity}(t, 0))) = \operatorname{bxor}(\operatorname{bxor}(\operatorname{leftresponse}(t + P + 1, k), \operatorname{leftresponse}(t + 1, k)), \operatorname{bxor}(\operatorname{leftresponse}(t + P, k), \operatorname{leftresponse}(t, k)))\).
\end{theorem}
\begin{proof}
Assume \(hP\), \(h\).
The goal follows from \(\operatorname{periodicmaskparityzerodefectdyadicjumpcone}(T, P, hP, h)\).
\end{proof}
\begin{theorem}
\label{ll:paper:periodic-late-maskparity-shift-p}
For every natural number \(T\) and natural number \(P\), if \(0 < P\), then if \(\operatorname{centerperiodic}(T, P)\), then for every natural number \(d\) and natural number \(k\), if \(T \le d\), then if \(d + P \le k\), then \(\operatorname{bxor}(\operatorname{maskparity}(d + P + \operatorname{twopow}(k), k), \operatorname{maskparity}(d + \operatorname{twopow}(k), k)) = \operatorname{bxor}(\operatorname{bxor}(\operatorname{leftresponse}(d + P, \operatorname{succ}(k)), \operatorname{leftresponse}(d + P, k)), \operatorname{bxor}(\operatorname{leftresponse}(d, \operatorname{succ}(k)), \operatorname{leftresponse}(d, k)))\).
\end{theorem}
\begin{proof}
Assume \(hP\), \(h\).
The goal follows from \(\operatorname{periodiclatesmaskparityshiftp}(T, P, hP, h)\).
\end{proof}
\begin{theorem}
\label{ll:paper:periodic-late-maskparity-zero-dyadic-jump}
For every natural number \(T\) and natural number \(P\), if \(0 < P\), then if \(\operatorname{centerperiodic}(T, P)\), then for every natural number \(d\) and natural number \(k\), if \(T \le d\), then if \(d + 1 \le k\), then \(\operatorname{bxor}(\operatorname{maskparity}(d + \operatorname{twopow}(\operatorname{succ}(k)), 0), \operatorname{maskparity}(d + \operatorname{twopow}(k), 0)) = \operatorname{bxor}(\operatorname{bxor}(\operatorname{bxor}(\operatorname{leftresponse}(d + 1, \operatorname{succ}(k)), \operatorname{leftresponse}(d + 1, k)), \operatorname{bxor}(\operatorname{leftresponse}(d, \operatorname{succ}(k)), \operatorname{leftresponse}(d, k))), \operatorname{bxor}(\operatorname{bxor}(\operatorname{center}(d + 1 + \operatorname{lag}(P, \operatorname{succ}(k))), \operatorname{center}(d + 1 + \operatorname{lag}(P, k))), \operatorname{bxor}(\operatorname{center}(d + \operatorname{lag}(P, \operatorname{succ}(k))), \operatorname{center}(d + \operatorname{lag}(P, k)))))\).
\end{theorem}
\begin{proof}
Assume \(hP\), \(h\).
The goal follows from \(\operatorname{periodiclatesmaskparityzerodyadicjump}(T, P, hP, h)\).
\end{proof}
\begin{theorem}
\label{ll:paper:periodic-late-maskparity-zero-defect-dyadic-jump}
For every natural number \(T\) and natural number \(P\), if \(0 < P\), then if \(\operatorname{centerperiodic}(T, P)\), then for every natural number \(d\) and natural number \(k\), if \(T \le d\), then if \(d + P + 1 \le k\), then \(\operatorname{bxor}(\operatorname{bxor}(\operatorname{maskparity}(d + P + \operatorname{twopow}(\operatorname{succ}(k)), 0), \operatorname{maskparity}(d + \operatorname{twopow}(\operatorname{succ}(k)), 0)), \operatorname{bxor}(\operatorname{maskparity}(d + P + \operatorname{twopow}(k), 0), \operatorname{maskparity}(d + \operatorname{twopow}(k), 0))) = \operatorname{bxor}(\operatorname{bxor}(\operatorname{bxor}(\operatorname{leftresponse}(d + P + 1, \operatorname{succ}(k)), \operatorname{leftresponse}(d + 1, \operatorname{succ}(k))), \operatorname{bxor}(\operatorname{leftresponse}(d + P, \operatorname{succ}(k)), \operatorname{leftresponse}(d, \operatorname{succ}(k)))), \operatorname{bxor}(\operatorname{bxor}(\operatorname{leftresponse}(d + P + 1, k), \operatorname{leftresponse}(d + 1, k)), \operatorname{bxor}(\operatorname{leftresponse}(d + P, k), \operatorname{leftresponse}(d, k))))\).
\end{theorem}
\begin{proof}
Assume \(hP\), \(h\).
The goal follows from \(\operatorname{periodiclatesmaskparityzerodefectdyadicjump}(T, P, hP, h)\).
\end{proof}
\begin{theorem}
\label{ll:paper:periodic-late-leftresponse-unsettled-decomp}
For every natural number \(T\) and natural number \(P\), if \(0 < P\), then if \(\operatorname{centerperiodic}(T, P)\), then for every natural number \(d\) and natural number \(k\), if \(T \le d\), then if \(d \le k\), then \(\operatorname{leftresponse}(d + \operatorname{twopow}(k), k) = \operatorname{bxor}(\operatorname{bxor}(\operatorname{center}(d + \operatorname{lag}(P, \operatorname{succ}(k))), \operatorname{center}(d + \operatorname{lag}(P, k))), \operatorname{bxor}(\operatorname{ldiag}(d, d + \operatorname{twopow}(k)), \operatorname{maskparity}(d + \operatorname{twopow}(k), k)))\).
\end{theorem}
\begin{proof}
Assume \(hP\), \(h\).
The goal follows from \(\operatorname{periodiclateleftresponseunsettleddecomp}(T, P, hP, h)\).
\end{proof}
\begin{theorem}
\label{ll:paper:periodic-zcenter-from-lag-late-left-mask}
For every natural number \(T\) and natural number \(P\), if \(0 < P\), then if \(\operatorname{centerperiodic}(T, P)\), then for every natural number \(d\) and natural number \(k\), if \(T \le d\), then if \(d \le k\), then if \(\operatorname{settledat}(d, k)\), then \(\operatorname{zcentereq}(d, \operatorname{bxor}(\operatorname{bxor}(\operatorname{center}(d + \operatorname{lag}(P, \operatorname{succ}(k))), \operatorname{center}(d + \operatorname{lag}(P, k))), \operatorname{bxor}(\operatorname{leftresponse}(d + \operatorname{twopow}(k), k), \operatorname{maskparity}(d + \operatorname{twopow}(k), k))))\).
\end{theorem}
\begin{proof}
Assume \(hP\), \(h\).
The goal follows from \(\operatorname{periodiczcenterfromlaglateleftmask}(T, P, hP, h)\).
\end{proof}
\begin{theorem}
\label{ll:paper:periodic-maskparity-zero-defect-colm1-defect}
For every natural number \(T\) and natural number \(P\), if \(0 < P\), then if \(\operatorname{centerperiodic}(T, P)\), then for every natural number \(t\), if \(T \le t\), then \(\operatorname{bxor}(\operatorname{maskparity}(t + P, 0), \operatorname{maskparity}(t, 0)) = \operatorname{bxor}(\operatorname{cell}(t + P, \operatorname{negpow}(0)), \operatorname{cell}(t, \operatorname{negpow}(0)))\).
\end{theorem}
\begin{proof}
Assume \(hP\), \(h\).
The goal follows from \(\operatorname{periodicmaskparityzerodefectcolm1defect}(T, P, hP, h)\).
\end{proof}
\begin{theorem}
\label{ll:paper:rdiag-late-eq-witt-xor-left-mask}
For every natural number \(r\) and natural number \(k\), \(\operatorname{rdiag}(r, r + \operatorname{twopow}(k)) = \operatorname{bxor}(\operatorname{wittresp}(r + \operatorname{twopow}(k), k), \operatorname{bxor}(\operatorname{leftresponse}(r + \operatorname{twopow}(k), k), \operatorname{maskparity}(r + \operatorname{twopow}(k), k)))\).
\end{theorem}
\begin{proof}
The goal follows from \(\operatorname{rdiaglateeqwittxorleftmask}(r, k)\).
\end{proof}
\begin{theorem}
\label{ll:paper:two-halves-witt-scale}
For every natural number \(r\) and natural number \(k\), \(\operatorname{bxor}(\operatorname{rdiag}(r, r + \operatorname{twopow}(k)), \operatorname{ldiag}(r, r + \operatorname{twopow}(k))) = \operatorname{bxor}(\operatorname{bxor}(\operatorname{center}(r + \operatorname{twopow}(k) + \operatorname{twopow}(k)), \operatorname{center}(r + \operatorname{twopow}(k))), \operatorname{wittresp}(r + \operatorname{twopow}(k), k))\).
\end{theorem}
\begin{proof}
The goal follows from \(\operatorname{twohalveswittscale}(r, k)\).
\end{proof}
\begin{theorem}
\label{ll:paper:ldiag-leftresponse-coboundary-cone}
For every natural number \(r\) and natural number \(k\), if \(r < \operatorname{twopow}(k)\), then \(\operatorname{ldiag}(r, r + \operatorname{twopow}(k)) = \operatorname{bxor}(\operatorname{bxor}(\operatorname{leftresponse}(r, \operatorname{succ}(k)), \operatorname{leftresponse}(r, k)), \operatorname{leftresponse}(r + \operatorname{twopow}(k), k))\).
\end{theorem}
\begin{proof}
Assume \(hr\).
The goal follows from \(\operatorname{ldiagleftresponsecoboundarycone}(r, k, hr)\).
\end{proof}
\begin{theorem}
\label{ll:paper:two-halves-witt-coboundary-cone}
For every natural number \(r\) and natural number \(k\), if \(r < \operatorname{twopow}(k)\), then \(\operatorname{bxor}(\operatorname{rdiag}(r, r + \operatorname{twopow}(k)), \operatorname{ldiag}(r, r + \operatorname{twopow}(k))) = \operatorname{bxor}(\operatorname{bxor}(\operatorname{wittresp}(r, \operatorname{succ}(k)), \operatorname{wittresp}(r, k)), \operatorname{wittresp}(r + \operatorname{twopow}(k), k))\).
\end{theorem}
\begin{proof}
Assume \(hr\).
The goal follows from \(\operatorname{twohalveswittcoboundarycone}(r, k, hr)\).
\end{proof}
\begin{theorem}
\label{ll:paper:periodic-maskparity-late-cone}
For every natural number \(T\) and natural number \(P\), if \(0 < P\), then if \(\operatorname{centerperiodic}(T, P)\), then for every natural number \(r\) and natural number \(k\), if \(T \le r\), then if \(r < \operatorname{twopow}(k)\), then \(\operatorname{maskparity}(r + \operatorname{twopow}(k), k) = \operatorname{bxor}(\operatorname{bxor}(\operatorname{leftresponse}(r, \operatorname{succ}(k)), \operatorname{leftresponse}(r, k)), \operatorname{bxor}(\operatorname{center}(r + \operatorname{lag}(P, \operatorname{succ}(k))), \operatorname{center}(r + \operatorname{lag}(P, k))))\).
\end{theorem}
\begin{proof}
Assume \(hP\), \(h\).
The goal follows from \(\operatorname{periodicmaskparitylatecone}(T, P, hP, h)\).
\end{proof}
\begin{theorem}
\label{ll:paper:periodic-two-halves-scale-eq}
For every natural number \(T\) and natural number \(P\), if \(0 < P\), then if \(\operatorname{centerperiodic}(T, P)\), then for every natural number \(r\) and natural number \(k\), if \(T \le r\), then if \(r < \operatorname{twopow}(k)\), then \(\operatorname{bxor}(\operatorname{rdiag}(r, r + \operatorname{twopow}(k)), \operatorname{ldiag}(r, r + \operatorname{twopow}(k))) = \operatorname{bxor}(\operatorname{bxor}(\operatorname{center}(r + \operatorname{lag}(P, \operatorname{succ}(k))), \operatorname{center}(r + \operatorname{lag}(P, k))), \operatorname{wittresp}(r + \operatorname{twopow}(k), k))\).
\end{theorem}
\begin{proof}
Assume \(hP\), \(h\).
The goal follows from \(\operatorname{periodictwohalvesscaleeq}(T, P, hP, h)\).
\end{proof}
\begin{theorem}
\label{ll:paper:periodic-rdiag-late-from-left-and-center}
For every natural number \(T\) and natural number \(P\), if \(0 < P\), then if \(\operatorname{centerperiodic}(T, P)\), then for every natural number \(r\) and natural number \(k\), if \(T \le r\), then if \(r < \operatorname{twopow}(k)\), then \(\operatorname{bxor}(\operatorname{rdiag}(r, r + \operatorname{twopow}(k)), \operatorname{wittresp}(r + \operatorname{twopow}(k), k)) = \operatorname{bxor}(\operatorname{bxor}(\operatorname{bxor}(\operatorname{leftresponse}(r, \operatorname{succ}(k)), \operatorname{leftresponse}(r, k)), \operatorname{leftresponse}(r + \operatorname{twopow}(k), k)), \operatorname{bxor}(\operatorname{center}(r + \operatorname{lag}(P, \operatorname{succ}(k))), \operatorname{center}(r + \operatorname{lag}(P, k))))\).
\end{theorem}
\begin{proof}
Assume \(hP\), \(h\).
The goal follows from \(\operatorname{periodicrdiaglatefromleftandcenter}(T, P, hP, h)\).
\end{proof}
\begin{theorem}
\label{ll:paper:rdiag-below-threshold-eq-center}
For every natural number \(r\) and natural number \(k\) and natural number \(J\), if \(\operatorname{rightthreshold}(k, J)\), then if \(r \le J\), then if \(r < \operatorname{twopow}(k)\), then \(\operatorname{rdiag}(r, r + \operatorname{twopow}(k)) = \operatorname{center}(r)\).
\end{theorem}
\begin{proof}
The goal follows from \(\operatorname{rdiagbelowthresholdeqcenter}(r, k, J)\).
\end{proof}
\begin{theorem}
\label{ll:paper:periodic-two-halves-below-threshold}
For every natural number \(T\) and natural number \(P\), if \(0 < P\), then if \(\operatorname{centerperiodic}(T, P)\), then for every natural number \(r\) and natural number \(k\) and natural number \(J\), if \(\operatorname{rightthreshold}(k, J)\), then if \(T \le r\), then if \(r \le J\), then if \(r < \operatorname{twopow}(k)\), then \(\operatorname{bxor}(\operatorname{center}(r), \operatorname{ldiag}(r, r + \operatorname{twopow}(k))) = \operatorname{bxor}(\operatorname{bxor}(\operatorname{center}(r + \operatorname{lag}(P, \operatorname{succ}(k))), \operatorname{center}(r + \operatorname{lag}(P, k))), \operatorname{wittresp}(r + \operatorname{twopow}(k), k))\).
\end{theorem}
\begin{proof}
Assume \(hP\), \(h\).
The goal follows from \(\operatorname{periodictwohalvesbelowthreshold}(T, P, hP, h)\).
\end{proof}
\begin{theorem}
\label{ll:paper:periodic-center-xor-witt-below-threshold}
For every natural number \(T\) and natural number \(P\), if \(0 < P\), then if \(\operatorname{centerperiodic}(T, P)\), then for every natural number \(r\) and natural number \(k\) and natural number \(J\), if \(\operatorname{rightthreshold}(k, J)\), then if \(T \le r\), then if \(r \le J\), then if \(r < \operatorname{twopow}(k)\), then \(\operatorname{bxor}(\operatorname{center}(r), \operatorname{wittresp}(r + \operatorname{twopow}(k), k)) = \operatorname{bxor}(\operatorname{bxor}(\operatorname{bxor}(\operatorname{leftresponse}(r, \operatorname{succ}(k)), \operatorname{leftresponse}(r, k)), \operatorname{leftresponse}(r + \operatorname{twopow}(k), k)), \operatorname{bxor}(\operatorname{center}(r + \operatorname{lag}(P, \operatorname{succ}(k))), \operatorname{center}(r + \operatorname{lag}(P, k))))\).
\end{theorem}
\begin{proof}
Assume \(hP\), \(h\).
The goal follows from \(\operatorname{periodiccenterxorwittbelowthreshold}(T, P, hP, h)\).
\end{proof}
\begin{theorem}
\label{ll:paper:periodic-two-halves-at-threshold}
For every natural number \(T\) and natural number \(P\), if \(0 < P\), then if \(\operatorname{centerperiodic}(T, P)\), then for every natural number \(k\) and natural number \(J\), if \(\operatorname{rightthreshold}(k, J)\), then if \(T \le J + 1\), then if \(J + 1 < \operatorname{twopow}(k)\), then \(\operatorname{bxor}(\operatorname{not}(\operatorname{center}(J + 1)), \operatorname{ldiag}(J + 1, J + 1 + \operatorname{twopow}(k))) = \operatorname{bxor}(\operatorname{bxor}(\operatorname{center}(J + 1 + \operatorname{lag}(P, \operatorname{succ}(k))), \operatorname{center}(J + 1 + \operatorname{lag}(P, k))), \operatorname{wittresp}(J + 1 + \operatorname{twopow}(k), k))\).
\end{theorem}
\begin{proof}
Assume \(hP\), \(h\).
The goal follows from \(\operatorname{periodictwohalvesatthreshold}(T, P, hP, h)\).
\end{proof}
\begin{theorem}
\label{ll:paper:periodic-not-center-xor-witt-at-threshold}
For every natural number \(T\) and natural number \(P\), if \(0 < P\), then if \(\operatorname{centerperiodic}(T, P)\), then for every natural number \(k\) and natural number \(J\), if \(\operatorname{rightthreshold}(k, J)\), then if \(T \le J + 1\), then if \(J + 1 < \operatorname{twopow}(k)\), then \(\operatorname{bxor}(\operatorname{not}(\operatorname{center}(J + 1)), \operatorname{wittresp}(J + 1 + \operatorname{twopow}(k), k)) = \operatorname{bxor}(\operatorname{bxor}(\operatorname{bxor}(\operatorname{leftresponse}(J + 1, \operatorname{succ}(k)), \operatorname{leftresponse}(J + 1, k)), \operatorname{leftresponse}(J + 1 + \operatorname{twopow}(k), k)), \operatorname{bxor}(\operatorname{center}(J + 1 + \operatorname{lag}(P, \operatorname{succ}(k))), \operatorname{center}(J + 1 + \operatorname{lag}(P, k))))\).
\end{theorem}
\begin{proof}
Assume \(hP\), \(h\).
The goal follows from \(\operatorname{periodicnotcenterxorwittatthreshold}(T, P, hP, h)\).
\end{proof}
\begin{theorem}
\label{ll:paper:threshold-center-left-mask-eq-not-rdiag}
For every natural number \(k\) and natural number \(J\), if \(\operatorname{rightthreshold}(k, J)\), then if \(J + 1 < \operatorname{twopow}(k)\), then \(\operatorname{bxor}(\operatorname{leftresponse}(J + 1, k), \operatorname{maskparity}(J + 1, k)) = \operatorname{not}(\operatorname{bxor}(\operatorname{center}(J + 1 + \operatorname{twopow}(k)), \operatorname{rdiag}(J + 1, J + 1 + \operatorname{twopow}(k))))\).
\end{theorem}
\begin{proof}
The goal follows from \(\operatorname{threshcenterleftmaskeqnotrdiag}(k, J)\).
\end{proof}
\begin{theorem}
\label{ll:paper:sector-one-below-threshold-eq-left-settled}
For every natural number \(d\) and natural number \(s\) and natural number \(k\) and natural number \(J\), if \(\operatorname{rightthreshold}(k, J)\), then if \(d \le k\), then if \(d + 2 \cdot s \le J\), then if \(d + 2 \cdot s < \operatorname{twopow}(k)\), then if \(\operatorname{settledby}(d, d + s)\), then \(\operatorname{bxor}(\operatorname{rdiag}(d + 2 \cdot s + \operatorname{twopow}(k), d + s + \operatorname{twopow}(k)), \operatorname{rdiag}(d + 2 \cdot s, d + s + \operatorname{twopow}(k))) = \operatorname{bxor}(\operatorname{ldiag}(d + \operatorname{twopow}(k), d + s + \operatorname{twopow}(k)), \operatorname{ldiag}(d, d + s + \operatorname{twopow}(k)))\).
\end{theorem}
\begin{proof}
The goal follows from \(\operatorname{sectoronebelowthresholdeqleftsettled}(d, s, k, J)\).
\end{proof}
\begin{theorem}
\label{ll:paper:sector-one-at-threshold-eq-left-settled}
For every natural number \(d\) and natural number \(s\) and natural number \(k\) and natural number \(J\), if \(\operatorname{rightthreshold}(k, J)\), then if \(d \le k\), then if \(d + 2 \cdot s = J + 1\), then if \(J + 1 < \operatorname{twopow}(k)\), then if \(\operatorname{settledby}(d, d + s)\), then \(\operatorname{not}(\operatorname{bxor}(\operatorname{rdiag}(J + 1 + \operatorname{twopow}(k), d + s + \operatorname{twopow}(k)), \operatorname{rdiag}(J + 1, d + s + \operatorname{twopow}(k)))) = \operatorname{bxor}(\operatorname{ldiag}(d + \operatorname{twopow}(k), d + s + \operatorname{twopow}(k)), \operatorname{ldiag}(d, d + s + \operatorname{twopow}(k)))\).
\end{theorem}
\begin{proof}
The goal follows from \(\operatorname{sectoroneatthresholdeqleftsettled}(d, s, k, J)\).
\end{proof}
\begin{theorem}
\label{ll:paper:rdiag-defect-step-below-threshold}
For every natural number \(t\) and natural number \(k\) and natural number \(J\) and natural number \(j\), if \(\operatorname{rightthreshold}(k, J)\), then if \(j \le J\), then \(\operatorname{bxor}(\operatorname{bxor}(\operatorname{rdiag}(j + 2, t + 1 + \operatorname{twopow}(k)), \operatorname{rdiag}(j + 2, t + 1)), \operatorname{bxor}(\operatorname{rdiag}(j + 2, t + \operatorname{twopow}(k)), \operatorname{rdiag}(j + 2, t))) = \operatorname{bxor}(\operatorname{rdiag}(j + 1, t + \operatorname{twopow}(k)), \operatorname{rdiag}(j + 1, t)) \operatorname{and} \operatorname{not}(\operatorname{rdiag}(j, t))\).
\end{theorem}
\begin{proof}
The goal follows from \(\operatorname{rdiagdefectstepbelowthreshold}(t, k, J, j)\).
\end{proof}
\begin{theorem}
\label{ll:paper:rdiag-threshold-plus-two-defect-step}
For every natural number \(t\) and natural number \(k\) and natural number \(J\), if \(\operatorname{rightthreshold}(k, J)\), then \(\operatorname{bxor}(\operatorname{bxor}(\operatorname{rdiag}(J + 2, t + 1 + \operatorname{twopow}(k)), \operatorname{rdiag}(J + 2, t + 1)), \operatorname{bxor}(\operatorname{rdiag}(J + 2, t + \operatorname{twopow}(k)), \operatorname{rdiag}(J + 2, t))) = \operatorname{not}(\operatorname{rdiag}(J, t))\).
\end{theorem}
\begin{proof}
The goal follows from \(\operatorname{rdiagthresholdplustwodefectstep}(t, k, J)\).
\end{proof}
\begin{theorem}
\label{ll:paper:rdiag-threshold-plus-two-defect-step-ldiag}
For every natural number \(d\) and natural number \(s\) and natural number \(k\) and natural number \(J\), if \(\operatorname{rightthreshold}(k, J)\), then if \(d + 2 \cdot s = J\), then \(\operatorname{bxor}(\operatorname{bxor}(\operatorname{rdiag}(J + 2, d + s + 1 + \operatorname{twopow}(k)), \operatorname{rdiag}(J + 2, d + s + 1)), \operatorname{bxor}(\operatorname{rdiag}(J + 2, d + s + \operatorname{twopow}(k)), \operatorname{rdiag}(J + 2, d + s))) = \operatorname{not}(\operatorname{ldiag}(d, d + s))\).
\end{theorem}
\begin{proof}
The goal follows from \(\operatorname{rdiagthresholdplustwodefectstepldiag}(d, s, k, J)\).
\end{proof}
\begin{theorem}
\label{ll:paper:rdiag-threshold-plus-two-center-decomp}
For every natural number \(k\) and natural number \(J\), if \(\operatorname{rightthreshold}(k, J)\), then \(\operatorname{rdiag}(J + 2, J + 2 + \operatorname{twopow}(k)) = \operatorname{bxor}(\operatorname{bxor}(\operatorname{center}(J + 2), \operatorname{bxor}(\operatorname{rdiag}(J + 2, J + \operatorname{twopow}(k)), \operatorname{rdiag}(J + 2, J))), \operatorname{bxor}(\operatorname{center}(J), \operatorname{rdiag}(J, J + 1)))\).
\end{theorem}
\begin{proof}
The goal follows from \(\operatorname{rdiagthresholdplustwocenterdecomp}(k, J)\).
\end{proof}
\begin{theorem}
\label{ll:paper:periodic-two-halves-at-threshold-plus-two}
For every natural number \(T\) and natural number \(P\), if \(0 < P\), then if \(\operatorname{centerperiodic}(T, P)\), then for every natural number \(k\) and natural number \(J\), if \(\operatorname{rightthreshold}(k, J)\), then if \(T \le J + 2\), then if \(J + 2 < \operatorname{twopow}(k)\), then \(\operatorname{bxor}(\operatorname{bxor}(\operatorname{bxor}(\operatorname{center}(J + 2), \operatorname{bxor}(\operatorname{rdiag}(J + 2, J + \operatorname{twopow}(k)), \operatorname{rdiag}(J + 2, J))), \operatorname{bxor}(\operatorname{center}(J), \operatorname{rdiag}(J, J + 1))), \operatorname{ldiag}(J + 2, J + 2 + \operatorname{twopow}(k))) = \operatorname{bxor}(\operatorname{bxor}(\operatorname{center}(J + 2 + \operatorname{lag}(P, \operatorname{succ}(k))), \operatorname{center}(J + 2 + \operatorname{lag}(P, k))), \operatorname{wittresp}(J + 2 + \operatorname{twopow}(k), k))\).
\end{theorem}
\begin{proof}
Assume \(hP\), \(h\).
The goal follows from \(\operatorname{periodictwohalvesatthresholdplustwo}(T, P, hP, h)\).
\end{proof}
\begin{theorem}
\label{ll:paper:periodic-center-xor-witt-at-threshold-plus-two}
For every natural number \(T\) and natural number \(P\), if \(0 < P\), then if \(\operatorname{centerperiodic}(T, P)\), then for every natural number \(k\) and natural number \(J\), if \(\operatorname{rightthreshold}(k, J)\), then if \(T \le J + 2\), then if \(J + 2 < \operatorname{twopow}(k)\), then \(\operatorname{bxor}(\operatorname{bxor}(\operatorname{bxor}(\operatorname{center}(J + 2), \operatorname{bxor}(\operatorname{rdiag}(J + 2, J + \operatorname{twopow}(k)), \operatorname{rdiag}(J + 2, J))), \operatorname{bxor}(\operatorname{center}(J), \operatorname{rdiag}(J, J + 1))), \operatorname{wittresp}(J + 2 + \operatorname{twopow}(k), k)) = \operatorname{bxor}(\operatorname{bxor}(\operatorname{bxor}(\operatorname{leftresponse}(J + 2, \operatorname{succ}(k)), \operatorname{leftresponse}(J + 2, k)), \operatorname{leftresponse}(J + 2 + \operatorname{twopow}(k), k)), \operatorname{bxor}(\operatorname{center}(J + 2 + \operatorname{lag}(P, \operatorname{succ}(k))), \operatorname{center}(J + 2 + \operatorname{lag}(P, k))))\).
\end{theorem}
\begin{proof}
Assume \(hP\), \(h\).
The goal follows from \(\operatorname{periodiccenterxorwittatthresholdplustwo}(T, P, hP, h)\).
\end{proof}
\begin{theorem}
\label{ll:paper:ldiag-damage-step-of-settled}
For every natural number \(d\) and natural number \(s\) and natural number \(k\), if \(d \le k\), then if \(\operatorname{settledby}(d, s)\), then \(\operatorname{bxor}(\operatorname{ldiag}(d + 2, s + 1 + \operatorname{twopow}(k)), \operatorname{ldiag}(d + 2, s + 1)) = \operatorname{bxor}(\operatorname{ldiag}(d + 1, s + \operatorname{twopow}(k)) \operatorname{or} \operatorname{ldiag}(d + 2, s + \operatorname{twopow}(k)), \operatorname{ldiag}(d + 1, s) \operatorname{or} \operatorname{ldiag}(d + 2, s))\).
\end{theorem}
\begin{proof}
The goal follows from \(\operatorname{ldiagdamagestepofsettled}(d, s, k)\).
\end{proof}
\begin{theorem}
\label{ll:paper:ldiag-damage-step-of-settled-pair}
For every natural number \(d\) and natural number \(s\) and natural number \(k\), if \(d + 1 \le k\), then if \(\operatorname{settledby}(d, s)\), then if \(\operatorname{settledby}(d + 1, s)\), then \(\operatorname{bxor}(\operatorname{ldiag}(d + 2, s + 1 + \operatorname{twopow}(k)), \operatorname{ldiag}(d + 2, s + 1)) = \operatorname{bxor}(\operatorname{ldiag}(d + 2, s + \operatorname{twopow}(k)), \operatorname{ldiag}(d + 2, s)) \operatorname{and} \operatorname{not}(\operatorname{ldiag}(d + 1, s))\).
\end{theorem}
\begin{proof}
The goal follows from \(\operatorname{ldiagdamagestepofsettledpair}(d, s, k)\).
\end{proof}
\begin{theorem}
\label{ll:paper:two-halves-off-center-witt-scale}
For every natural number \(d\) and natural number \(s\) and natural number \(k\), if \(d + 2 \cdot s < \operatorname{twopow}(k)\), then \(\operatorname{bitwinwitteq}(d + s + \operatorname{twopow}(k), \operatorname{twopow}(k), d + 2 \cdot s + \operatorname{twopow}(\operatorname{succ}(k)), \operatorname{bxor}(\operatorname{bxor}(\operatorname{rdiag}(d + 2 \cdot s, d + s + \operatorname{twopow}(k)), \operatorname{ldiag}(d, d + s + \operatorname{twopow}(k))), \operatorname{bxor}(\operatorname{ldiag}(d + \operatorname{twopow}(\operatorname{succ}(k)), d + s + \operatorname{twopow}(\operatorname{succ}(k))), \operatorname{ldiag}(d + \operatorname{twopow}(k), d + s + \operatorname{twopow}(k)))))\).
\end{theorem}
\begin{proof}
The goal follows from \(\operatorname{twohalvesoffcenterwittscale}(d, s, k)\).
\end{proof}
\begin{theorem}
\label{ll:paper:off-center-witt-scale-settled-below-threshold}
For every natural number \(d\) and natural number \(s\) and natural number \(k\) and natural number \(J\), if \(\operatorname{rightthreshold}(k, J)\), then if \(d \le k\), then if \(d + 2 \cdot s \le J\), then if \(d + 2 \cdot s < \operatorname{twopow}(k)\), then if \(\operatorname{settledby}(d, d + s)\), then \(\operatorname{bitwinwitteq}(d + s + \operatorname{twopow}(k), \operatorname{twopow}(k), d + 2 \cdot s + \operatorname{twopow}(\operatorname{succ}(k)), \operatorname{bxor}(\operatorname{ldiag}(d + \operatorname{twopow}(\operatorname{succ}(k)), d + s + \operatorname{twopow}(\operatorname{succ}(k))), \operatorname{ldiag}(d + \operatorname{twopow}(k), d + s + \operatorname{twopow}(k))))\).
\end{theorem}
\begin{proof}
The goal follows from \(\operatorname{offcenterwittscalesettledbelowthreshold}(d, s, k, J)\).
\end{proof}
\begin{theorem}
\label{ll:paper:off-center-witt-scale-settled-at-threshold}
For every natural number \(d\) and natural number \(s\) and natural number \(k\) and natural number \(J\), if \(\operatorname{rightthreshold}(k, J)\), then if \(d \le k\), then if \(d + 2 \cdot s = J + 1\), then if \(J + 1 < \operatorname{twopow}(k)\), then if \(\operatorname{settledby}(d, d + s)\), then \(\operatorname{bitwinwitteq}(d + s + \operatorname{twopow}(k), \operatorname{twopow}(k), J + 1 + \operatorname{twopow}(\operatorname{succ}(k)), \operatorname{not}(\operatorname{bxor}(\operatorname{ldiag}(d + \operatorname{twopow}(\operatorname{succ}(k)), d + s + \operatorname{twopow}(\operatorname{succ}(k))), \operatorname{ldiag}(d + \operatorname{twopow}(k), d + s + \operatorname{twopow}(k)))))\).
\end{theorem}
\begin{proof}
The goal follows from \(\operatorname{offcenterwittscalesettledatthreshold}(d, s, k, J)\).
\end{proof}
\begin{theorem}
\label{ll:paper:colm1-dyadic-jump-eq-maskparity-zero-jump}
For every natural number \(d\) and natural number \(k\), \(\operatorname{bxor}(\operatorname{ldiag}(d + \operatorname{twopow}(\operatorname{succ}(k)), d + 1 + \operatorname{twopow}(\operatorname{succ}(k))), \operatorname{ldiag}(d + \operatorname{twopow}(k), d + 1 + \operatorname{twopow}(k))) = \operatorname{bxor}(\operatorname{bxor}(\operatorname{maskparity}(d + 1 + \operatorname{twopow}(\operatorname{succ}(k)), 0), \operatorname{maskparity}(d + 1 + \operatorname{twopow}(k), 0)), \operatorname{bxor}(\operatorname{bxor}(\operatorname{center}(d + 1 + \operatorname{twopow}(\operatorname{succ}(k)) + 1), \operatorname{center}(d + 1 + \operatorname{twopow}(\operatorname{succ}(k)))), \operatorname{bxor}(\operatorname{center}(d + 1 + \operatorname{twopow}(k) + 1), \operatorname{center}(d + 1 + \operatorname{twopow}(k)))))\).
\end{theorem}
\begin{proof}
The goal follows from \(\operatorname{colm1dyadicjumpeqmaskparityzerojump}(d, k)\).
\end{proof}
\begin{theorem}
\label{ll:paper:periodic-colm1-dyadic-jump-eq-leftresponse-diff}
For every natural number \(T\) and natural number \(P\), if \(0 < P\), then if \(\operatorname{centerperiodic}(T, P)\), then for every natural number \(d\) and natural number \(k\), if \(T \le d + 1\), then if \(d + 2 \le k\), then \(\operatorname{bxor}(\operatorname{ldiag}(d + \operatorname{twopow}(\operatorname{succ}(k)), d + 1 + \operatorname{twopow}(\operatorname{succ}(k))), \operatorname{ldiag}(d + \operatorname{twopow}(k), d + 1 + \operatorname{twopow}(k))) = \operatorname{bxor}(\operatorname{bxor}(\operatorname{leftresponse}(d + 2, \operatorname{succ}(k)), \operatorname{leftresponse}(d + 2, k)), \operatorname{bxor}(\operatorname{leftresponse}(d + 1, \operatorname{succ}(k)), \operatorname{leftresponse}(d + 1, k)))\).
\end{theorem}
\begin{proof}
Assume \(hP\), \(h\).
The goal follows from \(\operatorname{periodiccolm1dyadicjumpeqleftresponsediff}(T, P, hP, h)\).
\end{proof}
\begin{theorem}
\label{ll:paper:periodic-off-center-witt-below-threshold-eq-leftresponse}
For every natural number \(T\) and natural number \(P\), if \(0 < P\), then if \(\operatorname{centerperiodic}(T, P)\), then for every natural number \(d\) and natural number \(k\) and natural number \(J\), if \(\operatorname{rightthreshold}(k, J)\), then if \(T \le d + 1\), then if \(d + 2 \le k\), then if \(d + 2 \le J\), then if \(\operatorname{settledby}(d, d + 1)\), then \(\operatorname{bitwinwitteq}(d + 1 + \operatorname{twopow}(k), \operatorname{twopow}(k), d + 2 + \operatorname{twopow}(\operatorname{succ}(k)), \operatorname{bxor}(\operatorname{bxor}(\operatorname{leftresponse}(d + 2, \operatorname{succ}(k)), \operatorname{leftresponse}(d + 2, k)), \operatorname{bxor}(\operatorname{leftresponse}(d + 1, \operatorname{succ}(k)), \operatorname{leftresponse}(d + 1, k))))\).
\end{theorem}
\begin{proof}
Assume \(hP\), \(h\).
The goal follows from \(\operatorname{periodicoffcenterwittbelowthresholdeqleftresponse}(T, P, hP, h)\).
\end{proof}
\begin{theorem}
\label{ll:paper:periodic-off-center-witt-at-threshold-eq-not-leftresponse}
For every natural number \(T\) and natural number \(P\), if \(0 < P\), then if \(\operatorname{centerperiodic}(T, P)\), then for every natural number \(d\) and natural number \(k\) and natural number \(J\), if \(\operatorname{rightthreshold}(k, J)\), then if \(T \le d + 1\), then if \(d + 2 \le k\), then if \(d + 2 = J + 1\), then if \(J + 1 < \operatorname{twopow}(k)\), then if \(\operatorname{settledby}(d, d + 1)\), then \(\operatorname{bitwinwitteq}(d + 1 + \operatorname{twopow}(k), \operatorname{twopow}(k), J + 1 + \operatorname{twopow}(\operatorname{succ}(k)), \operatorname{not}(\operatorname{bxor}(\operatorname{bxor}(\operatorname{leftresponse}(d + 2, \operatorname{succ}(k)), \operatorname{leftresponse}(d + 2, k)), \operatorname{bxor}(\operatorname{leftresponse}(d + 1, \operatorname{succ}(k)), \operatorname{leftresponse}(d + 1, k)))))\).
\end{theorem}
\begin{proof}
Assume \(hP\), \(h\).
The goal follows from \(\operatorname{periodicoffcenterwittatthresholdeqnotleftresponse}(T, P, hP, h)\).
\end{proof}
\begin{theorem}
\label{ll:paper:rdiag-threshold-plus-two-defect-cone}
For every natural number \(t0\) and natural number \(k\) and natural number \(J\), if \(2 \cdot t0 \le J + 1\), then \(\operatorname{bxor}(\operatorname{rdiag}(J + 2, t0 + \operatorname{twopow}(k)), \operatorname{rdiag}(J + 2, t0)) = \operatorname{rdiag}(J + 2, t0 + \operatorname{twopow}(k))\).
\end{theorem}
\begin{proof}
The goal follows from \(\operatorname{rdiagthresholdplustwodefectcone}(t0, k, J)\).
\end{proof}
\begin{theorem}
\label{ll:paper:rdiag-threshold-plus-two-defect-two-step}
For every natural number \(t\) and natural number \(k\) and natural number \(J\), if \(\operatorname{rightthreshold}(k, J)\), then \(\operatorname{bxor}(\operatorname{bxor}(\operatorname{rdiag}(J + 2, t + 2 + \operatorname{twopow}(k)), \operatorname{rdiag}(J + 2, t + 2)), \operatorname{bxor}(\operatorname{rdiag}(J + 2, t + \operatorname{twopow}(k)), \operatorname{rdiag}(J + 2, t))) = \operatorname{bxor}(\operatorname{rdiag}(J, t), \operatorname{rdiag}(J, t + 1))\).
\end{theorem}
\begin{proof}
The goal follows from \(\operatorname{rdiagthresholdplustwodefecttwostep}(t, k, J)\).
\end{proof}
\begin{theorem}
\label{ll:paper:carry-row-scale-decomp}
For every natural number \(t0\) and natural number \(k\), \(\operatorname{wittcarry}(\operatorname{wittrow}(t0 + \operatorname{twopow}(k))) = \operatorname{wittxor}(\operatorname{wittxor}(\operatorname{wittcarry}(\operatorname{wittscale}(k, \operatorname{wittrow}(t0))), \operatorname{wittcarry}(\operatorname{wittwin}(t0, \operatorname{twopow}(k)))), \operatorname{crosscarry}(\operatorname{wittscale}(k, \operatorname{wittrow}(t0)), \operatorname{wittwin}(t0, \operatorname{twopow}(k))))\).
\end{theorem}
\begin{proof}
The goal follows from \(\operatorname{carryrowscaledecomp}(t0, k)\).
\end{proof}
\begin{theorem}
\label{ll:paper:carry-row-scale-cone}
For every natural number \(t0\) and natural number \(k\), if \(2 \cdot t0 + 1 < \operatorname{twopow}(k)\), then \(\operatorname{wittcarry}(\operatorname{wittrow}(t0 + \operatorname{twopow}(k))) = \operatorname{wittxor}(\operatorname{wittxor}(\operatorname{wittscale}(k, \operatorname{wittcarry}(\operatorname{wittrow}(t0))), \operatorname{wittcarry}(\operatorname{wittwin}(t0, \operatorname{twopow}(k)))), \operatorname{crosscarry}(\operatorname{wittscale}(k, \operatorname{wittrow}(t0)), \operatorname{wittwin}(t0, \operatorname{twopow}(k))))\).
\end{theorem}
\begin{proof}
The goal follows from \(\operatorname{carryrowscalecone}(t0, k)\).
\end{proof}
\begin{theorem}
\label{ll:paper:carry-winwitt-late-decomp}
For every natural number \(t0\) and natural number \(k\), \(\operatorname{wittcarry}(\operatorname{wittwin}(t0 + \operatorname{twopow}(k), \operatorname{twopow}(k))) = \operatorname{wittxor}(\operatorname{wittxor}(\operatorname{wittcarry}(\operatorname{wittwin}(t0, \operatorname{twopow}(\operatorname{succ}(k)))), \operatorname{wittcarry}(\operatorname{wittscale}(k, \operatorname{wittwin}(t0, \operatorname{twopow}(k))))), \operatorname{crosscarry}(\operatorname{wittwin}(t0, \operatorname{twopow}(\operatorname{succ}(k))), \operatorname{wittscale}(k, \operatorname{wittwin}(t0, \operatorname{twopow}(k)))))\).
\end{theorem}
\begin{proof}
The goal follows from \(\operatorname{carrywinwittlatedecomp}(t0, k)\).
\end{proof}
\begin{theorem}
\label{ll:paper:periodic-crosscarry-center-slot-and}
For every natural number \(T\) and natural number \(P\), if \(0 < P\), then if \(\operatorname{centerperiodic}(T, P)\), then for every natural number \(t0\) and natural number \(k\), if \(T \le t0\), then if \(t0 + 1 < \operatorname{twopow}(k)\), then \(\operatorname{bitcrosscarryeq}(\operatorname{wittscale}(k, \operatorname{wittrow}(t0)), \operatorname{wittwin}(t0, \operatorname{twopow}(k)), t0 + 1 + \operatorname{twopow}(k), \operatorname{bxor}(\operatorname{rdiag}(t0 + 1, t0) \operatorname{and} \operatorname{bxor}(\operatorname{center}(t0 + \operatorname{lag}(P, k)), \operatorname{center}(t0)), \operatorname{bxor}(\operatorname{rdiag}(t0 + 1 + \operatorname{twopow}(k), t0 + \operatorname{twopow}(k)), \operatorname{rdiag}(t0 + 1, t0)) \operatorname{and} \operatorname{center}(t0)))\).
\end{theorem}
\begin{proof}
Assume \(hP\), \(h\).
The goal follows from \(\operatorname{periodiccrosscarrycenterslotand}(T, P, hP, h)\).
\end{proof}
\begin{theorem}
\label{ll:paper:periodic-crosscarry-colm1-and}
For every natural number \(T\) and natural number \(P\), if \(0 < P\), then if \(\operatorname{centerperiodic}(T, P)\), then for every natural number \(d\) and natural number \(k\), if \(T \le d + 1\), then if \(d + 2 < \operatorname{twopow}(k)\), then \(\operatorname{bitcrosscarryeq}(\operatorname{wittscale}(k, \operatorname{wittrow}(d + 1)), \operatorname{wittwin}(d + 1, \operatorname{twopow}(k)), d + 2 + \operatorname{twopow}(k), \operatorname{bxor}(\operatorname{ldiag}(d, d + 1) \operatorname{and} \operatorname{bxor}(\operatorname{center}(d + 1 + \operatorname{lag}(P, k)), \operatorname{center}(d + 1)), \operatorname{bxor}(\operatorname{ldiag}(d + \operatorname{twopow}(k), d + 1 + \operatorname{twopow}(k)), \operatorname{ldiag}(d, d + 1)) \operatorname{and} \operatorname{center}(d + 1)))\).
\end{theorem}
\begin{proof}
Assume \(hP\), \(h\).
The goal follows from \(\operatorname{periodiccrosscarrycolm1and}(T, P, hP, h)\).
\end{proof}
\begin{theorem}
\label{ll:paper:crosscarry-row-scale-at-threshold-eq-rdiag}
For every natural number \(t0\) and natural number \(k\) and natural number \(J\), if \(\operatorname{rightthreshold}(k, J)\), then if \(J + 1 < \operatorname{twopow}(k)\), then \(\operatorname{bitcrosscarryeq}(\operatorname{wittscale}(k, \operatorname{wittrow}(t0)), \operatorname{wittwin}(t0, \operatorname{twopow}(k)), J + 1, \operatorname{rdiag}(J, t0))\).
\end{theorem}
\begin{proof}
The goal follows from \(\operatorname{crosscarryrowscaleatthresholdeqrdiag}(t0, k, J)\).
\end{proof}
\begin{theorem}
\label{ll:paper:carry-row-scale-at-threshold}
For every natural number \(t0\) and natural number \(k\) and natural number \(J\), if \(\operatorname{rightthreshold}(k, J)\), then if \(J + 1 < \operatorname{twopow}(k)\), then \(\operatorname{bitcarryeq}(\operatorname{wittrow}(t0 + \operatorname{twopow}(k)), J + 1, \operatorname{bxor}(\operatorname{rdiag}(J + 1, t0) \operatorname{and} \operatorname{rdiag}(J, t0), \operatorname{rdiag}(J, t0)))\).
\end{theorem}
\begin{proof}
The goal follows from \(\operatorname{carryrowscaleatthreshold}(t0, k, J)\).
\end{proof}
\begin{theorem}
\label{ll:paper:periodic-winwitt-step-center-slot-expanded}
For every natural number \(T\) and natural number \(P\), if \(0 < P\), then if \(\operatorname{centerperiodic}(T, P)\), then for every natural number \(t0\) and natural number \(k\), if \(T \le t0\), then if \(2 \cdot t0 + 2 < \operatorname{twopow}(k)\), then \(\operatorname{bitwinwitteq}(t0 + 1, \operatorname{twopow}(k), t0 + 2 + \operatorname{twopow}(k), \operatorname{bxor}(\operatorname{bxor}(\operatorname{bxor}(\operatorname{rdiag}(t0 + 2 + \operatorname{twopow}(k), t0 + \operatorname{twopow}(k)), \operatorname{rdiag}(t0 + 2, t0)), \operatorname{bxor}(\operatorname{bxor}(\operatorname{rdiag}(t0 + 1 + \operatorname{twopow}(k), t0 + \operatorname{twopow}(k)), \operatorname{rdiag}(t0 + 1, t0)), \operatorname{bxor}(\operatorname{center}(t0 + \operatorname{lag}(P, k)), \operatorname{center}(t0)))), \operatorname{bxor}(\operatorname{bxor}(\operatorname{rdiag}(t0 + 1 + \operatorname{twopow}(k), t0 + \operatorname{twopow}(k)), \operatorname{rdiag}(t0 + 1, t0)) \operatorname{and} \operatorname{center}(t0 + \operatorname{lag}(P, k)), \operatorname{rdiag}(t0 + 1, t0) \operatorname{and} \operatorname{bxor}(\operatorname{center}(t0 + \operatorname{lag}(P, k)), \operatorname{center}(t0)))))\).
\end{theorem}
\begin{proof}
Assume \(hP\), \(h\).
The goal follows from \(\operatorname{periodicwinwittstepcenterslotexpanded}(T, P, hP, h)\).
\end{proof}
\begin{theorem}
\label{ll:paper:periodic-crosscarry-at-center-bit-and}
For every natural number \(T\) and natural number \(P\), if \(0 < P\), then if \(\operatorname{centerperiodic}(T, P)\), then for every natural number \(s\) and natural number \(k\), if \(T \le s + 1\), then if \(s + 2 < \operatorname{twopow}(k)\), then \(\operatorname{bitcrosscarryeq}(\operatorname{wittscale}(k, \operatorname{wittrow}(s + 1)), \operatorname{wittwin}(s + 1, \operatorname{twopow}(k)), s + 1 + \operatorname{twopow}(k), \operatorname{bxor}(\operatorname{center}(s + 1) \operatorname{and} \operatorname{bxor}(\operatorname{rdiag}(s + \operatorname{twopow}(k), s + 1 + \operatorname{twopow}(k)), \operatorname{rdiag}(s, s + 1)), \operatorname{bxor}(\operatorname{center}(s + 1 + \operatorname{lag}(P, k)), \operatorname{center}(s + 1)) \operatorname{and} \operatorname{rdiag}(s, s + 1)))\).
\end{theorem}
\begin{proof}
Assume \(hP\), \(h\).
The goal follows from \(\operatorname{periodiccrosscarryatcenterbitand}(T, P, hP, h)\).
\end{proof}
\begin{theorem}
\label{ll:paper:periodic-crosscarry-at-white-center-eq-maskparity}
For every natural number \(T\) and natural number \(P\), if \(0 < P\), then if \(\operatorname{centerperiodic}(T, P)\), then for every natural number \(s\) and natural number \(k\), if \(T \le s + 1\), then if \(s + 2 < \operatorname{twopow}(k)\), then if \(\operatorname{center}(s + 1) = \operatorname{false}\), then \(\operatorname{bitcrosscarryeq}(\operatorname{wittscale}(k, \operatorname{wittrow}(s + 1)), \operatorname{wittwin}(s + 1, \operatorname{twopow}(k)), s + 1 + \operatorname{twopow}(k), \operatorname{center}(s + 1 + \operatorname{lag}(P, k)) \operatorname{and} \operatorname{maskparity}(s + 1, 0))\).
\end{theorem}
\begin{proof}
Assume \(hP\), \(h\).
The goal follows from \(\operatorname{periodiccrosscarryatwhitecentereqmaskparity}(T, P, hP, h)\).
\end{proof}
\begin{theorem}
\label{ll:paper:carry-row-xor-maskparity-zero-eq-col1}
For every natural number \(s\), \(\operatorname{bitcarryeq}(\operatorname{wittrow}(s + 1), s + 1, \operatorname{bxor}(\operatorname{rdiag}(s, s + 1), \operatorname{maskparity}(s + 1, 0)))\).
\end{theorem}
\begin{proof}
The goal follows from \(\operatorname{carryrowxormaskparityzeroeqcol1}(s)\).
\end{proof}
\begin{theorem}
\label{ll:paper:periodic-crosscarry-colm1-at-white-center}
For every natural number \(T\) and natural number \(P\), if \(0 < P\), then if \(\operatorname{centerperiodic}(T, P)\), then for every natural number \(d\) and natural number \(k\), if \(T \le d + 1\), then if \(d + 2 < \operatorname{twopow}(k)\), then if \(\operatorname{center}(d + 1) = \operatorname{false}\), then \(\operatorname{bitcrosscarryeq}(\operatorname{wittscale}(k, \operatorname{wittrow}(d + 1)), \operatorname{wittwin}(d + 1, \operatorname{twopow}(k)), d + 2 + \operatorname{twopow}(k), \operatorname{ldiag}(d, d + 1) \operatorname{and} \operatorname{center}(d + 1 + \operatorname{lag}(P, k)))\).
\end{theorem}
\begin{proof}
Assume \(hP\), \(h\).
The goal follows from \(\operatorname{periodiccrosscarrycolm1atwhitecenter}(T, P, hP, h)\).
\end{proof}
\begin{theorem}
\label{ll:paper:periodic-crosscarry-colm1-white-center-lag-eq}
For every natural number \(T\) and natural number \(P\), if \(0 < P\), then if \(\operatorname{centerperiodic}(T, P)\), then for every natural number \(d\) and natural number \(k1\) and natural number \(k2\), if \(\operatorname{lag}(P, k1) = \operatorname{lag}(P, k2)\), then if \(T \le d + 1\), then if \(d + 2 < \operatorname{twopow}(k1)\), then if \(d + 2 < \operatorname{twopow}(k2)\), then if \(\operatorname{center}(d + 1) = \operatorname{false}\), then \(\operatorname{wittbit}(\operatorname{crosscarry}(\operatorname{wittscale}(k1, \operatorname{wittrow}(d + 1)), \operatorname{wittwin}(d + 1, \operatorname{twopow}(k1))), d + 2 + \operatorname{twopow}(k1)) = \operatorname{wittbit}(\operatorname{crosscarry}(\operatorname{wittscale}(k2, \operatorname{wittrow}(d + 1)), \operatorname{wittwin}(d + 1, \operatorname{twopow}(k2))), d + 2 + \operatorname{twopow}(k2))\).
\end{theorem}
\begin{proof}
Assume \(hP\), \(h\).
The goal follows from \(\operatorname{periodiccrosscarrycolm1whitecenterlageq}(T, P, hP, h)\).
\end{proof}
\begin{theorem}
\label{ll:paper:crosscarry-row-scale-sector-two-colm1-and}
For every natural number \(d\) and natural number \(k\), if \(d + 1 \le k\), then \(\operatorname{bitcrosscarryeq}(\operatorname{wittscale}(k, \operatorname{wittrow}(d + 1)), \operatorname{wittwin}(d + 1, \operatorname{twopow}(k)), d + 2 + \operatorname{twopow}(\operatorname{succ}(k)), \operatorname{bxor}(\operatorname{ldiag}(d, d + 1) \operatorname{and} \operatorname{bxor}(\operatorname{center}(d + 1), \operatorname{ldiag}(d + 1, d + 1 + \operatorname{twopow}(k))), \operatorname{bxor}(\operatorname{ldiag}(d, d + 1), \operatorname{ldiag}(d, d + 1 + \operatorname{twopow}(k))) \operatorname{and} \operatorname{center}(d + 1)))\).
\end{theorem}
\begin{proof}
The goal follows from \(\operatorname{crosscarryrowscalesectortwocolm1and}(d, k)\).
\end{proof}
\begin{theorem}
\label{ll:paper:crosscarry-row-scale-sector-two-colm1-at-white-center}
For every natural number \(d\) and natural number \(k\), if \(d + 1 \le k\), then if \(\operatorname{center}(d + 1) = \operatorname{false}\), then \(\operatorname{bitcrosscarryeq}(\operatorname{wittscale}(k, \operatorname{wittrow}(d + 1)), \operatorname{wittwin}(d + 1, \operatorname{twopow}(k)), d + 2 + \operatorname{twopow}(\operatorname{succ}(k)), \operatorname{ldiag}(d, d + 1) \operatorname{and} \operatorname{ldiag}(d + 1, d + 1 + \operatorname{twopow}(k)))\).
\end{theorem}
\begin{proof}
The goal follows from \(\operatorname{crosscarryrowscalesectortwocolm1atwhitecenter}(d, k)\).
\end{proof}
\begin{theorem}
\label{ll:paper:periodic-crosscarry-sector-one-xor-two-colm1-and}
For every natural number \(T\) and natural number \(P\), if \(0 < P\), then if \(\operatorname{centerperiodic}(T, P)\), then for every natural number \(d\) and natural number \(k\), if \(T \le d + 1\), then if \(d + 1 \le k\), then if \(d + 2 < \operatorname{twopow}(k)\), then \(\operatorname{bitcrosscarryxoreq}(\operatorname{wittscale}(k, \operatorname{wittrow}(d + 1)), \operatorname{wittwin}(d + 1, \operatorname{twopow}(k)), d + 2 + \operatorname{twopow}(k), d + 2 + \operatorname{twopow}(\operatorname{succ}(k)), \operatorname{bxor}(\operatorname{ldiag}(d, d + 1) \operatorname{and} \operatorname{bxor}(\operatorname{center}(d + 1 + \operatorname{lag}(P, k)), \operatorname{ldiag}(d + 1, d + 1 + \operatorname{twopow}(k))), \operatorname{bxor}(\operatorname{ldiag}(d + \operatorname{twopow}(k), d + 1 + \operatorname{twopow}(k)), \operatorname{ldiag}(d, d + 1 + \operatorname{twopow}(k))) \operatorname{and} \operatorname{center}(d + 1)))\).
\end{theorem}
\begin{proof}
Assume \(hP\), \(h\).
The goal follows from \(\operatorname{periodiccrosscarrysectoronexortwocolm1and}(T, P, hP, h)\).
\end{proof}
\begin{theorem}
\label{ll:paper:phiscale-winwitt-triality-low}
For every natural number \(d\) and natural number \(s\) and natural number \(k\), if \(d + 2 \cdot s \le k\), then \(\operatorname{wittbit}(\operatorname{wittscale}(k, \operatorname{wittwin}(d + s, \operatorname{twopow}(k))), d + 2 \cdot s + \operatorname{twopow}(\operatorname{succ}(k))) = \operatorname{bton}(\operatorname{bxor}(\operatorname{ldiag}(d + \operatorname{twopow}(k), d + s + \operatorname{twopow}(k)), \operatorname{ldiag}(d, d + s + \operatorname{twopow}(k))))\).
\end{theorem}
\begin{proof}
The goal follows from \(\operatorname{phiscalewinwitttrialitylow}(d, s, k)\).
\end{proof}
\begin{theorem}
\label{ll:paper:periodic-crosscarry-winwitt-cocycle-colm1-and}
For every natural number \(T\) and natural number \(P\), if \(0 < P\), then if \(\operatorname{centerperiodic}(T, P)\), then for every natural number \(d\) and natural number \(k\), if \(T \le d + 1\), then if \(d + 2 \le k\), then \(\operatorname{bitcrosscarryeq}(\operatorname{wittwin}(d + 1, \operatorname{twopow}(\operatorname{succ}(k))), \operatorname{wittscale}(k, \operatorname{wittwin}(d + 1, \operatorname{twopow}(k))), d + 2 + \operatorname{twopow}(\operatorname{succ}(k)), \operatorname{bxor}(\operatorname{bxor}(\operatorname{ldiag}(d + \operatorname{twopow}(\operatorname{succ}(k)), d + 1 + \operatorname{twopow}(\operatorname{succ}(k))), \operatorname{ldiag}(d, d + 1)) \operatorname{and} \operatorname{bxor}(\operatorname{center}(d + 1 + \operatorname{lag}(P, k)), \operatorname{ldiag}(d + 1, d + 1 + \operatorname{twopow}(k))), \operatorname{bxor}(\operatorname{ldiag}(d + \operatorname{twopow}(k), d + 1 + \operatorname{twopow}(k)), \operatorname{ldiag}(d, d + 1 + \operatorname{twopow}(k))) \operatorname{and} \operatorname{bxor}(\operatorname{center}(d + 1 + \operatorname{lag}(P, \operatorname{succ}(k))), \operatorname{center}(d + 1))))\).
\end{theorem}
\begin{proof}
Assume \(hP\), \(h\).
The goal follows from \(\operatorname{periodiccrosscarrywinwittcocyclecolm1and}(T, P, hP, h)\).
\end{proof}
\begin{theorem}
\label{ll:paper:carry-row-late-at-threshold-center-eq-maskparity}
For every natural number \(k\) and natural number \(J\), if \(\operatorname{rightthreshold}(k, J)\), then if \(J + 1 < \operatorname{twopow}(k)\), then \(\operatorname{bitcarryeq}(\operatorname{wittrow}(J + 1 + \operatorname{twopow}(k)), J + 1, \operatorname{maskparity}(J + 1, 0))\).
\end{theorem}
\begin{proof}
The goal follows from \(\operatorname{carryrowlateatthresholdcentereqmaskparity}(k, J)\).
\end{proof}
\begin{theorem}
\label{ll:paper:maskparity-zero-dyadic-jump-eq-witt-carry-crosscarry}
For every natural number \(s\) and natural number \(k\), if \(s + 2 < \operatorname{twopow}(k)\), then \(\operatorname{bitwinwitteq}(s + 1, \operatorname{twopow}(k), s + \operatorname{twopow}(k), \operatorname{bxor}(\operatorname{bxor}(\operatorname{maskparity}(s + 1 + \operatorname{twopow}(k), 0), \operatorname{maskparity}(s + 1, 0)), \operatorname{bxor}(\operatorname{bxor}(\operatorname{rdiag}(s + \operatorname{twopow}(k), s + 1 + \operatorname{twopow}(k)), \operatorname{rdiag}(s, s + 1)) \operatorname{and} \operatorname{wittresp}(s + 1, k), \operatorname{bxor}(\operatorname{center}(s + 1) \operatorname{and} \operatorname{bxor}(\operatorname{rdiag}(s + \operatorname{twopow}(k), s + 1 + \operatorname{twopow}(k)), \operatorname{rdiag}(s, s + 1)), \operatorname{wittresp}(s + 1, k) \operatorname{and} \operatorname{rdiag}(s, s + 1)))))\).
\end{theorem}
\begin{proof}
The goal follows from \(\operatorname{maskparityzerodyadicjumpeqwittcarrycrosscarry}(s, k)\).
\end{proof}
\begin{theorem}
\label{ll:paper:maskparity-zero-dyadic-jump-eq-not-witt-and-xor-crosscarry}
For every natural number \(s\) and natural number \(k\), if \(s + 2 < \operatorname{twopow}(k)\), then \(\operatorname{bitcrosscarryeq}(\operatorname{wittscale}(k, \operatorname{wittrow}(s + 1)), \operatorname{wittwin}(s + 1, \operatorname{twopow}(k)), s + 1 + \operatorname{twopow}(k), \operatorname{bxor}(\operatorname{bxor}(\operatorname{maskparity}(s + 1 + \operatorname{twopow}(k), 0), \operatorname{maskparity}(s + 1, 0)), \operatorname{not}(\operatorname{wittresp}(s + 1, k)) \operatorname{and} \operatorname{bxor}(\operatorname{rdiag}(s + \operatorname{twopow}(k), s + 1 + \operatorname{twopow}(k)), \operatorname{rdiag}(s, s + 1))))\).
\end{theorem}
\begin{proof}
The goal follows from \(\operatorname{maskparityzerodyadicjumpeqnotwittandxorcrosscarry}(s, k)\).
\end{proof}
\begin{theorem}
\label{ll:paper:maskparity-dyadic-jump-eq-crosscarry-of-wittresp-true}
For every natural number \(s\) and natural number \(k\), if \(s + 2 < \operatorname{twopow}(k)\), then if \(\operatorname{wittresp}(s + 1, k) = \operatorname{true}\), then \(\operatorname{bitcrosscarryeq}(\operatorname{wittscale}(k, \operatorname{wittrow}(s + 1)), \operatorname{wittwin}(s + 1, \operatorname{twopow}(k)), s + 1 + \operatorname{twopow}(k), \operatorname{bxor}(\operatorname{maskparity}(s + 1 + \operatorname{twopow}(k), 0), \operatorname{maskparity}(s + 1, 0)))\).
\end{theorem}
\begin{proof}
The goal follows from \(\operatorname{maskparitydyadicjumpeqcrosscarryofwittresptrue}(s, k)\).
\end{proof}
\begin{theorem}
\label{ll:paper:maskparity-dyadic-jump-eq-not-center-and-witt-of-wittresp-false}
For every natural number \(s\) and natural number \(k\), if \(s + 2 < \operatorname{twopow}(k)\), then if \(\operatorname{wittresp}(s + 1, k) = \operatorname{false}\), then \(\operatorname{bxor}(\operatorname{maskparity}(s + 1 + \operatorname{twopow}(k), 0), \operatorname{maskparity}(s + 1, 0)) = \operatorname{not}(\operatorname{center}(s + 1)) \operatorname{and} \operatorname{bxor}(\operatorname{rdiag}(s + \operatorname{twopow}(k), s + 1 + \operatorname{twopow}(k)), \operatorname{rdiag}(s, s + 1))\).
\end{theorem}
\begin{proof}
The goal follows from \(\operatorname{maskparitydyadicjumpeqnotcenterandwittofwittrespfalse}(s, k)\).
\end{proof}
\begin{theorem}
\label{ll:paper:periodic-leftresponse-step-diff-eq-crosscarry-of-out-of-phase}
For every natural number \(T\) and natural number \(P\), if \(0 < P\), then if \(\operatorname{centerperiodic}(T, P)\), then for every natural number \(s\) and natural number \(k\), if \(T \le s + 1\), then if \(s + 2 < \operatorname{twopow}(k)\), then if \(\operatorname{wittresp}(s + 1, k) = \operatorname{true}\), then \(\operatorname{bitcrosscarryeq}(\operatorname{wittscale}(k, \operatorname{wittrow}(s + 1)), \operatorname{wittwin}(s + 1, \operatorname{twopow}(k)), s + 1 + \operatorname{twopow}(k), \operatorname{bxor}(\operatorname{bxor}(\operatorname{leftresponse}(s + 2, k), \operatorname{leftresponse}(s + 1, k)), \operatorname{bxor}(\operatorname{bxor}(\operatorname{center}(s + 2 + \operatorname{lag}(P, k)), \operatorname{center}(s + 1 + \operatorname{lag}(P, k))), \operatorname{bxor}(\operatorname{center}(s + 2), \operatorname{center}(s + 1)))))\).
\end{theorem}
\begin{proof}
Assume \(hP\), \(h\).
The goal follows from \(\operatorname{periodicleftresponsestepdiffeqcrosscarryofoutofphase}(T, P, hP, h)\).
\end{proof}
\begin{theorem}
\label{ll:paper:crosscarry-winwitt-cocycle-at-threshold-succ}
For every natural number \(t0\) and natural number \(k\) and natural number \(J\), if \(\operatorname{rightthreshold}(k, J)\), then if \(J + 2 < \operatorname{twopow}(k)\), then \(\operatorname{bitcrosscarryeq}(\operatorname{wittwin}(t0, \operatorname{twopow}(\operatorname{succ}(k))), \operatorname{wittscale}(k, \operatorname{wittwin}(t0, \operatorname{twopow}(k))), J + 2, \operatorname{bxor}(\operatorname{rdiag}(J + 2, t0 + \operatorname{twopow}(\operatorname{succ}(k))), \operatorname{rdiag}(J + 2, t0)))\).
\end{theorem}
\begin{proof}
The goal follows from \(\operatorname{crosscarrywinwittcocycleatthresholdsucc}(t0, k, J)\).
\end{proof}
\begin{theorem}
\label{ll:paper:crosscarry-row-scale-eq-crosscarry-late-row}
For every natural number \(t0\) and natural number \(k\), \(\operatorname{crosscarry}(\operatorname{wittscale}(k, \operatorname{wittrow}(t0)), \operatorname{wittwin}(t0, \operatorname{twopow}(k))) = \operatorname{crosscarry}(\operatorname{wittscale}(k, \operatorname{wittrow}(t0)), \operatorname{wittrow}(t0 + \operatorname{twopow}(k)))\).
\end{theorem}
\begin{proof}
The goal follows from \(\operatorname{crosscarryrowscaleeqcrosscarrylaterow}(t0, k)\).
\end{proof}
\begin{theorem}
\label{ll:paper:periodic-crosscarry-colm1-wronskian}
For every natural number \(T\) and natural number \(P\), if \(0 < P\), then if \(\operatorname{centerperiodic}(T, P)\), then for every natural number \(d\) and natural number \(k\), if \(T \le d + 1\), then if \(d + 2 < \operatorname{twopow}(k)\), then \(\operatorname{bitcrosscarryeq}(\operatorname{wittscale}(k, \operatorname{wittrow}(d + 1)), \operatorname{wittwin}(d + 1, \operatorname{twopow}(k)), d + 2 + \operatorname{twopow}(k), \operatorname{bxor}(\operatorname{ldiag}(d, d + 1) \operatorname{and} \operatorname{center}(d + 1 + \operatorname{lag}(P, k)), \operatorname{ldiag}(d + \operatorname{twopow}(k), d + 1 + \operatorname{twopow}(k)) \operatorname{and} \operatorname{center}(d + 1)))\).
\end{theorem}
\begin{proof}
Assume \(hP\), \(h\).
The goal follows from \(\operatorname{periodiccrosscarrycolm1wronskian}(T, P, hP, h)\).
\end{proof}
\begin{theorem}
\label{ll:paper:crosscarry-row-scale-sector-two-colm1-wronskian}
For every natural number \(d\) and natural number \(k\), if \(d + 1 \le k\), then \(\operatorname{bitcrosscarryeq}(\operatorname{wittscale}(k, \operatorname{wittrow}(d + 1)), \operatorname{wittwin}(d + 1, \operatorname{twopow}(k)), d + 2 + \operatorname{twopow}(\operatorname{succ}(k)), \operatorname{bxor}(\operatorname{ldiag}(d, d + 1) \operatorname{and} \operatorname{ldiag}(d + 1, d + 1 + \operatorname{twopow}(k)), \operatorname{ldiag}(d, d + 1 + \operatorname{twopow}(k)) \operatorname{and} \operatorname{center}(d + 1)))\).
\end{theorem}
\begin{proof}
The goal follows from \(\operatorname{crosscarryrowscalesectortwocolm1wronskian}(d, k)\).
\end{proof}
\begin{theorem}
\label{ll:paper:maskparity-zero-and-center-eq-false}
For every natural number \(t\), \(\operatorname{maskparity}(t, 0) \operatorname{and} \operatorname{center}(t) = \operatorname{false}\).
\end{theorem}
\begin{proof}
The goal follows from \(\operatorname{maskparityzeroandcentereqfalse}(t)\).
\end{proof}
\begin{theorem}
\label{ll:paper:periodic-colm1-defect-and-center-eq-false}
For every natural number \(T\) and natural number \(P\), if \(0 < P\), then if \(\operatorname{centerperiodic}(T, P)\), then for every natural number \(d\), if \(T \le d + 1\), then \(\operatorname{bxor}(\operatorname{ldiag}(d + P, d + P + 1), \operatorname{ldiag}(d, d + 1)) \operatorname{and} \operatorname{center}(d + 1) = \operatorname{false}\).
\end{theorem}
\begin{proof}
Assume \(hP\), \(h\).
The goal follows from \(\operatorname{periodiccolm1defectandcentereqfalse}(T, P, hP, h)\).
\end{proof}
\begin{theorem}
\label{ll:paper:periodic-crosscarry-colm1-shift-p-eq-leftresponse-and-witt}
For every natural number \(T\) and natural number \(P\), if \(0 < P\), then if \(\operatorname{centerperiodic}(T, P)\), then for every natural number \(d\) and natural number \(k\), if \(T \le d + 1\), then if \(d + P + 2 < \operatorname{twopow}(k)\), then \(\operatorname{bitcrosscarryeq}(\operatorname{wittscale}(k, \operatorname{wittrow}(d + P + 1)), \operatorname{wittwin}(d + P + 1, \operatorname{twopow}(k)), d + P + 2 + \operatorname{twopow}(k), \operatorname{bxor}(\operatorname{bxor}(\operatorname{ldiag}(d, d + 1) \operatorname{and} \operatorname{center}(d + 1 + \operatorname{lag}(P, k)), \operatorname{ldiag}(d + \operatorname{twopow}(k), d + 1 + \operatorname{twopow}(k)) \operatorname{and} \operatorname{center}(d + 1)), \operatorname{bxor}(\operatorname{bxor}(\operatorname{leftresponse}(d + P + 2, k), \operatorname{leftresponse}(d + 2, k)), \operatorname{bxor}(\operatorname{leftresponse}(d + P + 1, k), \operatorname{leftresponse}(d + 1, k))) \operatorname{and} \operatorname{wittresp}(d + 1, k)))\).
\end{theorem}
\begin{proof}
Assume \(hP\), \(h\).
The goal follows from \(\operatorname{periodiccrosscarrycolm1shiftpeqleftresponseandwitt}(T, P, hP, h)\).
\end{proof}
\begin{theorem}
\label{ll:paper:periodic-crosscarry-colm1-shift-p-eq-of-wittresp-false}
For every natural number \(T\) and natural number \(P\), if \(0 < P\), then if \(\operatorname{centerperiodic}(T, P)\), then for every natural number \(d\) and natural number \(k\), if \(T \le d + 1\), then if \(d + P + 2 < \operatorname{twopow}(k)\), then if \(\operatorname{wittresp}(d + 1, k) = \operatorname{false}\), then \(\operatorname{wittbit}(\operatorname{crosscarry}(\operatorname{wittscale}(k, \operatorname{wittrow}(d + P + 1)), \operatorname{wittwin}(d + P + 1, \operatorname{twopow}(k))), d + P + 2 + \operatorname{twopow}(k)) = \operatorname{wittbit}(\operatorname{crosscarry}(\operatorname{wittscale}(k, \operatorname{wittrow}(d + 1)), \operatorname{wittwin}(d + 1, \operatorname{twopow}(k))), d + 2 + \operatorname{twopow}(k))\).
\end{theorem}
\begin{proof}
Assume \(hP\), \(h\).
The goal follows from \(\operatorname{periodiccrosscarrycolm1shiftpeqofwittrespfalse}(T, P, hP, h)\).
\end{proof}
\begin{theorem}
\label{ll:paper:periodic-leftresponse-step-diff-eq-witt-of-in-phase-black}
For every natural number \(T\) and natural number \(P\), if \(0 < P\), then if \(\operatorname{centerperiodic}(T, P)\), then for every natural number \(s\) and natural number \(k\), if \(T \le s + 1\), then if \(s + 2 < \operatorname{twopow}(k)\), then if \(\operatorname{center}(s + 1) = \operatorname{true}\), then if \(\operatorname{wittresp}(s + 1, k) = \operatorname{false}\), then \(\operatorname{bxor}(\operatorname{leftresponse}(s + 2, k), \operatorname{leftresponse}(s + 1, k)) = \operatorname{wittresp}(s + 2, k)\).
\end{theorem}
\begin{proof}
Assume \(hP\), \(h\).
The goal follows from \(\operatorname{periodicleftresponsestepdiffeqwittofinphaseblack}(T, P, hP, h)\).
\end{proof}
\begin{theorem}
\label{ll:paper:phiscale-row-testbit-outside-cone}
For every natural number \(t0\) and natural number \(k\) and natural number \(r\), if \(\operatorname{twopow}(k) \le r\), then \(\operatorname{wittbit}(\operatorname{wittscale}(k, \operatorname{wittrow}(t0)), r + \operatorname{twopow}(k)) = \operatorname{bton}(\operatorname{bxor}(\operatorname{bxor}(\operatorname{rdiag}(r + \operatorname{twopow}(k), t0), \operatorname{rdiag}(r, t0)), \operatorname{rdiag}(r - \operatorname{twopow}(k), t0)))\).
\end{theorem}
\begin{proof}
The goal follows from \(\operatorname{phiscalerowtestbitoutsidecone}(t0, k, r)\).
\end{proof}
\begin{theorem}
\label{ll:paper:winwitt-testbit-outside-cone}
For every natural number \(t0\) and natural number \(k\) and natural number \(r\), if \(\operatorname{twopow}(k) \le r\), then \(\operatorname{bitwinwitteq}(t0, \operatorname{twopow}(k), r + \operatorname{twopow}(k), \operatorname{bxor}(\operatorname{rdiag}(r + \operatorname{twopow}(k), t0 + \operatorname{twopow}(k)), \operatorname{bxor}(\operatorname{bxor}(\operatorname{rdiag}(r + \operatorname{twopow}(k), t0), \operatorname{rdiag}(r, t0)), \operatorname{rdiag}(r - \operatorname{twopow}(k), t0))))\).
\end{theorem}
\begin{proof}
The goal follows from \(\operatorname{winwitttestbitoutsidecone}(t0, k, r)\).
\end{proof}
\begin{theorem}
\label{ll:paper:periodic-crosscarry-colm1-outside-cone-and}
For every natural number \(T\) and natural number \(P\), if \(0 < P\), then if \(\operatorname{centerperiodic}(T, P)\), then for every natural number \(d\) and natural number \(k\), if \(T \le d + 1\), then if \(\operatorname{twopow}(k) \le d\), then \(\operatorname{bitcrosscarryeq}(\operatorname{wittscale}(k, \operatorname{wittrow}(d + 1)), \operatorname{wittwin}(d + 1, \operatorname{twopow}(k)), d + 2 + \operatorname{twopow}(k), \operatorname{bxor}(\operatorname{bxor}(\operatorname{bxor}(\operatorname{ldiag}(d - \operatorname{twopow}(k), d + 1), \operatorname{ldiag}(d, d + 1)), \operatorname{rdiag}(d + 2 - \operatorname{twopow}(k), d + 1)) \operatorname{and} \operatorname{bxor}(\operatorname{center}(d + 1 + \operatorname{lag}(P, k)), \operatorname{bxor}(\operatorname{bxor}(\operatorname{ldiag}(d + 1 - \operatorname{twopow}(k), d + 1), \operatorname{center}(d + 1)), \operatorname{rdiag}(d + 1 - \operatorname{twopow}(k), d + 1))), \operatorname{bxor}(\operatorname{ldiag}(d + \operatorname{twopow}(k), d + 1 + \operatorname{twopow}(k)), \operatorname{bxor}(\operatorname{bxor}(\operatorname{ldiag}(d - \operatorname{twopow}(k), d + 1), \operatorname{ldiag}(d, d + 1)), \operatorname{rdiag}(d + 2 - \operatorname{twopow}(k), d + 1))) \operatorname{and} \operatorname{bxor}(\operatorname{bxor}(\operatorname{ldiag}(d + 1 - \operatorname{twopow}(k), d + 1), \operatorname{center}(d + 1)), \operatorname{rdiag}(d + 1 - \operatorname{twopow}(k), d + 1))))\).
\end{theorem}
\begin{proof}
Assume \(hP\), \(h\).
The goal follows from \(\operatorname{periodiccrosscarrycolm1outsideconeand}(T, P, hP, h)\).
\end{proof}
\begin{theorem}
\label{ll:paper:periodic-crosscarry-colm1-outside-cone-wronskian}
For every natural number \(T\) and natural number \(P\), if \(0 < P\), then if \(\operatorname{centerperiodic}(T, P)\), then for every natural number \(d\) and natural number \(k\), if \(T \le d + 1\), then if \(\operatorname{twopow}(k) \le d\), then \(\operatorname{bitcrosscarryeq}(\operatorname{wittscale}(k, \operatorname{wittrow}(d + 1)), \operatorname{wittwin}(d + 1, \operatorname{twopow}(k)), d + 2 + \operatorname{twopow}(k), \operatorname{bxor}(\operatorname{bxor}(\operatorname{bxor}(\operatorname{ldiag}(d - \operatorname{twopow}(k), d + 1), \operatorname{ldiag}(d, d + 1)), \operatorname{rdiag}(d + 2 - \operatorname{twopow}(k), d + 1)) \operatorname{and} \operatorname{center}(d + 1 + \operatorname{lag}(P, k)), \operatorname{ldiag}(d + \operatorname{twopow}(k), d + 1 + \operatorname{twopow}(k)) \operatorname{and} \operatorname{bxor}(\operatorname{bxor}(\operatorname{ldiag}(d + 1 - \operatorname{twopow}(k), d + 1), \operatorname{center}(d + 1)), \operatorname{rdiag}(d + 1 - \operatorname{twopow}(k), d + 1))))\).
\end{theorem}
\begin{proof}
Assume \(hP\), \(h\).
The goal follows from \(\operatorname{periodiccrosscarrycolm1outsideconewronskian}(T, P, hP, h)\).
\end{proof}
\begin{theorem}
\label{ll:paper:periodic-crosscarry-colm1-outside-cone-at-white-center}
For every natural number \(T\) and natural number \(P\), if \(0 < P\), then if \(\operatorname{centerperiodic}(T, P)\), then for every natural number \(d\) and natural number \(k\), if \(T \le d + 1\), then if \(\operatorname{twopow}(k) \le d\), then if \(\operatorname{center}(d + 1) = \operatorname{false}\), then \(\operatorname{bitcrosscarryeq}(\operatorname{wittscale}(k, \operatorname{wittrow}(d + 1)), \operatorname{wittwin}(d + 1, \operatorname{twopow}(k)), d + 2 + \operatorname{twopow}(k), \operatorname{bxor}(\operatorname{bxor}(\operatorname{bxor}(\operatorname{ldiag}(d - \operatorname{twopow}(k), d + 1), \operatorname{ldiag}(d, d + 1)), \operatorname{rdiag}(d + 2 - \operatorname{twopow}(k), d + 1)) \operatorname{and} \operatorname{center}(d + 1 + \operatorname{lag}(P, k)), \operatorname{ldiag}(d + \operatorname{twopow}(k), d + 1 + \operatorname{twopow}(k)) \operatorname{and} \operatorname{bxor}(\operatorname{ldiag}(d + 1 - \operatorname{twopow}(k), d + 1), \operatorname{rdiag}(d + 1 - \operatorname{twopow}(k), d + 1))))\).
\end{theorem}
\begin{proof}
Assume \(hP\), \(h\).
The goal follows from \(\operatorname{periodiccrosscarrycolm1outsideconeatwhitecenter}(T, P, hP, h)\).
\end{proof}
\begin{theorem}
\label{ll:paper:periodic-crosscarry-colm1-outside-cone-at-both-white-centers}
For every natural number \(T\) and natural number \(P\), if \(0 < P\), then if \(\operatorname{centerperiodic}(T, P)\), then for every natural number \(d\) and natural number \(k\), if \(T \le d + 1\), then if \(\operatorname{twopow}(k) \le d\), then if \(\operatorname{center}(d + 1) = \operatorname{false}\), then if \(\operatorname{center}(d + 1 + \operatorname{lag}(P, k)) = \operatorname{false}\), then \(\operatorname{bitcrosscarryeq}(\operatorname{wittscale}(k, \operatorname{wittrow}(d + 1)), \operatorname{wittwin}(d + 1, \operatorname{twopow}(k)), d + 2 + \operatorname{twopow}(k), \operatorname{ldiag}(d + \operatorname{twopow}(k), d + 1 + \operatorname{twopow}(k)) \operatorname{and} \operatorname{bxor}(\operatorname{ldiag}(d + 1 - \operatorname{twopow}(k), d + 1), \operatorname{rdiag}(d + 1 - \operatorname{twopow}(k), d + 1)))\).
\end{theorem}
\begin{proof}
Assume \(hP\), \(h\).
The goal follows from \(\operatorname{periodiccrosscarrycolm1outsideconeatbothwhitecenters}(T, P, hP, h)\).
\end{proof}
\begin{theorem}
\label{ll:paper:crosscarry-colm1-at-white-late-center-eq-and-witt}
For every natural number \(d\) and natural number \(k\), if \(\operatorname{center}(d + 1 + \operatorname{twopow}(k)) = \operatorname{false}\), then \(\operatorname{bitcrosscarryeq}(\operatorname{wittscale}(k, \operatorname{wittrow}(d + 1)), \operatorname{wittwin}(d + 1, \operatorname{twopow}(k)), d + 2 + \operatorname{twopow}(k), \operatorname{ldiag}(d + \operatorname{twopow}(k), d + 1 + \operatorname{twopow}(k)) \operatorname{and} \operatorname{wittresp}(d + 1, k))\).
\end{theorem}
\begin{proof}
The goal follows from \(\operatorname{crosscarrycolm1atwhitelatecentereqandwitt}(d, k)\).
\end{proof}
\begin{theorem}
\label{ll:paper:crosscarry-colm1-outside-cone-at-white-late-center-via-response}
For every natural number \(d\) and natural number \(k\), if \(\operatorname{twopow}(k) \le d + 1\), then if \(\operatorname{center}(d + 1 + \operatorname{twopow}(k)) = \operatorname{false}\), then \(\operatorname{bitcrosscarryeq}(\operatorname{wittscale}(k, \operatorname{wittrow}(d + 1)), \operatorname{wittwin}(d + 1, \operatorname{twopow}(k)), d + 2 + \operatorname{twopow}(k), \operatorname{ldiag}(d + \operatorname{twopow}(k), d + 1 + \operatorname{twopow}(k)) \operatorname{and} \operatorname{bxor}(\operatorname{response}(d + 1, k), \operatorname{rdiag}(d + 1 - \operatorname{twopow}(k), d + 1)))\).
\end{theorem}
\begin{proof}
The goal follows from \(\operatorname{crosscarrycolm1outsideconeatwhitelatecenterviaresponse}(d, k)\).
\end{proof}
\begin{theorem}
\label{ll:paper:crosscarry-row-scale-at-center-bit-outside-cone-and}
For every natural number \(s\) and natural number \(k\), if \(\operatorname{twopow}(k) \le s\), then \(\operatorname{bitcrosscarryeq}(\operatorname{wittscale}(k, \operatorname{wittrow}(s + 1)), \operatorname{wittwin}(s + 1, \operatorname{twopow}(k)), s + 1 + \operatorname{twopow}(k), \operatorname{bxor}(\operatorname{bxor}(\operatorname{bxor}(\operatorname{ldiag}(s + 1 - \operatorname{twopow}(k), s + 1), \operatorname{center}(s + 1)), \operatorname{rdiag}(s + 1 - \operatorname{twopow}(k), s + 1)) \operatorname{and} \operatorname{bxor}(\operatorname{rdiag}(s + \operatorname{twopow}(k), s + 1 + \operatorname{twopow}(k)), \operatorname{bxor}(\operatorname{bxor}(\operatorname{ldiag}(s + 2 - \operatorname{twopow}(k), s + 1), \operatorname{rdiag}(s, s + 1)), \operatorname{rdiag}(s - \operatorname{twopow}(k), s + 1))), \operatorname{wittresp}(s + 1, k) \operatorname{and} \operatorname{bxor}(\operatorname{bxor}(\operatorname{ldiag}(s + 2 - \operatorname{twopow}(k), s + 1), \operatorname{rdiag}(s, s + 1)), \operatorname{rdiag}(s - \operatorname{twopow}(k), s + 1))))\).
\end{theorem}
\begin{proof}
The goal follows from \(\operatorname{crosscarryrowscaleatcenterbitoutsideconeand}(s, k)\).
\end{proof}
\begin{theorem}
\label{ll:paper:crosscarry-at-white-center-outside-cone-noncenter-and}
For every natural number \(s\) and natural number \(k\), if \(\operatorname{twopow}(k) \le s\), then if \(\operatorname{center}(s + 1) = \operatorname{false}\), then \(\operatorname{bitcrosscarryeq}(\operatorname{wittscale}(k, \operatorname{wittrow}(s + 1)), \operatorname{wittwin}(s + 1, \operatorname{twopow}(k)), s + 1 + \operatorname{twopow}(k), \operatorname{bxor}(\operatorname{bxor}(\operatorname{ldiag}(s + 1 - \operatorname{twopow}(k), s + 1), \operatorname{rdiag}(s + 1 - \operatorname{twopow}(k), s + 1)) \operatorname{and} \operatorname{bxor}(\operatorname{rdiag}(s + \operatorname{twopow}(k), s + 1 + \operatorname{twopow}(k)), \operatorname{bxor}(\operatorname{bxor}(\operatorname{ldiag}(s + 2 - \operatorname{twopow}(k), s + 1), \operatorname{rdiag}(s, s + 1)), \operatorname{rdiag}(s - \operatorname{twopow}(k), s + 1))), \operatorname{wittresp}(s + 1, k) \operatorname{and} \operatorname{bxor}(\operatorname{bxor}(\operatorname{ldiag}(s + 2 - \operatorname{twopow}(k), s + 1), \operatorname{rdiag}(s, s + 1)), \operatorname{rdiag}(s - \operatorname{twopow}(k), s + 1))))\).
\end{theorem}
\begin{proof}
The goal follows from \(\operatorname{crosscarryatwhitecenteroutsideconenoncenterand}(s, k)\).
\end{proof}
\begin{theorem}
\label{ll:paper:crosscarry-at-center-bit-outside-cone-wronskian}
For every natural number \(s\) and natural number \(k\), if \(\operatorname{twopow}(k) \le s\), then \(\operatorname{bitcrosscarryeq}(\operatorname{wittscale}(k, \operatorname{wittrow}(s + 1)), \operatorname{wittwin}(s + 1, \operatorname{twopow}(k)), s + 1 + \operatorname{twopow}(k), \operatorname{bxor}(\operatorname{bxor}(\operatorname{bxor}(\operatorname{ldiag}(s + 1 - \operatorname{twopow}(k), s + 1), \operatorname{center}(s + 1)), \operatorname{rdiag}(s + 1 - \operatorname{twopow}(k), s + 1)) \operatorname{and} \operatorname{rdiag}(s + \operatorname{twopow}(k), s + 1 + \operatorname{twopow}(k)), \operatorname{center}(s + 1 + \operatorname{twopow}(k)) \operatorname{and} \operatorname{bxor}(\operatorname{bxor}(\operatorname{ldiag}(s + 2 - \operatorname{twopow}(k), s + 1), \operatorname{rdiag}(s, s + 1)), \operatorname{rdiag}(s - \operatorname{twopow}(k), s + 1))))\).
\end{theorem}
\begin{proof}
The goal follows from \(\operatorname{crosscarryatcenterbitoutsideconewronskian}(s, k)\).
\end{proof}
\begin{theorem}
\label{ll:paper:periodic-crosscarry-at-both-white-centers-outside-cone-eq-maskparity}
For every natural number \(T\) and natural number \(P\), if \(0 < P\), then if \(\operatorname{centerperiodic}(T, P)\), then for every natural number \(s\) and natural number \(k\), if \(T \le s + 1\), then if \(\operatorname{twopow}(k) \le s\), then if \(\operatorname{center}(s + 1) = \operatorname{false}\), then if \(\operatorname{center}(s + 1 + \operatorname{lag}(P, k)) = \operatorname{false}\), then \(\operatorname{bitcrosscarryeq}(\operatorname{wittscale}(k, \operatorname{wittrow}(s + 1)), \operatorname{wittwin}(s + 1, \operatorname{twopow}(k)), s + 1 + \operatorname{twopow}(k), \operatorname{bxor}(\operatorname{ldiag}(s + 1 - \operatorname{twopow}(k), s + 1), \operatorname{rdiag}(s + 1 - \operatorname{twopow}(k), s + 1)) \operatorname{and} \operatorname{maskparity}(s + 1 + \operatorname{twopow}(k), 0))\).
\end{theorem}
\begin{proof}
Assume \(hP\), \(h\).
The goal follows from \(\operatorname{periodiccrosscarryatbothwhitecentersoutsideconeeqmaskparity}(T, P, hP, h)\).
\end{proof}
\begin{theorem}
\label{ll:paper:crosscarry-at-white-late-center-eq-witt-and-maskparity}
For every natural number \(s\) and natural number \(k\), if \(\operatorname{center}(s + 1 + \operatorname{twopow}(k)) = \operatorname{false}\), then \(\operatorname{bitcrosscarryeq}(\operatorname{wittscale}(k, \operatorname{wittrow}(s + 1)), \operatorname{wittwin}(s + 1, \operatorname{twopow}(k)), s + 1 + \operatorname{twopow}(k), \operatorname{wittresp}(s + 1, k) \operatorname{and} \operatorname{maskparity}(s + 1 + \operatorname{twopow}(k), 0))\).
\end{theorem}
\begin{proof}
The goal follows from \(\operatorname{crosscarryatwhitelatecentereqwittandmaskparity}(s, k)\).
\end{proof}
\begin{theorem}
\label{ll:paper:periodic-maskparity-zero-defect-mul-eq-colm1-defect-mul}
For every natural number \(T\) and natural number \(P\), if \(0 < P\), then if \(\operatorname{centerperiodic}(T, P)\), then for every natural number \(m\) and natural number \(d\), if \(T \le d + 1\), then \(\operatorname{bxor}(\operatorname{maskparity}(d + m \cdot P + 1, 0), \operatorname{maskparity}(d + 1, 0)) = \operatorname{bxor}(\operatorname{ldiag}(d + m \cdot P, d + m \cdot P + 1), \operatorname{ldiag}(d, d + 1))\).
\end{theorem}
\begin{proof}
Assume \(hP\), \(h\).
The goal follows from \(\operatorname{periodicmaskparityzerodefectmuleqcolm1defectmul}(T, P, hP, h)\).
\end{proof}
\begin{theorem}
\label{ll:paper:wittresp-xor-response-shift-mul-outside-cone}
For every natural number \(t\) and natural number \(k\) and natural number \(m\) and natural number \(P\), if \(\operatorname{twopow}(k) \le t\), then \(\operatorname{bxor}(\operatorname{bxor}(\operatorname{wittresp}(t + m \cdot P, k), \operatorname{wittresp}(t, k)), \operatorname{bxor}(\operatorname{response}(t + m \cdot P, k), \operatorname{response}(t, k))) = \operatorname{bxor}(\operatorname{rdiag}(t + m \cdot P - \operatorname{twopow}(k), t + m \cdot P), \operatorname{rdiag}(t - \operatorname{twopow}(k), t))\).
\end{theorem}
\begin{proof}
The goal follows from \(\operatorname{wittrespxorresponseshiftmuloutsidecone}(t, k, m, P)\).
\end{proof}
\begin{theorem}
\label{ll:paper:periodic-wittresp-shift-mul-outside-cone}
For every natural number \(T\) and natural number \(P\), if \(0 < P\), then if \(\operatorname{centerperiodic}(T, P)\), then for every natural number \(t\) and natural number \(k\) and natural number \(m\), if \(T \le t\), then if \(\operatorname{twopow}(k) \le t\), then \(\operatorname{bxor}(\operatorname{wittresp}(t + m \cdot P, k), \operatorname{wittresp}(t, k)) = \operatorname{bxor}(\operatorname{bxor}(\operatorname{ldiag}(t + m \cdot P - \operatorname{twopow}(k), t + m \cdot P), \operatorname{ldiag}(t - \operatorname{twopow}(k), t)), \operatorname{bxor}(\operatorname{rdiag}(t + m \cdot P - \operatorname{twopow}(k), t + m \cdot P), \operatorname{rdiag}(t - \operatorname{twopow}(k), t)))\).
\end{theorem}
\begin{proof}
Assume \(hP\), \(h\).
The goal follows from \(\operatorname{periodicwittrespshiftmuloutsidecone}(T, P, hP, h)\).
\end{proof}
\begin{theorem}
\label{ll:paper:periodic-crosscarry-colm1-shift-mul-at-white-late-center}
For every natural number \(T\) and natural number \(P\), if \(0 < P\), then if \(\operatorname{centerperiodic}(T, P)\), then for every natural number \(m\) and natural number \(d\) and natural number \(k\), if \(T \le d + 1\), then if \(\operatorname{center}(d + 1 + \operatorname{lag}(P, k)) = \operatorname{false}\), then \(\operatorname{bitcrosscarryeq}(\operatorname{wittscale}(k, \operatorname{wittrow}(d + m \cdot P + 1)), \operatorname{wittwin}(d + m \cdot P + 1, \operatorname{twopow}(k)), d + m \cdot P + 2 + \operatorname{twopow}(k), \operatorname{bxor}(\operatorname{ldiag}(d + \operatorname{twopow}(k), d + 1 + \operatorname{twopow}(k)) \operatorname{and} \operatorname{wittresp}(d + 1, k), \operatorname{bxor}(\operatorname{bxor}(\operatorname{maskparity}(d + m \cdot P + 1 + \operatorname{twopow}(k), 0), \operatorname{maskparity}(d + 1 + \operatorname{twopow}(k), 0)) \operatorname{and} \operatorname{wittresp}(d + m \cdot P + 1, k), \operatorname{ldiag}(d + \operatorname{twopow}(k), d + 1 + \operatorname{twopow}(k)) \operatorname{and} \operatorname{bxor}(\operatorname{wittresp}(d + m \cdot P + 1, k), \operatorname{wittresp}(d + 1, k)))))\).
\end{theorem}
\begin{proof}
Assume \(hP\), \(h\).
The goal follows from \(\operatorname{periodiccrosscarrycolm1shiftmulatwhitelatecenter}(T, P, hP, h)\).
\end{proof}
\begin{theorem}
\label{ll:paper:periodic-crosscarry-center-bit-shift-mul-at-white-late-center}
For every natural number \(T\) and natural number \(P\), if \(0 < P\), then if \(\operatorname{centerperiodic}(T, P)\), then for every natural number \(m\) and natural number \(s\) and natural number \(k\), if \(T \le s + 1\), then if \(\operatorname{center}(s + 1 + \operatorname{lag}(P, k)) = \operatorname{false}\), then \(\operatorname{bitcrosscarryeq}(\operatorname{wittscale}(k, \operatorname{wittrow}(s + m \cdot P + 1)), \operatorname{wittwin}(s + m \cdot P + 1, \operatorname{twopow}(k)), s + m \cdot P + 1 + \operatorname{twopow}(k), \operatorname{bxor}(\operatorname{wittresp}(s + 1, k) \operatorname{and} \operatorname{maskparity}(s + 1 + \operatorname{twopow}(k), 0), \operatorname{bxor}(\operatorname{bxor}(\operatorname{wittresp}(s + m \cdot P + 1, k), \operatorname{wittresp}(s + 1, k)) \operatorname{and} \operatorname{maskparity}(s + m \cdot P + 1 + \operatorname{twopow}(k), 0), \operatorname{wittresp}(s + 1, k) \operatorname{and} \operatorname{bxor}(\operatorname{maskparity}(s + m \cdot P + 1 + \operatorname{twopow}(k), 0), \operatorname{maskparity}(s + 1 + \operatorname{twopow}(k), 0)))))\).
\end{theorem}
\begin{proof}
Assume \(hP\), \(h\).
The goal follows from \(\operatorname{periodiccrosscarrycenterbitshiftmulatwhitelatecenter}(T, P, hP, h)\).
\end{proof}
\begin{theorem}
\label{ll:paper:frobdefect-row-cone-eq-zero}
For every natural number \(t0\) and natural number \(k\), if \(2 \cdot t0 + 1 < \operatorname{twopow}(k)\), then \(\operatorname{frobdefect}(k, \operatorname{wittrow}(t0)) = 0\).
\end{theorem}
\begin{proof}
The goal follows from \(\operatorname{frobdefectrowconeeqzero}(t0, k)\).
\end{proof}
\begin{theorem}
\label{ll:paper:winwitt-step-frobdefect}
For every natural number \(t0\) and natural number \(k\), \(\operatorname{wittwin}(t0 + 1, \operatorname{twopow}(k)) = \operatorname{wittxor}(\operatorname{wittxor}(\operatorname{wittphi}(\operatorname{wittwin}(t0, \operatorname{twopow}(k))), \operatorname{wittshl}(\operatorname{crosscarry}(\operatorname{wittscale}(k, \operatorname{wittrow}(t0)), \operatorname{wittwin}(t0, \operatorname{twopow}(k))), 1)), \operatorname{wittshl}(\operatorname{frobdefect}(k, \operatorname{wittrow}(t0)), 1))\).
\end{theorem}
\begin{proof}
The goal follows from \(\operatorname{winwittstepfrobdefect}(t0, k)\).
\end{proof}
\begin{theorem}
\label{ll:paper:winwitt-step-frobdefect-testbit-succ}
For every natural number \(t0\) and natural number \(k\) and natural number \(j\), \(\operatorname{bitwinwitteq}(t0 + 1, \operatorname{twopow}(k), j + 1, \operatorname{bxor}(\operatorname{bxor}(\operatorname{wittbitbool}(\operatorname{wittphi}(\operatorname{wittwin}(t0, \operatorname{twopow}(k))), j + 1), \operatorname{wittbitbool}(\operatorname{crosscarry}(\operatorname{wittscale}(k, \operatorname{wittrow}(t0)), \operatorname{wittwin}(t0, \operatorname{twopow}(k))), j)), \operatorname{wittbitbool}(\operatorname{frobdefect}(k, \operatorname{wittrow}(t0)), j)))\).
\end{theorem}
\begin{proof}
The goal follows from \(\operatorname{winwittstepfrobdefecttestbitsucc}(t0, k, j)\).
\end{proof}
\begin{theorem}
\label{ll:paper:frobdefect-row-testbit-outside-cone}
For every natural number \(t0\) and natural number \(k\) and natural number \(r\), if \(\operatorname{twopow}(k) \le r\), then \(\operatorname{bitfrobdefecteq}(k, \operatorname{wittrow}(t0), r + 1 + \operatorname{twopow}(k), \operatorname{bxor}(\operatorname{bxor}(\operatorname{rdiag}(r + 1 + \operatorname{twopow}(k), t0) \operatorname{and} \operatorname{bxor}(\operatorname{rdiag}(r, t0), \operatorname{rdiag}(r - \operatorname{twopow}(k), t0)), \operatorname{rdiag}(r + 1, t0) \operatorname{and} \operatorname{bxor}(\operatorname{rdiag}(r + \operatorname{twopow}(k), t0), \operatorname{rdiag}(r - \operatorname{twopow}(k), t0))), \operatorname{rdiag}(r + 1 - \operatorname{twopow}(k), t0) \operatorname{and} \operatorname{bxor}(\operatorname{rdiag}(r + \operatorname{twopow}(k), t0), \operatorname{rdiag}(r, t0))))\).
\end{theorem}
\begin{proof}
The goal follows from \(\operatorname{frobdefectrowtestbitoutsidecone}(t0, k, r)\).
\end{proof}
\begin{theorem}
\label{ll:paper:frobdefect-row-center-slot-outside-cone}
For every natural number \(s\) and natural number \(k\), if \(\operatorname{twopow}(k) \le s\), then \(\operatorname{bitfrobdefecteq}(k, \operatorname{wittrow}(s + 1), s + 1 + \operatorname{twopow}(k), \operatorname{bxor}(\operatorname{bxor}(\operatorname{ldiag}(s + 1 - \operatorname{twopow}(k), s + 1) \operatorname{and} \operatorname{bxor}(\operatorname{rdiag}(s, s + 1), \operatorname{rdiag}(s - \operatorname{twopow}(k), s + 1)), \operatorname{center}(s + 1) \operatorname{and} \operatorname{bxor}(\operatorname{ldiag}(s + 2 - \operatorname{twopow}(k), s + 1), \operatorname{rdiag}(s - \operatorname{twopow}(k), s + 1))), \operatorname{rdiag}(s + 1 - \operatorname{twopow}(k), s + 1) \operatorname{and} \operatorname{bxor}(\operatorname{ldiag}(s + 2 - \operatorname{twopow}(k), s + 1), \operatorname{rdiag}(s, s + 1))))\).
\end{theorem}
\begin{proof}
The goal follows from \(\operatorname{frobdefectrowcenterslotoutsidecone}(s, k)\).
\end{proof}
\begin{theorem}
\label{ll:paper:frobdefect-row-center-slot-at-white-center-outside-cone}
For every natural number \(s\) and natural number \(k\), if \(\operatorname{twopow}(k) \le s\), then if \(\operatorname{center}(s + 1) = \operatorname{false}\), then \(\operatorname{bitfrobdefecteq}(k, \operatorname{wittrow}(s + 1), s + 1 + \operatorname{twopow}(k), \operatorname{bxor}(\operatorname{ldiag}(s + 1 - \operatorname{twopow}(k), s + 1) \operatorname{and} \operatorname{bxor}(\operatorname{maskparity}(s + 1, 0), \operatorname{rdiag}(s - \operatorname{twopow}(k), s + 1)), \operatorname{rdiag}(s + 1 - \operatorname{twopow}(k), s + 1) \operatorname{and} \operatorname{bxor}(\operatorname{ldiag}(s + 2 - \operatorname{twopow}(k), s + 1), \operatorname{maskparity}(s + 1, 0))))\).
\end{theorem}
\begin{proof}
The goal follows from \(\operatorname{frobdefectrowcenterslotatwhitecenteroutsidecone}(s, k)\).
\end{proof}
\begin{theorem}
\label{ll:paper:frobdefect-row-colm1-slot-outside-cone}
For every natural number \(d\) and natural number \(k\), if \(\operatorname{twopow}(k) \le d\), then \(\operatorname{bitfrobdefecteq}(k, \operatorname{wittrow}(d + 1), d + 2 + \operatorname{twopow}(k), \operatorname{bxor}(\operatorname{bxor}(\operatorname{ldiag}(d - \operatorname{twopow}(k), d + 1) \operatorname{and} \operatorname{bxor}(\operatorname{center}(d + 1), \operatorname{rdiag}(d + 1 - \operatorname{twopow}(k), d + 1)), \operatorname{ldiag}(d, d + 1) \operatorname{and} \operatorname{bxor}(\operatorname{ldiag}(d + 1 - \operatorname{twopow}(k), d + 1), \operatorname{rdiag}(d + 1 - \operatorname{twopow}(k), d + 1))), \operatorname{rdiag}(d + 2 - \operatorname{twopow}(k), d + 1) \operatorname{and} \operatorname{bxor}(\operatorname{ldiag}(d + 1 - \operatorname{twopow}(k), d + 1), \operatorname{center}(d + 1))))\).
\end{theorem}
\begin{proof}
The goal follows from \(\operatorname{frobdefectrowcolm1slotoutsidecone}(d, k)\).
\end{proof}
\begin{theorem}
\label{ll:paper:frobdefect-row-colm1-slot-at-white-center-outside-cone}
For every natural number \(d\) and natural number \(k\), if \(\operatorname{twopow}(k) \le d\), then if \(\operatorname{center}(d + 1) = \operatorname{false}\), then \(\operatorname{bitfrobdefecteq}(k, \operatorname{wittrow}(d + 1), d + 2 + \operatorname{twopow}(k), \operatorname{bxor}(\operatorname{bxor}(\operatorname{ldiag}(d - \operatorname{twopow}(k), d + 1) \operatorname{and} \operatorname{rdiag}(d + 1 - \operatorname{twopow}(k), d + 1), \operatorname{ldiag}(d, d + 1) \operatorname{and} \operatorname{bxor}(\operatorname{ldiag}(d + 1 - \operatorname{twopow}(k), d + 1), \operatorname{rdiag}(d + 1 - \operatorname{twopow}(k), d + 1))), \operatorname{rdiag}(d + 2 - \operatorname{twopow}(k), d + 1) \operatorname{and} \operatorname{ldiag}(d + 1 - \operatorname{twopow}(k), d + 1)))\).
\end{theorem}
\begin{proof}
The goal follows from \(\operatorname{frobdefectrowcolm1slotatwhitecenteroutsidecone}(d, k)\).
\end{proof}
\begin{theorem}
\label{ll:paper:frobdefect-row-col1-slot-outside-cone}
For every natural number \(s\) and natural number \(k\), if \(\operatorname{twopow}(k) \le s\), then \(\operatorname{bitfrobdefecteq}(k, \operatorname{wittrow}(s + 2), s + 1 + \operatorname{twopow}(k), \operatorname{bxor}(\operatorname{bxor}(\operatorname{ldiag}(s + 3 - \operatorname{twopow}(k), s + 2) \operatorname{and} \operatorname{bxor}(\operatorname{rdiag}(s, s + 2), \operatorname{rdiag}(s - \operatorname{twopow}(k), s + 2)), \operatorname{rdiag}(s + 1, s + 2) \operatorname{and} \operatorname{bxor}(\operatorname{rdiag}(s + \operatorname{twopow}(k), s + 2), \operatorname{rdiag}(s - \operatorname{twopow}(k), s + 2))), \operatorname{rdiag}(s + 1 - \operatorname{twopow}(k), s + 2) \operatorname{and} \operatorname{bxor}(\operatorname{rdiag}(s + \operatorname{twopow}(k), s + 2), \operatorname{rdiag}(s, s + 2))))\).
\end{theorem}
\begin{proof}
The goal follows from \(\operatorname{frobdefectrowcol1slotoutsidecone}(s, k)\).
\end{proof}
\begin{theorem}
\label{ll:paper:frobdefect-row-col1-slot-at-white-center-outside-cone}
For every natural number \(s\) and natural number \(k\), if \(\operatorname{twopow}(k) \le s\), then if \(\operatorname{center}(s + 2) = \operatorname{false}\), then \(\operatorname{bitfrobdefecteq}(k, \operatorname{wittrow}(s + 2), s + 1 + \operatorname{twopow}(k), \operatorname{bxor}(\operatorname{bxor}(\operatorname{ldiag}(s + 3 - \operatorname{twopow}(k), s + 2) \operatorname{and} \operatorname{bxor}(\operatorname{rdiag}(s, s + 2), \operatorname{rdiag}(s - \operatorname{twopow}(k), s + 2)), \operatorname{maskparity}(s + 2, 0) \operatorname{and} \operatorname{bxor}(\operatorname{rdiag}(s + \operatorname{twopow}(k), s + 2), \operatorname{rdiag}(s - \operatorname{twopow}(k), s + 2))), \operatorname{rdiag}(s + 1 - \operatorname{twopow}(k), s + 2) \operatorname{and} \operatorname{bxor}(\operatorname{rdiag}(s + \operatorname{twopow}(k), s + 2), \operatorname{rdiag}(s, s + 2))))\).
\end{theorem}
\begin{proof}
The goal follows from \(\operatorname{frobdefectrowcol1slotatwhitecenteroutsidecone}(s, k)\).
\end{proof}
\begin{theorem}
\label{ll:paper:periodic-winwitt-step-center-slot-outside-cone}
For every natural number \(T\) and natural number \(P\), if \(0 < P\), then if \(\operatorname{centerperiodic}(T, P)\), then for every natural number \(s\) and natural number \(k\), if \(T \le s + 1\), then if \(\operatorname{twopow}(k) \le s\), then \(\operatorname{wittresp}(s + 2, k) = \operatorname{bxor}(\operatorname{bxor}(\operatorname{wittbitbool}(\operatorname{wittphi}(\operatorname{wittwin}(s + 1, \operatorname{twopow}(k))), s + 2 + \operatorname{twopow}(k)), \operatorname{bxor}(\operatorname{bxor}(\operatorname{bxor}(\operatorname{ldiag}(s + 1 - \operatorname{twopow}(k), s + 1), \operatorname{center}(s + 1)), \operatorname{rdiag}(s + 1 - \operatorname{twopow}(k), s + 1)) \operatorname{and} \operatorname{rdiag}(s + \operatorname{twopow}(k), s + 1 + \operatorname{twopow}(k)), \operatorname{center}(s + 1 + \operatorname{lag}(P, k)) \operatorname{and} \operatorname{bxor}(\operatorname{bxor}(\operatorname{ldiag}(s + 2 - \operatorname{twopow}(k), s + 1), \operatorname{rdiag}(s, s + 1)), \operatorname{rdiag}(s - \operatorname{twopow}(k), s + 1)))), \operatorname{bxor}(\operatorname{bxor}(\operatorname{ldiag}(s + 1 - \operatorname{twopow}(k), s + 1) \operatorname{and} \operatorname{bxor}(\operatorname{rdiag}(s, s + 1), \operatorname{rdiag}(s - \operatorname{twopow}(k), s + 1)), \operatorname{center}(s + 1) \operatorname{and} \operatorname{bxor}(\operatorname{ldiag}(s + 2 - \operatorname{twopow}(k), s + 1), \operatorname{rdiag}(s - \operatorname{twopow}(k), s + 1))), \operatorname{rdiag}(s + 1 - \operatorname{twopow}(k), s + 1) \operatorname{and} \operatorname{bxor}(\operatorname{ldiag}(s + 2 - \operatorname{twopow}(k), s + 1), \operatorname{rdiag}(s, s + 1))))\).
\end{theorem}
\begin{proof}
Assume \(hP\), \(h\).
The goal follows from \(\operatorname{periodicwinwittstepcenterslotoutsidecone}(T, P, hP, h)\).
\end{proof}
\begin{theorem}
\label{ll:paper:periodic-winwitt-step-center-slot-at-white-late-center-outside-cone}
For every natural number \(T\) and natural number \(P\), if \(0 < P\), then if \(\operatorname{centerperiodic}(T, P)\), then for every natural number \(s\) and natural number \(k\), if \(T \le s + 1\), then if \(\operatorname{twopow}(k) \le s\), then if \(\operatorname{center}(s + 1 + \operatorname{lag}(P, k)) = \operatorname{false}\), then \(\operatorname{wittresp}(s + 2, k) = \operatorname{bxor}(\operatorname{bxor}(\operatorname{wittbitbool}(\operatorname{wittphi}(\operatorname{wittwin}(s + 1, \operatorname{twopow}(k))), s + 2 + \operatorname{twopow}(k)), \operatorname{wittresp}(s + 1, k) \operatorname{and} \operatorname{maskparity}(s + 1 + \operatorname{twopow}(k), 0)), \operatorname{bxor}(\operatorname{bxor}(\operatorname{ldiag}(s + 1 - \operatorname{twopow}(k), s + 1) \operatorname{and} \operatorname{bxor}(\operatorname{rdiag}(s, s + 1), \operatorname{rdiag}(s - \operatorname{twopow}(k), s + 1)), \operatorname{center}(s + 1) \operatorname{and} \operatorname{bxor}(\operatorname{ldiag}(s + 2 - \operatorname{twopow}(k), s + 1), \operatorname{rdiag}(s - \operatorname{twopow}(k), s + 1))), \operatorname{rdiag}(s + 1 - \operatorname{twopow}(k), s + 1) \operatorname{and} \operatorname{bxor}(\operatorname{ldiag}(s + 2 - \operatorname{twopow}(k), s + 1), \operatorname{rdiag}(s, s + 1))))\).
\end{theorem}
\begin{proof}
Assume \(hP\), \(h\).
The goal follows from \(\operatorname{periodicwinwittstepcenterslotatwhitelatecenteroutsidecone}(T, P, hP, h)\).
\end{proof}
\begin{theorem}
\label{ll:paper:periodic-winwitt-step-center-slot-at-both-white-centers-outside-cone}
For every natural number \(T\) and natural number \(P\), if \(0 < P\), then if \(\operatorname{centerperiodic}(T, P)\), then for every natural number \(s\) and natural number \(k\), if \(T \le s + 1\), then if \(\operatorname{twopow}(k) \le s\), then if \(\operatorname{center}(s + 1) = \operatorname{false}\), then if \(\operatorname{center}(s + 1 + \operatorname{lag}(P, k)) = \operatorname{false}\), then \(\operatorname{wittresp}(s + 2, k) = \operatorname{bxor}(\operatorname{bxor}(\operatorname{wittbitbool}(\operatorname{wittphi}(\operatorname{wittwin}(s + 1, \operatorname{twopow}(k))), s + 2 + \operatorname{twopow}(k)), \operatorname{bxor}(\operatorname{ldiag}(s + 1 - \operatorname{twopow}(k), s + 1), \operatorname{rdiag}(s + 1 - \operatorname{twopow}(k), s + 1)) \operatorname{and} \operatorname{maskparity}(s + 1 + \operatorname{twopow}(k), 0)), \operatorname{bxor}(\operatorname{ldiag}(s + 1 - \operatorname{twopow}(k), s + 1) \operatorname{and} \operatorname{bxor}(\operatorname{maskparity}(s + 1, 0), \operatorname{rdiag}(s - \operatorname{twopow}(k), s + 1)), \operatorname{rdiag}(s + 1 - \operatorname{twopow}(k), s + 1) \operatorname{and} \operatorname{bxor}(\operatorname{ldiag}(s + 2 - \operatorname{twopow}(k), s + 1), \operatorname{maskparity}(s + 1, 0))))\).
\end{theorem}
\begin{proof}
Assume \(hP\), \(h\).
The goal follows from \(\operatorname{periodicwinwittstepcenterslotatbothwhitecentersoutsidecone}(T, P, hP, h)\).
\end{proof}
\begin{theorem}
\label{ll:paper:periodic-winwitt-step-colm1-slot-outside-cone}
For every natural number \(T\) and natural number \(P\), if \(0 < P\), then if \(\operatorname{centerperiodic}(T, P)\), then for every natural number \(d\) and natural number \(k\), if \(T \le d + 1\), then if \(\operatorname{twopow}(k) \le d\), then \(\operatorname{bitwinwitteq}(d + 2, \operatorname{twopow}(k), d + 3 + \operatorname{twopow}(k), \operatorname{bxor}(\operatorname{bxor}(\operatorname{wittbitbool}(\operatorname{wittphi}(\operatorname{wittwin}(d + 1, \operatorname{twopow}(k))), d + 3 + \operatorname{twopow}(k)), \operatorname{bxor}(\operatorname{bxor}(\operatorname{bxor}(\operatorname{ldiag}(d - \operatorname{twopow}(k), d + 1), \operatorname{ldiag}(d, d + 1)), \operatorname{rdiag}(d + 2 - \operatorname{twopow}(k), d + 1)) \operatorname{and} \operatorname{center}(d + 1 + \operatorname{lag}(P, k)), \operatorname{ldiag}(d + \operatorname{twopow}(k), d + 1 + \operatorname{twopow}(k)) \operatorname{and} \operatorname{bxor}(\operatorname{bxor}(\operatorname{ldiag}(d + 1 - \operatorname{twopow}(k), d + 1), \operatorname{center}(d + 1)), \operatorname{rdiag}(d + 1 - \operatorname{twopow}(k), d + 1)))), \operatorname{bxor}(\operatorname{bxor}(\operatorname{ldiag}(d - \operatorname{twopow}(k), d + 1) \operatorname{and} \operatorname{bxor}(\operatorname{center}(d + 1), \operatorname{rdiag}(d + 1 - \operatorname{twopow}(k), d + 1)), \operatorname{ldiag}(d, d + 1) \operatorname{and} \operatorname{bxor}(\operatorname{ldiag}(d + 1 - \operatorname{twopow}(k), d + 1), \operatorname{rdiag}(d + 1 - \operatorname{twopow}(k), d + 1))), \operatorname{rdiag}(d + 2 - \operatorname{twopow}(k), d + 1) \operatorname{and} \operatorname{bxor}(\operatorname{ldiag}(d + 1 - \operatorname{twopow}(k), d + 1), \operatorname{center}(d + 1)))))\).
\end{theorem}
\begin{proof}
Assume \(hP\), \(h\).
The goal follows from \(\operatorname{periodicwinwittstepcolm1slotoutsidecone}(T, P, hP, h)\).
\end{proof}
\begin{theorem}
\label{ll:paper:periodic-winwitt-step-colm1-slot-at-both-white-centers-outside-cone}
For every natural number \(T\) and natural number \(P\), if \(0 < P\), then if \(\operatorname{centerperiodic}(T, P)\), then for every natural number \(d\) and natural number \(k\), if \(T \le d + 1\), then if \(\operatorname{twopow}(k) \le d\), then if \(\operatorname{center}(d + 1) = \operatorname{false}\), then if \(\operatorname{center}(d + 1 + \operatorname{lag}(P, k)) = \operatorname{false}\), then \(\operatorname{bitwinwitteq}(d + 2, \operatorname{twopow}(k), d + 3 + \operatorname{twopow}(k), \operatorname{bxor}(\operatorname{bxor}(\operatorname{wittbitbool}(\operatorname{wittphi}(\operatorname{wittwin}(d + 1, \operatorname{twopow}(k))), d + 3 + \operatorname{twopow}(k)), \operatorname{ldiag}(d + \operatorname{twopow}(k), d + 1 + \operatorname{twopow}(k)) \operatorname{and} \operatorname{bxor}(\operatorname{ldiag}(d + 1 - \operatorname{twopow}(k), d + 1), \operatorname{rdiag}(d + 1 - \operatorname{twopow}(k), d + 1))), \operatorname{bxor}(\operatorname{bxor}(\operatorname{ldiag}(d - \operatorname{twopow}(k), d + 1) \operatorname{and} \operatorname{rdiag}(d + 1 - \operatorname{twopow}(k), d + 1), \operatorname{ldiag}(d, d + 1) \operatorname{and} \operatorname{bxor}(\operatorname{ldiag}(d + 1 - \operatorname{twopow}(k), d + 1), \operatorname{rdiag}(d + 1 - \operatorname{twopow}(k), d + 1))), \operatorname{rdiag}(d + 2 - \operatorname{twopow}(k), d + 1) \operatorname{and} \operatorname{ldiag}(d + 1 - \operatorname{twopow}(k), d + 1))))\).
\end{theorem}
\begin{proof}
Assume \(hP\), \(h\).
The goal follows from \(\operatorname{periodicwinwittstepcolm1slotatbothwhitecentersoutsidecone}(T, P, hP, h)\).
\end{proof}
\begin{theorem}
\label{ll:paper:periodic-winwitt-step-col1-slot-outside-cone}
For every natural number \(s\) and natural number \(k\), if \(\operatorname{twopow}(k) \le s\), then \(\operatorname{bitwinwitteq}(s + 3, \operatorname{twopow}(k), s + 2 + \operatorname{twopow}(k), \operatorname{bxor}(\operatorname{bxor}(\operatorname{wittbitbool}(\operatorname{wittphi}(\operatorname{wittwin}(s + 2, \operatorname{twopow}(k))), s + 2 + \operatorname{twopow}(k)), \operatorname{wittbitbool}(\operatorname{crosscarry}(\operatorname{wittscale}(k, \operatorname{wittrow}(s + 2)), \operatorname{wittwin}(s + 2, \operatorname{twopow}(k))), s + 1 + \operatorname{twopow}(k))), \operatorname{bxor}(\operatorname{bxor}(\operatorname{ldiag}(s + 3 - \operatorname{twopow}(k), s + 2) \operatorname{and} \operatorname{bxor}(\operatorname{rdiag}(s, s + 2), \operatorname{rdiag}(s - \operatorname{twopow}(k), s + 2)), \operatorname{rdiag}(s + 1, s + 2) \operatorname{and} \operatorname{bxor}(\operatorname{rdiag}(s + \operatorname{twopow}(k), s + 2), \operatorname{rdiag}(s - \operatorname{twopow}(k), s + 2))), \operatorname{rdiag}(s + 1 - \operatorname{twopow}(k), s + 2) \operatorname{and} \operatorname{bxor}(\operatorname{rdiag}(s + \operatorname{twopow}(k), s + 2), \operatorname{rdiag}(s, s + 2)))))\).
\end{theorem}
\begin{proof}
The goal follows from \(\operatorname{periodicwinwittstepcol1slotoutsidecone}(s, k)\).
\end{proof}
\begin{theorem}
\label{ll:paper:periodic-winwitt-step-center-slot-shift-mul-at-white-late-center}
For every natural number \(T\) and natural number \(P\), if \(0 < P\), then if \(\operatorname{centerperiodic}(T, P)\), then for every natural number \(m\) and natural number \(s\) and natural number \(k\), if \(T \le s + 1\), then if \(\operatorname{center}(s + 1 + \operatorname{lag}(P, k)) = \operatorname{false}\), then \(\operatorname{bxor}(\operatorname{bxor}(\operatorname{bxor}(\operatorname{wittresp}(s + m \cdot P + 2, k), \operatorname{wittbitbool}(\operatorname{wittphi}(\operatorname{wittwin}(s + m \cdot P + 1, \operatorname{twopow}(k))), s + m \cdot P + 2 + \operatorname{twopow}(k))), \operatorname{wittbitbool}(\operatorname{frobdefect}(k, \operatorname{wittrow}(s + m \cdot P + 1)), s + m \cdot P + 1 + \operatorname{twopow}(k))), \operatorname{bxor}(\operatorname{bxor}(\operatorname{wittresp}(s + 2, k), \operatorname{wittbitbool}(\operatorname{wittphi}(\operatorname{wittwin}(s + 1, \operatorname{twopow}(k))), s + 2 + \operatorname{twopow}(k))), \operatorname{wittbitbool}(\operatorname{frobdefect}(k, \operatorname{wittrow}(s + 1)), s + 1 + \operatorname{twopow}(k)))) = \operatorname{bxor}(\operatorname{bxor}(\operatorname{wittresp}(s + m \cdot P + 1, k), \operatorname{wittresp}(s + 1, k)) \operatorname{and} \operatorname{maskparity}(s + m \cdot P + 1 + \operatorname{twopow}(k), 0), \operatorname{wittresp}(s + 1, k) \operatorname{and} \operatorname{bxor}(\operatorname{maskparity}(s + m \cdot P + 1 + \operatorname{twopow}(k), 0), \operatorname{maskparity}(s + 1 + \operatorname{twopow}(k), 0)))\).
\end{theorem}
\begin{proof}
Assume \(hP\), \(h\).
The goal follows from \(\operatorname{periodicwinwittstepcenterslotshiftmulatwhitelatecenter}(T, P, hP, h)\).
\end{proof}
\begin{theorem}
\label{ll:paper:frobdefect-xor}
For every natural number \(k\) and natural number \(a\) and natural number \(b\), \(\operatorname{frobdefect}(k, \operatorname{wittxor}(a, b)) = \operatorname{wittxor}(\operatorname{wittxor}(\operatorname{frobdefect}(k, a), \operatorname{frobdefect}(k, b)), \operatorname{wittxor}(\operatorname{crosscarry}(\operatorname{wittscale}(k, a), \operatorname{wittscale}(k, b)), \operatorname{wittscale}(k, \operatorname{crosscarry}(a, b))))\).
\end{theorem}
\begin{proof}
The goal follows from \(\operatorname{frobdefectxor}(k, a, b)\).
\end{proof}
\begin{theorem}
\label{ll:paper:frobdefect-phiscale-row-cone-eq-zero}
For every natural number \(t0\) and natural number \(k\), if \(2 \cdot t0 + 1 < \operatorname{twopow}(k)\), then \(\operatorname{frobdefect}(k, \operatorname{wittscale}(k, \operatorname{wittrow}(t0))) = 0\).
\end{theorem}
\begin{proof}
The goal follows from \(\operatorname{frobdefectphiscalerowconeeqzero}(t0, k)\).
\end{proof}
\begin{theorem}
\label{ll:paper:frobdefect-row-late-from-cone}
For every natural number \(t0\) and natural number \(k\), if \(2 \cdot t0 + 1 < \operatorname{twopow}(k)\), then \(\operatorname{frobdefect}(k, \operatorname{wittrow}(t0 + \operatorname{twopow}(k))) = \operatorname{wittxor}(\operatorname{frobdefect}(k, \operatorname{wittwin}(t0, \operatorname{twopow}(k))), \operatorname{wittxor}(\operatorname{crosscarry}(\operatorname{wittscale}(\operatorname{succ}(k), \operatorname{wittrow}(t0)), \operatorname{wittscale}(k, \operatorname{wittwin}(t0, \operatorname{twopow}(k)))), \operatorname{wittscale}(k, \operatorname{crosscarry}(\operatorname{wittscale}(k, \operatorname{wittrow}(t0)), \operatorname{wittwin}(t0, \operatorname{twopow}(k))))))\).
\end{theorem}
\begin{proof}
The goal follows from \(\operatorname{frobdefectrowlatefromcone}(t0, k)\).
\end{proof}
\begin{theorem}
\label{ll:paper:frobdefect-row-late-center-slot-rperiodic}
For every natural number \(d\) and natural number \(k\), if \(d + 1 \le k\), then \(\operatorname{bitfrobdefecteq}(k, \operatorname{wittrow}(d + 1 + \operatorname{twopow}(k)), d + 1 + \operatorname{twopow}(k) + \operatorname{twopow}(k), \operatorname{bxor}(\operatorname{bxor}(\operatorname{ldiag}(d + 1, d + 1 + \operatorname{twopow}(k)) \operatorname{and} \operatorname{bxor}(\operatorname{rdiag}(d + \operatorname{twopow}(k), d + 1 + \operatorname{twopow}(k)), \operatorname{rdiag}(d, d + 1)), \operatorname{center}(d + 1 + \operatorname{twopow}(k)) \operatorname{and} \operatorname{bxor}(\operatorname{ldiag}(d + 2, d + 1 + \operatorname{twopow}(k)), \operatorname{rdiag}(d, d + 1))), \operatorname{center}(d + 1) \operatorname{and} \operatorname{bxor}(\operatorname{ldiag}(d + 2, d + 1 + \operatorname{twopow}(k)), \operatorname{rdiag}(d + \operatorname{twopow}(k), d + 1 + \operatorname{twopow}(k)))))\).
\end{theorem}
\begin{proof}
The goal follows from \(\operatorname{frobdefectrowlatecenterslotrperiodic}(d, k)\).
\end{proof}
\begin{theorem}
\label{ll:paper:frobdefect-row-late-center-slot-at-both-white-centers-eq-zcenter-and-maskparity-jump}
For every natural number \(d\) and natural number \(k\), if \(d + 1 \le k\), then if \(\operatorname{settledat}(d + 1, k)\), then if \(\operatorname{center}(d + 1) = \operatorname{false}\), then if \(\operatorname{center}(d + 1 + \operatorname{twopow}(k)) = \operatorname{false}\), then \(\operatorname{zcenterandeq}(d + 1, \operatorname{bxor}(\operatorname{maskparity}(d + 1 + \operatorname{twopow}(k), 0), \operatorname{maskparity}(d + 1, 0)), \operatorname{wittbitbool}(\operatorname{frobdefect}(k, \operatorname{wittrow}(d + 1 + \operatorname{twopow}(k))), d + 1 + \operatorname{twopow}(k) + \operatorname{twopow}(k)))\).
\end{theorem}
\begin{proof}
The goal follows from \(\operatorname{frobdefectrowlatecenterslotatbothwhitecenterseqzcenterandmaskparityjump}(d, k)\).
\end{proof}
\begin{theorem}
\label{ll:paper:frobdefect-row-late-colm1-slot-rperiodic}
For every natural number \(d\) and natural number \(k\), if \(d + 2 \le k\), then \(\operatorname{bitfrobdefecteq}(k, \operatorname{wittrow}(d + 1 + \operatorname{twopow}(k)), d + 2 + \operatorname{twopow}(k) + \operatorname{twopow}(k), \operatorname{bxor}(\operatorname{bxor}(\operatorname{ldiag}(d, d + 1 + \operatorname{twopow}(k)) \operatorname{and} \operatorname{bxor}(\operatorname{center}(d + 1 + \operatorname{twopow}(k)), \operatorname{center}(d + 1)), \operatorname{ldiag}(d + \operatorname{twopow}(k), d + 1 + \operatorname{twopow}(k)) \operatorname{and} \operatorname{bxor}(\operatorname{ldiag}(d + 1, d + 1 + \operatorname{twopow}(k)), \operatorname{center}(d + 1))), \operatorname{ldiag}(d, d + 1) \operatorname{and} \operatorname{bxor}(\operatorname{ldiag}(d + 1, d + 1 + \operatorname{twopow}(k)), \operatorname{center}(d + 1 + \operatorname{twopow}(k)))))\).
\end{theorem}
\begin{proof}
The goal follows from \(\operatorname{frobdefectrowlatecolm1slotrperiodic}(d, k)\).
\end{proof}
\begin{theorem}
\label{ll:paper:frobdefect-row-late-colm1-slot-of-in-phase-eq-colm1-jump-and-damage}
For every natural number \(d\) and natural number \(k\), if \(d + 2 \le k\), then if \(\operatorname{settledat}(d + 1, k)\), then if \(\operatorname{center}(d + 1 + \operatorname{twopow}(k)) = \operatorname{center}(d + 1)\), then \(\operatorname{zcenterandeq}(d + 1, \operatorname{bxor}(\operatorname{ldiag}(d + \operatorname{twopow}(k), d + 1 + \operatorname{twopow}(k)), \operatorname{ldiag}(d, d + 1)), \operatorname{bxor}(\operatorname{wittbitbool}(\operatorname{frobdefect}(k, \operatorname{wittrow}(d + 1 + \operatorname{twopow}(k))), d + 2 + \operatorname{twopow}(k) + \operatorname{twopow}(k)), \operatorname{center}(d + 1) \operatorname{and} \operatorname{bxor}(\operatorname{ldiag}(d + \operatorname{twopow}(k), d + 1 + \operatorname{twopow}(k)), \operatorname{ldiag}(d, d + 1))))\).
\end{theorem}
\begin{proof}
The goal follows from \(\operatorname{frobdefectrowlatecolm1slotofinphaseeqcolm1jumpanddamage}(d, k)\).
\end{proof}
\begin{theorem}
\label{ll:paper:frobdefect-row-late-colm1-xor-center-slot-at-both-white-centers}
For every natural number \(d\) and natural number \(k\), if \(d + 2 \le k\), then if \(\operatorname{settledat}(d + 1, k)\), then if \(\operatorname{center}(d + 1) = \operatorname{false}\), then if \(\operatorname{center}(d + 1 + \operatorname{twopow}(k)) = \operatorname{false}\), then \(\operatorname{zcenterandeq}(d + 1, \operatorname{bxor}(\operatorname{center}(d + 2 + \operatorname{twopow}(k)), \operatorname{center}(d + 2)), \operatorname{bxor}(\operatorname{wittbitbool}(\operatorname{frobdefect}(k, \operatorname{wittrow}(d + 1 + \operatorname{twopow}(k))), d + 2 + \operatorname{twopow}(k) + \operatorname{twopow}(k)), \operatorname{wittbitbool}(\operatorname{frobdefect}(k, \operatorname{wittrow}(d + 1 + \operatorname{twopow}(k))), d + 1 + \operatorname{twopow}(k) + \operatorname{twopow}(k))))\).
\end{theorem}
\begin{proof}
The goal follows from \(\operatorname{frobdefectrowlatecolm1xorcenterslotatbothwhitecenters}(d, k)\).
\end{proof}
\begin{theorem}
\label{ll:paper:periodic-winwitt-step-late-center-slot-at-white-centers-eq-zcenter-and-maskparity-triality}
For every natural number \(T\) and natural number \(P\), if \(0 < P\), then if \(\operatorname{centerperiodic}(T, P)\), then for every natural number \(d\) and natural number \(k\), if \(T \le d + 1\), then if \(d + 1 \le k\), then if \(\operatorname{settledat}(d + 1, k)\), then if \(\operatorname{center}(d + 1) = \operatorname{false}\), then if \(\operatorname{center}(d + 1 + \operatorname{lag}(P, k)) = \operatorname{false}\), then if \(\operatorname{center}(d + 1 + \operatorname{lag}(P, \operatorname{succ}(k))) = \operatorname{false}\), then \(\operatorname{zcenterandeq}(d + 1, \operatorname{bxor}(\operatorname{maskparity}(d + 1 + \operatorname{twopow}(\operatorname{succ}(k)), 0), \operatorname{bxor}(\operatorname{maskparity}(d + 1 + \operatorname{twopow}(k), 0), \operatorname{maskparity}(d + 1, 0))), \operatorname{bxor}(\operatorname{wittresp}(d + 2 + \operatorname{twopow}(k), k), \operatorname{wittbitbool}(\operatorname{wittphi}(\operatorname{wittwin}(d + 1 + \operatorname{twopow}(k), \operatorname{twopow}(k))), d + 2 + \operatorname{twopow}(k) + \operatorname{twopow}(k))))\).
\end{theorem}
\begin{proof}
Assume \(hP\), \(h\).
The goal follows from \(\operatorname{periodicwinwittsteplatecenterslotatwhitecenterseqzcenterandmaskparitytriality}(T, P, hP, h)\).
\end{proof}
\begin{theorem}
\label{ll:paper:periodic-winwitt-step-late-colm1-slot-at-white-centers-eq-zcenter-and-colm1-triality}
For every natural number \(T\) and natural number \(P\), if \(0 < P\), then if \(\operatorname{centerperiodic}(T, P)\), then for every natural number \(d\) and natural number \(k\), if \(T \le d + 1\), then if \(d + 2 \le k\), then if \(\operatorname{settledat}(d + 1, k)\), then if \(\operatorname{center}(d + 1) = \operatorname{false}\), then if \(\operatorname{center}(d + 1 + \operatorname{lag}(P, k)) = \operatorname{false}\), then if \(\operatorname{center}(d + 1 + \operatorname{lag}(P, \operatorname{succ}(k))) = \operatorname{false}\), then \(\operatorname{zcenterandeq}(d + 1, \operatorname{bxor}(\operatorname{ldiag}(d + \operatorname{twopow}(\operatorname{succ}(k)), d + 1 + \operatorname{twopow}(\operatorname{succ}(k))), \operatorname{bxor}(\operatorname{ldiag}(d + \operatorname{twopow}(k), d + 1 + \operatorname{twopow}(k)), \operatorname{ldiag}(d, d + 1))), \operatorname{bxor}(\operatorname{wittbitbool}(\operatorname{wittwin}(d + 2 + \operatorname{twopow}(k), \operatorname{twopow}(k)), d + 3 + \operatorname{twopow}(k) + \operatorname{twopow}(k)), \operatorname{wittbitbool}(\operatorname{wittphi}(\operatorname{wittwin}(d + 1 + \operatorname{twopow}(k), \operatorname{twopow}(k))), d + 3 + \operatorname{twopow}(k) + \operatorname{twopow}(k))))\).
\end{theorem}
\begin{proof}
Assume \(hP\), \(h\).
The goal follows from \(\operatorname{periodicwinwittsteplatecolm1slotatwhitecenterseqzcenterandcolm1triality}(T, P, hP, h)\).
\end{proof}
\begin{theorem}
\label{ll:paper:wittresp-late-eq-zcenter-xor-center-triality}
For every natural number \(d\) and natural number \(k\), if \(d + 1 \le k\), then if \(\operatorname{settledat}(d + 1, k)\), then \(\operatorname{zcentereq}(d + 1, \operatorname{bxor}(\operatorname{wittresp}(d + 1 + \operatorname{twopow}(k), k), \operatorname{bxor}(\operatorname{center}(d + 1 + \operatorname{twopow}(\operatorname{succ}(k))), \operatorname{bxor}(\operatorname{center}(d + 1 + \operatorname{twopow}(k)), \operatorname{center}(d + 1)))))\).
\end{theorem}
\begin{proof}
The goal follows from \(\operatorname{wittresplateeqzcenterxorcentertriality}(d, k)\).
\end{proof}
\begin{theorem}
\label{ll:paper:wittresp-late-at-white-centers-eq-zcenter}
For every natural number \(d\) and natural number \(k\), if \(d + 1 \le k\), then if \(\operatorname{settledat}(d + 1, k)\), then if \(\operatorname{center}(d + 1) = \operatorname{false}\), then if \(\operatorname{center}(d + 1 + \operatorname{twopow}(k)) = \operatorname{false}\), then if \(\operatorname{center}(d + 1 + \operatorname{twopow}(\operatorname{succ}(k))) = \operatorname{false}\), then \(\operatorname{zcentereq}(d + 1, \operatorname{wittresp}(d + 1 + \operatorname{twopow}(k), k))\).
\end{theorem}
\begin{proof}
The goal follows from \(\operatorname{wittresplateatwhitecenterseqzcenter}(d, k)\).
\end{proof}
\begin{theorem}
\label{ll:paper:periodic-wittresp-late-shift-mul-eq-zcenter-xor}
For every natural number \(T\) and natural number \(P\), if \(0 < P\), then if \(\operatorname{centerperiodic}(T, P)\), then for every natural number \(m\) and natural number \(d\) and natural number \(k\), if \(T \le d + 1\), then if \(d + 1 \le k\), then if \(\operatorname{settledat}(d + 1, k)\), then \(\operatorname{zcentereq}(d + 1, \operatorname{bxor}(\operatorname{bxor}(\operatorname{wittresp}(d + m \cdot P + 1 + \operatorname{twopow}(k), k), \operatorname{wittresp}(d + 1 + \operatorname{twopow}(k), k)), \operatorname{bxor}(\operatorname{ldiag}(d + m \cdot P + 1, d + m \cdot P + 1 + \operatorname{twopow}(k)), \operatorname{bxor}(\operatorname{rdiag}(d + m \cdot P + 1, d + m \cdot P + 1 + \operatorname{twopow}(k)), \operatorname{center}(d + 1)))))\).
\end{theorem}
\begin{proof}
Assume \(hP\), \(h\).
The goal follows from \(\operatorname{periodicwittresplateshiftmuleqzcenterxor}(T, P, hP, h)\).
\end{proof}
\begin{theorem}
\label{ll:paper:periodic-crosscarry-late-center-bit-shift-mul-at-white-centers}
For every natural number \(T\) and natural number \(P\), if \(0 < P\), then if \(\operatorname{centerperiodic}(T, P)\), then for every natural number \(m\) and natural number \(d\) and natural number \(k\), if \(T \le d + 1\), then if \(d + 1 \le k\), then if \(\operatorname{settledat}(d + 1, k)\), then if \(\operatorname{center}(d + 1) = \operatorname{false}\), then if \(\operatorname{center}(d + 1 + \operatorname{lag}(P, k)) = \operatorname{false}\), then if \(\operatorname{center}(d + 1 + \operatorname{lag}(P, \operatorname{succ}(k))) = \operatorname{false}\), then \(\operatorname{zcenterandeq}(d + 1, \operatorname{maskparity}(d + 1 + \operatorname{twopow}(\operatorname{succ}(k)), 0), \operatorname{bxor}(\operatorname{bxor}(\operatorname{wittbitbool}(\operatorname{crosscarry}(\operatorname{wittscale}(k, \operatorname{wittrow}(d + m \cdot P + 1 + \operatorname{twopow}(k))), \operatorname{wittwin}(d + m \cdot P + 1 + \operatorname{twopow}(k), \operatorname{twopow}(k))), d + m \cdot P + 1 + \operatorname{twopow}(k) + \operatorname{twopow}(k)), \operatorname{wittbitbool}(\operatorname{crosscarry}(\operatorname{wittscale}(k, \operatorname{wittrow}(d + 1 + \operatorname{twopow}(k))), \operatorname{wittwin}(d + 1 + \operatorname{twopow}(k), \operatorname{twopow}(k))), d + 1 + \operatorname{twopow}(k) + \operatorname{twopow}(k))), \operatorname{wittresp}(d + m \cdot P + 1 + \operatorname{twopow}(k), k) \operatorname{and} \operatorname{maskparity}(d + m \cdot P + 1 + \operatorname{twopow}(\operatorname{succ}(k)), 0)))\).
\end{theorem}
\begin{proof}
Assume \(hP\), \(h\).
The goal follows from \(\operatorname{periodiccrosscarrylatecenterbitshiftmulatwhitecenters}(T, P, hP, h)\).
\end{proof}
\begin{theorem}
\label{ll:paper:periodic-crosscarry-late-colm1-shift-mul-at-white-centers}
For every natural number \(T\) and natural number \(P\), if \(0 < P\), then if \(\operatorname{centerperiodic}(T, P)\), then for every natural number \(m\) and natural number \(d\) and natural number \(k\), if \(T \le d + 1\), then if \(d + 1 \le k\), then if \(\operatorname{settledat}(d + 1, k)\), then if \(\operatorname{center}(d + 1) = \operatorname{false}\), then if \(\operatorname{center}(d + 1 + \operatorname{lag}(P, k)) = \operatorname{false}\), then if \(\operatorname{center}(d + 1 + \operatorname{lag}(P, \operatorname{succ}(k))) = \operatorname{false}\), then \(\operatorname{zcenterandeq}(d + 1, \operatorname{ldiag}(d + \operatorname{twopow}(\operatorname{succ}(k)), d + 1 + \operatorname{twopow}(\operatorname{succ}(k))), \operatorname{bxor}(\operatorname{bxor}(\operatorname{wittbitbool}(\operatorname{crosscarry}(\operatorname{wittscale}(k, \operatorname{wittrow}(d + m \cdot P + 1 + \operatorname{twopow}(k))), \operatorname{wittwin}(d + m \cdot P + 1 + \operatorname{twopow}(k), \operatorname{twopow}(k))), d + m \cdot P + 2 + \operatorname{twopow}(k) + \operatorname{twopow}(k)), \operatorname{wittbitbool}(\operatorname{crosscarry}(\operatorname{wittscale}(k, \operatorname{wittrow}(d + 1 + \operatorname{twopow}(k))), \operatorname{wittwin}(d + 1 + \operatorname{twopow}(k), \operatorname{twopow}(k))), d + 2 + \operatorname{twopow}(k) + \operatorname{twopow}(k))), \operatorname{bxor}(\operatorname{bxor}(\operatorname{maskparity}(d + m \cdot P + 1 + \operatorname{twopow}(\operatorname{succ}(k)), 0), \operatorname{maskparity}(d + 1 + \operatorname{twopow}(\operatorname{succ}(k)), 0)) \operatorname{and} \operatorname{wittresp}(d + m \cdot P + 1 + \operatorname{twopow}(k), k), \operatorname{ldiag}(d + \operatorname{twopow}(\operatorname{succ}(k)), d + 1 + \operatorname{twopow}(\operatorname{succ}(k))) \operatorname{and} \operatorname{wittresp}(d + m \cdot P + 1 + \operatorname{twopow}(k), k))))\).
\end{theorem}
\begin{proof}
Assume \(hP\), \(h\).
The goal follows from \(\operatorname{periodiccrosscarrylatecolm1shiftmulatwhitecenters}(T, P, hP, h)\).
\end{proof}
\begin{theorem}
\label{ll:paper:periodic-winwitt-step-late-center-slot-shift-mul-at-white-centers}
For every natural number \(T\) and natural number \(P\), if \(0 < P\), then if \(\operatorname{centerperiodic}(T, P)\), then for every natural number \(m\) and natural number \(d\) and natural number \(k\), if \(T \le d + 1\), then if \(d + 1 \le k\), then if \(\operatorname{settledat}(d + 1, k)\), then if \(\operatorname{center}(d + 1) = \operatorname{false}\), then if \(\operatorname{center}(d + 1 + \operatorname{lag}(P, k)) = \operatorname{false}\), then if \(\operatorname{center}(d + 1 + \operatorname{lag}(P, \operatorname{succ}(k))) = \operatorname{false}\), then \(\operatorname{zcenterandeq}(d + 1, \operatorname{bxor}(\operatorname{maskparity}(d + 1 + \operatorname{twopow}(\operatorname{succ}(k)), 0), \operatorname{bxor}(\operatorname{maskparity}(d + 1 + \operatorname{twopow}(k), 0), \operatorname{maskparity}(d + 1, 0))), \operatorname{bxor}(\operatorname{bxor}(\operatorname{wittresp}(d + m \cdot P + 2 + \operatorname{twopow}(k), k), \operatorname{wittresp}(d + 2 + \operatorname{twopow}(k), k)), \operatorname{bxor}(\operatorname{bxor}(\operatorname{wittbitbool}(\operatorname{wittphi}(\operatorname{wittwin}(d + m \cdot P + 1 + \operatorname{twopow}(k), \operatorname{twopow}(k))), d + m \cdot P + 2 + \operatorname{twopow}(k) + \operatorname{twopow}(k)), \operatorname{wittbitbool}(\operatorname{wittphi}(\operatorname{wittwin}(d + 1 + \operatorname{twopow}(k), \operatorname{twopow}(k))), d + 2 + \operatorname{twopow}(k) + \operatorname{twopow}(k))), \operatorname{bxor}(\operatorname{wittresp}(d + m \cdot P + 1 + \operatorname{twopow}(k), k) \operatorname{and} \operatorname{maskparity}(d + m \cdot P + 1 + \operatorname{twopow}(\operatorname{succ}(k)), 0), \operatorname{wittbitbool}(\operatorname{frobdefect}(k, \operatorname{wittrow}(d + m \cdot P + 1 + \operatorname{twopow}(k))), d + m \cdot P + 1 + \operatorname{twopow}(k) + \operatorname{twopow}(k))))))\).
\end{theorem}
\begin{proof}
Assume \(hP\), \(h\).
The goal follows from \(\operatorname{periodicwinwittsteplatecenterslotshiftmulatwhitecenters}(T, P, hP, h)\).
\end{proof}
\begin{theorem}
\label{ll:paper:rdiag-restart}
For every natural number \(k\) and natural number \(j\) and natural number \(u\) and natural number \(s\), if \(j \le k\), then \(\operatorname{rdiag}(j, s + u \cdot \operatorname{twopow}(k)) = \operatorname{rdiag}(j, s)\).
\end{theorem}
\begin{proof}
The goal follows from \(\operatorname{rdiagrestart}(k, j, u, s)\).
\end{proof}
\begin{theorem}
\label{ll:paper:rdiag-restart-of-threshold}
For every natural number \(k\) and natural number \(J\) and natural number \(j\) and natural number \(u\) and natural number \(s\), if \(\operatorname{rightthreshold}(k, J)\), then if \(j \le J\), then \(\operatorname{rdiag}(j, s + u \cdot \operatorname{twopow}(k)) = \operatorname{rdiag}(j, s)\).
\end{theorem}
\begin{proof}
The goal follows from \(\operatorname{rdiagrestartofthreshold}(k, J, j, u, s)\).
\end{proof}
\begin{theorem}
\label{ll:paper:center-restart-rdiag}
For every natural number \(k\) and natural number \(s\) and natural number \(u\), if \(s \le k\), then \(\operatorname{rdiag}(s, s + u \cdot \operatorname{twopow}(k)) = \operatorname{center}(s)\).
\end{theorem}
\begin{proof}
The goal follows from \(\operatorname{centerrestartrdiag}(k, s, u)\).
\end{proof}
\begin{theorem}
\label{ll:paper:frobdefect-row-lt-twopow}
For every natural number \(t0\) and natural number \(k\) and natural number \(j\), if \(j < \operatorname{twopow}(k)\), then \(\operatorname{bitfrobdefecteq}(k, \operatorname{wittrow}(t0), j, \operatorname{false})\).
\end{theorem}
\begin{proof}
The goal follows from \(\operatorname{frobdefectrowlttwopow}(t0, k, j)\).
\end{proof}
\begin{theorem}
\label{ll:paper:frobdefect-row-sector-one}
For every natural number \(t0\) and natural number \(k\) and natural number \(r\), if \(r + 1 < \operatorname{twopow}(k)\), then \(\operatorname{bitfrobdefecteq}(k, \operatorname{wittrow}(t0), r + 1 + \operatorname{twopow}(k), \operatorname{bxor}(\operatorname{rdiag}(r + 1 + \operatorname{twopow}(k), t0) \operatorname{and} \operatorname{rdiag}(r, t0), \operatorname{rdiag}(r + 1, t0) \operatorname{and} \operatorname{rdiag}(r + \operatorname{twopow}(k), t0)))\).
\end{theorem}
\begin{proof}
The goal follows from \(\operatorname{frobdefectrowsectorone}(t0, k, r)\).
\end{proof}
\begin{theorem}
\label{ll:paper:frobdefect-row-sector-one-restart}
For every natural number \(d\) and natural number \(k\) and natural number \(u\), if \(d + 1 \le k\), then \(\operatorname{bitfrobdefecteq}(k, \operatorname{wittrow}(d + 1 + u \cdot \operatorname{twopow}(k)), d + 1 + \operatorname{twopow}(k), \operatorname{bxor}(\operatorname{rdiag}(d + 1 + \operatorname{twopow}(k), d + 1 + u \cdot \operatorname{twopow}(k)) \operatorname{and} \operatorname{rdiag}(d, d + 1), \operatorname{center}(d + 1) \operatorname{and} \operatorname{rdiag}(d + \operatorname{twopow}(k), d + 1 + u \cdot \operatorname{twopow}(k))))\).
\end{theorem}
\begin{proof}
The goal follows from \(\operatorname{frobdefectrowsectorerestart}(d, k, u)\).
\end{proof}
\begin{theorem}
\label{ll:paper:frobdefect-row-sector-two-restart}
For every natural number \(d\) and natural number \(k\) and natural number \(u\), if \(d + 1 \le k\), then \(\operatorname{bitfrobdefecteq}(k, \operatorname{wittrow}(d + 1 + u \cdot \operatorname{twopow}(k)), d + 1 + \operatorname{twopow}(k) + \operatorname{twopow}(k), \operatorname{bxor}(\operatorname{bxor}(\operatorname{rdiag}(d + 1 + \operatorname{twopow}(\operatorname{succ}(k)), d + 1 + u \cdot \operatorname{twopow}(k)) \operatorname{and} \operatorname{bxor}(\operatorname{rdiag}(d + \operatorname{twopow}(k), d + 1 + u \cdot \operatorname{twopow}(k)), \operatorname{rdiag}(d, d + 1)), \operatorname{rdiag}(d + 1 + \operatorname{twopow}(k), d + 1 + u \cdot \operatorname{twopow}(k)) \operatorname{and} \operatorname{bxor}(\operatorname{rdiag}(d + \operatorname{twopow}(\operatorname{succ}(k)), d + 1 + u \cdot \operatorname{twopow}(k)), \operatorname{rdiag}(d, d + 1))), \operatorname{center}(d + 1) \operatorname{and} \operatorname{bxor}(\operatorname{rdiag}(d + \operatorname{twopow}(\operatorname{succ}(k)), d + 1 + u \cdot \operatorname{twopow}(k)), \operatorname{rdiag}(d + \operatorname{twopow}(k), d + 1 + u \cdot \operatorname{twopow}(k)))))\).
\end{theorem}
\begin{proof}
The goal follows from \(\operatorname{frobdefectrowsectortworestart}(d, k, u)\).
\end{proof}
\begin{theorem}
\label{ll:paper:frobdefect-winwitt-sector-two-wronskian}
For every natural number \(t0\) and natural number \(k\) and natural number \(r\), if \(r + 1 \le k\), then \(\operatorname{bitfrobdefecteq}(k, \operatorname{wittwin}(t0, \operatorname{twopow}(k)), r + 1 + \operatorname{twopow}(\operatorname{succ}(k)), \operatorname{bxor}(\operatorname{wittbitbool}(\operatorname{wittwin}(t0, \operatorname{twopow}(k)), r + 1 + \operatorname{twopow}(\operatorname{succ}(k))) \operatorname{and} \operatorname{wittbitbool}(\operatorname{wittwin}(t0, \operatorname{twopow}(k)), r + \operatorname{twopow}(k)), \operatorname{wittbitbool}(\operatorname{wittwin}(t0, \operatorname{twopow}(k)), r + 1 + \operatorname{twopow}(k)) \operatorname{and} \operatorname{wittbitbool}(\operatorname{wittwin}(t0, \operatorname{twopow}(k)), r + \operatorname{twopow}(\operatorname{succ}(k)))))\).
\end{theorem}
\begin{proof}
The goal follows from \(\operatorname{frobdefectwinwittsectortwowronskian}(t0, k, r)\).
\end{proof}
\begin{theorem}
\label{ll:paper:crosscarry-scale-row-winwitt-sector-two-cone-eq}
For every natural number \(t0\) and natural number \(k\) and natural number \(r\), if \(2 \cdot t0 + 1 < \operatorname{twopow}(k)\), then if \(r + 1 \le k\), then \(\operatorname{wittbitbool}(\operatorname{crosscarry}(\operatorname{wittscale}(\operatorname{succ}(k), \operatorname{wittrow}(t0)), \operatorname{wittscale}(k, \operatorname{wittwin}(t0, \operatorname{twopow}(k)))), r + 1 + \operatorname{twopow}(\operatorname{succ}(k))) = \operatorname{wittbitbool}(\operatorname{wittscale}(k, \operatorname{crosscarry}(\operatorname{wittscale}(k, \operatorname{wittrow}(t0)), \operatorname{wittwin}(t0, \operatorname{twopow}(k)))), r + 1 + \operatorname{twopow}(\operatorname{succ}(k)))\).
\end{theorem}
\begin{proof}
The goal follows from \(\operatorname{crosscarryscalerowwinwittsectortwoconeeq}(t0, k, r)\).
\end{proof}
\begin{theorem}
\label{ll:paper:frobdefect-row-late-eq-frobdefect-winwitt-sector-two}
For every natural number \(t0\) and natural number \(k\) and natural number \(r\), if \(2 \cdot t0 + 1 < \operatorname{twopow}(k)\), then if \(r + 1 \le k\), then \(\operatorname{bitfrobdefecteq}(k, \operatorname{wittrow}(t0 + \operatorname{twopow}(k)), r + 1 + \operatorname{twopow}(\operatorname{succ}(k)), \operatorname{wittbitbool}(\operatorname{frobdefect}(k, \operatorname{wittwin}(t0, \operatorname{twopow}(k))), r + 1 + \operatorname{twopow}(\operatorname{succ}(k))))\).
\end{theorem}
\begin{proof}
The goal follows from \(\operatorname{frobdefectrowlateeqfrobdefectwinwittsectortwo}(t0, k, r)\).
\end{proof}
\begin{theorem}
\label{ll:paper:frobdefect-row-late-eq-winwitt-sector-two-wronskian}
For every natural number \(t0\) and natural number \(k\) and natural number \(r\), if \(2 \cdot t0 + 1 < \operatorname{twopow}(k)\), then if \(r + 1 \le k\), then \(\operatorname{bitfrobdefecteq}(k, \operatorname{wittrow}(t0 + \operatorname{twopow}(k)), r + 1 + \operatorname{twopow}(\operatorname{succ}(k)), \operatorname{bxor}(\operatorname{wittbitbool}(\operatorname{wittwin}(t0, \operatorname{twopow}(k)), r + 1 + \operatorname{twopow}(\operatorname{succ}(k))) \operatorname{and} \operatorname{wittbitbool}(\operatorname{wittwin}(t0, \operatorname{twopow}(k)), r + \operatorname{twopow}(k)), \operatorname{wittbitbool}(\operatorname{wittwin}(t0, \operatorname{twopow}(k)), r + 1 + \operatorname{twopow}(k)) \operatorname{and} \operatorname{wittbitbool}(\operatorname{wittwin}(t0, \operatorname{twopow}(k)), r + \operatorname{twopow}(\operatorname{succ}(k)))))\).
\end{theorem}
\begin{proof}
The goal follows from \(\operatorname{frobdefectrowlateeqwinwittsectortwowronskian}(t0, k, r)\).
\end{proof}
\begin{theorem}
\label{ll:paper:periodic-ldiag-shift-mul-eq-response-shift-mul}
For every natural number \(T\) and natural number \(P\), if \(0 < P\), then if \(\operatorname{centerperiodic}(T, P)\), then for every natural number \(m\) and natural number \(t\) and natural number \(k\), if \(T \le t\), then if \(\operatorname{twopow}(k) \le t\), then \(\operatorname{bxor}(\operatorname{ldiag}(t + m \cdot P - \operatorname{twopow}(k), t + m \cdot P), \operatorname{ldiag}(t - \operatorname{twopow}(k), t)) = \operatorname{bxor}(\operatorname{response}(t + m \cdot P, k), \operatorname{response}(t, k))\).
\end{theorem}
\begin{proof}
Assume \(hP\), \(h\).
The goal follows from \(\operatorname{periodicldiagshiftmuleqresponseshiftmul}(T, P, hP, h)\).
\end{proof}
\begin{theorem}
\label{ll:paper:not-periodic-maskparity-zero-mul}
For every natural number \(T\) and natural number \(P\), if \(0 < P\), then if \(\operatorname{centerperiodic}(T, P)\), then for every natural number \(T1\) and natural number \(m\), if \(0 < m\), then not \(\operatorname{maskparityzeroperiodic}(T1, m \cdot P)\).
\end{theorem}
\begin{proof}
Assume \(hP\), \(h\), \(T1\), \(m\), \(hm\).
The goal follows from \(\operatorname{notperiodicmaskparityzeromul}(T, P, hP, h, T1, m, hm)\).
\end{proof}
\begin{theorem}
\label{ll:paper:not-periodic-colm1-mul}
For every natural number \(T\) and natural number \(P\), if \(0 < P\), then if \(\operatorname{centerperiodic}(T, P)\), then for every natural number \(T1\) and natural number \(m\), if \(0 < m\), then not \(\operatorname{colnegoneperiodic}(T1, m \cdot P)\).
\end{theorem}
\begin{proof}
Assume \(hP\), \(h\), \(T1\), \(m\), \(hm\).
The goal follows from \(\operatorname{notperiodiccolm1mul}(T, P, hP, h, T1, m, hm)\).
\end{proof}
\begin{theorem}
\label{ll:paper:phiscale-row-late-sector-one-eq-winwitt-cone}
For every natural number \(t0\) and natural number \(k\) and natural number \(j\), if \(2 \cdot t0 + 1 < \operatorname{twopow}(k)\), then if \(j \le k\), then \(\operatorname{wittbitbool}(\operatorname{wittscale}(k, \operatorname{wittrow}(t0 + \operatorname{twopow}(k))), j + \operatorname{twopow}(k)) = \operatorname{wittbitbool}(\operatorname{wittwin}(t0, \operatorname{twopow}(k)), j + \operatorname{twopow}(k))\).
\end{theorem}
\begin{proof}
The goal follows from \(\operatorname{phiscalerowlatesectoroneeqwinwittcone}(t0, k, j)\).
\end{proof}
\begin{theorem}
\label{ll:paper:frobdefect-row-late-sector-one-eq-crosscarry-cone}
For every natural number \(t0\) and natural number \(k\) and natural number \(r\), if \(2 \cdot t0 + 1 < \operatorname{twopow}(k)\), then if \(r + 1 \le k\), then \(\operatorname{bitfrobdefecteq}(k, \operatorname{wittrow}(t0 + \operatorname{twopow}(k)), r + 1 + \operatorname{twopow}(k), \operatorname{wittbitbool}(\operatorname{crosscarry}(\operatorname{wittscale}(k, \operatorname{wittrow}(t0)), \operatorname{wittwin}(t0, \operatorname{twopow}(k))), r + 1 + \operatorname{twopow}(k)))\).
\end{theorem}
\begin{proof}
The goal follows from \(\operatorname{frobdefectrowlatesectoroneeqcrosscarrycone}(t0, k, r)\).
\end{proof}
\begin{theorem}
\label{ll:paper:frobdefect-row-late-sector-one-center-slot-wronskian}
For every natural number \(d\) and natural number \(k\), if \(d + 1 \le k\), then \(\operatorname{bitfrobdefecteq}(k, \operatorname{wittrow}(d + 1 + \operatorname{twopow}(k)), d + 1 + \operatorname{twopow}(k), \operatorname{bxor}(\operatorname{center}(d + 1 + \operatorname{twopow}(k)) \operatorname{and} \operatorname{rdiag}(d, d + 1), \operatorname{center}(d + 1) \operatorname{and} \operatorname{rdiag}(d + \operatorname{twopow}(k), d + 1 + \operatorname{twopow}(k))))\).
\end{theorem}
\begin{proof}
The goal follows from \(\operatorname{frobdefectrowlatesectoronecenterslotwronskian}(d, k)\).
\end{proof}
\begin{theorem}
\label{ll:paper:crosscarry-row-winwitt-sector-two-center-slot}
For every natural number \(d\) and natural number \(k\), if \(d + 1 \le k\), then if \(\operatorname{settledat}(d + 1, k)\), then if \(\operatorname{center}(d + 1) = \operatorname{false}\), then \(\operatorname{zcenterandeq}(d + 1, \operatorname{maskparity}(d + 1, 0), \operatorname{wittbitbool}(\operatorname{crosscarry}(\operatorname{wittscale}(k, \operatorname{wittrow}(d + 1)), \operatorname{wittwin}(d + 1, \operatorname{twopow}(k))), d + 1 + \operatorname{twopow}(\operatorname{succ}(k))))\).
\end{theorem}
\begin{proof}
The goal follows from \(\operatorname{crosscarryrowwinwittsectortwocenterslot}(d, k)\).
\end{proof}
\begin{theorem}
\label{ll:paper:ldiag-settled-restart}
For every natural number \(k\) and natural number \(j\) and natural number \(v\) and natural number \(t\), if \(j \le k\), then if \(\operatorname{settledby}(j, t)\), then \(\operatorname{ldiag}(j, t + v \cdot \operatorname{twopow}(k)) = \operatorname{ldiag}(j, t)\).
\end{theorem}
\begin{proof}
The goal follows from \(\operatorname{ldiagsettledrestart}(k, j, v, t)\).
\end{proof}
\begin{theorem}
\label{ll:paper:zcenter-restart-ldiag}
For every natural number \(k\) and natural number \(d\) and natural number \(u\), if \(d \le k\), then if \(1 \le u\), then if \(\operatorname{settledat}(d, k)\), then \(\operatorname{zcentereq}(d, \operatorname{ldiag}(d, d + u \cdot \operatorname{twopow}(k)))\).
\end{theorem}
\begin{proof}
The goal follows from \(\operatorname{zcenterrestartldiag}(k, d, u)\).
\end{proof}
\begin{theorem}
\label{ll:paper:frobdefect-row-late-center-slot-of-zcenter-true}
For every natural number \(d\) and natural number \(k\), if \(d + 1 \le k\), then if \(\operatorname{settledat}(d + 1, k)\), then if \(\operatorname{zcentereq}(d + 1, \operatorname{true})\), then if \(\operatorname{center}(d + 1) = \operatorname{false}\), then if \(\operatorname{center}(d + 1 + \operatorname{twopow}(k)) = \operatorname{false}\), then \(\operatorname{bitfrobdefecteq}(k, \operatorname{wittrow}(d + 1 + \operatorname{twopow}(k)), d + 1 + \operatorname{twopow}(k) + \operatorname{twopow}(k), \operatorname{bxor}(\operatorname{maskparity}(d + 1 + \operatorname{twopow}(k), 0), \operatorname{maskparity}(d + 1, 0)))\).
\end{theorem}
\begin{proof}
The goal follows from \(\operatorname{frobdefectrowlatecenterslotofzcentertrue}(d, k)\).
\end{proof}
\begin{theorem}
\label{ll:paper:periodic-winwitt-step-late-center-slot-of-zcenter-true}
For every natural number \(T\) and natural number \(P\), if \(0 < P\), then if \(\operatorname{centerperiodic}(T, P)\), then for every natural number \(d\) and natural number \(k\), if \(T \le d + 1\), then if \(d + 1 \le k\), then if \(\operatorname{settledat}(d + 1, k)\), then if \(\operatorname{zcentereq}(d + 1, \operatorname{true})\), then if \(\operatorname{center}(d + 1) = \operatorname{false}\), then if \(\operatorname{center}(d + 1 + \operatorname{lag}(P, k)) = \operatorname{false}\), then if \(\operatorname{center}(d + 1 + \operatorname{lag}(P, \operatorname{succ}(k))) = \operatorname{false}\), then \(\operatorname{bxor}(\operatorname{maskparity}(d + 1 + \operatorname{twopow}(\operatorname{succ}(k)), 0), \operatorname{bxor}(\operatorname{maskparity}(d + 1 + \operatorname{twopow}(k), 0), \operatorname{maskparity}(d + 1, 0))) = \operatorname{bxor}(\operatorname{wittresp}(d + 2 + \operatorname{twopow}(k), k), \operatorname{wittbitbool}(\operatorname{wittphi}(\operatorname{wittwin}(d + 1 + \operatorname{twopow}(k), \operatorname{twopow}(k))), d + 2 + \operatorname{twopow}(k) + \operatorname{twopow}(k)))\).
\end{theorem}
\begin{proof}
Assume \(hP\), \(h\).
The goal follows from \(\operatorname{periodicwinwittsteplatecenterslotofzcentertrue}(T, P, hP, h)\).
\end{proof}
\begin{theorem}
\label{ll:paper:periodic-winwitt-step-late-colm1-slot-of-zcenter-true}
For every natural number \(T\) and natural number \(P\), if \(0 < P\), then if \(\operatorname{centerperiodic}(T, P)\), then for every natural number \(d\) and natural number \(k\), if \(T \le d + 1\), then if \(d + 2 \le k\), then if \(\operatorname{settledat}(d + 1, k)\), then if \(\operatorname{zcentereq}(d + 1, \operatorname{true})\), then if \(\operatorname{center}(d + 1) = \operatorname{false}\), then if \(\operatorname{center}(d + 1 + \operatorname{lag}(P, k)) = \operatorname{false}\), then if \(\operatorname{center}(d + 1 + \operatorname{lag}(P, \operatorname{succ}(k))) = \operatorname{false}\), then \(\operatorname{bxor}(\operatorname{ldiag}(d + \operatorname{twopow}(\operatorname{succ}(k)), d + 1 + \operatorname{twopow}(\operatorname{succ}(k))), \operatorname{bxor}(\operatorname{ldiag}(d + \operatorname{twopow}(k), d + 1 + \operatorname{twopow}(k)), \operatorname{ldiag}(d, d + 1))) = \operatorname{bxor}(\operatorname{wittbitbool}(\operatorname{wittwin}(d + 2 + \operatorname{twopow}(k), \operatorname{twopow}(k)), d + 3 + \operatorname{twopow}(k) + \operatorname{twopow}(k)), \operatorname{wittbitbool}(\operatorname{wittphi}(\operatorname{wittwin}(d + 1 + \operatorname{twopow}(k), \operatorname{twopow}(k))), d + 3 + \operatorname{twopow}(k) + \operatorname{twopow}(k)))\).
\end{theorem}
\begin{proof}
Assume \(hP\), \(h\).
The goal follows from \(\operatorname{periodicwinwittsteplatecolm1slotofzcentertrue}(T, P, hP, h)\).
\end{proof}
\begin{theorem}
\label{ll:paper:rowagree-mono}
For every natural number \(t\) and natural number \(Q\) and natural number \(a\) and natural number \(L\) and natural number \(a1\) and natural number \(L1\), if \(a \le a1\), then if \(a1 + L1 \le a + L\), then if \(\operatorname{rowagree}(t, Q, a, L)\), then \(\operatorname{rowagree}(t, Q, a1, L1)\).
\end{theorem}
\begin{proof}
The goal follows from \(\operatorname{rowagreemono}(t, Q, a, L, a1, L1)\).
\end{proof}
\begin{theorem}
\label{ll:paper:winwittagree-mono}
For every natural number \(t\) and natural number \(k\) and natural number \(Q\) and natural number \(a\) and natural number \(L\) and natural number \(a1\) and natural number \(L1\), if \(a \le a1\), then if \(a1 + L1 \le a + L\), then if \(\operatorname{winwittagree}(t, k, Q, a, L)\), then \(\operatorname{winwittagree}(t, k, Q, a1, L1)\).
\end{theorem}
\begin{proof}
The goal follows from \(\operatorname{winwittagreemono}(t, k, Q, a, L, a1, L1)\).
\end{proof}
\begin{theorem}
\label{ll:paper:rowagree-step}
For every natural number \(t\) and natural number \(Q\) and natural number \(a\) and natural number \(L\), if \(\operatorname{rowagree}(t, Q, a, L + 2)\), then \(\operatorname{rowagree}(t + 1, Q, a + 2, L)\).
\end{theorem}
\begin{proof}
The goal follows from \(\operatorname{rowagreestep}(t, Q, a, L)\).
\end{proof}
\begin{theorem}
\label{ll:paper:rowagree-steps}
For every natural number \(t\) and natural number \(Q\) and natural number \(a\) and natural number \(L\) and natural number \(s\), if \(\operatorname{rowagree}(t, Q, a, L + 2 \cdot s)\), then \(\operatorname{rowagree}(t + s, Q, a + 2 \cdot s, L)\).
\end{theorem}
\begin{proof}
The goal follows from \(\operatorname{rowagreesteps}(t, Q, a, L, s)\).
\end{proof}
\begin{theorem}
\label{ll:paper:phiscale-row-testbit-eq-of-rowagree}
For every natural number \(t\) and natural number \(k\) and natural number \(Q\) and natural number \(a\) and natural number \(L\) and natural number \(i\), if \(\operatorname{twopow}(\operatorname{succ}(k)) \le a\), then if \(i \le L\), then if \(\operatorname{rowagree}(t, Q, a - \operatorname{twopow}(\operatorname{succ}(k)), L)\), then if \(\operatorname{rowagree}(t, Q, a - \operatorname{twopow}(k), L)\), then if \(\operatorname{rowagree}(t, Q, a, L)\), then \(\operatorname{wittbitbool}(\operatorname{wittscale}(k, \operatorname{wittrow}(t + Q)), a + i + Q) = \operatorname{wittbitbool}(\operatorname{wittscale}(k, \operatorname{wittrow}(t)), a + i)\).
\end{theorem}
\begin{proof}
The goal follows from \(\operatorname{phiscalerowtestbiteqofrowagree}(t, k, Q, a, L, i)\).
\end{proof}
\begin{theorem}
\label{ll:paper:frobdefect-row-testbit-eq-of-rowagree}
For every natural number \(t\) and natural number \(k\) and natural number \(Q\) and natural number \(a\) and natural number \(L\) and natural number \(i\), if \(\operatorname{twopow}(\operatorname{succ}(k)) \le a\), then if \(i \le L\), then if \(\operatorname{rowagree}(t, Q, a - \operatorname{twopow}(\operatorname{succ}(k)), L + 1)\), then if \(\operatorname{rowagree}(t, Q, a - \operatorname{twopow}(k), L + 1)\), then if \(\operatorname{rowagree}(t, Q, a, L + 1)\), then \(\operatorname{bitfrobdefecteq}(k, \operatorname{wittrow}(t + Q), a + i + 1 + Q, \operatorname{wittbitbool}(\operatorname{frobdefect}(k, \operatorname{wittrow}(t)), a + i + 1))\).
\end{theorem}
\begin{proof}
The goal follows from \(\operatorname{frobdefectrowtestbiteqofrowagree}(t, k, Q, a, L, i)\).
\end{proof}
\begin{theorem}
\label{ll:paper:crosscarry-testbit-eq-of-rowagree-winwittagree}
For every natural number \(t\) and natural number \(k\) and natural number \(Q\) and natural number \(a\) and natural number \(L\) and natural number \(i\), if \(\operatorname{twopow}(\operatorname{succ}(k)) \le a\), then if \(i \le L\), then if \(\operatorname{winwittagree}(t, k, Q, a, L + 1)\), then if \(\operatorname{rowagree}(t, Q, a - \operatorname{twopow}(\operatorname{succ}(k)), L + 1)\), then if \(\operatorname{rowagree}(t, Q, a - \operatorname{twopow}(k), L + 1)\), then if \(\operatorname{rowagree}(t, Q, a, L + 1)\), then \(\operatorname{bitcrosscarryeq}(\operatorname{wittscale}(k, \operatorname{wittrow}(t + Q)), \operatorname{wittwin}(t + Q, \operatorname{twopow}(k)), a + i + 1 + Q, \operatorname{wittbitbool}(\operatorname{crosscarry}(\operatorname{wittscale}(k, \operatorname{wittrow}(t)), \operatorname{wittwin}(t, \operatorname{twopow}(k))), a + i + 1))\).
\end{theorem}
\begin{proof}
The goal follows from \(\operatorname{crosscarrytestbiteqofrowagreewinwittagree}(t, k, Q, a, L, i)\).
\end{proof}
\begin{theorem}
\label{ll:paper:winwittagree-step}
For every natural number \(t\) and natural number \(k\) and natural number \(Q\) and natural number \(a\) and natural number \(L\), if \(\operatorname{twopow}(\operatorname{succ}(k)) \le a\), then if \(\operatorname{winwittagree}(t, k, Q, a, L + 2)\), then if \(\operatorname{rowagree}(t, Q, a - \operatorname{twopow}(\operatorname{succ}(k)), L + 1)\), then if \(\operatorname{rowagree}(t, Q, a - \operatorname{twopow}(k), L + 1)\), then if \(\operatorname{rowagree}(t, Q, a, L + 1)\), then \(\operatorname{winwittagree}(t + 1, k, Q, a + 2, L)\).
\end{theorem}
\begin{proof}
The goal follows from \(\operatorname{winwittagreestep}(t, k, Q, a, L)\).
\end{proof}
\begin{theorem}
\label{ll:paper:winwittagree-steps}
For every natural number \(t\) and natural number \(k\) and natural number \(Q\) and natural number \(a\) and natural number \(L\) and natural number \(s\), if \(\operatorname{twopow}(\operatorname{succ}(k)) \le a\), then if \(\operatorname{winwittagree}(t, k, Q, a, L + 2 \cdot s)\), then if \(\operatorname{rowagree}(t, Q, a - \operatorname{twopow}(\operatorname{succ}(k)), L + 2 \cdot s)\), then if \(\operatorname{rowagree}(t, Q, a - \operatorname{twopow}(k), L + 2 \cdot s)\), then if \(\operatorname{rowagree}(t, Q, a, L + 2 \cdot s)\), then \(\operatorname{winwittagree}(t + s, k, Q, a + 2 \cdot s, L)\).
\end{theorem}
\begin{proof}
The goal follows from \(\operatorname{winwittagreesteps}(t, k, Q, a, L, s)\).
\end{proof}
\begin{theorem}
\label{ll:paper:wittresp-eq-of-winwittagree-steps}
For every natural number \(t0\) and natural number \(k\) and natural number \(Q\) and natural number \(s\), if \(\operatorname{twopow}(k) + s \le t0\), then if \(\operatorname{winwittagree}(t0, k, Q, t0 + \operatorname{twopow}(k) - s, 2 \cdot s)\), then if \(\operatorname{rowagree}(t0, Q, t0 + \operatorname{twopow}(k) - s - \operatorname{twopow}(\operatorname{succ}(k)), 2 \cdot s)\), then if \(\operatorname{rowagree}(t0, Q, t0 + \operatorname{twopow}(k) - s - \operatorname{twopow}(k), 2 \cdot s)\), then if \(\operatorname{rowagree}(t0, Q, t0 + \operatorname{twopow}(k) - s, 2 \cdot s)\), then \(\operatorname{wittresp}(t0 + s + Q, k) = \operatorname{wittresp}(t0 + s, k)\).
\end{theorem}
\begin{proof}
The goal follows from \(\operatorname{wittrespeqofwinwittagreesteps}(t0, k, Q, s)\).
\end{proof}
\begin{theorem}
\label{ll:paper:colm1-eq-of-rowagree-steps}
For every natural number \(t0\) and natural number \(Q\) and natural number \(s\), if \(s \le t0 + 2\), then if \(\operatorname{rowagree}(t0 + 1, Q, t0 + 2 - s, 2 \cdot s)\), then \(\operatorname{ldiag}(t0 + s + Q, t0 + s + 1 + Q) = \operatorname{ldiag}(t0 + s, t0 + s + 1)\).
\end{theorem}
\begin{proof}
The goal follows from \(\operatorname{colm1eqofrowagreesteps}(t0, Q, s)\).
\end{proof}
\begin{theorem}
\label{ll:paper:maskparity-zero-eq-of-rowagree-steps}
For every natural number \(T\) and natural number \(P\), if \(0 < P\), then if \(\operatorname{centerperiodic}(T, P)\), then for every natural number \(m\) and natural number \(t0\) and natural number \(s\), if \(T \le t0 + s + 1\), then if \(s \le t0 + 2\), then if \(\operatorname{rowagree}(t0 + 1, m \cdot P, t0 + 2 - s, 2 \cdot s)\), then \(\operatorname{maskparity}(t0 + s + 1 + m \cdot P, 0) = \operatorname{maskparity}(t0 + s + 1, 0)\).
\end{theorem}
\begin{proof}
Assume \(hP\), \(hper\).
The goal follows from \(\operatorname{maskparityzeroeqofrowagreesteps}(T, P, hP, hper)\).
\end{proof}
\begin{theorem}
\label{ll:paper:colnegoneperiodic-of-sector-rowagree}
For every natural number \(k\) and natural number \(d\) and natural number \(U\) and natural number \(Q\), if \(1 \le U\), then if \(\operatorname{sectorrowagree}(k, d, U, Q)\), then \(\operatorname{colnegoneperiodic}(d + U \cdot \operatorname{twopow}(k), Q)\).
\end{theorem}
\begin{proof}
The goal follows from \(\operatorname{colnegoneperiodicofsectorrowagree}(k, d, U, Q)\).
\end{proof}
\begin{theorem}
\label{ll:paper:not-sector-rowagree-mul}
For every natural number \(T\) and natural number \(P\), if \(0 < P\), then if \(\operatorname{centerperiodic}(T, P)\), then for every natural number \(k\) and natural number \(d\) and natural number \(U\) and natural number \(m\), if \(1 \le U\), then if \(0 < m\), then not \(\operatorname{sectorrowagree}(k, d, U, m \cdot P)\).
\end{theorem}
\begin{proof}
Assume \(hP\), \(h\), \(k\), \(d\), \(U\), \(m\), \(hU\), \(hm\).
The goal follows from \(\operatorname{notsectorrowagreemul}(T, P, hP, h, k, d, U, m, hU, hm)\).
\end{proof}
\begin{theorem}
\label{ll:paper:rdiag-lcmperiod-eq}
For every natural number \(P\) and natural number \(k\) and natural number \(j\) and natural number \(t\), if \(j \le k\), then \(\operatorname{rdiag}(j, t + \operatorname{lcmperiod}(P, k)) = \operatorname{rdiag}(j, t)\).
\end{theorem}
\begin{proof}
The goal follows from \(\operatorname{rdiaglcmperiodeq}(P, k, j, t)\).
\end{proof}
\begin{theorem}
\label{ll:paper:ldiag-lcmperiod-eq}
For every natural number \(P\) and natural number \(k\) and natural number \(j\) and natural number \(t\), if \(j \le k\), then if \(\operatorname{settledby}(j, t)\), then \(\operatorname{ldiag}(j, t + \operatorname{lcmperiod}(P, k)) = \operatorname{ldiag}(j, t)\).
\end{theorem}
\begin{proof}
The goal follows from \(\operatorname{ldiaglcmperiodeq}(P, k, j, t)\).
\end{proof}
\begin{theorem}
\label{ll:paper:center-lcmperiod-eq}
For every natural number \(T\) and natural number \(P\), if \(0 < P\), then if \(\operatorname{centerperiodic}(T, P)\), then for every natural number \(k\) and natural number \(t\), if \(T \le t\), then \(\operatorname{center}(t + \operatorname{lcmperiod}(P, k)) = \operatorname{center}(t)\).
\end{theorem}
\begin{proof}
Assume \(hP\), \(h\).
The goal follows from \(\operatorname{centerlcmperiodeq}(T, P, hP, h)\).
\end{proof}
\begin{theorem}
\label{ll:paper:not-sector-rowagree-lcmperiod}
For every natural number \(T\) and natural number \(P\), if \(0 < P\), then if \(\operatorname{centerperiodic}(T, P)\), then for every natural number \(k\) and natural number \(d\) and natural number \(U\), if \(1 \le U\), then not \(\operatorname{sectorrowagree}(k, d, U, \operatorname{lcmperiod}(P, k))\).
\end{theorem}
\begin{proof}
Assume \(hP\), \(h\), \(k\), \(d\), \(U\), \(hU\).
The goal follows from \(\operatorname{notsectorrowagreelcmperiod}(T, P, hP, h, k, d, U, hU)\).
\end{proof}
\begin{theorem}
\label{ll:paper:rowagree-extend-right}
For every natural number \(t\) and natural number \(Q\) and natural number \(a\) and natural number \(L\), if \(\operatorname{rowagree}(t, Q, a, L + 1)\), then if \(\operatorname{wittbitbool}(\operatorname{wittrow}(t + 1 + Q), a + L + 2 + Q) = \operatorname{wittbitbool}(\operatorname{wittrow}(t + 1), a + L + 2)\), then \(\operatorname{rowagree}(t, Q, a, L + 2)\).
\end{theorem}
\begin{proof}
The goal follows from \(\operatorname{rowagreeextendright}(t, Q, a, L)\).
\end{proof}
\begin{theorem}
\label{ll:paper:periodic-rowagree-center}
For every natural number \(T\) and natural number \(P\), if \(0 < P\), then if \(\operatorname{centerperiodic}(T, P)\), then for every natural number \(m\) and natural number \(t\), if \(T \le t\), then \(\operatorname{rowagree}(t, m \cdot P, t, 0)\).
\end{theorem}
\begin{proof}
Assume \(hP\), \(hper\).
The goal follows from \(\operatorname{periodicrowagreecenter}(T, P, hP, hper)\).
\end{proof}
\begin{theorem}
\label{ll:paper:periodic-rowagree-black-center}
For every natural number \(T\) and natural number \(P\), if \(0 < P\), then if \(\operatorname{centerperiodic}(T, P)\), then for every natural number \(m\) and natural number \(t\), if \(T \le t + 1\), then if \(\operatorname{center}(t + 1) = \operatorname{true}\), then \(\operatorname{rowagree}(t + 1, m \cdot P, t + 1, 1)\).
\end{theorem}
\begin{proof}
Assume \(hP\), \(hper\).
The goal follows from \(\operatorname{periodicrowagreeblackcenter}(T, P, hP, hper)\).
\end{proof}
\begin{theorem}
\label{ll:paper:periodic-rowagree-symm-of-right}
For every natural number \(T\) and natural number \(P\), if \(0 < P\), then if \(\operatorname{centerperiodic}(T, P)\), then for every natural number \(m\) and natural number \(R\) and natural number \(t\), if \(T \le t\), then if \(R \le t\), then if \(\operatorname{rowagree}(t, m \cdot P, t - R, R)\), then \(\operatorname{rowagree}(t, m \cdot P, t - R, 2 \cdot R)\).
\end{theorem}
\begin{proof}
Assume \(hP\), \(hper\).
The goal follows from \(\operatorname{periodicrowagreesymmofright}(T, P, hP, hper)\).
\end{proof}
\begin{theorem}
\label{ll:paper:periodic-rowagree-lcmperiod-symm-of-right}
For every natural number \(T\) and natural number \(P\), if \(0 < P\), then if \(\operatorname{centerperiodic}(T, P)\), then for every natural number \(k\) and natural number \(R\) and natural number \(t\), if \(T \le t\), then if \(R \le t\), then if \(\operatorname{rowagree}(t, \operatorname{lcmperiod}(P, k), t - R, R)\), then \(\operatorname{rowagree}(t, \operatorname{lcmperiod}(P, k), t - R, 2 \cdot R)\).
\end{theorem}
\begin{proof}
Assume \(hP\), \(hper\).
The goal follows from \(\operatorname{periodicrowagreelcmperiodsymmofright}(T, P, hP, hper)\).
\end{proof}
\begin{theorem}
\label{ll:paper:colnegoneperiodic-of-sector-right-rowagree}
For every natural number \(T\) and natural number \(P\), if \(0 < P\), then if \(\operatorname{centerperiodic}(T, P)\), then for every natural number \(k\) and natural number \(d\) and natural number \(U\) and natural number \(m\), if \(1 \le U\), then if \(1 \le d\), then if \(T \le d + U \cdot \operatorname{twopow}(k)\), then if \(\operatorname{sectorrightrowagree}(k, d, U, m \cdot P)\), then \(\operatorname{colnegoneperiodic}(d + U \cdot \operatorname{twopow}(k), m \cdot P)\).
\end{theorem}
\begin{proof}
Assume \(hP\), \(hper\).
The goal follows from \(\operatorname{colnegoneperiodicofsectorrightrowagree}(T, P, hP, hper)\).
\end{proof}
\begin{theorem}
\label{ll:paper:not-sector-right-rowagree-lcmperiod}
For every natural number \(T\) and natural number \(P\), if \(0 < P\), then if \(\operatorname{centerperiodic}(T, P)\), then for every natural number \(k\) and natural number \(d\) and natural number \(U\), if \(1 \le U\), then if \(1 \le d\), then if \(T \le d + U \cdot \operatorname{twopow}(k)\), then not \(\operatorname{sectorrightrowagree}(k, d, U, \operatorname{lcmperiod}(P, k))\).
\end{theorem}
\begin{proof}
Assume \(hP\), \(hper\), \(k\), \(d\), \(U\), \(hU\), \(hd\), \(ht\).
The goal follows from \(\operatorname{notsectorrightrowagreelcmperiod}(T, P, hP, hper, k, d, U, hU, hd, ht)\).
\end{proof}
\begin{theorem}
\label{ll:paper:winwittagree-extend-right}
For every natural number \(t\) and natural number \(k\) and natural number \(Q\) and natural number \(a\) and natural number \(L\), if \(\operatorname{twopow}(\operatorname{succ}(k)) \le a\), then if \(\operatorname{winwittagree}(t, k, Q, a, L + 1)\), then if \(\operatorname{rowagree}(t, Q, a - \operatorname{twopow}(\operatorname{succ}(k)), L + 1)\), then if \(\operatorname{rowagree}(t, Q, a - \operatorname{twopow}(k), L + 1)\), then if \(\operatorname{rowagree}(t, Q, a, L + 1)\), then if \(\operatorname{wittbitbool}(\operatorname{wittwin}(t + 1 + Q, \operatorname{twopow}(k)), a + L + 2 + Q) = \operatorname{wittbitbool}(\operatorname{wittwin}(t + 1, \operatorname{twopow}(k)), a + L + 2)\), then \(\operatorname{winwittagree}(t, k, Q, a, L + 2)\).
\end{theorem}
\begin{proof}
The goal follows from \(\operatorname{winwittagreeextendright}(t, k, Q, a, L)\).
\end{proof}
\begin{theorem}
\label{ll:paper:winwittagree-of-rowagree-three-and-late}
For every natural number \(t\) and natural number \(k\) and natural number \(Q\) and natural number \(a\) and natural number \(L\), if \(\operatorname{twopow}(\operatorname{succ}(k)) \le a\), then if \(\operatorname{rowagree}(t, Q, a - \operatorname{twopow}(\operatorname{succ}(k)), L)\), then if \(\operatorname{rowagree}(t, Q, a - \operatorname{twopow}(k), L)\), then if \(\operatorname{rowagree}(t, Q, a, L)\), then if \(\operatorname{rowagree}(t + \operatorname{twopow}(k), Q, a, L)\), then \(\operatorname{winwittagree}(t, k, Q, a, L)\).
\end{theorem}
\begin{proof}
The goal follows from \(\operatorname{winwittagreeofrowagreethreeandlate}(t, k, Q, a, L)\).
\end{proof}
\begin{theorem}
\label{ll:paper:rowagree-late-of-winwittagree-and-three}
For every natural number \(t\) and natural number \(k\) and natural number \(Q\) and natural number \(a\) and natural number \(L\), if \(\operatorname{twopow}(\operatorname{succ}(k)) \le a\), then if \(\operatorname{winwittagree}(t, k, Q, a, L)\), then if \(\operatorname{rowagree}(t, Q, a - \operatorname{twopow}(\operatorname{succ}(k)), L)\), then if \(\operatorname{rowagree}(t, Q, a - \operatorname{twopow}(k), L)\), then if \(\operatorname{rowagree}(t, Q, a, L)\), then \(\operatorname{rowagree}(t + \operatorname{twopow}(k), Q, a, L)\).
\end{theorem}
\begin{proof}
The goal follows from \(\operatorname{rowagreelateofwinwittagreeandthree}(t, k, Q, a, L)\).
\end{proof}
\begin{theorem}
\label{ll:paper:wittrespbool-periodic-of-sector-winwittagree}
For every natural number \(k\) and natural number \(d\) and natural number \(U\) and natural number \(Q\), if \(2 \le U\), then if \(\operatorname{sectorwinwittagree}(k, d, U, Q)\), then for every natural number \(t\), if \(d + U \cdot \operatorname{twopow}(k) \le t\), then \(\operatorname{wittresp}(t + Q, k) = \operatorname{wittresp}(t, k)\).
\end{theorem}
\begin{proof}
The goal follows from \(\operatorname{wittrespboolperiodicofsectorwinwittagree}(k, d, U, Q)\).
\end{proof}
\begin{theorem}
\label{ll:paper:right-boundary-row-lcmperiod-eq}
For every natural number \(P\) and natural number \(k\) and natural number \(j\) and natural number \(t\), if \(j \le k\), then \(\operatorname{wittbitbool}(\operatorname{wittrow}(t + \operatorname{lcmperiod}(P, k)), j) = \operatorname{wittbitbool}(\operatorname{wittrow}(t), j)\).
\end{theorem}
\begin{proof}
The goal follows from \(\operatorname{rightboundaryrowlcmperiodeq}(P, k, j, t)\).
\end{proof}
\begin{theorem}
\label{ll:paper:left-boundary-row-lcmperiod-eq}
For every natural number \(P\) and natural number \(k\) and natural number \(j\) and natural number \(t\), if \(j \le k\), then if \(j \le t\), then if \(\operatorname{settledby}(j, t)\), then \(\operatorname{wittbitbool}(\operatorname{wittrow}(t + \operatorname{lcmperiod}(P, k)), 2 \cdot (t + \operatorname{lcmperiod}(P, k)) - j) = \operatorname{wittbitbool}(\operatorname{wittrow}(t), 2 \cdot t - j)\).
\end{theorem}
\begin{proof}
The goal follows from \(\operatorname{leftboundaryrowlcmperiodeq}(P, k, j, t)\).
\end{proof}
\begin{theorem}
\label{ll:paper:frobdefect-row-sector-one-lcmperiod-xor}
For every natural number \(P\) and natural number \(d\) and natural number \(k\) and natural number \(u\), if \(d + 1 \le k\), then \(\operatorname{bxor}(\operatorname{wittbitbool}(\operatorname{frobdefect}(k, \operatorname{wittrow}(d + 1 + u \cdot \operatorname{twopow}(k) + \operatorname{lcmperiod}(P, k))), d + 1 + \operatorname{twopow}(k)), \operatorname{wittbitbool}(\operatorname{frobdefect}(k, \operatorname{wittrow}(d + 1 + u \cdot \operatorname{twopow}(k))), d + 1 + \operatorname{twopow}(k))) = \operatorname{bxor}(\operatorname{bxor}(\operatorname{rdiag}(d + 1 + \operatorname{twopow}(k), d + 1 + u \cdot \operatorname{twopow}(k) + \operatorname{lcmperiod}(P, k)), \operatorname{rdiag}(d + 1 + \operatorname{twopow}(k), d + 1 + u \cdot \operatorname{twopow}(k))) \operatorname{and} \operatorname{rdiag}(d, d + 1), \operatorname{center}(d + 1) \operatorname{and} \operatorname{bxor}(\operatorname{rdiag}(d + \operatorname{twopow}(k), d + 1 + u \cdot \operatorname{twopow}(k) + \operatorname{lcmperiod}(P, k)), \operatorname{rdiag}(d + \operatorname{twopow}(k), d + 1 + u \cdot \operatorname{twopow}(k))))\).
\end{theorem}
\begin{proof}
The goal follows from \(\operatorname{frobdefectrowsectoronelcmperiodxor}(P, d, k, u)\).
\end{proof}
\begin{theorem}
\label{ll:paper:rhcenter-eq-rdiag}
For every natural number \(t\) and natural number \(p\), if \(p \le t\), then \(\operatorname{rhcenter}(t, p) = \operatorname{rdiag}(t - p, t)\).
\end{theorem}
\begin{proof}
The goal follows from \(\operatorname{rhcentereqrdiag}(t, p)\).
\end{proof}
\begin{theorem}
\label{ll:paper:rhcenter-eq-row-testbit}
For every natural number \(t\) and natural number \(p\), if \(p \le t\), then \(\operatorname{rhcenter}(t, p) = \operatorname{wittbitbool}(\operatorname{wittrow}(t), t - p)\).
\end{theorem}
\begin{proof}
The goal follows from \(\operatorname{rhcentereqrowtestbit}(t, p)\).
\end{proof}
\begin{theorem}
\label{ll:paper:lhcenter-eq-ldiag}
For every natural number \(t\) and natural number \(p\), if \(p \le t\), then \(\operatorname{lhcenter}(t, p) = \operatorname{ldiag}(t - p, t)\).
\end{theorem}
\begin{proof}
The goal follows from \(\operatorname{lhcentereqldiag}(t, p)\).
\end{proof}
\begin{theorem}
\label{ll:paper:lhcenter-eq-row-testbit}
For every natural number \(t\) and natural number \(p\), if \(p \le t\), then \(\operatorname{lhcenter}(t, p) = \operatorname{wittbitbool}(\operatorname{wittrow}(t), t + p)\).
\end{theorem}
\begin{proof}
The goal follows from \(\operatorname{lhcentereqrowtestbit}(t, p)\).
\end{proof}
\begin{theorem}
\label{ll:paper:rhcenter-zero-eq-center}
For every natural number \(t\), \(\operatorname{rhcenter}(t, 0) = \operatorname{center}(t)\).
\end{theorem}
\begin{proof}
The goal follows from \(\operatorname{rhcenterzeroeqcenter}(t)\).
\end{proof}
\begin{theorem}
\label{ll:paper:rhcenter-self-eq-true}
For every natural number \(t\), \(\operatorname{rhcenter}(t, t) = \operatorname{true}\).
\end{theorem}
\begin{proof}
The goal follows from \(\operatorname{rhcenterselfeqtrue}(t)\).
\end{proof}
\begin{theorem}
\label{ll:paper:rhcenter-gt-eq-false}
For every natural number \(t\) and natural number \(p\), if \(t < p\), then \(\operatorname{rhcenter}(t, p) = \operatorname{false}\).
\end{theorem}
\begin{proof}
The goal follows from \(\operatorname{rhcentergteqfalse}(t, p)\).
\end{proof}
\begin{theorem}
\label{ll:paper:phicenter-eq-lhcenter-rhcenter}
For every natural number \(t\), \(\operatorname{phicenter}(t) = \operatorname{bxor}(\operatorname{lhcenter}(t, 1), \operatorname{center}(t) \operatorname{or} \operatorname{rhcenter}(t, 1))\).
\end{theorem}
\begin{proof}
The goal follows from \(\operatorname{phicentereqlhcenterrhcenter}(t)\).
\end{proof}
\begin{theorem}
\label{ll:paper:lhcenter-one-eq-of-center-fixed}
For every natural number \(t\), \(\operatorname{lhcenter}(t, 1) = \operatorname{bxor}(\operatorname{center}(t + 1), \operatorname{center}(t) \operatorname{or} \operatorname{rhcenter}(t, 1))\).
\end{theorem}
\begin{proof}
The goal follows from \(\operatorname{lhcenteroneeqofcenterfixed}(t)\).
\end{proof}
\begin{theorem}
\label{ll:paper:periodic-lhcenter-one-defect-eq-not-center-and-rhcenter-one-defect}
For every natural number \(T\) and natural number \(P\), if \(0 < P\), then if \(\operatorname{centerperiodic}(T, P)\), then for every natural number \(m\) and natural number \(t\), if \(T \le t\), then \(\operatorname{bxor}(\operatorname{lhcenter}(t + m \cdot P, 1), \operatorname{lhcenter}(t, 1)) = \operatorname{not}(\operatorname{center}(t)) \operatorname{and} \operatorname{bxor}(\operatorname{rhcenter}(t + m \cdot P, 1), \operatorname{rhcenter}(t, 1))\).
\end{theorem}
\begin{proof}
Assume \(hP\), \(hper\).
The goal follows from \(\operatorname{periodiclhcenteronedefecteqnotcenterandrhcenteronedefect}(T, P, hP, hper)\).
\end{proof}
\begin{theorem}
\label{ll:paper:rhcenter-step-succ}
For every natural number \(t\) and natural number \(p\), \(\operatorname{rhcenter}(t + 1, p + 1) = \operatorname{bxor}(\operatorname{rhcenter}(t, p), \operatorname{rhcenter}(t, p + 1) \operatorname{or} \operatorname{rhcenter}(t, p + 2))\).
\end{theorem}
\begin{proof}
The goal follows from \(\operatorname{rhcenterstepsucc}(t, p)\).
\end{proof}
\begin{theorem}
\label{ll:paper:rhcenter-inward-step}
For every natural number \(t\) and natural number \(p\), \(\operatorname{rhcenter}(t, p) = \operatorname{bxor}(\operatorname{rhcenter}(t + 1, p + 1), \operatorname{rhcenter}(t, p + 1) \operatorname{or} \operatorname{rhcenter}(t, p + 2))\).
\end{theorem}
\begin{proof}
The goal follows from \(\operatorname{rhcenterinwardstep}(t, p)\).
\end{proof}
\begin{theorem}
\label{ll:paper:rhbandagree-extend-inward}
For every natural number \(t\) and natural number \(Q\) and natural number \(p0\) and natural number \(L\), if \(\operatorname{rhbandagree}(t, Q, p0 + 1, L + 1)\), then if \(\operatorname{rhcenter}(t + 1 + Q, p0 + 1) = \operatorname{rhcenter}(t + 1, p0 + 1)\), then \(\operatorname{rhbandagree}(t, Q, p0, L + 2)\).
\end{theorem}
\begin{proof}
The goal follows from \(\operatorname{rhbandagreeextendinward}(t, Q, p0, L)\).
\end{proof}
\begin{theorem}
\label{ll:paper:rhbandagree-step}
For every natural number \(t\) and natural number \(Q\) and natural number \(p0\) and natural number \(L\), if \(\operatorname{rhbandagree}(t, Q, p0, L + 2)\), then \(\operatorname{rhbandagree}(t + 1, Q, p0 + 1, L)\).
\end{theorem}
\begin{proof}
The goal follows from \(\operatorname{rhbandagreestep}(t, Q, p0, L)\).
\end{proof}
\begin{theorem}
\label{ll:paper:rowagree-right-of-rhbandagree}
For every natural number \(t\) and natural number \(Q\) and natural number \(p0\) and natural number \(L\), if \(p0 + L \le t\), then if \(\operatorname{rhbandagree}(t, Q, p0, L)\), then \(\operatorname{rowagree}(t, Q, t - (p0 + L), L)\).
\end{theorem}
\begin{proof}
The goal follows from \(\operatorname{rowagreerightofrhbandagree}(t, Q, p0, L)\).
\end{proof}
\begin{theorem}
\label{ll:paper:rhbandagree-of-rowagree-right}
For every natural number \(t\) and natural number \(Q\) and natural number \(p0\) and natural number \(L\), if \(p0 + L \le t\), then if \(\operatorname{rowagree}(t, Q, t - (p0 + L), L)\), then \(\operatorname{rhbandagree}(t, Q, p0, L)\).
\end{theorem}
\begin{proof}
The goal follows from \(\operatorname{rhbandagreeofrowagreeright}(t, Q, p0, L)\).
\end{proof}
\begin{theorem}
\label{ll:paper:sectorrightrowagree-of-rhbandagree}
For every natural number \(k\) and natural number \(d\) and natural number \(U\) and natural number \(Q\), if \(1 \le U\), then if \(1 \le d\), then if \(\operatorname{sectorrhbandagree}(k, d, U, Q)\), then \(\operatorname{sectorrightrowagree}(k, d, U, Q)\).
\end{theorem}
\begin{proof}
The goal follows from \(\operatorname{sectorrightrowagreeofrhbandagree}(k, d, U, Q)\).
\end{proof}
\begin{theorem}
\label{ll:paper:not-sector-rhbandagree-lcmperiod}
For every natural number \(T\) and natural number \(P\), if \(0 < P\), then if \(\operatorname{centerperiodic}(T, P)\), then for every natural number \(k\) and natural number \(d\) and natural number \(U\), if \(1 \le U\), then if \(1 \le d\), then if \(T \le d + U \cdot \operatorname{twopow}(k)\), then not \(\operatorname{sectorrhbandagree}(k, d, U, \operatorname{lcmperiod}(P, k))\).
\end{theorem}
\begin{proof}
Assume \(hP\), \(hper\), \(k\), \(d\), \(U\), \(hU\), \(hd\), \(ht\).
The goal follows from \(\operatorname{notsectorrhbandagreelcmperiod}(T, P, hP, hper, k, d, U, hU, hd, ht)\).
\end{proof}
\begin{theorem}
\label{ll:paper:rhcenter-right-edge-lcmperiod-eq}
For every natural number \(P\) and natural number \(k\) and natural number \(j\) and natural number \(t\), if \(j \le k\), then if \(j \le t\), then \(\operatorname{rhcenter}(t + \operatorname{lcmperiod}(P, k), t + \operatorname{lcmperiod}(P, k) - j) = \operatorname{rhcenter}(t, t - j)\).
\end{theorem}
\begin{proof}
The goal follows from \(\operatorname{rhcenterrightedgelcmperiodeq}(P, k, j, t)\).
\end{proof}
\begin{theorem}
\label{ll:paper:colnegoneperiodic-of-sector-winwittagree}
For every natural number \(k\) and natural number \(d\) and natural number \(U\) and natural number \(Q\), if \(2 \le U\), then if \(\operatorname{sectorwinwittagree}(k, d, U, Q)\), then \(\operatorname{colnegoneperiodic}(d + (U + 1) \cdot \operatorname{twopow}(k), Q)\).
\end{theorem}
\begin{proof}
The goal follows from \(\operatorname{colnegoneperiodicofsectorwinwittagree}(k, d, U, Q)\).
\end{proof}
\begin{theorem}
\label{ll:paper:not-sector-winwittagree-lcmperiod}
For every natural number \(T\) and natural number \(P\), if \(0 < P\), then if \(\operatorname{centerperiodic}(T, P)\), then for every natural number \(k\) and natural number \(d\) and natural number \(U\), if \(2 \le U\), then not \(\operatorname{sectorwinwittagree}(k, d, U, \operatorname{lcmperiod}(P, k))\).
\end{theorem}
\begin{proof}
Assume \(hP\), \(hper\), \(k\), \(d\), \(U\), \(hU\).
The goal follows from \(\operatorname{notsectorwinwittagreelcmperiod}(T, P, hP, hper, k, d, U, hU)\).
\end{proof}
\begin{theorem}
\label{ll:paper:sectorwinwittagree-of-sectorwiderowagree}
For every natural number \(k\) and natural number \(d\) and natural number \(U\) and natural number \(Q\), if \(2 \le U\), then if \(\operatorname{sectorwiderowagree}(k, d, U, Q)\), then \(\operatorname{sectorwinwittagree}(k, d, U, Q)\).
\end{theorem}
\begin{proof}
The goal follows from \(\operatorname{sectorwinwittagreeofsectorwiderowagree}(k, d, U, Q)\).
\end{proof}
\begin{theorem}
\label{ll:paper:sectorwinwittagree-of-sectorwiderightrowagree}
For every natural number \(T\) and natural number \(P\), if \(0 < P\), then if \(\operatorname{centerperiodic}(T, P)\), then for every natural number \(k\) and natural number \(d\) and natural number \(U\) and natural number \(m\), if \(2 \le U\), then if \(T \le d + U \cdot \operatorname{twopow}(k)\), then if \(\operatorname{sectorwiderightrowagree}(k, d, U, m \cdot P)\), then \(\operatorname{sectorwinwittagree}(k, d, U, m \cdot P)\).
\end{theorem}
\begin{proof}
Assume \(hP\), \(hper\).
The goal follows from \(\operatorname{sectorwinwittagreeofsectorwiderightrowagree}(T, P, hP, hper)\).
\end{proof}
\begin{theorem}
\label{ll:paper:not-sectorwiderightrowagree-lcmperiod}
For every natural number \(T\) and natural number \(P\), if \(0 < P\), then if \(\operatorname{centerperiodic}(T, P)\), then for every natural number \(k\) and natural number \(d\) and natural number \(U\), if \(2 \le U\), then if \(T \le d + U \cdot \operatorname{twopow}(k)\), then not \(\operatorname{sectorwiderightrowagree}(k, d, U, \operatorname{lcmperiod}(P, k))\).
\end{theorem}
\begin{proof}
Assume \(hP\), \(hper\), \(k\), \(d\), \(U\), \(hU\), \(ht\).
The goal follows from \(\operatorname{notsectorwiderightrowagreelcmperiod}(T, P, hP, hper, k, d, U, hU, ht)\).
\end{proof}
\begin{theorem}
\label{ll:paper:lhcenter-zero-eq-center}
For every natural number \(t\), \(\operatorname{lhcenter}(t, 0) = \operatorname{center}(t)\).
\end{theorem}
\begin{proof}
The goal follows from \(\operatorname{lhcenterzeroeqcenter}(t)\).
\end{proof}
\begin{theorem}
\label{ll:paper:lhcenter-self-eq-true}
For every natural number \(t\), \(\operatorname{lhcenter}(t, t) = \operatorname{true}\).
\end{theorem}
\begin{proof}
The goal follows from \(\operatorname{lhcenterselfeqtrue}(t)\).
\end{proof}
\begin{theorem}
\label{ll:paper:lhcenter-gt-eq-false}
For every natural number \(t\) and natural number \(p\), if \(t < p\), then \(\operatorname{lhcenter}(t, p) = \operatorname{false}\).
\end{theorem}
\begin{proof}
The goal follows from \(\operatorname{lhcentergteqfalse}(t, p)\).
\end{proof}
\begin{theorem}
\label{ll:paper:lhcenter-step-succ}
For every natural number \(t\) and natural number \(p\), \(\operatorname{lhcenter}(t + 1, p + 1) = \operatorname{bxor}(\operatorname{lhcenter}(t, p + 2), \operatorname{lhcenter}(t, p + 1) \operatorname{or} \operatorname{lhcenter}(t, p))\).
\end{theorem}
\begin{proof}
The goal follows from \(\operatorname{lhcenterstepsucc}(t, p)\).
\end{proof}
\begin{theorem}
\label{ll:paper:lhcenter-outward-step}
For every natural number \(t\) and natural number \(p\), \(\operatorname{lhcenter}(t, p + 2) = \operatorname{bxor}(\operatorname{lhcenter}(t + 1, p + 1), \operatorname{lhcenter}(t, p + 1) \operatorname{or} \operatorname{lhcenter}(t, p))\).
\end{theorem}
\begin{proof}
The goal follows from \(\operatorname{lhcenteroutwardstep}(t, p)\).
\end{proof}
\begin{theorem}
\label{ll:paper:center-eq-rhcenter-inward-zero}
For every natural number \(t\), \(\operatorname{center}(t) = \operatorname{bxor}(\operatorname{rhcenter}(t + 1, 1), \operatorname{rhcenter}(t, 1) \operatorname{or} \operatorname{rhcenter}(t, 2))\).
\end{theorem}
\begin{proof}
The goal follows from \(\operatorname{centereqrhcenterinwardzero}(t)\).
\end{proof}
\begin{theorem}
\label{ll:paper:periodic-rhcenter-step-one-eq-of-or-eq}
For every natural number \(T\) and natural number \(P\), if \(0 < P\), then if \(\operatorname{centerperiodic}(T, P)\), then for every natural number \(m\) and natural number \(t\), if \(T \le t\), then if \((\operatorname{rhcenter}(t + m \cdot P, 1) \operatorname{or} \operatorname{rhcenter}(t + m \cdot P, 2)) = (\operatorname{rhcenter}(t, 1) \operatorname{or} \operatorname{rhcenter}(t, 2))\), then \(\operatorname{rhcenter}(t + 1 + m \cdot P, 1) = \operatorname{rhcenter}(t + 1, 1)\).
\end{theorem}
\begin{proof}
Assume \(hP\), \(hper\).
The goal follows from \(\operatorname{periodicrhcentersteponeeqoforeq}(T, P, hP, hper)\).
\end{proof}
\begin{theorem}
\label{ll:paper:periodic-rhcenter-or-eq-of-step-one-eq}
For every natural number \(T\) and natural number \(P\), if \(0 < P\), then if \(\operatorname{centerperiodic}(T, P)\), then for every natural number \(m\) and natural number \(t\), if \(T \le t\), then if \(\operatorname{rhcenter}(t + 1 + m \cdot P, 1) = \operatorname{rhcenter}(t + 1, 1)\), then \((\operatorname{rhcenter}(t + m \cdot P, 1) \operatorname{or} \operatorname{rhcenter}(t + m \cdot P, 2)) = (\operatorname{rhcenter}(t, 1) \operatorname{or} \operatorname{rhcenter}(t, 2))\).
\end{theorem}
\begin{proof}
Assume \(hP\), \(hper\).
The goal follows from \(\operatorname{periodicrhcenteroreqofsteponeeq}(T, P, hP, hper)\).
\end{proof}
\begin{theorem}
\label{ll:paper:rhbandagree-steps}
For every natural number \(t\) and natural number \(Q\) and natural number \(p0\) and natural number \(L\) and natural number \(s\), if \(\operatorname{rhbandagree}(t, Q, p0, L + 2 \cdot s)\), then \(\operatorname{rhbandagree}(t + s, Q, p0 + s, L)\).
\end{theorem}
\begin{proof}
The goal follows from \(\operatorname{rhbandagreesteps}(t, Q, p0, L, s)\).
\end{proof}
\begin{theorem}
\label{ll:paper:periodic-rhbandagree-zero-of-one}
For every natural number \(T\) and natural number \(P\), if \(0 < P\), then if \(\operatorname{centerperiodic}(T, P)\), then for every natural number \(m\) and natural number \(t\) and natural number \(L\), if \(T \le t\), then if \(\operatorname{rhbandagree}(t, m \cdot P, 1, L)\), then \(\operatorname{rhbandagree}(t, m \cdot P, 0, L + 1)\).
\end{theorem}
\begin{proof}
Assume \(hP\), \(hper\).
The goal follows from \(\operatorname{periodicrhbandagreezeroofone}(T, P, hP, hper)\).
\end{proof}
\begin{theorem}
\label{ll:paper:periodic-rhbandagree-zero-step}
For every natural number \(T\) and natural number \(P\), if \(0 < P\), then if \(\operatorname{centerperiodic}(T, P)\), then for every natural number \(m\) and natural number \(t\) and natural number \(L\), if \(T \le t\), then if \(\operatorname{rhbandagree}(t, m \cdot P, 0, L + 1)\), then \(\operatorname{rhbandagree}(t + 1, m \cdot P, 0, L)\).
\end{theorem}
\begin{proof}
Assume \(hP\), \(hper\).
The goal follows from \(\operatorname{periodicrhbandagreezerostep}(T, P, hP, hper)\).
\end{proof}
\begin{theorem}
\label{ll:paper:periodic-rhbandagree-zero-steps}
For every natural number \(T\) and natural number \(P\), if \(0 < P\), then if \(\operatorname{centerperiodic}(T, P)\), then for every natural number \(m\) and natural number \(t\) and natural number \(L\) and natural number \(s\), if \(T \le t\), then if \(\operatorname{rhbandagree}(t, m \cdot P, 0, L + s)\), then \(\operatorname{rhbandagree}(t + s, m \cdot P, 0, L)\).
\end{theorem}
\begin{proof}
Assume \(hP\), \(hper\).
The goal follows from \(\operatorname{periodicrhbandagreezerosteps}(T, P, hP, hper)\).
\end{proof}
\begin{theorem}
\label{ll:paper:rhbandagree-extend-inward-steps}
For every natural number \(t\) and natural number \(Q\) and natural number \(p0\) and natural number \(S\) and natural number \(L\), if \(\operatorname{rhbandagree}(t, Q, p0 + S, L + 1)\), then if \(\operatorname{rhbandagree}(t + 1, Q, p0 + 1, S)\), then \(\operatorname{rhbandagree}(t, Q, p0, L + 1 + S)\).
\end{theorem}
\begin{proof}
The goal follows from \(\operatorname{rhbandagreeextendinwardsteps}(t, Q, p0, S, L)\).
\end{proof}
\begin{theorem}
\label{ll:paper:rhbandagree-zero-of-edge-and-succ}
For every natural number \(t\) and natural number \(Q\) and natural number \(S\) and natural number \(L\), if \(\operatorname{rhbandagree}(t, Q, S, L + 1)\), then if \(\operatorname{rhbandagree}(t + 1, Q, 1, S)\), then \(\operatorname{rhbandagree}(t, Q, 0, L + 1 + S)\).
\end{theorem}
\begin{proof}
The goal follows from \(\operatorname{rhbandagreezeroofedgeandsucc}(t, Q, S, L)\).
\end{proof}
\begin{theorem}
\label{ll:paper:periodic-rhbandagree-bounce-zero}
For every natural number \(T\) and natural number \(P\), if \(0 < P\), then if \(\operatorname{centerperiodic}(T, P)\), then for every natural number \(m\) and natural number \(t\) and natural number \(S\) and natural number \(L\), if \(T \le t\), then if \(\operatorname{rhbandagree}(t, m \cdot P, S + 1, L + 1)\), then if \(\operatorname{rhbandagree}(t + 1, m \cdot P, 2, S)\), then \(\operatorname{rhbandagree}(t, m \cdot P, 0, L + 2 + S)\).
\end{theorem}
\begin{proof}
Assume \(hP\), \(hper\).
The goal follows from \(\operatorname{periodicrhbandagreebouncezero}(T, P, hP, hper)\).
\end{proof}
\begin{theorem}
\label{ll:paper:periodic-rhbandagree-bounce-succ}
For every natural number \(T\) and natural number \(P\), if \(0 < P\), then if \(\operatorname{centerperiodic}(T, P)\), then for every natural number \(m\) and natural number \(t\) and natural number \(S\) and natural number \(L\), if \(T \le t\), then if \(\operatorname{rhbandagree}(t, m \cdot P, S + 1, L + 1)\), then if \(\operatorname{rhbandagree}(t + 1, m \cdot P, 2, S)\), then \(\operatorname{rhbandagree}(t + 1, m \cdot P, 0, L + 1 + S)\).
\end{theorem}
\begin{proof}
Assume \(hP\), \(hper\).
The goal follows from \(\operatorname{periodicrhbandagreebouncesucc}(T, P, hP, hper)\).
\end{proof}
\begin{theorem}
\label{ll:paper:rhcenter-one-periodic-of-sector-rhbandagree}
For every natural number \(T\) and natural number \(P\), if \(0 < P\), then if \(\operatorname{centerperiodic}(T, P)\), then for every natural number \(k\) and natural number \(d\) and natural number \(U\) and natural number \(m\), if \(T \le d + U \cdot \operatorname{twopow}(k)\), then if \(\operatorname{sectorrhbandagree}(k, d, U, m \cdot P)\), then for every natural number \(t\), if \(d + U \cdot \operatorname{twopow}(k) + 1 \le t\), then \(\operatorname{rhcenter}(t + m \cdot P, 1) = \operatorname{rhcenter}(t, 1)\).
\end{theorem}
\begin{proof}
Assume \(hP\), \(hper\).
The goal follows from \(\operatorname{rhcenteroneperiodicofsectorrhbandagree}(T, P, hP, hper)\).
\end{proof}
\begin{theorem}
\label{ll:paper:lhcenter-one-periodic-of-sector-rhbandagree}
For every natural number \(T\) and natural number \(P\), if \(0 < P\), then if \(\operatorname{centerperiodic}(T, P)\), then for every natural number \(k\) and natural number \(d\) and natural number \(U\) and natural number \(m\), if \(T \le d + U \cdot \operatorname{twopow}(k)\), then if \(\operatorname{sectorrhbandagree}(k, d, U, m \cdot P)\), then for every natural number \(t\), if \(d + U \cdot \operatorname{twopow}(k) + 1 \le t\), then \(\operatorname{lhcenter}(t + m \cdot P, 1) = \operatorname{lhcenter}(t, 1)\).
\end{theorem}
\begin{proof}
Assume \(hP\), \(hper\).
The goal follows from \(\operatorname{lhcenteroneperiodicofsectorrhbandagree}(T, P, hP, hper)\).
\end{proof}
\begin{theorem}
\label{ll:paper:rhcenter-diag-shift-twopow-eq}
For every natural number \(k\) and natural number \(t\) and natural number \(p\) and natural number \(v\), if \(p \le t\), then if \(t - p \le k\), then \(\operatorname{rhcenter}(t + v \cdot \operatorname{twopow}(k), p + v \cdot \operatorname{twopow}(k)) = \operatorname{rhcenter}(t, p)\).
\end{theorem}
\begin{proof}
The goal follows from \(\operatorname{rhcenterdiagshifttwopoweq}(k, t, p, v)\).
\end{proof}
\begin{theorem}
\label{ll:paper:rhcenter-sector-diag-eq-center}
For every natural number \(j\) and natural number \(k\) and natural number \(u\), if \(j \le k\), then \(\operatorname{rhcenter}(j + u \cdot \operatorname{twopow}(k), u \cdot \operatorname{twopow}(k)) = \operatorname{center}(j)\).
\end{theorem}
\begin{proof}
The goal follows from \(\operatorname{rhcentersectordiageqcenter}(j, k, u)\).
\end{proof}
\begin{theorem}
\label{ll:paper:rhcenter-scale-step}
For every natural number \(t\) and natural number \(k\) and natural number \(p\), if \(p + \operatorname{twopow}(\operatorname{succ}(k)) \le t\), then \(\operatorname{rhcenter}(t + \operatorname{twopow}(k), p + \operatorname{twopow}(k)) = \operatorname{bxor}(\operatorname{bxor}(\operatorname{bxor}(\operatorname{rhcenter}(t, p), \operatorname{rhcenter}(t, p + \operatorname{twopow}(k))), \operatorname{rhcenter}(t, p + \operatorname{twopow}(\operatorname{succ}(k)))), \operatorname{wittbitbool}(\operatorname{wittwin}(t, \operatorname{twopow}(k)), t - p))\).
\end{theorem}
\begin{proof}
The goal follows from \(\operatorname{rhcenterscalestep}(t, k, p)\).
\end{proof}
\begin{theorem}
\label{ll:paper:rhcenter-scale-inward-step}
For every natural number \(t\) and natural number \(k\) and natural number \(p\), if \(p + \operatorname{twopow}(\operatorname{succ}(k)) \le t\), then \(\operatorname{rhcenter}(t, p) = \operatorname{bxor}(\operatorname{bxor}(\operatorname{bxor}(\operatorname{rhcenter}(t + \operatorname{twopow}(k), p + \operatorname{twopow}(k)), \operatorname{rhcenter}(t, p + \operatorname{twopow}(k))), \operatorname{rhcenter}(t, p + \operatorname{twopow}(\operatorname{succ}(k)))), \operatorname{wittbitbool}(\operatorname{wittwin}(t, \operatorname{twopow}(k)), t - p))\).
\end{theorem}
\begin{proof}
The goal follows from \(\operatorname{rhcenterscaleinwardstep}(t, k, p)\).
\end{proof}
\begin{theorem}
\label{ll:paper:rhbandagree-mono}
For every natural number \(t\) and natural number \(Q\) and natural number \(p0\) and natural number \(L\) and natural number \(p0p\) and natural number \(Lp\), if \(p0 \le p0p\), then if \(p0p + Lp \le p0 + L\), then if \(\operatorname{rhbandagree}(t, Q, p0, L)\), then \(\operatorname{rhbandagree}(t, Q, p0p, Lp)\).
\end{theorem}
\begin{proof}
The goal follows from \(\operatorname{rhbandagreemono}(t, Q, p0, L, p0p, Lp)\).
\end{proof}
\begin{theorem}
\label{ll:paper:rhcolagree-mono}
For every natural number \(t\) and natural number \(Q\) and natural number \(p0\) and natural number \(S\) and natural number \(s0\) and natural number \(Sp\), if \(s0 + Sp \le S\), then if \(\operatorname{rhcolagree}(t, Q, p0, S)\), then \(\operatorname{rhcolagree}(t + s0, Q, p0, Sp)\).
\end{theorem}
\begin{proof}
The goal follows from \(\operatorname{rhcolagreemono}(t, Q, p0, S, s0, Sp)\).
\end{proof}
\begin{theorem}
\label{ll:paper:rhoragree-mono}
For every natural number \(t\) and natural number \(Q\) and natural number \(p0\) and natural number \(S\) and natural number \(s0\) and natural number \(Sp\), if \(s0 + Sp \le S\), then if \(\operatorname{rhoragree}(t, Q, p0, S)\), then \(\operatorname{rhoragree}(t + s0, Q, p0, Sp)\).
\end{theorem}
\begin{proof}
The goal follows from \(\operatorname{rhoragreemono}(t, Q, p0, S, s0, Sp)\).
\end{proof}
\begin{theorem}
\label{ll:paper:rhoragree-of-col-agree}
For every natural number \(t\) and natural number \(Q\) and natural number \(p0\) and natural number \(S\), if \(\operatorname{rhcolagree}(t, Q, p0, S)\), then if for every natural number \(s\), if \(s < S\), then \(\operatorname{rhcenter}(t + (s + 1) + Q, p0 + 1) = \operatorname{rhcenter}(t + (s + 1), p0 + 1)\), then \(\operatorname{rhoragree}(t, Q, p0 + 1, S)\).
\end{theorem}
\begin{proof}
The goal follows from \(\operatorname{rhoragreeofcolagree}(t, Q, p0, S)\).
\end{proof}
\begin{theorem}
\label{ll:paper:periodic-rhoragree-one-of-col-one}
For every natural number \(T\) and natural number \(P\), if \(0 < P\), then if \(\operatorname{centerperiodic}(T, P)\), then for every natural number \(m\) and natural number \(t\) and natural number \(S\), if \(T \le t\), then if for every natural number \(s\), if \(s < S\), then \(\operatorname{rhcenter}(t + (s + 1) + m \cdot P, 1) = \operatorname{rhcenter}(t + (s + 1), 1)\), then \(\operatorname{rhoragree}(t, m \cdot P, 1, S)\).
\end{theorem}
\begin{proof}
Assume \(hP\), \(hper\).
The goal follows from \(\operatorname{periodicrhoragreeoneofcolone}(T, P, hP, hper)\).
\end{proof}
\begin{theorem}
\label{ll:paper:periodic-col-one-of-rhoragree-one}
For every natural number \(T\) and natural number \(P\), if \(0 < P\), then if \(\operatorname{centerperiodic}(T, P)\), then for every natural number \(m\) and natural number \(t\) and natural number \(S\), if \(T \le t\), then if \(\operatorname{rhoragree}(t, m \cdot P, 1, S)\), then for every natural number \(s\), if \(s < S\), then \(\operatorname{rhcenter}(t + (s + 1) + m \cdot P, 1) = \operatorname{rhcenter}(t + (s + 1), 1)\).
\end{theorem}
\begin{proof}
Assume \(hP\), \(hper\).
The goal follows from \(\operatorname{periodiccoloneofrhoragreeone}(T, P, hP, hper)\).
\end{proof}
\begin{theorem}
\label{ll:paper:rhcenter-inward-determine-step}
For every natural number \(t\) and natural number \(Q\) and natural number \(p0\) and natural number \(S\), if \(\operatorname{rhcolagree}(t, Q, p0 + 1, S + 1)\), then if \(\operatorname{rhoragree}(t, Q, p0 + 1, S + 1)\), then \(\operatorname{rhcolagree}(t, Q, p0, S)\) and \(\operatorname{rhoragree}(t, Q, p0, S)\).
\end{theorem}
\begin{proof}
The goal follows from \(\operatorname{rhcenterinwarddeterminestep}(t, Q, p0, S)\).
\end{proof}
\begin{theorem}
\label{ll:paper:rhcenter-inward-determine}
For every natural number \(t\) and natural number \(Q\) and natural number \(p0\) and natural number \(S\) and natural number \(K\), if \(\operatorname{rhcolagree}(t, Q, p0 + K, S + K)\), then if \(\operatorname{rhoragree}(t, Q, p0 + K, S + K)\), then \(\operatorname{rhcolagree}(t, Q, p0, S)\) and \(\operatorname{rhoragree}(t, Q, p0, S)\).
\end{theorem}
\begin{proof}
The goal follows from \(\operatorname{rhcenterinwarddetermine}(t, Q, p0, S, K)\).
\end{proof}
\begin{theorem}
\label{ll:paper:rhbandagree-of-inward-determine}
For every natural number \(t\) and natural number \(Q\) and natural number \(p0\) and natural number \(L\), if \(\operatorname{rhcolagree}(t, Q, p0 + L, L)\), then if \(\operatorname{rhoragree}(t, Q, p0 + L, L)\), then \(\operatorname{rhbandagree}(t, Q, p0, L)\).
\end{theorem}
\begin{proof}
The goal follows from \(\operatorname{rhbandagreeofinwarddetermine}(t, Q, p0, L)\).
\end{proof}
\begin{theorem}
\label{ll:paper:rhbandagree-outside-cone}
For every natural number \(t\) and natural number \(Q\) and natural number \(p0\) and natural number \(L\), if \(t + Q < p0\), then \(\operatorname{rhbandagree}(t, Q, p0, L)\).
\end{theorem}
\begin{proof}
The goal follows from \(\operatorname{rhbandagreeoutsidecone}(t, Q, p0, L)\).
\end{proof}
\begin{theorem}
\label{ll:paper:rhbandagree-scale-step}
For every natural number \(t\) and natural number \(k\) and natural number \(Q\) and natural number \(p0\) and natural number \(L\), if \(p0 + L + \operatorname{twopow}(\operatorname{succ}(k)) \le t\), then if \(\operatorname{rhbandagree}(t, Q, p0, L)\), then if \(\operatorname{rhbandagree}(t, Q, p0 + \operatorname{twopow}(k), L)\), then if \(\operatorname{rhbandagree}(t, Q, p0 + \operatorname{twopow}(\operatorname{succ}(k)), L)\), then if \(\operatorname{winwittagree}(t, k, Q, t - (p0 + L), L)\), then \(\operatorname{rhbandagree}(t + \operatorname{twopow}(k), Q, p0 + \operatorname{twopow}(k), L)\).
\end{theorem}
\begin{proof}
The goal follows from \(\operatorname{rhbandagreescalestep}(t, k, Q, p0, L)\).
\end{proof}
\begin{theorem}
\label{ll:paper:rhbandagree-scale-inward-step}
For every natural number \(t\) and natural number \(k\) and natural number \(Q\) and natural number \(p0\) and natural number \(L\), if \(p0 + L + \operatorname{twopow}(\operatorname{succ}(k)) \le t\), then if \(\operatorname{rhbandagree}(t + \operatorname{twopow}(k), Q, p0 + \operatorname{twopow}(k), L)\), then if \(\operatorname{rhbandagree}(t, Q, p0 + \operatorname{twopow}(k), L)\), then if \(\operatorname{rhbandagree}(t, Q, p0 + \operatorname{twopow}(\operatorname{succ}(k)), L)\), then if \(\operatorname{winwittagree}(t, k, Q, t - (p0 + L), L)\), then \(\operatorname{rhbandagree}(t, Q, p0, L)\).
\end{theorem}
\begin{proof}
The goal follows from \(\operatorname{rhbandagreescaleinwardstep}(t, k, Q, p0, L)\).
\end{proof}
\begin{theorem}
\label{ll:paper:winwittagree-of-rhbandagree-scale}
For every natural number \(t\) and natural number \(k\) and natural number \(Q\) and natural number \(p0\) and natural number \(L\), if \(p0 + L + \operatorname{twopow}(\operatorname{succ}(k)) \le t\), then if \(\operatorname{rhbandagree}(t, Q, p0, L)\), then if \(\operatorname{rhbandagree}(t, Q, p0 + \operatorname{twopow}(k), L)\), then if \(\operatorname{rhbandagree}(t, Q, p0 + \operatorname{twopow}(\operatorname{succ}(k)), L)\), then if \(\operatorname{rhbandagree}(t + \operatorname{twopow}(k), Q, p0 + \operatorname{twopow}(k), L)\), then \(\operatorname{winwittagree}(t, k, Q, t - (p0 + L), L)\).
\end{theorem}
\begin{proof}
The goal follows from \(\operatorname{winwittagreeofrhbandagreescale}(t, k, Q, p0, L)\).
\end{proof}
\begin{theorem}
\label{ll:paper:sectorrhbandagree-of-inward-determine}
For every natural number \(k\) and natural number \(d\) and natural number \(U\) and natural number \(Q\), if for every natural number \(u\), if \(U \le u\), then \(\operatorname{rhcolagree}(d + u \cdot \operatorname{twopow}(k) + 1, Q, \operatorname{twopow}(k) + 1, \operatorname{twopow}(k) + 1)\), then if for every natural number \(u\), if \(U \le u\), then \(\operatorname{rhoragree}(d + u \cdot \operatorname{twopow}(k) + 1, Q, \operatorname{twopow}(k) + 1, \operatorname{twopow}(k) + 1)\), then \(\operatorname{sectorrhbandagree}(k, d, U, Q)\).
\end{theorem}
\begin{proof}
The goal follows from \(\operatorname{sectorrhbandagreeofinwarddetermine}(k, d, U, Q)\).
\end{proof}
\begin{theorem}
\label{ll:paper:not-sector-inward-determine-lcmperiod}
For every natural number \(T\) and natural number \(P\), if \(0 < P\), then if \(\operatorname{centerperiodic}(T, P)\), then for every natural number \(k\) and natural number \(d\) and natural number \(U\), if \(1 \le U\), then if \(1 \le d\), then if \(T \le d + U \cdot \operatorname{twopow}(k)\), then if for every natural number \(u\), if \(U \le u\), then \(\operatorname{rhcolagree}(d + u \cdot \operatorname{twopow}(k) + 1, \operatorname{lcmperiod}(P, k), \operatorname{twopow}(k) + 1, \operatorname{twopow}(k) + 1)\), then not for every natural number \(u\), if \(U \le u\), then \(\operatorname{rhoragree}(d + u \cdot \operatorname{twopow}(k) + 1, \operatorname{lcmperiod}(P, k), \operatorname{twopow}(k) + 1, \operatorname{twopow}(k) + 1)\).
\end{theorem}
\begin{proof}
Assume \(hP\), \(hper\), \(k\), \(d\), \(U\), \(hU\), \(hd\), \(ht\), \(hcol\).
The goal follows from \(\operatorname{notsectorinwarddeterminelcmperiod}(T, P, hP, hper, k, d, U, hU, hd, ht, hcol)\).
\end{proof}
\begin{theorem}
\label{ll:paper:rhcenter-scale-step-le}
For every natural number \(t\) and natural number \(k\) and natural number \(p\), if \(p \le t\), then \(\operatorname{rhcenter}(t + \operatorname{twopow}(k), p + \operatorname{twopow}(k)) = \operatorname{bxor}(\operatorname{bxor}(\operatorname{bxor}(\operatorname{rhcenter}(t, p), \operatorname{rhcenter}(t, p + \operatorname{twopow}(k))), \operatorname{rhcenter}(t, p + \operatorname{twopow}(\operatorname{succ}(k)))), \operatorname{wittbitbool}(\operatorname{wittwin}(t, \operatorname{twopow}(k)), t - p))\).
\end{theorem}
\begin{proof}
The goal follows from \(\operatorname{rhcenterscalesteple}(t, k, p)\).
\end{proof}
\begin{theorem}
\label{ll:paper:rhcenter-scale-inward-step-le}
For every natural number \(t\) and natural number \(k\) and natural number \(p\), if \(p \le t\), then \(\operatorname{rhcenter}(t, p) = \operatorname{bxor}(\operatorname{bxor}(\operatorname{bxor}(\operatorname{rhcenter}(t + \operatorname{twopow}(k), p + \operatorname{twopow}(k)), \operatorname{rhcenter}(t, p + \operatorname{twopow}(k))), \operatorname{rhcenter}(t, p + \operatorname{twopow}(\operatorname{succ}(k)))), \operatorname{wittbitbool}(\operatorname{wittwin}(t, \operatorname{twopow}(k)), t - p))\).
\end{theorem}
\begin{proof}
The goal follows from \(\operatorname{rhcenterscaleinwardsteple}(t, k, p)\).
\end{proof}
\begin{theorem}
\label{ll:paper:rhbandagree-scale-step-le}
For every natural number \(t\) and natural number \(k\) and natural number \(Q\) and natural number \(p0\) and natural number \(L\), if \(p0 + L \le t\), then if \(\operatorname{rhbandagree}(t, Q, p0, L)\), then if \(\operatorname{rhbandagree}(t, Q, p0 + \operatorname{twopow}(k), L)\), then if \(\operatorname{rhbandagree}(t, Q, p0 + \operatorname{twopow}(\operatorname{succ}(k)), L)\), then if \(\operatorname{winwittagree}(t, k, Q, t - (p0 + L), L)\), then \(\operatorname{rhbandagree}(t + \operatorname{twopow}(k), Q, p0 + \operatorname{twopow}(k), L)\).
\end{theorem}
\begin{proof}
The goal follows from \(\operatorname{rhbandagreescalesteple}(t, k, Q, p0, L)\).
\end{proof}
\begin{theorem}
\label{ll:paper:rhbandagree-scale-inward-step-le}
For every natural number \(t\) and natural number \(k\) and natural number \(Q\) and natural number \(p0\) and natural number \(L\), if \(p0 + L \le t\), then if \(\operatorname{rhbandagree}(t + \operatorname{twopow}(k), Q, p0 + \operatorname{twopow}(k), L)\), then if \(\operatorname{rhbandagree}(t, Q, p0 + \operatorname{twopow}(k), L)\), then if \(\operatorname{rhbandagree}(t, Q, p0 + \operatorname{twopow}(\operatorname{succ}(k)), L)\), then if \(\operatorname{winwittagree}(t, k, Q, t - (p0 + L), L)\), then \(\operatorname{rhbandagree}(t, Q, p0, L)\).
\end{theorem}
\begin{proof}
The goal follows from \(\operatorname{rhbandagreescaleinwardsteple}(t, k, Q, p0, L)\).
\end{proof}
\begin{theorem}
\label{ll:paper:winwittagree-of-rhbandagree-scale-le}
For every natural number \(t\) and natural number \(k\) and natural number \(Q\) and natural number \(p0\) and natural number \(L\), if \(p0 + L \le t\), then if \(\operatorname{rhbandagree}(t, Q, p0, L)\), then if \(\operatorname{rhbandagree}(t, Q, p0 + \operatorname{twopow}(k), L)\), then if \(\operatorname{rhbandagree}(t, Q, p0 + \operatorname{twopow}(\operatorname{succ}(k)), L)\), then if \(\operatorname{rhbandagree}(t + \operatorname{twopow}(k), Q, p0 + \operatorname{twopow}(k), L)\), then \(\operatorname{winwittagree}(t, k, Q, t - (p0 + L), L)\).
\end{theorem}
\begin{proof}
The goal follows from \(\operatorname{winwittagreeofrhbandagreescalele}(t, k, Q, p0, L)\).
\end{proof}
\begin{theorem}
\label{ll:paper:rhbandagree-scale-inward-steps}
For every natural number \(t\) and natural number \(k\) and natural number \(Q\) and natural number \(p0\) and natural number \(L\) and natural number \(M\), if \(p0 + M \cdot \operatorname{twopow}(k) + L \le t\), then if \(\operatorname{rhbandagree}(t, Q, p0 + M \cdot \operatorname{twopow}(k), L)\), then if \(\operatorname{rhbandagree}(t, Q, p0 + (M + 1) \cdot \operatorname{twopow}(k), L)\), then if for every natural number \(m\), if \(m < M\), then \(\operatorname{rhbandagree}(t + \operatorname{twopow}(k), Q, p0 + (m + 1) \cdot \operatorname{twopow}(k), L)\), then if for every natural number \(m\), if \(m < M\), then \(\operatorname{winwittagree}(t, k, Q, t - (p0 + m \cdot \operatorname{twopow}(k) + L), L)\), then \(\operatorname{rhbandagree}(t, Q, p0, L)\).
\end{theorem}
\begin{proof}
The goal follows from \(\operatorname{rhbandagreescaleinwardsteps}(t, k, Q, p0, L, M)\).
\end{proof}
\begin{theorem}
\label{ll:paper:rhcenter-inward-of-black-left}
For every natural number \(t\) and natural number \(p\), if \(\operatorname{rhcenter}(t, p + 1) = \operatorname{true}\), then \(\operatorname{rhcenter}(t, p) = \operatorname{not}(\operatorname{rhcenter}(t + 1, p + 1))\).
\end{theorem}
\begin{proof}
The goal follows from \(\operatorname{rhcenterinwardofblackleft}(t, p)\).
\end{proof}
\begin{theorem}
\label{ll:paper:rhcenter-inward-of-black-right}
For every natural number \(t\) and natural number \(p\), if \(\operatorname{rhcenter}(t, p + 2) = \operatorname{true}\), then \(\operatorname{rhcenter}(t, p) = \operatorname{not}(\operatorname{rhcenter}(t + 1, p + 1))\).
\end{theorem}
\begin{proof}
The goal follows from \(\operatorname{rhcenterinwardofblackright}(t, p)\).
\end{proof}
\begin{theorem}
\label{ll:paper:rhcenter-agree-of-black-left}
For every natural number \(t\) and natural number \(Q\) and natural number \(p\), if \(\operatorname{rhcenter}(t + Q, p + 1) = \operatorname{true}\), then if \(\operatorname{rhcenter}(t, p + 1) = \operatorname{true}\), then if \(\operatorname{rhcenter}(t + 1 + Q, p + 1) = \operatorname{rhcenter}(t + 1, p + 1)\), then \(\operatorname{rhcenter}(t + Q, p) = \operatorname{rhcenter}(t, p)\).
\end{theorem}
\begin{proof}
The goal follows from \(\operatorname{rhcenteragreeofblackleft}(t, Q, p)\).
\end{proof}
\begin{theorem}
\label{ll:paper:rhcenter-white-pair-of-diag-eq}
For every natural number \(t\) and natural number \(p\), if \(\operatorname{rhcenter}(t + 1, p + 1) = \operatorname{rhcenter}(t, p)\), then \(\operatorname{rhcenter}(t, p + 1) = \operatorname{false}\) and \(\operatorname{rhcenter}(t, p + 2) = \operatorname{false}\).
\end{theorem}
\begin{proof}
The goal follows from \(\operatorname{rhcenterwhitepairofdiageq}(t, p)\).
\end{proof}
\begin{theorem}
\label{ll:paper:rhcenter-diag-eq-of-white-pair}
For every natural number \(t\) and natural number \(p\), if \(\operatorname{rhcenter}(t, p + 1) = \operatorname{false}\), then if \(\operatorname{rhcenter}(t, p + 2) = \operatorname{false}\), then \(\operatorname{rhcenter}(t + 1, p + 1) = \operatorname{rhcenter}(t, p)\).
\end{theorem}
\begin{proof}
The goal follows from \(\operatorname{rhcenterdiageqofwhitepair}(t, p)\).
\end{proof}
\begin{theorem}
\label{ll:paper:rhcenter-diag-shift-twopow-eq-of-threshold}
For every natural number \(k\) and natural number \(J\) and natural number \(p\) and natural number \(s\) and natural number \(u\), if \(\operatorname{rightthreshold}(k, J)\), then if \(s \le J\), then \(\operatorname{rhcenter}(p + s + u \cdot \operatorname{twopow}(k), p + u \cdot \operatorname{twopow}(k)) = \operatorname{rhcenter}(p + s, p)\).
\end{theorem}
\begin{proof}
The goal follows from \(\operatorname{rhcenterdiagshifttwopoweqofthreshold}(k, J, p, s, u)\).
\end{proof}
\begin{theorem}
\label{ll:paper:rhcenter-sector-diag-eq-center-of-threshold}
For every natural number \(k\) and natural number \(J\) and natural number \(s\) and natural number \(u\), if \(\operatorname{rightthreshold}(k, J)\), then if \(s \le J\), then \(\operatorname{rhcenter}(s + u \cdot \operatorname{twopow}(k), u \cdot \operatorname{twopow}(k)) = \operatorname{center}(s)\).
\end{theorem}
\begin{proof}
The goal follows from \(\operatorname{rhcentersectordiageqcenterofthreshold}(k, J, s, u)\).
\end{proof}
\begin{theorem}
\label{ll:paper:rhcenter-zero-gap-of-threshold}
For every natural number \(k\) and natural number \(J\) and natural number \(u\) and natural number \(j\), if \(\operatorname{rightthreshold}(k, J)\), then if \(1 \le j\), then if \(j \le J\), then if \(j \le u \cdot \operatorname{twopow}(k)\), then \(\operatorname{rhcenter}(u \cdot \operatorname{twopow}(k), u \cdot \operatorname{twopow}(k) - j) = \operatorname{false}\).
\end{theorem}
\begin{proof}
The goal follows from \(\operatorname{rhcenterzerogapofthreshold}(k, J, u, j)\).
\end{proof}
\begin{theorem}
\label{ll:paper:rhcenter-last-one-of-threshold}
For every natural number \(k\) and natural number \(J\) and natural number \(u\) and natural number \(j\), if \(\operatorname{rightthreshold}(k, J)\), then if \(j \le J\), then if \(j \le (u + 1) \cdot \operatorname{twopow}(k) - 1\), then \(\operatorname{rhcenter}((u + 1) \cdot \operatorname{twopow}(k) - 1, (u + 1) \cdot \operatorname{twopow}(k) - 1 - j) = \operatorname{true}\).
\end{theorem}
\begin{proof}
The goal follows from \(\operatorname{rhcenterlastoneofthreshold}(k, J, u, j)\).
\end{proof}
\begin{theorem}
\label{ll:paper:rhcenter-scale-step-uncond}
For every natural number \(t\) and natural number \(k\) and natural number \(p\), \(\operatorname{rhcenter}(t + \operatorname{twopow}(k), p + \operatorname{twopow}(k)) = \operatorname{bxor}(\operatorname{bxor}(\operatorname{bxor}(\operatorname{rhcenter}(t, p), \operatorname{rhcenter}(t, p + \operatorname{twopow}(k))), \operatorname{rhcenter}(t, p + \operatorname{twopow}(\operatorname{succ}(k)))), \operatorname{wittbitbool}(\operatorname{wittwin}(t, \operatorname{twopow}(k)), t - p))\).
\end{theorem}
\begin{proof}
The goal follows from \(\operatorname{rhcenterscalestepuncond}(t, k, p)\).
\end{proof}
\begin{theorem}
\label{ll:paper:rhcenter-scale-inward-step-uncond}
For every natural number \(t\) and natural number \(k\) and natural number \(p\), \(\operatorname{rhcenter}(t, p) = \operatorname{bxor}(\operatorname{bxor}(\operatorname{bxor}(\operatorname{rhcenter}(t + \operatorname{twopow}(k), p + \operatorname{twopow}(k)), \operatorname{rhcenter}(t, p + \operatorname{twopow}(k))), \operatorname{rhcenter}(t, p + \operatorname{twopow}(\operatorname{succ}(k)))), \operatorname{wittbitbool}(\operatorname{wittwin}(t, \operatorname{twopow}(k)), t - p))\).
\end{theorem}
\begin{proof}
The goal follows from \(\operatorname{rhcenterscaleinwardstepuncond}(t, k, p)\).
\end{proof}
\begin{theorem}
\label{ll:paper:rhwinwittagree-of-winwittagree}
For every natural number \(t\) and natural number \(k\) and natural number \(Q\) and natural number \(p0\) and natural number \(L\), if \(p0 + L \le t\), then if \(\operatorname{winwittagree}(t, k, Q, t - (p0 + L), L)\), then \(\operatorname{rhwinwittagree}(t, k, Q, p0, L)\).
\end{theorem}
\begin{proof}
The goal follows from \(\operatorname{rhwinwittagreeofwinwittagree}(t, k, Q, p0, L)\).
\end{proof}
\begin{theorem}
\label{ll:paper:rhwinwittagree-outside-cone}
For every natural number \(t\) and natural number \(k\) and natural number \(Q\) and natural number \(p0\) and natural number \(L\), if \(t + Q \le p0\), then \(\operatorname{rhwinwittagree}(t, k, Q, p0, L)\).
\end{theorem}
\begin{proof}
The goal follows from \(\operatorname{rhwinwittagreeoutsidecone}(t, k, Q, p0, L)\).
\end{proof}
\begin{theorem}
\label{ll:paper:rhbandagree-scale-inward-step-uncond}
For every natural number \(t\) and natural number \(k\) and natural number \(Q\) and natural number \(p0\) and natural number \(L\), if \(\operatorname{rhbandagree}(t + \operatorname{twopow}(k), Q, p0 + \operatorname{twopow}(k), L)\), then if \(\operatorname{rhbandagree}(t, Q, p0 + \operatorname{twopow}(k), L)\), then if \(\operatorname{rhbandagree}(t, Q, p0 + \operatorname{twopow}(\operatorname{succ}(k)), L)\), then if \(\operatorname{rhwinwittagree}(t, k, Q, p0, L)\), then \(\operatorname{rhbandagree}(t, Q, p0, L)\).
\end{theorem}
\begin{proof}
The goal follows from \(\operatorname{rhbandagreescaleinwardstepuncond}(t, k, Q, p0, L)\).
\end{proof}
\begin{theorem}
\label{ll:paper:rhbandagree-scale-inward-steps-uncond}
For every natural number \(t\) and natural number \(k\) and natural number \(Q\) and natural number \(p0\) and natural number \(L\) and natural number \(M\), if \(\operatorname{rhbandagree}(t, Q, p0 + M \cdot \operatorname{twopow}(k), L)\), then if \(\operatorname{rhbandagree}(t, Q, p0 + (M + 1) \cdot \operatorname{twopow}(k), L)\), then if for every natural number \(m\), if \(m < M\), then \(\operatorname{rhbandagree}(t + \operatorname{twopow}(k), Q, p0 + (m + 1) \cdot \operatorname{twopow}(k), L)\), then if for every natural number \(m\), if \(m < M\), then \(\operatorname{rhwinwittagree}(t, k, Q, p0 + m \cdot \operatorname{twopow}(k), L)\), then \(\operatorname{rhbandagree}(t, Q, p0, L)\).
\end{theorem}
\begin{proof}
The goal follows from \(\operatorname{rhbandagreescaleinwardstepsuncond}(t, k, Q, p0, L, M)\).
\end{proof}
\begin{theorem}
\label{ll:paper:rhbandagree-of-scale-inward-cone}
For every natural number \(t\) and natural number \(k\) and natural number \(Q\) and natural number \(p0\) and natural number \(L\) and natural number \(M\), if \(t + Q < p0 + M \cdot \operatorname{twopow}(k)\), then if for every natural number \(m\), if \(m < M\), then \(\operatorname{rhbandagree}(t + \operatorname{twopow}(k), Q, p0 + (m + 1) \cdot \operatorname{twopow}(k), L)\), then if for every natural number \(m\), if \(m < M\), then \(\operatorname{rhwinwittagree}(t, k, Q, p0 + m \cdot \operatorname{twopow}(k), L)\), then \(\operatorname{rhbandagree}(t, Q, p0, L)\).
\end{theorem}
\begin{proof}
The goal follows from \(\operatorname{rhbandagreeofscaleinwardcone}(t, k, Q, p0, L, M)\).
\end{proof}
\begin{theorem}
\label{ll:paper:rhcenter-diag-diff-eq-or}
For every natural number \(t\) and natural number \(p\), \(\operatorname{bxor}(\operatorname{rhcenter}(t + 1, p + 1), \operatorname{rhcenter}(t, p)) = (\operatorname{rhcenter}(t, p + 1) \operatorname{or} \operatorname{rhcenter}(t, p + 2))\).
\end{theorem}
\begin{proof}
The goal follows from \(\operatorname{rhcenterdiagdiffeqor}(t, p)\).
\end{proof}
\begin{theorem}
\label{ll:paper:rhcenter-defect-diag-step}
For every natural number \(t\) and natural number \(Q\) and natural number \(p\), \(\operatorname{bxor}(\operatorname{bxor}(\operatorname{rhcenter}(t + 1 + Q, p + 1), \operatorname{rhcenter}(t + 1, p + 1)), \operatorname{bxor}(\operatorname{rhcenter}(t + Q, p), \operatorname{rhcenter}(t, p))) = \operatorname{bxor}(\operatorname{rhcenter}(t + Q, p + 1) \operatorname{or} \operatorname{rhcenter}(t + Q, p + 2), \operatorname{rhcenter}(t, p + 1) \operatorname{or} \operatorname{rhcenter}(t, p + 2))\).
\end{theorem}
\begin{proof}
The goal follows from \(\operatorname{rhcenterdefectdiagstep}(t, Q, p)\).
\end{proof}
\begin{theorem}
\label{ll:paper:rhcenter-defect-diag-step-of-col-eq}
For every natural number \(t\) and natural number \(Q\) and natural number \(p\), if \(\operatorname{rhcenter}(t + Q, p + 1) = \operatorname{rhcenter}(t, p + 1)\), then \(\operatorname{bxor}(\operatorname{bxor}(\operatorname{rhcenter}(t + 1 + Q, p + 1), \operatorname{rhcenter}(t + 1, p + 1)), \operatorname{bxor}(\operatorname{rhcenter}(t + Q, p), \operatorname{rhcenter}(t, p))) = \operatorname{not}(\operatorname{rhcenter}(t, p + 1)) \operatorname{and} \operatorname{bxor}(\operatorname{rhcenter}(t + Q, p + 2), \operatorname{rhcenter}(t, p + 2))\).
\end{theorem}
\begin{proof}
The goal follows from \(\operatorname{rhcenterdefectdiagstepofcoleq}(t, Q, p)\).
\end{proof}
\begin{theorem}
\label{ll:paper:rhcenter-defect-inward-of-col-eq-two}
For every natural number \(t\) and natural number \(Q\) and natural number \(p\), if \(\operatorname{rhcenter}(t + Q, p + 1) = \operatorname{rhcenter}(t, p + 1)\), then if \(\operatorname{rhcenter}(t + 1 + Q, p + 1) = \operatorname{rhcenter}(t + 1, p + 1)\), then \(\operatorname{bxor}(\operatorname{rhcenter}(t + Q, p), \operatorname{rhcenter}(t, p)) = \operatorname{not}(\operatorname{rhcenter}(t, p + 1)) \operatorname{and} \operatorname{bxor}(\operatorname{rhcenter}(t + Q, p + 2), \operatorname{rhcenter}(t, p + 2))\).
\end{theorem}
\begin{proof}
The goal follows from \(\operatorname{rhcenterdefectinwardofcoleqtwo}(t, Q, p)\).
\end{proof}
\begin{theorem}
\label{ll:paper:rhcenter-outward-of-white-mid}
For every natural number \(t\) and natural number \(p\), if \(\operatorname{rhcenter}(t, p + 1) = \operatorname{false}\), then \(\operatorname{rhcenter}(t, p + 2) = \operatorname{bxor}(\operatorname{rhcenter}(t + 1, p + 1), \operatorname{rhcenter}(t, p))\).
\end{theorem}
\begin{proof}
The goal follows from \(\operatorname{rhcenteroutwardofwhitemid}(t, p)\).
\end{proof}
\begin{theorem}
\label{ll:paper:rhcenter-agree-outward-of-white-mid}
For every natural number \(t\) and natural number \(Q\) and natural number \(p\), if \(\operatorname{rhcenter}(t, p + 1) = \operatorname{false}\), then if \(\operatorname{rhcenter}(t + Q, p) = \operatorname{rhcenter}(t, p)\), then if \(\operatorname{rhcenter}(t + Q, p + 1) = \operatorname{rhcenter}(t, p + 1)\), then if \(\operatorname{rhcenter}(t + 1 + Q, p + 1) = \operatorname{rhcenter}(t + 1, p + 1)\), then \(\operatorname{rhcenter}(t + Q, p + 2) = \operatorname{rhcenter}(t, p + 2)\).
\end{theorem}
\begin{proof}
The goal follows from \(\operatorname{rhcenteragreeoutwardofwhitemid}(t, Q, p)\).
\end{proof}
\begin{theorem}
\label{ll:paper:rhcenter-defect-inward-of-white-mid}
For every natural number \(t\) and natural number \(Q\) and natural number \(p\), if \(\operatorname{rhcenter}(t, p + 1) = \operatorname{false}\), then if \(\operatorname{rhcenter}(t + Q, p + 1) = \operatorname{rhcenter}(t, p + 1)\), then if \(\operatorname{rhcenter}(t + 1 + Q, p + 1) = \operatorname{rhcenter}(t + 1, p + 1)\), then \(\operatorname{bxor}(\operatorname{rhcenter}(t + Q, p), \operatorname{rhcenter}(t, p)) = \operatorname{bxor}(\operatorname{rhcenter}(t + Q, p + 2), \operatorname{rhcenter}(t, p + 2))\).
\end{theorem}
\begin{proof}
The goal follows from \(\operatorname{rhcenterdefectinwardofwhitemid}(t, Q, p)\).
\end{proof}
\begin{theorem}
\label{ll:paper:rhcenter-agree-inward-of-col-eq-two}
For every natural number \(t\) and natural number \(Q\) and natural number \(p\), if \(\operatorname{rhcenter}(t + Q, p + 1) = \operatorname{rhcenter}(t, p + 1)\), then if \(\operatorname{rhcenter}(t + 1 + Q, p + 1) = \operatorname{rhcenter}(t + 1, p + 1)\), then if \(\operatorname{not}(\operatorname{rhcenter}(t, p + 1)) \operatorname{and} \operatorname{bxor}(\operatorname{rhcenter}(t + Q, p + 2), \operatorname{rhcenter}(t, p + 2)) = \operatorname{false}\), then \(\operatorname{rhcenter}(t + Q, p) = \operatorname{rhcenter}(t, p)\).
\end{theorem}
\begin{proof}
The goal follows from \(\operatorname{rhcenteragreeinwardofcoleqtwo}(t, Q, p)\).
\end{proof}
\begin{theorem}
\label{ll:paper:not-lhbandagree-of-lt-add}
For every natural number \(t\) and natural number \(Q\) and natural number \(L\), if \(0 < Q\), then if \(t + Q \le L\), then not \(\operatorname{lhbandagree}(t, Q, 0, L)\).
\end{theorem}
\begin{proof}
The goal follows from \(\operatorname{notlhbandagreeofltadd}(t, Q, L)\).
\end{proof}
\begin{theorem}
\label{ll:paper:lhcenter-outward-determine-step}
For every natural number \(t\) and natural number \(Q\) and natural number \(p0\) and natural number \(S\), if \(\operatorname{lhcolagree}(t, Q, p0, S + 1)\), then if \(\operatorname{lhcolagree}(t, Q, p0 + 1, S + 1)\), then \(\operatorname{lhcolagree}(t, Q, p0 + 2, S)\).
\end{theorem}
\begin{proof}
The goal follows from \(\operatorname{lhcenteroutwarddeterminestep}(t, Q, p0, S)\).
\end{proof}
\begin{theorem}
\label{ll:paper:lhcenter-outward-determine}
For every natural number \(t\) and natural number \(Q\) and natural number \(p0\) and natural number \(S\) and natural number \(K\), if \(\operatorname{lhcolagree}(t, Q, p0, S + K)\), then if \(\operatorname{lhcolagree}(t, Q, p0 + 1, S + K)\), then \(\operatorname{lhcolagree}(t, Q, p0 + K, S)\) and \(\operatorname{lhcolagree}(t, Q, p0 + K + 1, S)\).
\end{theorem}
\begin{proof}
The goal follows from \(\operatorname{lhcenteroutwarddetermine}(t, Q, p0, S, K)\).
\end{proof}
\begin{theorem}
\label{ll:paper:lhbandagree-of-outward-determine}
For every natural number \(t\) and natural number \(Q\) and natural number \(p0\) and natural number \(L\), if \(\operatorname{lhcolagree}(t, Q, p0, L)\), then if \(\operatorname{lhcolagree}(t, Q, p0 + 1, L)\), then \(\operatorname{lhbandagree}(t, Q, p0, L + 1)\).
\end{theorem}
\begin{proof}
The goal follows from \(\operatorname{lhbandagreeofoutwarddetermine}(t, Q, p0, L)\).
\end{proof}
\begin{theorem}
\label{ll:paper:not-lhcolagree-zero-one-of-pos}
For every natural number \(t\) and natural number \(Q\) and natural number \(S\), if \(0 < Q\), then if \(t + Q \le S + 1\), then if \(\operatorname{lhcolagree}(t, Q, 0, S)\), then not \(\operatorname{lhcolagree}(t, Q, 1, S)\).
\end{theorem}
\begin{proof}
The goal follows from \(\operatorname{notlhcolagreezerooneofpos}(t, Q, S)\).
\end{proof}
\begin{theorem}
\label{ll:paper:periodic-lhcolagree-zero}
For every natural number \(T\) and natural number \(P\), if \(0 < P\), then if \(\operatorname{centerperiodic}(T, P)\), then for every natural number \(m\) and natural number \(t\) and natural number \(S\), if \(T \le t\), then \(\operatorname{lhcolagree}(t, m \cdot P, 0, S)\).
\end{theorem}
\begin{proof}
Assume \(hP\), \(hper\).
The goal follows from \(\operatorname{periodiclhcolagreezero}(T, P, hP, hper)\).
\end{proof}
\begin{theorem}
\label{ll:paper:not-periodic-lhcolagree-one-horizon}
For every natural number \(T\) and natural number \(P\), if \(0 < P\), then if \(\operatorname{centerperiodic}(T, P)\), then for every natural number \(m\) and natural number \(t\) and natural number \(S\), if \(0 < m\), then if \(T \le t\), then if \(t + m \cdot P \le S + 1\), then not \(\operatorname{lhcolagree}(t, m \cdot P, 1, S)\).
\end{theorem}
\begin{proof}
Assume \(hP\), \(hper\).
The goal follows from \(\operatorname{notperiodiclhcolagreeonehorizon}(T, P, hP, hper)\).
\end{proof}
\begin{theorem}
\label{ll:paper:periodic-lhcolagree-one-of-rhcolagree-one}
For every natural number \(T\) and natural number \(P\), if \(0 < P\), then if \(\operatorname{centerperiodic}(T, P)\), then for every natural number \(m\) and natural number \(t\) and natural number \(S\), if \(T \le t\), then if \(\operatorname{rhcolagree}(t, m \cdot P, 1, S)\), then \(\operatorname{lhcolagree}(t, m \cdot P, 1, S)\).
\end{theorem}
\begin{proof}
Assume \(hP\), \(hper\).
The goal follows from \(\operatorname{periodiclhcolagreeoneofrhcolagreeone}(T, P, hP, hper)\).
\end{proof}
\begin{theorem}
\label{ll:paper:not-periodic-rhcolagree-one-horizon}
For every natural number \(T\) and natural number \(P\), if \(0 < P\), then if \(\operatorname{centerperiodic}(T, P)\), then for every natural number \(m\) and natural number \(t\) and natural number \(S\), if \(0 < m\), then if \(T \le t\), then if \(t + m \cdot P \le S + 1\), then not \(\operatorname{rhcolagree}(t, m \cdot P, 1, S)\).
\end{theorem}
\begin{proof}
Assume \(hP\), \(hper\).
The goal follows from \(\operatorname{notperiodicrhcolagreeonehorizon}(T, P, hP, hper)\).
\end{proof}
\begin{theorem}
\label{ll:paper:periodic-rhcolagree-one-of-rhbandagree-zero}
For every natural number \(T\) and natural number \(P\), if \(0 < P\), then if \(\operatorname{centerperiodic}(T, P)\), then for every natural number \(m\) and natural number \(t\) and natural number \(S\), if \(T \le t\), then if \(\operatorname{rhbandagree}(t, m \cdot P, 0, S + 1)\), then \(\operatorname{rhcolagree}(t, m \cdot P, 1, S)\).
\end{theorem}
\begin{proof}
Assume \(hP\), \(hper\).
The goal follows from \(\operatorname{periodicrhcolagreeoneofrhbandagreezero}(T, P, hP, hper)\).
\end{proof}
\begin{theorem}
\label{ll:paper:periodic-rhcolagree-one-of-finite-sectors}
For every natural number \(T\) and natural number \(P\), if \(0 < P\), then if \(\operatorname{centerperiodic}(T, P)\), then for every natural number \(k\) and natural number \(d\) and natural number \(U\) and natural number \(M\) and natural number \(m\), if \(1 \le M\), then if \(T \le d + U \cdot \operatorname{twopow}(k)\), then if for every natural number \(j\), if \(j < M\), then \(\operatorname{rhbandagree}(d + (U + j) \cdot \operatorname{twopow}(k) + 1, m \cdot P, 0, \operatorname{twopow}(k) + 1)\), then \(\operatorname{rhcolagree}(d + U \cdot \operatorname{twopow}(k) + 1, m \cdot P, 1, M \cdot \operatorname{twopow}(k))\).
\end{theorem}
\begin{proof}
Assume \(hP\), \(hper\).
The goal follows from \(\operatorname{periodicrhcolagreeoneoffinitesectors}(T, P, hP, hper)\).
\end{proof}
\begin{theorem}
\label{ll:paper:not-periodic-finite-sector-rhbandagree-mul}
For every natural number \(T\) and natural number \(P\), if \(0 < P\), then if \(\operatorname{centerperiodic}(T, P)\), then for every natural number \(k\) and natural number \(d\) and natural number \(U\) and natural number \(M\) and natural number \(m\), if \(1 \le M\), then if \(0 < m\), then if \(T \le d + U \cdot \operatorname{twopow}(k)\), then if \(d + U \cdot \operatorname{twopow}(k) + m \cdot P \le M \cdot \operatorname{twopow}(k)\), then not for every natural number \(j\), if \(j < M\), then \(\operatorname{rhbandagree}(d + (U + j) \cdot \operatorname{twopow}(k) + 1, m \cdot P, 0, \operatorname{twopow}(k) + 1)\).
\end{theorem}
\begin{proof}
Assume \(hP\), \(hper\).
The goal follows from \(\operatorname{notperiodicfinitesectorrhbandagreemul}(T, P, hP, hper)\).
\end{proof}
\begin{theorem}
\label{ll:paper:not-periodic-finite-sector-rhbandagree-lcmperiod}
For every natural number \(T\) and natural number \(P\), if \(0 < P\), then if \(\operatorname{centerperiodic}(T, P)\), then for every natural number \(k\) and natural number \(d\) and natural number \(U\) and natural number \(M\), if \(1 \le M\), then if \(T \le d + U \cdot \operatorname{twopow}(k)\), then if \(d + U \cdot \operatorname{twopow}(k) + \operatorname{lcmperiod}(P, k) \le M \cdot \operatorname{twopow}(k)\), then not for every natural number \(j\), if \(j < M\), then \(\operatorname{rhbandagree}(d + (U + j) \cdot \operatorname{twopow}(k) + 1, \operatorname{lcmperiod}(P, k), 0, \operatorname{twopow}(k) + 1)\).
\end{theorem}
\begin{proof}
Assume \(hP\), \(hper\).
The goal follows from \(\operatorname{notperiodicfinitesectorrhbandagreelcmperiod}(T, P, hP, hper)\).
\end{proof}
\begin{theorem}
\label{ll:paper:periodic-lhcenter-or-zero-defect-eq}
For every natural number \(T\) and natural number \(P\), if \(0 < P\), then if \(\operatorname{centerperiodic}(T, P)\), then for every natural number \(m\) and natural number \(t\), if \(T \le t\), then \(\operatorname{bxor}(\operatorname{lhcenter}(t + m \cdot P, 1) \operatorname{or} \operatorname{lhcenter}(t + m \cdot P, 0), \operatorname{lhcenter}(t, 1) \operatorname{or} \operatorname{lhcenter}(t, 0)) = \operatorname{bxor}(\operatorname{lhcenter}(t + m \cdot P, 1), \operatorname{lhcenter}(t, 1))\).
\end{theorem}
\begin{proof}
Assume \(hP\), \(hper\).
The goal follows from \(\operatorname{periodiclhcenterorzerodefecteq}(T, P, hP, hper)\).
\end{proof}
\begin{theorem}
\label{ll:paper:periodic-lhcenter-two-defect-eq}
For every natural number \(T\) and natural number \(P\), if \(0 < P\), then if \(\operatorname{centerperiodic}(T, P)\), then for every natural number \(m\) and natural number \(t\), if \(T \le t\), then \(\operatorname{bxor}(\operatorname{lhcenter}(t + m \cdot P, 2), \operatorname{lhcenter}(t, 2)) = \operatorname{bxor}(\operatorname{bxor}(\operatorname{lhcenter}(t + 1 + m \cdot P, 1), \operatorname{lhcenter}(t + 1, 1)), \operatorname{bxor}(\operatorname{lhcenter}(t + m \cdot P, 1), \operatorname{lhcenter}(t, 1)))\).
\end{theorem}
\begin{proof}
Assume \(hP\), \(hper\).
The goal follows from \(\operatorname{periodiclhcentertwodefecteq}(T, P, hP, hper)\).
\end{proof}
\begin{theorem}
\label{ll:paper:lhcenter-eq-row-testbit-uncond}
For every natural number \(t\) and natural number \(p\), \(\operatorname{lhcenter}(t, p) = \operatorname{wittbitbool}(\operatorname{wittrow}(t), t + p)\).
\end{theorem}
\begin{proof}
The goal follows from \(\operatorname{lhcentereqrowtestbituncond}(t, p)\).
\end{proof}
\begin{theorem}
\label{ll:paper:lhcenter-scale-step-uncond}
For every natural number \(t\) and natural number \(k\) and natural number \(p\), \(\operatorname{lhcenter}(t + \operatorname{twopow}(k), p + \operatorname{twopow}(k)) = \operatorname{bxor}(\operatorname{bxor}(\operatorname{bxor}(\operatorname{lhcenter}(t, p), \operatorname{lhcenter}(t, p + \operatorname{twopow}(k))), \operatorname{lhcenter}(t, p + \operatorname{twopow}(k + 1))), \operatorname{wittbitbool}(\operatorname{wittwin}(t, \operatorname{twopow}(k)), t + p + \operatorname{twopow}(k + 1)))\).
\end{theorem}
\begin{proof}
The goal follows from \(\operatorname{lhcenterscalestepuncond}(t, k, p)\).
\end{proof}
\begin{theorem}
\label{ll:paper:lhcenter-scale-outward-step-uncond}
For every natural number \(t\) and natural number \(k\) and natural number \(p\), \(\operatorname{lhcenter}(t, p + \operatorname{twopow}(k + 1)) = \operatorname{bxor}(\operatorname{bxor}(\operatorname{bxor}(\operatorname{lhcenter}(t + \operatorname{twopow}(k), p + \operatorname{twopow}(k)), \operatorname{lhcenter}(t, p)), \operatorname{lhcenter}(t, p + \operatorname{twopow}(k))), \operatorname{wittbitbool}(\operatorname{wittwin}(t, \operatorname{twopow}(k)), t + p + \operatorname{twopow}(k + 1)))\).
\end{theorem}
\begin{proof}
The goal follows from \(\operatorname{lhcenterscaleoutwardstepuncond}(t, k, p)\).
\end{proof}
\begin{theorem}
\label{ll:paper:lhcenter-scale-cross-step}
For every natural number \(t\) and natural number \(k\) and natural number \(r\), if \(r \le \operatorname{twopow}(k)\), then \(\operatorname{lhcenter}(t + \operatorname{twopow}(k), \operatorname{twopow}(k) - r) = \operatorname{bxor}(\operatorname{bxor}(\operatorname{bxor}(\operatorname{rhcenter}(t, r), \operatorname{lhcenter}(t, \operatorname{twopow}(k) - r)), \operatorname{lhcenter}(t, \operatorname{twopow}(k + 1) - r)), \operatorname{wittbitbool}(\operatorname{wittwin}(t, \operatorname{twopow}(k)), t + \operatorname{twopow}(k + 1) - r))\).
\end{theorem}
\begin{proof}
The goal follows from \(\operatorname{lhcenterscalecrossstep}(t, k, r)\).
\end{proof}
\begin{theorem}
\label{ll:paper:rhcenter-eq-lhcenter-scale-cross}
For every natural number \(t\) and natural number \(k\) and natural number \(r\), if \(r \le \operatorname{twopow}(k)\), then \(\operatorname{rhcenter}(t, r) = \operatorname{bxor}(\operatorname{bxor}(\operatorname{bxor}(\operatorname{lhcenter}(t + \operatorname{twopow}(k), \operatorname{twopow}(k) - r), \operatorname{lhcenter}(t, \operatorname{twopow}(k) - r)), \operatorname{lhcenter}(t, \operatorname{twopow}(k + 1) - r)), \operatorname{wittbitbool}(\operatorname{wittwin}(t, \operatorname{twopow}(k)), t + \operatorname{twopow}(k + 1) - r))\).
\end{theorem}
\begin{proof}
The goal follows from \(\operatorname{rhcentereqlhcenterscalecross}(t, k, r)\).
\end{proof}
\begin{theorem}
\label{ll:paper:rhbandagree-of-lhbandagree-scale-cross}
For every natural number \(t\) and natural number \(k\) and natural number \(Q\) and natural number \(L\), if \(L \le \operatorname{twopow}(k)\), then if \(\operatorname{lhbandagree}(t + \operatorname{twopow}(k), Q, \operatorname{twopow}(k) - L, L)\), then if \(\operatorname{lhbandagree}(t, Q, \operatorname{twopow}(k) - L, L)\), then if \(\operatorname{lhbandagree}(t, Q, \operatorname{twopow}(k + 1) - L, L)\), then if \(\operatorname{winwittagree}(t, k, Q, t + \operatorname{twopow}(k + 1) - L, L)\), then \(\operatorname{rhbandagree}(t, Q, 0, L)\).
\end{theorem}
\begin{proof}
The goal follows from \(\operatorname{rhbandagreeoflhbandagreescalecross}(t, k, Q, L)\).
\end{proof}
\begin{theorem}
\label{ll:paper:winwittagree-of-rhbandagree-lhbandagree-scale-cross}
For every natural number \(t\) and natural number \(k\) and natural number \(Q\) and natural number \(L\), if \(L \le \operatorname{twopow}(k)\), then if \(\operatorname{rhbandagree}(t, Q, 0, L)\), then if \(\operatorname{lhbandagree}(t + \operatorname{twopow}(k), Q, \operatorname{twopow}(k) - L, L)\), then if \(\operatorname{lhbandagree}(t, Q, \operatorname{twopow}(k) - L, L)\), then if \(\operatorname{lhbandagree}(t, Q, \operatorname{twopow}(k + 1) - L, L)\), then \(\operatorname{winwittagree}(t, k, Q, t + \operatorname{twopow}(k + 1) - L, L)\).
\end{theorem}
\begin{proof}
The goal follows from \(\operatorname{winwittagreeofrhbandagreelhbanscalecross}(t, k, Q, L)\).
\end{proof}
\begin{theorem}
\label{ll:paper:periodic-rhbandagree-of-lhcolagree-one-and-winwittagree}
For every natural number \(T\) and natural number \(P\), if \(0 < P\), then if \(\operatorname{centerperiodic}(T, P)\), then for every natural number \(m\) and natural number \(t\) and natural number \(k\), if \(T \le t\), then if \(\operatorname{lhcolagree}(t, m \cdot P, 1, \operatorname{twopow}(k + 1))\), then if \(\operatorname{winwittagree}(t, k, m \cdot P, t + \operatorname{twopow}(k), \operatorname{twopow}(k))\), then \(\operatorname{rhbandagree}(t, m \cdot P, 0, \operatorname{twopow}(k))\).
\end{theorem}
\begin{proof}
Assume \(hP\), \(hper\).
The goal follows from \(\operatorname{periodicrhbandagreeoflhcolagreeoneandwinwittagree}(T, P, hP, hper)\).
\end{proof}
\begin{theorem}
\label{ll:paper:periodic-winwittagree-of-lhcolagree-one-and-rhbandagree}
For every natural number \(T\) and natural number \(P\), if \(0 < P\), then if \(\operatorname{centerperiodic}(T, P)\), then for every natural number \(m\) and natural number \(t\) and natural number \(k\), if \(T \le t\), then if \(\operatorname{lhcolagree}(t, m \cdot P, 1, \operatorname{twopow}(k + 1))\), then if \(\operatorname{rhbandagree}(t, m \cdot P, 0, \operatorname{twopow}(k))\), then \(\operatorname{winwittagree}(t, k, m \cdot P, t + \operatorname{twopow}(k), \operatorname{twopow}(k))\).
\end{theorem}
\begin{proof}
Assume \(hP\), \(hper\).
The goal follows from \(\operatorname{periodicwinwittagreeoflhcolagreeoneandrhbandagree}(T, P, hP, hper)\).
\end{proof}
\begin{theorem}
\label{ll:paper:lhbandagree-scale-step-uncond}
For every natural number \(t\) and natural number \(k\) and natural number \(Q\) and natural number \(p0\) and natural number \(L\), if \(\operatorname{lhbandagree}(t, Q, p0, L)\), then if \(\operatorname{lhbandagree}(t, Q, p0 + \operatorname{twopow}(k), L)\), then if \(\operatorname{lhbandagree}(t, Q, p0 + \operatorname{twopow}(k + 1), L)\), then if \(\operatorname{winwittagree}(t, k, Q, t + p0 + \operatorname{twopow}(k + 1), L)\), then \(\operatorname{lhbandagree}(t + \operatorname{twopow}(k), Q, p0 + \operatorname{twopow}(k), L)\).
\end{theorem}
\begin{proof}
The goal follows from \(\operatorname{lhbandagreescalestepuncond}(t, k, Q, p0, L)\).
\end{proof}
\begin{theorem}
\label{ll:paper:lhbandagree-scale-outward-step-uncond}
For every natural number \(t\) and natural number \(k\) and natural number \(Q\) and natural number \(p0\) and natural number \(L\), if \(\operatorname{lhbandagree}(t + \operatorname{twopow}(k), Q, p0 + \operatorname{twopow}(k), L)\), then if \(\operatorname{lhbandagree}(t, Q, p0, L)\), then if \(\operatorname{lhbandagree}(t, Q, p0 + \operatorname{twopow}(k), L)\), then if \(\operatorname{winwittagree}(t, k, Q, t + p0 + \operatorname{twopow}(k + 1), L)\), then \(\operatorname{lhbandagree}(t, Q, p0 + \operatorname{twopow}(k + 1), L)\).
\end{theorem}
\begin{proof}
The goal follows from \(\operatorname{lhbandagreescaleoutwardstepuncond}(t, k, Q, p0, L)\).
\end{proof}
\begin{theorem}
\label{ll:paper:lhcenter-diag-diff-eq-or}
For every natural number \(t\) and natural number \(p\), \(\operatorname{bxor}(\operatorname{lhcenter}(t + 1, p + 1), \operatorname{lhcenter}(t, p + 2)) = (\operatorname{lhcenter}(t, p + 1) \operatorname{or} \operatorname{lhcenter}(t, p))\).
\end{theorem}
\begin{proof}
The goal follows from \(\operatorname{lhcenterdiagdiffeqor}(t, p)\).
\end{proof}
\begin{theorem}
\label{ll:paper:lhcenter-defect-outward-step}
For every natural number \(t\) and natural number \(Q\) and natural number \(p\), \(\operatorname{bxor}(\operatorname{lhcenter}(t + Q, p + 2), \operatorname{lhcenter}(t, p + 2)) = \operatorname{bxor}(\operatorname{bxor}(\operatorname{lhcenter}(t + 1 + Q, p + 1), \operatorname{lhcenter}(t + 1, p + 1)), \operatorname{bxor}(\operatorname{lhcenter}(t + Q, p + 1) \operatorname{or} \operatorname{lhcenter}(t + Q, p), \operatorname{lhcenter}(t, p + 1) \operatorname{or} \operatorname{lhcenter}(t, p)))\).
\end{theorem}
\begin{proof}
The goal follows from \(\operatorname{lhcenterdefectoutwardstep}(t, Q, p)\).
\end{proof}
\begin{theorem}
\label{ll:paper:lhcenter-defect-outward-of-col-eq}
For every natural number \(t\) and natural number \(Q\) and natural number \(p\), if \(\operatorname{lhcenter}(t + Q, p) = \operatorname{lhcenter}(t, p)\), then \(\operatorname{bxor}(\operatorname{lhcenter}(t + Q, p + 2), \operatorname{lhcenter}(t, p + 2)) = \operatorname{bxor}(\operatorname{bxor}(\operatorname{lhcenter}(t + 1 + Q, p + 1), \operatorname{lhcenter}(t + 1, p + 1)), \operatorname{not}(\operatorname{lhcenter}(t, p)) \operatorname{and} \operatorname{bxor}(\operatorname{lhcenter}(t + Q, p + 1), \operatorname{lhcenter}(t, p + 1)))\).
\end{theorem}
\begin{proof}
The goal follows from \(\operatorname{lhcenterdefectoutwardofcoleq}(t, Q, p)\).
\end{proof}
\begin{theorem}
\label{ll:paper:lhcenter-defect-outward-of-black}
For every natural number \(t\) and natural number \(Q\) and natural number \(p\), if \(\operatorname{lhcenter}(t + Q, p) = \operatorname{true}\), then if \(\operatorname{lhcenter}(t, p) = \operatorname{true}\), then \(\operatorname{bxor}(\operatorname{lhcenter}(t + Q, p + 2), \operatorname{lhcenter}(t, p + 2)) = \operatorname{bxor}(\operatorname{lhcenter}(t + 1 + Q, p + 1), \operatorname{lhcenter}(t + 1, p + 1))\).
\end{theorem}
\begin{proof}
The goal follows from \(\operatorname{lhcenterdefectoutwardofblack}(t, Q, p)\).
\end{proof}
\begin{theorem}
\label{ll:paper:winwitt-testbit-left-edge}
For every natural number \(t\) and natural number \(k\), \(\operatorname{wittbitbool}(\operatorname{wittwin}(t, \operatorname{twopow}(k)), t + t + \operatorname{twopow}(k + 1)) = \operatorname{false}\).
\end{theorem}
\begin{proof}
The goal follows from \(\operatorname{winwitttestbitleftedge}(t, k)\).
\end{proof}
\begin{theorem}
\label{ll:paper:winwitt-testbit-gt-left-edge}
For every natural number \(t\) and natural number \(k\) and natural number \(p\), if \(t < p\), then \(\operatorname{wittbitbool}(\operatorname{wittwin}(t, \operatorname{twopow}(k)), t + p + \operatorname{twopow}(k + 1)) = \operatorname{false}\).
\end{theorem}
\begin{proof}
The goal follows from \(\operatorname{winwitttestbitgtleftedge}(t, k, p)\).
\end{proof}
\begin{theorem}
\label{ll:paper:winwitt-testbit-ge-left-edge}
For every natural number \(t\) and natural number \(k\) and natural number \(p\), if \(t \le p\), then \(\operatorname{wittbitbool}(\operatorname{wittwin}(t, \operatorname{twopow}(k)), t + p + \operatorname{twopow}(k + 1)) = \operatorname{false}\).
\end{theorem}
\begin{proof}
The goal follows from \(\operatorname{winwitttestbitgeleftedge}(t, k, p)\).
\end{proof}
\begin{theorem}
\label{ll:paper:winwitt-testbit-settled-left}
For every natural number \(t\) and natural number \(k\) and natural number \(j\), if \(j \le k\), then if \(j \le t\), then if \(\operatorname{settledby}(j, t)\), then \(\operatorname{wittbitbool}(\operatorname{wittwin}(t, \operatorname{twopow}(k)), t + (t - j) + \operatorname{twopow}(k + 1)) = \operatorname{false}\).
\end{theorem}
\begin{proof}
The goal follows from \(\operatorname{winwitttestbitsetledleft}(t, k, j)\).
\end{proof}
\begin{theorem}
\label{ll:paper:lhbandagree-outside-cone}
For every natural number \(t\) and natural number \(Q\) and natural number \(p0\) and natural number \(L\), if \(t + Q < p0\), then \(\operatorname{lhbandagree}(t, Q, p0, L)\).
\end{theorem}
\begin{proof}
The goal follows from \(\operatorname{lhbandagreeoutsidecone}(t, Q, p0, L)\).
\end{proof}
\begin{theorem}
\label{ll:paper:lhbandagree-mono}
For every natural number \(t\) and natural number \(Q\) and natural number \(p0\) and natural number \(L\) and natural number \(p0p\) and natural number \(Lp\), if \(p0 \le p0p\), then if \(p0p + Lp \le p0 + L\), then if \(\operatorname{lhbandagree}(t, Q, p0, L)\), then \(\operatorname{lhbandagree}(t, Q, p0p, Lp)\).
\end{theorem}
\begin{proof}
The goal follows from \(\operatorname{lhbandagreemono}(t, Q, p0, L, p0p, Lp)\).
\end{proof}
\begin{theorem}
\label{ll:paper:lhcolagree-mono}
For every natural number \(t\) and natural number \(Q\) and natural number \(p0\) and natural number \(S\) and natural number \(s0\) and natural number \(Sp\), if \(s0 + Sp \le S\), then if \(\operatorname{lhcolagree}(t, Q, p0, S)\), then \(\operatorname{lhcolagree}(t + s0, Q, p0, Sp)\).
\end{theorem}
\begin{proof}
The goal follows from \(\operatorname{lhcolagreemono}(t, Q, p0, S, s0, Sp)\).
\end{proof}
\begin{theorem}
\label{ll:paper:lhbandagree-step}
For every natural number \(t\) and natural number \(Q\) and natural number \(p0\) and natural number \(L\), if \(\operatorname{lhbandagree}(t, Q, p0, L + 2)\), then \(\operatorname{lhbandagree}(t + 1, Q, p0 + 1, L)\).
\end{theorem}
\begin{proof}
The goal follows from \(\operatorname{lhbandagreestep}(t, Q, p0, L)\).
\end{proof}
\begin{theorem}
\label{ll:paper:lhbandagree-steps}
For every natural number \(t\) and natural number \(Q\) and natural number \(p0\) and natural number \(L\) and natural number \(s\), if \(\operatorname{lhbandagree}(t, Q, p0, L + 2 \cdot s)\), then \(\operatorname{lhbandagree}(t + s, Q, p0 + s, L)\).
\end{theorem}
\begin{proof}
The goal follows from \(\operatorname{lhbandagreesteps}(t, Q, p0, L, s)\).
\end{proof}
\begin{theorem}
\label{ll:paper:periodic-lhbandagree-zero-of-one}
For every natural number \(T\) and natural number \(P\), if \(0 < P\), then if \(\operatorname{centerperiodic}(T, P)\), then for every natural number \(m\) and natural number \(t\) and natural number \(L\), if \(T \le t\), then if \(\operatorname{lhbandagree}(t, m \cdot P, 1, L)\), then \(\operatorname{lhbandagree}(t, m \cdot P, 0, L + 1)\).
\end{theorem}
\begin{proof}
Assume \(hP\), \(hper\).
The goal follows from \(\operatorname{periodiclhbandagreezeroofone}(T, P, hP, hper)\).
\end{proof}
\begin{theorem}
\label{ll:paper:periodic-rhbandagree-of-rhcolagree-one-and-winwittagree}
For every natural number \(T\) and natural number \(P\), if \(0 < P\), then if \(\operatorname{centerperiodic}(T, P)\), then for every natural number \(m\) and natural number \(t\) and natural number \(k\), if \(T \le t\), then if \(\operatorname{rhcolagree}(t, m \cdot P, 1, \operatorname{twopow}(k + 1))\), then if \(\operatorname{winwittagree}(t, k, m \cdot P, t + \operatorname{twopow}(k), \operatorname{twopow}(k))\), then \(\operatorname{rhbandagree}(t, m \cdot P, 0, \operatorname{twopow}(k))\).
\end{theorem}
\begin{proof}
Assume \(hP\), \(hper\).
The goal follows from \(\operatorname{periodicrhbandagreeofrhcolagreeoneandwinwittagree}(T, P, hP, hper)\).
\end{proof}
\begin{theorem}
\label{ll:paper:periodic-winwittagree-of-rhcolagree-one-and-rhbandagree}
For every natural number \(T\) and natural number \(P\), if \(0 < P\), then if \(\operatorname{centerperiodic}(T, P)\), then for every natural number \(m\) and natural number \(t\) and natural number \(k\), if \(T \le t\), then if \(\operatorname{rhcolagree}(t, m \cdot P, 1, \operatorname{twopow}(k + 1))\), then if \(\operatorname{rhbandagree}(t, m \cdot P, 0, \operatorname{twopow}(k))\), then \(\operatorname{winwittagree}(t, k, m \cdot P, t + \operatorname{twopow}(k), \operatorname{twopow}(k))\).
\end{theorem}
\begin{proof}
Assume \(hP\), \(hper\).
The goal follows from \(\operatorname{periodicwinwittagreeofrhcolagreeoneandrhbandagree}(T, P, hP, hper)\).
\end{proof}
\begin{theorem}
\label{ll:paper:lhcenter-scale-inward-step-uncond}
For every natural number \(t\) and natural number \(k\) and natural number \(p\), \(\operatorname{lhcenter}(t, p) = \operatorname{bxor}(\operatorname{bxor}(\operatorname{bxor}(\operatorname{lhcenter}(t + \operatorname{twopow}(k), p + \operatorname{twopow}(k)), \operatorname{lhcenter}(t, p + \operatorname{twopow}(k))), \operatorname{lhcenter}(t, p + \operatorname{twopow}(k + 1))), \operatorname{wittbitbool}(\operatorname{wittwin}(t, \operatorname{twopow}(k)), t + p + \operatorname{twopow}(k + 1)))\).
\end{theorem}
\begin{proof}
The goal follows from \(\operatorname{lhcenterscaleinwardstepuncond}(t, k, p)\).
\end{proof}
\begin{theorem}
\label{ll:paper:lhbandagree-scale-inward-step-uncond}
For every natural number \(t\) and natural number \(k\) and natural number \(Q\) and natural number \(p0\) and natural number \(L\), if \(\operatorname{lhbandagree}(t + \operatorname{twopow}(k), Q, p0 + \operatorname{twopow}(k), L)\), then if \(\operatorname{lhbandagree}(t, Q, p0 + \operatorname{twopow}(k), L)\), then if \(\operatorname{lhbandagree}(t, Q, p0 + \operatorname{twopow}(k + 1), L)\), then if \(\operatorname{winwittagree}(t, k, Q, t + p0 + \operatorname{twopow}(k + 1), L)\), then \(\operatorname{lhbandagree}(t, Q, p0, L)\).
\end{theorem}
\begin{proof}
The goal follows from \(\operatorname{lhbandagreescaleinwardstepuncond}(t, k, Q, p0, L)\).
\end{proof}
\begin{theorem}
\label{ll:paper:winwittagree-outside-left-cone}
For every natural number \(t\) and natural number \(k\) and natural number \(Q\) and natural number \(p0\) and natural number \(L\), if \(t + Q \le p0\), then \(\operatorname{winwittagree}(t, k, Q, t + p0 + \operatorname{twopow}(k + 1), L)\).
\end{theorem}
\begin{proof}
The goal follows from \(\operatorname{winwittagreeoutsideleftcone}(t, k, Q, p0, L)\).
\end{proof}
\begin{theorem}
\label{ll:paper:rdiag-add-eq-lhcenter-uncond}
For every natural number \(t\) and natural number \(p\), \(\operatorname{rdiag}(t + p, t) = \operatorname{lhcenter}(t, p)\).
\end{theorem}
\begin{proof}
The goal follows from \(\operatorname{rdiagaddeqlhcenteruncond}(t, p)\).
\end{proof}
\begin{theorem}
\label{ll:paper:winwitt-testbit-eq-lhcenter-scale-uncond}
For every natural number \(t\) and natural number \(k\) and natural number \(p\), \(\operatorname{wittbitbool}(\operatorname{wittwin}(t, \operatorname{twopow}(k)), t + p + \operatorname{twopow}(k + 1)) = \operatorname{bxor}(\operatorname{bxor}(\operatorname{bxor}(\operatorname{lhcenter}(t + \operatorname{twopow}(k), p + \operatorname{twopow}(k)), \operatorname{lhcenter}(t, p)), \operatorname{lhcenter}(t, p + \operatorname{twopow}(k))), \operatorname{lhcenter}(t, p + \operatorname{twopow}(k + 1)))\).
\end{theorem}
\begin{proof}
The goal follows from \(\operatorname{winwitttestbiteqlhcenterscaleuncond}(t, k, p)\).
\end{proof}
\begin{theorem}
\label{ll:paper:winwittagree-of-lhbandagree-scale-uncond}
For every natural number \(t\) and natural number \(k\) and natural number \(Q\) and natural number \(p0\) and natural number \(L\), if \(\operatorname{lhbandagree}(t + \operatorname{twopow}(k), Q, p0 + \operatorname{twopow}(k), L)\), then if \(\operatorname{lhbandagree}(t, Q, p0, L)\), then if \(\operatorname{lhbandagree}(t, Q, p0 + \operatorname{twopow}(k), L)\), then if \(\operatorname{lhbandagree}(t, Q, p0 + \operatorname{twopow}(k + 1), L)\), then \(\operatorname{winwittagree}(t, k, Q, t + p0 + \operatorname{twopow}(k + 1), L)\).
\end{theorem}
\begin{proof}
The goal follows from \(\operatorname{winwittagreeoflhbandagreescaleuncond}(t, k, Q, p0, L)\).
\end{proof}
\begin{theorem}
\label{ll:paper:lhcenter-scale-mid-step-uncond}
For every natural number \(t\) and natural number \(k\) and natural number \(p\), \(\operatorname{lhcenter}(t, p + \operatorname{twopow}(k)) = \operatorname{bxor}(\operatorname{bxor}(\operatorname{bxor}(\operatorname{lhcenter}(t + \operatorname{twopow}(k), p + \operatorname{twopow}(k)), \operatorname{lhcenter}(t, p)), \operatorname{lhcenter}(t, p + \operatorname{twopow}(k + 1))), \operatorname{wittbitbool}(\operatorname{wittwin}(t, \operatorname{twopow}(k)), t + p + \operatorname{twopow}(k + 1)))\).
\end{theorem}
\begin{proof}
The goal follows from \(\operatorname{lhcenterscalemidstepuncond}(t, k, p)\).
\end{proof}
\begin{theorem}
\label{ll:paper:lhbandagree-scale-mid-step-uncond}
For every natural number \(t\) and natural number \(k\) and natural number \(Q\) and natural number \(p0\) and natural number \(L\), if \(\operatorname{lhbandagree}(t + \operatorname{twopow}(k), Q, p0 + \operatorname{twopow}(k), L)\), then if \(\operatorname{lhbandagree}(t, Q, p0, L)\), then if \(\operatorname{lhbandagree}(t, Q, p0 + \operatorname{twopow}(k + 1), L)\), then if \(\operatorname{winwittagree}(t, k, Q, t + p0 + \operatorname{twopow}(k + 1), L)\), then \(\operatorname{lhbandagree}(t, Q, p0 + \operatorname{twopow}(k), L)\).
\end{theorem}
\begin{proof}
The goal follows from \(\operatorname{lhbandagreescalemidstepuncond}(t, k, Q, p0, L)\).
\end{proof}
\begin{theorem}
\label{ll:paper:lhcenter-scale-step-outer-strip}
For every natural number \(t\) and natural number \(k\) and natural number \(p\), if \(t < p + \operatorname{twopow}(k)\), then \(\operatorname{lhcenter}(t + \operatorname{twopow}(k), p + \operatorname{twopow}(k)) = \operatorname{bxor}(\operatorname{lhcenter}(t, p), \operatorname{wittbitbool}(\operatorname{wittwin}(t, \operatorname{twopow}(k)), t + p + \operatorname{twopow}(k + 1)))\).
\end{theorem}
\begin{proof}
The goal follows from \(\operatorname{lhcenterscalestepouterstrip}(t, k, p)\).
\end{proof}
\begin{theorem}
\label{ll:paper:lhcenter-scale-inward-outer-strip}
For every natural number \(t\) and natural number \(k\) and natural number \(p\), if \(t < p + \operatorname{twopow}(k)\), then \(\operatorname{lhcenter}(t, p) = \operatorname{bxor}(\operatorname{lhcenter}(t + \operatorname{twopow}(k), p + \operatorname{twopow}(k)), \operatorname{wittbitbool}(\operatorname{wittwin}(t, \operatorname{twopow}(k)), t + p + \operatorname{twopow}(k + 1)))\).
\end{theorem}
\begin{proof}
The goal follows from \(\operatorname{lhcenterscaleinwardouterstrip}(t, k, p)\).
\end{proof}
\begin{theorem}
\label{ll:paper:winwitt-testbit-outer-left-strip}
For every natural number \(t\) and natural number \(k\) and natural number \(p\), if \(t < p + \operatorname{twopow}(k)\), then \(\operatorname{wittbitbool}(\operatorname{wittwin}(t, \operatorname{twopow}(k)), t + p + \operatorname{twopow}(k + 1)) = \operatorname{bxor}(\operatorname{lhcenter}(t + \operatorname{twopow}(k), p + \operatorname{twopow}(k)), \operatorname{lhcenter}(t, p))\).
\end{theorem}
\begin{proof}
The goal follows from \(\operatorname{winwitttestbitouterleftstrip}(t, k, p)\).
\end{proof}
\begin{theorem}
\label{ll:paper:winwitt-testbit-outer-left-strip-eq-ldiag-xor}
For every natural number \(t\) and natural number \(k\) and natural number \(p\), if \(p \le t\), then if \(t < p + \operatorname{twopow}(k)\), then \(\operatorname{wittbitbool}(\operatorname{wittwin}(t, \operatorname{twopow}(k)), t + p + \operatorname{twopow}(k + 1)) = \operatorname{bxor}(\operatorname{ldiag}(t - p, t + \operatorname{twopow}(k)), \operatorname{ldiag}(t - p, t))\).
\end{theorem}
\begin{proof}
The goal follows from \(\operatorname{winwitttestbitouterleftstripeqldiagxor}(t, k, p)\).
\end{proof}
\begin{theorem}
\label{ll:paper:lhbandagree-scale-inward-outer-strip}
For every natural number \(t\) and natural number \(k\) and natural number \(Q\) and natural number \(p0\) and natural number \(L\), if \(t + Q < p0 + \operatorname{twopow}(k)\), then if \(\operatorname{lhbandagree}(t + \operatorname{twopow}(k), Q, p0 + \operatorname{twopow}(k), L)\), then if \(\operatorname{winwittagree}(t, k, Q, t + p0 + \operatorname{twopow}(k + 1), L)\), then \(\operatorname{lhbandagree}(t, Q, p0, L)\).
\end{theorem}
\begin{proof}
The goal follows from \(\operatorname{lhbandagreescaleinwardouterstrip}(t, k, Q, p0, L)\).
\end{proof}
\begin{theorem}
\label{ll:paper:lhbandagree-scale-step-outer-strip}
For every natural number \(t\) and natural number \(k\) and natural number \(Q\) and natural number \(p0\) and natural number \(L\), if \(t + Q < p0 + \operatorname{twopow}(k)\), then if \(\operatorname{lhbandagree}(t, Q, p0, L)\), then if \(\operatorname{winwittagree}(t, k, Q, t + p0 + \operatorname{twopow}(k + 1), L)\), then \(\operatorname{lhbandagree}(t + \operatorname{twopow}(k), Q, p0 + \operatorname{twopow}(k), L)\).
\end{theorem}
\begin{proof}
The goal follows from \(\operatorname{lhbandagreescalestepouterstrip}(t, k, Q, p0, L)\).
\end{proof}
\begin{theorem}
\label{ll:paper:winwittagree-outer-left-strip}
For every natural number \(t\) and natural number \(k\) and natural number \(Q\) and natural number \(p0\) and natural number \(L\), if \(t + Q < p0 + \operatorname{twopow}(k)\), then if \(\operatorname{lhbandagree}(t + \operatorname{twopow}(k), Q, p0 + \operatorname{twopow}(k), L)\), then if \(\operatorname{lhbandagree}(t, Q, p0, L)\), then \(\operatorname{winwittagree}(t, k, Q, t + p0 + \operatorname{twopow}(k + 1), L)\).
\end{theorem}
\begin{proof}
The goal follows from \(\operatorname{winwittagreeouterleftstrip}(t, k, Q, p0, L)\).
\end{proof}
\begin{theorem}
\label{ll:paper:not-lhbandagree-and-winwittagree-of-lt-twopow}
For every natural number \(t\) and natural number \(k\) and natural number \(Q\) and natural number \(L\), if \(0 < Q\), then if \(t + Q \le L\), then if \(t + Q < \operatorname{twopow}(k)\), then if \(\operatorname{lhbandagree}(t + \operatorname{twopow}(k), Q, \operatorname{twopow}(k), L)\), then not \(\operatorname{winwittagree}(t, k, Q, t + \operatorname{twopow}(k + 1), L)\).
\end{theorem}
\begin{proof}
The goal follows from \(\operatorname{notlhbandagreeandwinwittagreeoflttwopow}(t, k, Q, L)\).
\end{proof}
\begin{theorem}
\label{ll:paper:frobdefect-row-testbit-lhcenter-uncond}
For every natural number \(t\) and natural number \(k\) and natural number \(p\), \(\operatorname{wittbitbool}(\operatorname{frobdefect}(k, \operatorname{wittrow}(t)), t + p + 1 + \operatorname{twopow}(k + 1)) = \operatorname{bxor}(\operatorname{bxor}(\operatorname{lhcenter}(t, p + 1 + \operatorname{twopow}(k + 1)) \operatorname{and} \operatorname{bxor}(\operatorname{lhcenter}(t, p + \operatorname{twopow}(k)), \operatorname{lhcenter}(t, p)), \operatorname{lhcenter}(t, p + 1 + \operatorname{twopow}(k)) \operatorname{and} \operatorname{bxor}(\operatorname{lhcenter}(t, p + \operatorname{twopow}(k + 1)), \operatorname{lhcenter}(t, p))), \operatorname{lhcenter}(t, p + 1) \operatorname{and} \operatorname{bxor}(\operatorname{lhcenter}(t, p + \operatorname{twopow}(k + 1)), \operatorname{lhcenter}(t, p + \operatorname{twopow}(k))))\).
\end{theorem}
\begin{proof}
The goal follows from \(\operatorname{frobdefectrowtestbitlhcenteruncond}(t, k, p)\).
\end{proof}
\begin{theorem}
\label{ll:paper:frobdefect-row-testbit-lhcenter-outer-strip-eq-false}
For every natural number \(t\) and natural number \(k\) and natural number \(p\), if \(t < p + \operatorname{twopow}(k)\), then \(\operatorname{wittbitbool}(\operatorname{frobdefect}(k, \operatorname{wittrow}(t)), t + p + 1 + \operatorname{twopow}(k + 1)) = \operatorname{false}\).
\end{theorem}
\begin{proof}
The goal follows from \(\operatorname{frobdefectrowtestbitlhcenterouterstripeqfalse}(t, k, p)\).
\end{proof}
\begin{theorem}
\label{ll:paper:phiscale-row-testbit-lhcenter-uncond}
For every natural number \(t\) and natural number \(k\) and natural number \(p\), \(\operatorname{wittbitbool}(\operatorname{wittscale}(k, \operatorname{wittrow}(t)), t + p + \operatorname{twopow}(k + 1)) = \operatorname{bxor}(\operatorname{bxor}(\operatorname{lhcenter}(t, p + \operatorname{twopow}(k + 1)), \operatorname{lhcenter}(t, p + \operatorname{twopow}(k))), \operatorname{lhcenter}(t, p))\).
\end{theorem}
\begin{proof}
The goal follows from \(\operatorname{phiscalerowtestbitlhcenteruncond}(t, k, p)\).
\end{proof}
\begin{theorem}
\label{ll:paper:phiscale-row-testbit-lhcenter-outer-strip}
For every natural number \(t\) and natural number \(k\) and natural number \(p\), if \(t < p + \operatorname{twopow}(k)\), then \(\operatorname{wittbitbool}(\operatorname{wittscale}(k, \operatorname{wittrow}(t)), t + p + \operatorname{twopow}(k + 1)) = \operatorname{lhcenter}(t, p)\).
\end{theorem}
\begin{proof}
The goal follows from \(\operatorname{phiscalerowtestbitlhcenterouterstrip}(t, k, p)\).
\end{proof}
\begin{theorem}
\label{ll:paper:crosscarry-row-scale-lhcenter-outer-strip}
For every natural number \(t\) and natural number \(k\) and natural number \(p\), if \(t < p + \operatorname{twopow}(k)\), then \(\operatorname{wittbitbool}(\operatorname{crosscarry}(\operatorname{wittscale}(k, \operatorname{wittrow}(t)), \operatorname{wittwin}(t, \operatorname{twopow}(k))), t + p + 1 + \operatorname{twopow}(k + 1)) = \operatorname{bxor}(\operatorname{lhcenter}(t, p + 1) \operatorname{and} \operatorname{wittbitbool}(\operatorname{wittwin}(t, \operatorname{twopow}(k)), t + p + \operatorname{twopow}(k + 1)), \operatorname{wittbitbool}(\operatorname{wittwin}(t, \operatorname{twopow}(k)), t + (p + 1) + \operatorname{twopow}(k + 1)) \operatorname{and} \operatorname{lhcenter}(t, p))\).
\end{theorem}
\begin{proof}
The goal follows from \(\operatorname{crosscarryrowscalelhcenterouterstrip}(t, k, p)\).
\end{proof}
\begin{theorem}
\label{ll:paper:winwitt-step-outer-left-strip}
For every natural number \(t\) and natural number \(k\) and natural number \(p\), if \(t < p + \operatorname{twopow}(k)\), then \(\operatorname{wittbitbool}(\operatorname{wittwin}(t + 1, \operatorname{twopow}(k)), t + 1 + (p + 1) + \operatorname{twopow}(k + 1)) = \operatorname{bxor}(\operatorname{wittbitbool}(\operatorname{wittphi}(\operatorname{wittwin}(t, \operatorname{twopow}(k))), t + 1 + (p + 1) + \operatorname{twopow}(k + 1)), \operatorname{bxor}(\operatorname{lhcenter}(t, p + 1) \operatorname{and} \operatorname{wittbitbool}(\operatorname{wittwin}(t, \operatorname{twopow}(k)), t + p + \operatorname{twopow}(k + 1)), \operatorname{wittbitbool}(\operatorname{wittwin}(t, \operatorname{twopow}(k)), t + (p + 1) + \operatorname{twopow}(k + 1)) \operatorname{and} \operatorname{lhcenter}(t, p)))\).
\end{theorem}
\begin{proof}
The goal follows from \(\operatorname{winwittstepouterleftstrip}(t, k, p)\).
\end{proof}
\begin{theorem}
\label{ll:paper:winwittagree-step-outer-left-strip}
For every natural number \(t\) and natural number \(k\) and natural number \(Q\) and natural number \(p0\) and natural number \(L\), if \(t + Q < p0 + \operatorname{twopow}(k)\), then if \(\operatorname{lhbandagree}(t, Q, p0, L + 1)\), then if \(\operatorname{winwittagree}(t, k, Q, t + p0 + \operatorname{twopow}(k + 1), L + 2)\), then \(\operatorname{winwittagree}(t + 1, k, Q, t + 1 + (p0 + 1) + \operatorname{twopow}(k + 1), L)\).
\end{theorem}
\begin{proof}
The goal follows from \(\operatorname{winwittagreestepouterleftstrip}(t, k, Q, p0, L)\).
\end{proof}
\begin{theorem}
\label{ll:paper:winwittagree-steps-outer-left-strip}
For every natural number \(t\) and natural number \(k\) and natural number \(Q\) and natural number \(p0\) and natural number \(L\) and natural number \(s\), if \(t + Q < p0 + \operatorname{twopow}(k)\), then if \(\operatorname{lhbandagree}(t, Q, p0, L + 2 \cdot s)\), then if \(\operatorname{winwittagree}(t, k, Q, t + p0 + \operatorname{twopow}(k + 1), L + 2 \cdot s)\), then \(\operatorname{winwittagree}(t + s, k, Q, t + s + (p0 + s) + \operatorname{twopow}(k + 1), L)\).
\end{theorem}
\begin{proof}
The goal follows from \(\operatorname{winwittagreestepsouterleftstrip}(t, k, Q, p0, L, s)\).
\end{proof}
\begin{theorem}
\label{ll:paper:winwitt-outward-step-outer-left-strip}
For every natural number \(t\) and natural number \(k\) and natural number \(p\), if \(t < p + \operatorname{twopow}(k)\), then \(\operatorname{wittbitbool}(\operatorname{wittwin}(t, \operatorname{twopow}(k)), t + (p + 2) + \operatorname{twopow}(k + 1)) = \operatorname{bxor}(\operatorname{bxor}(\operatorname{wittbitbool}(\operatorname{wittwin}(t + 1, \operatorname{twopow}(k)), t + 1 + (p + 1) + \operatorname{twopow}(k + 1)), \operatorname{wittbitbool}(\operatorname{wittwin}(t, \operatorname{twopow}(k)), t + (p + 1) + \operatorname{twopow}(k + 1)) \operatorname{or} \operatorname{wittbitbool}(\operatorname{wittwin}(t, \operatorname{twopow}(k)), t + p + \operatorname{twopow}(k + 1))), \operatorname{bxor}(\operatorname{lhcenter}(t, p + 1) \operatorname{and} \operatorname{wittbitbool}(\operatorname{wittwin}(t, \operatorname{twopow}(k)), t + p + \operatorname{twopow}(k + 1)), \operatorname{wittbitbool}(\operatorname{wittwin}(t, \operatorname{twopow}(k)), t + (p + 1) + \operatorname{twopow}(k + 1)) \operatorname{and} \operatorname{lhcenter}(t, p)))\).
\end{theorem}
\begin{proof}
The goal follows from \(\operatorname{winwittoutwardstepouterleftstrip}(t, k, p)\).
\end{proof}
\begin{theorem}
\label{ll:paper:winwittagree-extend-outer-left-strip}
For every natural number \(t\) and natural number \(k\) and natural number \(Q\) and natural number \(p0\) and natural number \(L\), if \(t + Q < p0 + \operatorname{twopow}(k)\), then if \(\operatorname{lhbandagree}(t, Q, p0, L + 1)\), then if \(\operatorname{winwittagree}(t, k, Q, t + p0 + \operatorname{twopow}(k + 1), L + 1)\), then if \(\operatorname{wittbitbool}(\operatorname{wittwin}(t + 1 + Q, \operatorname{twopow}(k)), t + 1 + (p0 + L + 1) + \operatorname{twopow}(k + 1) + Q) = \operatorname{wittbitbool}(\operatorname{wittwin}(t + 1, \operatorname{twopow}(k)), t + 1 + (p0 + L + 1) + \operatorname{twopow}(k + 1))\), then \(\operatorname{winwittagree}(t, k, Q, t + p0 + \operatorname{twopow}(k + 1), L + 2)\).
\end{theorem}
\begin{proof}
The goal follows from \(\operatorname{winwittagreeextendouterleftstrip}(t, k, Q, p0, L)\).
\end{proof}
\begin{theorem}
\label{ll:paper:winwittagree-outward-step-outer-left-strip}
For every natural number \(t\) and natural number \(k\) and natural number \(Q\) and natural number \(p0\) and natural number \(L\), if \(t + Q < p0 + \operatorname{twopow}(k)\), then if \(\operatorname{lhbandagree}(t, Q, p0, L + 1)\), then if \(\operatorname{winwittagree}(t, k, Q, t + p0 + \operatorname{twopow}(k + 1), L + 1)\), then if \(\operatorname{winwittagree}(t + 1, k, Q, t + 1 + (p0 + 1) + \operatorname{twopow}(k + 1), L)\), then \(\operatorname{winwittagree}(t, k, Q, t + (p0 + 2) + \operatorname{twopow}(k + 1), L)\).
\end{theorem}
\begin{proof}
The goal follows from \(\operatorname{winwittagreeoutwardstepouterleftstrip}(t, k, Q, p0, L)\).
\end{proof}
\begin{theorem}
\label{ll:paper:rhcenter-eq-lhcenter-scale-cross-outer}
For every natural number \(t\) and natural number \(k\) and natural number \(r\), if \(r \le \operatorname{twopow}(k)\), then if \(t + r < \operatorname{twopow}(k + 1)\), then \(\operatorname{rhcenter}(t, r) = \operatorname{bxor}(\operatorname{bxor}(\operatorname{lhcenter}(t + \operatorname{twopow}(k), \operatorname{twopow}(k) - r), \operatorname{lhcenter}(t, \operatorname{twopow}(k) - r)), \operatorname{wittbitbool}(\operatorname{wittwin}(t, \operatorname{twopow}(k)), t + \operatorname{twopow}(k + 1) - r))\).
\end{theorem}
\begin{proof}
The goal follows from \(\operatorname{rhcentereqlhcenterscalecrossouter}(t, k, r)\).
\end{proof}
\begin{theorem}
\label{ll:paper:rhbandagree-of-lhbandagree-scale-cross-outer}
For every natural number \(t\) and natural number \(k\) and natural number \(Q\) and natural number \(L\), if \(L \le \operatorname{twopow}(k)\), then if \(t + Q + L < \operatorname{twopow}(k + 1)\), then if \(\operatorname{lhbandagree}(t + \operatorname{twopow}(k), Q, \operatorname{twopow}(k) - L, L)\), then if \(\operatorname{lhbandagree}(t, Q, \operatorname{twopow}(k) - L, L)\), then if \(\operatorname{winwittagree}(t, k, Q, t + \operatorname{twopow}(k + 1) - L, L)\), then \(\operatorname{rhbandagree}(t, Q, 0, L)\).
\end{theorem}
\begin{proof}
The goal follows from \(\operatorname{rhbandagreeoflhbandagreescalecrossouter}(t, k, Q, L)\).
\end{proof}
\begin{theorem}
\label{ll:paper:winwittagree-of-rhbandagree-lhbandagree-scale-cross-outer}
For every natural number \(t\) and natural number \(k\) and natural number \(Q\) and natural number \(L\), if \(L \le \operatorname{twopow}(k)\), then if \(t + Q + L < \operatorname{twopow}(k + 1)\), then if \(\operatorname{rhbandagree}(t, Q, 0, L)\), then if \(\operatorname{lhbandagree}(t + \operatorname{twopow}(k), Q, \operatorname{twopow}(k) - L, L)\), then if \(\operatorname{lhbandagree}(t, Q, \operatorname{twopow}(k) - L, L)\), then \(\operatorname{winwittagree}(t, k, Q, t + \operatorname{twopow}(k + 1) - L, L)\).
\end{theorem}
\begin{proof}
The goal follows from \(\operatorname{winwittagreeofrhbandagreelhbaragreescalecrossouter}(t, k, Q, L)\).
\end{proof}
\begin{theorem}
\label{ll:paper:periodic-rhbandagree-of-lhcolagree-one-and-winwittagree-outer}
For every natural number \(T\) and natural number \(P\), if \(0 < P\), then if \(\operatorname{centerperiodic}(T, P)\), then for every natural number \(m\) and natural number \(t\) and natural number \(k\), if \(T \le t\), then if \(t + m \cdot P < \operatorname{twopow}(k)\), then if \(\operatorname{lhcolagree}(t, m \cdot P, 1, \operatorname{twopow}(k))\), then if \(\operatorname{lhcolagree}(t + \operatorname{twopow}(k), m \cdot P, 1, \operatorname{twopow}(k))\), then if \(\operatorname{winwittagree}(t, k, m \cdot P, t + \operatorname{twopow}(k), \operatorname{twopow}(k))\), then \(\operatorname{rhbandagree}(t, m \cdot P, 0, \operatorname{twopow}(k))\).
\end{theorem}
\begin{proof}
Assume \(hP\), \(hper\).
The goal follows from \(\operatorname{periodicrhbandagreeoflhcolagreeoneandwinwittagreeouter}(T, P, hP, hper)\).
\end{proof}
\begin{theorem}
\label{ll:paper:periodic-winwittagree-of-lhcolagree-one-and-rhbandagree-outer}
For every natural number \(T\) and natural number \(P\), if \(0 < P\), then if \(\operatorname{centerperiodic}(T, P)\), then for every natural number \(m\) and natural number \(t\) and natural number \(k\), if \(T \le t\), then if \(t + m \cdot P < \operatorname{twopow}(k)\), then if \(\operatorname{lhcolagree}(t, m \cdot P, 1, \operatorname{twopow}(k))\), then if \(\operatorname{lhcolagree}(t + \operatorname{twopow}(k), m \cdot P, 1, \operatorname{twopow}(k))\), then if \(\operatorname{rhbandagree}(t, m \cdot P, 0, \operatorname{twopow}(k))\), then \(\operatorname{winwittagree}(t, k, m \cdot P, t + \operatorname{twopow}(k), \operatorname{twopow}(k))\).
\end{theorem}
\begin{proof}
Assume \(hP\), \(hper\).
The goal follows from \(\operatorname{periodicwinwittagreeoflhcolagreeoneandrhbandagreeouter}(T, P, hP, hper)\).
\end{proof}
\begin{theorem}
\label{ll:paper:rhcenter-eq-lhcenter-scale-cross-cone}
For every natural number \(t\) and natural number \(k\) and natural number \(r\), if \(t + r < \operatorname{twopow}(k)\), then \(\operatorname{rhcenter}(t, r) = \operatorname{bxor}(\operatorname{lhcenter}(t + \operatorname{twopow}(k), \operatorname{twopow}(k) - r), \operatorname{wittbitbool}(\operatorname{wittwin}(t, \operatorname{twopow}(k)), t + \operatorname{twopow}(k + 1) - r))\).
\end{theorem}
\begin{proof}
The goal follows from \(\operatorname{rhcentereqlhcenterscalecrosscone}(t, k, r)\).
\end{proof}
\begin{theorem}
\label{ll:paper:lhcenter-scale-cross-cone}
For every natural number \(t\) and natural number \(k\) and natural number \(r\), if \(t + r < \operatorname{twopow}(k)\), then \(\operatorname{lhcenter}(t + \operatorname{twopow}(k), \operatorname{twopow}(k) - r) = \operatorname{bxor}(\operatorname{rhcenter}(t, r), \operatorname{wittbitbool}(\operatorname{wittwin}(t, \operatorname{twopow}(k)), t + \operatorname{twopow}(k + 1) - r))\).
\end{theorem}
\begin{proof}
The goal follows from \(\operatorname{lhcenterscalecrosscone}(t, k, r)\).
\end{proof}
\begin{theorem}
\label{ll:paper:winwitt-testbit-scale-cross-cone}
For every natural number \(t\) and natural number \(k\) and natural number \(r\), if \(t + r < \operatorname{twopow}(k)\), then \(\operatorname{wittbitbool}(\operatorname{wittwin}(t, \operatorname{twopow}(k)), t + \operatorname{twopow}(k + 1) - r) = \operatorname{bxor}(\operatorname{lhcenter}(t + \operatorname{twopow}(k), \operatorname{twopow}(k) - r), \operatorname{rhcenter}(t, r))\).
\end{theorem}
\begin{proof}
The goal follows from \(\operatorname{winwitttestbitscalecrosscone}(t, k, r)\).
\end{proof}
\begin{theorem}
\label{ll:paper:winwitt-testbit-scale-cross-zero-gap}
For every natural number \(t\) and natural number \(k\) and natural number \(r\), if \(t < r\), then if \(t + r < \operatorname{twopow}(k)\), then \(\operatorname{wittbitbool}(\operatorname{wittwin}(t, \operatorname{twopow}(k)), t + \operatorname{twopow}(k + 1) - r) = \operatorname{lhcenter}(t + \operatorname{twopow}(k), \operatorname{twopow}(k) - r)\).
\end{theorem}
\begin{proof}
The goal follows from \(\operatorname{winwitttestbitscalecrosszerogap}(t, k, r)\).
\end{proof}
\begin{theorem}
\label{ll:paper:winwitt-testbit-zero-gap-eq-lhcenter-late}
For every natural number \(t\) and natural number \(k\) and natural number \(p\), if \(t < p\), then if \(t + p < \operatorname{twopow}(k)\), then \(\operatorname{wittbitbool}(\operatorname{wittwin}(t, \operatorname{twopow}(k)), t + p + \operatorname{twopow}(k)) = \operatorname{lhcenter}(t + \operatorname{twopow}(k), p)\).
\end{theorem}
\begin{proof}
The goal follows from \(\operatorname{winwitttestbitzerogapeqlhcenterlate}(t, k, p)\).
\end{proof}
\begin{theorem}
\label{ll:paper:lhbandagree-of-winwittagree-zero-gap}
For every natural number \(t\) and natural number \(k\) and natural number \(Q\) and natural number \(p0\) and natural number \(L\), if \(t + Q < p0\), then if \(t + Q + p0 + L < \operatorname{twopow}(k)\), then if \(\operatorname{winwittagree}(t, k, Q, t + p0 + \operatorname{twopow}(k), L)\), then \(\operatorname{lhbandagree}(t + \operatorname{twopow}(k), Q, p0, L)\).
\end{theorem}
\begin{proof}
The goal follows from \(\operatorname{lhbandagreeofwinwittagreezerogap}(t, k, Q, p0, L)\).
\end{proof}
\begin{theorem}
\label{ll:paper:winwittagree-of-lhbandagree-zero-gap}
For every natural number \(t\) and natural number \(k\) and natural number \(Q\) and natural number \(p0\) and natural number \(L\), if \(t + Q < p0\), then if \(t + Q + p0 + L < \operatorname{twopow}(k)\), then if \(\operatorname{lhbandagree}(t + \operatorname{twopow}(k), Q, p0, L)\), then \(\operatorname{winwittagree}(t, k, Q, t + p0 + \operatorname{twopow}(k), L)\).
\end{theorem}
\begin{proof}
The goal follows from \(\operatorname{winwittagreeoflhbandagreezerogap}(t, k, Q, p0, L)\).
\end{proof}
\begin{theorem}
\label{ll:paper:rhbandagree-of-lhbandagree-and-winwittagree-cross-cone}
For every natural number \(t\) and natural number \(k\) and natural number \(Q\) and natural number \(L\), if \(t + Q + L < \operatorname{twopow}(k)\), then if \(\operatorname{lhbandagree}(t + \operatorname{twopow}(k), Q, \operatorname{twopow}(k) - L, L)\), then if \(\operatorname{winwittagree}(t, k, Q, t + \operatorname{twopow}(k + 1) - L, L)\), then \(\operatorname{rhbandagree}(t, Q, 0, L)\).
\end{theorem}
\begin{proof}
The goal follows from \(\operatorname{rhbandagreeoflhbandagreeandwinwittagreecrosscone}(t, k, Q, L)\).
\end{proof}
\begin{theorem}
\label{ll:paper:winwittagree-of-rhbandagree-and-lhbandagree-cross-cone}
For every natural number \(t\) and natural number \(k\) and natural number \(Q\) and natural number \(L\), if \(t + Q + L < \operatorname{twopow}(k)\), then if \(\operatorname{rhbandagree}(t, Q, 0, L)\), then if \(\operatorname{lhbandagree}(t + \operatorname{twopow}(k), Q, \operatorname{twopow}(k) - L, L)\), then \(\operatorname{winwittagree}(t, k, Q, t + \operatorname{twopow}(k + 1) - L, L)\).
\end{theorem}
\begin{proof}
The goal follows from \(\operatorname{winwittagreeofrhbandagreeandlhbandagreecrosscone}(t, k, Q, L)\).
\end{proof}
\begin{theorem}
\label{ll:paper:lhbandagree-of-rhbandagree-and-winwittagree-cross-cone}
For every natural number \(t\) and natural number \(k\) and natural number \(Q\) and natural number \(L\), if \(t + Q + L < \operatorname{twopow}(k)\), then if \(\operatorname{rhbandagree}(t, Q, 0, L)\), then if \(\operatorname{winwittagree}(t, k, Q, t + \operatorname{twopow}(k + 1) - L, L)\), then \(\operatorname{lhbandagree}(t + \operatorname{twopow}(k), Q, \operatorname{twopow}(k) - L, L)\).
\end{theorem}
\begin{proof}
The goal follows from \(\operatorname{lhbandagreeofrhbandagreeandwinwittagreecrosscone}(t, k, Q, L)\).
\end{proof}
\begin{theorem}
\label{ll:paper:periodic-rhbandagree-of-lhcolagree-late-and-winwittagree-cone}
For every natural number \(T\) and natural number \(P\), if \(0 < P\), then if \(\operatorname{centerperiodic}(T, P)\), then for every natural number \(m\) and natural number \(t\) and natural number \(k\) and natural number \(L\), if \(T \le t\), then if \(t + m \cdot P + L < \operatorname{twopow}(k)\), then if \(\operatorname{lhcolagree}(t + \operatorname{twopow}(k), m \cdot P, 1, \operatorname{twopow}(k))\), then if \(\operatorname{winwittagree}(t, k, m \cdot P, t + \operatorname{twopow}(k + 1) - L, L)\), then \(\operatorname{rhbandagree}(t, m \cdot P, 0, L)\).
\end{theorem}
\begin{proof}
Assume \(hP\), \(hper\).
The goal follows from \(\operatorname{periodicrhbandagreeoflhcolagreelateandwinwittagreecone}(T, P, hP, hper)\).
\end{proof}
\begin{theorem}
\label{ll:paper:periodic-winwittagree-of-lhcolagree-late-and-rhbandagree-cone}
For every natural number \(T\) and natural number \(P\), if \(0 < P\), then if \(\operatorname{centerperiodic}(T, P)\), then for every natural number \(m\) and natural number \(t\) and natural number \(k\) and natural number \(L\), if \(T \le t\), then if \(t + m \cdot P + L < \operatorname{twopow}(k)\), then if \(\operatorname{lhcolagree}(t + \operatorname{twopow}(k), m \cdot P, 1, \operatorname{twopow}(k))\), then if \(\operatorname{rhbandagree}(t, m \cdot P, 0, L)\), then \(\operatorname{winwittagree}(t, k, m \cdot P, t + \operatorname{twopow}(k + 1) - L, L)\).
\end{theorem}
\begin{proof}
Assume \(hP\), \(hper\).
The goal follows from \(\operatorname{periodicwinwittagreeoflhcolagreelateandrhbandagreecone}(T, P, hP, hper)\).
\end{proof}
\begin{theorem}
\label{ll:paper:lhcenter-step-outer-strip-from-winwitt}
For every natural number \(t\) and natural number \(k\) and natural number \(p\), if \(t < p + \operatorname{twopow}(k)\), then \(\operatorname{lhcenter}(t + 1 + \operatorname{twopow}(k), p + 1 + \operatorname{twopow}(k)) = \operatorname{bxor}(\operatorname{bxor}(\operatorname{lhcenter}(t, p + 2), \operatorname{wittbitbool}(\operatorname{wittwin}(t, \operatorname{twopow}(k)), t + (p + 2) + \operatorname{twopow}(k + 1))), \operatorname{bxor}(\operatorname{lhcenter}(t, p + 1), \operatorname{wittbitbool}(\operatorname{wittwin}(t, \operatorname{twopow}(k)), t + (p + 1) + \operatorname{twopow}(k + 1))) \operatorname{or} \operatorname{bxor}(\operatorname{lhcenter}(t, p), \operatorname{wittbitbool}(\operatorname{wittwin}(t, \operatorname{twopow}(k)), t + p + \operatorname{twopow}(k + 1))))\).
\end{theorem}
\begin{proof}
The goal follows from \(\operatorname{lhcenterstepouterstripfromwinwitt}(t, k, p)\).
\end{proof}
\begin{theorem}
\label{ll:paper:winwitt-step-lhcenter-defect-uncond}
For every natural number \(t\) and natural number \(k\) and natural number \(p\), \(\operatorname{wittbitbool}(\operatorname{wittwin}(t + 1, \operatorname{twopow}(k)), t + 1 + (p + 1) + \operatorname{twopow}(k + 1)) = \operatorname{bxor}(\operatorname{wittbitbool}(\operatorname{wittwin}(t, \operatorname{twopow}(k)), t + (p + 2) + \operatorname{twopow}(k + 1)), \operatorname{bxor}(\operatorname{bxor}(\operatorname{bxor}(\operatorname{lhcenter}(t + \operatorname{twopow}(k), p + 1 + \operatorname{twopow}(k)) \operatorname{or} \operatorname{lhcenter}(t + \operatorname{twopow}(k), p + \operatorname{twopow}(k)), \operatorname{lhcenter}(t, p + 1) \operatorname{or} \operatorname{lhcenter}(t, p)), \operatorname{lhcenter}(t, p + 1 + \operatorname{twopow}(k)) \operatorname{or} \operatorname{lhcenter}(t, p + \operatorname{twopow}(k))), \operatorname{lhcenter}(t, p + 1 + \operatorname{twopow}(k + 1)) \operatorname{or} \operatorname{lhcenter}(t, p + \operatorname{twopow}(k + 1))))\).
\end{theorem}
\begin{proof}
The goal follows from \(\operatorname{winwittsteplhcenterdefectuncond}(t, k, p)\).
\end{proof}
\begin{theorem}
\label{ll:paper:winwittagree-step-of-lhbandagree-scale-uncond}
For every natural number \(t\) and natural number \(k\) and natural number \(Q\) and natural number \(p0\) and natural number \(L\), if \(\operatorname{lhbandagree}(t + \operatorname{twopow}(k), Q, p0 + \operatorname{twopow}(k), L + 1)\), then if \(\operatorname{lhbandagree}(t, Q, p0, L + 1)\), then if \(\operatorname{lhbandagree}(t, Q, p0 + \operatorname{twopow}(k), L + 1)\), then if \(\operatorname{lhbandagree}(t, Q, p0 + \operatorname{twopow}(k + 1), L + 1)\), then if \(\operatorname{winwittagree}(t, k, Q, t + (p0 + 2) + \operatorname{twopow}(k + 1), L)\), then \(\operatorname{winwittagree}(t + 1, k, Q, t + 1 + (p0 + 1) + \operatorname{twopow}(k + 1), L)\).
\end{theorem}
\begin{proof}
The goal follows from \(\operatorname{winwittagreestepoflhbandagreescaleuncond}(t, k, Q, p0, L)\).
\end{proof}
\begin{theorem}
\label{ll:paper:winwittagree-outward-of-lhbandagree-scale-uncond}
For every natural number \(t\) and natural number \(k\) and natural number \(Q\) and natural number \(p0\) and natural number \(L\), if \(\operatorname{lhbandagree}(t + \operatorname{twopow}(k), Q, p0 + \operatorname{twopow}(k), L + 1)\), then if \(\operatorname{lhbandagree}(t, Q, p0, L + 1)\), then if \(\operatorname{lhbandagree}(t, Q, p0 + \operatorname{twopow}(k), L + 1)\), then if \(\operatorname{lhbandagree}(t, Q, p0 + \operatorname{twopow}(k + 1), L + 1)\), then if \(\operatorname{winwittagree}(t + 1, k, Q, t + 1 + (p0 + 1) + \operatorname{twopow}(k + 1), L)\), then \(\operatorname{winwittagree}(t, k, Q, t + (p0 + 2) + \operatorname{twopow}(k + 1), L)\).
\end{theorem}
\begin{proof}
The goal follows from \(\operatorname{winwittagreeoutwardoflhbandagreescaleuncond}(t, k, Q, p0, L)\).
\end{proof}
\begin{theorem}
\label{ll:paper:phiscale-row-testbit-lhcenter-sector-one}
For every natural number \(t\) and natural number \(k\) and natural number \(p\), if \(t + p < \operatorname{twopow}(k)\), then \(\operatorname{wittbitbool}(\operatorname{wittscale}(k, \operatorname{wittrow}(t)), t + p + \operatorname{twopow}(k)) = \operatorname{lhcenter}(t, p)\).
\end{theorem}
\begin{proof}
The goal follows from \(\operatorname{phiscalerowtestbitlhcentersectorone}(t, k, p)\).
\end{proof}
\begin{theorem}
\label{ll:paper:lhcenter-scale-same-col-cone}
For every natural number \(t\) and natural number \(k\) and natural number \(p\), if \(t + p < \operatorname{twopow}(k)\), then \(\operatorname{lhcenter}(t + \operatorname{twopow}(k), p) = \operatorname{bxor}(\operatorname{lhcenter}(t, p), \operatorname{wittbitbool}(\operatorname{wittwin}(t, \operatorname{twopow}(k)), t + p + \operatorname{twopow}(k)))\).
\end{theorem}
\begin{proof}
The goal follows from \(\operatorname{lhcenterscalesamecolcone}(t, k, p)\).
\end{proof}
\begin{theorem}
\label{ll:paper:lhcenter-scale-inward-same-col-cone}
For every natural number \(t\) and natural number \(k\) and natural number \(p\), if \(t + p < \operatorname{twopow}(k)\), then \(\operatorname{lhcenter}(t, p) = \operatorname{bxor}(\operatorname{lhcenter}(t + \operatorname{twopow}(k), p), \operatorname{wittbitbool}(\operatorname{wittwin}(t, \operatorname{twopow}(k)), t + p + \operatorname{twopow}(k)))\).
\end{theorem}
\begin{proof}
The goal follows from \(\operatorname{lhcenterscaleinwardsamecolcone}(t, k, p)\).
\end{proof}
\begin{theorem}
\label{ll:paper:winwitt-testbit-sector-one-lhcenter}
For every natural number \(t\) and natural number \(k\) and natural number \(p\), if \(t + p < \operatorname{twopow}(k)\), then \(\operatorname{wittbitbool}(\operatorname{wittwin}(t, \operatorname{twopow}(k)), t + p + \operatorname{twopow}(k)) = \operatorname{bxor}(\operatorname{lhcenter}(t + \operatorname{twopow}(k), p), \operatorname{lhcenter}(t, p))\).
\end{theorem}
\begin{proof}
The goal follows from \(\operatorname{winwitttestbitsectoronelhcenter}(t, k, p)\).
\end{proof}
\begin{theorem}
\label{ll:paper:lhbandagree-scale-same-col-cone}
For every natural number \(t\) and natural number \(k\) and natural number \(Q\) and natural number \(p0\) and natural number \(L\), if \(t + Q + p0 + L < \operatorname{twopow}(k)\), then if \(\operatorname{lhbandagree}(t, Q, p0, L)\), then if \(\operatorname{winwittagree}(t, k, Q, t + p0 + \operatorname{twopow}(k), L)\), then \(\operatorname{lhbandagree}(t + \operatorname{twopow}(k), Q, p0, L)\).
\end{theorem}
\begin{proof}
The goal follows from \(\operatorname{lhbandagreescalesamecolcone}(t, k, Q, p0, L)\).
\end{proof}
\begin{theorem}
\label{ll:paper:lhbandagree-scale-inward-same-col-cone}
For every natural number \(t\) and natural number \(k\) and natural number \(Q\) and natural number \(p0\) and natural number \(L\), if \(t + Q + p0 + L < \operatorname{twopow}(k)\), then if \(\operatorname{lhbandagree}(t + \operatorname{twopow}(k), Q, p0, L)\), then if \(\operatorname{winwittagree}(t, k, Q, t + p0 + \operatorname{twopow}(k), L)\), then \(\operatorname{lhbandagree}(t, Q, p0, L)\).
\end{theorem}
\begin{proof}
The goal follows from \(\operatorname{lhbandagreescaleinwardsamecolcone}(t, k, Q, p0, L)\).
\end{proof}
\begin{theorem}
\label{ll:paper:winwittagree-sector-one-lhcenter}
For every natural number \(t\) and natural number \(k\) and natural number \(Q\) and natural number \(p0\) and natural number \(L\), if \(t + Q + p0 + L < \operatorname{twopow}(k)\), then if \(\operatorname{lhbandagree}(t + \operatorname{twopow}(k), Q, p0, L)\), then if \(\operatorname{lhbandagree}(t, Q, p0, L)\), then \(\operatorname{winwittagree}(t, k, Q, t + p0 + \operatorname{twopow}(k), L)\).
\end{theorem}
\begin{proof}
The goal follows from \(\operatorname{winwittagreesectoronelhcenter}(t, k, Q, p0, L)\).
\end{proof}
\begin{theorem}
\label{ll:paper:not-lhbandagree-late-and-winwittagree-sector-one}
For every natural number \(t\) and natural number \(k\) and natural number \(Q\) and natural number \(L\), if \(0 < Q\), then if \(t + Q \le L\), then if \(t + Q + L < \operatorname{twopow}(k)\), then if \(\operatorname{lhbandagree}(t + \operatorname{twopow}(k), Q, 0, L)\), then not \(\operatorname{winwittagree}(t, k, Q, t + \operatorname{twopow}(k), L)\).
\end{theorem}
\begin{proof}
The goal follows from \(\operatorname{notlhbandagreelateandwinwittagreesectorone}(t, k, Q, L)\).
\end{proof}
\begin{theorem}
\label{ll:paper:frobdefect-row-testbit-lhcenter-sector-one-eq-false}
For every natural number \(t\) and natural number \(k\) and natural number \(p\), if \(t + p + 1 < \operatorname{twopow}(k)\), then \(\operatorname{wittbitbool}(\operatorname{frobdefect}(k, \operatorname{wittrow}(t)), t + p + 1 + \operatorname{twopow}(k)) = \operatorname{false}\).
\end{theorem}
\begin{proof}
The goal follows from \(\operatorname{frobdefectrowtestbitlhcentersectoroneeqfalse}(t, k, p)\).
\end{proof}
\begin{theorem}
\label{ll:paper:crosscarry-row-scale-lhcenter-sector-one}
For every natural number \(t\) and natural number \(k\) and natural number \(p\), if \(t + p + 1 < \operatorname{twopow}(k)\), then \(\operatorname{wittbitbool}(\operatorname{crosscarry}(\operatorname{wittscale}(k, \operatorname{wittrow}(t)), \operatorname{wittwin}(t, \operatorname{twopow}(k))), t + p + 1 + \operatorname{twopow}(k)) = \operatorname{bxor}(\operatorname{lhcenter}(t, p + 1) \operatorname{and} \operatorname{wittbitbool}(\operatorname{wittwin}(t, \operatorname{twopow}(k)), t + p + \operatorname{twopow}(k)), \operatorname{wittbitbool}(\operatorname{wittwin}(t, \operatorname{twopow}(k)), t + (p + 1) + \operatorname{twopow}(k)) \operatorname{and} \operatorname{lhcenter}(t, p))\).
\end{theorem}
\begin{proof}
The goal follows from \(\operatorname{crosscarryrowscalelhcentersectorone}(t, k, p)\).
\end{proof}
\begin{theorem}
\label{ll:paper:winwitt-step-sector-one-lhcenter}
For every natural number \(t\) and natural number \(k\) and natural number \(p\), if \(t + p + 1 < \operatorname{twopow}(k)\), then \(\operatorname{wittbitbool}(\operatorname{wittwin}(t + 1, \operatorname{twopow}(k)), t + 1 + (p + 1) + \operatorname{twopow}(k)) = \operatorname{bxor}(\operatorname{wittbitbool}(\operatorname{wittphi}(\operatorname{wittwin}(t, \operatorname{twopow}(k))), t + 1 + (p + 1) + \operatorname{twopow}(k)), \operatorname{bxor}(\operatorname{lhcenter}(t, p + 1) \operatorname{and} \operatorname{wittbitbool}(\operatorname{wittwin}(t, \operatorname{twopow}(k)), t + p + \operatorname{twopow}(k)), \operatorname{wittbitbool}(\operatorname{wittwin}(t, \operatorname{twopow}(k)), t + (p + 1) + \operatorname{twopow}(k)) \operatorname{and} \operatorname{lhcenter}(t, p)))\).
\end{theorem}
\begin{proof}
The goal follows from \(\operatorname{winwittstepsectoronelhcenter}(t, k, p)\).
\end{proof}
\begin{theorem}
\label{ll:paper:winwittagree-step-sector-one-lhcenter}
For every natural number \(t\) and natural number \(k\) and natural number \(Q\) and natural number \(p0\) and natural number \(L\), if \(t + Q + p0 + L + 1 < \operatorname{twopow}(k)\), then if \(\operatorname{lhbandagree}(t, Q, p0, L + 1)\), then if \(\operatorname{winwittagree}(t, k, Q, t + p0 + \operatorname{twopow}(k), L + 2)\), then \(\operatorname{winwittagree}(t + 1, k, Q, t + 1 + (p0 + 1) + \operatorname{twopow}(k), L)\).
\end{theorem}
\begin{proof}
The goal follows from \(\operatorname{winwittagreestepsectoronelhcenter}(t, k, Q, p0, L)\).
\end{proof}
\begin{theorem}
\label{ll:paper:winwittagree-outward-step-sector-one-lhcenter}
For every natural number \(t\) and natural number \(k\) and natural number \(Q\) and natural number \(p0\) and natural number \(L\), if \(t + Q + p0 + L + 1 < \operatorname{twopow}(k)\), then if \(\operatorname{lhbandagree}(t, Q, p0, L + 1)\), then if \(\operatorname{winwittagree}(t, k, Q, t + p0 + \operatorname{twopow}(k), L + 1)\), then if \(\operatorname{winwittagree}(t + 1, k, Q, t + 1 + (p0 + 1) + \operatorname{twopow}(k), L)\), then \(\operatorname{winwittagree}(t, k, Q, t + p0 + \operatorname{twopow}(k), L + 2)\).
\end{theorem}
\begin{proof}
The goal follows from \(\operatorname{winwittagreeoutwardstepsectoronelhcenter}(t, k, Q, p0, L)\).
\end{proof}
\begin{theorem}
\label{ll:paper:periodic-winwitt-sector-one-zero-eq}
For every natural number \(T\) and natural number \(P\), if \(0 < P\), then if \(\operatorname{centerperiodic}(T, P)\), then for every natural number \(m\) and natural number \(t\) and natural number \(k\), if \(T \le t\), then if \(t + m \cdot P < \operatorname{twopow}(k)\), then \(\operatorname{wittbitbool}(\operatorname{wittwin}(t + m \cdot P, \operatorname{twopow}(k)), t + m \cdot P + \operatorname{twopow}(k)) = \operatorname{wittbitbool}(\operatorname{wittwin}(t, \operatorname{twopow}(k)), t + \operatorname{twopow}(k))\).
\end{theorem}
\begin{proof}
Assume \(hP\), \(hper\).
The goal follows from \(\operatorname{periodicwinwittsectoronezeroeq}(T, P, hP, hper)\).
\end{proof}
\begin{theorem}
\label{ll:paper:periodic-winwitt-sector-one-two-defect-eq}
For every natural number \(T\) and natural number \(P\), if \(0 < P\), then if \(\operatorname{centerperiodic}(T, P)\), then for every natural number \(m\) and natural number \(t\) and natural number \(k\), if \(T \le t\), then if \(t + m \cdot P + 2 < \operatorname{twopow}(k)\), then \(\operatorname{bxor}(\operatorname{wittbitbool}(\operatorname{wittwin}(t + m \cdot P, \operatorname{twopow}(k)), t + m \cdot P + 2 + \operatorname{twopow}(k)), \operatorname{wittbitbool}(\operatorname{wittwin}(t, \operatorname{twopow}(k)), t + 2 + \operatorname{twopow}(k))) = \operatorname{bxor}(\operatorname{bxor}(\operatorname{wittbitbool}(\operatorname{wittwin}(t + 1 + m \cdot P, \operatorname{twopow}(k)), t + 1 + m \cdot P + 1 + \operatorname{twopow}(k)), \operatorname{wittbitbool}(\operatorname{wittwin}(t + 1, \operatorname{twopow}(k)), t + 1 + 1 + \operatorname{twopow}(k))), \operatorname{bxor}(\operatorname{wittbitbool}(\operatorname{wittwin}(t + m \cdot P, \operatorname{twopow}(k)), t + m \cdot P + 1 + \operatorname{twopow}(k)), \operatorname{wittbitbool}(\operatorname{wittwin}(t, \operatorname{twopow}(k)), t + 1 + \operatorname{twopow}(k))))\).
\end{theorem}
\begin{proof}
Assume \(hP\), \(hper\).
The goal follows from \(\operatorname{periodicwinwittsectoronetwodefecteq}(T, P, hP, hper)\).
\end{proof}
\begin{theorem}
\label{ll:paper:not-rhbandagree-of-lt-add}
For every natural number \(t\) and natural number \(Q\) and natural number \(L\), if \(0 < Q\), then if \(t + Q \le L\), then not \(\operatorname{rhbandagree}(t, Q, 0, L)\).
\end{theorem}
\begin{proof}
The goal follows from \(\operatorname{notrhbandagreeofltadd}(t, Q, L)\).
\end{proof}
\begin{theorem}
\label{ll:paper:phiscale-row-testbit-rhcenter-cone}
For every natural number \(t\) and natural number \(k\) and natural number \(p\), if \(t + p < \operatorname{twopow}(k)\), then \(\operatorname{wittbitbool}(\operatorname{wittscale}(k, \operatorname{wittrow}(t)), t + \operatorname{twopow}(k) - p) = \operatorname{rhcenter}(t, p)\).
\end{theorem}
\begin{proof}
The goal follows from \(\operatorname{phiscalerowtestbitrhcentercone}(t, k, p)\).
\end{proof}
\begin{theorem}
\label{ll:paper:rhcenter-scale-same-col-cone}
For every natural number \(t\) and natural number \(k\) and natural number \(p\), if \(t + p < \operatorname{twopow}(k)\), then \(\operatorname{rhcenter}(t + \operatorname{twopow}(k), p) = \operatorname{bxor}(\operatorname{rhcenter}(t, p), \operatorname{wittbitbool}(\operatorname{wittwin}(t, \operatorname{twopow}(k)), t + \operatorname{twopow}(k) - p))\).
\end{theorem}
\begin{proof}
The goal follows from \(\operatorname{rhcenterscalesamecolcone}(t, k, p)\).
\end{proof}
\begin{theorem}
\label{ll:paper:rhcenter-scale-inward-same-col-cone}
For every natural number \(t\) and natural number \(k\) and natural number \(p\), if \(t + p < \operatorname{twopow}(k)\), then \(\operatorname{rhcenter}(t, p) = \operatorname{bxor}(\operatorname{rhcenter}(t + \operatorname{twopow}(k), p), \operatorname{wittbitbool}(\operatorname{wittwin}(t, \operatorname{twopow}(k)), t + \operatorname{twopow}(k) - p))\).
\end{theorem}
\begin{proof}
The goal follows from \(\operatorname{rhcenterscaleinwardsamecolcone}(t, k, p)\).
\end{proof}
\begin{theorem}
\label{ll:paper:winwitt-testbit-rhcenter-same-col-cone}
For every natural number \(t\) and natural number \(k\) and natural number \(p\), if \(t + p < \operatorname{twopow}(k)\), then \(\operatorname{wittbitbool}(\operatorname{wittwin}(t, \operatorname{twopow}(k)), t + \operatorname{twopow}(k) - p) = \operatorname{bxor}(\operatorname{rhcenter}(t + \operatorname{twopow}(k), p), \operatorname{rhcenter}(t, p))\).
\end{theorem}
\begin{proof}
The goal follows from \(\operatorname{winwitttestbitrhcentersamecolcone}(t, k, p)\).
\end{proof}
\begin{theorem}
\label{ll:paper:winwitt-testbit-zero-gap-eq-rhcenter-late}
For every natural number \(t\) and natural number \(k\) and natural number \(p\), if \(t < p\), then if \(t + p < \operatorname{twopow}(k)\), then \(\operatorname{wittbitbool}(\operatorname{wittwin}(t, \operatorname{twopow}(k)), t + \operatorname{twopow}(k) - p) = \operatorname{rhcenter}(t + \operatorname{twopow}(k), p)\).
\end{theorem}
\begin{proof}
The goal follows from \(\operatorname{winwitttestbitzerogapeqrhcenterlate}(t, k, p)\).
\end{proof}
\begin{theorem}
\label{ll:paper:rhbandagree-scale-same-col-cone}
For every natural number \(t\) and natural number \(k\) and natural number \(Q\) and natural number \(p0\) and natural number \(L\), if \(t + Q + p0 + L < \operatorname{twopow}(k)\), then if \(\operatorname{rhbandagree}(t, Q, p0, L)\), then if \(\operatorname{winwittagree}(t, k, Q, t + \operatorname{twopow}(k) - (p0 + L), L)\), then \(\operatorname{rhbandagree}(t + \operatorname{twopow}(k), Q, p0, L)\).
\end{theorem}
\begin{proof}
The goal follows from \(\operatorname{rhbandagreescalesamecolcone}(t, k, Q, p0, L)\).
\end{proof}
\begin{theorem}
\label{ll:paper:rhbandagree-scale-inward-same-col-cone}
For every natural number \(t\) and natural number \(k\) and natural number \(Q\) and natural number \(p0\) and natural number \(L\), if \(t + Q + p0 + L < \operatorname{twopow}(k)\), then if \(\operatorname{rhbandagree}(t + \operatorname{twopow}(k), Q, p0, L)\), then if \(\operatorname{winwittagree}(t, k, Q, t + \operatorname{twopow}(k) - (p0 + L), L)\), then \(\operatorname{rhbandagree}(t, Q, p0, L)\).
\end{theorem}
\begin{proof}
The goal follows from \(\operatorname{rhbandagreescaleinwardsamecolcone}(t, k, Q, p0, L)\).
\end{proof}
\begin{theorem}
\label{ll:paper:winwittagree-rhcenter-same-col-cone}
For every natural number \(t\) and natural number \(k\) and natural number \(Q\) and natural number \(p0\) and natural number \(L\), if \(t + Q + p0 + L < \operatorname{twopow}(k)\), then if \(\operatorname{rhbandagree}(t + \operatorname{twopow}(k), Q, p0, L)\), then if \(\operatorname{rhbandagree}(t, Q, p0, L)\), then \(\operatorname{winwittagree}(t, k, Q, t + \operatorname{twopow}(k) - (p0 + L), L)\).
\end{theorem}
\begin{proof}
The goal follows from \(\operatorname{winwittagreerhcentersamecolcone}(t, k, Q, p0, L)\).
\end{proof}
\begin{theorem}
\label{ll:paper:not-rhbandagree-late-and-winwittagree-cone}
For every natural number \(t\) and natural number \(k\) and natural number \(Q\) and natural number \(L\), if \(0 < Q\), then if \(t + Q \le L\), then if \(t + Q + L < \operatorname{twopow}(k)\), then if \(\operatorname{rhbandagree}(t + \operatorname{twopow}(k), Q, 0, L)\), then not \(\operatorname{winwittagree}(t, k, Q, t + \operatorname{twopow}(k) - L, L)\).
\end{theorem}
\begin{proof}
The goal follows from \(\operatorname{notrhbandagreelateandwinwittagreecone}(t, k, Q, L)\).
\end{proof}
\begin{theorem}
\label{ll:paper:lhbandagree-of-rowagree-symm}
For every natural number \(t\) and natural number \(Q\) and natural number \(R\), if \(R \le t\), then if \(\operatorname{rowagree}(t, Q, t - R, 2 \cdot R)\), then \(\operatorname{lhbandagree}(t, Q, 0, R)\).
\end{theorem}
\begin{proof}
The goal follows from \(\operatorname{lhbandagreeofrowagreesymm}(t, Q, R)\).
\end{proof}
\begin{theorem}
\label{ll:paper:periodic-lhbandagree-zero-of-rhbandagree-zero}
For every natural number \(T\) and natural number \(P\), if \(0 < P\), then if \(\operatorname{centerperiodic}(T, P)\), then for every natural number \(m\) and natural number \(t\) and natural number \(R\), if \(T \le t\), then if \(R \le t\), then if \(\operatorname{rhbandagree}(t, m \cdot P, 0, R)\), then \(\operatorname{lhbandagree}(t, m \cdot P, 0, R)\).
\end{theorem}
\begin{proof}
Assume \(hP\), \(hper\).
The goal follows from \(\operatorname{periodiclhbandagreezeroofrhbandagreezero}(T, P, hP, hper)\).
\end{proof}
\begin{theorem}
\label{ll:paper:periodic-lhbandagree-late-of-rhbandagree-late}
For every natural number \(T\) and natural number \(P\), if \(0 < P\), then if \(\operatorname{centerperiodic}(T, P)\), then for every natural number \(m\) and natural number \(t\) and natural number \(k\) and natural number \(R\), if \(T \le t\), then if \(R \le \operatorname{twopow}(k)\), then if \(\operatorname{rhbandagree}(t + \operatorname{twopow}(k), m \cdot P, 0, R)\), then \(\operatorname{lhbandagree}(t + \operatorname{twopow}(k), m \cdot P, 0, R)\).
\end{theorem}
\begin{proof}
Assume \(hP\), \(hper\).
The goal follows from \(\operatorname{periodiclhbandagreelateofrhbandagreelate}(T, P, hP, hper)\).
\end{proof}
\begin{theorem}
\label{ll:paper:not-periodic-rhbandagree-late-and-winwittagree-sector-one}
For every natural number \(T\) and natural number \(P\), if \(0 < P\), then if \(\operatorname{centerperiodic}(T, P)\), then for every natural number \(m\) and natural number \(t\) and natural number \(k\) and natural number \(L\), if \(0 < m\), then if \(T \le t\), then if \(t + m \cdot P \le L\), then if \(t + m \cdot P + L < \operatorname{twopow}(k)\), then if \(\operatorname{rhbandagree}(t + \operatorname{twopow}(k), m \cdot P, 0, L)\), then not \(\operatorname{winwittagree}(t, k, m \cdot P, t + \operatorname{twopow}(k), L)\).
\end{theorem}
\begin{proof}
Assume \(hP\), \(hper\).
The goal follows from \(\operatorname{notperiodicrhbandagreelateandwinwittagreesectorone}(T, P, hP, hper)\).
\end{proof}
\begin{theorem}
\label{ll:paper:periodic-winwittagree-sector-one-left-of-right}
For every natural number \(T\) and natural number \(P\), if \(0 < P\), then if \(\operatorname{centerperiodic}(T, P)\), then for every natural number \(m\) and natural number \(t\) and natural number \(k\) and natural number \(R\), if \(T \le t\), then if \(R \le t\), then if \(t + m \cdot P + R < \operatorname{twopow}(k)\), then if \(\operatorname{rhbandagree}(t, m \cdot P, 0, R)\), then if \(\operatorname{winwittagree}(t, k, m \cdot P, t + \operatorname{twopow}(k) - R, R)\), then \(\operatorname{winwittagree}(t, k, m \cdot P, t + \operatorname{twopow}(k), R)\).
\end{theorem}
\begin{proof}
Assume \(hP\), \(hper\).
The goal follows from \(\operatorname{periodicwinwittagreesectoroneleftofright}(T, P, hP, hper)\).
\end{proof}
\begin{theorem}
\label{ll:paper:frobdefect-row-testbit-gt-two-mul-eq-false}
For every natural number \(t\) and natural number \(k\) and natural number \(b\), if \(2 \cdot t + 1 < b\), then if \(b < \operatorname{twopow}(k) + \operatorname{twopow}(k)\), then \(\operatorname{wittbitbool}(\operatorname{frobdefect}(k, \operatorname{wittrow}(t)), b) = \operatorname{false}\).
\end{theorem}
\begin{proof}
The goal follows from \(\operatorname{frobdefectrowtestbitgttwomuleqfalse}(t, k, b)\).
\end{proof}
\begin{theorem}
\label{ll:paper:frobdefect-row-testbit-rhcenter-cone-eq-false}
For every natural number \(t\) and natural number \(k\) and natural number \(p\), if \(t + p + 1 < \operatorname{twopow}(k)\), then \(\operatorname{wittbitbool}(\operatorname{frobdefect}(k, \operatorname{wittrow}(t)), t + \operatorname{twopow}(k) - p) = \operatorname{false}\).
\end{theorem}
\begin{proof}
The goal follows from \(\operatorname{frobdefectrowtestbitrhcenterconeeqfalse}(t, k, p)\).
\end{proof}
\begin{theorem}
\label{ll:paper:crosscarry-row-scale-rhcenter-cone}
For every natural number \(t\) and natural number \(k\) and natural number \(p\), if \(t + p + 1 < \operatorname{twopow}(k)\), then \(\operatorname{wittbitbool}(\operatorname{crosscarry}(\operatorname{wittscale}(k, \operatorname{wittrow}(t)), \operatorname{wittwin}(t, \operatorname{twopow}(k))), t + \operatorname{twopow}(k) - p) = \operatorname{bxor}(\operatorname{rhcenter}(t, p) \operatorname{and} \operatorname{wittbitbool}(\operatorname{wittwin}(t, \operatorname{twopow}(k)), t + \operatorname{twopow}(k) - (p + 1)), \operatorname{wittbitbool}(\operatorname{wittwin}(t, \operatorname{twopow}(k)), t + \operatorname{twopow}(k) - p) \operatorname{and} \operatorname{rhcenter}(t, p + 1))\).
\end{theorem}
\begin{proof}
The goal follows from \(\operatorname{crosscarryrowscalerhcentercone}(t, k, p)\).
\end{proof}
\begin{theorem}
\label{ll:paper:winwitt-step-rhcenter-cone}
For every natural number \(t\) and natural number \(k\) and natural number \(p\), if \(t + p + 1 < \operatorname{twopow}(k)\), then \(\operatorname{wittbitbool}(\operatorname{wittwin}(t + 1, \operatorname{twopow}(k)), t + 1 + \operatorname{twopow}(k) - p) = \operatorname{bxor}(\operatorname{wittbitbool}(\operatorname{wittphi}(\operatorname{wittwin}(t, \operatorname{twopow}(k))), t + 1 + \operatorname{twopow}(k) - p), \operatorname{bxor}(\operatorname{rhcenter}(t, p) \operatorname{and} \operatorname{wittbitbool}(\operatorname{wittwin}(t, \operatorname{twopow}(k)), t + \operatorname{twopow}(k) - (p + 1)), \operatorname{wittbitbool}(\operatorname{wittwin}(t, \operatorname{twopow}(k)), t + \operatorname{twopow}(k) - p) \operatorname{and} \operatorname{rhcenter}(t, p + 1)))\).
\end{theorem}
\begin{proof}
The goal follows from \(\operatorname{winwittsteprhcentercone}(t, k, p)\).
\end{proof}
\begin{theorem}
\label{ll:paper:winwitt-step-zero-gap-eq-phi}
For every natural number \(t\) and natural number \(k\) and natural number \(p\), if \(t < p\), then if \(t + p + 1 < \operatorname{twopow}(k)\), then \(\operatorname{wittbitbool}(\operatorname{wittwin}(t + 1, \operatorname{twopow}(k)), t + 1 + \operatorname{twopow}(k) - p) = \operatorname{wittbitbool}(\operatorname{wittphi}(\operatorname{wittwin}(t, \operatorname{twopow}(k))), t + 1 + \operatorname{twopow}(k) - p)\).
\end{theorem}
\begin{proof}
The goal follows from \(\operatorname{winwittstepzerogapeqphi}(t, k, p)\).
\end{proof}
\begin{theorem}
\label{ll:paper:winwitt-step-lhcenter-zero-gap-eq-phi}
For every natural number \(t\) and natural number \(k\) and natural number \(p\), if \(t < p\), then if \(t + p + 1 < \operatorname{twopow}(k)\), then \(\operatorname{wittbitbool}(\operatorname{wittwin}(t + 1, \operatorname{twopow}(k)), t + 1 + (p + 1) + \operatorname{twopow}(k)) = \operatorname{wittbitbool}(\operatorname{wittphi}(\operatorname{wittwin}(t, \operatorname{twopow}(k))), t + 1 + (p + 1) + \operatorname{twopow}(k))\).
\end{theorem}
\begin{proof}
The goal follows from \(\operatorname{winwittsteplhcenterzerogapeqphi}(t, k, p)\).
\end{proof}
\begin{theorem}
\label{ll:paper:winwittagree-step-rhcenter-cone}
For every natural number \(t\) and natural number \(k\) and natural number \(Q\) and natural number \(p0\) and natural number \(L\), if \(t + Q + p0 + L + 1 < \operatorname{twopow}(k)\), then if \(\operatorname{rhbandagree}(t, Q, p0, L + 1)\), then if \(\operatorname{winwittagree}(t, k, Q, t + \operatorname{twopow}(k) - (p0 + L + 1), L + 2)\), then \(\operatorname{winwittagree}(t + 1, k, Q, t + 1 + \operatorname{twopow}(k) - (p0 + L), L)\).
\end{theorem}
\begin{proof}
The goal follows from \(\operatorname{winwittagreesteprhcentercone}(t, k, Q, p0, L)\).
\end{proof}
\begin{theorem}
\label{ll:paper:winwittagree-inward-step-rhcenter-cone}
For every natural number \(t\) and natural number \(k\) and natural number \(Q\) and natural number \(p0\) and natural number \(L\), if \(t + Q + p0 + L + 1 < \operatorname{twopow}(k)\), then if \(\operatorname{rhbandagree}(t, Q, p0, L + 1)\), then if \(\operatorname{winwittagree}(t, k, Q, t + \operatorname{twopow}(k) - (p0 + L + 1), L + 1)\), then if \(\operatorname{winwittagree}(t + 1, k, Q, t + 1 + \operatorname{twopow}(k) - (p0 + L), L)\), then \(\operatorname{winwittagree}(t, k, Q, t + \operatorname{twopow}(k) - (p0 + L + 1), L + 2)\).
\end{theorem}
\begin{proof}
The goal follows from \(\operatorname{winwittagreeinwardsteprhcentercone}(t, k, Q, p0, L)\).
\end{proof}
\begin{theorem}
\label{ll:paper:winwitt-step-sector-one-center-slot}
For every natural number \(t\) and natural number \(k\), if \(t + 1 < \operatorname{twopow}(k)\), then \(\operatorname{wittbitbool}(\operatorname{wittwin}(t + 1, \operatorname{twopow}(k)), t + 1 + \operatorname{twopow}(k)) = \operatorname{bxor}(\operatorname{wittbitbool}(\operatorname{wittwin}(t, \operatorname{twopow}(k)), t + 1 + \operatorname{twopow}(k)), \operatorname{bxor}(\operatorname{wittbitbool}(\operatorname{wittwin}(t, \operatorname{twopow}(k)), t + \operatorname{twopow}(k)) \operatorname{or} \operatorname{wittbitbool}(\operatorname{wittwin}(t, \operatorname{twopow}(k)), t + \operatorname{twopow}(k) - 1), \operatorname{bxor}(\operatorname{center}(t) \operatorname{and} \operatorname{wittbitbool}(\operatorname{wittwin}(t, \operatorname{twopow}(k)), t + \operatorname{twopow}(k) - 1), \operatorname{wittbitbool}(\operatorname{wittwin}(t, \operatorname{twopow}(k)), t + \operatorname{twopow}(k)) \operatorname{and} \operatorname{rhcenter}(t, 1))))\).
\end{theorem}
\begin{proof}
The goal follows from \(\operatorname{winwittstepsectoronecenterslot}(t, k)\).
\end{proof}
\begin{theorem}
\label{ll:paper:periodic-winwitt-sector-one-one-defect-eq-late}
For every natural number \(T\) and natural number \(P\), if \(0 < P\), then if \(\operatorname{centerperiodic}(T, P)\), then for every natural number \(m\) and natural number \(t\) and natural number \(k\), if \(T \le t\), then if \(t + m \cdot P + 1 < \operatorname{twopow}(k)\), then \(\operatorname{bxor}(\operatorname{wittbitbool}(\operatorname{wittwin}(t + m \cdot P, \operatorname{twopow}(k)), t + m \cdot P + 1 + \operatorname{twopow}(k)), \operatorname{wittbitbool}(\operatorname{wittwin}(t, \operatorname{twopow}(k)), t + 1 + \operatorname{twopow}(k))) = \operatorname{bxor}(\operatorname{wittbitbool}(\operatorname{wittwin}(t, \operatorname{twopow}(k)), t + \operatorname{twopow}(k)) \operatorname{and} \operatorname{bxor}(\operatorname{rhcenter}(t + m \cdot P, 1), \operatorname{rhcenter}(t, 1)), \operatorname{not}(\operatorname{center}(t + \operatorname{twopow}(k))) \operatorname{and} \operatorname{bxor}(\operatorname{wittbitbool}(\operatorname{wittwin}(t + m \cdot P, \operatorname{twopow}(k)), t + m \cdot P + \operatorname{twopow}(k) - 1), \operatorname{wittbitbool}(\operatorname{wittwin}(t, \operatorname{twopow}(k)), t + \operatorname{twopow}(k) - 1)))\).
\end{theorem}
\begin{proof}
Assume \(hP\), \(hper\).
The goal follows from \(\operatorname{periodicwinwittsectoroneonedefecteqlate}(T, P, hP, hper)\).
\end{proof}
\begin{theorem}
\label{ll:paper:periodic-winwitt-sector-one-one-defect-eq-early}
For every natural number \(T\) and natural number \(P\), if \(0 < P\), then if \(\operatorname{centerperiodic}(T, P)\), then for every natural number \(m\) and natural number \(t\) and natural number \(k\), if \(T \le t\), then if \(t + m \cdot P + 1 < \operatorname{twopow}(k)\), then \(\operatorname{bxor}(\operatorname{wittbitbool}(\operatorname{wittwin}(t + m \cdot P, \operatorname{twopow}(k)), t + m \cdot P + 1 + \operatorname{twopow}(k)), \operatorname{wittbitbool}(\operatorname{wittwin}(t, \operatorname{twopow}(k)), t + 1 + \operatorname{twopow}(k))) = \operatorname{bxor}(\operatorname{wittbitbool}(\operatorname{wittwin}(t, \operatorname{twopow}(k)), t + \operatorname{twopow}(k)) \operatorname{and} \operatorname{bxor}(\operatorname{rhcenter}(t + \operatorname{twopow}(k) + m \cdot P, 1), \operatorname{rhcenter}(t + \operatorname{twopow}(k), 1)), \operatorname{not}(\operatorname{center}(t)) \operatorname{and} \operatorname{bxor}(\operatorname{wittbitbool}(\operatorname{wittwin}(t + m \cdot P, \operatorname{twopow}(k)), t + m \cdot P + \operatorname{twopow}(k) - 1), \operatorname{wittbitbool}(\operatorname{wittwin}(t, \operatorname{twopow}(k)), t + \operatorname{twopow}(k) - 1)))\).
\end{theorem}
\begin{proof}
Assume \(hP\), \(hper\).
The goal follows from \(\operatorname{periodicwinwittsectoroneonedefecteqearly}(T, P, hP, hper)\).
\end{proof}
\begin{theorem}
\label{ll:paper:periodic-winwitt-sector-one-one-defect-of-witt-zero-false}
For every natural number \(T\) and natural number \(P\), if \(0 < P\), then if \(\operatorname{centerperiodic}(T, P)\), then for every natural number \(m\) and natural number \(t\) and natural number \(k\), if \(T \le t\), then if \(t + m \cdot P + 1 < \operatorname{twopow}(k)\), then if \(\operatorname{wittbitbool}(\operatorname{wittwin}(t, \operatorname{twopow}(k)), t + \operatorname{twopow}(k)) = \operatorname{false}\), then \(\operatorname{bxor}(\operatorname{wittbitbool}(\operatorname{wittwin}(t + m \cdot P, \operatorname{twopow}(k)), t + m \cdot P + 1 + \operatorname{twopow}(k)), \operatorname{wittbitbool}(\operatorname{wittwin}(t, \operatorname{twopow}(k)), t + 1 + \operatorname{twopow}(k))) = \operatorname{not}(\operatorname{center}(t)) \operatorname{and} \operatorname{bxor}(\operatorname{wittbitbool}(\operatorname{wittwin}(t + m \cdot P, \operatorname{twopow}(k)), t + m \cdot P + \operatorname{twopow}(k) - 1), \operatorname{wittbitbool}(\operatorname{wittwin}(t, \operatorname{twopow}(k)), t + \operatorname{twopow}(k) - 1))\).
\end{theorem}
\begin{proof}
Assume \(hP\), \(hper\).
The goal follows from \(\operatorname{periodicwinwittsectoroneonedefectofwittzerofalse}(T, P, hP, hper)\).
\end{proof}
\begin{theorem}
\label{ll:paper:periodic-winwitt-sector-one-one-defect-of-witt-right-eq}
For every natural number \(T\) and natural number \(P\), if \(0 < P\), then if \(\operatorname{centerperiodic}(T, P)\), then for every natural number \(m\) and natural number \(t\) and natural number \(k\), if \(T \le t\), then if \(t + m \cdot P + 1 < \operatorname{twopow}(k)\), then if \(\operatorname{wittbitbool}(\operatorname{wittwin}(t + m \cdot P, \operatorname{twopow}(k)), t + m \cdot P + \operatorname{twopow}(k) - 1) = \operatorname{wittbitbool}(\operatorname{wittwin}(t, \operatorname{twopow}(k)), t + \operatorname{twopow}(k) - 1)\), then \(\operatorname{bxor}(\operatorname{wittbitbool}(\operatorname{wittwin}(t + m \cdot P, \operatorname{twopow}(k)), t + m \cdot P + 1 + \operatorname{twopow}(k)), \operatorname{wittbitbool}(\operatorname{wittwin}(t, \operatorname{twopow}(k)), t + 1 + \operatorname{twopow}(k))) = \operatorname{wittbitbool}(\operatorname{wittwin}(t, \operatorname{twopow}(k)), t + \operatorname{twopow}(k)) \operatorname{and} \operatorname{bxor}(\operatorname{rhcenter}(t + m \cdot P, 1), \operatorname{rhcenter}(t, 1))\).
\end{theorem}
\begin{proof}
Assume \(hP\), \(hper\).
The goal follows from \(\operatorname{periodicwinwittsectoroneonedefectofwittrighteq}(T, P, hP, hper)\).
\end{proof}
\begin{theorem}
\label{ll:paper:periodic-winwittagree-sector-one-two-of-right-one}
For every natural number \(T\) and natural number \(P\), if \(0 < P\), then if \(\operatorname{centerperiodic}(T, P)\), then for every natural number \(m\) and natural number \(t\) and natural number \(k\), if \(T \le t\), then if \(t + m \cdot P + 1 < \operatorname{twopow}(k)\), then if \(\operatorname{rhcenter}(t + m \cdot P, 1) = \operatorname{rhcenter}(t, 1)\), then if \(\operatorname{wittbitbool}(\operatorname{wittwin}(t + m \cdot P, \operatorname{twopow}(k)), t + m \cdot P + \operatorname{twopow}(k) - 1) = \operatorname{wittbitbool}(\operatorname{wittwin}(t, \operatorname{twopow}(k)), t + \operatorname{twopow}(k) - 1)\), then \(\operatorname{winwittagree}(t, k, m \cdot P, t + \operatorname{twopow}(k), 1)\).
\end{theorem}
\begin{proof}
Assume \(hP\), \(hper\).
The goal follows from \(\operatorname{periodicwinwittagreesectoronetwoofrightone}(T, P, hP, hper)\).
\end{proof}
\begin{theorem}
\label{ll:paper:periodic-winwittagree-sector-one-three-of-right-one-step}
For every natural number \(T\) and natural number \(P\), if \(0 < P\), then if \(\operatorname{centerperiodic}(T, P)\), then for every natural number \(m\) and natural number \(t\) and natural number \(k\), if \(T \le t\), then if \(t + m \cdot P + 2 < \operatorname{twopow}(k)\), then if \(\operatorname{rhcenter}(t + m \cdot P, 1) = \operatorname{rhcenter}(t, 1)\), then if \(\operatorname{rhcenter}(t + 1 + m \cdot P, 1) = \operatorname{rhcenter}(t + 1, 1)\), then if \(\operatorname{wittbitbool}(\operatorname{wittwin}(t + m \cdot P, \operatorname{twopow}(k)), t + m \cdot P + \operatorname{twopow}(k) - 1) = \operatorname{wittbitbool}(\operatorname{wittwin}(t, \operatorname{twopow}(k)), t + \operatorname{twopow}(k) - 1)\), then if \(\operatorname{wittbitbool}(\operatorname{wittwin}(t + 1 + m \cdot P, \operatorname{twopow}(k)), t + 1 + m \cdot P + \operatorname{twopow}(k) - 1) = \operatorname{wittbitbool}(\operatorname{wittwin}(t + 1, \operatorname{twopow}(k)), t + 1 + \operatorname{twopow}(k) - 1)\), then \(\operatorname{winwittagree}(t, k, m \cdot P, t + \operatorname{twopow}(k), 2)\).
\end{theorem}
\begin{proof}
Assume \(hP\), \(hper\).
The goal follows from \(\operatorname{periodicwinwittagreesectoronethreeofrightonestep}(T, P, hP, hper)\).
\end{proof}
\begin{theorem}
\label{ll:paper:winwitt-inward-step-frobdefect-testbit-succ}
For every natural number \(t\) and natural number \(k\) and natural number \(b\), if \(1 \le b\), then \(\operatorname{wittbitbool}(\operatorname{wittwin}(t, \operatorname{twopow}(k)), b + 1) = \operatorname{bxor}(\operatorname{wittbitbool}(\operatorname{wittwin}(t + 1, \operatorname{twopow}(k)), b + 1), \operatorname{bxor}(\operatorname{wittbitbool}(\operatorname{wittwin}(t, \operatorname{twopow}(k)), b) \operatorname{or} \operatorname{wittbitbool}(\operatorname{wittwin}(t, \operatorname{twopow}(k)), b - 1), \operatorname{bxor}(\operatorname{wittbitbool}(\operatorname{crosscarry}(\operatorname{wittscale}(k, \operatorname{wittrow}(t)), \operatorname{wittwin}(t, \operatorname{twopow}(k))), b), \operatorname{wittbitbool}(\operatorname{frobdefect}(k, \operatorname{wittrow}(t)), b))))\).
\end{theorem}
\begin{proof}
The goal follows from \(\operatorname{winwittinwardstepfrobdefecttestbitsucc}(t, k, b)\).
\end{proof}
\begin{theorem}
\label{ll:paper:winwitt-inward-step-gt-two-mul}
For every natural number \(t\) and natural number \(k\) and natural number \(b\), if \(2 \cdot t + 1 < b\), then if \(b < \operatorname{twopow}(k) + \operatorname{twopow}(k)\), then \(\operatorname{wittbitbool}(\operatorname{wittwin}(t, \operatorname{twopow}(k)), b + 1) = \operatorname{bxor}(\operatorname{wittbitbool}(\operatorname{wittwin}(t + 1, \operatorname{twopow}(k)), b + 1), \operatorname{bxor}(\operatorname{wittbitbool}(\operatorname{wittwin}(t, \operatorname{twopow}(k)), b) \operatorname{or} \operatorname{wittbitbool}(\operatorname{wittwin}(t, \operatorname{twopow}(k)), b - 1), \operatorname{wittbitbool}(\operatorname{crosscarry}(\operatorname{wittscale}(k, \operatorname{wittrow}(t)), \operatorname{wittwin}(t, \operatorname{twopow}(k))), b)))\).
\end{theorem}
\begin{proof}
The goal follows from \(\operatorname{winwittinwardstepgttwomul}(t, k, b)\).
\end{proof}
\begin{theorem}
\label{ll:paper:winwitt-sector-one-one-from-center-and-right}
For every natural number \(t\) and natural number \(k\), if \(t + 1 < \operatorname{twopow}(k)\), then \(\operatorname{wittbitbool}(\operatorname{wittwin}(t, \operatorname{twopow}(k)), t + 1 + \operatorname{twopow}(k)) = \operatorname{bxor}(\operatorname{wittbitbool}(\operatorname{wittwin}(t + 1, \operatorname{twopow}(k)), t + 1 + \operatorname{twopow}(k)), \operatorname{bxor}(\operatorname{wittbitbool}(\operatorname{wittwin}(t, \operatorname{twopow}(k)), t + \operatorname{twopow}(k)) \operatorname{or} \operatorname{wittbitbool}(\operatorname{wittwin}(t, \operatorname{twopow}(k)), t + \operatorname{twopow}(k) - 1), \operatorname{bxor}(\operatorname{center}(t) \operatorname{and} \operatorname{wittbitbool}(\operatorname{wittwin}(t, \operatorname{twopow}(k)), t + \operatorname{twopow}(k) - 1), \operatorname{wittbitbool}(\operatorname{wittwin}(t, \operatorname{twopow}(k)), t + \operatorname{twopow}(k)) \operatorname{and} \operatorname{rhcenter}(t, 1))))\).
\end{theorem}
\begin{proof}
The goal follows from \(\operatorname{winwittsectoroneonefromcenterandright}(t, k)\).
\end{proof}
\begin{theorem}
\label{ll:paper:winwitt-inward-step-sector-one-lhcenter}
For every natural number \(t\) and natural number \(k\) and natural number \(p\), if \(t + p + 1 < \operatorname{twopow}(k)\), then \(\operatorname{wittbitbool}(\operatorname{wittwin}(t, \operatorname{twopow}(k)), t + (p + 2) + \operatorname{twopow}(k)) = \operatorname{bxor}(\operatorname{wittbitbool}(\operatorname{wittwin}(t + 1, \operatorname{twopow}(k)), t + 1 + (p + 1) + \operatorname{twopow}(k)), \operatorname{bxor}(\operatorname{wittbitbool}(\operatorname{wittwin}(t, \operatorname{twopow}(k)), t + (p + 1) + \operatorname{twopow}(k)) \operatorname{or} \operatorname{wittbitbool}(\operatorname{wittwin}(t, \operatorname{twopow}(k)), t + p + \operatorname{twopow}(k)), \operatorname{bxor}(\operatorname{lhcenter}(t, p + 1) \operatorname{and} \operatorname{wittbitbool}(\operatorname{wittwin}(t, \operatorname{twopow}(k)), t + p + \operatorname{twopow}(k)), \operatorname{wittbitbool}(\operatorname{wittwin}(t, \operatorname{twopow}(k)), t + (p + 1) + \operatorname{twopow}(k)) \operatorname{and} \operatorname{lhcenter}(t, p))))\).
\end{theorem}
\begin{proof}
The goal follows from \(\operatorname{winwittinwardstepsectoronelhcenter}(t, k, p)\).
\end{proof}
\begin{theorem}
\label{ll:paper:winwitt-inward-step-rhcenter-cone}
For every natural number \(t\) and natural number \(k\) and natural number \(p\), if \(t + p + 2 < \operatorname{twopow}(k)\), then \(\operatorname{wittbitbool}(\operatorname{wittwin}(t, \operatorname{twopow}(k)), t + \operatorname{twopow}(k) - p) = \operatorname{bxor}(\operatorname{wittbitbool}(\operatorname{wittwin}(t + 1, \operatorname{twopow}(k)), t + 1 + \operatorname{twopow}(k) - (p + 1)), \operatorname{bxor}(\operatorname{wittbitbool}(\operatorname{wittwin}(t, \operatorname{twopow}(k)), t + \operatorname{twopow}(k) - (p + 1)) \operatorname{or} \operatorname{wittbitbool}(\operatorname{wittwin}(t, \operatorname{twopow}(k)), t + \operatorname{twopow}(k) - (p + 2)), \operatorname{bxor}(\operatorname{rhcenter}(t, p + 1) \operatorname{and} \operatorname{wittbitbool}(\operatorname{wittwin}(t, \operatorname{twopow}(k)), t + \operatorname{twopow}(k) - (p + 2)), \operatorname{wittbitbool}(\operatorname{wittwin}(t, \operatorname{twopow}(k)), t + \operatorname{twopow}(k) - (p + 1)) \operatorname{and} \operatorname{rhcenter}(t, p + 2))))\).
\end{theorem}
\begin{proof}
The goal follows from \(\operatorname{winwittinwardsteprhcentercone}(t, k, p)\).
\end{proof}
\begin{theorem}
\label{ll:paper:winwitt-inward-step-zero-gap}
For every natural number \(t\) and natural number \(k\) and natural number \(p\), if \(t \le p\), then if \(t + p + 2 < \operatorname{twopow}(k)\), then \(\operatorname{wittbitbool}(\operatorname{wittwin}(t, \operatorname{twopow}(k)), t + \operatorname{twopow}(k) - p) = \operatorname{bxor}(\operatorname{wittbitbool}(\operatorname{wittwin}(t + 1, \operatorname{twopow}(k)), t + 1 + \operatorname{twopow}(k) - (p + 1)), \operatorname{wittbitbool}(\operatorname{wittwin}(t, \operatorname{twopow}(k)), t + \operatorname{twopow}(k) - (p + 1)) \operatorname{or} \operatorname{wittbitbool}(\operatorname{wittwin}(t, \operatorname{twopow}(k)), t + \operatorname{twopow}(k) - (p + 2)))\).
\end{theorem}
\begin{proof}
The goal follows from \(\operatorname{winwittinwardstepzerogap}(t, k, p)\).
\end{proof}
\begin{theorem}
\label{ll:paper:winwittagree-step-zero-gap}
For every natural number \(t\) and natural number \(k\) and natural number \(Q\) and natural number \(p0\) and natural number \(L\), if \(t + Q < p0\), then if \(t + Q + p0 + L + 1 < \operatorname{twopow}(k)\), then if \(\operatorname{winwittagree}(t, k, Q, t + \operatorname{twopow}(k) - (p0 + L + 1), L + 2)\), then \(\operatorname{winwittagree}(t + 1, k, Q, t + 1 + \operatorname{twopow}(k) - (p0 + L), L)\).
\end{theorem}
\begin{proof}
The goal follows from \(\operatorname{winwittagreestepzerogap}(t, k, Q, p0, L)\).
\end{proof}
\begin{theorem}
\label{ll:paper:winwittagree-inward-step-zero-gap}
For every natural number \(t\) and natural number \(k\) and natural number \(Q\) and natural number \(p0\) and natural number \(L\), if \(t + Q < p0\), then if \(t + Q + p0 + L + 1 < \operatorname{twopow}(k)\), then if \(\operatorname{winwittagree}(t, k, Q, t + \operatorname{twopow}(k) - (p0 + L + 1), L + 1)\), then if \(\operatorname{winwittagree}(t + 1, k, Q, t + 1 + \operatorname{twopow}(k) - (p0 + L), L)\), then \(\operatorname{winwittagree}(t, k, Q, t + \operatorname{twopow}(k) - (p0 + L + 1), L + 2)\).
\end{theorem}
\begin{proof}
The goal follows from \(\operatorname{winwittagreeinwardstepzerogap}(t, k, Q, p0, L)\).
\end{proof}
\begin{theorem}
\label{ll:paper:winwittagree-step-lhcenter-zero-gap}
For every natural number \(t\) and natural number \(k\) and natural number \(Q\) and natural number \(p0\) and natural number \(L\), if \(t + Q < p0\), then if \(t + Q + p0 + L + 1 < \operatorname{twopow}(k)\), then if \(\operatorname{winwittagree}(t, k, Q, t + p0 + \operatorname{twopow}(k), L + 2)\), then \(\operatorname{winwittagree}(t + 1, k, Q, t + 1 + (p0 + 1) + \operatorname{twopow}(k), L)\).
\end{theorem}
\begin{proof}
The goal follows from \(\operatorname{winwittagreesteplhcenterzerogap}(t, k, Q, p0, L)\).
\end{proof}
\begin{theorem}
\label{ll:paper:winwittagree-outward-step-lhcenter-zero-gap}
For every natural number \(t\) and natural number \(k\) and natural number \(Q\) and natural number \(p0\) and natural number \(L\), if \(t + Q < p0\), then if \(t + Q + p0 + L + 1 < \operatorname{twopow}(k)\), then if \(\operatorname{winwittagree}(t, k, Q, t + p0 + \operatorname{twopow}(k), L + 1)\), then if \(\operatorname{winwittagree}(t + 1, k, Q, t + 1 + (p0 + 1) + \operatorname{twopow}(k), L)\), then \(\operatorname{winwittagree}(t, k, Q, t + p0 + \operatorname{twopow}(k), L + 2)\).
\end{theorem}
\begin{proof}
The goal follows from \(\operatorname{winwittagreeoutwardsteplhcenterzerogap}(t, k, Q, p0, L)\).
\end{proof}
\begin{theorem}
\label{ll:paper:winwittagree-steps-rhcenter-cone}
For every natural number \(t\) and natural number \(k\) and natural number \(Q\) and natural number \(p0\) and natural number \(L\) and natural number \(s\), if \(t + Q + p0 + L + 2 \cdot s < \operatorname{twopow}(k)\), then if \(\operatorname{rhbandagree}(t, Q, p0, L + 2 \cdot s)\), then if \(\operatorname{winwittagree}(t, k, Q, t + \operatorname{twopow}(k) - (p0 + L + 2 \cdot s), L + 2 \cdot s)\), then \(\operatorname{winwittagree}(t + s, k, Q, t + s + \operatorname{twopow}(k) - (p0 + s + L), L)\).
\end{theorem}
\begin{proof}
The goal follows from \(\operatorname{winwittagreestepsrhcentercone}(t, k, Q, p0, L, s)\).
\end{proof}
\begin{theorem}
\label{ll:paper:rhbandagree-late-steps-cone}
For every natural number \(t\) and natural number \(k\) and natural number \(Q\) and natural number \(p0\) and natural number \(L\) and natural number \(s\), if \(t + Q + p0 + L + 2 \cdot s < \operatorname{twopow}(k)\), then if \(\operatorname{rhbandagree}(t, Q, p0, L + 2 \cdot s)\), then if \(\operatorname{winwittagree}(t, k, Q, t + \operatorname{twopow}(k) - (p0 + L + 2 \cdot s), L + 2 \cdot s)\), then \(\operatorname{rhbandagree}(t + s + \operatorname{twopow}(k), Q, p0 + s, L)\).
\end{theorem}
\begin{proof}
The goal follows from \(\operatorname{rhbandagreelatestepscone}(t, k, Q, p0, L, s)\).
\end{proof}
\begin{theorem}
\label{ll:paper:winwittagree-steps-zero-gap}
For every natural number \(t\) and natural number \(k\) and natural number \(Q\) and natural number \(p0\) and natural number \(L\) and natural number \(s\), if \(t + Q < p0\), then if \(t + Q + p0 + L + 2 \cdot s < \operatorname{twopow}(k)\), then if \(\operatorname{winwittagree}(t, k, Q, t + \operatorname{twopow}(k) - (p0 + L + 2 \cdot s), L + 2 \cdot s)\), then \(\operatorname{winwittagree}(t + s, k, Q, t + s + \operatorname{twopow}(k) - (p0 + s + L), L)\).
\end{theorem}
\begin{proof}
The goal follows from \(\operatorname{winwittagreestepszerogap}(t, k, Q, p0, L, s)\).
\end{proof}
\begin{theorem}
\label{ll:paper:winwittagree-steps-sector-one-lhcenter}
For every natural number \(t\) and natural number \(k\) and natural number \(Q\) and natural number \(p0\) and natural number \(L\) and natural number \(s\), if \(t + Q + p0 + L + 2 \cdot s < \operatorname{twopow}(k)\), then if \(\operatorname{lhbandagree}(t, Q, p0, L + 2 \cdot s)\), then if \(\operatorname{winwittagree}(t, k, Q, t + p0 + \operatorname{twopow}(k), L + 2 \cdot s)\), then \(\operatorname{winwittagree}(t + s, k, Q, t + s + (p0 + s) + \operatorname{twopow}(k), L)\).
\end{theorem}
\begin{proof}
The goal follows from \(\operatorname{winwittagreestepssectoronelhcenter}(t, k, Q, p0, L, s)\).
\end{proof}
\begin{theorem}
\label{ll:paper:lhbandagree-late-steps-sector-one}
For every natural number \(t\) and natural number \(k\) and natural number \(Q\) and natural number \(p0\) and natural number \(L\) and natural number \(s\), if \(t + Q + p0 + L + 2 \cdot s < \operatorname{twopow}(k)\), then if \(\operatorname{lhbandagree}(t, Q, p0, L + 2 \cdot s)\), then if \(\operatorname{winwittagree}(t, k, Q, t + p0 + \operatorname{twopow}(k), L + 2 \cdot s)\), then \(\operatorname{lhbandagree}(t + s + \operatorname{twopow}(k), Q, p0 + s, L)\).
\end{theorem}
\begin{proof}
The goal follows from \(\operatorname{lhbandagreelatestepssectorone}(t, k, Q, p0, L, s)\).
\end{proof}
\begin{theorem}
\label{ll:paper:phiscale-row-testbit-inter-cone-gap-eq-false}
For every natural number \(t\) and natural number \(k\) and natural number \(b\), if \(2 \cdot t < b\), then if \(b < \operatorname{twopow}(k)\), then \(\operatorname{wittbitbool}(\operatorname{wittscale}(k, \operatorname{wittrow}(t)), b) = \operatorname{false}\).
\end{theorem}
\begin{proof}
The goal follows from \(\operatorname{phiscalerowtestbitinterconegapeqfalse}(t, k, b)\).
\end{proof}
\begin{theorem}
\label{ll:paper:winwitt-testbit-inter-cone-gap-eq-row-late}
For every natural number \(t\) and natural number \(k\) and natural number \(b\), if \(2 \cdot t < b\), then if \(b < \operatorname{twopow}(k)\), then \(\operatorname{wittbitbool}(\operatorname{wittwin}(t, \operatorname{twopow}(k)), b) = \operatorname{wittbitbool}(\operatorname{wittrow}(t + \operatorname{twopow}(k)), b)\).
\end{theorem}
\begin{proof}
The goal follows from \(\operatorname{winwitttestbitinterconegapeqrowlate}(t, k, b)\).
\end{proof}
\begin{theorem}
\label{ll:paper:crosscarry-row-scale-inter-cone-gap-eq-false}
For every natural number \(t\) and natural number \(k\) and natural number \(b\), if \(2 \cdot t + 1 < b\), then if \(b < \operatorname{twopow}(k)\), then \(\operatorname{wittbitbool}(\operatorname{crosscarry}(\operatorname{wittscale}(k, \operatorname{wittrow}(t)), \operatorname{wittwin}(t, \operatorname{twopow}(k))), b) = \operatorname{false}\).
\end{theorem}
\begin{proof}
The goal follows from \(\operatorname{crosscarryrowscaleinterconegapeqfalse}(t, k, b)\).
\end{proof}
\begin{theorem}
\label{ll:paper:winwitt-step-inter-cone-gap-eq-phi}
For every natural number \(t\) and natural number \(k\) and natural number \(b\), if \(2 \cdot t + 1 < b\), then if \(b < \operatorname{twopow}(k)\), then \(\operatorname{wittbitbool}(\operatorname{wittwin}(t + 1, \operatorname{twopow}(k)), b + 1) = \operatorname{wittbitbool}(\operatorname{wittphi}(\operatorname{wittwin}(t, \operatorname{twopow}(k))), b + 1)\).
\end{theorem}
\begin{proof}
The goal follows from \(\operatorname{winwittstepinterconegapeqphi}(t, k, b)\).
\end{proof}
\begin{theorem}
\label{ll:paper:winwitt-inward-step-inter-cone-gap}
For every natural number \(t\) and natural number \(k\) and natural number \(b\), if \(2 \cdot t + 1 < b\), then if \(b < \operatorname{twopow}(k)\), then \(\operatorname{wittbitbool}(\operatorname{wittwin}(t, \operatorname{twopow}(k)), b + 1) = \operatorname{bxor}(\operatorname{wittbitbool}(\operatorname{wittwin}(t + 1, \operatorname{twopow}(k)), b + 1), \operatorname{wittbitbool}(\operatorname{wittwin}(t, \operatorname{twopow}(k)), b) \operatorname{or} \operatorname{wittbitbool}(\operatorname{wittwin}(t, \operatorname{twopow}(k)), b - 1))\).
\end{theorem}
\begin{proof}
The goal follows from \(\operatorname{winwittinwardstepinterconegap}(t, k, b)\).
\end{proof}
\begin{theorem}
\label{ll:paper:phiscale-row-testbit-sector-one-eq-row}
For every natural number \(t\) and natural number \(k\) and natural number \(d\), if \(2 \cdot t < \operatorname{twopow}(k)\), then if \(d < \operatorname{twopow}(k)\), then \(\operatorname{wittbitbool}(\operatorname{wittscale}(k, \operatorname{wittrow}(t)), d + \operatorname{twopow}(k)) = \operatorname{wittbitbool}(\operatorname{wittrow}(t), d)\).
\end{theorem}
\begin{proof}
The goal follows from \(\operatorname{phiscalerowtestbitsectoroneeqrow}(t, k, d)\).
\end{proof}
\begin{theorem}
\label{ll:paper:phiscale-row-testbit-sector-two-eq-row}
For every natural number \(t\) and natural number \(k\) and natural number \(d\), if \(2 \cdot t < \operatorname{twopow}(k)\), then \(\operatorname{wittbitbool}(\operatorname{wittscale}(k, \operatorname{wittrow}(t)), d + \operatorname{twopow}(k + 1)) = \operatorname{wittbitbool}(\operatorname{wittrow}(t), d)\).
\end{theorem}
\begin{proof}
The goal follows from \(\operatorname{phiscalerowtestbitsectortwoeqrow}(t, k, d)\).
\end{proof}
\begin{theorem}
\label{ll:paper:winwitt-testbit-sector-one-eq-row-xor}
For every natural number \(t\) and natural number \(k\) and natural number \(d\), if \(2 \cdot t < \operatorname{twopow}(k)\), then if \(d < \operatorname{twopow}(k)\), then \(\operatorname{wittbitbool}(\operatorname{wittwin}(t, \operatorname{twopow}(k)), d + \operatorname{twopow}(k)) = \operatorname{bxor}(\operatorname{wittbitbool}(\operatorname{wittrow}(t + \operatorname{twopow}(k)), d + \operatorname{twopow}(k)), \operatorname{wittbitbool}(\operatorname{wittrow}(t), d))\).
\end{theorem}
\begin{proof}
The goal follows from \(\operatorname{winwitttestbitsectoroneeqrowxor}(t, k, d)\).
\end{proof}
\begin{theorem}
\label{ll:paper:winwitt-testbit-sector-two-eq-row-xor}
For every natural number \(t\) and natural number \(k\) and natural number \(d\), if \(2 \cdot t < \operatorname{twopow}(k)\), then \(\operatorname{wittbitbool}(\operatorname{wittwin}(t, \operatorname{twopow}(k)), d + \operatorname{twopow}(k + 1)) = \operatorname{bxor}(\operatorname{wittbitbool}(\operatorname{wittrow}(t + \operatorname{twopow}(k)), d + \operatorname{twopow}(k + 1)), \operatorname{wittbitbool}(\operatorname{wittrow}(t), d))\).
\end{theorem}
\begin{proof}
The goal follows from \(\operatorname{winwitttestbitsectortwoeqrowxor}(t, k, d)\).
\end{proof}
\begin{theorem}
\label{ll:paper:winwitt-sector-one-xor-sector-two-rhcenter-cone}
For every natural number \(t\) and natural number \(k\) and natural number \(p\), if \(t + p < \operatorname{twopow}(k)\), then \(\operatorname{bxor}(\operatorname{wittbitbool}(\operatorname{wittwin}(t, \operatorname{twopow}(k)), t + \operatorname{twopow}(k) - p), \operatorname{wittbitbool}(\operatorname{wittwin}(t, \operatorname{twopow}(k)), t + \operatorname{twopow}(k + 1) - p)) = \operatorname{bxor}(\operatorname{rhcenter}(t + \operatorname{twopow}(k), p), \operatorname{lhcenter}(t + \operatorname{twopow}(k), \operatorname{twopow}(k) - p))\).
\end{theorem}
\begin{proof}
The goal follows from \(\operatorname{winwittsectoronexorsectortworhcentercone}(t, k, p)\).
\end{proof}
\begin{theorem}
\label{ll:paper:winwitt-sector-one-xor-sector-two-lhcenter-cone}
For every natural number \(t\) and natural number \(k\) and natural number \(p\), if \(t + p < \operatorname{twopow}(k)\), then \(\operatorname{bxor}(\operatorname{wittbitbool}(\operatorname{wittwin}(t, \operatorname{twopow}(k)), t + p + \operatorname{twopow}(k)), \operatorname{wittbitbool}(\operatorname{wittwin}(t, \operatorname{twopow}(k)), t + p + \operatorname{twopow}(k + 1))) = \operatorname{bxor}(\operatorname{lhcenter}(t + \operatorname{twopow}(k), p), \operatorname{lhcenter}(t + \operatorname{twopow}(k), p + \operatorname{twopow}(k)))\).
\end{theorem}
\begin{proof}
The goal follows from \(\operatorname{winwittsectoronexorsectortwolhcentercone}(t, k, p)\).
\end{proof}
\begin{theorem}
\label{ll:paper:lhbandagree-late-shift-of-winwittagree-sectors}
For every natural number \(t\) and natural number \(k\) and natural number \(Q\) and natural number \(p0\) and natural number \(L\), if \(t + Q + p0 + L < \operatorname{twopow}(k)\), then if \(\operatorname{lhbandagree}(t + \operatorname{twopow}(k), Q, p0, L)\), then if \(\operatorname{winwittagree}(t, k, Q, t + p0 + \operatorname{twopow}(k), L)\), then if \(\operatorname{winwittagree}(t, k, Q, t + p0 + \operatorname{twopow}(k + 1), L)\), then \(\operatorname{lhbandagree}(t + \operatorname{twopow}(k), Q, p0 + \operatorname{twopow}(k), L)\).
\end{theorem}
\begin{proof}
The goal follows from \(\operatorname{lhbandagreelateshiftofwinwittagreesectors}(t, k, Q, p0, L)\).
\end{proof}
\begin{theorem}
\label{ll:paper:rhbandagree-late-of-lhbandagree-late-and-winwittagree-sectors}
For every natural number \(t\) and natural number \(k\) and natural number \(Q\) and natural number \(L\), if \(t + Q + L < \operatorname{twopow}(k)\), then if \(\operatorname{lhbandagree}(t + \operatorname{twopow}(k), Q, \operatorname{twopow}(k) - L, L)\), then if \(\operatorname{winwittagree}(t, k, Q, t + \operatorname{twopow}(k) - L, L)\), then if \(\operatorname{winwittagree}(t, k, Q, t + \operatorname{twopow}(k + 1) - L, L)\), then \(\operatorname{rhbandagree}(t + \operatorname{twopow}(k), Q, 0, L)\).
\end{theorem}
\begin{proof}
The goal follows from \(\operatorname{rhbandagreelateoflhbandagreelateandwinwittagreesectors}(t, k, Q, L)\).
\end{proof}
\begin{theorem}
\label{ll:paper:winwitt-testbit-twopow-eq-not-rhcenter-late}
For every natural number \(t\) and natural number \(k\), if \(2 \cdot t < \operatorname{twopow}(k)\), then \(\operatorname{wittbitbool}(\operatorname{wittwin}(t, \operatorname{twopow}(k)), \operatorname{twopow}(k)) = \operatorname{not}(\operatorname{rhcenter}(t + \operatorname{twopow}(k), t))\).
\end{theorem}
\begin{proof}
The goal follows from \(\operatorname{winwitttestbittwopoweqnotrhcenterlate}(t, k)\).
\end{proof}
\begin{theorem}
\label{ll:paper:rhcenter-late-self-from-zero-gap-step}
For every natural number \(t\) and natural number \(k\), if \(2 \cdot t + 2 < \operatorname{twopow}(k)\), then \(\operatorname{rhcenter}(t + \operatorname{twopow}(k), t) = \operatorname{bxor}(\operatorname{rhcenter}(t + 1 + \operatorname{twopow}(k), t + 1), \operatorname{rhcenter}(t + \operatorname{twopow}(k), t + 1) \operatorname{or} \operatorname{rhcenter}(t + \operatorname{twopow}(k), t + 2))\).
\end{theorem}
\begin{proof}
The goal follows from \(\operatorname{rhcenterlateselffromzerogapstep}(t, k)\).
\end{proof}
\begin{theorem}
\label{ll:paper:lhcenter-late-shift-eq-ldiag}
For every natural number \(t\) and natural number \(k\) and natural number \(p\), if \(p \le t\), then \(\operatorname{lhcenter}(t + \operatorname{twopow}(k), p + \operatorname{twopow}(k)) = \operatorname{ldiag}(t - p, t + \operatorname{twopow}(k))\).
\end{theorem}
\begin{proof}
The goal follows from \(\operatorname{lhcenterlateshifteqldiag}(t, k, p)\).
\end{proof}
\begin{theorem}
\label{ll:paper:lhcenter-late-cross-eq-ldiag}
For every natural number \(t\) and natural number \(k\) and natural number \(r\), if \(r \le \operatorname{twopow}(k)\), then \(\operatorname{lhcenter}(t + \operatorname{twopow}(k), \operatorname{twopow}(k) - r) = \operatorname{ldiag}(t + r, t + \operatorname{twopow}(k))\).
\end{theorem}
\begin{proof}
The goal follows from \(\operatorname{lhcenterlatecrosseqldiag}(t, k, r)\).
\end{proof}
\begin{theorem}
\label{ll:paper:lhcenter-late-twopow-eq-zcenter}
For every natural number \(t\) and natural number \(k\), if \(t \le k\), then if \(\operatorname{settledat}(t, k)\), then \(\operatorname{zcentereq}(t, \operatorname{lhcenter}(t + \operatorname{twopow}(k), \operatorname{twopow}(k)))\).
\end{theorem}
\begin{proof}
The goal follows from \(\operatorname{lhcenterlatetwopoweqzcenter}(t, k)\).
\end{proof}
\begin{theorem}
\label{ll:paper:winwitt-testbit-sector-two-center-slot-zcenter}
For every natural number \(t\) and natural number \(k\), if \(t \le k\), then if \(\operatorname{settledat}(t, k)\), then \(\operatorname{zcentereq}(t, \operatorname{bxor}(\operatorname{wittbitbool}(\operatorname{wittwin}(t, \operatorname{twopow}(k)), t + \operatorname{twopow}(k + 1)), \operatorname{center}(t)))\).
\end{theorem}
\begin{proof}
The goal follows from \(\operatorname{winwitttestbitsectortwocenterslotzcenter}(t, k)\).
\end{proof}
\begin{theorem}
\label{ll:paper:winwitt-sector-one-xor-sector-two-center-slot-zcenter}
For every natural number \(t\) and natural number \(k\), if \(t \le k\), then if \(\operatorname{settledat}(t, k)\), then \(\operatorname{zcentereq}(t, \operatorname{bxor}(\operatorname{center}(t + \operatorname{twopow}(k)), \operatorname{bxor}(\operatorname{wittbitbool}(\operatorname{wittwin}(t, \operatorname{twopow}(k)), t + \operatorname{twopow}(k)), \operatorname{wittbitbool}(\operatorname{wittwin}(t, \operatorname{twopow}(k)), t + \operatorname{twopow}(k + 1)))))\).
\end{theorem}
\begin{proof}
The goal follows from \(\operatorname{winwittsectoronexorsectortwocenterslotzcenter}(t, k)\).
\end{proof}
\begin{theorem}
\label{ll:paper:lhcenter-late-shift-of-zero-diag}
For every natural number \(t\) and natural number \(k\) and natural number \(p\), if \(p \le t\), then if \(t - p \le k\), then if \(\operatorname{zerodiagonal}(t - p)\), then if \(\operatorname{settledby}(t - p, t + \operatorname{twopow}(k))\), then \(\operatorname{lhcenter}(t + \operatorname{twopow}(k), p + \operatorname{twopow}(k)) = \operatorname{false}\).
\end{theorem}
\begin{proof}
The goal follows from \(\operatorname{lhcenterlateshiftofzerodiag}(t, k, p)\).
\end{proof}
\begin{theorem}
\label{ll:paper:lhcenter-late-cross-of-zero-diag}
For every natural number \(t\) and natural number \(k\) and natural number \(r\), if \(t + r \le k\), then if \(\operatorname{zerodiagonal}(t + r)\), then if \(\operatorname{settledby}(t + r, t + \operatorname{twopow}(k))\), then \(\operatorname{lhcenter}(t + \operatorname{twopow}(k), \operatorname{twopow}(k) - r) = \operatorname{false}\).
\end{theorem}
\begin{proof}
The goal follows from \(\operatorname{lhcenterlatecrossofzerodiag}(t, k, r)\).
\end{proof}
\begin{theorem}
\label{ll:paper:winwitt-testbit-outer-left-strip-of-zero-diag}
For every natural number \(t\) and natural number \(k\) and natural number \(p\), if \(p \le t\), then if \(t - p \le k\), then if \(\operatorname{zerodiagonal}(t - p)\), then if \(\operatorname{settledby}(t - p, t + \operatorname{twopow}(k))\), then \(\operatorname{wittbitbool}(\operatorname{wittwin}(t, \operatorname{twopow}(k)), t + p + \operatorname{twopow}(k + 1)) = \operatorname{lhcenter}(t, p)\).
\end{theorem}
\begin{proof}
The goal follows from \(\operatorname{winwitttestbitouterleftstripofzerodiag}(t, k, p)\).
\end{proof}
\begin{theorem}
\label{ll:paper:winwitt-testbit-scale-cross-of-zero-diag}
For every natural number \(t\) and natural number \(k\) and natural number \(r\), if \(t + r \le k\), then if \(\operatorname{zerodiagonal}(t + r)\), then if \(\operatorname{settledby}(t + r, t + \operatorname{twopow}(k))\), then \(\operatorname{wittbitbool}(\operatorname{wittwin}(t, \operatorname{twopow}(k)), t + \operatorname{twopow}(k + 1) - r) = \operatorname{rhcenter}(t, r)\).
\end{theorem}
\begin{proof}
The goal follows from \(\operatorname{winwitttestbitscalecrossofzerodiag}(t, k, r)\).
\end{proof}
\begin{theorem}
\label{ll:paper:winwitt-testbit-scale-cross-zero-gap-of-zero-diag}
For every natural number \(t\) and natural number \(k\) and natural number \(r\), if \(t < r\), then if \(t + r \le k\), then if \(\operatorname{zerodiagonal}(t + r)\), then if \(\operatorname{settledby}(t + r, t + \operatorname{twopow}(k))\), then \(\operatorname{wittbitbool}(\operatorname{wittwin}(t, \operatorname{twopow}(k)), t + \operatorname{twopow}(k + 1) - r) = \operatorname{false}\).
\end{theorem}
\begin{proof}
The goal follows from \(\operatorname{winwitttestbitscalecrosszerogapofzerodiag}(t, k, r)\).
\end{proof}
\begin{theorem}
\label{ll:paper:winwitt-sector-one-xor-sector-two-lhcenter-of-zero-diag}
For every natural number \(t\) and natural number \(k\) and natural number \(p\), if \(p \le t\), then if \(t + p < \operatorname{twopow}(k)\), then if \(t - p \le k\), then if \(\operatorname{zerodiagonal}(t - p)\), then if \(\operatorname{settledby}(t - p, t + \operatorname{twopow}(k))\), then \(\operatorname{bxor}(\operatorname{wittbitbool}(\operatorname{wittwin}(t, \operatorname{twopow}(k)), t + p + \operatorname{twopow}(k)), \operatorname{wittbitbool}(\operatorname{wittwin}(t, \operatorname{twopow}(k)), t + p + \operatorname{twopow}(k + 1))) = \operatorname{lhcenter}(t + \operatorname{twopow}(k), p)\).
\end{theorem}
\begin{proof}
The goal follows from \(\operatorname{winwittsectoronexorsectortwolhcenterofzerodiag}(t, k, p)\).
\end{proof}
\begin{theorem}
\label{ll:paper:winwitt-sector-one-xor-sector-two-rhcenter-of-zero-diag}
For every natural number \(t\) and natural number \(k\) and natural number \(r\), if \(t + r \le k\), then if \(\operatorname{zerodiagonal}(t + r)\), then if \(\operatorname{settledby}(t + r, t + \operatorname{twopow}(k))\), then \(\operatorname{bxor}(\operatorname{wittbitbool}(\operatorname{wittwin}(t, \operatorname{twopow}(k)), t + \operatorname{twopow}(k) - r), \operatorname{wittbitbool}(\operatorname{wittwin}(t, \operatorname{twopow}(k)), t + \operatorname{twopow}(k + 1) - r)) = \operatorname{rhcenter}(t + \operatorname{twopow}(k), r)\).
\end{theorem}
\begin{proof}
The goal follows from \(\operatorname{winwittsectoronexorsectortworhcenterofzerodiag}(t, k, r)\).
\end{proof}
\begin{theorem}
\label{ll:paper:lhcenter-late-twopow-of-zero-diag}
For every natural number \(z\) and natural number \(k\), if \(z \le k\), then if \(\operatorname{zerodiagonal}(z)\), then if \(\operatorname{settledat}(z, k)\), then \(\operatorname{lhcenter}(z + \operatorname{twopow}(k), \operatorname{twopow}(k)) = \operatorname{false}\).
\end{theorem}
\begin{proof}
The goal follows from \(\operatorname{lhcenterlatetwopowofzerodiag}(z, k)\).
\end{proof}
\begin{theorem}
\label{ll:paper:winwitt-testbit-sector-two-center-slot-of-zero-diag}
For every natural number \(z\) and natural number \(k\), if \(z \le k\), then if \(\operatorname{zerodiagonal}(z)\), then if \(\operatorname{settledat}(z, k)\), then \(\operatorname{wittbitbool}(\operatorname{wittwin}(z, \operatorname{twopow}(k)), z + \operatorname{twopow}(k + 1)) = \operatorname{center}(z)\).
\end{theorem}
\begin{proof}
The goal follows from \(\operatorname{winwitttestbitsectortwocenterslotofzerodiag}(z, k)\).
\end{proof}
\begin{theorem}
\label{ll:paper:winwitt-sector-one-xor-sector-two-center-slot-of-zero-diag}
For every natural number \(z\) and natural number \(k\), if \(z \le k\), then if \(\operatorname{zerodiagonal}(z)\), then if \(\operatorname{settledat}(z, k)\), then \(\operatorname{bxor}(\operatorname{wittbitbool}(\operatorname{wittwin}(z, \operatorname{twopow}(k)), z + \operatorname{twopow}(k)), \operatorname{wittbitbool}(\operatorname{wittwin}(z, \operatorname{twopow}(k)), z + \operatorname{twopow}(k + 1))) = \operatorname{center}(z + \operatorname{twopow}(k))\).
\end{theorem}
\begin{proof}
The goal follows from \(\operatorname{winwittsectoronexorsectortwocenterslotofzerodiag}(z, k)\).
\end{proof}
\begin{theorem}
\label{ll:paper:lhcenter-late-shift-after-zero-diag}
For every natural number \(t\) and natural number \(k\) and natural number \(p\), if \(p \le t\), then if \(t - p + 2 \le k\), then if \(\operatorname{zerodiagonal}(t - p)\), then if \(\operatorname{settledby}(t - p + 2, t + 2 + \operatorname{twopow}(k))\), then \(\operatorname{lhcenter}(t + 2 + \operatorname{twopow}(k), p + \operatorname{twopow}(k)) = \operatorname{true}\).
\end{theorem}
\begin{proof}
The goal follows from \(\operatorname{lhcenterlateshiftafterzerodiag}(t, k, p)\).
\end{proof}
\begin{theorem}
\label{ll:paper:lhcenter-late-cross-after-zero-diag}
For every natural number \(t\) and natural number \(k\) and natural number \(r\), if \(t + r + 2 \le k\), then if \(\operatorname{zerodiagonal}(t + r)\), then if \(\operatorname{settledby}(t + r + 2, t + \operatorname{twopow}(k))\), then \(\operatorname{lhcenter}(t + \operatorname{twopow}(k), \operatorname{twopow}(k) - (r + 2)) = \operatorname{true}\).
\end{theorem}
\begin{proof}
The goal follows from \(\operatorname{lhcenterlatecrossafterzerodiag}(t, k, r)\).
\end{proof}
\begin{theorem}
\label{ll:paper:winwitt-testbit-outer-left-strip-after-zero-diag}
For every natural number \(t\) and natural number \(k\) and natural number \(p\), if \(p \le t\), then if \(t - p + 2 \le k\), then if \(\operatorname{zerodiagonal}(t - p)\), then if \(\operatorname{settledby}(t - p + 2, t + 2 + \operatorname{twopow}(k))\), then \(\operatorname{wittbitbool}(\operatorname{wittwin}(t + 2, \operatorname{twopow}(k)), t + 2 + p + \operatorname{twopow}(k + 1)) = \operatorname{not}(\operatorname{lhcenter}(t + 2, p))\).
\end{theorem}
\begin{proof}
The goal follows from \(\operatorname{winwitttestbitouterleftstripafterzerodiag}(t, k, p)\).
\end{proof}
\begin{theorem}
\label{ll:paper:winwitt-testbit-scale-cross-after-zero-diag}
For every natural number \(t\) and natural number \(k\) and natural number \(r\), if \(t + r + 2 \le k\), then if \(\operatorname{zerodiagonal}(t + r)\), then if \(\operatorname{settledby}(t + r + 2, t + \operatorname{twopow}(k))\), then \(\operatorname{wittbitbool}(\operatorname{wittwin}(t, \operatorname{twopow}(k)), t + \operatorname{twopow}(k + 1) - (r + 2)) = \operatorname{not}(\operatorname{rhcenter}(t, r + 2))\).
\end{theorem}
\begin{proof}
The goal follows from \(\operatorname{winwitttestbitscalecrossafterzerodiag}(t, k, r)\).
\end{proof}
\begin{theorem}
\label{ll:paper:winwitt-testbit-scale-cross-zero-gap-after-zero-diag}
For every natural number \(t\) and natural number \(k\) and natural number \(r\), if \(t < r + 2\), then if \(t + r + 2 \le k\), then if \(\operatorname{zerodiagonal}(t + r)\), then if \(\operatorname{settledby}(t + r + 2, t + \operatorname{twopow}(k))\), then \(\operatorname{wittbitbool}(\operatorname{wittwin}(t, \operatorname{twopow}(k)), t + \operatorname{twopow}(k + 1) - (r + 2)) = \operatorname{true}\).
\end{theorem}
\begin{proof}
The goal follows from \(\operatorname{winwitttestbitscalecrosszerogapafterzerodiag}(t, k, r)\).
\end{proof}
\begin{theorem}
\label{ll:paper:winwitt-sector-one-xor-sector-two-lhcenter-after-zero-diag}
For every natural number \(t\) and natural number \(k\) and natural number \(p\), if \(p \le t\), then if \(t + 2 + p < \operatorname{twopow}(k)\), then if \(t - p + 2 \le k\), then if \(\operatorname{zerodiagonal}(t - p)\), then if \(\operatorname{settledby}(t - p + 2, t + 2 + \operatorname{twopow}(k))\), then \(\operatorname{bxor}(\operatorname{wittbitbool}(\operatorname{wittwin}(t + 2, \operatorname{twopow}(k)), t + 2 + p + \operatorname{twopow}(k)), \operatorname{wittbitbool}(\operatorname{wittwin}(t + 2, \operatorname{twopow}(k)), t + 2 + p + \operatorname{twopow}(k + 1))) = \operatorname{not}(\operatorname{lhcenter}(t + 2 + \operatorname{twopow}(k), p))\).
\end{theorem}
\begin{proof}
The goal follows from \(\operatorname{winwittsectoronexorsectortwolhcenterafterzerodiag}(t, k, p)\).
\end{proof}
\begin{theorem}
\label{ll:paper:winwitt-sector-one-xor-sector-two-rhcenter-after-zero-diag}
For every natural number \(t\) and natural number \(k\) and natural number \(r\), if \(t + r + 2 \le k\), then if \(\operatorname{zerodiagonal}(t + r)\), then if \(\operatorname{settledby}(t + r + 2, t + \operatorname{twopow}(k))\), then \(\operatorname{bxor}(\operatorname{wittbitbool}(\operatorname{wittwin}(t, \operatorname{twopow}(k)), t + \operatorname{twopow}(k) - (r + 2)), \operatorname{wittbitbool}(\operatorname{wittwin}(t, \operatorname{twopow}(k)), t + \operatorname{twopow}(k + 1) - (r + 2))) = \operatorname{not}(\operatorname{rhcenter}(t + \operatorname{twopow}(k), r + 2))\).
\end{theorem}
\begin{proof}
The goal follows from \(\operatorname{winwittsectoronexorsectortworhcenterafterzerodiag}(t, k, r)\).
\end{proof}
\begin{theorem}
\label{ll:paper:lhcenter-late-twopow-after-zero-diag}
For every natural number \(z\) and natural number \(k\), if \(z + 2 \le k\), then if \(\operatorname{zerodiagonal}(z)\), then if \(\operatorname{settledat}(z + 2, k)\), then \(\operatorname{lhcenter}(z + 2 + \operatorname{twopow}(k), \operatorname{twopow}(k)) = \operatorname{true}\).
\end{theorem}
\begin{proof}
The goal follows from \(\operatorname{lhcenterlatetwopowafterzerodiag}(z, k)\).
\end{proof}
\begin{theorem}
\label{ll:paper:winwitt-testbit-sector-two-center-slot-after-zero-diag}
For every natural number \(z\) and natural number \(k\), if \(z + 2 \le k\), then if \(\operatorname{zerodiagonal}(z)\), then if \(\operatorname{settledat}(z + 2, k)\), then \(\operatorname{wittbitbool}(\operatorname{wittwin}(z + 2, \operatorname{twopow}(k)), z + 2 + \operatorname{twopow}(k + 1)) = \operatorname{not}(\operatorname{center}(z + 2))\).
\end{theorem}
\begin{proof}
The goal follows from \(\operatorname{winwitttestbitsectortwocenterslotafterzerodiag}(z, k)\).
\end{proof}
\begin{theorem}
\label{ll:paper:winwitt-sector-one-xor-sector-two-center-slot-after-zero-diag}
For every natural number \(z\) and natural number \(k\), if \(z + 2 \le k\), then if \(\operatorname{zerodiagonal}(z)\), then if \(\operatorname{settledat}(z + 2, k)\), then \(\operatorname{bxor}(\operatorname{wittbitbool}(\operatorname{wittwin}(z + 2, \operatorname{twopow}(k)), z + 2 + \operatorname{twopow}(k)), \operatorname{wittbitbool}(\operatorname{wittwin}(z + 2, \operatorname{twopow}(k)), z + 2 + \operatorname{twopow}(k + 1))) = \operatorname{not}(\operatorname{center}(z + 2 + \operatorname{twopow}(k)))\).
\end{theorem}
\begin{proof}
The goal follows from \(\operatorname{winwittsectoronexorsectortwocenterslotafterzerodiag}(z, k)\).
\end{proof}
\begin{theorem}
\label{ll:paper:lhcenter-late-cross-reset-after-zero-diag}
For every natural number \(t\) and natural number \(k\) and natural number \(r\), if \(t + r + 2 \le k\), then if \(\operatorname{zerodiagonal}(t + r)\), then if \(\operatorname{settledby}(t + r + 2, t + \operatorname{twopow}(k))\), then \(\operatorname{lhcenter}(t + 1 + \operatorname{twopow}(k), \operatorname{twopow}(k) - (r + 2)) = \operatorname{not}(\operatorname{lhcenter}(t + \operatorname{twopow}(k), \operatorname{twopow}(k) - (r + 1)))\).
\end{theorem}
\begin{proof}
The goal follows from \(\operatorname{lhcenterlatecrossresetafterzerodiag}(t, k, r)\).
\end{proof}
\begin{theorem}
\label{ll:paper:lhcenter-late-shift-reset-after-zero-diag}
For every natural number \(t\) and natural number \(k\) and natural number \(p\), if \(p \le t\), then if \(t - p + 2 \le k\), then if \(\operatorname{zerodiagonal}(t - p)\), then if \(\operatorname{settledby}(t - p + 2, t + 2 + \operatorname{twopow}(k))\), then \(\operatorname{lhcenter}(t + 3 + \operatorname{twopow}(k), p + 1 + \operatorname{twopow}(k)) = \operatorname{not}(\operatorname{lhcenter}(t + 2 + \operatorname{twopow}(k), p + 2 + \operatorname{twopow}(k)))\).
\end{theorem}
\begin{proof}
The goal follows from \(\operatorname{lhcenterlateshiftresetafterzerodiag}(t, k, p)\).
\end{proof}
\begin{theorem}
\label{ll:paper:winwitt-sector-two-step-complement-after-zero-diag}
For every natural number \(t\) and natural number \(k\) and natural number \(r\), if \(t + r + 2 \le k\), then if \(\operatorname{zerodiagonal}(t + r)\), then if \(\operatorname{settledby}(t + r + 2, t + \operatorname{twopow}(k))\), then \(\operatorname{bxor}(\operatorname{wittbitbool}(\operatorname{wittwin}(t, \operatorname{twopow}(k)), t + \operatorname{twopow}(k + 1) - (r + 1)), \operatorname{wittbitbool}(\operatorname{wittwin}(t + 1, \operatorname{twopow}(k)), t + \operatorname{twopow}(k + 1) - (r + 1))) = \operatorname{not}(\operatorname{bxor}(\operatorname{rhcenter}(t, r + 1), \operatorname{rhcenter}(t + 1, r + 2)))\).
\end{theorem}
\begin{proof}
The goal follows from \(\operatorname{winwittsectortwostepcomplementafterzerodiag}(t, k, r)\).
\end{proof}
\begin{theorem}
\label{ll:paper:winwitt-sector-two-step-flip-zero-gap-after-zero-diag}
For every natural number \(t\) and natural number \(k\) and natural number \(r\), if \(t \le r\), then if \(t + r + 2 \le k\), then if \(\operatorname{zerodiagonal}(t + r)\), then if \(\operatorname{settledby}(t + r + 2, t + \operatorname{twopow}(k))\), then \(\operatorname{wittbitbool}(\operatorname{wittwin}(t + 1, \operatorname{twopow}(k)), t + \operatorname{twopow}(k + 1) - (r + 1)) = \operatorname{not}(\operatorname{wittbitbool}(\operatorname{wittwin}(t, \operatorname{twopow}(k)), t + \operatorname{twopow}(k + 1) - (r + 1)))\).
\end{theorem}
\begin{proof}
The goal follows from \(\operatorname{winwittsectortwostepflipzerogapafterzerodiag}(t, k, r)\).
\end{proof}
\begin{theorem}
\label{ll:paper:winwitt-sector-one-xor-sector-two-step-complement-after-zero-diag}
For every natural number \(t\) and natural number \(k\) and natural number \(r\), if \(t + r + 2 \le k\), then if \(\operatorname{zerodiagonal}(t + r)\), then if \(\operatorname{settledby}(t + r + 2, t + \operatorname{twopow}(k))\), then \(\operatorname{bxor}(\operatorname{bxor}(\operatorname{wittbitbool}(\operatorname{wittwin}(t, \operatorname{twopow}(k)), t + \operatorname{twopow}(k) - (r + 1)), \operatorname{wittbitbool}(\operatorname{wittwin}(t, \operatorname{twopow}(k)), t + \operatorname{twopow}(k + 1) - (r + 1))), \operatorname{bxor}(\operatorname{wittbitbool}(\operatorname{wittwin}(t + 1, \operatorname{twopow}(k)), t + \operatorname{twopow}(k) - (r + 1)), \operatorname{wittbitbool}(\operatorname{wittwin}(t + 1, \operatorname{twopow}(k)), t + \operatorname{twopow}(k + 1) - (r + 1)))) = \operatorname{not}(\operatorname{bxor}(\operatorname{rhcenter}(t + \operatorname{twopow}(k), r + 1), \operatorname{rhcenter}(t + 1 + \operatorname{twopow}(k), r + 2)))\).
\end{theorem}
\begin{proof}
The goal follows from \(\operatorname{winwittsectoronexorsectortwostepcomplementafterzerodiag}(t, k, r)\).
\end{proof}
\begin{theorem}
\label{ll:paper:lhcenter-late-cross-step-one-of-zero-diag}
For every natural number \(t\) and natural number \(k\) and natural number \(r\), if \(t + r + 2 \le k\), then if \(\operatorname{zerodiagonal}(t + r)\), then if \(\operatorname{settledby}(t + r + 2, t + 1 + \operatorname{twopow}(k))\), then \(\operatorname{lhcenter}(t + 1 + \operatorname{twopow}(k), \operatorname{twopow}(k) - (r + 1)) = \operatorname{true}\).
\end{theorem}
\begin{proof}
The goal follows from \(\operatorname{lhcenterlatecrosssteponeofzerodiag}(t, k, r)\).
\end{proof}
\begin{theorem}
\label{ll:paper:winwitt-testbit-scale-cross-step-one-of-zero-diag}
For every natural number \(t\) and natural number \(k\) and natural number \(r\), if \(t + r + 2 \le k\), then if \(\operatorname{zerodiagonal}(t + r)\), then if \(\operatorname{settledby}(t + r + 2, t + 1 + \operatorname{twopow}(k))\), then \(\operatorname{wittbitbool}(\operatorname{wittwin}(t + 1, \operatorname{twopow}(k)), t + \operatorname{twopow}(k + 1) - r) = \operatorname{not}(\operatorname{rhcenter}(t + 1, r + 1))\).
\end{theorem}
\begin{proof}
The goal follows from \(\operatorname{winwitttestbitscalecrosssteponeofzerodiag}(t, k, r)\).
\end{proof}
\begin{theorem}
\label{ll:paper:winwitt-testbit-scale-cross-zero-gap-step-one-of-zero-diag}
For every natural number \(t\) and natural number \(k\) and natural number \(r\), if \(t < r\), then if \(t + r + 2 \le k\), then if \(\operatorname{zerodiagonal}(t + r)\), then if \(\operatorname{settledby}(t + r + 2, t + 1 + \operatorname{twopow}(k))\), then \(\operatorname{wittbitbool}(\operatorname{wittwin}(t + 1, \operatorname{twopow}(k)), t + \operatorname{twopow}(k + 1) - r) = \operatorname{true}\).
\end{theorem}
\begin{proof}
The goal follows from \(\operatorname{winwitttestbitscalecrosszerogapsteponeofzerodiag}(t, k, r)\).
\end{proof}
\begin{theorem}
\label{ll:paper:winwitt-sector-two-zero-to-black-step-of-zero-diag}
For every natural number \(t\) and natural number \(k\) and natural number \(r\), if \(t + r + 2 \le k\), then if \(\operatorname{zerodiagonal}(t + r)\), then if \(\operatorname{settledby}(t + r, t + \operatorname{twopow}(k))\), then if \(\operatorname{settledby}(t + r + 2, t + 1 + \operatorname{twopow}(k))\), then \(\operatorname{bxor}(\operatorname{wittbitbool}(\operatorname{wittwin}(t, \operatorname{twopow}(k)), t + \operatorname{twopow}(k + 1) - r), \operatorname{wittbitbool}(\operatorname{wittwin}(t + 1, \operatorname{twopow}(k)), t + \operatorname{twopow}(k + 1) - r)) = \operatorname{not}(\operatorname{bxor}(\operatorname{rhcenter}(t, r), \operatorname{rhcenter}(t + 1, r + 1)))\).
\end{theorem}
\begin{proof}
The goal follows from \(\operatorname{winwittsectortwozerotoblackstepofzerodiag}(t, k, r)\).
\end{proof}
\begin{theorem}
\label{ll:paper:winwitt-sector-one-xor-sector-two-rhcenter-step-one-of-zero-diag}
For every natural number \(t\) and natural number \(k\) and natural number \(r\), if \(t + r + 2 \le k\), then if \(\operatorname{zerodiagonal}(t + r)\), then if \(\operatorname{settledby}(t + r + 2, t + 1 + \operatorname{twopow}(k))\), then \(\operatorname{bxor}(\operatorname{wittbitbool}(\operatorname{wittwin}(t + 1, \operatorname{twopow}(k)), t + \operatorname{twopow}(k) - r), \operatorname{wittbitbool}(\operatorname{wittwin}(t + 1, \operatorname{twopow}(k)), t + \operatorname{twopow}(k + 1) - r)) = \operatorname{not}(\operatorname{rhcenter}(t + 1 + \operatorname{twopow}(k), r + 1))\).
\end{theorem}
\begin{proof}
The goal follows from \(\operatorname{winwittsectoronexorsectortworhcentersteponeofzerodiag}(t, k, r)\).
\end{proof}
\begin{theorem}
\label{ll:paper:winwitt-sector-one-xor-sector-two-zero-to-black-step-of-zero-diag}
For every natural number \(t\) and natural number \(k\) and natural number \(r\), if \(t + r + 2 \le k\), then if \(\operatorname{zerodiagonal}(t + r)\), then if \(\operatorname{settledby}(t + r, t + \operatorname{twopow}(k))\), then if \(\operatorname{settledby}(t + r + 2, t + 1 + \operatorname{twopow}(k))\), then \(\operatorname{bxor}(\operatorname{bxor}(\operatorname{wittbitbool}(\operatorname{wittwin}(t, \operatorname{twopow}(k)), t + \operatorname{twopow}(k) - r), \operatorname{wittbitbool}(\operatorname{wittwin}(t, \operatorname{twopow}(k)), t + \operatorname{twopow}(k + 1) - r)), \operatorname{bxor}(\operatorname{wittbitbool}(\operatorname{wittwin}(t + 1, \operatorname{twopow}(k)), t + \operatorname{twopow}(k) - r), \operatorname{wittbitbool}(\operatorname{wittwin}(t + 1, \operatorname{twopow}(k)), t + \operatorname{twopow}(k + 1) - r))) = \operatorname{not}(\operatorname{bxor}(\operatorname{rhcenter}(t + \operatorname{twopow}(k), r), \operatorname{rhcenter}(t + 1 + \operatorname{twopow}(k), r + 1)))\).
\end{theorem}
\begin{proof}
The goal follows from \(\operatorname{winwittsectoronexorsectortwozerotoblackstepofzerodiag}(t, k, r)\).
\end{proof}
\begin{theorem}
\label{ll:paper:phiscale-row-testbit-scale-cross-cone}
For every natural number \(t\) and natural number \(k\) and natural number \(r\), if \(t + r < \operatorname{twopow}(k)\), then \(\operatorname{wittbitbool}(\operatorname{wittscale}(k, \operatorname{wittrow}(t)), t + \operatorname{twopow}(k + 1) - r) = \operatorname{rhcenter}(t, r)\).
\end{theorem}
\begin{proof}
The goal follows from \(\operatorname{phiscalerowtestbitscalecrosscone}(t, k, r)\).
\end{proof}
\begin{theorem}
\label{ll:paper:phiscale-row-testbit-scale-cross-zero-gap-eq-false}
For every natural number \(t\) and natural number \(k\) and natural number \(r\), if \(t < r\), then if \(t + r < \operatorname{twopow}(k)\), then \(\operatorname{wittbitbool}(\operatorname{wittscale}(k, \operatorname{wittrow}(t)), t + \operatorname{twopow}(k + 1) - r) = \operatorname{false}\).
\end{theorem}
\begin{proof}
The goal follows from \(\operatorname{phiscalerowtestbitscalecrosszerogapeqfalse}(t, k, r)\).
\end{proof}
\begin{theorem}
\label{ll:paper:crosscarry-row-scale-cross-zero-gap-eq-false}
For every natural number \(t\) and natural number \(k\) and natural number \(r\), if \(t \le r\), then if \(t + r + 2 < \operatorname{twopow}(k)\), then \(\operatorname{wittbitbool}(\operatorname{crosscarry}(\operatorname{wittscale}(k, \operatorname{wittrow}(t)), \operatorname{wittwin}(t, \operatorname{twopow}(k))), t + \operatorname{twopow}(k + 1) - (r + 1)) = \operatorname{false}\).
\end{theorem}
\begin{proof}
The goal follows from \(\operatorname{crosscarryrowscalecrosszerogapeqfalse}(t, k, r)\).
\end{proof}
\begin{theorem}
\label{ll:paper:winwitt-step-scale-cross-zero-gap-eq-phi}
For every natural number \(t\) and natural number \(k\) and natural number \(r\), if \(t \le r\), then if \(t + r + 2 < \operatorname{twopow}(k)\), then \(\operatorname{wittbitbool}(\operatorname{wittwin}(t + 1, \operatorname{twopow}(k)), t + \operatorname{twopow}(k + 1) - r) = \operatorname{wittbitbool}(\operatorname{wittphi}(\operatorname{wittwin}(t, \operatorname{twopow}(k))), t + \operatorname{twopow}(k + 1) - r)\).
\end{theorem}
\begin{proof}
The goal follows from \(\operatorname{winwittstepscalecrosszerogapeqphi}(t, k, r)\).
\end{proof}
\begin{theorem}
\label{ll:paper:winwitt-inward-step-scale-cross-zero-gap}
For every natural number \(t\) and natural number \(k\) and natural number \(r\), if \(t \le r\), then if \(t + r + 2 < \operatorname{twopow}(k)\), then \(\operatorname{wittbitbool}(\operatorname{wittwin}(t, \operatorname{twopow}(k)), t + \operatorname{twopow}(k + 1) - r) = \operatorname{bxor}(\operatorname{wittbitbool}(\operatorname{wittwin}(t + 1, \operatorname{twopow}(k)), t + \operatorname{twopow}(k + 1) - r), \operatorname{wittbitbool}(\operatorname{wittwin}(t, \operatorname{twopow}(k)), t + \operatorname{twopow}(k + 1) - (r + 1)) \operatorname{or} \operatorname{wittbitbool}(\operatorname{wittwin}(t, \operatorname{twopow}(k)), t + \operatorname{twopow}(k + 1) - (r + 2)))\).
\end{theorem}
\begin{proof}
The goal follows from \(\operatorname{winwittinwardstepscalecrosszerogap}(t, k, r)\).
\end{proof}
\begin{theorem}
\label{ll:paper:lhcenter-late-cross-linear-step-one-before-zero-diag}
For every natural number \(t\) and natural number \(k\) and natural number \(r\), if \(t + r + 2 \le k\), then if \(\operatorname{zerodiagonal}(t + r + 1)\), then if \(\operatorname{settledby}(t + r + 1, t + \operatorname{twopow}(k))\), then \(\operatorname{lhcenter}(t + 1 + \operatorname{twopow}(k), \operatorname{twopow}(k) - (r + 1)) = \operatorname{bxor}(\operatorname{lhcenter}(t + \operatorname{twopow}(k), \operatorname{twopow}(k) - r), \operatorname{lhcenter}(t + \operatorname{twopow}(k), \operatorname{twopow}(k) - (r + 2)))\).
\end{theorem}
\begin{proof}
The goal follows from \(\operatorname{lhcenterlatecrosslinearsteponebeforezerodiag}(t, k, r)\).
\end{proof}
\begin{theorem}
\label{ll:paper:lhcenter-late-cross-linear-step-two-before-zero-diag}
For every natural number \(t\) and natural number \(k\) and natural number \(r\), if \(t + r + 2 \le k\), then if \(\operatorname{zerodiagonal}(t + r + 2)\), then if \(\operatorname{settledby}(t + r + 2, t + \operatorname{twopow}(k))\), then \(\operatorname{lhcenter}(t + 1 + \operatorname{twopow}(k), \operatorname{twopow}(k) - (r + 1)) = \operatorname{bxor}(\operatorname{lhcenter}(t + \operatorname{twopow}(k), \operatorname{twopow}(k) - r), \operatorname{lhcenter}(t + \operatorname{twopow}(k), \operatorname{twopow}(k) - (r + 1)))\).
\end{theorem}
\begin{proof}
The goal follows from \(\operatorname{lhcenterlatecrosslinearsteptwobeforezerodiag}(t, k, r)\).
\end{proof}
\begin{theorem}
\label{ll:paper:winwitt-sector-two-linear-step-one-before-zero-diag}
For every natural number \(t\) and natural number \(k\) and natural number \(r\), if \(t \le r\), then if \(t + r + 2 \le k\), then if \(\operatorname{zerodiagonal}(t + r + 1)\), then if \(\operatorname{settledby}(t + r + 1, t + \operatorname{twopow}(k))\), then \(\operatorname{wittbitbool}(\operatorname{wittwin}(t + 1, \operatorname{twopow}(k)), t + \operatorname{twopow}(k + 1) - r) = \operatorname{bxor}(\operatorname{wittbitbool}(\operatorname{wittwin}(t, \operatorname{twopow}(k)), t + \operatorname{twopow}(k + 1) - r), \operatorname{wittbitbool}(\operatorname{wittwin}(t, \operatorname{twopow}(k)), t + \operatorname{twopow}(k + 1) - (r + 2)))\).
\end{theorem}
\begin{proof}
The goal follows from \(\operatorname{winwittsectortwolinearsteponebeforezerodiag}(t, k, r)\).
\end{proof}
\begin{theorem}
\label{ll:paper:winwitt-sector-two-linear-step-two-before-zero-diag}
For every natural number \(t\) and natural number \(k\) and natural number \(r\), if \(t \le r\), then if \(t + r + 2 \le k\), then if \(\operatorname{zerodiagonal}(t + r + 2)\), then if \(\operatorname{settledby}(t + r + 2, t + \operatorname{twopow}(k))\), then \(\operatorname{wittbitbool}(\operatorname{wittwin}(t + 1, \operatorname{twopow}(k)), t + \operatorname{twopow}(k + 1) - r) = \operatorname{bxor}(\operatorname{wittbitbool}(\operatorname{wittwin}(t, \operatorname{twopow}(k)), t + \operatorname{twopow}(k + 1) - r), \operatorname{wittbitbool}(\operatorname{wittwin}(t, \operatorname{twopow}(k)), t + \operatorname{twopow}(k + 1) - (r + 1)))\).
\end{theorem}
\begin{proof}
The goal follows from \(\operatorname{winwittsectortwolinearsteptwobeforezerodiag}(t, k, r)\).
\end{proof}
\begin{theorem}
\label{ll:paper:lhcenter-late-cross-twin-before-zero-diag}
For every natural number \(t\) and natural number \(k\) and natural number \(r\), if \(t + r + 2 \le k\), then if \(\operatorname{zerodiagonal}(t + r + 2)\), then if \(\operatorname{settledby}(t + r + 2, t + \operatorname{twopow}(k))\), then \(\operatorname{lhcenter}(t + \operatorname{twopow}(k), \operatorname{twopow}(k) - r) = \operatorname{lhcenter}(t + \operatorname{twopow}(k), \operatorname{twopow}(k) - (r + 1))\).
\end{theorem}
\begin{proof}
The goal follows from \(\operatorname{lhcenterlatecrosstwinbeforezerodiag}(t, k, r)\).
\end{proof}
\begin{theorem}
\label{ll:paper:lhcenter-late-shift-twin-before-zero-diag}
For every natural number \(t\) and natural number \(k\) and natural number \(p\), if \(p \le t\), then if \(t - p + 2 \le k\), then if \(\operatorname{zerodiagonal}(t - p + 2)\), then if \(\operatorname{settledby}(t - p + 2, t + 2 + \operatorname{twopow}(k))\), then \(\operatorname{lhcenter}(t + 2 + \operatorname{twopow}(k), p + 2 + \operatorname{twopow}(k)) = \operatorname{lhcenter}(t + 2 + \operatorname{twopow}(k), p + 1 + \operatorname{twopow}(k))\).
\end{theorem}
\begin{proof}
The goal follows from \(\operatorname{lhcenterlateshifttwinbeforezerodiag}(t, k, p)\).
\end{proof}
\begin{theorem}
\label{ll:paper:winwitt-sector-two-twin-before-zero-diag}
For every natural number \(t\) and natural number \(k\) and natural number \(r\), if \(t + r + 2 \le k\), then if \(\operatorname{zerodiagonal}(t + r + 2)\), then if \(\operatorname{settledby}(t + r + 2, t + \operatorname{twopow}(k))\), then \(\operatorname{bxor}(\operatorname{wittbitbool}(\operatorname{wittwin}(t, \operatorname{twopow}(k)), t + \operatorname{twopow}(k + 1) - r), \operatorname{wittbitbool}(\operatorname{wittwin}(t, \operatorname{twopow}(k)), t + \operatorname{twopow}(k + 1) - (r + 1))) = \operatorname{bxor}(\operatorname{rhcenter}(t, r), \operatorname{rhcenter}(t, r + 1))\).
\end{theorem}
\begin{proof}
The goal follows from \(\operatorname{winwittsectortwotwinbeforezerodiag}(t, k, r)\).
\end{proof}
\begin{theorem}
\label{ll:paper:winwitt-sector-two-twin-zero-gap-before-zero-diag}
For every natural number \(t\) and natural number \(k\) and natural number \(r\), if \(t < r\), then if \(t + r + 2 \le k\), then if \(\operatorname{zerodiagonal}(t + r + 2)\), then if \(\operatorname{settledby}(t + r + 2, t + \operatorname{twopow}(k))\), then \(\operatorname{wittbitbool}(\operatorname{wittwin}(t, \operatorname{twopow}(k)), t + \operatorname{twopow}(k + 1) - r) = \operatorname{wittbitbool}(\operatorname{wittwin}(t, \operatorname{twopow}(k)), t + \operatorname{twopow}(k + 1) - (r + 1))\).
\end{theorem}
\begin{proof}
The goal follows from \(\operatorname{winwittsectortwotwinzerogapbeforezerodiag}(t, k, r)\).
\end{proof}
\begin{theorem}
\label{ll:paper:winwitt-sector-one-xor-sector-two-twin-before-zero-diag}
For every natural number \(t\) and natural number \(k\) and natural number \(r\), if \(t + r + 2 \le k\), then if \(\operatorname{zerodiagonal}(t + r + 2)\), then if \(\operatorname{settledby}(t + r + 2, t + \operatorname{twopow}(k))\), then \(\operatorname{bxor}(\operatorname{bxor}(\operatorname{wittbitbool}(\operatorname{wittwin}(t, \operatorname{twopow}(k)), t + \operatorname{twopow}(k) - r), \operatorname{wittbitbool}(\operatorname{wittwin}(t, \operatorname{twopow}(k)), t + \operatorname{twopow}(k + 1) - r)), \operatorname{bxor}(\operatorname{wittbitbool}(\operatorname{wittwin}(t, \operatorname{twopow}(k)), t + \operatorname{twopow}(k) - (r + 1)), \operatorname{wittbitbool}(\operatorname{wittwin}(t, \operatorname{twopow}(k)), t + \operatorname{twopow}(k + 1) - (r + 1)))) = \operatorname{bxor}(\operatorname{rhcenter}(t + \operatorname{twopow}(k), r), \operatorname{rhcenter}(t + \operatorname{twopow}(k), r + 1))\).
\end{theorem}
\begin{proof}
The goal follows from \(\operatorname{winwittsectoronexorsectortwotwinbeforezerodiag}(t, k, r)\).
\end{proof}
\begin{theorem}
\label{ll:paper:lhcenter-late-shift-antiperiodic-of-doubles}
For every natural number \(P\) and natural number \(t\) and natural number \(k\) and natural number \(p\), if \(p \le t\), then if \(t - p + 1 \le k\), then if \(\operatorname{zerodiagonal}(t - p)\), then if \(\operatorname{doublesat}(P, t - p)\), then if \(\operatorname{settledby}(t - p + 1, t + 1 + \operatorname{twopow}(k))\), then \(\operatorname{lhcenter}(t + 1 + P + \operatorname{twopow}(k), p + P + \operatorname{twopow}(k)) = \operatorname{not}(\operatorname{lhcenter}(t + 1 + \operatorname{twopow}(k), p + \operatorname{twopow}(k)))\).
\end{theorem}
\begin{proof}
The goal follows from \(\operatorname{lhcenterlateshiftantiperiodicofdoubles}(P, t, k, p)\).
\end{proof}
\begin{theorem}
\label{ll:paper:lhcenter-late-cross-antiperiodic-of-doubles}
For every natural number \(P\) and natural number \(t\) and natural number \(k\) and natural number \(r\), if \(t + P + r + 1 \le k\), then if \(\operatorname{zerodiagonal}(t + P + r)\), then if \(\operatorname{doublesat}(P, t + P + r)\), then if \(\operatorname{settledby}(t + P + r + 1, t + \operatorname{twopow}(k))\), then \(\operatorname{lhcenter}(t + P + \operatorname{twopow}(k), \operatorname{twopow}(k) - (r + 1)) = \operatorname{not}(\operatorname{lhcenter}(t + \operatorname{twopow}(k), \operatorname{twopow}(k) - (r + P + 1)))\).
\end{theorem}
\begin{proof}
The goal follows from \(\operatorname{lhcenterlatecrossantiperiodicofdoubles}(P, t, k, r)\).
\end{proof}
\begin{theorem}
\label{ll:paper:winwitt-outer-left-strip-antiperiodic-of-doubles}
For every natural number \(P\) and natural number \(t\) and natural number \(k\) and natural number \(p\), if \(p \le t\), then if \(t + 1 < p + \operatorname{twopow}(k)\), then if \(t - p + 1 \le k\), then if \(\operatorname{zerodiagonal}(t - p)\), then if \(\operatorname{doublesat}(P, t - p)\), then if \(\operatorname{settledby}(t - p + 1, t + 1 + \operatorname{twopow}(k))\), then \(\operatorname{bxor}(\operatorname{wittbitbool}(\operatorname{wittwin}(t + 1, \operatorname{twopow}(k)), t + 1 + p + \operatorname{twopow}(k + 1)), \operatorname{wittbitbool}(\operatorname{wittwin}(t + 1 + P, \operatorname{twopow}(k)), t + 1 + P + (p + P) + \operatorname{twopow}(k + 1))) = \operatorname{not}(\operatorname{bxor}(\operatorname{lhcenter}(t + 1, p), \operatorname{lhcenter}(t + 1 + P, p + P)))\).
\end{theorem}
\begin{proof}
The goal follows from \(\operatorname{winwittouterleftstripantiperiodicofdoubles}(P, t, k, p)\).
\end{proof}
\begin{theorem}
\label{ll:paper:winwitt-scale-cross-antiperiodic-of-doubles}
For every natural number \(P\) and natural number \(t\) and natural number \(k\) and natural number \(r\), if \(t + P + r + 1 \le k\), then if \(\operatorname{zerodiagonal}(t + P + r)\), then if \(\operatorname{doublesat}(P, t + P + r)\), then if \(\operatorname{settledby}(t + P + r + 1, t + \operatorname{twopow}(k))\), then \(\operatorname{bxor}(\operatorname{wittbitbool}(\operatorname{wittwin}(t, \operatorname{twopow}(k)), t + \operatorname{twopow}(k + 1) - (r + P + 1)), \operatorname{wittbitbool}(\operatorname{wittwin}(t + P, \operatorname{twopow}(k)), t + P + \operatorname{twopow}(k + 1) - (r + 1))) = \operatorname{not}(\operatorname{bxor}(\operatorname{rhcenter}(t, r + P + 1), \operatorname{rhcenter}(t + P, r + 1)))\).
\end{theorem}
\begin{proof}
The goal follows from \(\operatorname{winwittscalecrossantiperiodicofdoubles}(P, t, k, r)\).
\end{proof}
\begin{theorem}
\label{ll:paper:winwitt-scale-cross-zero-gap-antiperiodic-of-doubles}
For every natural number \(P\) and natural number \(t\) and natural number \(k\) and natural number \(r\), if \(t + P \le r\), then if \(t + P + r + 1 \le k\), then if \(\operatorname{zerodiagonal}(t + P + r)\), then if \(\operatorname{doublesat}(P, t + P + r)\), then if \(\operatorname{settledby}(t + P + r + 1, t + \operatorname{twopow}(k))\), then \(\operatorname{wittbitbool}(\operatorname{wittwin}(t + P, \operatorname{twopow}(k)), t + P + \operatorname{twopow}(k + 1) - (r + 1)) = \operatorname{not}(\operatorname{wittbitbool}(\operatorname{wittwin}(t, \operatorname{twopow}(k)), t + \operatorname{twopow}(k + 1) - (r + P + 1)))\).
\end{theorem}
\begin{proof}
The goal follows from \(\operatorname{winwittscalecrosszerogapantiperiodicofdoubles}(P, t, k, r)\).
\end{proof}
\begin{theorem}
\label{ll:paper:winwitt-sector-one-xor-sector-two-antiperiodic-of-doubles}
For every natural number \(P\) and natural number \(t\) and natural number \(k\) and natural number \(r\), if \(t + P + r + 1 \le k\), then if \(\operatorname{zerodiagonal}(t + P + r)\), then if \(\operatorname{doublesat}(P, t + P + r)\), then if \(\operatorname{settledby}(t + P + r + 1, t + \operatorname{twopow}(k))\), then \(\operatorname{bxor}(\operatorname{bxor}(\operatorname{wittbitbool}(\operatorname{wittwin}(t, \operatorname{twopow}(k)), t + \operatorname{twopow}(k) - (r + P + 1)), \operatorname{wittbitbool}(\operatorname{wittwin}(t, \operatorname{twopow}(k)), t + \operatorname{twopow}(k + 1) - (r + P + 1))), \operatorname{bxor}(\operatorname{wittbitbool}(\operatorname{wittwin}(t + P, \operatorname{twopow}(k)), t + P + \operatorname{twopow}(k) - (r + 1)), \operatorname{wittbitbool}(\operatorname{wittwin}(t + P, \operatorname{twopow}(k)), t + P + \operatorname{twopow}(k + 1) - (r + 1)))) = \operatorname{not}(\operatorname{bxor}(\operatorname{rhcenter}(t + \operatorname{twopow}(k), r + P + 1), \operatorname{rhcenter}(t + P + \operatorname{twopow}(k), r + 1)))\).
\end{theorem}
\begin{proof}
The goal follows from \(\operatorname{winwittsectoronexorsectortwoantiperiodicofdoubles}(P, t, k, r)\).
\end{proof}
\begin{theorem}
\label{ll:paper:ldiag-periodic-of-doubles-le}
For every natural number \(P\) and natural number \(z\) and natural number \(d\) and natural number \(s\) and natural number \(k\), if \(\operatorname{doublesat}(P, z)\), then if \(d \le z\), then if \(d \le k\), then if \(\operatorname{settledby}(d, s + \operatorname{twopow}(k))\), then \(\operatorname{ldiag}(d, s + P + \operatorname{twopow}(k)) = \operatorname{ldiag}(d, s + \operatorname{twopow}(k))\).
\end{theorem}
\begin{proof}
The goal follows from \(\operatorname{ldiagperiodicofdoublesle}(P, z, d, s, k)\).
\end{proof}
\begin{theorem}
\label{ll:paper:lhcenter-late-cross-periodic-of-doubles-le}
For every natural number \(P\) and natural number \(z\) and natural number \(t\) and natural number \(k\) and natural number \(r\), if \(t + P + r \le z\), then if \(z \le k\), then if \(\operatorname{doublesat}(P, z)\), then if \(\operatorname{settledby}(t + P + r, t + \operatorname{twopow}(k))\), then \(\operatorname{lhcenter}(t + P + \operatorname{twopow}(k), \operatorname{twopow}(k) - r) = \operatorname{lhcenter}(t + \operatorname{twopow}(k), \operatorname{twopow}(k) - (r + P))\).
\end{theorem}
\begin{proof}
The goal follows from \(\operatorname{lhcenterlatecrossperiodicofdoublesle}(P, z, t, k, r)\).
\end{proof}
\begin{theorem}
\label{ll:paper:winwitt-scale-cross-periodic-of-doubles-le}
For every natural number \(P\) and natural number \(z\) and natural number \(t\) and natural number \(k\) and natural number \(r\), if \(t + P + r \le z\), then if \(z \le k\), then if \(\operatorname{doublesat}(P, z)\), then if \(\operatorname{settledby}(t + P + r, t + \operatorname{twopow}(k))\), then \(\operatorname{bxor}(\operatorname{wittbitbool}(\operatorname{wittwin}(t, \operatorname{twopow}(k)), t + \operatorname{twopow}(k + 1) - (r + P)), \operatorname{wittbitbool}(\operatorname{wittwin}(t + P, \operatorname{twopow}(k)), t + P + \operatorname{twopow}(k + 1) - r)) = \operatorname{bxor}(\operatorname{rhcenter}(t, r + P), \operatorname{rhcenter}(t + P, r))\).
\end{theorem}
\begin{proof}
The goal follows from \(\operatorname{winwittscalecrossperiodicofdoublesle}(P, z, t, k, r)\).
\end{proof}
\begin{theorem}
\label{ll:paper:winwitt-scale-cross-zero-gap-periodic-of-doubles-le}
For every natural number \(P\) and natural number \(z\) and natural number \(t\) and natural number \(k\) and natural number \(r\), if \(t + P < r\), then if \(t + P + r \le z\), then if \(z \le k\), then if \(\operatorname{doublesat}(P, z)\), then if \(\operatorname{settledby}(t + P + r, t + \operatorname{twopow}(k))\), then \(\operatorname{wittbitbool}(\operatorname{wittwin}(t + P, \operatorname{twopow}(k)), t + P + \operatorname{twopow}(k + 1) - r) = \operatorname{wittbitbool}(\operatorname{wittwin}(t, \operatorname{twopow}(k)), t + \operatorname{twopow}(k + 1) - (r + P))\).
\end{theorem}
\begin{proof}
The goal follows from \(\operatorname{winwittscalecrosszerogapperiodicofdoublesle}(P, z, t, k, r)\).
\end{proof}
\begin{theorem}
\label{ll:paper:lhcenter-late-cross-doubles-eq-reset-step}
For every natural number \(P\) and natural number \(t\) and natural number \(k\) and natural number \(r\), if \(t + P + r + 2 \le k\), then if \(\operatorname{zerodiagonal}(t + P + r)\), then if \(\operatorname{doublesat}(P, t + P + r)\), then if \(\operatorname{settledby}(t + P + r + 1, t + \operatorname{twopow}(k))\), then if \(\operatorname{settledby}(t + P + r + 2, t + \operatorname{twopow}(k))\), then \(\operatorname{lhcenter}(t + P + \operatorname{twopow}(k), \operatorname{twopow}(k) - (r + 1)) = \operatorname{lhcenter}(t + 1 + \operatorname{twopow}(k), \operatorname{twopow}(k) - (r + P + 2))\).
\end{theorem}
\begin{proof}
The goal follows from \(\operatorname{lhcenterlatecrossdoubleseqresetstep}(P, t, k, r)\).
\end{proof}
\begin{theorem}
\label{ll:paper:lhcenter-late-cross-reset-antiperiodic-of-doubles}
For every natural number \(P\) and natural number \(t\) and natural number \(k\) and natural number \(r\), if \(t + P + r + 2 \le k\), then if \(\operatorname{zerodiagonal}(t + P + r)\), then if \(\operatorname{doublesat}(P, t + P + r)\), then if \(\operatorname{settledby}(t + P + r + 1, t + \operatorname{twopow}(k))\), then if \(\operatorname{settledby}(t + P + r + 2, t + \operatorname{twopow}(k))\), then \(\operatorname{lhcenter}(t + 1 + P + \operatorname{twopow}(k), \operatorname{twopow}(k) - (r + 2)) = \operatorname{not}(\operatorname{lhcenter}(t + 1 + \operatorname{twopow}(k), \operatorname{twopow}(k) - (r + P + 2)))\).
\end{theorem}
\begin{proof}
The goal follows from \(\operatorname{lhcenterlatecrossresetantiperiodicofdoubles}(P, t, k, r)\).
\end{proof}
\begin{theorem}
\label{ll:paper:winwitt-scale-cross-doubles-eq-reset-step}
For every natural number \(P\) and natural number \(t\) and natural number \(k\) and natural number \(r\), if \(t + P + r + 2 \le k\), then if \(\operatorname{zerodiagonal}(t + P + r)\), then if \(\operatorname{doublesat}(P, t + P + r)\), then if \(\operatorname{settledby}(t + P + r + 1, t + \operatorname{twopow}(k))\), then if \(\operatorname{settledby}(t + P + r + 2, t + \operatorname{twopow}(k))\), then \(\operatorname{bxor}(\operatorname{wittbitbool}(\operatorname{wittwin}(t + P, \operatorname{twopow}(k)), t + P + \operatorname{twopow}(k + 1) - (r + 1)), \operatorname{wittbitbool}(\operatorname{wittwin}(t + 1, \operatorname{twopow}(k)), t + \operatorname{twopow}(k + 1) - (r + P + 1))) = \operatorname{bxor}(\operatorname{rhcenter}(t + P, r + 1), \operatorname{rhcenter}(t + 1, r + P + 2))\).
\end{theorem}
\begin{proof}
The goal follows from \(\operatorname{winwittscalecrossdoubleseqresetstep}(P, t, k, r)\).
\end{proof}
\begin{theorem}
\label{ll:paper:winwitt-scale-cross-zero-gap-doubles-eq-reset-step}
For every natural number \(P\) and natural number \(t\) and natural number \(k\) and natural number \(r\), if \(t + P \le r\), then if \(t + P + r + 2 \le k\), then if \(\operatorname{zerodiagonal}(t + P + r)\), then if \(\operatorname{doublesat}(P, t + P + r)\), then if \(\operatorname{settledby}(t + P + r + 1, t + \operatorname{twopow}(k))\), then if \(\operatorname{settledby}(t + P + r + 2, t + \operatorname{twopow}(k))\), then \(\operatorname{wittbitbool}(\operatorname{wittwin}(t + P, \operatorname{twopow}(k)), t + P + \operatorname{twopow}(k + 1) - (r + 1)) = \operatorname{wittbitbool}(\operatorname{wittwin}(t + 1, \operatorname{twopow}(k)), t + \operatorname{twopow}(k + 1) - (r + P + 1))\).
\end{theorem}
\begin{proof}
The goal follows from \(\operatorname{winwittscalecrosszerogapdoubleseqresetstep}(P, t, k, r)\).
\end{proof}
\begin{theorem}
\label{ll:paper:winwitt-scale-cross-reset-antiperiodic-of-doubles}
For every natural number \(P\) and natural number \(t\) and natural number \(k\) and natural number \(r\), if \(t + P + r + 2 \le k\), then if \(\operatorname{zerodiagonal}(t + P + r)\), then if \(\operatorname{doublesat}(P, t + P + r)\), then if \(\operatorname{settledby}(t + P + r + 1, t + \operatorname{twopow}(k))\), then if \(\operatorname{settledby}(t + P + r + 2, t + \operatorname{twopow}(k))\), then \(\operatorname{bxor}(\operatorname{wittbitbool}(\operatorname{wittwin}(t + 1, \operatorname{twopow}(k)), t + \operatorname{twopow}(k + 1) - (r + P + 1)), \operatorname{wittbitbool}(\operatorname{wittwin}(t + 1 + P, \operatorname{twopow}(k)), t + P + \operatorname{twopow}(k + 1) - (r + 1))) = \operatorname{not}(\operatorname{bxor}(\operatorname{rhcenter}(t + 1, r + P + 2), \operatorname{rhcenter}(t + 1 + P, r + 2)))\).
\end{theorem}
\begin{proof}
The goal follows from \(\operatorname{winwittscalecrossresetantiperiodicofdoubles}(P, t, k, r)\).
\end{proof}
\begin{theorem}
\label{ll:paper:winwitt-scale-cross-zero-gap-reset-antiperiodic-of-doubles}
For every natural number \(P\) and natural number \(t\) and natural number \(k\) and natural number \(r\), if \(t + P \le r\), then if \(t + P + r + 2 \le k\), then if \(\operatorname{zerodiagonal}(t + P + r)\), then if \(\operatorname{doublesat}(P, t + P + r)\), then if \(\operatorname{settledby}(t + P + r + 1, t + \operatorname{twopow}(k))\), then if \(\operatorname{settledby}(t + P + r + 2, t + \operatorname{twopow}(k))\), then \(\operatorname{wittbitbool}(\operatorname{wittwin}(t + 1 + P, \operatorname{twopow}(k)), t + P + \operatorname{twopow}(k + 1) - (r + 1)) = \operatorname{not}(\operatorname{wittbitbool}(\operatorname{wittwin}(t + 1, \operatorname{twopow}(k)), t + \operatorname{twopow}(k + 1) - (r + P + 1)))\).
\end{theorem}
\begin{proof}
The goal follows from \(\operatorname{winwittscalecrosszerogapresetantiperiodicofdoubles}(P, t, k, r)\).
\end{proof}
\begin{theorem}
\label{ll:paper:winwitt-sector-one-xor-sector-two-doubles-eq-reset-step}
For every natural number \(P\) and natural number \(t\) and natural number \(k\) and natural number \(r\), if \(t + P + r + 2 \le k\), then if \(\operatorname{zerodiagonal}(t + P + r)\), then if \(\operatorname{doublesat}(P, t + P + r)\), then if \(\operatorname{settledby}(t + P + r + 1, t + \operatorname{twopow}(k))\), then if \(\operatorname{settledby}(t + P + r + 2, t + \operatorname{twopow}(k))\), then \(\operatorname{bxor}(\operatorname{bxor}(\operatorname{wittbitbool}(\operatorname{wittwin}(t + P, \operatorname{twopow}(k)), t + P + \operatorname{twopow}(k) - (r + 1)), \operatorname{wittbitbool}(\operatorname{wittwin}(t + P, \operatorname{twopow}(k)), t + P + \operatorname{twopow}(k + 1) - (r + 1))), \operatorname{bxor}(\operatorname{wittbitbool}(\operatorname{wittwin}(t + 1, \operatorname{twopow}(k)), t + \operatorname{twopow}(k) - (r + P + 1)), \operatorname{wittbitbool}(\operatorname{wittwin}(t + 1, \operatorname{twopow}(k)), t + \operatorname{twopow}(k + 1) - (r + P + 1)))) = \operatorname{bxor}(\operatorname{rhcenter}(t + P + \operatorname{twopow}(k), r + 1), \operatorname{rhcenter}(t + 1 + \operatorname{twopow}(k), r + P + 2))\).
\end{theorem}
\begin{proof}
The goal follows from \(\operatorname{winwittsectoronexorsectortwodoubleseqresetstep}(P, t, k, r)\).
\end{proof}
\begin{theorem}
\label{ll:paper:winwitt-sector-one-xor-sector-two-reset-antiperiodic-of-doubles}
For every natural number \(P\) and natural number \(t\) and natural number \(k\) and natural number \(r\), if \(t + P + r + 2 \le k\), then if \(\operatorname{zerodiagonal}(t + P + r)\), then if \(\operatorname{doublesat}(P, t + P + r)\), then if \(\operatorname{settledby}(t + P + r + 1, t + \operatorname{twopow}(k))\), then if \(\operatorname{settledby}(t + P + r + 2, t + \operatorname{twopow}(k))\), then \(\operatorname{bxor}(\operatorname{bxor}(\operatorname{wittbitbool}(\operatorname{wittwin}(t + 1, \operatorname{twopow}(k)), t + \operatorname{twopow}(k) - (r + P + 1)), \operatorname{wittbitbool}(\operatorname{wittwin}(t + 1, \operatorname{twopow}(k)), t + \operatorname{twopow}(k + 1) - (r + P + 1))), \operatorname{bxor}(\operatorname{wittbitbool}(\operatorname{wittwin}(t + 1 + P, \operatorname{twopow}(k)), t + P + \operatorname{twopow}(k) - (r + 1)), \operatorname{wittbitbool}(\operatorname{wittwin}(t + 1 + P, \operatorname{twopow}(k)), t + P + \operatorname{twopow}(k + 1) - (r + 1)))) = \operatorname{not}(\operatorname{bxor}(\operatorname{rhcenter}(t + 1 + \operatorname{twopow}(k), r + P + 2), \operatorname{rhcenter}(t + 1 + P + \operatorname{twopow}(k), r + 2)))\).
\end{theorem}
\begin{proof}
The goal follows from \(\operatorname{winwittsectoronexorsectortworesetantiperiodicofdoubles}(P, t, k, r)\).
\end{proof}
\begin{theorem}
\label{ll:paper:lhcenter-late-twopow-sub-one-eq-zcenter-before-zero-diag}
For every natural number \(t\) and natural number \(k\), if \(t + 2 \le k\), then if \(\operatorname{zerodiagonal}(t + 2)\), then if \(\operatorname{settledby}(t, t + \operatorname{twopow}(k))\), then if \(\operatorname{settledby}(t + 2, t + \operatorname{twopow}(k))\), then \(\operatorname{zcentereq}(t, \operatorname{lhcenter}(t + \operatorname{twopow}(k), \operatorname{twopow}(k) - 1))\).
\end{theorem}
\begin{proof}
The goal follows from \(\operatorname{lhcenterlatetwopowsuboneeqzcenterbeforezerodiag}(t, k)\).
\end{proof}
\begin{theorem}
\label{ll:paper:winwitt-testbit-sector-two-sub-one-zcenter-before-zero-diag}
For every natural number \(t\) and natural number \(k\), if \(t + 2 \le k\), then if \(\operatorname{zerodiagonal}(t + 2)\), then if \(\operatorname{settledby}(t, t + \operatorname{twopow}(k))\), then if \(\operatorname{settledby}(t + 2, t + \operatorname{twopow}(k))\), then \(\operatorname{zcentereq}(t, \operatorname{bxor}(\operatorname{wittbitbool}(\operatorname{wittwin}(t, \operatorname{twopow}(k)), t + \operatorname{twopow}(k + 1) - 1), \operatorname{rhcenter}(t, 1)))\).
\end{theorem}
\begin{proof}
The goal follows from \(\operatorname{winwitttestbitsectortwosubonezcenterbeforezerodiag}(t, k)\).
\end{proof}
\begin{theorem}
\label{ll:paper:winwitt-sector-one-xor-sector-two-sub-one-zcenter-before-zero-diag}
For every natural number \(t\) and natural number \(k\), if \(t + 2 \le k\), then if \(\operatorname{zerodiagonal}(t + 2)\), then if \(\operatorname{settledby}(t, t + \operatorname{twopow}(k))\), then if \(\operatorname{settledby}(t + 2, t + \operatorname{twopow}(k))\), then \(\operatorname{zcentereq}(t, \operatorname{bxor}(\operatorname{rhcenter}(t + \operatorname{twopow}(k), 1), \operatorname{bxor}(\operatorname{wittbitbool}(\operatorname{wittwin}(t, \operatorname{twopow}(k)), t + \operatorname{twopow}(k) - 1), \operatorname{wittbitbool}(\operatorname{wittwin}(t, \operatorname{twopow}(k)), t + \operatorname{twopow}(k + 1) - 1))))\).
\end{theorem}
\begin{proof}
The goal follows from \(\operatorname{winwittsectoronexorsectortwosubonezcenterbeforezerodiag}(t, k)\).
\end{proof}
\begin{theorem}
\label{ll:paper:lhcenter-late-shift-twopow-antiperiodic-zcenter-of-doubles}
For every natural number \(P\) and natural number \(z\) and natural number \(k\), if \(z + 1 \le k\), then if \(\operatorname{zerodiagonal}(z)\), then if \(\operatorname{doublesat}(P, z)\), then if \(\operatorname{settledby}(z + 1, z + 1 + \operatorname{twopow}(k))\), then \(\operatorname{zcentereq}(z + 1, \operatorname{not}(\operatorname{lhcenter}(z + 1 + P + \operatorname{twopow}(k), P + \operatorname{twopow}(k))))\).
\end{theorem}
\begin{proof}
The goal follows from \(\operatorname{lhcenterlateshifttwopowantiperiodiczcenterofdoubles}(P, z, k)\).
\end{proof}
\begin{theorem}
\label{ll:paper:winwitt-outer-left-strip-zcenter-antiperiodic-of-doubles}
For every natural number \(P\) and natural number \(z\) and natural number \(k\), if \(z + 1 \le k\), then if \(\operatorname{zerodiagonal}(z)\), then if \(\operatorname{doublesat}(P, z)\), then if \(\operatorname{settledby}(z + 1, z + 1 + \operatorname{twopow}(k))\), then \(\operatorname{zcentereq}(z + 1, \operatorname{not}(\operatorname{bxor}(\operatorname{wittbitbool}(\operatorname{wittwin}(z + 1 + P, \operatorname{twopow}(k)), z + 1 + P + P + \operatorname{twopow}(k + 1)), \operatorname{lhcenter}(z + 1 + P, P))))\).
\end{theorem}
\begin{proof}
The goal follows from \(\operatorname{winwittouterleftstripzcenterantiperiodicofdoubles}(P, z, k)\).
\end{proof}
\begin{theorem}
\label{ll:paper:winwitt-sector-one-xor-sector-two-zcenter-antiperiodic-of-doubles}
For every natural number \(P\) and natural number \(z\) and natural number \(k\), if \(z + 1 + P + P < \operatorname{twopow}(k)\), then if \(z + 1 \le k\), then if \(\operatorname{zerodiagonal}(z)\), then if \(\operatorname{doublesat}(P, z)\), then if \(\operatorname{settledby}(z + 1, z + 1 + \operatorname{twopow}(k))\), then \(\operatorname{zcentereq}(z + 1, \operatorname{not}(\operatorname{bxor}(\operatorname{lhcenter}(z + 1 + P + \operatorname{twopow}(k), P), \operatorname{bxor}(\operatorname{wittbitbool}(\operatorname{wittwin}(z + 1 + P, \operatorname{twopow}(k)), z + 1 + P + P + \operatorname{twopow}(k)), \operatorname{wittbitbool}(\operatorname{wittwin}(z + 1 + P, \operatorname{twopow}(k)), z + 1 + P + P + \operatorname{twopow}(k + 1))))))\).
\end{theorem}
\begin{proof}
The goal follows from \(\operatorname{winwittsectoronexorsectortwozcenterantiperiodicofdoubles}(P, z, k)\).
\end{proof}
\begin{theorem}
\label{ll:paper:lhcenter-late-shift-twopow-periodic-zcenter-of-doubles-le}
For every natural number \(P\) and natural number \(z\) and natural number \(d\) and natural number \(k\), if \(\operatorname{doublesat}(P, z)\), then if \(d \le z\), then if \(d \le k\), then if \(\operatorname{settledby}(d, d + \operatorname{twopow}(k))\), then \(\operatorname{zcentereq}(d, \operatorname{lhcenter}(d + P + \operatorname{twopow}(k), P + \operatorname{twopow}(k)))\).
\end{theorem}
\begin{proof}
The goal follows from \(\operatorname{lhcenterlateshifttwopowperiodiczcenterofdoublesle}(P, z, d, k)\).
\end{proof}
\begin{theorem}
\label{ll:paper:lhcenter-late-twopow-sub-one-reset-after-zero-diag}
For every natural number \(z\) and natural number \(k\), if \(z + 2 \le k\), then if \(\operatorname{zerodiagonal}(z)\), then if \(\operatorname{settledby}(z + 2, z + 1 + \operatorname{twopow}(k))\), then \(\operatorname{lhcenter}(z + 2 + \operatorname{twopow}(k), \operatorname{twopow}(k) - 1) = \operatorname{not}(\operatorname{lhcenter}(z + 1 + \operatorname{twopow}(k), \operatorname{twopow}(k)))\).
\end{theorem}
\begin{proof}
The goal follows from \(\operatorname{lhcenterlatetwopowsuboneresetafterzerodiag}(z, k)\).
\end{proof}
\begin{theorem}
\label{ll:paper:winwitt-sector-two-center-to-sub-one-step-complement-after-zero-diag}
For every natural number \(z\) and natural number \(k\), if \(z + 2 \le k\), then if \(\operatorname{zerodiagonal}(z)\), then if \(\operatorname{settledby}(z + 2, z + 1 + \operatorname{twopow}(k))\), then \(\operatorname{bxor}(\operatorname{wittbitbool}(\operatorname{wittwin}(z + 1, \operatorname{twopow}(k)), z + 1 + \operatorname{twopow}(k + 1)), \operatorname{wittbitbool}(\operatorname{wittwin}(z + 2, \operatorname{twopow}(k)), z + 2 + \operatorname{twopow}(k + 1) - 1)) = \operatorname{not}(\operatorname{bxor}(\operatorname{center}(z + 1), \operatorname{rhcenter}(z + 2, 1)))\).
\end{theorem}
\begin{proof}
The goal follows from \(\operatorname{winwittsectortwocentertosubonestepcomplementafterzerodiag}(z, k)\).
\end{proof}
\begin{theorem}
\label{ll:paper:winwitt-sector-one-xor-sector-two-center-to-sub-one-step-complement-after-zero-diag}
For every natural number \(z\) and natural number \(k\), if \(z + 2 \le k\), then if \(\operatorname{zerodiagonal}(z)\), then if \(\operatorname{settledby}(z + 2, z + 1 + \operatorname{twopow}(k))\), then \(\operatorname{bxor}(\operatorname{bxor}(\operatorname{wittbitbool}(\operatorname{wittwin}(z + 1, \operatorname{twopow}(k)), z + 1 + \operatorname{twopow}(k)), \operatorname{wittbitbool}(\operatorname{wittwin}(z + 1, \operatorname{twopow}(k)), z + 1 + \operatorname{twopow}(k + 1))), \operatorname{bxor}(\operatorname{wittbitbool}(\operatorname{wittwin}(z + 2, \operatorname{twopow}(k)), z + 2 + \operatorname{twopow}(k) - 1), \operatorname{wittbitbool}(\operatorname{wittwin}(z + 2, \operatorname{twopow}(k)), z + 2 + \operatorname{twopow}(k + 1) - 1))) = \operatorname{not}(\operatorname{bxor}(\operatorname{center}(z + 1 + \operatorname{twopow}(k)), \operatorname{rhcenter}(z + 2 + \operatorname{twopow}(k), 1)))\).
\end{theorem}
\begin{proof}
The goal follows from \(\operatorname{winwittsectoronexorsectortwocentertosubonestepcomplementafterzerodiag}(z, k)\).
\end{proof}
\begin{theorem}
\label{ll:paper:lhcenter-late-twopow-sub-one-eq-not-zcenter-after-zero-diag}
For every natural number \(z\) and natural number \(k\), if \(z + 2 \le k\), then if \(\operatorname{zerodiagonal}(z)\), then if \(\operatorname{settledby}(z + 1, z + 1 + \operatorname{twopow}(k))\), then if \(\operatorname{settledby}(z + 2, z + 1 + \operatorname{twopow}(k))\), then \(\operatorname{zcentereq}(z + 1, \operatorname{not}(\operatorname{lhcenter}(z + 2 + \operatorname{twopow}(k), \operatorname{twopow}(k) - 1)))\).
\end{theorem}
\begin{proof}
The goal follows from \(\operatorname{lhcenterlatetwopowsuboneeqnotzcenterafterzerodiag}(z, k)\).
\end{proof}
\begin{theorem}
\label{ll:paper:winwitt-testbit-sector-two-sub-one-not-zcenter-after-zero-diag}
For every natural number \(z\) and natural number \(k\), if \(z + 2 \le k\), then if \(\operatorname{zerodiagonal}(z)\), then if \(\operatorname{settledby}(z + 1, z + 1 + \operatorname{twopow}(k))\), then if \(\operatorname{settledby}(z + 2, z + 1 + \operatorname{twopow}(k))\), then \(\operatorname{zcentereq}(z + 1, \operatorname{not}(\operatorname{bxor}(\operatorname{wittbitbool}(\operatorname{wittwin}(z + 2, \operatorname{twopow}(k)), z + 2 + \operatorname{twopow}(k + 1) - 1), \operatorname{rhcenter}(z + 2, 1))))\).
\end{theorem}
\begin{proof}
The goal follows from \(\operatorname{winwitttestbitsectortwosubonenotzcenterafterzerodiag}(z, k)\).
\end{proof}
\begin{theorem}
\label{ll:paper:winwitt-sector-one-xor-sector-two-sub-one-not-zcenter-after-zero-diag}
For every natural number \(z\) and natural number \(k\), if \(z + 2 \le k\), then if \(\operatorname{zerodiagonal}(z)\), then if \(\operatorname{settledby}(z + 1, z + 1 + \operatorname{twopow}(k))\), then if \(\operatorname{settledby}(z + 2, z + 1 + \operatorname{twopow}(k))\), then \(\operatorname{zcentereq}(z + 1, \operatorname{not}(\operatorname{bxor}(\operatorname{rhcenter}(z + 2 + \operatorname{twopow}(k), 1), \operatorname{bxor}(\operatorname{wittbitbool}(\operatorname{wittwin}(z + 2, \operatorname{twopow}(k)), z + 2 + \operatorname{twopow}(k) - 1), \operatorname{wittbitbool}(\operatorname{wittwin}(z + 2, \operatorname{twopow}(k)), z + 2 + \operatorname{twopow}(k + 1) - 1)))))\).
\end{theorem}
\begin{proof}
The goal follows from \(\operatorname{winwittsectoronexorsectortwosubonenotzcenterafterzerodiag}(z, k)\).
\end{proof}
\begin{theorem}
\label{ll:paper:lhcenter-late-shift-twopow-doubles-eq-after-zero-diag-sub-one}
For every natural number \(P\) and natural number \(z\) and natural number \(k\), if \(z + 2 \le k\), then if \(\operatorname{zerodiagonal}(z)\), then if \(\operatorname{doublesat}(P, z)\), then if \(\operatorname{settledby}(z + 1, z + 1 + \operatorname{twopow}(k))\), then if \(\operatorname{settledby}(z + 2, z + 1 + \operatorname{twopow}(k))\), then \(\operatorname{lhcenter}(z + 1 + P + \operatorname{twopow}(k), P + \operatorname{twopow}(k)) = \operatorname{lhcenter}(z + 2 + \operatorname{twopow}(k), \operatorname{twopow}(k) - 1)\).
\end{theorem}
\begin{proof}
The goal follows from \(\operatorname{lhcenterlateshifttwopowdoubleseqafterzerodiagsubone}(P, z, k)\).
\end{proof}
\begin{theorem}
\label{ll:paper:lhcenter-late-twopow-reset-three-after-zero-diag}
For every natural number \(z\) and natural number \(k\), if \(z + 2 \le k\), then if \(\operatorname{zerodiagonal}(z)\), then if \(\operatorname{settledby}(z + 2, z + 2 + \operatorname{twopow}(k))\), then \(\operatorname{lhcenter}(z + 3 + \operatorname{twopow}(k), \operatorname{twopow}(k)) = \operatorname{not}(\operatorname{lhcenter}(z + 2 + \operatorname{twopow}(k), 1 + \operatorname{twopow}(k)))\).
\end{theorem}
\begin{proof}
The goal follows from \(\operatorname{lhcenterlatetwopowresetthreeafterzerodiag}(z, k)\).
\end{proof}
\begin{theorem}
\label{ll:paper:winwitt-sector-two-one-to-center-step-complement-after-zero-diag}
For every natural number \(z\) and natural number \(k\), if \(z + 2 \le k\), then if \(\operatorname{zerodiagonal}(z)\), then if \(\operatorname{settledby}(z + 2, z + 2 + \operatorname{twopow}(k))\), then \(\operatorname{bxor}(\operatorname{wittbitbool}(\operatorname{wittwin}(z + 2, \operatorname{twopow}(k)), z + 2 + 1 + \operatorname{twopow}(k + 1)), \operatorname{wittbitbool}(\operatorname{wittwin}(z + 3, \operatorname{twopow}(k)), z + 3 + \operatorname{twopow}(k + 1))) = \operatorname{not}(\operatorname{bxor}(\operatorname{lhcenter}(z + 2, 1), \operatorname{center}(z + 3)))\).
\end{theorem}
\begin{proof}
The goal follows from \(\operatorname{winwittsectortwoonetocenterstepcomplementafterzerodiag}(z, k)\).
\end{proof}
\begin{theorem}
\label{ll:paper:winwitt-sector-one-xor-sector-two-one-to-center-step-complement-after-zero-diag}
For every natural number \(z\) and natural number \(k\), if \(z + 2 \le k\), then if \(\operatorname{zerodiagonal}(z)\), then if \(\operatorname{settledby}(z + 2, z + 2 + \operatorname{twopow}(k))\), then \(\operatorname{bxor}(\operatorname{bxor}(\operatorname{wittbitbool}(\operatorname{wittwin}(z + 2, \operatorname{twopow}(k)), z + 2 + 1 + \operatorname{twopow}(k)), \operatorname{wittbitbool}(\operatorname{wittwin}(z + 2, \operatorname{twopow}(k)), z + 2 + 1 + \operatorname{twopow}(k + 1))), \operatorname{bxor}(\operatorname{wittbitbool}(\operatorname{wittwin}(z + 3, \operatorname{twopow}(k)), z + 3 + \operatorname{twopow}(k)), \operatorname{wittbitbool}(\operatorname{wittwin}(z + 3, \operatorname{twopow}(k)), z + 3 + \operatorname{twopow}(k + 1)))) = \operatorname{not}(\operatorname{bxor}(\operatorname{lhcenter}(z + 2 + \operatorname{twopow}(k), 1), \operatorname{center}(z + 3 + \operatorname{twopow}(k))))\).
\end{theorem}
\begin{proof}
The goal follows from \(\operatorname{winwittsectoronexorsectortwoonetocenterstepcomplementafterzerodiag}(z, k)\).
\end{proof}
\begin{theorem}
\label{ll:paper:lhcenter-late-cross-eq-of-winwittagree-zero-gap}
For every natural number \(t\) and natural number \(k\) and natural number \(Q\) and natural number \(r\) and natural number \(a\) and natural number \(L\), if \(t + Q < r\), then if \(t + Q + r < \operatorname{twopow}(k)\), then if \(a \le t + \operatorname{twopow}(k + 1) - r\), then if \(t + \operatorname{twopow}(k + 1) - r \le a + L\), then if \(\operatorname{winwittagree}(t, k, Q, a, L)\), then \(\operatorname{lhcenter}(t + Q + \operatorname{twopow}(k), \operatorname{twopow}(k) - r) = \operatorname{lhcenter}(t + \operatorname{twopow}(k), \operatorname{twopow}(k) - r)\).
\end{theorem}
\begin{proof}
The goal follows from \(\operatorname{lhcenterlatecrosseqofwinwittagreezerogap}(t, k, Q, r, a, L)\).
\end{proof}
\begin{theorem}
\label{ll:paper:not-winwittagree-two-scale-cross-zero-gap-of-zero-diag}
For every natural number \(t\) and natural number \(k\) and natural number \(r\) and natural number \(a\) and natural number \(L\), if \(t + 2 < r\), then if \(t + r + 2 \le k\), then if \(\operatorname{zerodiagonal}(t + r)\), then if \(\operatorname{settledby}(t + r, t + \operatorname{twopow}(k))\), then if \(\operatorname{settledby}(t + r + 2, t + 2 + \operatorname{twopow}(k))\), then if \(a \le t + \operatorname{twopow}(k + 1) - r\), then if \(t + \operatorname{twopow}(k + 1) - r \le a + L\), then not \(\operatorname{winwittagree}(t, k, 2, a, L)\).
\end{theorem}
\begin{proof}
The goal follows from \(\operatorname{notwinwittagreetwoscalecrosszerogapofzerodiag}(t, k, r, a, L)\).
\end{proof}
\begin{theorem}
\label{ll:paper:not-winwittagree-scale-cross-zero-gap-of-zero-and-after-zero-diag}
For every natural number \(t\) and natural number \(k\) and natural number \(Q\) and natural number \(r\) and natural number \(z\) and natural number \(a\) and natural number \(L\), if \(t + Q < r\), then if \(t + Q + r \le k\), then if \(\operatorname{zerodiagonal}(t + r)\), then if \(\operatorname{zerodiagonal}(z)\), then if \(t + Q + r = z + 2\), then if \(\operatorname{settledby}(t + r, t + \operatorname{twopow}(k))\), then if \(\operatorname{settledby}(z + 2, t + Q + \operatorname{twopow}(k))\), then if \(a \le t + \operatorname{twopow}(k + 1) - r\), then if \(t + \operatorname{twopow}(k + 1) - r \le a + L\), then not \(\operatorname{winwittagree}(t, k, Q, a, L)\).
\end{theorem}
\begin{proof}
The goal follows from \(\operatorname{notwinwittagreescalecrosszerogapofzeroandafterzerodiag}(t, k, Q, r, z, a, L)\).
\end{proof}
\begin{theorem}
\label{ll:paper:winwitt-scale-cross-zero-gap-shift-flip-of-doubles-and-winwittagree}
For every natural number \(Q\) and natural number \(t\) and natural number \(k\) and natural number \(r\) and natural number \(a\) and natural number \(L\), if \(t + Q \le r\), then if \(t + Q + r + 1 \le k\), then if \(\operatorname{zerodiagonal}(t + Q + r)\), then if \(\operatorname{doublesat}(Q, t + Q + r)\), then if \(\operatorname{settledby}(t + Q + r + 1, t + \operatorname{twopow}(k))\), then if \(a \le t + \operatorname{twopow}(k + 1) - (r + 1)\), then if \(t + \operatorname{twopow}(k + 1) - (r + 1) \le a + L\), then if \(\operatorname{winwittagree}(t, k, Q, a, L)\), then \(\operatorname{wittbitbool}(\operatorname{wittwin}(t, \operatorname{twopow}(k)), t + \operatorname{twopow}(k + 1) - (r + 1)) = \operatorname{not}(\operatorname{wittbitbool}(\operatorname{wittwin}(t, \operatorname{twopow}(k)), t + \operatorname{twopow}(k + 1) - (r + Q + 1)))\).
\end{theorem}
\begin{proof}
The goal follows from \(\operatorname{winwittscalecrosszerogapshiftflipofdoublesandwinwittagree}(Q, t, k, r, a, L)\).
\end{proof}
\begin{theorem}
\label{ll:paper:lhcenter-late-cross-shift-flip-of-doubles-and-winwittagree-zero-gap}
For every natural number \(Q\) and natural number \(t\) and natural number \(k\) and natural number \(r\) and natural number \(a\) and natural number \(L\), if \(t + Q \le r\), then if \(t + Q + r + 1 \le k\), then if \(\operatorname{zerodiagonal}(t + Q + r)\), then if \(\operatorname{doublesat}(Q, t + Q + r)\), then if \(\operatorname{settledby}(t + Q + r + 1, t + \operatorname{twopow}(k))\), then if \(a \le t + \operatorname{twopow}(k + 1) - (r + 1)\), then if \(t + \operatorname{twopow}(k + 1) - (r + 1) \le a + L\), then if \(\operatorname{winwittagree}(t, k, Q, a, L)\), then \(\operatorname{lhcenter}(t + \operatorname{twopow}(k), \operatorname{twopow}(k) - (r + 1)) = \operatorname{not}(\operatorname{lhcenter}(t + \operatorname{twopow}(k), \operatorname{twopow}(k) - (r + Q + 1)))\).
\end{theorem}
\begin{proof}
The goal follows from \(\operatorname{lhcenterlatecrossshiftflipofdoublesandwinwittagreezerogap}(Q, t, k, r, a, L)\).
\end{proof}
\begin{theorem}
\label{ll:paper:lhcenter-late-shift-eq-of-winwittagree-outer-left-strip}
For every natural number \(t\) and natural number \(k\) and natural number \(Q\) and natural number \(p\) and natural number \(a\) and natural number \(L\), if \(t + Q + p < \operatorname{twopow}(k)\), then if \(a \le t + p + \operatorname{twopow}(k + 1)\), then if \(t + p + \operatorname{twopow}(k + 1) \le a + L\), then if \(\operatorname{lhcenter}(t + Q, p) = \operatorname{lhcenter}(t, p)\), then if \(\operatorname{winwittagree}(t, k, Q, a, L)\), then \(\operatorname{lhcenter}(t + Q + \operatorname{twopow}(k), p + \operatorname{twopow}(k)) = \operatorname{lhcenter}(t + \operatorname{twopow}(k), p + \operatorname{twopow}(k))\).
\end{theorem}
\begin{proof}
The goal follows from \(\operatorname{lhcenterlateshifteqofwinwittagreeouterleftstrip}(t, k, Q, p, a, L)\).
\end{proof}
\begin{theorem}
\label{ll:paper:not-winwittagree-two-outer-left-strip-of-zero-diag}
For every natural number \(t\) and natural number \(k\) and natural number \(p\) and natural number \(a\) and natural number \(L\), if \(p \le t\), then if \(t - p + 2 \le k\), then if \(\operatorname{zerodiagonal}(t - p)\), then if \(\operatorname{settledby}(t - p, t + \operatorname{twopow}(k))\), then if \(\operatorname{settledby}(t - p + 2, t + 2 + \operatorname{twopow}(k))\), then if \(a \le t + p + \operatorname{twopow}(k + 1)\), then if \(t + p + \operatorname{twopow}(k + 1) \le a + L\), then if \(\operatorname{lhcenter}(t + 2, p) = \operatorname{lhcenter}(t, p)\), then not \(\operatorname{winwittagree}(t, k, 2, a, L)\).
\end{theorem}
\begin{proof}
The goal follows from \(\operatorname{notwinwittagreetwoouterleftstripofzerodiag}(t, k, p, a, L)\).
\end{proof}
\begin{theorem}
\label{ll:paper:zcenter-shift-xor-eq-center-xor-of-winwittagree-sector-two-center-slot}
For every natural number \(t\) and natural number \(k\) and natural number \(Q\) and natural number \(a\) and natural number \(L\), if \(t + Q \le k\), then if \(\operatorname{settledby}(t, t + \operatorname{twopow}(k))\), then if \(\operatorname{settledby}(t + Q, t + Q + \operatorname{twopow}(k))\), then if \(a \le t + \operatorname{twopow}(k + 1)\), then if \(t + \operatorname{twopow}(k + 1) \le a + L\), then if \(\operatorname{winwittagree}(t, k, Q, a, L)\), then \(\operatorname{zcenterxoreq}(t + Q, t, \operatorname{bxor}(\operatorname{center}(t + Q), \operatorname{center}(t)))\).
\end{theorem}
\begin{proof}
The goal follows from \(\operatorname{zcentershiftxoreqcenterxorofwinwittagreesectortwocenterslot}(t, k, Q, a, L)\).
\end{proof}
\begin{theorem}
\label{ll:paper:zcenter-shift-eq-of-winwittagree-sector-two-center-slot}
For every natural number \(t\) and natural number \(k\) and natural number \(Q\) and natural number \(a\) and natural number \(L\), if \(t + Q \le k\), then if \(\operatorname{settledby}(t, t + \operatorname{twopow}(k))\), then if \(\operatorname{settledby}(t + Q, t + Q + \operatorname{twopow}(k))\), then if \(a \le t + \operatorname{twopow}(k + 1)\), then if \(t + \operatorname{twopow}(k + 1) \le a + L\), then if \(\operatorname{center}(t + Q) = \operatorname{center}(t)\), then if \(\operatorname{winwittagree}(t, k, Q, a, L)\), then \(\operatorname{zcenterxoreq}(t + Q, t, \operatorname{false})\).
\end{theorem}
\begin{proof}
The goal follows from \(\operatorname{zcentershifteqofwinwittagreesectortwocenterslot}(t, k, Q, a, L)\).
\end{proof}
\begin{theorem}
\label{ll:paper:periodic-zcenter-shift-eq-of-winwittagree-sector-two-center-slot}
For every natural number \(T\) and natural number \(P\), if \(0 < P\), then if \(\operatorname{centerperiodic}(T, P)\), then for every natural number \(m\) and natural number \(t\) and natural number \(k\) and natural number \(a\) and natural number \(L\), if \(T \le t\), then if \(t + m \cdot P \le k\), then if \(\operatorname{settledby}(t, t + \operatorname{twopow}(k))\), then if \(\operatorname{settledby}(t + m \cdot P, t + m \cdot P + \operatorname{twopow}(k))\), then if \(a \le t + \operatorname{twopow}(k + 1)\), then if \(t + \operatorname{twopow}(k + 1) \le a + L\), then if \(\operatorname{winwittagree}(t, k, m \cdot P, a, L)\), then \(\operatorname{zcenterxoreq}(t + m \cdot P, t, \operatorname{false})\).
\end{theorem}
\begin{proof}
Assume \(hP\), \(hper\).
The goal follows from \(\operatorname{periodiczcentershifteqofwinwittagreesectortwocenterslot}(T, P, hP, hper)\).
\end{proof}
\begin{theorem}
\label{ll:paper:not-winwittagree-two-sector-two-center-slot-of-zero-diag}
For every natural number \(z\) and natural number \(k\) and natural number \(a\) and natural number \(L\), if \(z + 2 \le k\), then if \(\operatorname{zerodiagonal}(z)\), then if \(\operatorname{settledby}(z, z + \operatorname{twopow}(k))\), then if \(\operatorname{settledby}(z + 2, z + 2 + \operatorname{twopow}(k))\), then if \(a \le z + \operatorname{twopow}(k + 1)\), then if \(z + \operatorname{twopow}(k + 1) \le a + L\), then if \(\operatorname{center}(z + 2) = \operatorname{center}(z)\), then not \(\operatorname{winwittagree}(z, k, 2, a, L)\).
\end{theorem}
\begin{proof}
The goal follows from \(\operatorname{notwinwittagreetwosectortwocenterslotofzerodiag}(z, k, a, L)\).
\end{proof}
\begin{theorem}
\label{ll:paper:not-periodic-winwittagree-sector-two-center-slot-of-zero-and-after-zero-diag}
For every natural number \(T\) and natural number \(P\), if \(0 < P\), then if \(\operatorname{centerperiodic}(T, P)\), then for every natural number \(m\) and natural number \(z\_1\) and natural number \(z\_2\) and natural number \(k\) and natural number \(a\) and natural number \(L\), if \(T \le z\_1\), then if \(z\_1 + m \cdot P \le k\), then if \(\operatorname{zerodiagonal}(z\_1)\), then if \(\operatorname{zerodiagonal}(z\_2)\), then if \(z\_1 + m \cdot P = z\_2 + 2\), then if \(\operatorname{settledby}(z\_1, z\_1 + \operatorname{twopow}(k))\), then if \(\operatorname{settledby}(z\_2 + 2, z\_2 + 2 + \operatorname{twopow}(k))\), then if \(a \le z\_1 + \operatorname{twopow}(k + 1)\), then if \(z\_1 + \operatorname{twopow}(k + 1) \le a + L\), then not \(\operatorname{winwittagree}(z\_1, k, m \cdot P, a, L)\).
\end{theorem}
\begin{proof}
Assume \(hP\), \(hper\).
The goal follows from \(\operatorname{notperiodicwinwittagreesectortwocenterslotofzeroandafterzerodiag}(T, P, hP, hper)\).
\end{proof}
\begin{theorem}
\label{ll:paper:winwitt-scale-cross-zero-gap-late-shift-flip-of-doubles-and-winwittagree}
For every natural number \(Q\) and natural number \(t\) and natural number \(k\) and natural number \(r\) and natural number \(a\) and natural number \(L\), if \(t + Q \le r\), then if \(t + Q + r + 1 \le k\), then if \(\operatorname{zerodiagonal}(t + Q + r)\), then if \(\operatorname{doublesat}(Q, t + Q + r)\), then if \(\operatorname{settledby}(t + Q + r + 1, t + \operatorname{twopow}(k))\), then if \(a \le t + \operatorname{twopow}(k + 1) - (r + Q + 1)\), then if \(t + \operatorname{twopow}(k + 1) - (r + Q + 1) \le a + L\), then if \(\operatorname{winwittagree}(t, k, Q, a, L)\), then \(\operatorname{wittbitbool}(\operatorname{wittwin}(t + Q, \operatorname{twopow}(k)), t + Q + \operatorname{twopow}(k + 1) - (r + 1)) = \operatorname{not}(\operatorname{wittbitbool}(\operatorname{wittwin}(t + Q, \operatorname{twopow}(k)), t + \operatorname{twopow}(k + 1) - (r + 1)))\).
\end{theorem}
\begin{proof}
The goal follows from \(\operatorname{winwittscalecrosszerogaplateshiftflipofdoublesandwinwittagree}(Q, t, k, r, a, L)\).
\end{proof}
\begin{theorem}
\label{ll:paper:lhcenter-late-cross-late-shift-flip-of-doubles-and-winwittagree-zero-gap}
For every natural number \(Q\) and natural number \(t\) and natural number \(k\) and natural number \(r\) and natural number \(a\) and natural number \(L\), if \(t + Q \le r\), then if \(t + Q + r + 1 \le k\), then if \(\operatorname{zerodiagonal}(t + Q + r)\), then if \(\operatorname{doublesat}(Q, t + Q + r)\), then if \(\operatorname{settledby}(t + Q + r + 1, t + \operatorname{twopow}(k))\), then if \(a \le t + \operatorname{twopow}(k + 1) - (r + Q + 1)\), then if \(t + \operatorname{twopow}(k + 1) - (r + Q + 1) \le a + L\), then if \(\operatorname{winwittagree}(t, k, Q, a, L)\), then \(\operatorname{lhcenter}(t + Q + \operatorname{twopow}(k), \operatorname{twopow}(k) - (r + 1)) = \operatorname{not}(\operatorname{lhcenter}(t + Q + \operatorname{twopow}(k), \operatorname{twopow}(k) - (r + Q + 1)))\).
\end{theorem}
\begin{proof}
The goal follows from \(\operatorname{lhcenterlatecrosslateshiftflipofdoublesandwinwittagreezerogap}(Q, t, k, r, a, L)\).
\end{proof}
\begin{theorem}
\label{ll:paper:winwitt-scale-cross-zero-gap-same-bit-flip-of-doubles-and-winwittagree}
For every natural number \(Q\) and natural number \(t\) and natural number \(k\) and natural number \(r\) and natural number \(a\) and natural number \(L\), if \(t + Q \le r\), then if \(t + Q + r + 1 \le k\), then if \(\operatorname{zerodiagonal}(t + Q + r)\), then if \(\operatorname{doublesat}(Q, t + Q + r)\), then if \(\operatorname{settledby}(t + Q + r + 1, t + \operatorname{twopow}(k))\), then if \(a \le t + \operatorname{twopow}(k + 1) - (r + Q + 1)\), then if \(t + \operatorname{twopow}(k + 1) - (r + 1) \le a + L\), then if \(\operatorname{winwittagree}(t, k, Q, a, L)\), then \(\operatorname{wittbitbool}(\operatorname{wittwin}(t + Q, \operatorname{twopow}(k)), t + \operatorname{twopow}(k + 1) - (r + 1)) = \operatorname{not}(\operatorname{wittbitbool}(\operatorname{wittwin}(t, \operatorname{twopow}(k)), t + \operatorname{twopow}(k + 1) - (r + 1)))\).
\end{theorem}
\begin{proof}
The goal follows from \(\operatorname{winwittscalecrosszerogapsamebitflipofdoublesandwinwittagree}(Q, t, k, r, a, L)\).
\end{proof}
\begin{theorem}
\label{ll:paper:winwitt-scale-cross-zero-gap-reset-shift-flip-of-doubles-and-winwittagree}
For every natural number \(Q\) and natural number \(t\) and natural number \(k\) and natural number \(r\) and natural number \(a\) and natural number \(L\), if \(t + Q \le r\), then if \(t + Q + r + 2 \le k\), then if \(\operatorname{zerodiagonal}(t + Q + r)\), then if \(\operatorname{doublesat}(Q, t + Q + r)\), then if \(\operatorname{settledby}(t + Q + r + 1, t + \operatorname{twopow}(k))\), then if \(\operatorname{settledby}(t + Q + r + 2, t + \operatorname{twopow}(k))\), then if \(a \le t + \operatorname{twopow}(k + 1) - (r + 1)\), then if \(t + \operatorname{twopow}(k + 1) - (r + 1) \le a + L\), then if \(\operatorname{winwittagree}(t + 1, k, Q, a, L)\), then \(\operatorname{wittbitbool}(\operatorname{wittwin}(t + 1, \operatorname{twopow}(k)), t + \operatorname{twopow}(k + 1) - (r + 1)) = \operatorname{not}(\operatorname{wittbitbool}(\operatorname{wittwin}(t + 1, \operatorname{twopow}(k)), t + \operatorname{twopow}(k + 1) - (r + Q + 1)))\).
\end{theorem}
\begin{proof}
The goal follows from \(\operatorname{winwittscalecrosszerogapresetshiftflipofdoublesandwinwittagree}(Q, t, k, r, a, L)\).
\end{proof}
\begin{theorem}
\label{ll:paper:lhcenter-late-cross-reset-shift-flip-of-doubles-and-winwittagree-zero-gap}
For every natural number \(Q\) and natural number \(t\) and natural number \(k\) and natural number \(r\) and natural number \(a\) and natural number \(L\), if \(t + Q \le r\), then if \(t + Q + r + 2 \le k\), then if \(\operatorname{zerodiagonal}(t + Q + r)\), then if \(\operatorname{doublesat}(Q, t + Q + r)\), then if \(\operatorname{settledby}(t + Q + r + 1, t + \operatorname{twopow}(k))\), then if \(\operatorname{settledby}(t + Q + r + 2, t + \operatorname{twopow}(k))\), then if \(a \le t + \operatorname{twopow}(k + 1) - (r + 1)\), then if \(t + \operatorname{twopow}(k + 1) - (r + 1) \le a + L\), then if \(\operatorname{winwittagree}(t + 1, k, Q, a, L)\), then \(\operatorname{lhcenter}(t + 1 + \operatorname{twopow}(k), \operatorname{twopow}(k) - (r + 2)) = \operatorname{not}(\operatorname{lhcenter}(t + 1 + \operatorname{twopow}(k), \operatorname{twopow}(k) - (r + Q + 2)))\).
\end{theorem}
\begin{proof}
The goal follows from \(\operatorname{lhcenterlatecrossresetshiftflipofdoublesandwinwittagreezerogap}(Q, t, k, r, a, L)\).
\end{proof}
\begin{theorem}
\label{ll:paper:winwitt-scale-cross-zero-gap-shift-eq-step-one-of-doubles-and-winwittagree}
For every natural number \(Q\) and natural number \(t\) and natural number \(k\) and natural number \(r\) and natural number \(a\) and natural number \(L\), if \(t + Q \le r\), then if \(t + Q + r + 2 \le k\), then if \(\operatorname{zerodiagonal}(t + Q + r)\), then if \(\operatorname{doublesat}(Q, t + Q + r)\), then if \(\operatorname{settledby}(t + Q + r + 1, t + \operatorname{twopow}(k))\), then if \(\operatorname{settledby}(t + Q + r + 2, t + \operatorname{twopow}(k))\), then if \(a \le t + \operatorname{twopow}(k + 1) - (r + 1)\), then if \(t + \operatorname{twopow}(k + 1) - (r + 1) \le a + L\), then if \(\operatorname{winwittagree}(t, k, Q, a, L)\), then \(\operatorname{wittbitbool}(\operatorname{wittwin}(t, \operatorname{twopow}(k)), t + \operatorname{twopow}(k + 1) - (r + 1)) = \operatorname{wittbitbool}(\operatorname{wittwin}(t + 1, \operatorname{twopow}(k)), t + \operatorname{twopow}(k + 1) - (r + Q + 1))\).
\end{theorem}
\begin{proof}
The goal follows from \(\operatorname{winwittscalecrosszerogapshifteqsteponeofdoublesandwinwittagree}(Q, t, k, r, a, L)\).
\end{proof}
\begin{theorem}
\label{ll:paper:winwitt-scale-cross-zero-gap-step-flip-at-shift-of-doubles-and-winwittagree}
For every natural number \(Q\) and natural number \(t\) and natural number \(k\) and natural number \(r\) and natural number \(a\) and natural number \(L\), if \(t + Q \le r\), then if \(t + Q + r + 2 \le k\), then if \(\operatorname{zerodiagonal}(t + Q + r)\), then if \(\operatorname{doublesat}(Q, t + Q + r)\), then if \(\operatorname{settledby}(t + Q + r + 1, t + \operatorname{twopow}(k))\), then if \(\operatorname{settledby}(t + Q + r + 2, t + \operatorname{twopow}(k))\), then if \(a \le t + \operatorname{twopow}(k + 1) - (r + 1)\), then if \(t + \operatorname{twopow}(k + 1) - (r + 1) \le a + L\), then if \(\operatorname{winwittagree}(t, k, Q, a, L)\), then if \(\operatorname{winwittagree}(t + 1, k, Q, a, L)\), then \(\operatorname{wittbitbool}(\operatorname{wittwin}(t + 1, \operatorname{twopow}(k)), t + \operatorname{twopow}(k + 1) - (r + 1)) = \operatorname{not}(\operatorname{wittbitbool}(\operatorname{wittwin}(t, \operatorname{twopow}(k)), t + \operatorname{twopow}(k + 1) - (r + 1)))\).
\end{theorem}
\begin{proof}
The goal follows from \(\operatorname{winwittscalecrosszerogapstepflipatshiftofdoublesandwinwittagree}(Q, t, k, r, a, L)\).
\end{proof}
\begin{theorem}
\label{ll:paper:lhcenter-late-cross-step-flip-at-shift-of-doubles-and-winwittagree-zero-gap}
For every natural number \(Q\) and natural number \(t\) and natural number \(k\) and natural number \(r\) and natural number \(a\) and natural number \(L\), if \(t + Q \le r\), then if \(t + Q + r + 2 \le k\), then if \(\operatorname{zerodiagonal}(t + Q + r)\), then if \(\operatorname{doublesat}(Q, t + Q + r)\), then if \(\operatorname{settledby}(t + Q + r + 1, t + \operatorname{twopow}(k))\), then if \(\operatorname{settledby}(t + Q + r + 2, t + \operatorname{twopow}(k))\), then if \(a \le t + \operatorname{twopow}(k + 1) - (r + 1)\), then if \(t + \operatorname{twopow}(k + 1) - (r + 1) \le a + L\), then if \(\operatorname{winwittagree}(t, k, Q, a, L)\), then if \(\operatorname{winwittagree}(t + 1, k, Q, a, L)\), then \(\operatorname{lhcenter}(t + 1 + \operatorname{twopow}(k), \operatorname{twopow}(k) - (r + 2)) = \operatorname{not}(\operatorname{lhcenter}(t + \operatorname{twopow}(k), \operatorname{twopow}(k) - (r + 1)))\).
\end{theorem}
\begin{proof}
The goal follows from \(\operatorname{lhcenterlatecrossstepflipatshiftofdoublesandwinwittagreezerogap}(Q, t, k, r, a, L)\).
\end{proof}
\begin{theorem}
\label{ll:paper:lhcenter-late-cross-step}
For every natural number \(t\) and natural number \(k\) and natural number \(r\), if \(r + 3 \le \operatorname{twopow}(k)\), then \(\operatorname{lhcenter}(t + 1 + \operatorname{twopow}(k), \operatorname{twopow}(k) - (r + 2)) = \operatorname{bxor}(\operatorname{lhcenter}(t + \operatorname{twopow}(k), \operatorname{twopow}(k) - (r + 1)), \operatorname{lhcenter}(t + \operatorname{twopow}(k), \operatorname{twopow}(k) - (r + 2)) \operatorname{or} \operatorname{lhcenter}(t + \operatorname{twopow}(k), \operatorname{twopow}(k) - (r + 3)))\).
\end{theorem}
\begin{proof}
The goal follows from \(\operatorname{lhcenterlatecrossstep}(t, k, r)\).
\end{proof}
\begin{theorem}
\label{ll:paper:lhcenter-late-cross-or-eq-true-of-step-flip}
For every natural number \(t\) and natural number \(k\) and natural number \(r\), if \(r + 3 \le \operatorname{twopow}(k)\), then if \(\operatorname{lhcenter}(t + 1 + \operatorname{twopow}(k), \operatorname{twopow}(k) - (r + 2)) = \operatorname{not}(\operatorname{lhcenter}(t + \operatorname{twopow}(k), \operatorname{twopow}(k) - (r + 1)))\), then \((\operatorname{lhcenter}(t + \operatorname{twopow}(k), \operatorname{twopow}(k) - (r + 2)) \operatorname{or} \operatorname{lhcenter}(t + \operatorname{twopow}(k), \operatorname{twopow}(k) - (r + 3))) = \operatorname{true}\).
\end{theorem}
\begin{proof}
The goal follows from \(\operatorname{lhcenterlatecrossoreqtrueofstepflip}(t, k, r)\).
\end{proof}
\begin{theorem}
\label{ll:paper:lhcenter-late-cross-or-eq-true-at-shift-of-doubles-and-winwittagree-zero-gap}
For every natural number \(Q\) and natural number \(t\) and natural number \(k\) and natural number \(r\) and natural number \(a\) and natural number \(L\), if \(t + Q \le r\), then if \(t + Q + r + 2 \le k\), then if \(\operatorname{zerodiagonal}(t + Q + r)\), then if \(\operatorname{doublesat}(Q, t + Q + r)\), then if \(\operatorname{settledby}(t + Q + r + 1, t + \operatorname{twopow}(k))\), then if \(\operatorname{settledby}(t + Q + r + 2, t + \operatorname{twopow}(k))\), then if \(a \le t + \operatorname{twopow}(k + 1) - (r + 1)\), then if \(t + \operatorname{twopow}(k + 1) - (r + 1) \le a + L\), then if \(\operatorname{winwittagree}(t, k, Q, a, L)\), then if \(\operatorname{winwittagree}(t + 1, k, Q, a, L)\), then \((\operatorname{lhcenter}(t + \operatorname{twopow}(k), \operatorname{twopow}(k) - (r + 2)) \operatorname{or} \operatorname{lhcenter}(t + \operatorname{twopow}(k), \operatorname{twopow}(k) - (r + 3))) = \operatorname{true}\).
\end{theorem}
\begin{proof}
The goal follows from \(\operatorname{lhcenterlatecrossoreqtrueatshiftofdoublesandwinwittagreezerogap}(Q, t, k, r, a, L)\).
\end{proof}
\begin{theorem}
\label{ll:paper:winwitt-scale-cross-zero-gap-or-eq-true-at-shift-of-doubles-and-winwittagree}
For every natural number \(Q\) and natural number \(t\) and natural number \(k\) and natural number \(r\) and natural number \(a\) and natural number \(L\), if \(t + Q \le r\), then if \(t + Q + r + 2 \le k\), then if \(\operatorname{zerodiagonal}(t + Q + r)\), then if \(\operatorname{doublesat}(Q, t + Q + r)\), then if \(\operatorname{settledby}(t + Q + r + 1, t + \operatorname{twopow}(k))\), then if \(\operatorname{settledby}(t + Q + r + 2, t + \operatorname{twopow}(k))\), then if \(a \le t + \operatorname{twopow}(k + 1) - (r + 1)\), then if \(t + \operatorname{twopow}(k + 1) - (r + 1) \le a + L\), then if \(\operatorname{winwittagree}(t, k, Q, a, L)\), then if \(\operatorname{winwittagree}(t + 1, k, Q, a, L)\), then \((\operatorname{wittbitbool}(\operatorname{wittwin}(t, \operatorname{twopow}(k)), t + \operatorname{twopow}(k + 1) - (r + 2)) \operatorname{or} \operatorname{wittbitbool}(\operatorname{wittwin}(t, \operatorname{twopow}(k)), t + \operatorname{twopow}(k + 1) - (r + 3))) = \operatorname{true}\).
\end{theorem}
\begin{proof}
The goal follows from \(\operatorname{winwittscalecrosszerogaporeqtrueatshiftofdoublesandwinwittagree}(Q, t, k, r, a, L)\).
\end{proof}
\begin{theorem}
\label{ll:paper:lhcenter-late-cross-pin-r-plus-three-of-doubles-and-winwittagree-and-zero-diag-two}
For every natural number \(Q\) and natural number \(t\) and natural number \(k\) and natural number \(r\) and natural number \(a\) and natural number \(L\), if \(t + Q \le r\), then if \(t + Q + r + 2 \le k\), then if \(\operatorname{zerodiagonal}(t + Q + r)\), then if \(\operatorname{doublesat}(Q, t + Q + r)\), then if \(\operatorname{zerodiagonal}(t + r + 2)\), then if \(\operatorname{settledby}(t + r + 2, t + \operatorname{twopow}(k))\), then if \(\operatorname{settledby}(t + Q + r + 1, t + \operatorname{twopow}(k))\), then if \(\operatorname{settledby}(t + Q + r + 2, t + \operatorname{twopow}(k))\), then if \(a \le t + \operatorname{twopow}(k + 1) - (r + 1)\), then if \(t + \operatorname{twopow}(k + 1) - (r + 1) \le a + L\), then if \(\operatorname{winwittagree}(t, k, Q, a, L)\), then if \(\operatorname{winwittagree}(t + 1, k, Q, a, L)\), then \(\operatorname{lhcenter}(t + \operatorname{twopow}(k), \operatorname{twopow}(k) - (r + 3)) = \operatorname{true}\).
\end{theorem}
\begin{proof}
The goal follows from \(\operatorname{lhcenterlatecrosspinrplusthreeofdoublesandwinwittagreeandzerodiagtwo}(Q, t, k, r, a, L)\).
\end{proof}
\begin{theorem}
\label{ll:paper:lhcenter-late-cross-pin-r-plus-two-of-doubles-and-winwittagree-and-zero-diag-three}
For every natural number \(Q\) and natural number \(t\) and natural number \(k\) and natural number \(r\) and natural number \(a\) and natural number \(L\), if \(1 \le Q\), then if \(t + Q \le r\), then if \(t + Q + r + 2 \le k\), then if \(\operatorname{zerodiagonal}(t + Q + r)\), then if \(\operatorname{doublesat}(Q, t + Q + r)\), then if \(\operatorname{zerodiagonal}(t + r + 3)\), then if \(\operatorname{settledby}(t + r + 3, t + \operatorname{twopow}(k))\), then if \(\operatorname{settledby}(t + Q + r + 1, t + \operatorname{twopow}(k))\), then if \(\operatorname{settledby}(t + Q + r + 2, t + \operatorname{twopow}(k))\), then if \(a \le t + \operatorname{twopow}(k + 1) - (r + 1)\), then if \(t + \operatorname{twopow}(k + 1) - (r + 1) \le a + L\), then if \(\operatorname{winwittagree}(t, k, Q, a, L)\), then if \(\operatorname{winwittagree}(t + 1, k, Q, a, L)\), then \(\operatorname{lhcenter}(t + \operatorname{twopow}(k), \operatorname{twopow}(k) - (r + 2)) = \operatorname{true}\).
\end{theorem}
\begin{proof}
The goal follows from \(\operatorname{lhcenterlatecrosspinrplustwoofdoublesandwinwittagreeandzerodiagthree}(Q, t, k, r, a, L)\).
\end{proof}
\begin{theorem}
\label{ll:paper:lhcenter-late-cross-pin-r-plus-one-of-doubles-and-winwittagree-and-zero-diag-three}
For every natural number \(Q\) and natural number \(t\) and natural number \(k\) and natural number \(r\) and natural number \(a\) and natural number \(L\), if \(1 \le Q\), then if \(t + Q \le r\), then if \(t + Q + r + 2 \le k\), then if \(\operatorname{zerodiagonal}(t + Q + r)\), then if \(\operatorname{doublesat}(Q, t + Q + r)\), then if \(\operatorname{zerodiagonal}(t + r + 3)\), then if \(\operatorname{settledby}(t + r + 3, t + \operatorname{twopow}(k))\), then if \(\operatorname{settledby}(t + Q + r + 1, t + \operatorname{twopow}(k))\), then if \(\operatorname{settledby}(t + Q + r + 2, t + \operatorname{twopow}(k))\), then if \(a \le t + \operatorname{twopow}(k + 1) - (r + 1)\), then if \(t + \operatorname{twopow}(k + 1) - (r + 1) \le a + L\), then if \(\operatorname{winwittagree}(t, k, Q, a, L)\), then if \(\operatorname{winwittagree}(t + 1, k, Q, a, L)\), then \(\operatorname{lhcenter}(t + \operatorname{twopow}(k), \operatorname{twopow}(k) - (r + 1)) = \operatorname{true}\).
\end{theorem}
\begin{proof}
The goal follows from \(\operatorname{lhcenterlatecrosspinrplusoneofdoublesandwinwittagreeandzerodiagthree}(Q, t, k, r, a, L)\).
\end{proof}
\begin{theorem}
\label{ll:paper:lhcenter-late-cross-pin-doubling-diag-false-of-doubles-and-winwittagree-and-zero-diag-three}
For every natural number \(Q\) and natural number \(t\) and natural number \(k\) and natural number \(r\) and natural number \(a\) and natural number \(L\), if \(1 \le Q\), then if \(t + Q \le r\), then if \(t + Q + r + 2 \le k\), then if \(\operatorname{zerodiagonal}(t + Q + r)\), then if \(\operatorname{doublesat}(Q, t + Q + r)\), then if \(\operatorname{zerodiagonal}(t + r + 3)\), then if \(\operatorname{settledby}(t + r + 3, t + \operatorname{twopow}(k))\), then if \(\operatorname{settledby}(t + Q + r + 1, t + \operatorname{twopow}(k))\), then if \(\operatorname{settledby}(t + Q + r + 2, t + \operatorname{twopow}(k))\), then if \(a \le t + \operatorname{twopow}(k + 1) - (r + 1)\), then if \(t + \operatorname{twopow}(k + 1) - (r + 1) \le a + L\), then if \(\operatorname{winwittagree}(t, k, Q, a, L)\), then if \(\operatorname{winwittagree}(t + 1, k, Q, a, L)\), then \(\operatorname{lhcenter}(t + \operatorname{twopow}(k), \operatorname{twopow}(k) - (r + Q + 1)) = \operatorname{false}\).
\end{theorem}
\begin{proof}
The goal follows from \(\operatorname{lhcenterlatecrosspindoublingdiagfalseofdoublesandwinwittagreeandzerodiagthree}(Q, t, k, r, a, L)\).
\end{proof}
\begin{theorem}
\label{ll:paper:lhcenter-late-cross-pin-before-zero-diag-of-doubles-three-and-winwittagree}
For every natural number \(t\) and natural number \(k\) and natural number \(r\) and natural number \(a\) and natural number \(L\), if \(t + 3 \le r\), then if \(t + 3 + r + 2 \le k\), then if \(\operatorname{zerodiagonal}(t + 3 + r)\), then if \(\operatorname{doublesat}(3, t + 3 + r)\), then if \(\operatorname{settledby}(t + 3 + r, t + \operatorname{twopow}(k))\), then if \(\operatorname{settledby}(t + 3 + r + 1, t + \operatorname{twopow}(k))\), then if \(\operatorname{settledby}(t + 3 + r + 2, t + \operatorname{twopow}(k))\), then if \(a \le t + \operatorname{twopow}(k + 1) - (r + 1)\), then if \(t + \operatorname{twopow}(k + 1) - (r + 1) \le a + L\), then if \(\operatorname{winwittagree}(t, k, 3, a, L)\), then if \(\operatorname{winwittagree}(t + 1, k, 3, a, L)\), then \(\operatorname{lhcenter}(t + \operatorname{twopow}(k), \operatorname{twopow}(k) - (r + 1)) = \operatorname{true}\).
\end{theorem}
\begin{proof}
The goal follows from \(\operatorname{lhcenterlatecrosspinbeforezerodiagofdoublesthreeandwinwittagree}(t, k, r, a, L)\).
\end{proof}
\begin{theorem}
\label{ll:paper:not-winwittagree-step-pair-at-shift-of-doubles-and-both-false}
For every natural number \(Q\) and natural number \(t\) and natural number \(k\) and natural number \(r\) and natural number \(a\) and natural number \(L\), if \(t + Q \le r\), then if \(t + Q + r + 2 \le k\), then if \(\operatorname{zerodiagonal}(t + Q + r)\), then if \(\operatorname{doublesat}(Q, t + Q + r)\), then if \(\operatorname{settledby}(t + Q + r + 1, t + \operatorname{twopow}(k))\), then if \(\operatorname{settledby}(t + Q + r + 2, t + \operatorname{twopow}(k))\), then if \(a \le t + \operatorname{twopow}(k + 1) - (r + 1)\), then if \(t + \operatorname{twopow}(k + 1) - (r + 1) \le a + L\), then if \(\operatorname{lhcenter}(t + \operatorname{twopow}(k), \operatorname{twopow}(k) - (r + 2)) = \operatorname{false}\), then if \(\operatorname{lhcenter}(t + \operatorname{twopow}(k), \operatorname{twopow}(k) - (r + 3)) = \operatorname{false}\), then if \(\operatorname{winwittagree}(t, k, Q, a, L)\), then not \(\operatorname{winwittagree}(t + 1, k, Q, a, L)\).
\end{theorem}
\begin{proof}
The goal follows from \(\operatorname{notwinwittagreesteppairatshiftofdoublesandbothfalse}(Q, t, k, r, a, L)\).
\end{proof}
\begin{theorem}
\label{ll:paper:winwitt-scale-cross-zero-gap-reset-late-shift-flip-of-doubles-and-winwittagree}
For every natural number \(Q\) and natural number \(t\) and natural number \(k\) and natural number \(r\) and natural number \(a\) and natural number \(L\), if \(t + Q \le r\), then if \(t + Q + r + 2 \le k\), then if \(\operatorname{zerodiagonal}(t + Q + r)\), then if \(\operatorname{doublesat}(Q, t + Q + r)\), then if \(\operatorname{settledby}(t + Q + r + 1, t + \operatorname{twopow}(k))\), then if \(\operatorname{settledby}(t + Q + r + 2, t + \operatorname{twopow}(k))\), then if \(a \le t + \operatorname{twopow}(k + 1) - (r + Q + 1)\), then if \(t + \operatorname{twopow}(k + 1) - (r + Q + 1) \le a + L\), then if \(\operatorname{winwittagree}(t + 1, k, Q, a, L)\), then \(\operatorname{wittbitbool}(\operatorname{wittwin}(t + 1 + Q, \operatorname{twopow}(k)), t + Q + \operatorname{twopow}(k + 1) - (r + 1)) = \operatorname{not}(\operatorname{wittbitbool}(\operatorname{wittwin}(t + 1 + Q, \operatorname{twopow}(k)), t + \operatorname{twopow}(k + 1) - (r + 1)))\).
\end{theorem}
\begin{proof}
The goal follows from \(\operatorname{winwittscalecrosszerogapresetlateshiftflipofdoublesandwinwittagree}(Q, t, k, r, a, L)\).
\end{proof}
\begin{theorem}
\label{ll:paper:lhcenter-late-cross-reset-late-shift-flip-of-doubles-and-winwittagree-zero-gap}
For every natural number \(Q\) and natural number \(t\) and natural number \(k\) and natural number \(r\) and natural number \(a\) and natural number \(L\), if \(t + Q \le r\), then if \(t + Q + r + 2 \le k\), then if \(\operatorname{zerodiagonal}(t + Q + r)\), then if \(\operatorname{doublesat}(Q, t + Q + r)\), then if \(\operatorname{settledby}(t + Q + r + 1, t + \operatorname{twopow}(k))\), then if \(\operatorname{settledby}(t + Q + r + 2, t + \operatorname{twopow}(k))\), then if \(a \le t + \operatorname{twopow}(k + 1) - (r + Q + 1)\), then if \(t + \operatorname{twopow}(k + 1) - (r + Q + 1) \le a + L\), then if \(\operatorname{winwittagree}(t + 1, k, Q, a, L)\), then \(\operatorname{lhcenter}(t + 1 + Q + \operatorname{twopow}(k), \operatorname{twopow}(k) - (r + 2)) = \operatorname{not}(\operatorname{lhcenter}(t + 1 + Q + \operatorname{twopow}(k), \operatorname{twopow}(k) - (r + Q + 2)))\).
\end{theorem}
\begin{proof}
The goal follows from \(\operatorname{lhcenterlatecrossresetlateshiftflipofdoublesandwinwittagreezerogap}(Q, t, k, r, a, L)\).
\end{proof}
\begin{theorem}
\label{ll:paper:winwitt-scale-cross-zero-gap-reset-same-bit-flip-of-doubles-and-winwittagree}
For every natural number \(Q\) and natural number \(t\) and natural number \(k\) and natural number \(r\) and natural number \(a\) and natural number \(L\), if \(t + Q \le r\), then if \(t + Q + r + 2 \le k\), then if \(\operatorname{zerodiagonal}(t + Q + r)\), then if \(\operatorname{doublesat}(Q, t + Q + r)\), then if \(\operatorname{settledby}(t + Q + r + 1, t + \operatorname{twopow}(k))\), then if \(\operatorname{settledby}(t + Q + r + 2, t + \operatorname{twopow}(k))\), then if \(a \le t + \operatorname{twopow}(k + 1) - (r + Q + 1)\), then if \(t + \operatorname{twopow}(k + 1) - (r + 1) \le a + L\), then if \(\operatorname{winwittagree}(t + 1, k, Q, a, L)\), then \(\operatorname{wittbitbool}(\operatorname{wittwin}(t + 1 + Q, \operatorname{twopow}(k)), t + \operatorname{twopow}(k + 1) - (r + 1)) = \operatorname{not}(\operatorname{wittbitbool}(\operatorname{wittwin}(t + 1, \operatorname{twopow}(k)), t + \operatorname{twopow}(k + 1) - (r + 1)))\).
\end{theorem}
\begin{proof}
The goal follows from \(\operatorname{winwittscalecrosszerogapresetsamebitflipofdoublesandwinwittagree}(Q, t, k, r, a, L)\).
\end{proof}
\begin{theorem}
\label{ll:paper:winwitt-scale-cross-zero-gap-late-step-flip-at-shift-of-doubles-and-winwittagree}
For every natural number \(Q\) and natural number \(t\) and natural number \(k\) and natural number \(r\) and natural number \(a\) and natural number \(L\), if \(t + Q \le r\), then if \(t + Q + r + 2 \le k\), then if \(\operatorname{zerodiagonal}(t + Q + r)\), then if \(\operatorname{doublesat}(Q, t + Q + r)\), then if \(\operatorname{settledby}(t + Q + r + 1, t + \operatorname{twopow}(k))\), then if \(\operatorname{settledby}(t + Q + r + 2, t + \operatorname{twopow}(k))\), then if \(a \le t + \operatorname{twopow}(k + 1) - (r + Q + 1)\), then if \(t + \operatorname{twopow}(k + 1) - (r + Q + 1) \le a + L\), then if \(\operatorname{winwittagree}(t, k, Q, a, L)\), then if \(\operatorname{winwittagree}(t + 1, k, Q, a, L)\), then \(\operatorname{wittbitbool}(\operatorname{wittwin}(t + 1 + Q, \operatorname{twopow}(k)), t + \operatorname{twopow}(k + 1) - (r + 1)) = \operatorname{not}(\operatorname{wittbitbool}(\operatorname{wittwin}(t + Q, \operatorname{twopow}(k)), t + \operatorname{twopow}(k + 1) - (r + 1)))\).
\end{theorem}
\begin{proof}
The goal follows from \(\operatorname{winwittscalecrosszerogaplatestepflipatshiftofdoublesandwinwittagree}(Q, t, k, r, a, L)\).
\end{proof}
\begin{theorem}
\label{ll:paper:lhcenter-late-cross-late-step-flip-at-shift-of-doubles-and-winwittagree-zero-gap}
For every natural number \(Q\) and natural number \(t\) and natural number \(k\) and natural number \(r\) and natural number \(a\) and natural number \(L\), if \(t + Q \le r\), then if \(t + Q + r + 2 \le k\), then if \(\operatorname{zerodiagonal}(t + Q + r)\), then if \(\operatorname{doublesat}(Q, t + Q + r)\), then if \(\operatorname{settledby}(t + Q + r + 1, t + \operatorname{twopow}(k))\), then if \(\operatorname{settledby}(t + Q + r + 2, t + \operatorname{twopow}(k))\), then if \(a \le t + \operatorname{twopow}(k + 1) - (r + Q + 1)\), then if \(t + \operatorname{twopow}(k + 1) - (r + Q + 1) \le a + L\), then if \(\operatorname{winwittagree}(t, k, Q, a, L)\), then if \(\operatorname{winwittagree}(t + 1, k, Q, a, L)\), then \(\operatorname{lhcenter}(t + 1 + Q + \operatorname{twopow}(k), \operatorname{twopow}(k) - (r + Q + 2)) = \operatorname{not}(\operatorname{lhcenter}(t + Q + \operatorname{twopow}(k), \operatorname{twopow}(k) - (r + Q + 1)))\).
\end{theorem}
\begin{proof}
The goal follows from \(\operatorname{lhcenterlatecrosslatestepflipatshiftofdoublesandwinwittagreezerogap}(Q, t, k, r, a, L)\).
\end{proof}
\begin{theorem}
\label{ll:paper:lhcenter-late-cross-late-or-eq-true-at-shift-of-doubles-and-winwittagree-zero-gap}
For every natural number \(Q\) and natural number \(t\) and natural number \(k\) and natural number \(r\) and natural number \(a\) and natural number \(L\), if \(t + Q \le r\), then if \(t + Q + r + 2 \le k\), then if \(\operatorname{zerodiagonal}(t + Q + r)\), then if \(\operatorname{doublesat}(Q, t + Q + r)\), then if \(\operatorname{settledby}(t + Q + r + 1, t + \operatorname{twopow}(k))\), then if \(\operatorname{settledby}(t + Q + r + 2, t + \operatorname{twopow}(k))\), then if \(a \le t + \operatorname{twopow}(k + 1) - (r + Q + 1)\), then if \(t + \operatorname{twopow}(k + 1) - (r + Q + 1) \le a + L\), then if \(\operatorname{winwittagree}(t, k, Q, a, L)\), then if \(\operatorname{winwittagree}(t + 1, k, Q, a, L)\), then \((\operatorname{lhcenter}(t + Q + \operatorname{twopow}(k), \operatorname{twopow}(k) - (r + Q + 2)) \operatorname{or} \operatorname{lhcenter}(t + Q + \operatorname{twopow}(k), \operatorname{twopow}(k) - (r + Q + 3))) = \operatorname{true}\).
\end{theorem}
\begin{proof}
The goal follows from \(\operatorname{lhcenterlatecrosslateoreqtrueatshiftofdoublesandwinwittagreezerogap}(Q, t, k, r, a, L)\).
\end{proof}
\begin{theorem}
\label{ll:paper:winwitt-scale-cross-zero-gap-late-or-eq-true-at-shift-of-doubles-and-winwittagree}
For every natural number \(Q\) and natural number \(t\) and natural number \(k\) and natural number \(r\) and natural number \(a\) and natural number \(L\), if \(t + Q \le r\), then if \(t + Q + r + 2 \le k\), then if \(\operatorname{zerodiagonal}(t + Q + r)\), then if \(\operatorname{doublesat}(Q, t + Q + r)\), then if \(\operatorname{settledby}(t + Q + r + 1, t + \operatorname{twopow}(k))\), then if \(\operatorname{settledby}(t + Q + r + 2, t + \operatorname{twopow}(k))\), then if \(a \le t + \operatorname{twopow}(k + 1) - (r + Q + 1)\), then if \(t + \operatorname{twopow}(k + 1) - (r + Q + 1) \le a + L\), then if \(\operatorname{winwittagree}(t, k, Q, a, L)\), then if \(\operatorname{winwittagree}(t + 1, k, Q, a, L)\), then \((\operatorname{wittbitbool}(\operatorname{wittwin}(t + Q, \operatorname{twopow}(k)), t + \operatorname{twopow}(k + 1) - (r + 2)) \operatorname{or} \operatorname{wittbitbool}(\operatorname{wittwin}(t + Q, \operatorname{twopow}(k)), t + \operatorname{twopow}(k + 1) - (r + 3))) = \operatorname{true}\).
\end{theorem}
\begin{proof}
The goal follows from \(\operatorname{winwittscalecrosszerogaplateoreqtrueatshiftofdoublesandwinwittagree}(Q, t, k, r, a, L)\).
\end{proof}
\begin{theorem}
\label{ll:paper:lhcenter-late-cross-pin-r-plus-three-of-doubles-two-and-winwittagree}
For every natural number \(t\) and natural number \(k\) and natural number \(r\) and natural number \(a\) and natural number \(L\), if \(t + 2 \le r\), then if \(t + 2 + r + 2 \le k\), then if \(\operatorname{zerodiagonal}(t + 2 + r)\), then if \(\operatorname{doublesat}(2, t + 2 + r)\), then if \(\operatorname{settledby}(t + 2 + r, t + \operatorname{twopow}(k))\), then if \(\operatorname{settledby}(t + 2 + r + 1, t + \operatorname{twopow}(k))\), then if \(\operatorname{settledby}(t + 2 + r + 2, t + \operatorname{twopow}(k))\), then if \(a \le t + \operatorname{twopow}(k + 1) - (r + 1)\), then if \(t + \operatorname{twopow}(k + 1) - (r + 1) \le a + L\), then if \(\operatorname{winwittagree}(t, k, 2, a, L)\), then if \(\operatorname{winwittagree}(t + 1, k, 2, a, L)\), then \(\operatorname{lhcenter}(t + \operatorname{twopow}(k), \operatorname{twopow}(k) - (r + 3)) = \operatorname{true}\).
\end{theorem}
\begin{proof}
The goal follows from \(\operatorname{lhcenterlatecrosspinrplusthreeofdoublestwoandwinwittagree}(t, k, r, a, L)\).
\end{proof}
\begin{theorem}
\label{ll:paper:lhcenter-late-cross-pin-r-plus-one-false-of-doubles-two-and-winwittagree}
For every natural number \(t\) and natural number \(k\) and natural number \(r\) and natural number \(a\) and natural number \(L\), if \(t + 2 \le r\), then if \(t + 2 + r + 2 \le k\), then if \(\operatorname{zerodiagonal}(t + 2 + r)\), then if \(\operatorname{doublesat}(2, t + 2 + r)\), then if \(\operatorname{settledby}(t + 2 + r, t + \operatorname{twopow}(k))\), then if \(\operatorname{settledby}(t + 2 + r + 1, t + \operatorname{twopow}(k))\), then if \(\operatorname{settledby}(t + 2 + r + 2, t + \operatorname{twopow}(k))\), then if \(a \le t + \operatorname{twopow}(k + 1) - (r + 1)\), then if \(t + \operatorname{twopow}(k + 1) - (r + 1) \le a + L\), then if \(\operatorname{winwittagree}(t, k, 2, a, L)\), then if \(\operatorname{winwittagree}(t + 1, k, 2, a, L)\), then \(\operatorname{lhcenter}(t + \operatorname{twopow}(k), \operatorname{twopow}(k) - (r + 1)) = \operatorname{false}\).
\end{theorem}
\begin{proof}
The goal follows from \(\operatorname{lhcenterlatecrosspinrplusonefalseofdoublestwoandwinwittagree}(t, k, r, a, L)\).
\end{proof}
\begin{theorem}
\label{ll:paper:lhcenter-late-cross-pin-r-false-of-doubles-two-and-winwittagree}
For every natural number \(t\) and natural number \(k\) and natural number \(r\) and natural number \(a\) and natural number \(L\), if \(t + 2 \le r\), then if \(t + 2 + r + 2 \le k\), then if \(\operatorname{zerodiagonal}(t + 2 + r)\), then if \(\operatorname{doublesat}(2, t + 2 + r)\), then if \(\operatorname{settledby}(t + 2 + r, t + \operatorname{twopow}(k))\), then if \(\operatorname{settledby}(t + 2 + r + 1, t + \operatorname{twopow}(k))\), then if \(\operatorname{settledby}(t + 2 + r + 2, t + \operatorname{twopow}(k))\), then if \(a \le t + \operatorname{twopow}(k + 1) - (r + 1)\), then if \(t + \operatorname{twopow}(k + 1) - (r + 1) \le a + L\), then if \(\operatorname{winwittagree}(t, k, 2, a, L)\), then if \(\operatorname{winwittagree}(t + 1, k, 2, a, L)\), then \(\operatorname{lhcenter}(t + \operatorname{twopow}(k), \operatorname{twopow}(k) - r) = \operatorname{false}\).
\end{theorem}
\begin{proof}
The goal follows from \(\operatorname{lhcenterlatecrosspinrfalseofdoublestwoandwinwittagree}(t, k, r, a, L)\).
\end{proof}
\begin{theorem}
\label{ll:paper:winwitt-scale-cross-zero-gap-shift-eq-of-doubles-le-and-winwittagree}
For every natural number \(Q\) and natural number \(z\) and natural number \(t\) and natural number \(k\) and natural number \(r\) and natural number \(a\) and natural number \(L\), if \(t + Q < r\), then if \(t + Q + r \le z\), then if \(z \le k\), then if \(\operatorname{doublesat}(Q, z)\), then if \(\operatorname{settledby}(t + Q + r, t + \operatorname{twopow}(k))\), then if \(a \le t + \operatorname{twopow}(k + 1) - r\), then if \(t + \operatorname{twopow}(k + 1) - r \le a + L\), then if \(\operatorname{winwittagree}(t, k, Q, a, L)\), then \(\operatorname{wittbitbool}(\operatorname{wittwin}(t, \operatorname{twopow}(k)), t + \operatorname{twopow}(k + 1) - r) = \operatorname{wittbitbool}(\operatorname{wittwin}(t, \operatorname{twopow}(k)), t + \operatorname{twopow}(k + 1) - (r + Q))\).
\end{theorem}
\begin{proof}
The goal follows from \(\operatorname{winwittscalecrosszerogapshifteqofdoublesleandwinwittagree}(Q, z, t, k, r, a, L)\).
\end{proof}
\begin{theorem}
\label{ll:paper:lhcenter-late-cross-shift-eq-of-doubles-le-and-winwittagree-zero-gap}
For every natural number \(Q\) and natural number \(z\) and natural number \(t\) and natural number \(k\) and natural number \(r\) and natural number \(a\) and natural number \(L\), if \(t + Q < r\), then if \(t + Q + r \le z\), then if \(z \le k\), then if \(\operatorname{doublesat}(Q, z)\), then if \(\operatorname{settledby}(t + Q + r, t + \operatorname{twopow}(k))\), then if \(a \le t + \operatorname{twopow}(k + 1) - r\), then if \(t + \operatorname{twopow}(k + 1) - r \le a + L\), then if \(\operatorname{winwittagree}(t, k, Q, a, L)\), then \(\operatorname{lhcenter}(t + \operatorname{twopow}(k), \operatorname{twopow}(k) - r) = \operatorname{lhcenter}(t + \operatorname{twopow}(k), \operatorname{twopow}(k) - (r + Q))\).
\end{theorem}
\begin{proof}
The goal follows from \(\operatorname{lhcenterlatecrossshifteqofdoublesleandwinwittagreezerogap}(Q, z, t, k, r, a, L)\).
\end{proof}
\begin{theorem}
\label{ll:paper:winwitt-scale-cross-zero-gap-same-bit-eq-of-doubles-le-and-winwittagree}
For every natural number \(Q\) and natural number \(z\) and natural number \(t\) and natural number \(k\) and natural number \(r\) and natural number \(a\) and natural number \(L\), if \(t + Q < r\), then if \(t + Q + r \le z\), then if \(z \le k\), then if \(\operatorname{doublesat}(Q, z)\), then if \(\operatorname{settledby}(t + Q + r, t + \operatorname{twopow}(k))\), then if \(a \le t + \operatorname{twopow}(k + 1) - (r + Q)\), then if \(t + \operatorname{twopow}(k + 1) - r \le a + L\), then if \(\operatorname{winwittagree}(t, k, Q, a, L)\), then \(\operatorname{wittbitbool}(\operatorname{wittwin}(t + Q, \operatorname{twopow}(k)), t + \operatorname{twopow}(k + 1) - r) = \operatorname{wittbitbool}(\operatorname{wittwin}(t, \operatorname{twopow}(k)), t + \operatorname{twopow}(k + 1) - r)\).
\end{theorem}
\begin{proof}
The goal follows from \(\operatorname{winwittscalecrosszerogapsamebiteqofdoublesleandwinwittagree}(Q, z, t, k, r, a, L)\).
\end{proof}
\begin{theorem}
\label{ll:paper:lhcenter-late-cross-eq-false-at-shift-of-zero-diag-and-winwittagree-zero-gap}
For every natural number \(Q\) and natural number \(t\) and natural number \(k\) and natural number \(r\) and natural number \(a\) and natural number \(L\), if \(t + Q < r\), then if \(t + Q + r \le k\), then if \(\operatorname{zerodiagonal}(t + Q + r)\), then if \(\operatorname{settledby}(t + Q + r, t + Q + \operatorname{twopow}(k))\), then if \(a \le t + \operatorname{twopow}(k + 1) - r\), then if \(t + \operatorname{twopow}(k + 1) - r \le a + L\), then if \(\operatorname{winwittagree}(t, k, Q, a, L)\), then \(\operatorname{lhcenter}(t + \operatorname{twopow}(k), \operatorname{twopow}(k) - r) = \operatorname{false}\).
\end{theorem}
\begin{proof}
The goal follows from \(\operatorname{lhcenterlatecrosseqfalseatshiftofzerodiagandwinwittagreezerogap}(Q, t, k, r, a, L)\).
\end{proof}
\begin{theorem}
\label{ll:paper:lhcenter-late-cross-eq-true-at-shift-after-zero-diag-and-winwittagree-zero-gap}
For every natural number \(Q\) and natural number \(t\) and natural number \(k\) and natural number \(r\) and natural number \(a\) and natural number \(L\), if \(t + Q \le r + 1\), then if \(t + Q + r + 2 \le k\), then if \(\operatorname{zerodiagonal}(t + Q + r)\), then if \(\operatorname{settledby}(t + Q + r + 2, t + Q + \operatorname{twopow}(k))\), then if \(a \le t + \operatorname{twopow}(k + 1) - (r + 2)\), then if \(t + \operatorname{twopow}(k + 1) - (r + 2) \le a + L\), then if \(\operatorname{winwittagree}(t, k, Q, a, L)\), then \(\operatorname{lhcenter}(t + \operatorname{twopow}(k), \operatorname{twopow}(k) - (r + 2)) = \operatorname{true}\).
\end{theorem}
\begin{proof}
The goal follows from \(\operatorname{lhcenterlatecrosseqtrueatshiftafterzerodiagandwinwittagreezerogap}(Q, t, k, r, a, L)\).
\end{proof}
\begin{theorem}
\label{ll:paper:lhcenter-late-cross-step-flip-at-shift-after-zero-diag-of-single-winwittagree-zero-gap}
For every natural number \(Q\) and natural number \(t\) and natural number \(k\) and natural number \(r\) and natural number \(a\) and natural number \(L\), if \(t + Q \le r + 1\), then if \(t + Q + r + 2 \le k\), then if \(\operatorname{zerodiagonal}(t + Q + r)\), then if \(\operatorname{settledby}(t + Q + r + 2, t + Q + \operatorname{twopow}(k))\), then if \(a \le t + \operatorname{twopow}(k + 1) - (r + 2)\), then if \(t + \operatorname{twopow}(k + 1) - (r + 2) \le a + L\), then if \(\operatorname{winwittagree}(t, k, Q, a, L)\), then \(\operatorname{lhcenter}(t + 1 + \operatorname{twopow}(k), \operatorname{twopow}(k) - (r + 2)) = \operatorname{not}(\operatorname{lhcenter}(t + \operatorname{twopow}(k), \operatorname{twopow}(k) - (r + 1)))\).
\end{theorem}
\begin{proof}
The goal follows from \(\operatorname{lhcenterlatecrossstepflipatshiftafterzerodiagofsinglewinwittagreezerogap}(Q, t, k, r, a, L)\).
\end{proof}
\begin{theorem}
\label{ll:paper:winwitt-scale-cross-zero-gap-step-flip-at-shift-after-zero-diag-of-single-winwittagree}
For every natural number \(Q\) and natural number \(t\) and natural number \(k\) and natural number \(r\) and natural number \(a\) and natural number \(L\), if \(t \le r\), then if \(t + Q \le r + 1\), then if \(t + Q + r + 2 \le k\), then if \(\operatorname{zerodiagonal}(t + Q + r)\), then if \(\operatorname{settledby}(t + Q + r + 2, t + Q + \operatorname{twopow}(k))\), then if \(a \le t + \operatorname{twopow}(k + 1) - (r + 2)\), then if \(t + \operatorname{twopow}(k + 1) - (r + 2) \le a + L\), then if \(\operatorname{winwittagree}(t, k, Q, a, L)\), then \(\operatorname{wittbitbool}(\operatorname{wittwin}(t + 1, \operatorname{twopow}(k)), t + \operatorname{twopow}(k + 1) - (r + 1)) = \operatorname{not}(\operatorname{wittbitbool}(\operatorname{wittwin}(t, \operatorname{twopow}(k)), t + \operatorname{twopow}(k + 1) - (r + 1)))\).
\end{theorem}
\begin{proof}
The goal follows from \(\operatorname{winwittscalecrosszerogapstepflipatshiftafterzerodiagofsinglewinwittagree}(Q, t, k, r, a, L)\).
\end{proof}
\begin{theorem}
\label{ll:paper:lhcenter-late-cross-step-one-eq-true-at-shift-of-zero-diag-and-single-winwittagree-zero-gap}
For every natural number \(Q\) and natural number \(t\) and natural number \(k\) and natural number \(r\) and natural number \(a\) and natural number \(L\), if \(t + Q < r\), then if \(t + Q + r + 2 \le k\), then if \(\operatorname{zerodiagonal}(t + Q + r)\), then if \(\operatorname{settledby}(t + Q + r, t + Q + \operatorname{twopow}(k))\), then if \(\operatorname{settledby}(t + Q + r + 2, t + Q + \operatorname{twopow}(k))\), then if \(a \le t + \operatorname{twopow}(k + 1) - (r + 2)\), then if \(t + \operatorname{twopow}(k + 1) - r \le a + L\), then if \(\operatorname{winwittagree}(t, k, Q, a, L)\), then \(\operatorname{lhcenter}(t + 1 + \operatorname{twopow}(k), \operatorname{twopow}(k) - (r + 1)) = \operatorname{true}\).
\end{theorem}
\begin{proof}
The goal follows from \(\operatorname{lhcenterlatecrosssteponeeqtrueatshiftofzerodiagandsinglewinwittagreezerogap}(Q, t, k, r, a, L)\).
\end{proof}
\begin{theorem}
\label{ll:paper:winwitt-scale-cross-zero-gap-reset-shift-flip-of-doubles-and-single-winwittagree}
For every natural number \(Q\) and natural number \(t\) and natural number \(k\) and natural number \(r\) and natural number \(a\) and natural number \(L\), if \(t + Q \le r\), then if \(t + Q + r + 2 \le k\), then if \(\operatorname{zerodiagonal}(t + Q + r)\), then if \(\operatorname{doublesat}(Q, t + Q + r)\), then if \(\operatorname{settledby}(t + Q + r + 1, t + \operatorname{twopow}(k))\), then if \(\operatorname{settledby}(t + Q + r + 2, t + \operatorname{twopow}(k))\), then if \(a \le t + \operatorname{twopow}(k + 1) - (r + 2)\), then if \(t + \operatorname{twopow}(k + 1) - (r + 1) \le a + L\), then if \(\operatorname{winwittagree}(t, k, Q, a, L)\), then \(\operatorname{wittbitbool}(\operatorname{wittwin}(t + 1, \operatorname{twopow}(k)), t + \operatorname{twopow}(k + 1) - (r + 1)) = \operatorname{not}(\operatorname{wittbitbool}(\operatorname{wittwin}(t + 1, \operatorname{twopow}(k)), t + \operatorname{twopow}(k + 1) - (r + Q + 1)))\).
\end{theorem}
\begin{proof}
The goal follows from \(\operatorname{winwittscalecrosszerogapresetshiftflipofdoublesandsinglewinwittagree}(Q, t, k, r, a, L)\).
\end{proof}
\begin{theorem}
\label{ll:paper:lhcenter-late-cross-reset-shift-flip-of-doubles-and-single-winwittagree-zero-gap}
For every natural number \(Q\) and natural number \(t\) and natural number \(k\) and natural number \(r\) and natural number \(a\) and natural number \(L\), if \(t + Q \le r\), then if \(t + Q + r + 2 \le k\), then if \(\operatorname{zerodiagonal}(t + Q + r)\), then if \(\operatorname{doublesat}(Q, t + Q + r)\), then if \(\operatorname{settledby}(t + Q + r + 1, t + \operatorname{twopow}(k))\), then if \(\operatorname{settledby}(t + Q + r + 2, t + \operatorname{twopow}(k))\), then if \(a \le t + \operatorname{twopow}(k + 1) - (r + 2)\), then if \(t + \operatorname{twopow}(k + 1) - (r + 1) \le a + L\), then if \(\operatorname{winwittagree}(t, k, Q, a, L)\), then \(\operatorname{lhcenter}(t + 1 + \operatorname{twopow}(k), \operatorname{twopow}(k) - (r + 2)) = \operatorname{not}(\operatorname{lhcenter}(t + 1 + \operatorname{twopow}(k), \operatorname{twopow}(k) - (r + Q + 2)))\).
\end{theorem}
\begin{proof}
The goal follows from \(\operatorname{lhcenterlatecrossresetshiftflipofdoublesandsinglewinwittagreezerogap}(Q, t, k, r, a, L)\).
\end{proof}
\begin{theorem}
\label{ll:paper:winwitt-scale-cross-zero-gap-eq-false-at-shift-of-zero-diag-and-winwittagree}
For every natural number \(Q\) and natural number \(t\) and natural number \(k\) and natural number \(r\) and natural number \(a\) and natural number \(L\), if \(t + Q < r\), then if \(t + Q + r \le k\), then if \(\operatorname{zerodiagonal}(t + Q + r)\), then if \(\operatorname{settledby}(t + Q + r, t + Q + \operatorname{twopow}(k))\), then if \(a \le t + \operatorname{twopow}(k + 1) - r\), then if \(t + \operatorname{twopow}(k + 1) - r \le a + L\), then if \(\operatorname{winwittagree}(t, k, Q, a, L)\), then \(\operatorname{wittbitbool}(\operatorname{wittwin}(t, \operatorname{twopow}(k)), t + \operatorname{twopow}(k + 1) - r) = \operatorname{false}\).
\end{theorem}
\begin{proof}
The goal follows from \(\operatorname{winwittscalecrosszerogapeqfalseatshiftofzerodiagandwinwittagree}(Q, t, k, r, a, L)\).
\end{proof}
\begin{theorem}
\label{ll:paper:winwitt-scale-cross-zero-gap-eq-true-at-shift-after-zero-diag-and-winwittagree}
For every natural number \(Q\) and natural number \(t\) and natural number \(k\) and natural number \(r\) and natural number \(a\) and natural number \(L\), if \(t \le r\), then if \(t + Q \le r + 1\), then if \(t + Q + r + 2 \le k\), then if \(\operatorname{zerodiagonal}(t + Q + r)\), then if \(\operatorname{settledby}(t + Q + r + 2, t + Q + \operatorname{twopow}(k))\), then if \(a \le t + \operatorname{twopow}(k + 1) - (r + 2)\), then if \(t + \operatorname{twopow}(k + 1) - (r + 2) \le a + L\), then if \(\operatorname{winwittagree}(t, k, Q, a, L)\), then \(\operatorname{wittbitbool}(\operatorname{wittwin}(t, \operatorname{twopow}(k)), t + \operatorname{twopow}(k + 1) - (r + 2)) = \operatorname{true}\).
\end{theorem}
\begin{proof}
The goal follows from \(\operatorname{winwittscalecrosszerogapeqtrueatshiftafterzerodiagandwinwittagree}(Q, t, k, r, a, L)\).
\end{proof}
\begin{theorem}
\label{ll:paper:winwitt-scale-cross-zero-gap-step-one-eq-true-at-shift-of-zero-diag-and-single-winwittagree}
For every natural number \(Q\) and natural number \(t\) and natural number \(k\) and natural number \(r\) and natural number \(a\) and natural number \(L\), if \(t + Q < r\), then if \(t + Q + r + 2 \le k\), then if \(\operatorname{zerodiagonal}(t + Q + r)\), then if \(\operatorname{settledby}(t + Q + r, t + Q + \operatorname{twopow}(k))\), then if \(\operatorname{settledby}(t + Q + r + 2, t + Q + \operatorname{twopow}(k))\), then if \(a \le t + \operatorname{twopow}(k + 1) - (r + 2)\), then if \(t + \operatorname{twopow}(k + 1) - r \le a + L\), then if \(\operatorname{winwittagree}(t, k, Q, a, L)\), then \(\operatorname{wittbitbool}(\operatorname{wittwin}(t + 1, \operatorname{twopow}(k)), t + \operatorname{twopow}(k + 1) - r) = \operatorname{true}\).
\end{theorem}
\begin{proof}
The goal follows from \(\operatorname{winwittscalecrosszerogapsteponeeqtrueatshiftofzerodiagandsinglewinwittagree}(Q, t, k, r, a, L)\).
\end{proof}
\begin{theorem}
\label{ll:paper:winwitt-scale-cross-zero-gap-step-flip-at-shift-of-zero-diag-and-single-winwittagree}
For every natural number \(Q\) and natural number \(t\) and natural number \(k\) and natural number \(r\) and natural number \(a\) and natural number \(L\), if \(t + Q < r\), then if \(t + Q + r + 2 \le k\), then if \(\operatorname{zerodiagonal}(t + Q + r)\), then if \(\operatorname{settledby}(t + Q + r, t + Q + \operatorname{twopow}(k))\), then if \(\operatorname{settledby}(t + Q + r + 2, t + Q + \operatorname{twopow}(k))\), then if \(a \le t + \operatorname{twopow}(k + 1) - (r + 2)\), then if \(t + \operatorname{twopow}(k + 1) - r \le a + L\), then if \(\operatorname{winwittagree}(t, k, Q, a, L)\), then \(\operatorname{wittbitbool}(\operatorname{wittwin}(t + 1, \operatorname{twopow}(k)), t + \operatorname{twopow}(k + 1) - r) = \operatorname{not}(\operatorname{wittbitbool}(\operatorname{wittwin}(t, \operatorname{twopow}(k)), t + \operatorname{twopow}(k + 1) - r))\).
\end{theorem}
\begin{proof}
The goal follows from \(\operatorname{winwittscalecrosszerogapstepflipatshiftofzerodiagandsinglewinwittagree}(Q, t, k, r, a, L)\).
\end{proof}
\begin{theorem}
\label{ll:paper:not-winwittagree-at-shift-after-zero-diag-of-zero-diag-two}
For every natural number \(Q\) and natural number \(t\) and natural number \(k\) and natural number \(r\) and natural number \(a\) and natural number \(L\), if \(t + Q \le r + 1\), then if \(t + Q + r + 2 \le k\), then if \(\operatorname{zerodiagonal}(t + Q + r)\), then if \(\operatorname{zerodiagonal}(t + r + 2)\), then if \(\operatorname{settledby}(t + r + 2, t + \operatorname{twopow}(k))\), then if \(\operatorname{settledby}(t + Q + r + 2, t + Q + \operatorname{twopow}(k))\), then if \(a \le t + \operatorname{twopow}(k + 1) - (r + 2)\), then if \(t + \operatorname{twopow}(k + 1) - (r + 2) \le a + L\), then not \(\operatorname{winwittagree}(t, k, Q, a, L)\).
\end{theorem}
\begin{proof}
The goal follows from \(\operatorname{notwinwittagreeatshiftafterzerodiagofzerodiagtwo}(Q, t, k, r, a, L)\).
\end{proof}
\begin{theorem}
\label{ll:paper:lhcenter-late-cross-pin-r-plus-one-of-single-winwittagree-and-zero-diag-three}
For every natural number \(Q\) and natural number \(t\) and natural number \(k\) and natural number \(r\) and natural number \(a\) and natural number \(L\), if \(t + Q \le r + 1\), then if \(t + Q + r + 3 \le k\), then if \(\operatorname{zerodiagonal}(t + Q + r)\), then if \(\operatorname{zerodiagonal}(t + r + 3)\), then if \(\operatorname{settledby}(t + r + 3, t + \operatorname{twopow}(k))\), then if \(\operatorname{settledby}(t + Q + r + 2, t + Q + \operatorname{twopow}(k))\), then if \(a \le t + \operatorname{twopow}(k + 1) - (r + 2)\), then if \(t + \operatorname{twopow}(k + 1) - (r + 2) \le a + L\), then if \(\operatorname{winwittagree}(t, k, Q, a, L)\), then \(\operatorname{lhcenter}(t + \operatorname{twopow}(k), \operatorname{twopow}(k) - (r + 1)) = \operatorname{true}\).
\end{theorem}
\begin{proof}
The goal follows from \(\operatorname{lhcenterlatecrosspinrplusoneofsinglewinwittagreeandzerodiagthree}(Q, t, k, r, a, L)\).
\end{proof}
\begin{theorem}
\label{ll:paper:lhcenter-late-cross-step-one-pin-r-plus-two-false-of-single-winwittagree-and-zero-diag-three}
For every natural number \(Q\) and natural number \(t\) and natural number \(k\) and natural number \(r\) and natural number \(a\) and natural number \(L\), if \(t + Q \le r + 1\), then if \(t + Q + r + 3 \le k\), then if \(\operatorname{zerodiagonal}(t + Q + r)\), then if \(\operatorname{zerodiagonal}(t + r + 3)\), then if \(\operatorname{settledby}(t + r + 3, t + \operatorname{twopow}(k))\), then if \(\operatorname{settledby}(t + Q + r + 2, t + Q + \operatorname{twopow}(k))\), then if \(a \le t + \operatorname{twopow}(k + 1) - (r + 2)\), then if \(t + \operatorname{twopow}(k + 1) - (r + 2) \le a + L\), then if \(\operatorname{winwittagree}(t, k, Q, a, L)\), then \(\operatorname{lhcenter}(t + 1 + \operatorname{twopow}(k), \operatorname{twopow}(k) - (r + 2)) = \operatorname{false}\).
\end{theorem}
\begin{proof}
The goal follows from \(\operatorname{lhcenterlatecrosssteponepinrplustwofalseofsinglewinwittagreeandzerodiagthree}(Q, t, k, r, a, L)\).
\end{proof}
\begin{theorem}
\label{ll:paper:lhcenter-late-cross-pin-doubling-diag-false-of-doubles-and-single-winwittagree-and-zero-diag-three}
For every natural number \(Q\) and natural number \(t\) and natural number \(k\) and natural number \(r\) and natural number \(a\) and natural number \(L\), if \(t + Q \le r\), then if \(t + Q + r + 3 \le k\), then if \(\operatorname{zerodiagonal}(t + Q + r)\), then if \(\operatorname{doublesat}(Q, t + Q + r)\), then if \(\operatorname{zerodiagonal}(t + r + 3)\), then if \(\operatorname{settledby}(t + r + 3, t + \operatorname{twopow}(k))\), then if \(\operatorname{settledby}(t + Q + r + 1, t + \operatorname{twopow}(k))\), then if \(\operatorname{settledby}(t + Q + r + 2, t + \operatorname{twopow}(k))\), then if \(a \le t + \operatorname{twopow}(k + 1) - (r + 2)\), then if \(t + \operatorname{twopow}(k + 1) - (r + 1) \le a + L\), then if \(\operatorname{winwittagree}(t, k, Q, a, L)\), then \(\operatorname{lhcenter}(t + \operatorname{twopow}(k), \operatorname{twopow}(k) - (r + Q + 1)) = \operatorname{false}\).
\end{theorem}
\begin{proof}
The goal follows from \(\operatorname{lhcenterlatecrosspindoublingdiagfalseofdoublesandsinglewinwittagreeandzerodiagthree}(Q, t, k, r, a, L)\).
\end{proof}
\begin{theorem}
\label{ll:paper:lhcenter-late-cross-step-one-pin-reset-diag-true-of-doubles-and-single-winwittagree-and-zero-diag-three}
For every natural number \(Q\) and natural number \(t\) and natural number \(k\) and natural number \(r\) and natural number \(a\) and natural number \(L\), if \(t + Q \le r\), then if \(t + Q + r + 3 \le k\), then if \(\operatorname{zerodiagonal}(t + Q + r)\), then if \(\operatorname{doublesat}(Q, t + Q + r)\), then if \(\operatorname{zerodiagonal}(t + r + 3)\), then if \(\operatorname{settledby}(t + r + 3, t + \operatorname{twopow}(k))\), then if \(\operatorname{settledby}(t + Q + r + 1, t + \operatorname{twopow}(k))\), then if \(\operatorname{settledby}(t + Q + r + 2, t + \operatorname{twopow}(k))\), then if \(a \le t + \operatorname{twopow}(k + 1) - (r + 2)\), then if \(t + \operatorname{twopow}(k + 1) - (r + 1) \le a + L\), then if \(\operatorname{winwittagree}(t, k, Q, a, L)\), then \(\operatorname{lhcenter}(t + 1 + \operatorname{twopow}(k), \operatorname{twopow}(k) - (r + Q + 2)) = \operatorname{true}\).
\end{theorem}
\begin{proof}
The goal follows from \(\operatorname{lhcenterlatecrosssteponepinresetdiagtrueofdoublesandsinglewinwittagreeandzerodiagthree}(Q, t, k, r, a, L)\).
\end{proof}
\begin{theorem}
\label{ll:paper:lhcenter-late-cross-pin-before-zero-diag-of-zero-diag-three-and-single-winwittagree}
For every natural number \(t\) and natural number \(k\) and natural number \(r\) and natural number \(a\) and natural number \(L\), if \(t + 2 \le r\), then if \(t + 3 + r + 2 \le k\), then if \(\operatorname{zerodiagonal}(t + 3 + r)\), then if \(\operatorname{settledby}(t + 3 + r, t + \operatorname{twopow}(k))\), then if \(\operatorname{settledby}(t + 3 + r + 2, t + 3 + \operatorname{twopow}(k))\), then if \(a \le t + \operatorname{twopow}(k + 1) - (r + 2)\), then if \(t + \operatorname{twopow}(k + 1) - (r + 2) \le a + L\), then if \(\operatorname{winwittagree}(t, k, 3, a, L)\), then \(\operatorname{lhcenter}(t + \operatorname{twopow}(k), \operatorname{twopow}(k) - (r + 1)) = \operatorname{true}\).
\end{theorem}
\begin{proof}
The goal follows from \(\operatorname{lhcenterlatecrosspinbeforezerodiagofzerodiagthreeandsinglewinwittagree}(t, k, r, a, L)\).
\end{proof}
\begin{theorem}
\label{ll:paper:lhcenter-late-cross-step-one-pin-r-plus-two-false-of-zero-diag-three-and-single-winwittagree}
For every natural number \(t\) and natural number \(k\) and natural number \(r\) and natural number \(a\) and natural number \(L\), if \(t + 2 \le r\), then if \(t + 3 + r + 2 \le k\), then if \(\operatorname{zerodiagonal}(t + 3 + r)\), then if \(\operatorname{settledby}(t + 3 + r, t + \operatorname{twopow}(k))\), then if \(\operatorname{settledby}(t + 3 + r + 2, t + 3 + \operatorname{twopow}(k))\), then if \(a \le t + \operatorname{twopow}(k + 1) - (r + 2)\), then if \(t + \operatorname{twopow}(k + 1) - (r + 2) \le a + L\), then if \(\operatorname{winwittagree}(t, k, 3, a, L)\), then \(\operatorname{lhcenter}(t + 1 + \operatorname{twopow}(k), \operatorname{twopow}(k) - (r + 2)) = \operatorname{false}\).
\end{theorem}
\begin{proof}
The goal follows from \(\operatorname{lhcenterlatecrosssteponepinrplustwofalseofzerodiagthreeandsinglewinwittagree}(t, k, r, a, L)\).
\end{proof}
\begin{theorem}
\label{ll:paper:lhcenter-late-cross-pin-doubling-diag-false-of-doubles-three-and-single-winwittagree}
For every natural number \(t\) and natural number \(k\) and natural number \(r\) and natural number \(a\) and natural number \(L\), if \(t + 3 \le r\), then if \(t + 3 + r + 2 \le k\), then if \(\operatorname{zerodiagonal}(t + 3 + r)\), then if \(\operatorname{doublesat}(3, t + 3 + r)\), then if \(\operatorname{settledby}(t + 3 + r, t + \operatorname{twopow}(k))\), then if \(\operatorname{settledby}(t + 3 + r + 1, t + \operatorname{twopow}(k))\), then if \(\operatorname{settledby}(t + 3 + r + 2, t + \operatorname{twopow}(k))\), then if \(a \le t + \operatorname{twopow}(k + 1) - (r + 2)\), then if \(t + \operatorname{twopow}(k + 1) - (r + 1) \le a + L\), then if \(\operatorname{winwittagree}(t, k, 3, a, L)\), then \(\operatorname{lhcenter}(t + \operatorname{twopow}(k), \operatorname{twopow}(k) - (r + 4)) = \operatorname{false}\).
\end{theorem}
\begin{proof}
The goal follows from \(\operatorname{lhcenterlatecrosspindoublingdiagfalseofdoublesthreeandsinglewinwittagree}(t, k, r, a, L)\).
\end{proof}
\begin{theorem}
\label{ll:paper:lhcenter-late-cross-step-one-pin-reset-diag-true-of-doubles-three-and-single-winwittagree}
For every natural number \(t\) and natural number \(k\) and natural number \(r\) and natural number \(a\) and natural number \(L\), if \(t + 3 \le r\), then if \(t + 3 + r + 2 \le k\), then if \(\operatorname{zerodiagonal}(t + 3 + r)\), then if \(\operatorname{doublesat}(3, t + 3 + r)\), then if \(\operatorname{settledby}(t + 3 + r, t + \operatorname{twopow}(k))\), then if \(\operatorname{settledby}(t + 3 + r + 1, t + \operatorname{twopow}(k))\), then if \(\operatorname{settledby}(t + 3 + r + 2, t + \operatorname{twopow}(k))\), then if \(a \le t + \operatorname{twopow}(k + 1) - (r + 2)\), then if \(t + \operatorname{twopow}(k + 1) - (r + 1) \le a + L\), then if \(\operatorname{winwittagree}(t, k, 3, a, L)\), then \(\operatorname{lhcenter}(t + 1 + \operatorname{twopow}(k), \operatorname{twopow}(k) - (r + 5)) = \operatorname{true}\).
\end{theorem}
\begin{proof}
The goal follows from \(\operatorname{lhcenterlatecrosssteponepinresetdiagtrueofdoublesthreeandsinglewinwittagree}(t, k, r, a, L)\).
\end{proof}
\begin{theorem}
\label{ll:paper:not-winwittagree-two-at-shift-after-zero-diag}
For every natural number \(t\) and natural number \(k\) and natural number \(r\) and natural number \(a\) and natural number \(L\), if \(t + 1 \le r\), then if \(t + 2 + r + 2 \le k\), then if \(\operatorname{zerodiagonal}(t + 2 + r)\), then if \(\operatorname{settledby}(t + 2 + r, t + \operatorname{twopow}(k))\), then if \(\operatorname{settledby}(t + 2 + r + 2, t + 2 + \operatorname{twopow}(k))\), then if \(a \le t + \operatorname{twopow}(k + 1) - (r + 2)\), then if \(t + \operatorname{twopow}(k + 1) - (r + 2) \le a + L\), then not \(\operatorname{winwittagree}(t, k, 2, a, L)\).
\end{theorem}
\begin{proof}
The goal follows from \(\operatorname{notwinwittagreetwoatshiftafterzerodiag}(t, k, r, a, L)\).
\end{proof}
\begin{theorem}
\label{ll:paper:not-winwittagree-of-doubles-and-same-bit-eq}
For every natural number \(Q\) and natural number \(t\) and natural number \(k\) and natural number \(r\) and natural number \(a\) and natural number \(L\), if \(t + Q \le r\), then if \(t + Q + r + 1 \le k\), then if \(\operatorname{zerodiagonal}(t + Q + r)\), then if \(\operatorname{doublesat}(Q, t + Q + r)\), then if \(\operatorname{settledby}(t + Q + r + 1, t + \operatorname{twopow}(k))\), then if \(a \le t + \operatorname{twopow}(k + 1) - (r + Q + 1)\), then if \(t + \operatorname{twopow}(k + 1) - (r + 1) \le a + L\), then if \(\operatorname{wittbitbool}(\operatorname{wittwin}(t + Q, \operatorname{twopow}(k)), t + \operatorname{twopow}(k + 1) - (r + 1)) = \operatorname{wittbitbool}(\operatorname{wittwin}(t, \operatorname{twopow}(k)), t + \operatorname{twopow}(k + 1) - (r + 1))\), then not \(\operatorname{winwittagree}(t, k, Q, a, L)\).
\end{theorem}
\begin{proof}
The goal follows from \(\operatorname{notwinwittagreeofdoublesandsamebiteq}(Q, t, k, r, a, L)\).
\end{proof}
\begin{theorem}
\label{ll:paper:winwitt-scale-cross-zero-gap-pin-r-plus-one-of-single-winwittagree-and-zero-diag-three}
For every natural number \(Q\) and natural number \(t\) and natural number \(k\) and natural number \(r\) and natural number \(a\) and natural number \(L\), if \(t \le r\), then if \(t + Q \le r + 1\), then if \(t + Q + r + 3 \le k\), then if \(\operatorname{zerodiagonal}(t + Q + r)\), then if \(\operatorname{zerodiagonal}(t + r + 3)\), then if \(\operatorname{settledby}(t + r + 3, t + \operatorname{twopow}(k))\), then if \(\operatorname{settledby}(t + Q + r + 2, t + Q + \operatorname{twopow}(k))\), then if \(a \le t + \operatorname{twopow}(k + 1) - (r + 2)\), then if \(t + \operatorname{twopow}(k + 1) - (r + 2) \le a + L\), then if \(\operatorname{winwittagree}(t, k, Q, a, L)\), then \(\operatorname{wittbitbool}(\operatorname{wittwin}(t, \operatorname{twopow}(k)), t + \operatorname{twopow}(k + 1) - (r + 1)) = \operatorname{true}\).
\end{theorem}
\begin{proof}
The goal follows from \(\operatorname{winwittscalecrosszerogappinrplusoneofsinglewinwittagreeandzerodiagthree}(Q, t, k, r, a, L)\).
\end{proof}
\begin{theorem}
\label{ll:paper:winwitt-scale-cross-zero-gap-step-one-pin-r-plus-one-false-of-single-winwittagree-and-zero-diag-three}
For every natural number \(Q\) and natural number \(t\) and natural number \(k\) and natural number \(r\) and natural number \(a\) and natural number \(L\), if \(t \le r\), then if \(t + Q \le r + 1\), then if \(t + Q + r + 3 \le k\), then if \(\operatorname{zerodiagonal}(t + Q + r)\), then if \(\operatorname{zerodiagonal}(t + r + 3)\), then if \(\operatorname{settledby}(t + r + 3, t + \operatorname{twopow}(k))\), then if \(\operatorname{settledby}(t + Q + r + 2, t + Q + \operatorname{twopow}(k))\), then if \(a \le t + \operatorname{twopow}(k + 1) - (r + 2)\), then if \(t + \operatorname{twopow}(k + 1) - (r + 2) \le a + L\), then if \(\operatorname{winwittagree}(t, k, Q, a, L)\), then \(\operatorname{wittbitbool}(\operatorname{wittwin}(t + 1, \operatorname{twopow}(k)), t + \operatorname{twopow}(k + 1) - (r + 1)) = \operatorname{false}\).
\end{theorem}
\begin{proof}
The goal follows from \(\operatorname{winwittscalecrosszerogapsteponepinrplusonefalseofsinglewinwittagreeandzerodiagthree}(Q, t, k, r, a, L)\).
\end{proof}
\begin{theorem}
\label{ll:paper:winwitt-scale-cross-zero-gap-pin-doubling-diag-false-of-doubles-and-single-winwittagree-and-zero-diag-three}
For every natural number \(Q\) and natural number \(t\) and natural number \(k\) and natural number \(r\) and natural number \(a\) and natural number \(L\), if \(t + Q \le r\), then if \(t + Q + r + 3 \le k\), then if \(\operatorname{zerodiagonal}(t + Q + r)\), then if \(\operatorname{doublesat}(Q, t + Q + r)\), then if \(\operatorname{zerodiagonal}(t + r + 3)\), then if \(\operatorname{settledby}(t + r + 3, t + \operatorname{twopow}(k))\), then if \(\operatorname{settledby}(t + Q + r + 1, t + \operatorname{twopow}(k))\), then if \(\operatorname{settledby}(t + Q + r + 2, t + \operatorname{twopow}(k))\), then if \(a \le t + \operatorname{twopow}(k + 1) - (r + 2)\), then if \(t + \operatorname{twopow}(k + 1) - (r + 1) \le a + L\), then if \(\operatorname{winwittagree}(t, k, Q, a, L)\), then \(\operatorname{wittbitbool}(\operatorname{wittwin}(t, \operatorname{twopow}(k)), t + \operatorname{twopow}(k + 1) - (r + Q + 1)) = \operatorname{false}\).
\end{theorem}
\begin{proof}
The goal follows from \(\operatorname{winwittscalecrosszerogappindoublingdiagfalseofdoublesandsinglewinwittagreeandzerodiagthree}(Q, t, k, r, a, L)\).
\end{proof}
\begin{theorem}
\label{ll:paper:lhcenter-late-cross-late-pin-doubling-diag-true-of-single-winwittagree-and-zero-diag-three}
For every natural number \(Q\) and natural number \(t\) and natural number \(k\) and natural number \(r\) and natural number \(a\) and natural number \(L\), if \(t + Q \le r\), then if \(t + Q + r + 3 \le k\), then if \(\operatorname{zerodiagonal}(t + Q + r)\), then if \(\operatorname{zerodiagonal}(t + r + 3)\), then if \(\operatorname{settledby}(t + r + 3, t + \operatorname{twopow}(k))\), then if \(\operatorname{settledby}(t + Q + r + 2, t + Q + \operatorname{twopow}(k))\), then if \(a \le t + \operatorname{twopow}(k + 1) - (r + 2)\), then if \(t + \operatorname{twopow}(k + 1) - (r + 1) \le a + L\), then if \(\operatorname{winwittagree}(t, k, Q, a, L)\), then \(\operatorname{lhcenter}(t + Q + \operatorname{twopow}(k), \operatorname{twopow}(k) - (r + 1)) = \operatorname{true}\).
\end{theorem}
\begin{proof}
The goal follows from \(\operatorname{lhcenterlatecrosslatepindoublingdiagtrueofsinglewinwittagreeandzerodiagthree}(Q, t, k, r, a, L)\).
\end{proof}
\begin{theorem}
\label{ll:paper:winwitt-scale-cross-zero-gap-step-one-pin-reset-diag-true-of-doubles-and-single-winwittagree-and-zero-diag-three}
For every natural number \(Q\) and natural number \(t\) and natural number \(k\) and natural number \(r\) and natural number \(a\) and natural number \(L\), if \(t + Q \le r\), then if \(t + Q + r + 3 \le k\), then if \(\operatorname{zerodiagonal}(t + Q + r)\), then if \(\operatorname{doublesat}(Q, t + Q + r)\), then if \(\operatorname{zerodiagonal}(t + r + 3)\), then if \(\operatorname{settledby}(t + r + 3, t + \operatorname{twopow}(k))\), then if \(\operatorname{settledby}(t + Q + r + 1, t + \operatorname{twopow}(k))\), then if \(\operatorname{settledby}(t + Q + r + 2, t + \operatorname{twopow}(k))\), then if \(a \le t + \operatorname{twopow}(k + 1) - (r + 2)\), then if \(t + \operatorname{twopow}(k + 1) - (r + 1) \le a + L\), then if \(\operatorname{winwittagree}(t, k, Q, a, L)\), then \(\operatorname{wittbitbool}(\operatorname{wittwin}(t + 1, \operatorname{twopow}(k)), t + \operatorname{twopow}(k + 1) - (r + Q + 1)) = \operatorname{true}\).
\end{theorem}
\begin{proof}
The goal follows from \(\operatorname{winwittscalecrosszerogapsteponepinresetdiagtrueofdoublesandsinglewinwittagreeandzerodiagthree}(Q, t, k, r, a, L)\).
\end{proof}
\begin{theorem}
\label{ll:paper:winwitt-scale-cross-zero-gap-pin-before-zero-diag-of-zero-diag-three-and-single-winwittagree}
For every natural number \(t\) and natural number \(k\) and natural number \(r\) and natural number \(a\) and natural number \(L\), if \(t + 2 \le r\), then if \(t + 3 + r + 2 \le k\), then if \(\operatorname{zerodiagonal}(t + 3 + r)\), then if \(\operatorname{settledby}(t + 3 + r, t + \operatorname{twopow}(k))\), then if \(\operatorname{settledby}(t + 3 + r + 2, t + 3 + \operatorname{twopow}(k))\), then if \(a \le t + \operatorname{twopow}(k + 1) - (r + 2)\), then if \(t + \operatorname{twopow}(k + 1) - (r + 2) \le a + L\), then if \(\operatorname{winwittagree}(t, k, 3, a, L)\), then \(\operatorname{wittbitbool}(\operatorname{wittwin}(t, \operatorname{twopow}(k)), t + \operatorname{twopow}(k + 1) - (r + 1)) = \operatorname{true}\).
\end{theorem}
\begin{proof}
The goal follows from \(\operatorname{winwittscalecrosszerogappinbeforezerodiagofzerodiagthreeandsinglewinwittagree}(t, k, r, a, L)\).
\end{proof}
\begin{theorem}
\label{ll:paper:winwitt-scale-cross-zero-gap-step-one-pin-before-zero-diag-false-of-zero-diag-three-and-single-winwittagree}
For every natural number \(t\) and natural number \(k\) and natural number \(r\) and natural number \(a\) and natural number \(L\), if \(t + 2 \le r\), then if \(t + 3 + r + 2 \le k\), then if \(\operatorname{zerodiagonal}(t + 3 + r)\), then if \(\operatorname{settledby}(t + 3 + r, t + \operatorname{twopow}(k))\), then if \(\operatorname{settledby}(t + 3 + r + 2, t + 3 + \operatorname{twopow}(k))\), then if \(a \le t + \operatorname{twopow}(k + 1) - (r + 2)\), then if \(t + \operatorname{twopow}(k + 1) - (r + 2) \le a + L\), then if \(\operatorname{winwittagree}(t, k, 3, a, L)\), then \(\operatorname{wittbitbool}(\operatorname{wittwin}(t + 1, \operatorname{twopow}(k)), t + \operatorname{twopow}(k + 1) - (r + 1)) = \operatorname{false}\).
\end{theorem}
\begin{proof}
The goal follows from \(\operatorname{winwittscalecrosszerogapsteponepinbeforezerodiagfalseofzerodiagthreeandsinglewinwittagree}(t, k, r, a, L)\).
\end{proof}
\begin{theorem}
\label{ll:paper:winwitt-scale-cross-zero-gap-pin-doubling-diag-false-of-doubles-three-and-single-winwittagree}
For every natural number \(t\) and natural number \(k\) and natural number \(r\) and natural number \(a\) and natural number \(L\), if \(t + 3 \le r\), then if \(t + 3 + r + 2 \le k\), then if \(\operatorname{zerodiagonal}(t + 3 + r)\), then if \(\operatorname{doublesat}(3, t + 3 + r)\), then if \(\operatorname{settledby}(t + 3 + r, t + \operatorname{twopow}(k))\), then if \(\operatorname{settledby}(t + 3 + r + 1, t + \operatorname{twopow}(k))\), then if \(\operatorname{settledby}(t + 3 + r + 2, t + \operatorname{twopow}(k))\), then if \(a \le t + \operatorname{twopow}(k + 1) - (r + 2)\), then if \(t + \operatorname{twopow}(k + 1) - (r + 1) \le a + L\), then if \(\operatorname{winwittagree}(t, k, 3, a, L)\), then \(\operatorname{wittbitbool}(\operatorname{wittwin}(t, \operatorname{twopow}(k)), t + \operatorname{twopow}(k + 1) - (r + 4)) = \operatorname{false}\).
\end{theorem}
\begin{proof}
The goal follows from \(\operatorname{winwittscalecrosszerogappindoublingdiagfalseofdoublesthreeandsinglewinwittagree}(t, k, r, a, L)\).
\end{proof}
\begin{theorem}
\label{ll:paper:winwitt-scale-cross-zero-gap-step-one-pin-reset-diag-true-of-doubles-three-and-single-winwittagree}
For every natural number \(t\) and natural number \(k\) and natural number \(r\) and natural number \(a\) and natural number \(L\), if \(t + 3 \le r\), then if \(t + 3 + r + 2 \le k\), then if \(\operatorname{zerodiagonal}(t + 3 + r)\), then if \(\operatorname{doublesat}(3, t + 3 + r)\), then if \(\operatorname{settledby}(t + 3 + r, t + \operatorname{twopow}(k))\), then if \(\operatorname{settledby}(t + 3 + r + 1, t + \operatorname{twopow}(k))\), then if \(\operatorname{settledby}(t + 3 + r + 2, t + \operatorname{twopow}(k))\), then if \(a \le t + \operatorname{twopow}(k + 1) - (r + 2)\), then if \(t + \operatorname{twopow}(k + 1) - (r + 1) \le a + L\), then if \(\operatorname{winwittagree}(t, k, 3, a, L)\), then \(\operatorname{wittbitbool}(\operatorname{wittwin}(t + 1, \operatorname{twopow}(k)), t + \operatorname{twopow}(k + 1) - (r + 4)) = \operatorname{true}\).
\end{theorem}
\begin{proof}
The goal follows from \(\operatorname{winwittscalecrosszerogapsteponepinresetdiagtrueofdoublesthreeandsinglewinwittagree}(t, k, r, a, L)\).
\end{proof}
\begin{theorem}
\label{ll:paper:lhcenter-late-cross-pin-doubling-diag-true-of-doubles-one-and-single-winwittagree}
For every natural number \(t\) and natural number \(k\) and natural number \(r\) and natural number \(a\) and natural number \(L\), if \(t + 1 \le r\), then if \(t + 1 + r + 1 \le k\), then if \(\operatorname{zerodiagonal}(t + 1 + r)\), then if \(\operatorname{doublesat}(1, t + 1 + r)\), then if \(\operatorname{settledby}(t + 1 + r, t + \operatorname{twopow}(k))\), then if \(\operatorname{settledby}(t + 1 + r + 1, t + \operatorname{twopow}(k))\), then if \(a \le t + \operatorname{twopow}(k + 1) - (r + 1)\), then if \(t + \operatorname{twopow}(k + 1) - (r + 1) \le a + L\), then if \(\operatorname{winwittagree}(t, k, 1, a, L)\), then \(\operatorname{lhcenter}(t + \operatorname{twopow}(k), \operatorname{twopow}(k) - (r + 2)) = \operatorname{true}\).
\end{theorem}
\begin{proof}
The goal follows from \(\operatorname{lhcenterlatecrosspindoublingdiagtrueofdoublesoneandsinglewinwittagree}(t, k, r, a, L)\).
\end{proof}
\begin{theorem}
\label{ll:paper:winwitt-scale-cross-zero-gap-pin-doubling-diag-true-of-doubles-one-and-single-winwittagree}
For every natural number \(t\) and natural number \(k\) and natural number \(r\) and natural number \(a\) and natural number \(L\), if \(t + 1 \le r\), then if \(t + 1 + r + 1 \le k\), then if \(\operatorname{zerodiagonal}(t + 1 + r)\), then if \(\operatorname{doublesat}(1, t + 1 + r)\), then if \(\operatorname{settledby}(t + 1 + r, t + \operatorname{twopow}(k))\), then if \(\operatorname{settledby}(t + 1 + r + 1, t + \operatorname{twopow}(k))\), then if \(a \le t + \operatorname{twopow}(k + 1) - (r + 1)\), then if \(t + \operatorname{twopow}(k + 1) - (r + 1) \le a + L\), then if \(\operatorname{winwittagree}(t, k, 1, a, L)\), then \(\operatorname{wittbitbool}(\operatorname{wittwin}(t, \operatorname{twopow}(k)), t + \operatorname{twopow}(k + 1) - (r + 2)) = \operatorname{true}\).
\end{theorem}
\begin{proof}
The goal follows from \(\operatorname{winwittscalecrosszerogappindoublingdiagtrueofdoublesoneandsinglewinwittagree}(t, k, r, a, L)\).
\end{proof}
\begin{theorem}
\label{ll:paper:lhcenter-late-cross-step-one-pin-reset-diag-false-of-doubles-one-and-single-winwittagree}
For every natural number \(t\) and natural number \(k\) and natural number \(r\) and natural number \(a\) and natural number \(L\), if \(t + 1 \le r\), then if \(t + 1 + r + 2 \le k\), then if \(\operatorname{zerodiagonal}(t + 1 + r)\), then if \(\operatorname{doublesat}(1, t + 1 + r)\), then if \(\operatorname{settledby}(t + 1 + r + 1, t + \operatorname{twopow}(k))\), then if \(\operatorname{settledby}(t + 1 + r + 2, t + \operatorname{twopow}(k))\), then if \(a \le t + \operatorname{twopow}(k + 1) - (r + 2)\), then if \(t + \operatorname{twopow}(k + 1) - (r + 1) \le a + L\), then if \(\operatorname{winwittagree}(t, k, 1, a, L)\), then \(\operatorname{lhcenter}(t + 1 + \operatorname{twopow}(k), \operatorname{twopow}(k) - (r + 3)) = \operatorname{false}\).
\end{theorem}
\begin{proof}
The goal follows from \(\operatorname{lhcenterlatecrosssteponepinresetdiagfalseofdoublesoneandsinglewinwittagree}(t, k, r, a, L)\).
\end{proof}
\begin{theorem}
\label{ll:paper:winwitt-scale-cross-zero-gap-step-one-pin-reset-diag-false-of-doubles-one-and-single-winwittagree}
For every natural number \(t\) and natural number \(k\) and natural number \(r\) and natural number \(a\) and natural number \(L\), if \(t + 1 \le r\), then if \(t + 1 + r + 2 \le k\), then if \(\operatorname{zerodiagonal}(t + 1 + r)\), then if \(\operatorname{doublesat}(1, t + 1 + r)\), then if \(\operatorname{settledby}(t + 1 + r + 1, t + \operatorname{twopow}(k))\), then if \(\operatorname{settledby}(t + 1 + r + 2, t + \operatorname{twopow}(k))\), then if \(a \le t + \operatorname{twopow}(k + 1) - (r + 2)\), then if \(t + \operatorname{twopow}(k + 1) - (r + 1) \le a + L\), then if \(\operatorname{winwittagree}(t, k, 1, a, L)\), then \(\operatorname{wittbitbool}(\operatorname{wittwin}(t + 1, \operatorname{twopow}(k)), t + \operatorname{twopow}(k + 1) - (r + 2)) = \operatorname{false}\).
\end{theorem}
\begin{proof}
The goal follows from \(\operatorname{winwittscalecrosszerogapsteponepinresetdiagfalseofdoublesoneandsinglewinwittagree}(t, k, r, a, L)\).
\end{proof}
\begin{theorem}
\label{ll:paper:not-winwittagree-one-of-doubles}
For every natural number \(t\) and natural number \(k\) and natural number \(r\) and natural number \(a\) and natural number \(L\), if \(t + 1 \le r\), then if \(t + 1 + r + 2 \le k\), then if \(\operatorname{zerodiagonal}(t + 1 + r)\), then if \(\operatorname{doublesat}(1, t + 1 + r)\), then if \(\operatorname{settledby}(t + 1 + r + 1, t + \operatorname{twopow}(k))\), then if \(\operatorname{settledby}(t + 1 + r + 2, t + \operatorname{twopow}(k))\), then if \(a \le t + \operatorname{twopow}(k + 1) - (r + 3)\), then if \(t + \operatorname{twopow}(k + 1) - (r + 1) \le a + L\), then not \(\operatorname{winwittagree}(t, k, 1, a, L)\).
\end{theorem}
\begin{proof}
The goal follows from \(\operatorname{notwinwittagreeoneofdoubles}(t, k, r, a, L)\).
\end{proof}
\begin{theorem}
\label{ll:paper:not-winwittagree-of-doubles-and-zero-diag-three-and-after-three}
For every natural number \(Q\) and natural number \(t\) and natural number \(k\) and natural number \(r\) and natural number \(a\) and natural number \(L\), if \(t + Q \le r\), then if \(t + Q + r + 3 \le k\), then if \(\operatorname{zerodiagonal}(t + Q + r)\), then if \(\operatorname{doublesat}(Q, t + Q + r)\), then if \(\operatorname{zerodiagonal}(t + r + 3)\), then if \(\operatorname{zerodiagonal}(t + Q + r + 3)\), then if \(\operatorname{settledby}(t + r + 3, t + \operatorname{twopow}(k))\), then if \(\operatorname{settledby}(t + Q + r + 1, t + \operatorname{twopow}(k))\), then if \(\operatorname{settledby}(t + Q + r + 2, t + \operatorname{twopow}(k))\), then if \(\operatorname{settledby}(t + Q + r + 3, t + 1 + \operatorname{twopow}(k))\), then if \(a \le t + \operatorname{twopow}(k + 1) - (r + 2)\), then if \(t + \operatorname{twopow}(k + 1) - (r + 1) \le a + L\), then not \(\operatorname{winwittagree}(t, k, Q, a, L)\).
\end{theorem}
\begin{proof}
The goal follows from \(\operatorname{notwinwittagreeofdoublesandzerodiagthreeandafterthree}(Q, t, k, r, a, L)\).
\end{proof}
\begin{theorem}
\label{ll:paper:lhcenter-late-cross-twin-at-shift-of-zero-diag-and-winwittagree-zero-gap}
For every natural number \(Q\) and natural number \(t\) and natural number \(k\) and natural number \(r\) and natural number \(a\) and natural number \(L\), if \(t + Q < r\), then if \(t + Q + r + 2 \le k\), then if \(\operatorname{zerodiagonal}(t + Q + r + 2)\), then if \(\operatorname{settledby}(t + Q + r + 2, t + Q + \operatorname{twopow}(k))\), then if \(a \le t + \operatorname{twopow}(k + 1) - (r + 1)\), then if \(t + \operatorname{twopow}(k + 1) - r \le a + L\), then if \(\operatorname{winwittagree}(t, k, Q, a, L)\), then \(\operatorname{lhcenter}(t + \operatorname{twopow}(k), \operatorname{twopow}(k) - r) = \operatorname{lhcenter}(t + \operatorname{twopow}(k), \operatorname{twopow}(k) - (r + 1))\).
\end{theorem}
\begin{proof}
The goal follows from \(\operatorname{lhcenterlatecrosstwinatshiftofzerodiagandwinwittagreezerogap}(Q, t, k, r, a, L)\).
\end{proof}
\begin{theorem}
\label{ll:paper:winwitt-scale-cross-zero-gap-twin-at-shift-of-zero-diag-and-winwittagree}
For every natural number \(Q\) and natural number \(t\) and natural number \(k\) and natural number \(r\) and natural number \(a\) and natural number \(L\), if \(t + Q < r\), then if \(t + Q + r + 2 \le k\), then if \(\operatorname{zerodiagonal}(t + Q + r + 2)\), then if \(\operatorname{settledby}(t + Q + r + 2, t + Q + \operatorname{twopow}(k))\), then if \(a \le t + \operatorname{twopow}(k + 1) - (r + 1)\), then if \(t + \operatorname{twopow}(k + 1) - r \le a + L\), then if \(\operatorname{winwittagree}(t, k, Q, a, L)\), then \(\operatorname{wittbitbool}(\operatorname{wittwin}(t, \operatorname{twopow}(k)), t + \operatorname{twopow}(k + 1) - r) = \operatorname{wittbitbool}(\operatorname{wittwin}(t, \operatorname{twopow}(k)), t + \operatorname{twopow}(k + 1) - (r + 1))\).
\end{theorem}
\begin{proof}
The goal follows from \(\operatorname{winwittscalecrosszerogaptwinatshiftofzerodiagandwinwittagree}(Q, t, k, r, a, L)\).
\end{proof}
\begin{theorem}
\label{ll:paper:lhcenter-late-cross-linear-step-one-at-shift-of-zero-diag-and-winwittagree-zero-gap}
For every natural number \(Q\) and natural number \(t\) and natural number \(k\) and natural number \(r\) and natural number \(a\) and natural number \(L\), if \(t + Q \le r + 2\), then if \(t + Q + r + 3 \le k\), then if \(\operatorname{zerodiagonal}(t + Q + r + 3)\), then if \(\operatorname{settledby}(t + Q + r + 3, t + Q + \operatorname{twopow}(k))\), then if \(a \le t + \operatorname{twopow}(k + 1) - (r + 3)\), then if \(t + \operatorname{twopow}(k + 1) - (r + 3) \le a + L\), then if \(\operatorname{winwittagree}(t, k, Q, a, L)\), then \(\operatorname{lhcenter}(t + 1 + \operatorname{twopow}(k), \operatorname{twopow}(k) - (r + 2)) = \operatorname{bxor}(\operatorname{lhcenter}(t + \operatorname{twopow}(k), \operatorname{twopow}(k) - (r + 1)), \operatorname{lhcenter}(t + \operatorname{twopow}(k), \operatorname{twopow}(k) - (r + 2)))\).
\end{theorem}
\begin{proof}
The goal follows from \(\operatorname{lhcenterlatecrosslinearsteponeatshiftofzerodiagandwinwittagreezerogap}(Q, t, k, r, a, L)\).
\end{proof}
\begin{theorem}
\label{ll:paper:lhcenter-late-cross-linear-step-two-at-shift-of-zero-diag-and-winwittagree-zero-gap}
For every natural number \(Q\) and natural number \(t\) and natural number \(k\) and natural number \(r\) and natural number \(a\) and natural number \(L\), if \(t + Q \le r + 1\), then if \(t + Q + r + 3 \le k\), then if \(\operatorname{zerodiagonal}(t + Q + r + 2)\), then if \(\operatorname{settledby}(t + Q + r + 2, t + Q + \operatorname{twopow}(k))\), then if \(a \le t + \operatorname{twopow}(k + 1) - (r + 2)\), then if \(t + \operatorname{twopow}(k + 1) - (r + 2) \le a + L\), then if \(\operatorname{winwittagree}(t, k, Q, a, L)\), then \(\operatorname{lhcenter}(t + 1 + \operatorname{twopow}(k), \operatorname{twopow}(k) - (r + 2)) = \operatorname{bxor}(\operatorname{lhcenter}(t + \operatorname{twopow}(k), \operatorname{twopow}(k) - (r + 1)), \operatorname{lhcenter}(t + \operatorname{twopow}(k), \operatorname{twopow}(k) - (r + 3)))\).
\end{theorem}
\begin{proof}
The goal follows from \(\operatorname{lhcenterlatecrosslinearsteptwoatshiftofzerodiagandwinwittagreezerogap}(Q, t, k, r, a, L)\).
\end{proof}
\begin{theorem}
\label{ll:paper:lhcenter-late-cross-step-one-eq-false-at-shift-of-zero-diag-and-single-winwittagree-zero-gap}
For every natural number \(Q\) and natural number \(t\) and natural number \(k\) and natural number \(r\) and natural number \(a\) and natural number \(L\), if \(t + Q \le r\), then if \(t + Q + r + 3 \le k\), then if \(\operatorname{zerodiagonal}(t + Q + r + 3)\), then if \(\operatorname{settledby}(t + Q + r + 3, t + Q + \operatorname{twopow}(k))\), then if \(a \le t + \operatorname{twopow}(k + 1) - (r + 3)\), then if \(t + \operatorname{twopow}(k + 1) - (r + 1) \le a + L\), then if \(\operatorname{winwittagree}(t, k, Q, a, L)\), then \(\operatorname{lhcenter}(t + 1 + \operatorname{twopow}(k), \operatorname{twopow}(k) - (r + 2)) = \operatorname{false}\).
\end{theorem}
\begin{proof}
The goal follows from \(\operatorname{lhcenterlatecrosssteponeeqfalseatshiftofzerodiagandsinglewinwittagreezerogap}(Q, t, k, r, a, L)\).
\end{proof}
\begin{theorem}
\label{ll:paper:winwitt-scale-cross-zero-gap-step-one-eq-false-at-shift-of-zero-diag-and-single-winwittagree}
For every natural number \(Q\) and natural number \(t\) and natural number \(k\) and natural number \(r\) and natural number \(a\) and natural number \(L\), if \(t + Q \le r\), then if \(t + Q + r + 3 \le k\), then if \(\operatorname{zerodiagonal}(t + Q + r + 3)\), then if \(\operatorname{settledby}(t + Q + r + 3, t + Q + \operatorname{twopow}(k))\), then if \(a \le t + \operatorname{twopow}(k + 1) - (r + 3)\), then if \(t + \operatorname{twopow}(k + 1) - (r + 1) \le a + L\), then if \(\operatorname{winwittagree}(t, k, Q, a, L)\), then \(\operatorname{wittbitbool}(\operatorname{wittwin}(t + 1, \operatorname{twopow}(k)), t + \operatorname{twopow}(k + 1) - (r + 1)) = \operatorname{false}\).
\end{theorem}
\begin{proof}
The goal follows from \(\operatorname{winwittscalecrosszerogapsteponeeqfalseatshiftofzerodiagandsinglewinwittagree}(Q, t, k, r, a, L)\).
\end{proof}
\begin{theorem}
\label{ll:paper:lhcenter-late-cross-step-one-eq-of-winwittagree-three-zero-gap}
For every natural number \(Q\) and natural number \(t\) and natural number \(k\) and natural number \(r\) and natural number \(a\) and natural number \(L\), if \(t + Q \le r\), then if \(t + Q + r + 3 < \operatorname{twopow}(k)\), then if \(a \le t + \operatorname{twopow}(k + 1) - (r + 3)\), then if \(t + \operatorname{twopow}(k + 1) - (r + 1) \le a + L\), then if \(\operatorname{winwittagree}(t, k, Q, a, L)\), then \(\operatorname{lhcenter}(t + 1 + Q + \operatorname{twopow}(k), \operatorname{twopow}(k) - (r + 2)) = \operatorname{lhcenter}(t + 1 + \operatorname{twopow}(k), \operatorname{twopow}(k) - (r + 2))\).
\end{theorem}
\begin{proof}
The goal follows from \(\operatorname{lhcenterlatecrosssteponeeqofwinwittagreethreezerogap}(Q, t, k, r, a, L)\).
\end{proof}
\begin{theorem}
\label{ll:paper:winwitt-scale-cross-zero-gap-step-one-eq-of-winwittagree-three}
For every natural number \(Q\) and natural number \(t\) and natural number \(k\) and natural number \(r\) and natural number \(a\) and natural number \(L\), if \(t + Q \le r\), then if \(t + Q + r + 3 < \operatorname{twopow}(k)\), then if \(a \le t + \operatorname{twopow}(k + 1) - (r + 3)\), then if \(t + \operatorname{twopow}(k + 1) - (r + 1) \le a + L\), then if \(\operatorname{winwittagree}(t, k, Q, a, L)\), then \(\operatorname{wittbitbool}(\operatorname{wittwin}(t + 1 + Q, \operatorname{twopow}(k)), t + Q + \operatorname{twopow}(k + 1) - (r + 1)) = \operatorname{wittbitbool}(\operatorname{wittwin}(t + 1, \operatorname{twopow}(k)), t + \operatorname{twopow}(k + 1) - (r + 1))\).
\end{theorem}
\begin{proof}
The goal follows from \(\operatorname{winwittscalecrosszerogapsteponeeqofwinwittagreethree}(Q, t, k, r, a, L)\).
\end{proof}
\begin{theorem}
\label{ll:paper:lhcenter-late-cross-step-one-shift-flip-on-doubling-diag-of-doubles-and-single-winwittagree-zero-gap}
For every natural number \(Q\) and natural number \(t\) and natural number \(k\) and natural number \(r\) and natural number \(a\) and natural number \(L\), if \(t + Q \le r\), then if \(t + Q + r + 3 \le k\), then if \(\operatorname{zerodiagonal}(t + Q + r + 2)\), then if \(\operatorname{doublesat}(Q, t + Q + r + 2)\), then if \(\operatorname{settledby}(t + Q + r + 1, t + \operatorname{twopow}(k))\), then if \(\operatorname{settledby}(t + Q + r + 2, t + \operatorname{twopow}(k))\), then if \(\operatorname{settledby}(t + Q + r + 3, t + \operatorname{twopow}(k))\), then if \(a \le t + \operatorname{twopow}(k + 1) - (r + 3)\), then if \(t + \operatorname{twopow}(k + 1) - (r + 1) \le a + L\), then if \(\operatorname{winwittagree}(t, k, Q, a, L)\), then \(\operatorname{lhcenter}(t + 1 + \operatorname{twopow}(k), \operatorname{twopow}(k) - (r + Q + 2)) = \operatorname{not}(\operatorname{lhcenter}(t + 1 + \operatorname{twopow}(k), \operatorname{twopow}(k) - (r + 2)))\).
\end{theorem}
\begin{proof}
The goal follows from \(\operatorname{lhcenterlatecrosssteponeshiftflipondoublingdiagofdoublesandsinglewinwittagreezerogap}(Q, t, k, r, a, L)\).
\end{proof}
\begin{theorem}
\label{ll:paper:winwitt-scale-cross-zero-gap-step-one-shift-flip-on-doubling-diag-of-doubles-and-single-winwittagree}
For every natural number \(Q\) and natural number \(t\) and natural number \(k\) and natural number \(r\) and natural number \(a\) and natural number \(L\), if \(t + Q \le r\), then if \(t + Q + r + 3 \le k\), then if \(\operatorname{zerodiagonal}(t + Q + r + 2)\), then if \(\operatorname{doublesat}(Q, t + Q + r + 2)\), then if \(\operatorname{settledby}(t + Q + r + 1, t + \operatorname{twopow}(k))\), then if \(\operatorname{settledby}(t + Q + r + 2, t + \operatorname{twopow}(k))\), then if \(\operatorname{settledby}(t + Q + r + 3, t + \operatorname{twopow}(k))\), then if \(a \le t + \operatorname{twopow}(k + 1) - (r + 3)\), then if \(t + \operatorname{twopow}(k + 1) - (r + 1) \le a + L\), then if \(\operatorname{winwittagree}(t, k, Q, a, L)\), then \(\operatorname{wittbitbool}(\operatorname{wittwin}(t + 1, \operatorname{twopow}(k)), t + \operatorname{twopow}(k + 1) - (r + Q + 1)) = \operatorname{not}(\operatorname{wittbitbool}(\operatorname{wittwin}(t + 1, \operatorname{twopow}(k)), t + \operatorname{twopow}(k + 1) - (r + 1)))\).
\end{theorem}
\begin{proof}
The goal follows from \(\operatorname{winwittscalecrosszerogapsteponeshiftflipondoublingdiagofdoublesandsinglewinwittagree}(Q, t, k, r, a, L)\).
\end{proof}
\begin{theorem}
\label{ll:paper:lhcenter-late-cross-late-eq-false-at-shift-of-zero-diag-and-winwittagree-zero-gap}
For every natural number \(Q\) and natural number \(t\) and natural number \(k\) and natural number \(r\) and natural number \(a\) and natural number \(L\), if \(t < r\), then if \(t + Q + r + 1 \le k\), then if \(\operatorname{zerodiagonal}(t + Q + r)\), then if \(\operatorname{settledby}(t + Q + r, t + \operatorname{twopow}(k))\), then if \(a \le t + \operatorname{twopow}(k + 1) - (r + Q)\), then if \(t + \operatorname{twopow}(k + 1) - (r + Q) \le a + L\), then if \(\operatorname{winwittagree}(t, k, Q, a, L)\), then \(\operatorname{lhcenter}(t + Q + \operatorname{twopow}(k), \operatorname{twopow}(k) - (r + Q)) = \operatorname{false}\).
\end{theorem}
\begin{proof}
The goal follows from \(\operatorname{lhcenterlatecrosslateeqfalseatshiftofzerodiagandwinwittagreezerogap}(Q, t, k, r, a, L)\).
\end{proof}
\begin{theorem}
\label{ll:paper:lhcenter-late-cross-late-eq-true-at-shift-after-zero-diag-and-winwittagree-zero-gap}
For every natural number \(Q\) and natural number \(t\) and natural number \(k\) and natural number \(r\) and natural number \(a\) and natural number \(L\), if \(t \le r + 1\), then if \(t + Q + r + 2 \le k\), then if \(\operatorname{zerodiagonal}(t + Q + r)\), then if \(\operatorname{settledby}(t + Q + r + 2, t + \operatorname{twopow}(k))\), then if \(a \le t + \operatorname{twopow}(k + 1) - (r + Q + 2)\), then if \(t + \operatorname{twopow}(k + 1) - (r + Q + 2) \le a + L\), then if \(\operatorname{winwittagree}(t, k, Q, a, L)\), then \(\operatorname{lhcenter}(t + Q + \operatorname{twopow}(k), \operatorname{twopow}(k) - (r + Q + 2)) = \operatorname{true}\).
\end{theorem}
\begin{proof}
The goal follows from \(\operatorname{lhcenterlatecrosslateeqtrueatshiftafterzerodiagandwinwittagreezerogap}(Q, t, k, r, a, L)\).
\end{proof}
\begin{theorem}
\label{ll:paper:lhcenter-late-cross-late-step-flip-at-shift-after-zero-diag-of-single-winwittagree-zero-gap}
For every natural number \(Q\) and natural number \(t\) and natural number \(k\) and natural number \(r\) and natural number \(a\) and natural number \(L\), if \(t \le r + 1\), then if \(t + Q + r + 2 \le k\), then if \(\operatorname{zerodiagonal}(t + Q + r)\), then if \(\operatorname{settledby}(t + Q + r + 2, t + \operatorname{twopow}(k))\), then if \(a \le t + \operatorname{twopow}(k + 1) - (r + Q + 2)\), then if \(t + \operatorname{twopow}(k + 1) - (r + Q + 2) \le a + L\), then if \(\operatorname{winwittagree}(t, k, Q, a, L)\), then \(\operatorname{lhcenter}(t + 1 + Q + \operatorname{twopow}(k), \operatorname{twopow}(k) - (r + Q + 2)) = \operatorname{not}(\operatorname{lhcenter}(t + Q + \operatorname{twopow}(k), \operatorname{twopow}(k) - (r + Q + 1)))\).
\end{theorem}
\begin{proof}
The goal follows from \(\operatorname{lhcenterlatecrosslatestepflipatshiftafterzerodiagofsinglewinwittagreezerogap}(Q, t, k, r, a, L)\).
\end{proof}
\begin{theorem}
\label{ll:paper:winwitt-scale-cross-zero-gap-late-step-flip-at-shift-after-zero-diag-of-single-winwittagree}
For every natural number \(Q\) and natural number \(t\) and natural number \(k\) and natural number \(r\) and natural number \(a\) and natural number \(L\), if \(t \le r\), then if \(t + Q + r + 2 \le k\), then if \(\operatorname{zerodiagonal}(t + Q + r)\), then if \(\operatorname{settledby}(t + Q + r + 2, t + \operatorname{twopow}(k))\), then if \(a \le t + \operatorname{twopow}(k + 1) - (r + Q + 2)\), then if \(t + \operatorname{twopow}(k + 1) - (r + Q + 2) \le a + L\), then if \(\operatorname{winwittagree}(t, k, Q, a, L)\), then \(\operatorname{wittbitbool}(\operatorname{wittwin}(t + 1 + Q, \operatorname{twopow}(k)), t + \operatorname{twopow}(k + 1) - (r + 1)) = \operatorname{not}(\operatorname{wittbitbool}(\operatorname{wittwin}(t + Q, \operatorname{twopow}(k)), t + \operatorname{twopow}(k + 1) - (r + 1)))\).
\end{theorem}
\begin{proof}
The goal follows from \(\operatorname{winwittscalecrosszerogaplatestepflipatshiftafterzerodiagofsinglewinwittagree}(Q, t, k, r, a, L)\).
\end{proof}
\begin{theorem}
\label{ll:paper:not-winwittagree-of-zero-diag-one-and-shift-three}
For every natural number \(Q\) and natural number \(t\) and natural number \(k\) and natural number \(r\) and natural number \(a\) and natural number \(L\), if \(t + Q \le r\), then if \(t + Q + r + 3 \le k\), then if \(\operatorname{zerodiagonal}(t + r + 1)\), then if \(\operatorname{zerodiagonal}(t + Q + r + 3)\), then if \(\operatorname{settledby}(t + r + 1 + 2, t + 1 + \operatorname{twopow}(k))\), then if \(\operatorname{settledby}(t + Q + r + 3, t + Q + \operatorname{twopow}(k))\), then if \(a \le t + \operatorname{twopow}(k + 1) - (r + 3)\), then if \(t + \operatorname{twopow}(k + 1) - (r + 1) \le a + L\), then not \(\operatorname{winwittagree}(t, k, Q, a, L)\).
\end{theorem}
\begin{proof}
The goal follows from \(\operatorname{notwinwittagreeofzerodiagoneandshiftthree}(Q, t, k, r, a, L)\).
\end{proof}
\begin{theorem}
\label{ll:paper:zero-diag-ge-two}
For every natural number \(z\), if \(\operatorname{zerodiagonal}(z)\), then \(2 \le z\).
\end{theorem}
\begin{proof}
The goal follows from \(\operatorname{zerodiaggetwo}(z)\).
\end{proof}
\begin{theorem}
\label{ll:paper:not-iszerodiag-add-two}
For every natural number \(z\), if \(\operatorname{zerodiagonal}(z)\), then not \(\operatorname{zerodiagonal}(z + 2)\).
\end{theorem}
\begin{proof}
The goal follows from \(\operatorname{notiszerodiagaddtwo}(z)\).
\end{proof}
\begin{theorem}
\label{ll:paper:not-isblackdiag-add-one-of-zero-diag}
For every natural number \(z\), if \(\operatorname{zerodiagonal}(z)\), then not \(\operatorname{blackdiagonal}(z + 1)\).
\end{theorem}
\begin{proof}
The goal follows from \(\operatorname{notisblackdiagaddoneofzerodiag}(z)\).
\end{proof}
\begin{theorem}
\label{ll:paper:not-iszerodiag-add-three}
For every natural number \(z\), if \(\operatorname{zerodiagonal}(z)\), then not \(\operatorname{zerodiagonal}(z + 3)\).
\end{theorem}
\begin{proof}
The goal follows from \(\operatorname{notiszerodiagaddthree}(z)\).
\end{proof}
\begin{theorem}
\label{ll:paper:not-isblackdiag-add-three-of-zero-diag}
For every natural number \(z\), if \(\operatorname{zerodiagonal}(z)\), then not \(\operatorname{blackdiagonal}(z + 3)\).
\end{theorem}
\begin{proof}
The goal follows from \(\operatorname{notisblackdiagaddthreeofzerodiag}(z)\).
\end{proof}
\begin{theorem}
\label{ll:paper:not-iszerodiag-add-four}
For every natural number \(z\), if \(\operatorname{zerodiagonal}(z)\), then not \(\operatorname{zerodiagonal}(z + 4)\).
\end{theorem}
\begin{proof}
The goal follows from \(\operatorname{notiszerodiagaddfour}(z)\).
\end{proof}
\begin{theorem}
\label{ll:paper:zero-diag-gap-ge-five}
For every natural number \(z1\) and natural number \(z2\), if \(\operatorname{zerodiagonal}(z1)\), then if \(\operatorname{zerodiagonal}(z2)\), then if \(z1 < z2\), then \(z1 + 5 \le z2\).
\end{theorem}
\begin{proof}
The goal follows from \(\operatorname{zerodiaggapgefive}(z1, z2)\).
\end{proof}
\begin{theorem}
\label{ll:paper:winwitt-scale-cross-zero-gap-late-eq-false-at-shift-of-zero-diag-and-winwittagree}
For every natural number \(Q\) and natural number \(t\) and natural number \(k\) and natural number \(r\) and natural number \(a\) and natural number \(L\), if \(t < r\), then if \(t + Q + r + 1 \le k\), then if \(\operatorname{zerodiagonal}(t + Q + r)\), then if \(\operatorname{settledby}(t + Q + r, t + \operatorname{twopow}(k))\), then if \(a \le t + \operatorname{twopow}(k + 1) - (r + Q)\), then if \(t + \operatorname{twopow}(k + 1) - (r + Q) \le a + L\), then if \(\operatorname{winwittagree}(t, k, Q, a, L)\), then \(\operatorname{wittbitbool}(\operatorname{wittwin}(t + Q, \operatorname{twopow}(k)), t + \operatorname{twopow}(k + 1) - r) = \operatorname{false}\).
\end{theorem}
\begin{proof}
The goal follows from \(\operatorname{winwittscalecrosszerogaplateeqfalseatshiftofzerodiagandwinwittagree}(Q, t, k, r, a, L)\).
\end{proof}
\begin{theorem}
\label{ll:paper:winwitt-scale-cross-zero-gap-late-eq-true-at-shift-after-zero-diag-and-winwittagree}
For every natural number \(Q\) and natural number \(t\) and natural number \(k\) and natural number \(r\) and natural number \(a\) and natural number \(L\), if \(t \le r + 1\), then if \(t + Q + r + 2 \le k\), then if \(\operatorname{zerodiagonal}(t + Q + r)\), then if \(\operatorname{settledby}(t + Q + r + 2, t + \operatorname{twopow}(k))\), then if \(a \le t + \operatorname{twopow}(k + 1) - (r + Q + 2)\), then if \(t + \operatorname{twopow}(k + 1) - (r + Q + 2) \le a + L\), then if \(\operatorname{winwittagree}(t, k, Q, a, L)\), then \(\operatorname{wittbitbool}(\operatorname{wittwin}(t + Q, \operatorname{twopow}(k)), t + \operatorname{twopow}(k + 1) - (r + 2)) = \operatorname{true}\).
\end{theorem}
\begin{proof}
The goal follows from \(\operatorname{winwittscalecrosszerogaplateeqtrueatshiftafterzerodiagandwinwittagree}(Q, t, k, r, a, L)\).
\end{proof}
\begin{theorem}
\label{ll:paper:lhcenter-late-cross-late-step-one-eq-true-at-shift-of-zero-diag-and-single-winwittagree-zero-gap}
For every natural number \(Q\) and natural number \(t\) and natural number \(k\) and natural number \(r\) and natural number \(a\) and natural number \(L\), if \(t < r\), then if \(t + Q + r + 2 \le k\), then if \(\operatorname{zerodiagonal}(t + Q + r)\), then if \(\operatorname{settledby}(t + Q + r, t + \operatorname{twopow}(k))\), then if \(\operatorname{settledby}(t + Q + r + 2, t + \operatorname{twopow}(k))\), then if \(a \le t + \operatorname{twopow}(k + 1) - (r + Q + 2)\), then if \(t + \operatorname{twopow}(k + 1) - (r + Q) \le a + L\), then if \(\operatorname{winwittagree}(t, k, Q, a, L)\), then \(\operatorname{lhcenter}(t + 1 + Q + \operatorname{twopow}(k), \operatorname{twopow}(k) - (r + Q + 1)) = \operatorname{true}\).
\end{theorem}
\begin{proof}
The goal follows from \(\operatorname{lhcenterlatecrosslatesteponeeqtrueatshiftofzerodiagandsinglewinwittagreezerogap}(Q, t, k, r, a, L)\).
\end{proof}
\begin{theorem}
\label{ll:paper:winwitt-scale-cross-zero-gap-late-step-one-eq-true-at-shift-of-zero-diag-and-single-winwittagree}
For every natural number \(Q\) and natural number \(t\) and natural number \(k\) and natural number \(r\) and natural number \(a\) and natural number \(L\), if \(t < r\), then if \(t + Q + r + 2 \le k\), then if \(\operatorname{zerodiagonal}(t + Q + r)\), then if \(\operatorname{settledby}(t + Q + r, t + \operatorname{twopow}(k))\), then if \(\operatorname{settledby}(t + Q + r + 2, t + \operatorname{twopow}(k))\), then if \(a \le t + \operatorname{twopow}(k + 1) - (r + Q + 2)\), then if \(t + \operatorname{twopow}(k + 1) - (r + Q) \le a + L\), then if \(\operatorname{winwittagree}(t, k, Q, a, L)\), then \(\operatorname{wittbitbool}(\operatorname{wittwin}(t + 1 + Q, \operatorname{twopow}(k)), t + \operatorname{twopow}(k + 1) - r) = \operatorname{true}\).
\end{theorem}
\begin{proof}
The goal follows from \(\operatorname{winwittscalecrosszerogaplatesteponeeqtrueatshiftofzerodiagandsinglewinwittagree}(Q, t, k, r, a, L)\).
\end{proof}
\begin{theorem}
\label{ll:paper:lhcenter-late-cross-step-one-pin-r-plus-three-true-of-doubles-three-and-single-winwittagree}
For every natural number \(t\) and natural number \(k\) and natural number \(r\) and natural number \(a\) and natural number \(L\), if \(t + 3 \le r\), then if \(t + 3 + r + 2 \le k\), then if \(\operatorname{zerodiagonal}(t + 3 + r)\), then if \(\operatorname{doublesat}(3, t + 3 + r)\), then if \(\operatorname{settledby}(t + 3 + r, t + \operatorname{twopow}(k))\), then if \(\operatorname{settledby}(t + 3 + r + 1, t + \operatorname{twopow}(k))\), then if \(\operatorname{settledby}(t + 3 + r + 2, t + \operatorname{twopow}(k))\), then if \(a \le t + \operatorname{twopow}(k + 1) - (r + 2)\), then if \(t + \operatorname{twopow}(k + 1) - (r + 1) \le a + L\), then if \(\operatorname{winwittagree}(t, k, 3, a, L)\), then \(\operatorname{lhcenter}(t + 1 + \operatorname{twopow}(k), \operatorname{twopow}(k) - (r + 3)) = \operatorname{true}\).
\end{theorem}
\begin{proof}
The goal follows from \(\operatorname{lhcenterlatecrosssteponepinrplusthreetrueofdoublesthreeandsinglewinwittagree}(t, k, r, a, L)\).
\end{proof}
\begin{theorem}
\label{ll:paper:lhcenter-late-cross-step-two-pin-r-plus-two-false-of-doubles-three-and-single-winwittagree}
For every natural number \(t\) and natural number \(k\) and natural number \(r\) and natural number \(a\) and natural number \(L\), if \(t + 3 < r\), then if \(t + 3 + r + 2 \le k\), then if \(\operatorname{zerodiagonal}(t + 3 + r)\), then if \(\operatorname{doublesat}(3, t + 3 + r)\), then if \(\operatorname{settledby}(t + 3 + r, t + \operatorname{twopow}(k))\), then if \(\operatorname{settledby}(t + 3 + r + 1, t + \operatorname{twopow}(k))\), then if \(\operatorname{settledby}(t + 3 + r + 2, t + \operatorname{twopow}(k))\), then if \(a \le t + \operatorname{twopow}(k + 1) - (r + 2)\), then if \(t + \operatorname{twopow}(k + 1) - r \le a + L\), then if \(\operatorname{winwittagree}(t, k, 3, a, L)\), then \(\operatorname{lhcenter}(t + 2 + \operatorname{twopow}(k), \operatorname{twopow}(k) - (r + 2)) = \operatorname{false}\).
\end{theorem}
\begin{proof}
The goal follows from \(\operatorname{lhcenterlatecrosssteptwopinrplustwofalseofdoublesthreeandsinglewinwittagree}(t, k, r, a, L)\).
\end{proof}
\begin{theorem}
\label{ll:paper:lhcenter-late-cross-step-three-pin-r-plus-three-true-of-doubles-three-and-single-winwittagree}
For every natural number \(t\) and natural number \(k\) and natural number \(r\) and natural number \(a\) and natural number \(L\), if \(t + 3 < r\), then if \(t + 3 + r + 2 \le k\), then if \(\operatorname{zerodiagonal}(t + 3 + r)\), then if \(\operatorname{doublesat}(3, t + 3 + r)\), then if \(\operatorname{settledby}(t + 3 + r, t + \operatorname{twopow}(k))\), then if \(\operatorname{settledby}(t + 3 + r + 1, t + \operatorname{twopow}(k))\), then if \(\operatorname{settledby}(t + 3 + r + 2, t + \operatorname{twopow}(k))\), then if \(a \le t + \operatorname{twopow}(k + 1) - (r + 2)\), then if \(t + \operatorname{twopow}(k + 1) - r \le a + L\), then if \(\operatorname{winwittagree}(t, k, 3, a, L)\), then \(\operatorname{lhcenter}(t + 3 + \operatorname{twopow}(k), \operatorname{twopow}(k) - (r + 3)) = \operatorname{true}\).
\end{theorem}
\begin{proof}
The goal follows from \(\operatorname{lhcenterlatecrossstepthreepinrplusthreetrueofdoublesthreeandsinglewinwittagree}(t, k, r, a, L)\).
\end{proof}
\begin{theorem}
\label{ll:paper:not-winwittagree-three-of-doubles}
For every natural number \(t\) and natural number \(k\) and natural number \(r\) and natural number \(a\) and natural number \(L\), if \(t + 3 < r\), then if \(t + 3 + r + 2 \le k\), then if \(\operatorname{zerodiagonal}(t + 3 + r)\), then if \(\operatorname{doublesat}(3, t + 3 + r)\), then if \(\operatorname{settledby}(t + 3 + r, t + \operatorname{twopow}(k))\), then if \(\operatorname{settledby}(t + 3 + r + 1, t + \operatorname{twopow}(k))\), then if \(\operatorname{settledby}(t + 3 + r + 2, t + \operatorname{twopow}(k))\), then if \(a \le t + \operatorname{twopow}(k + 1) - (r + 3)\), then if \(t + \operatorname{twopow}(k + 1) - r \le a + L\), then not \(\operatorname{winwittagree}(t, k, 3, a, L)\).
\end{theorem}
\begin{proof}
The goal follows from \(\operatorname{notwinwittagreethreeofdoubles}(t, k, r, a, L)\).
\end{proof}
\begin{theorem}
\label{ll:paper:winwitt-scale-cross-zero-gap-step-one-pin-r-plus-two-true-of-doubles-three-and-single-winwittagree}
For every natural number \(t\) and natural number \(k\) and natural number \(r\) and natural number \(a\) and natural number \(L\), if \(t + 3 \le r\), then if \(t + 3 + r + 2 \le k\), then if \(\operatorname{zerodiagonal}(t + 3 + r)\), then if \(\operatorname{doublesat}(3, t + 3 + r)\), then if \(\operatorname{settledby}(t + 3 + r, t + \operatorname{twopow}(k))\), then if \(\operatorname{settledby}(t + 3 + r + 1, t + \operatorname{twopow}(k))\), then if \(\operatorname{settledby}(t + 3 + r + 2, t + \operatorname{twopow}(k))\), then if \(a \le t + \operatorname{twopow}(k + 1) - (r + 2)\), then if \(t + \operatorname{twopow}(k + 1) - (r + 1) \le a + L\), then if \(\operatorname{winwittagree}(t, k, 3, a, L)\), then \(\operatorname{wittbitbool}(\operatorname{wittwin}(t + 1, \operatorname{twopow}(k)), t + \operatorname{twopow}(k + 1) - (r + 2)) = \operatorname{true}\).
\end{theorem}
\begin{proof}
The goal follows from \(\operatorname{winwittscalecrosszerogapsteponepinrplustwotrueofdoublesthreeandsinglewinwittagree}(t, k, r, a, L)\).
\end{proof}
\begin{theorem}
\label{ll:paper:winwitt-scale-cross-zero-gap-step-two-pin-r-false-of-doubles-three-and-single-winwittagree}
For every natural number \(t\) and natural number \(k\) and natural number \(r\) and natural number \(a\) and natural number \(L\), if \(t + 3 < r\), then if \(t + 3 + r + 2 \le k\), then if \(\operatorname{zerodiagonal}(t + 3 + r)\), then if \(\operatorname{doublesat}(3, t + 3 + r)\), then if \(\operatorname{settledby}(t + 3 + r, t + \operatorname{twopow}(k))\), then if \(\operatorname{settledby}(t + 3 + r + 1, t + \operatorname{twopow}(k))\), then if \(\operatorname{settledby}(t + 3 + r + 2, t + \operatorname{twopow}(k))\), then if \(a \le t + \operatorname{twopow}(k + 1) - (r + 2)\), then if \(t + \operatorname{twopow}(k + 1) - r \le a + L\), then if \(\operatorname{winwittagree}(t, k, 3, a, L)\), then \(\operatorname{wittbitbool}(\operatorname{wittwin}(t + 2, \operatorname{twopow}(k)), t + \operatorname{twopow}(k + 1) - r) = \operatorname{false}\).
\end{theorem}
\begin{proof}
The goal follows from \(\operatorname{winwittscalecrosszerogapsteptwopinrfalseofdoublesthreeandsinglewinwittagree}(t, k, r, a, L)\).
\end{proof}
\begin{theorem}
\label{ll:paper:winwitt-scale-cross-zero-gap-step-three-pin-r-true-of-doubles-three-and-single-winwittagree}
For every natural number \(t\) and natural number \(k\) and natural number \(r\) and natural number \(a\) and natural number \(L\), if \(t + 3 < r\), then if \(t + 3 + r + 2 \le k\), then if \(\operatorname{zerodiagonal}(t + 3 + r)\), then if \(\operatorname{doublesat}(3, t + 3 + r)\), then if \(\operatorname{settledby}(t + 3 + r, t + \operatorname{twopow}(k))\), then if \(\operatorname{settledby}(t + 3 + r + 1, t + \operatorname{twopow}(k))\), then if \(\operatorname{settledby}(t + 3 + r + 2, t + \operatorname{twopow}(k))\), then if \(a \le t + \operatorname{twopow}(k + 1) - (r + 2)\), then if \(t + \operatorname{twopow}(k + 1) - r \le a + L\), then if \(\operatorname{winwittagree}(t, k, 3, a, L)\), then \(\operatorname{wittbitbool}(\operatorname{wittwin}(t + 3, \operatorname{twopow}(k)), t + \operatorname{twopow}(k + 1) - r) = \operatorname{true}\).
\end{theorem}
\begin{proof}
The goal follows from \(\operatorname{winwittscalecrosszerogapstepthreepinrtrueofdoublesthreeandsinglewinwittagree}(t, k, r, a, L)\).
\end{proof}
\begin{theorem}
\label{ll:paper:winwitt-scale-cross-zero-gap-same-bit-r-flip-three-of-doubles-three-and-single-winwittagree}
For every natural number \(t\) and natural number \(k\) and natural number \(r\) and natural number \(a\) and natural number \(L\), if \(t + 3 < r\), then if \(t + 3 + r + 2 \le k\), then if \(\operatorname{zerodiagonal}(t + 3 + r)\), then if \(\operatorname{doublesat}(3, t + 3 + r)\), then if \(\operatorname{settledby}(t + 3 + r, t + \operatorname{twopow}(k))\), then if \(\operatorname{settledby}(t + 3 + r + 1, t + \operatorname{twopow}(k))\), then if \(\operatorname{settledby}(t + 3 + r + 2, t + \operatorname{twopow}(k))\), then if \(a \le t + \operatorname{twopow}(k + 1) - (r + 2)\), then if \(t + \operatorname{twopow}(k + 1) - r \le a + L\), then if \(\operatorname{winwittagree}(t, k, 3, a, L)\), then \(\operatorname{wittbitbool}(\operatorname{wittwin}(t + 3, \operatorname{twopow}(k)), t + \operatorname{twopow}(k + 1) - r) = \operatorname{not}(\operatorname{wittbitbool}(\operatorname{wittwin}(t, \operatorname{twopow}(k)), t + \operatorname{twopow}(k + 1) - r))\).
\end{theorem}
\begin{proof}
The goal follows from \(\operatorname{winwittscalecrosszerogapsamebitrflipthreeofdoublesthreeandsinglewinwittagree}(t, k, r, a, L)\).
\end{proof}
\begin{theorem}
\label{ll:paper:lhcenter-late-cross-step-one-pin-r-plus-two-true-of-single-winwittagree-and-zero-diag-four}
For every natural number \(Q\) and natural number \(t\) and natural number \(k\) and natural number \(r\) and natural number \(a\) and natural number \(L\), if \(t + Q < r\), then if \(t + Q + r + 2 \le k\), then if \(t + r + 4 \le k\), then if \(\operatorname{zerodiagonal}(t + Q + r)\), then if \(\operatorname{zerodiagonal}(t + r + 4)\), then if \(\operatorname{settledby}(t + Q + r, t + Q + \operatorname{twopow}(k))\), then if \(\operatorname{settledby}(t + Q + r + 2, t + Q + \operatorname{twopow}(k))\), then if \(\operatorname{settledby}(t + r + 4, t + 1 + \operatorname{twopow}(k))\), then if \(a \le t + \operatorname{twopow}(k + 1) - (r + 2)\), then if \(t + \operatorname{twopow}(k + 1) - r \le a + L\), then if \(\operatorname{winwittagree}(t, k, Q, a, L)\), then \(\operatorname{lhcenter}(t + 1 + \operatorname{twopow}(k), \operatorname{twopow}(k) - (r + 2)) = \operatorname{true}\).
\end{theorem}
\begin{proof}
The goal follows from \(\operatorname{lhcenterlatecrosssteponepinrplustwotrueofsinglewinwittagreeandzerodiagfour}(Q, t, k, r, a, L)\).
\end{proof}
\begin{theorem}
\label{ll:paper:lhcenter-late-cross-pin-r-plus-one-false-of-single-winwittagree-and-zero-diag-four}
For every natural number \(Q\) and natural number \(t\) and natural number \(k\) and natural number \(r\) and natural number \(a\) and natural number \(L\), if \(t + Q < r\), then if \(t + Q + r + 2 \le k\), then if \(t + r + 4 \le k\), then if \(\operatorname{zerodiagonal}(t + Q + r)\), then if \(\operatorname{zerodiagonal}(t + r + 4)\), then if \(\operatorname{settledby}(t + Q + r, t + Q + \operatorname{twopow}(k))\), then if \(\operatorname{settledby}(t + Q + r + 2, t + Q + \operatorname{twopow}(k))\), then if \(\operatorname{settledby}(t + r + 4, t + 1 + \operatorname{twopow}(k))\), then if \(a \le t + \operatorname{twopow}(k + 1) - (r + 2)\), then if \(t + \operatorname{twopow}(k + 1) - r \le a + L\), then if \(\operatorname{winwittagree}(t, k, Q, a, L)\), then \(\operatorname{lhcenter}(t + \operatorname{twopow}(k), \operatorname{twopow}(k) - (r + 1)) = \operatorname{false}\).
\end{theorem}
\begin{proof}
The goal follows from \(\operatorname{lhcenterlatecrosspinrplusonefalseofsinglewinwittagreeandzerodiagfour}(Q, t, k, r, a, L)\).
\end{proof}
\begin{theorem}
\label{ll:paper:lhcenter-late-cross-pin-doubling-diag-true-of-doubles-and-single-winwittagree-and-zero-diag-four}
For every natural number \(Q\) and natural number \(t\) and natural number \(k\) and natural number \(r\) and natural number \(a\) and natural number \(L\), if \(t + Q < r\), then if \(t + Q + r + 2 \le k\), then if \(t + r + 4 \le k\), then if \(\operatorname{zerodiagonal}(t + Q + r)\), then if \(\operatorname{doublesat}(Q, t + Q + r)\), then if \(\operatorname{zerodiagonal}(t + r + 4)\), then if \(\operatorname{settledby}(t + Q + r, t + Q + \operatorname{twopow}(k))\), then if \(\operatorname{settledby}(t + Q + r + 1, t + \operatorname{twopow}(k))\), then if \(\operatorname{settledby}(t + Q + r + 2, t + Q + \operatorname{twopow}(k))\), then if \(\operatorname{settledby}(t + r + 4, t + 1 + \operatorname{twopow}(k))\), then if \(a \le t + \operatorname{twopow}(k + 1) - (r + 2)\), then if \(t + \operatorname{twopow}(k + 1) - r \le a + L\), then if \(\operatorname{winwittagree}(t, k, Q, a, L)\), then \(\operatorname{lhcenter}(t + \operatorname{twopow}(k), \operatorname{twopow}(k) - (r + Q + 1)) = \operatorname{true}\).
\end{theorem}
\begin{proof}
The goal follows from \(\operatorname{lhcenterlatecrosspindoublingdiagtrueofdoublesandsinglewinwittagreeandzerodiagfour}(Q, t, k, r, a, L)\).
\end{proof}
\begin{theorem}
\label{ll:paper:lhcenter-late-cross-pin-r-plus-three-true-of-zero-diag-four-and-single-winwittagree}
For every natural number \(t\) and natural number \(k\) and natural number \(r\) and natural number \(a\) and natural number \(L\), if \(t + 4 \le r + 1\), then if \(t + 4 + r + 2 \le k\), then if \(\operatorname{zerodiagonal}(t + 4 + r)\), then if \(\operatorname{settledby}(t + 4 + r, t + \operatorname{twopow}(k))\), then if \(\operatorname{settledby}(t + 4 + r + 2, t + 4 + \operatorname{twopow}(k))\), then if \(a \le t + \operatorname{twopow}(k + 1) - (r + 2)\), then if \(t + \operatorname{twopow}(k + 1) - (r + 2) \le a + L\), then if \(\operatorname{winwittagree}(t, k, 4, a, L)\), then \(\operatorname{lhcenter}(t + \operatorname{twopow}(k), \operatorname{twopow}(k) - (r + 3)) = \operatorname{true}\).
\end{theorem}
\begin{proof}
The goal follows from \(\operatorname{lhcenterlatecrosspinrplusthreetrueofzerodiagfourandsinglewinwittagree}(t, k, r, a, L)\).
\end{proof}
\begin{theorem}
\label{ll:paper:lhcenter-late-cross-step-one-pin-r-plus-two-true-of-zero-diag-four-and-single-winwittagree}
For every natural number \(t\) and natural number \(k\) and natural number \(r\) and natural number \(a\) and natural number \(L\), if \(t + 4 < r\), then if \(t + 4 + r + 2 \le k\), then if \(\operatorname{zerodiagonal}(t + 4 + r)\), then if \(\operatorname{settledby}(t + 4 + r, t + 1 + \operatorname{twopow}(k))\), then if \(\operatorname{settledby}(t + 4 + r + 2, t + 4 + \operatorname{twopow}(k))\), then if \(a \le t + \operatorname{twopow}(k + 1) - (r + 2)\), then if \(t + \operatorname{twopow}(k + 1) - r \le a + L\), then if \(\operatorname{winwittagree}(t, k, 4, a, L)\), then \(\operatorname{lhcenter}(t + 1 + \operatorname{twopow}(k), \operatorname{twopow}(k) - (r + 2)) = \operatorname{true}\).
\end{theorem}
\begin{proof}
The goal follows from \(\operatorname{lhcenterlatecrosssteponepinrplustwotrueofzerodiagfourandsinglewinwittagree}(t, k, r, a, L)\).
\end{proof}
\begin{theorem}
\label{ll:paper:lhcenter-late-cross-pin-r-plus-one-false-of-zero-diag-four-and-single-winwittagree}
For every natural number \(t\) and natural number \(k\) and natural number \(r\) and natural number \(a\) and natural number \(L\), if \(t + 4 < r\), then if \(t + 4 + r + 2 \le k\), then if \(\operatorname{zerodiagonal}(t + 4 + r)\), then if \(\operatorname{settledby}(t + 4 + r, t + 1 + \operatorname{twopow}(k))\), then if \(\operatorname{settledby}(t + 4 + r + 2, t + 4 + \operatorname{twopow}(k))\), then if \(a \le t + \operatorname{twopow}(k + 1) - (r + 2)\), then if \(t + \operatorname{twopow}(k + 1) - r \le a + L\), then if \(\operatorname{winwittagree}(t, k, 4, a, L)\), then \(\operatorname{lhcenter}(t + \operatorname{twopow}(k), \operatorname{twopow}(k) - (r + 1)) = \operatorname{false}\).
\end{theorem}
\begin{proof}
The goal follows from \(\operatorname{lhcenterlatecrosspinrplusonefalseofzerodiagfourandsinglewinwittagree}(t, k, r, a, L)\).
\end{proof}
\begin{theorem}
\label{ll:paper:lhcenter-late-cross-pin-doubling-diag-true-of-doubles-four-and-single-winwittagree}
For every natural number \(t\) and natural number \(k\) and natural number \(r\) and natural number \(a\) and natural number \(L\), if \(t + 4 < r\), then if \(t + 4 + r + 2 \le k\), then if \(\operatorname{zerodiagonal}(t + 4 + r)\), then if \(\operatorname{doublesat}(4, t + 4 + r)\), then if \(\operatorname{settledby}(t + 4 + r, t + 1 + \operatorname{twopow}(k))\), then if \(\operatorname{settledby}(t + 4 + r + 1, t + \operatorname{twopow}(k))\), then if \(\operatorname{settledby}(t + 4 + r + 2, t + 4 + \operatorname{twopow}(k))\), then if \(a \le t + \operatorname{twopow}(k + 1) - (r + 2)\), then if \(t + \operatorname{twopow}(k + 1) - r \le a + L\), then if \(\operatorname{winwittagree}(t, k, 4, a, L)\), then \(\operatorname{lhcenter}(t + \operatorname{twopow}(k), \operatorname{twopow}(k) - (r + 4 + 1)) = \operatorname{true}\).
\end{theorem}
\begin{proof}
The goal follows from \(\operatorname{lhcenterlatecrosspindoublingdiagtrueofdoublesfourandsinglewinwittagree}(t, k, r, a, L)\).
\end{proof}
\begin{theorem}
\label{ll:paper:lhcenter-late-cross-step-one-pin-r-plus-four-false-of-doubles-four-and-single-winwittagree}
For every natural number \(t\) and natural number \(k\) and natural number \(r\) and natural number \(a\) and natural number \(L\), if \(t + 4 < r\), then if \(t + 4 + r + 2 \le k\), then if \(\operatorname{zerodiagonal}(t + 4 + r)\), then if \(\operatorname{doublesat}(4, t + 4 + r)\), then if \(\operatorname{settledby}(t + 4 + r, t + \operatorname{twopow}(k))\), then if \(\operatorname{settledby}(t + 4 + r + 1, t + \operatorname{twopow}(k))\), then if \(\operatorname{settledby}(t + 4 + r + 2, t + 4 + \operatorname{twopow}(k))\), then if \(a \le t + \operatorname{twopow}(k + 1) - (r + 2)\), then if \(t + \operatorname{twopow}(k + 1) - r \le a + L\), then if \(\operatorname{winwittagree}(t, k, 4, a, L)\), then \(\operatorname{lhcenter}(t + 1 + \operatorname{twopow}(k), \operatorname{twopow}(k) - (r + 4)) = \operatorname{false}\).
\end{theorem}
\begin{proof}
The goal follows from \(\operatorname{lhcenterlatecrosssteponepinrplusfourfalseofdoublesfourandsinglewinwittagree}(t, k, r, a, L)\).
\end{proof}
\begin{theorem}
\label{ll:paper:lhcenter-late-cross-step-two-pin-r-plus-three-true-of-doubles-four-and-single-winwittagree}
For every natural number \(t\) and natural number \(k\) and natural number \(r\) and natural number \(a\) and natural number \(L\), if \(t + 4 < r\), then if \(t + 4 + r + 2 \le k\), then if \(\operatorname{zerodiagonal}(t + 4 + r)\), then if \(\operatorname{doublesat}(4, t + 4 + r)\), then if \(\operatorname{settledby}(t + 4 + r, t + \operatorname{twopow}(k))\), then if \(\operatorname{settledby}(t + 4 + r + 1, t + \operatorname{twopow}(k))\), then if \(\operatorname{settledby}(t + 4 + r + 2, t + 4 + \operatorname{twopow}(k))\), then if \(a \le t + \operatorname{twopow}(k + 1) - (r + 2)\), then if \(t + \operatorname{twopow}(k + 1) - r \le a + L\), then if \(\operatorname{winwittagree}(t, k, 4, a, L)\), then \(\operatorname{lhcenter}(t + 2 + \operatorname{twopow}(k), \operatorname{twopow}(k) - (r + 3)) = \operatorname{true}\).
\end{theorem}
\begin{proof}
The goal follows from \(\operatorname{lhcenterlatecrosssteptwopinrplusthreetrueofdoublesfourandsinglewinwittagree}(t, k, r, a, L)\).
\end{proof}
\begin{theorem}
\label{ll:paper:lhcenter-late-cross-step-three-pin-r-plus-four-false-of-doubles-four-and-single-winwittagree}
For every natural number \(t\) and natural number \(k\) and natural number \(r\) and natural number \(a\) and natural number \(L\), if \(t + 4 < r\), then if \(t + 4 + r + 2 \le k\), then if \(\operatorname{zerodiagonal}(t + 4 + r)\), then if \(\operatorname{doublesat}(4, t + 4 + r)\), then if \(\operatorname{settledby}(t + 4 + r, t + \operatorname{twopow}(k))\), then if \(\operatorname{settledby}(t + 4 + r + 1, t + \operatorname{twopow}(k))\), then if \(\operatorname{settledby}(t + 4 + r + 2, t + 2 + \operatorname{twopow}(k))\), then if \(a \le t + \operatorname{twopow}(k + 1) - (r + 2)\), then if \(t + \operatorname{twopow}(k + 1) - r \le a + L\), then if \(\operatorname{winwittagree}(t, k, 4, a, L)\), then \(\operatorname{lhcenter}(t + 3 + \operatorname{twopow}(k), \operatorname{twopow}(k) - (r + 4)) = \operatorname{false}\).
\end{theorem}
\begin{proof}
The goal follows from \(\operatorname{lhcenterlatecrossstepthreepinrplusfourfalseofdoublesfourandsinglewinwittagree}(t, k, r, a, L)\).
\end{proof}
\begin{theorem}
\label{ll:paper:lhcenter-late-cross-step-three-pin-r-plus-two-false-of-zero-diag-four-and-single-winwittagree}
For every natural number \(t\) and natural number \(k\) and natural number \(r\) and natural number \(a\) and natural number \(L\), if \(t + 4 \le r + 1\), then if \(t + 4 + r + 2 \le k\), then if \(\operatorname{zerodiagonal}(t + 4 + r)\), then if \(\operatorname{settledby}(t + 4 + r, t + \operatorname{twopow}(k))\), then if \(\operatorname{settledby}(t + 4 + r + 2, t + 3 + \operatorname{twopow}(k))\), then if \(a \le t + \operatorname{twopow}(k + 1) - (r + 3)\), then if \(t + \operatorname{twopow}(k + 1) - (r + 2) \le a + L\), then if \(\operatorname{winwittagree}(t, k, 4, a, L)\), then \(\operatorname{lhcenter}(t + 3 + \operatorname{twopow}(k), \operatorname{twopow}(k) - (r + 2)) = \operatorname{false}\).
\end{theorem}
\begin{proof}
The goal follows from \(\operatorname{lhcenterlatecrossstepthreepinrplustwofalseofzerodiagfourandsinglewinwittagree}(t, k, r, a, L)\).
\end{proof}
\begin{theorem}
\label{ll:paper:lhcenter-late-cross-step-absorb-false-after-zero-diag}
For every natural number \(t\) and natural number \(k\) and natural number \(r\), if \(t + r + 2 \le k\), then if \(\operatorname{zerodiagonal}(t + r)\), then if \(\operatorname{settledby}(t + r + 2, t + \operatorname{twopow}(k))\), then if \(\operatorname{lhcenter}(t + \operatorname{twopow}(k), \operatorname{twopow}(k) - (r + 3)) = \operatorname{true}\), then \(\operatorname{lhcenter}(t + 1 + \operatorname{twopow}(k), \operatorname{twopow}(k) - (r + 3)) = \operatorname{false}\).
\end{theorem}
\begin{proof}
The goal follows from \(\operatorname{lhcenterlatecrossstepabsorbfalseafterzerodiag}(t, k, r)\).
\end{proof}
\begin{theorem}
\label{ll:paper:lhcenter-late-cross-two-step-absorb-false-after-zero-diag}
For every natural number \(t\) and natural number \(k\) and natural number \(r\), if \(t + r + 2 \le k\), then if \(\operatorname{zerodiagonal}(t + r)\), then if \(\operatorname{settledby}(t + r + 2, t + \operatorname{twopow}(k))\), then if \(\operatorname{lhcenter}(t + \operatorname{twopow}(k), \operatorname{twopow}(k) - (r + 1)) = \operatorname{false}\), then \(\operatorname{lhcenter}(t + 2 + \operatorname{twopow}(k), \operatorname{twopow}(k) - (r + 2)) = \operatorname{false}\).
\end{theorem}
\begin{proof}
The goal follows from \(\operatorname{lhcenterlatecrosstwostepabsorbfalseafterzerodiag}(t, k, r)\).
\end{proof}
\begin{theorem}
\label{ll:paper:lhcenter-late-cross-step-two-pin-r-plus-five-true-of-doubles-four-and-single-winwittagree}
For every natural number \(t\) and natural number \(k\) and natural number \(r\) and natural number \(a\) and natural number \(L\), if \(t + 4 < r\), then if \(t + 4 + r + 2 \le k\), then if \(\operatorname{zerodiagonal}(t + 4 + r)\), then if \(\operatorname{doublesat}(4, t + 4 + r)\), then if \(\operatorname{settledby}(t + 4 + r, t + \operatorname{twopow}(k))\), then if \(\operatorname{settledby}(t + 4 + r + 1, t + \operatorname{twopow}(k))\), then if \(\operatorname{settledby}(t + 4 + r + 2, t + 1 + \operatorname{twopow}(k))\), then if \(a \le t + \operatorname{twopow}(k + 1) - (r + 2)\), then if \(t + \operatorname{twopow}(k + 1) - r \le a + L\), then if \(\operatorname{winwittagree}(t, k, 4, a, L)\), then \(\operatorname{lhcenter}(t + 2 + \operatorname{twopow}(k), \operatorname{twopow}(k) - (r + 5)) = \operatorname{true}\).
\end{theorem}
\begin{proof}
The goal follows from \(\operatorname{lhcenterlatecrosssteptwopinrplusfivetrueofdoublesfourandsinglewinwittagree}(t, k, r, a, L)\).
\end{proof}
\begin{theorem}
\label{ll:paper:lhcenter-late-cross-step-three-pin-r-plus-five-false-of-doubles-four-and-single-winwittagree}
For every natural number \(t\) and natural number \(k\) and natural number \(r\) and natural number \(a\) and natural number \(L\), if \(t + 4 < r\), then if \(t + 4 + r + 2 \le k\), then if \(\operatorname{zerodiagonal}(t + 4 + r)\), then if \(\operatorname{doublesat}(4, t + 4 + r)\), then if \(\operatorname{settledby}(t + 4 + r, t + \operatorname{twopow}(k))\), then if \(\operatorname{settledby}(t + 4 + r + 1, t + \operatorname{twopow}(k))\), then if \(\operatorname{settledby}(t + 4 + r + 2, t + 1 + \operatorname{twopow}(k))\), then if \(a \le t + \operatorname{twopow}(k + 1) - (r + 2)\), then if \(t + \operatorname{twopow}(k + 1) - r \le a + L\), then if \(\operatorname{winwittagree}(t, k, 4, a, L)\), then \(\operatorname{lhcenter}(t + 3 + \operatorname{twopow}(k), \operatorname{twopow}(k) - (r + 5)) = \operatorname{false}\).
\end{theorem}
\begin{proof}
The goal follows from \(\operatorname{lhcenterlatecrossstepthreepinrplusfivefalseofdoublesfourandsinglewinwittagree}(t, k, r, a, L)\).
\end{proof}
\begin{theorem}
\label{ll:paper:lhcenter-late-cross-step-four-pin-r-plus-four-true-of-doubles-four-and-single-winwittagree}
For every natural number \(t\) and natural number \(k\) and natural number \(r\) and natural number \(a\) and natural number \(L\), if \(t + 4 < r\), then if \(t + 4 + r + 2 \le k\), then if \(\operatorname{zerodiagonal}(t + 4 + r)\), then if \(\operatorname{doublesat}(4, t + 4 + r)\), then if \(\operatorname{settledby}(t + 4 + r, t + \operatorname{twopow}(k))\), then if \(\operatorname{settledby}(t + 4 + r + 1, t + \operatorname{twopow}(k))\), then if \(\operatorname{settledby}(t + 4 + r + 2, t + 1 + \operatorname{twopow}(k))\), then if \(a \le t + \operatorname{twopow}(k + 1) - (r + 2)\), then if \(t + \operatorname{twopow}(k + 1) - r \le a + L\), then if \(\operatorname{winwittagree}(t, k, 4, a, L)\), then \(\operatorname{lhcenter}(t + 4 + \operatorname{twopow}(k), \operatorname{twopow}(k) - (r + 4)) = \operatorname{true}\).
\end{theorem}
\begin{proof}
The goal follows from \(\operatorname{lhcenterlatecrossstepfourpinrplusfourtrueofdoublesfourandsinglewinwittagree}(t, k, r, a, L)\).
\end{proof}
\begin{theorem}
\label{ll:paper:not-winwittagree-four-of-doubles}
For every natural number \(t\) and natural number \(k\) and natural number \(r\) and natural number \(a\) and natural number \(L\), if \(t + 4 < r\), then if \(t + 4 + r + 2 \le k\), then if \(\operatorname{zerodiagonal}(t + 4 + r)\), then if \(\operatorname{doublesat}(4, t + 4 + r)\), then if \(\operatorname{settledby}(t + 4 + r, t + \operatorname{twopow}(k))\), then if \(\operatorname{settledby}(t + 4 + r + 1, t + \operatorname{twopow}(k))\), then if \(\operatorname{settledby}(t + 4 + r + 2, t + 1 + \operatorname{twopow}(k))\), then if \(a \le t + \operatorname{twopow}(k + 1) - (r + 4)\), then if \(t + \operatorname{twopow}(k + 1) - r \le a + L\), then not \(\operatorname{winwittagree}(t, k, 4, a, L)\).
\end{theorem}
\begin{proof}
The goal follows from \(\operatorname{notwinwittagreefourofdoubles}(t, k, r, a, L)\).
\end{proof}
\begin{theorem}
\label{ll:paper:winwitt-scale-cross-zero-gap-step-one-pin-r-plus-one-true-of-single-winwittagree-and-zero-diag-four}
For every natural number \(Q\) and natural number \(t\) and natural number \(k\) and natural number \(r\) and natural number \(a\) and natural number \(L\), if \(t + Q < r\), then if \(t + Q + r + 2 \le k\), then if \(t + r + 4 \le k\), then if \(\operatorname{zerodiagonal}(t + Q + r)\), then if \(\operatorname{zerodiagonal}(t + r + 4)\), then if \(\operatorname{settledby}(t + Q + r, t + Q + \operatorname{twopow}(k))\), then if \(\operatorname{settledby}(t + Q + r + 2, t + Q + \operatorname{twopow}(k))\), then if \(\operatorname{settledby}(t + r + 4, t + 1 + \operatorname{twopow}(k))\), then if \(a \le t + \operatorname{twopow}(k + 1) - (r + 2)\), then if \(t + \operatorname{twopow}(k + 1) - r \le a + L\), then if \(\operatorname{winwittagree}(t, k, Q, a, L)\), then \(\operatorname{wittbitbool}(\operatorname{wittwin}(t + 1, \operatorname{twopow}(k)), t + \operatorname{twopow}(k + 1) - (r + 1)) = \operatorname{true}\).
\end{theorem}
\begin{proof}
The goal follows from \(\operatorname{winwittscalecrosszerogapsteponepinrplusonetrueofsinglewinwittagreeandzerodiagfour}(Q, t, k, r, a, L)\).
\end{proof}
\begin{theorem}
\label{ll:paper:winwitt-scale-cross-zero-gap-pin-r-plus-one-false-of-single-winwittagree-and-zero-diag-four}
For every natural number \(Q\) and natural number \(t\) and natural number \(k\) and natural number \(r\) and natural number \(a\) and natural number \(L\), if \(t + Q < r\), then if \(t + Q + r + 2 \le k\), then if \(t + r + 4 \le k\), then if \(\operatorname{zerodiagonal}(t + Q + r)\), then if \(\operatorname{zerodiagonal}(t + r + 4)\), then if \(\operatorname{settledby}(t + Q + r, t + Q + \operatorname{twopow}(k))\), then if \(\operatorname{settledby}(t + Q + r + 2, t + Q + \operatorname{twopow}(k))\), then if \(\operatorname{settledby}(t + r + 4, t + 1 + \operatorname{twopow}(k))\), then if \(a \le t + \operatorname{twopow}(k + 1) - (r + 2)\), then if \(t + \operatorname{twopow}(k + 1) - r \le a + L\), then if \(\operatorname{winwittagree}(t, k, Q, a, L)\), then \(\operatorname{wittbitbool}(\operatorname{wittwin}(t, \operatorname{twopow}(k)), t + \operatorname{twopow}(k + 1) - (r + 1)) = \operatorname{false}\).
\end{theorem}
\begin{proof}
The goal follows from \(\operatorname{winwittscalecrosszerogappinrplusonefalseofsinglewinwittagreeandzerodiagfour}(Q, t, k, r, a, L)\).
\end{proof}
\begin{theorem}
\label{ll:paper:winwitt-scale-cross-zero-gap-pin-doubling-diag-true-of-doubles-and-single-winwittagree-and-zero-diag-four}
For every natural number \(Q\) and natural number \(t\) and natural number \(k\) and natural number \(r\) and natural number \(a\) and natural number \(L\), if \(t + Q < r\), then if \(t + Q + r + 2 \le k\), then if \(t + r + 4 \le k\), then if \(\operatorname{zerodiagonal}(t + Q + r)\), then if \(\operatorname{doublesat}(Q, t + Q + r)\), then if \(\operatorname{zerodiagonal}(t + r + 4)\), then if \(\operatorname{settledby}(t + Q + r, t + Q + \operatorname{twopow}(k))\), then if \(\operatorname{settledby}(t + Q + r + 1, t + \operatorname{twopow}(k))\), then if \(\operatorname{settledby}(t + Q + r + 2, t + Q + \operatorname{twopow}(k))\), then if \(\operatorname{settledby}(t + r + 4, t + 1 + \operatorname{twopow}(k))\), then if \(a \le t + \operatorname{twopow}(k + 1) - (r + 2)\), then if \(t + \operatorname{twopow}(k + 1) - r \le a + L\), then if \(\operatorname{winwittagree}(t, k, Q, a, L)\), then \(\operatorname{wittbitbool}(\operatorname{wittwin}(t, \operatorname{twopow}(k)), t + \operatorname{twopow}(k + 1) - (r + Q + 1)) = \operatorname{true}\).
\end{theorem}
\begin{proof}
The goal follows from \(\operatorname{winwittscalecrosszerogappindoublingdiagtrueofdoublesandsinglewinwittagreeandzerodiagfour}(Q, t, k, r, a, L)\).
\end{proof}
\begin{theorem}
\label{ll:paper:winwitt-scale-cross-zero-gap-pin-r-plus-three-true-of-zero-diag-four-and-single-winwittagree}
For every natural number \(t\) and natural number \(k\) and natural number \(r\) and natural number \(a\) and natural number \(L\), if \(t + 4 \le r + 1\), then if \(t + 4 + r + 2 \le k\), then if \(\operatorname{zerodiagonal}(t + 4 + r)\), then if \(\operatorname{settledby}(t + 4 + r, t + \operatorname{twopow}(k))\), then if \(\operatorname{settledby}(t + 4 + r + 2, t + 4 + \operatorname{twopow}(k))\), then if \(a \le t + \operatorname{twopow}(k + 1) - (r + 2)\), then if \(t + \operatorname{twopow}(k + 1) - (r + 2) \le a + L\), then if \(\operatorname{winwittagree}(t, k, 4, a, L)\), then \(\operatorname{wittbitbool}(\operatorname{wittwin}(t, \operatorname{twopow}(k)), t + \operatorname{twopow}(k + 1) - (r + 3)) = \operatorname{true}\).
\end{theorem}
\begin{proof}
The goal follows from \(\operatorname{winwittscalecrosszerogappinrplusthreetrueofzerodiagfourandsinglewinwittagree}(t, k, r, a, L)\).
\end{proof}
\begin{theorem}
\label{ll:paper:winwitt-scale-cross-zero-gap-step-one-pin-r-plus-one-true-of-zero-diag-four-and-single-winwittagree}
For every natural number \(t\) and natural number \(k\) and natural number \(r\) and natural number \(a\) and natural number \(L\), if \(t + 4 < r\), then if \(t + 4 + r + 2 \le k\), then if \(\operatorname{zerodiagonal}(t + 4 + r)\), then if \(\operatorname{settledby}(t + 4 + r, t + 1 + \operatorname{twopow}(k))\), then if \(\operatorname{settledby}(t + 4 + r + 2, t + 4 + \operatorname{twopow}(k))\), then if \(a \le t + \operatorname{twopow}(k + 1) - (r + 2)\), then if \(t + \operatorname{twopow}(k + 1) - r \le a + L\), then if \(\operatorname{winwittagree}(t, k, 4, a, L)\), then \(\operatorname{wittbitbool}(\operatorname{wittwin}(t + 1, \operatorname{twopow}(k)), t + \operatorname{twopow}(k + 1) - (r + 1)) = \operatorname{true}\).
\end{theorem}
\begin{proof}
The goal follows from \(\operatorname{winwittscalecrosszerogapsteponepinrplusonetrueofzerodiagfourandsinglewinwittagree}(t, k, r, a, L)\).
\end{proof}
\begin{theorem}
\label{ll:paper:winwitt-scale-cross-zero-gap-pin-r-plus-one-false-of-zero-diag-four-and-single-winwittagree}
For every natural number \(t\) and natural number \(k\) and natural number \(r\) and natural number \(a\) and natural number \(L\), if \(t + 4 < r\), then if \(t + 4 + r + 2 \le k\), then if \(\operatorname{zerodiagonal}(t + 4 + r)\), then if \(\operatorname{settledby}(t + 4 + r, t + 1 + \operatorname{twopow}(k))\), then if \(\operatorname{settledby}(t + 4 + r + 2, t + 4 + \operatorname{twopow}(k))\), then if \(a \le t + \operatorname{twopow}(k + 1) - (r + 2)\), then if \(t + \operatorname{twopow}(k + 1) - r \le a + L\), then if \(\operatorname{winwittagree}(t, k, 4, a, L)\), then \(\operatorname{wittbitbool}(\operatorname{wittwin}(t, \operatorname{twopow}(k)), t + \operatorname{twopow}(k + 1) - (r + 1)) = \operatorname{false}\).
\end{theorem}
\begin{proof}
The goal follows from \(\operatorname{winwittscalecrosszerogappinrplusonefalseofzerodiagfourandsinglewinwittagree}(t, k, r, a, L)\).
\end{proof}
\begin{theorem}
\label{ll:paper:winwitt-scale-cross-zero-gap-pin-doubling-diag-true-of-doubles-four-and-single-winwittagree}
For every natural number \(t\) and natural number \(k\) and natural number \(r\) and natural number \(a\) and natural number \(L\), if \(t + 4 < r\), then if \(t + 4 + r + 2 \le k\), then if \(\operatorname{zerodiagonal}(t + 4 + r)\), then if \(\operatorname{doublesat}(4, t + 4 + r)\), then if \(\operatorname{settledby}(t + 4 + r, t + 1 + \operatorname{twopow}(k))\), then if \(\operatorname{settledby}(t + 4 + r + 1, t + \operatorname{twopow}(k))\), then if \(\operatorname{settledby}(t + 4 + r + 2, t + 4 + \operatorname{twopow}(k))\), then if \(a \le t + \operatorname{twopow}(k + 1) - (r + 2)\), then if \(t + \operatorname{twopow}(k + 1) - r \le a + L\), then if \(\operatorname{winwittagree}(t, k, 4, a, L)\), then \(\operatorname{wittbitbool}(\operatorname{wittwin}(t, \operatorname{twopow}(k)), t + \operatorname{twopow}(k + 1) - (r + 4 + 1)) = \operatorname{true}\).
\end{theorem}
\begin{proof}
The goal follows from \(\operatorname{winwittscalecrosszerogappindoublingdiagtrueofdoublesfourandsinglewinwittagree}(t, k, r, a, L)\).
\end{proof}
\begin{theorem}
\label{ll:paper:winwitt-scale-cross-zero-gap-step-four-pin-r-true-of-doubles-four-and-single-winwittagree}
For every natural number \(t\) and natural number \(k\) and natural number \(r\) and natural number \(a\) and natural number \(L\), if \(t + 4 < r\), then if \(t + 4 + r + 2 \le k\), then if \(\operatorname{zerodiagonal}(t + 4 + r)\), then if \(\operatorname{doublesat}(4, t + 4 + r)\), then if \(\operatorname{settledby}(t + 4 + r, t + \operatorname{twopow}(k))\), then if \(\operatorname{settledby}(t + 4 + r + 1, t + \operatorname{twopow}(k))\), then if \(\operatorname{settledby}(t + 4 + r + 2, t + 1 + \operatorname{twopow}(k))\), then if \(a \le t + \operatorname{twopow}(k + 1) - (r + 2)\), then if \(t + \operatorname{twopow}(k + 1) - r \le a + L\), then if \(\operatorname{winwittagree}(t, k, 4, a, L)\), then \(\operatorname{wittbitbool}(\operatorname{wittwin}(t + 4, \operatorname{twopow}(k)), t + \operatorname{twopow}(k + 1) - r) = \operatorname{true}\).
\end{theorem}
\begin{proof}
The goal follows from \(\operatorname{winwittscalecrosszerogapstepfourpinrtrueofdoublesfourandsinglewinwittagree}(t, k, r, a, L)\).
\end{proof}
\begin{theorem}
\label{ll:paper:winwitt-scale-cross-zero-gap-same-bit-r-flip-four-of-doubles-four-and-single-winwittagree}
For every natural number \(t\) and natural number \(k\) and natural number \(r\) and natural number \(a\) and natural number \(L\), if \(t + 4 < r\), then if \(t + 4 + r + 2 \le k\), then if \(\operatorname{zerodiagonal}(t + 4 + r)\), then if \(\operatorname{doublesat}(4, t + 4 + r)\), then if \(\operatorname{settledby}(t + 4 + r, t + \operatorname{twopow}(k))\), then if \(\operatorname{settledby}(t + 4 + r + 1, t + \operatorname{twopow}(k))\), then if \(\operatorname{settledby}(t + 4 + r + 2, t + 1 + \operatorname{twopow}(k))\), then if \(a \le t + \operatorname{twopow}(k + 1) - (r + 2)\), then if \(t + \operatorname{twopow}(k + 1) - r \le a + L\), then if \(\operatorname{winwittagree}(t, k, 4, a, L)\), then \(\operatorname{wittbitbool}(\operatorname{wittwin}(t + 4, \operatorname{twopow}(k)), t + \operatorname{twopow}(k + 1) - r) = \operatorname{not}(\operatorname{wittbitbool}(\operatorname{wittwin}(t, \operatorname{twopow}(k)), t + \operatorname{twopow}(k + 1) - r))\).
\end{theorem}
\begin{proof}
The goal follows from \(\operatorname{winwittscalecrosszerogapsamebitrflipfourofdoublesfourandsinglewinwittagree}(t, k, r, a, L)\).
\end{proof}
\begin{theorem}
\label{ll:paper:not-isblackdiag-add-four-of-zero-diag}
For every natural number \(z\), if \(\operatorname{zerodiagonal}(z)\), then not \(\operatorname{blackdiagonal}(z + 4)\).
\end{theorem}
\begin{proof}
The goal follows from \(\operatorname{notisblackdiagaddfourofzerodiag}(z)\).
\end{proof}
\begin{theorem}
\label{ll:paper:not-isblackdiag-add-five-of-zero-diag}
For every natural number \(z\), if \(\operatorname{zerodiagonal}(z)\), then not \(\operatorname{blackdiagonal}(z + 5)\).
\end{theorem}
\begin{proof}
The goal follows from \(\operatorname{notisblackdiagaddfiveofzerodiag}(z)\).
\end{proof}
\begin{theorem}
\label{ll:paper:not-isblackdiag-add-six-of-zero-diag}
For every natural number \(z\), if \(\operatorname{zerodiagonal}(z)\), then not \(\operatorname{blackdiagonal}(z + 6)\).
\end{theorem}
\begin{proof}
The goal follows from \(\operatorname{notisblackdiagaddsixofzerodiag}(z)\).
\end{proof}
\begin{theorem}
\label{ll:paper:isblackdiag-sub-one-of-zero-diag-and-add-five}
For every natural number \(z\), if \(\operatorname{zerodiagonal}(z)\), then if \(\operatorname{zerodiagonal}(z + 5)\), then \(\operatorname{blackdiagonal}(z - 1)\).
\end{theorem}
\begin{proof}
The goal follows from \(\operatorname{isblackdiagsuboneofzerodiagandaddfive}(z)\).
\end{proof}
\begin{theorem}
\label{ll:paper:zero-diag-le-two-of-add-five}
For every natural number \(z\), if \(\operatorname{zerodiagonal}(z)\), then if \(\operatorname{zerodiagonal}(z + 5)\), then \(z \le 2\).
\end{theorem}
\begin{proof}
The goal follows from \(\operatorname{zerodiagletwoofaddfive}(z)\).
\end{proof}
\begin{theorem}
\label{ll:paper:not-iszerodiag-add-five-of-gt-two}
For every natural number \(z\), if \(2 < z\), then if \(\operatorname{zerodiagonal}(z)\), then not \(\operatorname{zerodiagonal}(z + 5)\).
\end{theorem}
\begin{proof}
The goal follows from \(\operatorname{notiszerodiagaddfiveofgttwo}(z)\).
\end{proof}
\begin{theorem}
\label{ll:paper:zero-diag-gap-ge-six-of-gt-two}
For every natural number \(z1\) and natural number \(z2\), if \(2 < z1\), then if \(\operatorname{zerodiagonal}(z1)\), then if \(\operatorname{zerodiagonal}(z2)\), then if \(z1 < z2\), then \(z1 + 6 \le z2\).
\end{theorem}
\begin{proof}
The goal follows from \(\operatorname{zerodiaggapgesixofgttwo}(z1, z2)\).
\end{proof}
\begin{theorem}
\label{ll:paper:not-isblackdiag-add-seven-of-gt-two-and-zero-diag}
For every natural number \(z\), if \(2 < z\), then if \(\operatorname{zerodiagonal}(z)\), then not \(\operatorname{blackdiagonal}(z + 7)\).
\end{theorem}
\begin{proof}
The goal follows from \(\operatorname{notisblackdiagaddsevenofgttwoandzerodiag}(z)\).
\end{proof}
\begin{theorem}
\label{ll:paper:winwitt-scale-cross-zero-gap-step-one-pin-r-plus-three-false-of-doubles-four-and-single-winwittagree}
For every natural number \(t\) and natural number \(k\) and natural number \(r\) and natural number \(a\) and natural number \(L\), if \(t + 4 < r\), then if \(t + 4 + r + 2 \le k\), then if \(\operatorname{zerodiagonal}(t + 4 + r)\), then if \(\operatorname{doublesat}(4, t + 4 + r)\), then if \(\operatorname{settledby}(t + 4 + r, t + \operatorname{twopow}(k))\), then if \(\operatorname{settledby}(t + 4 + r + 1, t + \operatorname{twopow}(k))\), then if \(\operatorname{settledby}(t + 4 + r + 2, t + 4 + \operatorname{twopow}(k))\), then if \(a \le t + \operatorname{twopow}(k + 1) - (r + 2)\), then if \(t + \operatorname{twopow}(k + 1) - r \le a + L\), then if \(\operatorname{winwittagree}(t, k, 4, a, L)\), then \(\operatorname{wittbitbool}(\operatorname{wittwin}(t + 1, \operatorname{twopow}(k)), t + \operatorname{twopow}(k + 1) - (r + 3)) = \operatorname{false}\).
\end{theorem}
\begin{proof}
The goal follows from \(\operatorname{winwittscalecrosszerogapsteponepinrplusthreefalseofdoublesfourandsinglewinwittagree}(t, k, r, a, L)\).
\end{proof}
\begin{theorem}
\label{ll:paper:winwitt-scale-cross-zero-gap-step-two-pin-r-plus-one-true-of-doubles-four-and-single-winwittagree}
For every natural number \(t\) and natural number \(k\) and natural number \(r\) and natural number \(a\) and natural number \(L\), if \(t + 4 < r\), then if \(t + 4 + r + 2 \le k\), then if \(\operatorname{zerodiagonal}(t + 4 + r)\), then if \(\operatorname{doublesat}(4, t + 4 + r)\), then if \(\operatorname{settledby}(t + 4 + r, t + \operatorname{twopow}(k))\), then if \(\operatorname{settledby}(t + 4 + r + 1, t + \operatorname{twopow}(k))\), then if \(\operatorname{settledby}(t + 4 + r + 2, t + 4 + \operatorname{twopow}(k))\), then if \(a \le t + \operatorname{twopow}(k + 1) - (r + 2)\), then if \(t + \operatorname{twopow}(k + 1) - r \le a + L\), then if \(\operatorname{winwittagree}(t, k, 4, a, L)\), then \(\operatorname{wittbitbool}(\operatorname{wittwin}(t + 2, \operatorname{twopow}(k)), t + \operatorname{twopow}(k + 1) - (r + 1)) = \operatorname{true}\).
\end{theorem}
\begin{proof}
The goal follows from \(\operatorname{winwittscalecrosszerogapsteptwopinrplusonetrueofdoublesfourandsinglewinwittagree}(t, k, r, a, L)\).
\end{proof}
\begin{theorem}
\label{ll:paper:winwitt-scale-cross-zero-gap-step-three-pin-r-plus-one-false-of-doubles-four-and-single-winwittagree}
For every natural number \(t\) and natural number \(k\) and natural number \(r\) and natural number \(a\) and natural number \(L\), if \(t + 4 < r\), then if \(t + 4 + r + 2 \le k\), then if \(\operatorname{zerodiagonal}(t + 4 + r)\), then if \(\operatorname{doublesat}(4, t + 4 + r)\), then if \(\operatorname{settledby}(t + 4 + r, t + \operatorname{twopow}(k))\), then if \(\operatorname{settledby}(t + 4 + r + 1, t + \operatorname{twopow}(k))\), then if \(\operatorname{settledby}(t + 4 + r + 2, t + 2 + \operatorname{twopow}(k))\), then if \(a \le t + \operatorname{twopow}(k + 1) - (r + 2)\), then if \(t + \operatorname{twopow}(k + 1) - r \le a + L\), then if \(\operatorname{winwittagree}(t, k, 4, a, L)\), then \(\operatorname{wittbitbool}(\operatorname{wittwin}(t + 3, \operatorname{twopow}(k)), t + \operatorname{twopow}(k + 1) - (r + 1)) = \operatorname{false}\).
\end{theorem}
\begin{proof}
The goal follows from \(\operatorname{winwittscalecrosszerogapstepthreepinrplusonefalseofdoublesfourandsinglewinwittagree}(t, k, r, a, L)\).
\end{proof}
\begin{theorem}
\label{ll:paper:winwitt-scale-cross-zero-gap-step-two-pin-r-plus-three-true-of-doubles-four-and-single-winwittagree}
For every natural number \(t\) and natural number \(k\) and natural number \(r\) and natural number \(a\) and natural number \(L\), if \(t + 4 < r\), then if \(t + 4 + r + 2 \le k\), then if \(\operatorname{zerodiagonal}(t + 4 + r)\), then if \(\operatorname{doublesat}(4, t + 4 + r)\), then if \(\operatorname{settledby}(t + 4 + r, t + \operatorname{twopow}(k))\), then if \(\operatorname{settledby}(t + 4 + r + 1, t + \operatorname{twopow}(k))\), then if \(\operatorname{settledby}(t + 4 + r + 2, t + 1 + \operatorname{twopow}(k))\), then if \(a \le t + \operatorname{twopow}(k + 1) - (r + 2)\), then if \(t + \operatorname{twopow}(k + 1) - r \le a + L\), then if \(\operatorname{winwittagree}(t, k, 4, a, L)\), then \(\operatorname{wittbitbool}(\operatorname{wittwin}(t + 2, \operatorname{twopow}(k)), t + \operatorname{twopow}(k + 1) - (r + 3)) = \operatorname{true}\).
\end{theorem}
\begin{proof}
The goal follows from \(\operatorname{winwittscalecrosszerogapsteptwopinrplusthreetrueofdoublesfourandsinglewinwittagree}(t, k, r, a, L)\).
\end{proof}
\begin{theorem}
\label{ll:paper:winwitt-scale-cross-zero-gap-step-three-pin-r-plus-two-false-of-doubles-four-and-single-winwittagree}
For every natural number \(t\) and natural number \(k\) and natural number \(r\) and natural number \(a\) and natural number \(L\), if \(t + 4 < r\), then if \(t + 4 + r + 2 \le k\), then if \(\operatorname{zerodiagonal}(t + 4 + r)\), then if \(\operatorname{doublesat}(4, t + 4 + r)\), then if \(\operatorname{settledby}(t + 4 + r, t + \operatorname{twopow}(k))\), then if \(\operatorname{settledby}(t + 4 + r + 1, t + \operatorname{twopow}(k))\), then if \(\operatorname{settledby}(t + 4 + r + 2, t + 1 + \operatorname{twopow}(k))\), then if \(a \le t + \operatorname{twopow}(k + 1) - (r + 2)\), then if \(t + \operatorname{twopow}(k + 1) - r \le a + L\), then if \(\operatorname{winwittagree}(t, k, 4, a, L)\), then \(\operatorname{wittbitbool}(\operatorname{wittwin}(t + 3, \operatorname{twopow}(k)), t + \operatorname{twopow}(k + 1) - (r + 2)) = \operatorname{false}\).
\end{theorem}
\begin{proof}
The goal follows from \(\operatorname{winwittscalecrosszerogapstepthreepinrplustwofalseofdoublesfourandsinglewinwittagree}(t, k, r, a, L)\).
\end{proof}
\begin{theorem}
\label{ll:paper:winwitt-scale-cross-zero-gap-same-bit-r-plus-two-flip-three-of-doubles-four-and-single-winwittagree}
For every natural number \(t\) and natural number \(k\) and natural number \(r\) and natural number \(a\) and natural number \(L\), if \(t + 4 < r\), then if \(t + 4 + r + 2 \le k\), then if \(\operatorname{zerodiagonal}(t + 4 + r)\), then if \(\operatorname{doublesat}(4, t + 4 + r)\), then if \(\operatorname{settledby}(t + 4 + r, t + \operatorname{twopow}(k))\), then if \(\operatorname{settledby}(t + 4 + r + 1, t + \operatorname{twopow}(k))\), then if \(\operatorname{settledby}(t + 4 + r + 2, t + 1 + \operatorname{twopow}(k))\), then if \(a \le t + \operatorname{twopow}(k + 1) - (r + 2)\), then if \(t + \operatorname{twopow}(k + 1) - r \le a + L\), then if \(\operatorname{winwittagree}(t, k, 4, a, L)\), then \(\operatorname{wittbitbool}(\operatorname{wittwin}(t + 3, \operatorname{twopow}(k)), t + \operatorname{twopow}(k + 1) - (r + 2)) = \operatorname{not}(\operatorname{wittbitbool}(\operatorname{wittwin}(t, \operatorname{twopow}(k)), t + \operatorname{twopow}(k + 1) - (r + 2)))\).
\end{theorem}
\begin{proof}
The goal follows from \(\operatorname{winwittscalecrosszerogapsamebitrplustwoflipthreeofdoublesfourandsinglewinwittagree}(t, k, r, a, L)\).
\end{proof}
\begin{theorem}
\label{ll:paper:lhcenter-late-cross-step-one-pin-reset-diag-false-of-doubles-and-single-winwittagree-and-zero-diag-four}
For every natural number \(Q\) and natural number \(t\) and natural number \(k\) and natural number \(r\) and natural number \(a\) and natural number \(L\), if \(t + Q < r\), then if \(t + Q + r + 2 \le k\), then if \(t + r + 4 \le k\), then if \(\operatorname{zerodiagonal}(t + Q + r)\), then if \(\operatorname{doublesat}(Q, t + Q + r)\), then if \(\operatorname{zerodiagonal}(t + r + 4)\), then if \(\operatorname{settledby}(t + Q + r, t + Q + \operatorname{twopow}(k))\), then if \(\operatorname{settledby}(t + Q + r + 1, t + \operatorname{twopow}(k))\), then if \(\operatorname{settledby}(t + Q + r + 2, t + \operatorname{twopow}(k))\), then if \(\operatorname{settledby}(t + r + 4, t + 1 + \operatorname{twopow}(k))\), then if \(a \le t + \operatorname{twopow}(k + 1) - (r + 2)\), then if \(t + \operatorname{twopow}(k + 1) - r \le a + L\), then if \(\operatorname{winwittagree}(t, k, Q, a, L)\), then \(\operatorname{lhcenter}(t + 1 + \operatorname{twopow}(k), \operatorname{twopow}(k) - (r + Q + 2)) = \operatorname{false}\).
\end{theorem}
\begin{proof}
The goal follows from \(\operatorname{lhcenterlatecrosssteponepinresetdiagfalseofdoublesandsinglewinwittagreeandzerodiagfour}(Q, t, k, r, a, L)\).
\end{proof}
\begin{theorem}
\label{ll:paper:winwitt-scale-cross-zero-gap-step-one-pin-reset-diag-false-of-doubles-and-single-winwittagree-and-zero-diag-four}
For every natural number \(Q\) and natural number \(t\) and natural number \(k\) and natural number \(r\) and natural number \(a\) and natural number \(L\), if \(t + Q < r\), then if \(t + Q + r + 2 \le k\), then if \(t + r + 4 \le k\), then if \(\operatorname{zerodiagonal}(t + Q + r)\), then if \(\operatorname{doublesat}(Q, t + Q + r)\), then if \(\operatorname{zerodiagonal}(t + r + 4)\), then if \(\operatorname{settledby}(t + Q + r, t + Q + \operatorname{twopow}(k))\), then if \(\operatorname{settledby}(t + Q + r + 1, t + \operatorname{twopow}(k))\), then if \(\operatorname{settledby}(t + Q + r + 2, t + \operatorname{twopow}(k))\), then if \(\operatorname{settledby}(t + r + 4, t + 1 + \operatorname{twopow}(k))\), then if \(a \le t + \operatorname{twopow}(k + 1) - (r + 2)\), then if \(t + \operatorname{twopow}(k + 1) - r \le a + L\), then if \(\operatorname{winwittagree}(t, k, Q, a, L)\), then \(\operatorname{wittbitbool}(\operatorname{wittwin}(t + 1, \operatorname{twopow}(k)), t + \operatorname{twopow}(k + 1) - (r + Q + 1)) = \operatorname{false}\).
\end{theorem}
\begin{proof}
The goal follows from \(\operatorname{winwittscalecrosszerogapsteponepinresetdiagfalseofdoublesandsinglewinwittagreeandzerodiagfour}(Q, t, k, r, a, L)\).
\end{proof}
\begin{theorem}
\label{ll:paper:winwitt-scale-cross-zero-gap-step-one-same-bit-doubling-flip-of-doubles-and-single-winwittagree-and-zero-diag-four}
For every natural number \(Q\) and natural number \(t\) and natural number \(k\) and natural number \(r\) and natural number \(a\) and natural number \(L\), if \(t + Q < r\), then if \(t + Q + r + 2 \le k\), then if \(t + r + 4 \le k\), then if \(\operatorname{zerodiagonal}(t + Q + r)\), then if \(\operatorname{doublesat}(Q, t + Q + r)\), then if \(\operatorname{zerodiagonal}(t + r + 4)\), then if \(\operatorname{settledby}(t + Q + r, t + Q + \operatorname{twopow}(k))\), then if \(\operatorname{settledby}(t + Q + r + 1, t + \operatorname{twopow}(k))\), then if \(\operatorname{settledby}(t + Q + r + 2, t + \operatorname{twopow}(k))\), then if \(\operatorname{settledby}(t + r + 4, t + 1 + \operatorname{twopow}(k))\), then if \(a \le t + \operatorname{twopow}(k + 1) - (r + 2)\), then if \(t + \operatorname{twopow}(k + 1) - r \le a + L\), then if \(\operatorname{winwittagree}(t, k, Q, a, L)\), then \(\operatorname{wittbitbool}(\operatorname{wittwin}(t + 1, \operatorname{twopow}(k)), t + \operatorname{twopow}(k + 1) - (r + Q + 1)) = \operatorname{not}(\operatorname{wittbitbool}(\operatorname{wittwin}(t, \operatorname{twopow}(k)), t + \operatorname{twopow}(k + 1) - (r + Q + 1)))\).
\end{theorem}
\begin{proof}
The goal follows from \(\operatorname{winwittscalecrosszerogapsteponesamebitdoublingflipofdoublesandsinglewinwittagreeandzerodiagfour}(Q, t, k, r, a, L)\).
\end{proof}
\begin{theorem}
\label{ll:paper:winwitt-scale-cross-zero-gap-step-one-same-bit-doubling-flip-of-doubles-and-single-winwittagree-and-zero-diag-three}
For every natural number \(Q\) and natural number \(t\) and natural number \(k\) and natural number \(r\) and natural number \(a\) and natural number \(L\), if \(t + Q \le r\), then if \(t + Q + r + 3 \le k\), then if \(\operatorname{zerodiagonal}(t + Q + r)\), then if \(\operatorname{doublesat}(Q, t + Q + r)\), then if \(\operatorname{zerodiagonal}(t + r + 3)\), then if \(\operatorname{settledby}(t + r + 3, t + \operatorname{twopow}(k))\), then if \(\operatorname{settledby}(t + Q + r + 1, t + \operatorname{twopow}(k))\), then if \(\operatorname{settledby}(t + Q + r + 2, t + \operatorname{twopow}(k))\), then if \(a \le t + \operatorname{twopow}(k + 1) - (r + 2)\), then if \(t + \operatorname{twopow}(k + 1) - (r + 1) \le a + L\), then if \(\operatorname{winwittagree}(t, k, Q, a, L)\), then \(\operatorname{wittbitbool}(\operatorname{wittwin}(t + 1, \operatorname{twopow}(k)), t + \operatorname{twopow}(k + 1) - (r + Q + 1)) = \operatorname{not}(\operatorname{wittbitbool}(\operatorname{wittwin}(t, \operatorname{twopow}(k)), t + \operatorname{twopow}(k + 1) - (r + Q + 1)))\).
\end{theorem}
\begin{proof}
The goal follows from \(\operatorname{winwittscalecrosszerogapsteponesamebitdoublingflipofdoublesandsinglewinwittagreeandzerodiagthree}(Q, t, k, r, a, L)\).
\end{proof}
\begin{theorem}
\label{ll:paper:lhcenter-late-cross-step-one-at-shift-plus-three-eq-not-of-single-winwittagree-and-zero-diag-five}
For every natural number \(Q\) and natural number \(t\) and natural number \(k\) and natural number \(r\) and natural number \(a\) and natural number \(L\), if \(t + Q \le r + 1\), then if \(t + Q + r + 2 \le k\), then if \(t + r + 5 \le k\), then if \(\operatorname{zerodiagonal}(t + Q + r)\), then if \(\operatorname{zerodiagonal}(t + r + 5)\), then if \(\operatorname{settledby}(t + Q + r + 2, t + Q + \operatorname{twopow}(k))\), then if \(\operatorname{settledby}(t + r + 5, t + \operatorname{twopow}(k))\), then if \(a \le t + \operatorname{twopow}(k + 1) - (r + 2)\), then if \(t + \operatorname{twopow}(k + 1) - (r + 2) \le a + L\), then if \(\operatorname{winwittagree}(t, k, Q, a, L)\), then \(\operatorname{lhcenter}(t + 1 + \operatorname{twopow}(k), \operatorname{twopow}(k) - (r + 3)) = \operatorname{not}(\operatorname{lhcenter}(t + \operatorname{twopow}(k), \operatorname{twopow}(k) - (r + 3)))\).
\end{theorem}
\begin{proof}
The goal follows from \(\operatorname{lhcenterlatecrosssteponeatshiftplusthreeeqnotofsinglewinwittagreeandzerodiagfive}(Q, t, k, r, a, L)\).
\end{proof}
\begin{theorem}
\label{ll:paper:lhcenter-late-cross-eq-r-plus-one-r-plus-three-of-single-winwittagree-and-zero-diag-five}
For every natural number \(Q\) and natural number \(t\) and natural number \(k\) and natural number \(r\) and natural number \(a\) and natural number \(L\), if \(t + Q \le r + 1\), then if \(t + Q + r + 2 \le k\), then if \(t + r + 5 \le k\), then if \(\operatorname{zerodiagonal}(t + Q + r)\), then if \(\operatorname{zerodiagonal}(t + r + 5)\), then if \(\operatorname{settledby}(t + Q + r + 2, t + Q + \operatorname{twopow}(k))\), then if \(\operatorname{settledby}(t + r + 5, t + \operatorname{twopow}(k))\), then if \(a \le t + \operatorname{twopow}(k + 1) - (r + 2)\), then if \(t + \operatorname{twopow}(k + 1) - (r + 2) \le a + L\), then if \(\operatorname{winwittagree}(t, k, Q, a, L)\), then \(\operatorname{lhcenter}(t + \operatorname{twopow}(k), \operatorname{twopow}(k) - (r + 1)) = \operatorname{lhcenter}(t + \operatorname{twopow}(k), \operatorname{twopow}(k) - (r + 3))\).
\end{theorem}
\begin{proof}
The goal follows from \(\operatorname{lhcenterlatecrosseqrplusonerplusthreeofsinglewinwittagreeandzerodiagfive}(Q, t, k, r, a, L)\).
\end{proof}
\begin{theorem}
\label{ll:paper:lhcenter-late-cross-step-two-eq-r-plus-one-of-single-winwittagree-and-zero-diag-five}
For every natural number \(Q\) and natural number \(t\) and natural number \(k\) and natural number \(r\) and natural number \(a\) and natural number \(L\), if \(t + Q < r\), then if \(t + Q + r + 2 \le k\), then if \(t + r + 5 \le k\), then if \(\operatorname{zerodiagonal}(t + Q + r)\), then if \(\operatorname{zerodiagonal}(t + r + 5)\), then if \(\operatorname{settledby}(t + Q + r, t + Q + \operatorname{twopow}(k))\), then if \(\operatorname{settledby}(t + Q + r + 2, t + Q + \operatorname{twopow}(k))\), then if \(\operatorname{settledby}(t + r + 5, t + \operatorname{twopow}(k))\), then if \(a \le t + \operatorname{twopow}(k + 1) - (r + 2)\), then if \(t + \operatorname{twopow}(k + 1) - r \le a + L\), then if \(\operatorname{winwittagree}(t, k, Q, a, L)\), then \(\operatorname{lhcenter}(t + 2 + \operatorname{twopow}(k), \operatorname{twopow}(k) - (r + 1)) = \operatorname{lhcenter}(t + \operatorname{twopow}(k), \operatorname{twopow}(k) - (r + 1))\).
\end{theorem}
\begin{proof}
The goal follows from \(\operatorname{lhcenterlatecrosssteptwoeqrplusoneofsinglewinwittagreeandzerodiagfive}(Q, t, k, r, a, L)\).
\end{proof}
\begin{theorem}
\label{ll:paper:lhcenter-late-cross-pin-r-sub-one-true-of-single-winwittagree-and-zero-diag-five}
For every natural number \(Q\) and natural number \(t\) and natural number \(k\) and natural number \(r\) and natural number \(a\) and natural number \(L\), if \(t + Q < r\), then if \(t + Q + r + 2 \le k\), then if \(t + r + 5 \le k\), then if \(\operatorname{zerodiagonal}(t + Q + r)\), then if \(\operatorname{zerodiagonal}(t + r + 5)\), then if \(\operatorname{settledby}(t + Q + r, t + Q + \operatorname{twopow}(k))\), then if \(\operatorname{settledby}(t + Q + r + 2, t + Q + \operatorname{twopow}(k))\), then if \(\operatorname{settledby}(t + r + 5, t + \operatorname{twopow}(k))\), then if \(a \le t + \operatorname{twopow}(k + 1) - (r + 2)\), then if \(t + \operatorname{twopow}(k + 1) - r \le a + L\), then if \(\operatorname{winwittagree}(t, k, Q, a, L)\), then \(\operatorname{lhcenter}(t + \operatorname{twopow}(k), \operatorname{twopow}(k) - (r - 1)) = \operatorname{true}\).
\end{theorem}
\begin{proof}
The goal follows from \(\operatorname{lhcenterlatecrosspinrsubonetrueofsinglewinwittagreeandzerodiagfive}(Q, t, k, r, a, L)\).
\end{proof}
\begin{theorem}
\label{ll:paper:lhcenter-late-cross-pin-r-plus-four-true-of-zero-diag-five-and-single-winwittagree}
For every natural number \(t\) and natural number \(k\) and natural number \(r\) and natural number \(a\) and natural number \(L\), if \(t + 6 < r\), then if \(t + 5 + r + 2 \le k\), then if \(\operatorname{zerodiagonal}(t + 5 + r)\), then if \(\operatorname{doublesat}(5, t + 5 + r)\), then if \(\operatorname{settledby}(t + 4 + r, t + \operatorname{twopow}(k))\), then if \(\operatorname{settledby}(t + 5 + r, t + \operatorname{twopow}(k))\), then if \(\operatorname{settledby}(t + 5 + r + 2, t + 5 + \operatorname{twopow}(k))\), then if \(a \le t + \operatorname{twopow}(k + 1) - (r + 2)\), then if \(t + \operatorname{twopow}(k + 1) - (r - 1) \le a + L\), then if \(\operatorname{winwittagree}(t, k, 5, a, L)\), then \(\operatorname{lhcenter}(t + \operatorname{twopow}(k), \operatorname{twopow}(k) - (r + 4)) = \operatorname{true}\).
\end{theorem}
\begin{proof}
The goal follows from \(\operatorname{lhcenterlatecrosspinrplusfourtrueofzerodiagfiveandsinglewinwittagree}(t, k, r, a, L)\).
\end{proof}
\begin{theorem}
\label{ll:paper:lhcenter-late-cross-pin-r-plus-three-true-of-zero-diag-five-and-single-winwittagree}
For every natural number \(t\) and natural number \(k\) and natural number \(r\) and natural number \(a\) and natural number \(L\), if \(t + 6 < r\), then if \(t + 5 + r + 2 \le k\), then if \(\operatorname{zerodiagonal}(t + 5 + r)\), then if \(\operatorname{doublesat}(5, t + 5 + r)\), then if \(\operatorname{settledby}(t + 4 + r, t + \operatorname{twopow}(k))\), then if \(\operatorname{settledby}(t + 5 + r, t + \operatorname{twopow}(k))\), then if \(\operatorname{settledby}(t + 5 + r + 2, t + 5 + \operatorname{twopow}(k))\), then if \(a \le t + \operatorname{twopow}(k + 1) - (r + 2)\), then if \(t + \operatorname{twopow}(k + 1) - (r - 1) \le a + L\), then if \(\operatorname{winwittagree}(t, k, 5, a, L)\), then \(\operatorname{lhcenter}(t + \operatorname{twopow}(k), \operatorname{twopow}(k) - (r + 3)) = \operatorname{true}\).
\end{theorem}
\begin{proof}
The goal follows from \(\operatorname{lhcenterlatecrosspinrplusthreetrueofzerodiagfiveandsinglewinwittagree}(t, k, r, a, L)\).
\end{proof}
\begin{theorem}
\label{ll:paper:lhcenter-late-cross-pin-r-plus-one-true-of-zero-diag-five-and-single-winwittagree}
For every natural number \(t\) and natural number \(k\) and natural number \(r\) and natural number \(a\) and natural number \(L\), if \(t + 6 < r\), then if \(t + 5 + r + 2 \le k\), then if \(\operatorname{zerodiagonal}(t + 5 + r)\), then if \(\operatorname{doublesat}(5, t + 5 + r)\), then if \(\operatorname{settledby}(t + 4 + r, t + \operatorname{twopow}(k))\), then if \(\operatorname{settledby}(t + 5 + r, t + \operatorname{twopow}(k))\), then if \(\operatorname{settledby}(t + 5 + r + 2, t + 5 + \operatorname{twopow}(k))\), then if \(a \le t + \operatorname{twopow}(k + 1) - (r + 2)\), then if \(t + \operatorname{twopow}(k + 1) - (r - 1) \le a + L\), then if \(\operatorname{winwittagree}(t, k, 5, a, L)\), then \(\operatorname{lhcenter}(t + \operatorname{twopow}(k), \operatorname{twopow}(k) - (r + 1)) = \operatorname{true}\).
\end{theorem}
\begin{proof}
The goal follows from \(\operatorname{lhcenterlatecrosspinrplusonetrueofzerodiagfiveandsinglewinwittagree}(t, k, r, a, L)\).
\end{proof}
\begin{theorem}
\label{ll:paper:lhcenter-late-cross-pin-doubling-diag-false-of-doubles-five-and-single-winwittagree}
For every natural number \(t\) and natural number \(k\) and natural number \(r\) and natural number \(a\) and natural number \(L\), if \(t + 6 < r\), then if \(t + 5 + r + 2 \le k\), then if \(\operatorname{zerodiagonal}(t + 5 + r)\), then if \(\operatorname{doublesat}(5, t + 5 + r)\), then if \(\operatorname{settledby}(t + 4 + r, t + \operatorname{twopow}(k))\), then if \(\operatorname{settledby}(t + 5 + r, t + \operatorname{twopow}(k))\), then if \(\operatorname{settledby}(t + 5 + r + 1, t + \operatorname{twopow}(k))\), then if \(\operatorname{settledby}(t + 5 + r + 2, t + 5 + \operatorname{twopow}(k))\), then if \(a \le t + \operatorname{twopow}(k + 1) - (r + 2)\), then if \(t + \operatorname{twopow}(k + 1) - (r - 1) \le a + L\), then if \(\operatorname{winwittagree}(t, k, 5, a, L)\), then \(\operatorname{lhcenter}(t + \operatorname{twopow}(k), \operatorname{twopow}(k) - (r + 6)) = \operatorname{false}\).
\end{theorem}
\begin{proof}
The goal follows from \(\operatorname{lhcenterlatecrosspindoublingdiagfalseofdoublesfiveandsinglewinwittagree}(t, k, r, a, L)\).
\end{proof}
\begin{theorem}
\label{ll:paper:lhcenter-late-cross-step-one-pin-r-plus-five-true-of-doubles-five-and-single-winwittagree}
For every natural number \(t\) and natural number \(k\) and natural number \(r\) and natural number \(a\) and natural number \(L\), if \(t + 6 < r\), then if \(t + 5 + r + 2 \le k\), then if \(\operatorname{zerodiagonal}(t + 5 + r)\), then if \(\operatorname{doublesat}(5, t + 5 + r)\), then if \(\operatorname{settledby}(t + 4 + r, t + \operatorname{twopow}(k))\), then if \(\operatorname{settledby}(t + 5 + r, t + \operatorname{twopow}(k))\), then if \(\operatorname{settledby}(t + 5 + r + 1, t + \operatorname{twopow}(k))\), then if \(\operatorname{settledby}(t + 5 + r + 2, t + 5 + \operatorname{twopow}(k))\), then if \(a \le t + \operatorname{twopow}(k + 1) - (r + 2)\), then if \(t + \operatorname{twopow}(k + 1) - (r - 1) \le a + L\), then if \(\operatorname{winwittagree}(t, k, 5, a, L)\), then \(\operatorname{lhcenter}(t + 1 + \operatorname{twopow}(k), \operatorname{twopow}(k) - (r + 5)) = \operatorname{true}\).
\end{theorem}
\begin{proof}
The goal follows from \(\operatorname{lhcenterlatecrosssteponepinrplusfivetrueofdoublesfiveandsinglewinwittagree}(t, k, r, a, L)\).
\end{proof}
\begin{theorem}
\label{ll:paper:lhcenter-late-cross-step-two-pin-r-plus-four-true-of-doubles-five-and-single-winwittagree}
For every natural number \(t\) and natural number \(k\) and natural number \(r\) and natural number \(a\) and natural number \(L\), if \(t + 6 < r\), then if \(t + 5 + r + 2 \le k\), then if \(\operatorname{zerodiagonal}(t + 5 + r)\), then if \(\operatorname{doublesat}(5, t + 5 + r)\), then if \(\operatorname{settledby}(t + 4 + r, t + \operatorname{twopow}(k))\), then if \(\operatorname{settledby}(t + 5 + r, t + \operatorname{twopow}(k))\), then if \(\operatorname{settledby}(t + 5 + r + 1, t + \operatorname{twopow}(k))\), then if \(\operatorname{settledby}(t + 5 + r + 2, t + 5 + \operatorname{twopow}(k))\), then if \(a \le t + \operatorname{twopow}(k + 1) - (r + 2)\), then if \(t + \operatorname{twopow}(k + 1) - (r - 1) \le a + L\), then if \(\operatorname{winwittagree}(t, k, 5, a, L)\), then \(\operatorname{lhcenter}(t + 2 + \operatorname{twopow}(k), \operatorname{twopow}(k) - (r + 4)) = \operatorname{true}\).
\end{theorem}
\begin{proof}
The goal follows from \(\operatorname{lhcenterlatecrosssteptwopinrplusfourtrueofdoublesfiveandsinglewinwittagree}(t, k, r, a, L)\).
\end{proof}
\begin{theorem}
\label{ll:paper:lhcenter-late-cross-step-three-pin-r-plus-three-false-of-doubles-five-and-single-winwittagree}
For every natural number \(t\) and natural number \(k\) and natural number \(r\) and natural number \(a\) and natural number \(L\), if \(t + 6 < r\), then if \(t + 5 + r + 2 \le k\), then if \(\operatorname{zerodiagonal}(t + 5 + r)\), then if \(\operatorname{doublesat}(5, t + 5 + r)\), then if \(\operatorname{settledby}(t + 4 + r, t + \operatorname{twopow}(k))\), then if \(\operatorname{settledby}(t + 5 + r, t + \operatorname{twopow}(k))\), then if \(\operatorname{settledby}(t + 5 + r + 1, t + \operatorname{twopow}(k))\), then if \(\operatorname{settledby}(t + 5 + r + 2, t + 5 + \operatorname{twopow}(k))\), then if \(a \le t + \operatorname{twopow}(k + 1) - (r + 2)\), then if \(t + \operatorname{twopow}(k + 1) - (r - 1) \le a + L\), then if \(\operatorname{winwittagree}(t, k, 5, a, L)\), then \(\operatorname{lhcenter}(t + 3 + \operatorname{twopow}(k), \operatorname{twopow}(k) - (r + 3)) = \operatorname{false}\).
\end{theorem}
\begin{proof}
The goal follows from \(\operatorname{lhcenterlatecrossstepthreepinrplusthreefalseofdoublesfiveandsinglewinwittagree}(t, k, r, a, L)\).
\end{proof}
\begin{theorem}
\label{ll:paper:not-winwittagree-five-of-doubles}
For every natural number \(t\) and natural number \(k\) and natural number \(r\) and natural number \(a\) and natural number \(L\), if \(t + 6 < r\), then if \(t + 5 + r + 2 \le k\), then if \(\operatorname{zerodiagonal}(t + 5 + r)\), then if \(\operatorname{doublesat}(5, t + 5 + r)\), then if \(\operatorname{settledby}(t + 4 + r, t + \operatorname{twopow}(k))\), then if \(\operatorname{settledby}(t + 5 + r, t + \operatorname{twopow}(k))\), then if \(\operatorname{settledby}(t + 5 + r + 1, t + \operatorname{twopow}(k))\), then if \(\operatorname{settledby}(t + 5 + r + 2, t + 3 + \operatorname{twopow}(k))\), then if \(a \le t + \operatorname{twopow}(k + 1) - (r + 4)\), then if \(t + \operatorname{twopow}(k + 1) - (r - 1) \le a + L\), then not \(\operatorname{winwittagree}(t, k, 5, a, L)\).
\end{theorem}
\begin{proof}
The goal follows from \(\operatorname{notwinwittagreefiveofdoubles}(t, k, r, a, L)\).
\end{proof}
\begin{theorem}
\label{ll:paper:lhcenter-late-cross-step-four-pin-r-plus-four-true-of-doubles-five-and-single-winwittagree}
For every natural number \(t\) and natural number \(k\) and natural number \(r\) and natural number \(a\) and natural number \(L\), if \(t + 6 < r\), then if \(t + 5 + r + 2 \le k\), then if \(\operatorname{zerodiagonal}(t + 5 + r)\), then if \(\operatorname{doublesat}(5, t + 5 + r)\), then if \(\operatorname{settledby}(t + 4 + r, t + \operatorname{twopow}(k))\), then if \(\operatorname{settledby}(t + 5 + r, t + \operatorname{twopow}(k))\), then if \(\operatorname{settledby}(t + 5 + r + 1, t + \operatorname{twopow}(k))\), then if \(\operatorname{settledby}(t + 5 + r + 2, t + 3 + \operatorname{twopow}(k))\), then if \(a \le t + \operatorname{twopow}(k + 1) - (r + 2)\), then if \(t + \operatorname{twopow}(k + 1) - (r - 1) \le a + L\), then if \(\operatorname{winwittagree}(t, k, 5, a, L)\), then \(\operatorname{lhcenter}(t + 4 + \operatorname{twopow}(k), \operatorname{twopow}(k) - (r + 4)) = \operatorname{true}\).
\end{theorem}
\begin{proof}
The goal follows from \(\operatorname{lhcenterlatecrossstepfourpinrplusfourtrueofdoublesfiveandsinglewinwittagree}(t, k, r, a, L)\).
\end{proof}
\begin{theorem}
\label{ll:paper:lhcenter-late-cross-step-five-pin-r-plus-four-false-of-doubles-five-and-single-winwittagree}
For every natural number \(t\) and natural number \(k\) and natural number \(r\) and natural number \(a\) and natural number \(L\), if \(t + 6 < r\), then if \(t + 5 + r + 2 \le k\), then if \(\operatorname{zerodiagonal}(t + 5 + r)\), then if \(\operatorname{doublesat}(5, t + 5 + r)\), then if \(\operatorname{settledby}(t + 4 + r, t + \operatorname{twopow}(k))\), then if \(\operatorname{settledby}(t + 5 + r, t + \operatorname{twopow}(k))\), then if \(\operatorname{settledby}(t + 5 + r + 1, t + \operatorname{twopow}(k))\), then if \(\operatorname{settledby}(t + 5 + r + 2, t + 3 + \operatorname{twopow}(k))\), then if \(a \le t + \operatorname{twopow}(k + 1) - (r + 2)\), then if \(t + \operatorname{twopow}(k + 1) - (r - 1) \le a + L\), then if \(\operatorname{winwittagree}(t, k, 5, a, L)\), then \(\operatorname{lhcenter}(t + 5 + \operatorname{twopow}(k), \operatorname{twopow}(k) - (r + 4)) = \operatorname{false}\).
\end{theorem}
\begin{proof}
The goal follows from \(\operatorname{lhcenterlatecrossstepfivepinrplusfourfalseofdoublesfiveandsinglewinwittagree}(t, k, r, a, L)\).
\end{proof}
\begin{theorem}
\label{ll:paper:winwitt-scale-cross-zero-gap-pin-r-sub-one-true-of-single-winwittagree-and-zero-diag-five}
For every natural number \(Q\) and natural number \(t\) and natural number \(k\) and natural number \(r\) and natural number \(a\) and natural number \(L\), if \(t + Q + 1 < r\), then if \(t + Q + r + 2 \le k\), then if \(t + r + 5 \le k\), then if \(\operatorname{zerodiagonal}(t + Q + r)\), then if \(\operatorname{zerodiagonal}(t + r + 5)\), then if \(\operatorname{settledby}(t + Q + r, t + Q + \operatorname{twopow}(k))\), then if \(\operatorname{settledby}(t + Q + r + 2, t + Q + \operatorname{twopow}(k))\), then if \(\operatorname{settledby}(t + r + 5, t + \operatorname{twopow}(k))\), then if \(a \le t + \operatorname{twopow}(k + 1) - (r + 2)\), then if \(t + \operatorname{twopow}(k + 1) - r \le a + L\), then if \(\operatorname{winwittagree}(t, k, Q, a, L)\), then \(\operatorname{wittbitbool}(\operatorname{wittwin}(t, \operatorname{twopow}(k)), t + \operatorname{twopow}(k + 1) - (r - 1)) = \operatorname{true}\).
\end{theorem}
\begin{proof}
The goal follows from \(\operatorname{winwittscalecrosszerogappinrsubonetrueofsinglewinwittagreeandzerodiagfive}(Q, t, k, r, a, L)\).
\end{proof}
\begin{theorem}
\label{ll:paper:winwitt-scale-cross-zero-gap-pin-r-plus-four-true-of-doubles-five-and-single-winwittagree}
For every natural number \(t\) and natural number \(k\) and natural number \(r\) and natural number \(a\) and natural number \(L\), if \(t + 6 < r\), then if \(t + 5 + r + 2 \le k\), then if \(\operatorname{zerodiagonal}(t + 5 + r)\), then if \(\operatorname{doublesat}(5, t + 5 + r)\), then if \(\operatorname{settledby}(t + 4 + r, t + \operatorname{twopow}(k))\), then if \(\operatorname{settledby}(t + 5 + r, t + \operatorname{twopow}(k))\), then if \(\operatorname{settledby}(t + 5 + r + 2, t + 5 + \operatorname{twopow}(k))\), then if \(a \le t + \operatorname{twopow}(k + 1) - (r + 2)\), then if \(t + \operatorname{twopow}(k + 1) - (r - 1) \le a + L\), then if \(\operatorname{winwittagree}(t, k, 5, a, L)\), then \(\operatorname{wittbitbool}(\operatorname{wittwin}(t, \operatorname{twopow}(k)), t + \operatorname{twopow}(k + 1) - (r + 4)) = \operatorname{true}\).
\end{theorem}
\begin{proof}
The goal follows from \(\operatorname{winwittscalecrosszerogappinrplusfourtrueofdoublesfiveandsinglewinwittagree}(t, k, r, a, L)\).
\end{proof}
\begin{theorem}
\label{ll:paper:winwitt-scale-cross-zero-gap-pin-r-plus-three-true-of-doubles-five-and-single-winwittagree}
For every natural number \(t\) and natural number \(k\) and natural number \(r\) and natural number \(a\) and natural number \(L\), if \(t + 6 < r\), then if \(t + 5 + r + 2 \le k\), then if \(\operatorname{zerodiagonal}(t + 5 + r)\), then if \(\operatorname{doublesat}(5, t + 5 + r)\), then if \(\operatorname{settledby}(t + 4 + r, t + \operatorname{twopow}(k))\), then if \(\operatorname{settledby}(t + 5 + r, t + \operatorname{twopow}(k))\), then if \(\operatorname{settledby}(t + 5 + r + 2, t + 5 + \operatorname{twopow}(k))\), then if \(a \le t + \operatorname{twopow}(k + 1) - (r + 2)\), then if \(t + \operatorname{twopow}(k + 1) - (r - 1) \le a + L\), then if \(\operatorname{winwittagree}(t, k, 5, a, L)\), then \(\operatorname{wittbitbool}(\operatorname{wittwin}(t, \operatorname{twopow}(k)), t + \operatorname{twopow}(k + 1) - (r + 3)) = \operatorname{true}\).
\end{theorem}
\begin{proof}
The goal follows from \(\operatorname{winwittscalecrosszerogappinrplusthreetrueofdoublesfiveandsinglewinwittagree}(t, k, r, a, L)\).
\end{proof}
\begin{theorem}
\label{ll:paper:winwitt-scale-cross-zero-gap-pin-r-plus-one-true-of-doubles-five-and-single-winwittagree}
For every natural number \(t\) and natural number \(k\) and natural number \(r\) and natural number \(a\) and natural number \(L\), if \(t + 6 < r\), then if \(t + 5 + r + 2 \le k\), then if \(\operatorname{zerodiagonal}(t + 5 + r)\), then if \(\operatorname{doublesat}(5, t + 5 + r)\), then if \(\operatorname{settledby}(t + 4 + r, t + \operatorname{twopow}(k))\), then if \(\operatorname{settledby}(t + 5 + r, t + \operatorname{twopow}(k))\), then if \(\operatorname{settledby}(t + 5 + r + 2, t + 5 + \operatorname{twopow}(k))\), then if \(a \le t + \operatorname{twopow}(k + 1) - (r + 2)\), then if \(t + \operatorname{twopow}(k + 1) - (r - 1) \le a + L\), then if \(\operatorname{winwittagree}(t, k, 5, a, L)\), then \(\operatorname{wittbitbool}(\operatorname{wittwin}(t, \operatorname{twopow}(k)), t + \operatorname{twopow}(k + 1) - (r + 1)) = \operatorname{true}\).
\end{theorem}
\begin{proof}
The goal follows from \(\operatorname{winwittscalecrosszerogappinrplusonetrueofdoublesfiveandsinglewinwittagree}(t, k, r, a, L)\).
\end{proof}
\begin{theorem}
\label{ll:paper:winwitt-scale-cross-zero-gap-pin-doubling-diag-false-of-doubles-five-and-single-winwittagree}
For every natural number \(t\) and natural number \(k\) and natural number \(r\) and natural number \(a\) and natural number \(L\), if \(t + 6 < r\), then if \(t + 5 + r + 2 \le k\), then if \(\operatorname{zerodiagonal}(t + 5 + r)\), then if \(\operatorname{doublesat}(5, t + 5 + r)\), then if \(\operatorname{settledby}(t + 4 + r, t + \operatorname{twopow}(k))\), then if \(\operatorname{settledby}(t + 5 + r, t + \operatorname{twopow}(k))\), then if \(\operatorname{settledby}(t + 5 + r + 1, t + \operatorname{twopow}(k))\), then if \(\operatorname{settledby}(t + 5 + r + 2, t + 5 + \operatorname{twopow}(k))\), then if \(a \le t + \operatorname{twopow}(k + 1) - (r + 2)\), then if \(t + \operatorname{twopow}(k + 1) - (r - 1) \le a + L\), then if \(\operatorname{winwittagree}(t, k, 5, a, L)\), then \(\operatorname{wittbitbool}(\operatorname{wittwin}(t, \operatorname{twopow}(k)), t + \operatorname{twopow}(k + 1) - (r + 6)) = \operatorname{false}\).
\end{theorem}
\begin{proof}
The goal follows from \(\operatorname{winwittscalecrosszerogappindoublingdiagfalseofdoublesfiveandsinglewinwittagree}(t, k, r, a, L)\).
\end{proof}
\begin{theorem}
\label{ll:paper:winwitt-scale-cross-zero-gap-step-one-pin-r-plus-four-true-of-doubles-five-and-single-winwittagree}
For every natural number \(t\) and natural number \(k\) and natural number \(r\) and natural number \(a\) and natural number \(L\), if \(t + 6 < r\), then if \(t + 5 + r + 2 \le k\), then if \(\operatorname{zerodiagonal}(t + 5 + r)\), then if \(\operatorname{doublesat}(5, t + 5 + r)\), then if \(\operatorname{settledby}(t + 4 + r, t + \operatorname{twopow}(k))\), then if \(\operatorname{settledby}(t + 5 + r, t + \operatorname{twopow}(k))\), then if \(\operatorname{settledby}(t + 5 + r + 1, t + \operatorname{twopow}(k))\), then if \(\operatorname{settledby}(t + 5 + r + 2, t + 5 + \operatorname{twopow}(k))\), then if \(a \le t + \operatorname{twopow}(k + 1) - (r + 2)\), then if \(t + \operatorname{twopow}(k + 1) - (r - 1) \le a + L\), then if \(\operatorname{winwittagree}(t, k, 5, a, L)\), then \(\operatorname{wittbitbool}(\operatorname{wittwin}(t + 1, \operatorname{twopow}(k)), t + \operatorname{twopow}(k + 1) - (r + 4)) = \operatorname{true}\).
\end{theorem}
\begin{proof}
The goal follows from \(\operatorname{winwittscalecrosszerogapsteponepinrplusfourtrueofdoublesfiveandsinglewinwittagree}(t, k, r, a, L)\).
\end{proof}
\begin{theorem}
\label{ll:paper:winwitt-scale-cross-zero-gap-step-two-pin-r-plus-two-true-of-doubles-five-and-single-winwittagree}
For every natural number \(t\) and natural number \(k\) and natural number \(r\) and natural number \(a\) and natural number \(L\), if \(t + 6 < r\), then if \(t + 5 + r + 2 \le k\), then if \(\operatorname{zerodiagonal}(t + 5 + r)\), then if \(\operatorname{doublesat}(5, t + 5 + r)\), then if \(\operatorname{settledby}(t + 4 + r, t + \operatorname{twopow}(k))\), then if \(\operatorname{settledby}(t + 5 + r, t + \operatorname{twopow}(k))\), then if \(\operatorname{settledby}(t + 5 + r + 1, t + \operatorname{twopow}(k))\), then if \(\operatorname{settledby}(t + 5 + r + 2, t + 5 + \operatorname{twopow}(k))\), then if \(a \le t + \operatorname{twopow}(k + 1) - (r + 2)\), then if \(t + \operatorname{twopow}(k + 1) - (r - 1) \le a + L\), then if \(\operatorname{winwittagree}(t, k, 5, a, L)\), then \(\operatorname{wittbitbool}(\operatorname{wittwin}(t + 2, \operatorname{twopow}(k)), t + \operatorname{twopow}(k + 1) - (r + 2)) = \operatorname{true}\).
\end{theorem}
\begin{proof}
The goal follows from \(\operatorname{winwittscalecrosszerogapsteptwopinrplustwotrueofdoublesfiveandsinglewinwittagree}(t, k, r, a, L)\).
\end{proof}
\begin{theorem}
\label{ll:paper:winwitt-scale-cross-zero-gap-step-three-pin-r-false-of-doubles-five-and-single-winwittagree}
For every natural number \(t\) and natural number \(k\) and natural number \(r\) and natural number \(a\) and natural number \(L\), if \(t + 6 < r\), then if \(t + 5 + r + 2 \le k\), then if \(\operatorname{zerodiagonal}(t + 5 + r)\), then if \(\operatorname{doublesat}(5, t + 5 + r)\), then if \(\operatorname{settledby}(t + 4 + r, t + \operatorname{twopow}(k))\), then if \(\operatorname{settledby}(t + 5 + r, t + \operatorname{twopow}(k))\), then if \(\operatorname{settledby}(t + 5 + r + 1, t + \operatorname{twopow}(k))\), then if \(\operatorname{settledby}(t + 5 + r + 2, t + 5 + \operatorname{twopow}(k))\), then if \(a \le t + \operatorname{twopow}(k + 1) - (r + 2)\), then if \(t + \operatorname{twopow}(k + 1) - (r - 1) \le a + L\), then if \(\operatorname{winwittagree}(t, k, 5, a, L)\), then \(\operatorname{wittbitbool}(\operatorname{wittwin}(t + 3, \operatorname{twopow}(k)), t + \operatorname{twopow}(k + 1) - r) = \operatorname{false}\).
\end{theorem}
\begin{proof}
The goal follows from \(\operatorname{winwittscalecrosszerogapstepthreepinrfalseofdoublesfiveandsinglewinwittagree}(t, k, r, a, L)\).
\end{proof}
\begin{theorem}
\label{ll:paper:winwitt-scale-cross-zero-gap-step-four-pin-r-true-of-doubles-five-and-single-winwittagree}
For every natural number \(t\) and natural number \(k\) and natural number \(r\) and natural number \(a\) and natural number \(L\), if \(t + 6 < r\), then if \(t + 5 + r + 2 \le k\), then if \(\operatorname{zerodiagonal}(t + 5 + r)\), then if \(\operatorname{doublesat}(5, t + 5 + r)\), then if \(\operatorname{settledby}(t + 4 + r, t + \operatorname{twopow}(k))\), then if \(\operatorname{settledby}(t + 5 + r, t + \operatorname{twopow}(k))\), then if \(\operatorname{settledby}(t + 5 + r + 1, t + \operatorname{twopow}(k))\), then if \(\operatorname{settledby}(t + 5 + r + 2, t + 3 + \operatorname{twopow}(k))\), then if \(a \le t + \operatorname{twopow}(k + 1) - (r + 2)\), then if \(t + \operatorname{twopow}(k + 1) - (r - 1) \le a + L\), then if \(\operatorname{winwittagree}(t, k, 5, a, L)\), then \(\operatorname{wittbitbool}(\operatorname{wittwin}(t + 4, \operatorname{twopow}(k)), t + \operatorname{twopow}(k + 1) - r) = \operatorname{true}\).
\end{theorem}
\begin{proof}
The goal follows from \(\operatorname{winwittscalecrosszerogapstepfourpinrtrueofdoublesfiveandsinglewinwittagree}(t, k, r, a, L)\).
\end{proof}
\begin{theorem}
\label{ll:paper:winwitt-scale-cross-zero-gap-step-five-pin-r-sub-one-false-of-doubles-five-and-single-winwittagree}
For every natural number \(t\) and natural number \(k\) and natural number \(r\) and natural number \(a\) and natural number \(L\), if \(t + 6 < r\), then if \(t + 5 + r + 2 \le k\), then if \(\operatorname{zerodiagonal}(t + 5 + r)\), then if \(\operatorname{doublesat}(5, t + 5 + r)\), then if \(\operatorname{settledby}(t + 4 + r, t + \operatorname{twopow}(k))\), then if \(\operatorname{settledby}(t + 5 + r, t + \operatorname{twopow}(k))\), then if \(\operatorname{settledby}(t + 5 + r + 1, t + \operatorname{twopow}(k))\), then if \(\operatorname{settledby}(t + 5 + r + 2, t + 3 + \operatorname{twopow}(k))\), then if \(a \le t + \operatorname{twopow}(k + 1) - (r + 2)\), then if \(t + \operatorname{twopow}(k + 1) - (r - 1) \le a + L\), then if \(\operatorname{winwittagree}(t, k, 5, a, L)\), then \(\operatorname{wittbitbool}(\operatorname{wittwin}(t + 5, \operatorname{twopow}(k)), t + \operatorname{twopow}(k + 1) - (r - 1)) = \operatorname{false}\).
\end{theorem}
\begin{proof}
The goal follows from \(\operatorname{winwittscalecrosszerogapstepfivepinrsubonefalseofdoublesfiveandsinglewinwittagree}(t, k, r, a, L)\).
\end{proof}
\begin{theorem}
\label{ll:paper:lhcenter-late-cross-step-two-pin-r-plus-two-true-of-doubles-five-and-single-winwittagree}
For every natural number \(t\) and natural number \(k\) and natural number \(r\) and natural number \(a\) and natural number \(L\), if \(t + 6 < r\), then if \(t + 5 + r + 2 \le k\), then if \(\operatorname{zerodiagonal}(t + 5 + r)\), then if \(\operatorname{doublesat}(5, t + 5 + r)\), then if \(\operatorname{settledby}(t + 4 + r, t + \operatorname{twopow}(k))\), then if \(\operatorname{settledby}(t + 5 + r, t + \operatorname{twopow}(k))\), then if \(\operatorname{settledby}(t + 5 + r + 2, t + 5 + \operatorname{twopow}(k))\), then if \(a \le t + \operatorname{twopow}(k + 1) - (r + 2)\), then if \(t + \operatorname{twopow}(k + 1) - (r - 1) \le a + L\), then if \(\operatorname{winwittagree}(t, k, 5, a, L)\), then \(\operatorname{lhcenter}(t + 2 + \operatorname{twopow}(k), \operatorname{twopow}(k) - (r + 2)) = \operatorname{true}\).
\end{theorem}
\begin{proof}
The goal follows from \(\operatorname{lhcenterlatecrosssteptwopinrplustwotrueofdoublesfiveandsinglewinwittagree}(t, k, r, a, L)\).
\end{proof}
\begin{theorem}
\label{ll:paper:winwitt-scale-cross-zero-gap-step-two-pin-r-true-of-doubles-five-and-single-winwittagree}
For every natural number \(t\) and natural number \(k\) and natural number \(r\) and natural number \(a\) and natural number \(L\), if \(t + 6 < r\), then if \(t + 5 + r + 2 \le k\), then if \(\operatorname{zerodiagonal}(t + 5 + r)\), then if \(\operatorname{doublesat}(5, t + 5 + r)\), then if \(\operatorname{settledby}(t + 4 + r, t + \operatorname{twopow}(k))\), then if \(\operatorname{settledby}(t + 5 + r, t + \operatorname{twopow}(k))\), then if \(\operatorname{settledby}(t + 5 + r + 2, t + 5 + \operatorname{twopow}(k))\), then if \(a \le t + \operatorname{twopow}(k + 1) - (r + 2)\), then if \(t + \operatorname{twopow}(k + 1) - (r - 1) \le a + L\), then if \(\operatorname{winwittagree}(t, k, 5, a, L)\), then \(\operatorname{wittbitbool}(\operatorname{wittwin}(t + 2, \operatorname{twopow}(k)), t + \operatorname{twopow}(k + 1) - r) = \operatorname{true}\).
\end{theorem}
\begin{proof}
The goal follows from \(\operatorname{winwittscalecrosszerogapsteptwopinrtrueofdoublesfiveandsinglewinwittagree}(t, k, r, a, L)\).
\end{proof}
\begin{theorem}
\label{ll:paper:winwitt-scale-cross-zero-gap-same-bit-r-sub-one-flip-five-of-doubles-five-and-single-winwittagree}
For every natural number \(t\) and natural number \(k\) and natural number \(r\) and natural number \(a\) and natural number \(L\), if \(t + 6 < r\), then if \(t + 5 + r + 2 \le k\), then if \(\operatorname{zerodiagonal}(t + 5 + r)\), then if \(\operatorname{doublesat}(5, t + 5 + r)\), then if \(\operatorname{settledby}(t + 4 + r, t + \operatorname{twopow}(k))\), then if \(\operatorname{settledby}(t + 5 + r, t + \operatorname{twopow}(k))\), then if \(\operatorname{settledby}(t + 5 + r + 1, t + \operatorname{twopow}(k))\), then if \(\operatorname{settledby}(t + 5 + r + 2, t + 3 + \operatorname{twopow}(k))\), then if \(a \le t + \operatorname{twopow}(k + 1) - (r + 2)\), then if \(t + \operatorname{twopow}(k + 1) - (r - 1) \le a + L\), then if \(\operatorname{winwittagree}(t, k, 5, a, L)\), then \(\operatorname{wittbitbool}(\operatorname{wittwin}(t + 5, \operatorname{twopow}(k)), t + \operatorname{twopow}(k + 1) - (r - 1)) = \operatorname{not}(\operatorname{wittbitbool}(\operatorname{wittwin}(t, \operatorname{twopow}(k)), t + \operatorname{twopow}(k + 1) - (r - 1)))\).
\end{theorem}
\begin{proof}
The goal follows from \(\operatorname{winwittscalecrosszerogapsamebitrsuboneflipfiveofdoublesfiveandsinglewinwittagree}(t, k, r, a, L)\).
\end{proof}
\begin{theorem}
\label{ll:paper:not-centerperiodic-of-rhcolagree-one-horizon}
For every natural number \(T\) and natural number \(P\) and natural number \(m\) and natural number \(t\) and natural number \(S\), if \(0 < P\), then if \(0 < m\), then if \(T \le t\), then if \(t + m \cdot P \le S + 1\), then if \(\operatorname{rhcolagree}(t, m \cdot P, 1, S)\), then not \(\operatorname{centerperiodic}(T, P)\).
\end{theorem}
\begin{proof}
The goal follows from \(\operatorname{notcenterperiodicofrhcolagreeonehorizon}(T, P, m, t, S)\).
\end{proof}
\begin{theorem}
\label{ll:paper:not-centerperiodic-of-lhcolagree-one-horizon}
For every natural number \(T\) and natural number \(P\) and natural number \(m\) and natural number \(t\) and natural number \(S\), if \(0 < P\), then if \(0 < m\), then if \(T \le t\), then if \(t + m \cdot P \le S + 1\), then if \(\operatorname{lhcolagree}(t, m \cdot P, 1, S)\), then not \(\operatorname{centerperiodic}(T, P)\).
\end{theorem}
\begin{proof}
The goal follows from \(\operatorname{notcenterperiodicoflhcolagreeonehorizon}(T, P, m, t, S)\).
\end{proof}
\begin{theorem}
\label{ll:paper:not-centerperiodic-of-rhbandagree-zero-horizon}
For every natural number \(T\) and natural number \(P\) and natural number \(m\) and natural number \(t\) and natural number \(S\), if \(0 < P\), then if \(0 < m\), then if \(T \le t\), then if \(t + m \cdot P \le S + 1\), then if \(\operatorname{rhbandagree}(t, m \cdot P, 0, S + 1)\), then not \(\operatorname{centerperiodic}(T, P)\).
\end{theorem}
\begin{proof}
The goal follows from \(\operatorname{notcenterperiodicofrhbandagreezerohorizon}(T, P, m, t, S)\).
\end{proof}
\begin{theorem}
\label{ll:paper:not-centerperiodic-of-rhcenter-inward-determine-horizon}
For every natural number \(T\) and natural number \(P\) and natural number \(m\) and natural number \(t\) and natural number \(S\) and natural number \(K\), if \(0 < P\), then if \(0 < m\), then if \(T \le t\), then if \(t + m \cdot P \le S + 1\), then if \(\operatorname{rhcolagree}(t, m \cdot P, 1 + K, S + K)\), then if \(\operatorname{rhoragree}(t, m \cdot P, 1 + K, S + K)\), then not \(\operatorname{centerperiodic}(T, P)\).
\end{theorem}
\begin{proof}
The goal follows from \(\operatorname{notcenterperiodicofrhcenterinwarddeterminehorizon}(T, P, m, t, S, K)\).
\end{proof}
\begin{theorem}
\label{ll:paper:periodic-rhcolagree-one-of-rhoragree-one}
For every natural number \(T\) and natural number \(P\), if \(0 < P\), then if \(\operatorname{centerperiodic}(T, P)\), then for every natural number \(m\) and natural number \(t\) and natural number \(S\), if \(T \le t\), then if \(\operatorname{rhcenter}(t + m \cdot P, 1) = \operatorname{rhcenter}(t, 1)\), then if \(\operatorname{rhoragree}(t, m \cdot P, 1, S)\), then \(\operatorname{rhcolagree}(t, m \cdot P, 1, S)\).
\end{theorem}
\begin{proof}
Assume \(hP\), \(hper\).
The goal follows from \(\operatorname{periodicrhcolagreeoneofrhoragreeone}(T, P, hP, hper)\).
\end{proof}
\begin{theorem}
\label{ll:paper:not-centerperiodic-of-rhoragree-one-horizon}
For every natural number \(T\) and natural number \(P\) and natural number \(m\) and natural number \(t\) and natural number \(S\), if \(0 < P\), then if \(0 < m\), then if \(T \le t\), then if \(t + m \cdot P \le S + 1\), then if \(\operatorname{rhcenter}(t + m \cdot P, 1) = \operatorname{rhcenter}(t, 1)\), then if \(\operatorname{rhoragree}(t, m \cdot P, 1, S)\), then not \(\operatorname{centerperiodic}(T, P)\).
\end{theorem}
\begin{proof}
The goal follows from \(\operatorname{notcenterperiodicofrhoragreeonehorizon}(T, P, m, t, S)\).
\end{proof}
\begin{theorem}
\label{ll:paper:ldiag-settled-periodic-of-doubles-le}
For every natural number \(Q\) and natural number \(z\) and natural number \(d\) and natural number \(t\), if \(\operatorname{doublesat}(Q, z)\), then if \(d \le z\), then if \(\operatorname{settledby}(d, t)\), then \(\operatorname{ldiag}(d, t + Q) = \operatorname{ldiag}(d, t)\).
\end{theorem}
\begin{proof}
The goal follows from \(\operatorname{ldiagsettledperiodicofdoublesle}(Q, z, d, t)\).
\end{proof}
\begin{theorem}
\label{ll:paper:ldiag-settled-antiperiodic-of-doubles}
For every natural number \(Q\) and natural number \(z\) and natural number \(t\), if \(\operatorname{zerodiagonal}(z)\), then if \(\operatorname{doublesat}(Q, z)\), then if \(\operatorname{settledby}(z + 1, t)\), then \(\operatorname{ldiag}(z + 1, t + Q) = \operatorname{not}(\operatorname{ldiag}(z + 1, t))\).
\end{theorem}
\begin{proof}
The goal follows from \(\operatorname{ldiagsettledantiperiodicofdoubles}(Q, z, t)\).
\end{proof}
\begin{theorem}
\label{ll:paper:lhcenter-diag-shift-periodic-of-doubles-le}
For every natural number \(Q\) and natural number \(z\) and natural number \(t\) and natural number \(p\), if \(\operatorname{doublesat}(Q, z)\), then if \(t - p \le z\), then if \(\operatorname{settledby}(t - p, t)\), then \(\operatorname{lhcenter}(t + Q, p + Q) = \operatorname{lhcenter}(t, p)\).
\end{theorem}
\begin{proof}
The goal follows from \(\operatorname{lhcenterdiagshiftperiodicofdoublesle}(Q, z, t, p)\).
\end{proof}
\begin{theorem}
\label{ll:paper:lhcenter-diag-shift-antiperiodic-of-doubles}
For every natural number \(Q\) and natural number \(t\) and natural number \(p\), if \(p \le t\), then if \(\operatorname{zerodiagonal}(t - p)\), then if \(\operatorname{doublesat}(Q, t - p)\), then if \(\operatorname{settledby}(t - p + 1, t + 1)\), then \(\operatorname{lhcenter}(t + 1 + Q, p + Q) = \operatorname{not}(\operatorname{lhcenter}(t + 1, p))\).
\end{theorem}
\begin{proof}
The goal follows from \(\operatorname{lhcenterdiagshiftantiperiodicofdoubles}(Q, t, p)\).
\end{proof}
\begin{theorem}
\label{ll:paper:lhcenter-spatial-periodic-of-doubles-le-and-lhbandagree}
For every natural number \(Q\) and natural number \(z\) and natural number \(t\) and natural number \(p\) and natural number \(p\_0\) and natural number \(L\), if \(\operatorname{doublesat}(Q, z)\), then if \(t - p \le z\), then if \(\operatorname{settledby}(t - p, t)\), then if \(p\_0 \le p + Q\), then if \(p + Q \le p\_0 + L\), then if \(\operatorname{lhbandagree}(t, Q, p\_0, L)\), then \(\operatorname{lhcenter}(t, p + Q) = \operatorname{lhcenter}(t, p)\).
\end{theorem}
\begin{proof}
The goal follows from \(\operatorname{lhcenterspatialperiodicofdoublesleandlhbandagree}(Q, z, t, p, p\_0, L)\).
\end{proof}
\begin{theorem}
\label{ll:paper:lhcenter-spatial-antiperiodic-of-doubles-and-lhbandagree}
For every natural number \(Q\) and natural number \(t\) and natural number \(p\) and natural number \(p\_0\) and natural number \(L\), if \(p \le t\), then if \(\operatorname{zerodiagonal}(t - p)\), then if \(\operatorname{doublesat}(Q, t - p)\), then if \(\operatorname{settledby}(t - p + 1, t + 1)\), then if \(p\_0 \le p + Q\), then if \(p + Q \le p\_0 + L\), then if \(\operatorname{lhbandagree}(t + 1, Q, p\_0, L)\), then \(\operatorname{lhcenter}(t + 1, p + Q) = \operatorname{not}(\operatorname{lhcenter}(t + 1, p))\).
\end{theorem}
\begin{proof}
The goal follows from \(\operatorname{lhcenterspatialantiperiodicofdoublesandlhbandagree}(Q, t, p, p\_0, L)\).
\end{proof}
\begin{theorem}
\label{ll:paper:winwitt-testbit-settled-left-ge}
For every natural number \(t\) and natural number \(k\) and natural number \(p\), if \(t - p \le k\), then if \(\operatorname{settledby}(t - p, t)\), then \(\operatorname{wittbitbool}(\operatorname{wittwin}(t, \operatorname{twopow}(k)), t + p + \operatorname{twopow}(k + 1)) = \operatorname{false}\).
\end{theorem}
\begin{proof}
The goal follows from \(\operatorname{winwitttestbitsettledleftge}(t, k, p)\).
\end{proof}
\begin{theorem}
\label{ll:paper:winwittagree-settled-left-strip}
For every natural number \(t\) and natural number \(k\) and natural number \(Q\) and natural number \(p\_0\) and natural number \(L\), if \(t + Q - p\_0 \le k\), then if for every natural number \(d\), if \(d \le t + Q - p\_0\), then \(\operatorname{settledby}(d, t)\), then \(\operatorname{winwittagree}(t, k, Q, t + p\_0 + \operatorname{twopow}(k + 1), L)\).
\end{theorem}
\begin{proof}
The goal follows from \(\operatorname{winwittagreesettledleftstrip}(t, k, Q, p\_0, L)\).
\end{proof}
\begin{theorem}
\label{ll:paper:lhbandagree-inward-of-settled-left-strip}
For every natural number \(t\) and natural number \(k\) and natural number \(Q\) and natural number \(p\_0\) and natural number \(L\), if \(t + Q - p\_0 \le k\), then if for every natural number \(d\), if \(d \le t + Q - p\_0\), then \(\operatorname{settledby}(d, t)\), then if \(\operatorname{lhbandagree}(t + \operatorname{twopow}(k), Q, p\_0 + \operatorname{twopow}(k), L)\), then \(\operatorname{lhbandagree}(t, Q, p\_0, L)\).
\end{theorem}
\begin{proof}
The goal follows from \(\operatorname{lhbandagreeinwardofsettledleftstrip}(t, k, Q, p\_0, L)\).
\end{proof}
\begin{theorem}
\label{ll:paper:lhbandagree-step-of-settled-left-strip}
For every natural number \(t\) and natural number \(k\) and natural number \(Q\) and natural number \(p\_0\) and natural number \(L\), if \(t + Q - p\_0 \le k\), then if for every natural number \(d\), if \(d \le t + Q - p\_0\), then \(\operatorname{settledby}(d, t)\), then if \(\operatorname{lhbandagree}(t, Q, p\_0, L)\), then \(\operatorname{lhbandagree}(t + \operatorname{twopow}(k), Q, p\_0 + \operatorname{twopow}(k), L)\).
\end{theorem}
\begin{proof}
The goal follows from \(\operatorname{lhbandagreestepofsettledleftstrip}(t, k, Q, p\_0, L)\).
\end{proof}
\begin{theorem}
\label{ll:paper:lhcenter-scale-step-settled-left-ge}
For every natural number \(t\) and natural number \(k\) and natural number \(p\), if \(t - p \le k\), then if \(\operatorname{settledby}(t - p, t)\), then \(\operatorname{lhcenter}(t + \operatorname{twopow}(k), p + \operatorname{twopow}(k)) = \operatorname{lhcenter}(t, p)\).
\end{theorem}
\begin{proof}
The goal follows from \(\operatorname{lhcenterscalestepsettledleftge}(t, k, p)\).
\end{proof}
\begin{theorem}
\label{ll:paper:ldiag-of-zero-diag}
For every natural number \(z\) and natural number \(t\), if \(\operatorname{zerodiagonal}(z)\), then if \(\operatorname{settledby}(z, t)\), then \(\operatorname{ldiag}(z, t) = \operatorname{false}\).
\end{theorem}
\begin{proof}
The goal follows from \(\operatorname{ldiagofzerodiag}(z, t)\).
\end{proof}
\begin{theorem}
\label{ll:paper:ldiag-after-zero-diag}
For every natural number \(z\) and natural number \(t\), if \(\operatorname{zerodiagonal}(z)\), then if \(\operatorname{settledby}(z + 2, t)\), then \(\operatorname{ldiag}(z + 2, t) = \operatorname{true}\).
\end{theorem}
\begin{proof}
The goal follows from \(\operatorname{ldiagafterzerodiag}(z, t)\).
\end{proof}
\begin{theorem}
\label{ll:paper:lhcenter-of-zero-diag}
For every natural number \(t\) and natural number \(p\), if \(\operatorname{zerodiagonal}(t - p)\), then if \(\operatorname{settledby}(t - p, t)\), then \(\operatorname{lhcenter}(t, p) = \operatorname{false}\).
\end{theorem}
\begin{proof}
The goal follows from \(\operatorname{lhcenterofzerodiag}(t, p)\).
\end{proof}
\begin{theorem}
\label{ll:paper:lhcenter-after-zero-diag}
For every natural number \(t\) and natural number \(p\), if \(\operatorname{zerodiagonal}(t - p)\), then if \(\operatorname{settledby}(t - p + 2, t + 2)\), then \(\operatorname{lhcenter}(t + 2, p) = \operatorname{true}\).
\end{theorem}
\begin{proof}
The goal follows from \(\operatorname{lhcenterafterzerodiag}(t, p)\).
\end{proof}
\begin{theorem}
\label{ll:paper:lhcenter-twin-before-zero-diag}
For every natural number \(t\) and natural number \(p\), if \(\operatorname{zerodiagonal}(t - p)\), then if \(\operatorname{settledby}(t - p, t)\), then \(\operatorname{lhcenter}(t, p + 2) = \operatorname{lhcenter}(t, p + 1)\).
\end{theorem}
\begin{proof}
The goal follows from \(\operatorname{lhcentertwinbeforezerodiag}(t, p)\).
\end{proof}
\begin{theorem}
\label{ll:paper:lhcenter-reset-after-zero-diag}
For every natural number \(t\) and natural number \(p\), if \(\operatorname{zerodiagonal}(t - (p + 1))\), then if \(\operatorname{settledby}(t - (p + 1) + 2, t + 2)\), then \(\operatorname{lhcenter}(t + 3, p + 1) = \operatorname{not}(\operatorname{lhcenter}(t + 2, p + 2))\).
\end{theorem}
\begin{proof}
The goal follows from \(\operatorname{lhcenterresetafterzerodiag}(t, p)\).
\end{proof}
\begin{theorem}
\label{ll:paper:lhcenter-two-step-same-col-flip-of-zero-diag}
For every natural number \(t\) and natural number \(p\), if \(\operatorname{zerodiagonal}(t - p)\), then if \(\operatorname{settledby}(t - p, t)\), then if \(\operatorname{settledby}(t - p + 2, t + 2)\), then \(\operatorname{lhcenter}(t + 2, p) = \operatorname{not}(\operatorname{lhcenter}(t, p))\).
\end{theorem}
\begin{proof}
The goal follows from \(\operatorname{lhcentertwostepsamecolflipofzerodiag}(t, p)\).
\end{proof}
\begin{theorem}
\label{ll:paper:not-lhbandagree-two-of-zero-diag}
For every natural number \(t\) and natural number \(p\) and natural number \(p\_0\) and natural number \(L\), if \(\operatorname{zerodiagonal}(t - p)\), then if \(\operatorname{settledby}(t - p, t)\), then if \(\operatorname{settledby}(t - p + 2, t + 2)\), then if \(p\_0 \le p\), then if \(p \le p\_0 + L\), then not \(\operatorname{lhbandagree}(t, 2, p\_0, L)\).
\end{theorem}
\begin{proof}
The goal follows from \(\operatorname{notlhbandagreetwoofzerodiag}(t, p, p\_0, L)\).
\end{proof}
\begin{theorem}
\label{ll:paper:not-lhbandagree-of-zero-and-after-zero-diag}
For every natural number \(Q\) and natural number \(t\) and natural number \(p\) and natural number \(z\) and natural number \(p\_0\) and natural number \(L\), if \(\operatorname{zerodiagonal}(t - p)\), then if \(\operatorname{zerodiagonal}(z)\), then if \(t + Q - p = z + 2\), then if \(\operatorname{settledby}(t - p, t)\), then if \(\operatorname{settledby}(z + 2, t + Q)\), then if \(p\_0 \le p\), then if \(p \le p\_0 + L\), then not \(\operatorname{lhbandagree}(t, Q, p\_0, L)\).
\end{theorem}
\begin{proof}
The goal follows from \(\operatorname{notlhbandagreeofzeroandafterzerodiag}(Q, t, p, z, p\_0, L)\).
\end{proof}
\begin{theorem}
\label{ll:paper:lhcenter-eq-false-at-shift-of-zero-diag-and-lhbandagree}
For every natural number \(Q\) and natural number \(t\) and natural number \(p\) and natural number \(p\_0\) and natural number \(L\), if \(\operatorname{zerodiagonal}(t + Q - p)\), then if \(\operatorname{settledby}(t + Q - p, t + Q)\), then if \(p\_0 \le p\), then if \(p \le p\_0 + L\), then if \(\operatorname{lhbandagree}(t, Q, p\_0, L)\), then \(\operatorname{lhcenter}(t, p) = \operatorname{false}\).
\end{theorem}
\begin{proof}
The goal follows from \(\operatorname{lhcentereqfalseatshiftofzerodiagandlhbandagree}(Q, t, p, p\_0, L)\).
\end{proof}
\begin{theorem}
\label{ll:paper:lhcenter-eq-true-at-shift-after-zero-diag-and-lhbandagree}
For every natural number \(Q\) and natural number \(t\) and natural number \(p\) and natural number \(p\_0\) and natural number \(L\), if \(\operatorname{zerodiagonal}(t + Q - (p + 2))\), then if \(\operatorname{settledby}(t + Q - (p + 2) + 2, t + Q)\), then if \(p\_0 \le p\), then if \(p \le p\_0 + L\), then if \(\operatorname{lhbandagree}(t, Q, p\_0, L)\), then \(\operatorname{lhcenter}(t, p) = \operatorname{true}\).
\end{theorem}
\begin{proof}
The goal follows from \(\operatorname{lhcentereqtrueatshiftafterzerodiagandlhbandagree}(Q, t, p, p\_0, L)\).
\end{proof}
\begin{theorem}
\label{ll:paper:lhcenter-step-flip-at-shift-after-zero-diag-of-single-lhbandagree}
For every natural number \(Q\) and natural number \(t\) and natural number \(p\) and natural number \(p\_0\) and natural number \(L\), if \(\operatorname{zerodiagonal}(t + Q - (p + 3))\), then if \(\operatorname{settledby}(t + Q - (p + 3) + 2, t + Q)\), then if \(p\_0 \le p + 1\), then if \(p + 1 \le p\_0 + L\), then if \(\operatorname{lhbandagree}(t, Q, p\_0, L)\), then \(\operatorname{lhcenter}(t + 1, p + 1) = \operatorname{not}(\operatorname{lhcenter}(t, p + 2))\).
\end{theorem}
\begin{proof}
The goal follows from \(\operatorname{lhcenterstepflipatshiftafterzerodiagofsinglelhbandagree}(Q, t, p, p\_0, L)\).
\end{proof}
\begin{theorem}
\label{ll:paper:lhcenter-step-one-eq-true-at-shift-of-zero-diag-and-single-lhbandagree}
For every natural number \(Q\) and natural number \(t\) and natural number \(p\) and natural number \(p\_0\) and natural number \(L\), if \(\operatorname{zerodiagonal}(t + Q - (p + 2))\), then if \(\operatorname{settledby}(t + Q - (p + 2), t + Q)\), then if \(\operatorname{settledby}(t + Q - (p + 2) + 2, t + Q)\), then if \(p\_0 \le p\), then if \(p + 2 \le p\_0 + L\), then if \(\operatorname{lhbandagree}(t, Q, p\_0, L)\), then \(\operatorname{lhcenter}(t + 1, p + 1) = \operatorname{true}\).
\end{theorem}
\begin{proof}
The goal follows from \(\operatorname{lhcentersteponeeqtrueatshiftofzerodiagandsinglelhbandagree}(Q, t, p, p\_0, L)\).
\end{proof}
\begin{theorem}
\label{ll:paper:periodic-lhbandagree-zero-of-rhcolagree-one}
For every natural number \(T\) and natural number \(P\), if \(0 < P\), then if \(\operatorname{centerperiodic}(T, P)\), then for every natural number \(m\) and natural number \(t\) and natural number \(S\), if \(T \le t\), then if \(\operatorname{rhcolagree}(t, m \cdot P, 1, S)\), then \(\operatorname{lhbandagree}(t, m \cdot P, 0, S + 1)\).
\end{theorem}
\begin{proof}
Assume \(hP\), \(hper\).
The goal follows from \(\operatorname{periodiclhbandagreezeroofrhcolagreeone}(T, P, hP, hper)\).
\end{proof}
\begin{theorem}
\label{ll:paper:periodic-lhbandagree-steps-of-rhcolagree-one}
For every natural number \(T\) and natural number \(P\), if \(0 < P\), then if \(\operatorname{centerperiodic}(T, P)\), then for every natural number \(m\) and natural number \(t\) and natural number \(L\) and natural number \(s\), if \(T \le t\), then if \(\operatorname{rhcolagree}(t, m \cdot P, 1, L + s)\), then \(\operatorname{lhbandagree}(t + s, m \cdot P, 0, L + 1)\).
\end{theorem}
\begin{proof}
Assume \(hP\), \(hper\).
The goal follows from \(\operatorname{periodiclhbandagreestepsofrhcolagreeone}(T, P, hP, hper)\).
\end{proof}
\begin{theorem}
\label{ll:paper:lhcenter-twin-at-shift-of-zero-diag-and-lhbandagree}
For every natural number \(Q\) and natural number \(t\) and natural number \(p\) and natural number \(p\_0\) and natural number \(L\), if \(\operatorname{zerodiagonal}(t + Q - p)\), then if \(\operatorname{settledby}(t + Q - p, t + Q)\), then if \(p\_0 \le p + 1\), then if \(p + 2 \le p\_0 + L\), then if \(\operatorname{lhbandagree}(t, Q, p\_0, L)\), then \(\operatorname{lhcenter}(t, p + 2) = \operatorname{lhcenter}(t, p + 1)\).
\end{theorem}
\begin{proof}
The goal follows from \(\operatorname{lhcentertwinatshiftofzerodiagandlhbandagree}(Q, t, p, p\_0, L)\).
\end{proof}
\begin{theorem}
\label{ll:paper:lhcenter-linear-step-one-at-shift-of-zero-diag-and-lhbandagree}
For every natural number \(Q\) and natural number \(t\) and natural number \(p\) and natural number \(p\_0\) and natural number \(L\), if \(\operatorname{zerodiagonal}(t + Q - p)\), then if \(\operatorname{settledby}(t + Q - p, t + Q)\), then if \(p\_0 \le p\), then if \(p \le p\_0 + L\), then if \(\operatorname{lhbandagree}(t, Q, p\_0, L)\), then \(\operatorname{lhcenter}(t + 1, p + 1) = \operatorname{bxor}(\operatorname{lhcenter}(t, p + 2), \operatorname{lhcenter}(t, p + 1))\).
\end{theorem}
\begin{proof}
The goal follows from \(\operatorname{lhcenterlinearsteponeatshiftofzerodiagandlhbandagree}(Q, t, p, p\_0, L)\).
\end{proof}
\begin{theorem}
\label{ll:paper:lhcenter-linear-step-two-at-shift-of-zero-diag-and-lhbandagree}
For every natural number \(Q\) and natural number \(t\) and natural number \(p\) and natural number \(p\_0\) and natural number \(L\), if \(\operatorname{zerodiagonal}(t + Q - (p + 1))\), then if \(\operatorname{settledby}(t + Q - (p + 1), t + Q)\), then if \(p\_0 \le p + 1\), then if \(p + 1 \le p\_0 + L\), then if \(\operatorname{lhbandagree}(t, Q, p\_0, L)\), then \(\operatorname{lhcenter}(t + 1, p + 1) = \operatorname{bxor}(\operatorname{lhcenter}(t, p + 2), \operatorname{lhcenter}(t, p))\).
\end{theorem}
\begin{proof}
The goal follows from \(\operatorname{lhcenterlinearsteptwoatshiftofzerodiagandlhbandagree}(Q, t, p, p\_0, L)\).
\end{proof}
\begin{theorem}
\label{ll:paper:lhcenter-step-one-eq-false-at-shift-of-zero-diag-and-single-lhbandagree}
For every natural number \(Q\) and natural number \(t\) and natural number \(p\) and natural number \(p\_0\) and natural number \(L\), if \(\operatorname{zerodiagonal}(t + Q - p)\), then if \(\operatorname{settledby}(t + Q - p, t + Q)\), then if \(p\_0 \le p\), then if \(p + 2 \le p\_0 + L\), then if \(\operatorname{lhbandagree}(t, Q, p\_0, L)\), then \(\operatorname{lhcenter}(t + 1, p + 1) = \operatorname{false}\).
\end{theorem}
\begin{proof}
The goal follows from \(\operatorname{lhcentersteponeeqfalseatshiftofzerodiagandsinglelhbandagree}(Q, t, p, p\_0, L)\).
\end{proof}
\begin{theorem}
\label{ll:paper:not-lhbandagree-of-zero-diag-two-and-shift-zero-diag}
For every natural number \(Q\) and natural number \(t\) and natural number \(p\) and natural number \(p\_0\) and natural number \(L\), if \(\operatorname{zerodiagonal}(t - (p + 2))\), then if \(\operatorname{zerodiagonal}(t + Q - p)\), then if \(\operatorname{settledby}(t - (p + 2) + 2, t + 1)\), then if \(\operatorname{settledby}(t + Q - p, t + Q)\), then if \(p\_0 \le p\), then if \(p + 2 \le p\_0 + L\), then not \(\operatorname{lhbandagree}(t, Q, p\_0, L)\).
\end{theorem}
\begin{proof}
The goal follows from \(\operatorname{notlhbandagreeofzerodiagtwoandshiftzerodiag}(Q, t, p, p\_0, L)\).
\end{proof}
\begin{theorem}
\label{ll:paper:lhcenter-step-absorb-false-after-zero-diag}
For every natural number \(t\) and natural number \(p\), if \(\operatorname{zerodiagonal}(t - (p + 2))\), then if \(\operatorname{settledby}(t - (p + 2) + 2, t + 2)\), then if \(\operatorname{lhcenter}(t + 2, p + 1) = \operatorname{true}\), then \(\operatorname{lhcenter}(t + 3, p + 1) = \operatorname{false}\).
\end{theorem}
\begin{proof}
The goal follows from \(\operatorname{lhcenterstepabsorbfalseafterzerodiag}(t, p)\).
\end{proof}
\begin{theorem}
\label{ll:paper:lhcenter-two-step-absorb-false-after-zero-diag}
For every natural number \(t\) and natural number \(p\), if \(\operatorname{zerodiagonal}(t - (p + 2))\), then if \(\operatorname{settledby}(t - (p + 2) + 2, t + 2)\), then if \(\operatorname{lhcenter}(t + 2, p + 3) = \operatorname{false}\), then \(\operatorname{lhcenter}(t + 4, p + 2) = \operatorname{false}\).
\end{theorem}
\begin{proof}
The goal follows from \(\operatorname{lhcentertwostepabsorbfalseafterzerodiag}(t, p)\).
\end{proof}
\begin{theorem}
\label{ll:paper:lhcenter-late-cross-of-zero-diag-uncond}
For every natural number \(t\) and natural number \(k\) and natural number \(r\), if \(r \le \operatorname{twopow}(k)\), then if \(\operatorname{zerodiagonal}(t + r)\), then if \(\operatorname{settledby}(t + r, t + \operatorname{twopow}(k))\), then \(\operatorname{lhcenter}(t + \operatorname{twopow}(k), \operatorname{twopow}(k) - r) = \operatorname{false}\).
\end{theorem}
\begin{proof}
The goal follows from \(\operatorname{lhcenterlatecrossofzerodiaguncond}(t, k, r)\).
\end{proof}
\begin{theorem}
\label{ll:paper:lhcenter-late-cross-after-zero-diag-uncond}
For every natural number \(t\) and natural number \(k\) and natural number \(r\), if \(r + 2 \le \operatorname{twopow}(k)\), then if \(\operatorname{zerodiagonal}(t + r)\), then if \(\operatorname{settledby}(t + r + 2, t + \operatorname{twopow}(k))\), then \(\operatorname{lhcenter}(t + \operatorname{twopow}(k), \operatorname{twopow}(k) - (r + 2)) = \operatorname{true}\).
\end{theorem}
\begin{proof}
The goal follows from \(\operatorname{lhcenterlatecrossafterzerodiaguncond}(t, k, r)\).
\end{proof}
\begin{theorem}
\label{ll:paper:lhcenter-late-cross-reset-after-zero-diag-uncond}
For every natural number \(t\) and natural number \(k\) and natural number \(r\), if \(r + 3 \le \operatorname{twopow}(k)\), then if \(\operatorname{zerodiagonal}(t + r)\), then if \(\operatorname{settledby}(t + r + 2, t + \operatorname{twopow}(k))\), then \(\operatorname{lhcenter}(t + 1 + \operatorname{twopow}(k), \operatorname{twopow}(k) - (r + 2)) = \operatorname{not}(\operatorname{lhcenter}(t + \operatorname{twopow}(k), \operatorname{twopow}(k) - (r + 1)))\).
\end{theorem}
\begin{proof}
The goal follows from \(\operatorname{lhcenterlatecrossresetafterzerodiaguncond}(t, k, r)\).
\end{proof}
\begin{theorem}
\label{ll:paper:lhcenter-late-cross-twin-before-zero-diag-uncond}
For every natural number \(t\) and natural number \(k\) and natural number \(r\), if \(r + 2 \le \operatorname{twopow}(k)\), then if \(\operatorname{zerodiagonal}(t + r + 2)\), then if \(\operatorname{settledby}(t + r + 2, t + \operatorname{twopow}(k))\), then \(\operatorname{lhcenter}(t + \operatorname{twopow}(k), \operatorname{twopow}(k) - r) = \operatorname{lhcenter}(t + \operatorname{twopow}(k), \operatorname{twopow}(k) - (r + 1))\).
\end{theorem}
\begin{proof}
The goal follows from \(\operatorname{lhcenterlatecrosstwinbeforezerodiaguncond}(t, k, r)\).
\end{proof}
\begin{theorem}
\label{ll:paper:lhcenter-late-cross-periodic-of-doubles-le-uncond}
For every natural number \(Q\) and natural number \(z\) and natural number \(t\) and natural number \(k\) and natural number \(r\), if \(r + Q \le \operatorname{twopow}(k)\), then if \(t + Q + r \le z\), then if \(\operatorname{doublesat}(Q, z)\), then if \(\operatorname{settledby}(t + Q + r, t + \operatorname{twopow}(k))\), then \(\operatorname{lhcenter}(t + Q + \operatorname{twopow}(k), \operatorname{twopow}(k) - r) = \operatorname{lhcenter}(t + \operatorname{twopow}(k), \operatorname{twopow}(k) - (r + Q))\).
\end{theorem}
\begin{proof}
The goal follows from \(\operatorname{lhcenterlatecrossperiodicofdoublesleuncond}(Q, z, t, k, r)\).
\end{proof}
\begin{theorem}
\label{ll:paper:lhcenter-late-cross-antiperiodic-of-doubles-uncond}
For every natural number \(Q\) and natural number \(t\) and natural number \(k\) and natural number \(r\), if \(r + Q + 1 \le \operatorname{twopow}(k)\), then if \(\operatorname{zerodiagonal}(t + Q + r)\), then if \(\operatorname{doublesat}(Q, t + Q + r)\), then if \(\operatorname{settledby}(t + Q + r + 1, t + \operatorname{twopow}(k))\), then \(\operatorname{lhcenter}(t + Q + \operatorname{twopow}(k), \operatorname{twopow}(k) - (r + 1)) = \operatorname{not}(\operatorname{lhcenter}(t + \operatorname{twopow}(k), \operatorname{twopow}(k) - (r + Q + 1)))\).
\end{theorem}
\begin{proof}
The goal follows from \(\operatorname{lhcenterlatecrossantiperiodicofdoublesuncond}(Q, t, k, r)\).
\end{proof}
\begin{theorem}
\label{ll:paper:winwitt-testbit-scale-cross-of-zero-diag-uncond}
For every natural number \(t\) and natural number \(k\) and natural number \(r\), if \(t + r < \operatorname{twopow}(k)\), then if \(\operatorname{zerodiagonal}(t + r)\), then if \(\operatorname{settledby}(t + r, t + \operatorname{twopow}(k))\), then \(\operatorname{wittbitbool}(\operatorname{wittwin}(t, \operatorname{twopow}(k)), t + \operatorname{twopow}(k + 1) - r) = \operatorname{rhcenter}(t, r)\).
\end{theorem}
\begin{proof}
The goal follows from \(\operatorname{winwitttestbitscalecrossofzerodiaguncond}(t, k, r)\).
\end{proof}
\begin{theorem}
\label{ll:paper:winwitt-testbit-scale-cross-zero-gap-of-zero-diag-uncond}
For every natural number \(t\) and natural number \(k\) and natural number \(r\), if \(t < r\), then if \(t + r < \operatorname{twopow}(k)\), then if \(\operatorname{zerodiagonal}(t + r)\), then if \(\operatorname{settledby}(t + r, t + \operatorname{twopow}(k))\), then \(\operatorname{wittbitbool}(\operatorname{wittwin}(t, \operatorname{twopow}(k)), t + \operatorname{twopow}(k + 1) - r) = \operatorname{false}\).
\end{theorem}
\begin{proof}
The goal follows from \(\operatorname{winwitttestbitscalecrosszerogapofzerodiaguncond}(t, k, r)\).
\end{proof}
\begin{theorem}
\label{ll:paper:winwitt-testbit-scale-cross-after-zero-diag-uncond}
For every natural number \(t\) and natural number \(k\) and natural number \(r\), if \(t + r + 2 < \operatorname{twopow}(k)\), then if \(\operatorname{zerodiagonal}(t + r)\), then if \(\operatorname{settledby}(t + r + 2, t + \operatorname{twopow}(k))\), then \(\operatorname{wittbitbool}(\operatorname{wittwin}(t, \operatorname{twopow}(k)), t + \operatorname{twopow}(k + 1) - (r + 2)) = \operatorname{not}(\operatorname{rhcenter}(t, r + 2))\).
\end{theorem}
\begin{proof}
The goal follows from \(\operatorname{winwitttestbitscalecrossafterzerodiaguncond}(t, k, r)\).
\end{proof}
\begin{theorem}
\label{ll:paper:winwitt-testbit-scale-cross-zero-gap-after-zero-diag-uncond}
For every natural number \(t\) and natural number \(k\) and natural number \(r\), if \(t < r + 2\), then if \(t + r + 2 < \operatorname{twopow}(k)\), then if \(\operatorname{zerodiagonal}(t + r)\), then if \(\operatorname{settledby}(t + r + 2, t + \operatorname{twopow}(k))\), then \(\operatorname{wittbitbool}(\operatorname{wittwin}(t, \operatorname{twopow}(k)), t + \operatorname{twopow}(k + 1) - (r + 2)) = \operatorname{true}\).
\end{theorem}
\begin{proof}
The goal follows from \(\operatorname{winwitttestbitscalecrosszerogapafterzerodiaguncond}(t, k, r)\).
\end{proof}
\begin{theorem}
\label{ll:paper:winwitt-testbit-scale-cross-zero-gap-twin-before-zero-diag-uncond}
For every natural number \(t\) and natural number \(k\) and natural number \(r\), if \(t < r\), then if \(t + r + 2 \le \operatorname{twopow}(k)\), then if \(\operatorname{zerodiagonal}(t + r + 2)\), then if \(\operatorname{settledby}(t + r + 2, t + \operatorname{twopow}(k))\), then \(\operatorname{wittbitbool}(\operatorname{wittwin}(t, \operatorname{twopow}(k)), t + \operatorname{twopow}(k + 1) - r) = \operatorname{wittbitbool}(\operatorname{wittwin}(t, \operatorname{twopow}(k)), t + \operatorname{twopow}(k + 1) - (r + 1))\).
\end{theorem}
\begin{proof}
The goal follows from \(\operatorname{winwitttestbitscalecrosszerogaptwinbeforezerodiaguncond}(t, k, r)\).
\end{proof}
\begin{theorem}
\label{ll:paper:winwitt-sector-two-step-flip-zero-gap-after-zero-diag-uncond}
For every natural number \(t\) and natural number \(k\) and natural number \(r\), if \(t \le r\), then if \(t + r + 3 < \operatorname{twopow}(k)\), then if \(\operatorname{zerodiagonal}(t + r)\), then if \(\operatorname{settledby}(t + r + 2, t + \operatorname{twopow}(k))\), then \(\operatorname{wittbitbool}(\operatorname{wittwin}(t + 1, \operatorname{twopow}(k)), t + \operatorname{twopow}(k + 1) - (r + 1)) = \operatorname{not}(\operatorname{wittbitbool}(\operatorname{wittwin}(t, \operatorname{twopow}(k)), t + \operatorname{twopow}(k + 1) - (r + 1)))\).
\end{theorem}
\begin{proof}
The goal follows from \(\operatorname{winwittsectortwostepflipzerogapafterzerodiaguncond}(t, k, r)\).
\end{proof}
\begin{theorem}
\label{ll:paper:winwitt-scale-cross-periodic-of-doubles-le-uncond}
For every natural number \(Q\) and natural number \(z\) and natural number \(t\) and natural number \(k\) and natural number \(r\), if \(t + Q + r < \operatorname{twopow}(k)\), then if \(t + Q + r \le z\), then if \(\operatorname{doublesat}(Q, z)\), then if \(\operatorname{settledby}(t + Q + r, t + \operatorname{twopow}(k))\), then \(\operatorname{bxor}(\operatorname{wittbitbool}(\operatorname{wittwin}(t, \operatorname{twopow}(k)), t + \operatorname{twopow}(k + 1) - (r + Q)), \operatorname{wittbitbool}(\operatorname{wittwin}(t + Q, \operatorname{twopow}(k)), t + Q + \operatorname{twopow}(k + 1) - r)) = \operatorname{bxor}(\operatorname{rhcenter}(t, r + Q), \operatorname{rhcenter}(t + Q, r))\).
\end{theorem}
\begin{proof}
The goal follows from \(\operatorname{winwittscalecrossperiodicofdoublesleuncond}(Q, z, t, k, r)\).
\end{proof}
\begin{theorem}
\label{ll:paper:winwitt-scale-cross-zero-gap-periodic-of-doubles-le-uncond}
For every natural number \(Q\) and natural number \(z\) and natural number \(t\) and natural number \(k\) and natural number \(r\), if \(t + Q < r\), then if \(t + Q + r < \operatorname{twopow}(k)\), then if \(t + Q + r \le z\), then if \(\operatorname{doublesat}(Q, z)\), then if \(\operatorname{settledby}(t + Q + r, t + \operatorname{twopow}(k))\), then \(\operatorname{wittbitbool}(\operatorname{wittwin}(t + Q, \operatorname{twopow}(k)), t + Q + \operatorname{twopow}(k + 1) - r) = \operatorname{wittbitbool}(\operatorname{wittwin}(t, \operatorname{twopow}(k)), t + \operatorname{twopow}(k + 1) - (r + Q))\).
\end{theorem}
\begin{proof}
The goal follows from \(\operatorname{winwittscalecrosszerogapperiodicofdoublesleuncond}(Q, z, t, k, r)\).
\end{proof}
\begin{theorem}
\label{ll:paper:winwitt-scale-cross-antiperiodic-of-doubles-uncond}
For every natural number \(Q\) and natural number \(t\) and natural number \(k\) and natural number \(r\), if \(t + Q + r + 1 < \operatorname{twopow}(k)\), then if \(\operatorname{zerodiagonal}(t + Q + r)\), then if \(\operatorname{doublesat}(Q, t + Q + r)\), then if \(\operatorname{settledby}(t + Q + r + 1, t + \operatorname{twopow}(k))\), then \(\operatorname{bxor}(\operatorname{wittbitbool}(\operatorname{wittwin}(t, \operatorname{twopow}(k)), t + \operatorname{twopow}(k + 1) - (r + Q + 1)), \operatorname{wittbitbool}(\operatorname{wittwin}(t + Q, \operatorname{twopow}(k)), t + Q + \operatorname{twopow}(k + 1) - (r + 1))) = \operatorname{not}(\operatorname{bxor}(\operatorname{rhcenter}(t, r + Q + 1), \operatorname{rhcenter}(t + Q, r + 1)))\).
\end{theorem}
\begin{proof}
The goal follows from \(\operatorname{winwittscalecrossantiperiodicofdoublesuncond}(Q, t, k, r)\).
\end{proof}
\begin{theorem}
\label{ll:paper:winwitt-scale-cross-zero-gap-antiperiodic-of-doubles-uncond}
For every natural number \(Q\) and natural number \(t\) and natural number \(k\) and natural number \(r\), if \(t + Q \le r\), then if \(t + Q + r + 1 < \operatorname{twopow}(k)\), then if \(\operatorname{zerodiagonal}(t + Q + r)\), then if \(\operatorname{doublesat}(Q, t + Q + r)\), then if \(\operatorname{settledby}(t + Q + r + 1, t + \operatorname{twopow}(k))\), then \(\operatorname{wittbitbool}(\operatorname{wittwin}(t + Q, \operatorname{twopow}(k)), t + Q + \operatorname{twopow}(k + 1) - (r + 1)) = \operatorname{not}(\operatorname{wittbitbool}(\operatorname{wittwin}(t, \operatorname{twopow}(k)), t + \operatorname{twopow}(k + 1) - (r + Q + 1)))\).
\end{theorem}
\begin{proof}
The goal follows from \(\operatorname{winwittscalecrosszerogapantiperiodicofdoublesuncond}(Q, t, k, r)\).
\end{proof}
\begin{theorem}
\label{ll:paper:lhcenter-at-t-after-zero-diag}
For every natural number \(t\) and natural number \(p\), if \(\operatorname{zerodiagonal}(t - (p + 2))\), then if \(\operatorname{settledby}(t - (p + 2) + 2, t)\), then \(\operatorname{lhcenter}(t, p) = \operatorname{true}\).
\end{theorem}
\begin{proof}
The goal follows from \(\operatorname{lhcenterattafterzerodiag}(t, p)\).
\end{proof}
\begin{theorem}
\label{ll:paper:lhcenter-step-one-reset-after-zero-diag}
For every natural number \(t\) and natural number \(p\), if \(\operatorname{zerodiagonal}(t - (p + 3))\), then if \(\operatorname{settledby}(t - (p + 3) + 2, t)\), then \(\operatorname{lhcenter}(t + 1, p + 1) = \operatorname{not}(\operatorname{lhcenter}(t, p + 2))\).
\end{theorem}
\begin{proof}
The goal follows from \(\operatorname{lhcentersteponeresetafterzerodiag}(t, p)\).
\end{proof}
\begin{theorem}
\label{ll:paper:lhcenter-diag-shift-antiperiodic-at-t-of-doubles}
For every natural number \(Q\) and natural number \(t\) and natural number \(p\), if \(\operatorname{zerodiagonal}(t - (p + 1))\), then if \(\operatorname{doublesat}(Q, t - (p + 1))\), then if \(\operatorname{settledby}(t - (p + 1) + 1, t)\), then \(\operatorname{lhcenter}(t + Q, p + Q) = \operatorname{not}(\operatorname{lhcenter}(t, p))\).
\end{theorem}
\begin{proof}
The goal follows from \(\operatorname{lhcenterdiagshiftantiperiodicattofdoubles}(Q, t, p)\).
\end{proof}
\begin{theorem}
\label{ll:paper:lhcenter-spatial-antiperiodic-at-t-of-doubles-and-lhbandagree}
For every natural number \(Q\) and natural number \(t\) and natural number \(p\) and natural number \(p\_0\) and natural number \(L\), if \(\operatorname{zerodiagonal}(t - (p + 1))\), then if \(\operatorname{doublesat}(Q, t - (p + 1))\), then if \(\operatorname{settledby}(t - (p + 1) + 1, t)\), then if \(p\_0 \le p + Q\), then if \(p + Q \le p\_0 + L\), then if \(\operatorname{lhbandagree}(t, Q, p\_0, L)\), then \(\operatorname{lhcenter}(t, p + Q) = \operatorname{not}(\operatorname{lhcenter}(t, p))\).
\end{theorem}
\begin{proof}
The goal follows from \(\operatorname{lhcenterspatialantiperiodicattofdoublesandlhbandagree}(Q, t, p, p\_0, L)\).
\end{proof}
\begin{theorem}
\label{ll:paper:lhcenter-spatial-doubles-eq-step-one-reset-of-lhbandagree}
For every natural number \(Q\) and natural number \(t\) and natural number \(p\) and natural number \(p\_0\) and natural number \(L\), if \(\operatorname{zerodiagonal}(t - (p + 3))\), then if \(\operatorname{doublesat}(Q, t - (p + 3))\), then if \(\operatorname{settledby}(t - (p + 3) + 1, t)\), then if \(\operatorname{settledby}(t - (p + 3) + 2, t)\), then if \(p\_0 \le p + 2 + Q\), then if \(p + 2 + Q \le p\_0 + L\), then if \(\operatorname{lhbandagree}(t, Q, p\_0, L)\), then \(\operatorname{lhcenter}(t, p + 2 + Q) = \operatorname{lhcenter}(t + 1, p + 1)\).
\end{theorem}
\begin{proof}
The goal follows from \(\operatorname{lhcenterspatialdoubleseqsteponeresetoflhbandagree}(Q, t, p, p\_0, L)\).
\end{proof}
\begin{theorem}
\label{ll:paper:lhcenter-doubles-eq-step-one-reset}
For every natural number \(Q\) and natural number \(t\) and natural number \(p\), if \(\operatorname{zerodiagonal}(t - (p + 3))\), then if \(\operatorname{doublesat}(Q, t - (p + 3))\), then if \(\operatorname{settledby}(t - (p + 3) + 1, t)\), then if \(\operatorname{settledby}(t - (p + 3) + 2, t)\), then \(\operatorname{lhcenter}(t + Q, p + 2 + Q) = \operatorname{lhcenter}(t + 1, p + 1)\).
\end{theorem}
\begin{proof}
The goal follows from \(\operatorname{lhcenterdoubleseqsteponereset}(Q, t, p)\).
\end{proof}
\begin{theorem}
\label{ll:paper:lhcenter-step-one-reset-antiperiodic-of-doubles}
For every natural number \(Q\) and natural number \(t\) and natural number \(p\), if \(\operatorname{zerodiagonal}(t - (p + 3))\), then if \(\operatorname{doublesat}(Q, t - (p + 3))\), then if \(\operatorname{settledby}(t - (p + 3) + 1, t)\), then if \(\operatorname{settledby}(t - (p + 3) + 2, t)\), then \(\operatorname{lhcenter}(t + Q + 1, p + Q + 1) = \operatorname{not}(\operatorname{lhcenter}(t + 1, p + 1))\).
\end{theorem}
\begin{proof}
The goal follows from \(\operatorname{lhcentersteponeresetantiperiodicofdoubles}(Q, t, p)\).
\end{proof}
\begin{theorem}
\label{ll:paper:lhcenter-step-one-reset-shift-flip-of-doubles-and-single-lhbandagree}
For every natural number \(Q\) and natural number \(t\) and natural number \(p\) and natural number \(p\_0\) and natural number \(L\), if \(\operatorname{zerodiagonal}(t - (p + 3))\), then if \(\operatorname{doublesat}(Q, t - (p + 3))\), then if \(\operatorname{settledby}(t - (p + 3) + 1, t)\), then if \(\operatorname{settledby}(t - (p + 3) + 2, t)\), then if \(p\_0 \le p + 1 + Q\), then if \(p + 2 + Q \le p\_0 + L\), then if \(\operatorname{lhbandagree}(t, Q, p\_0, L)\), then \(\operatorname{lhcenter}(t + 1, p + 1 + Q) = \operatorname{not}(\operatorname{lhcenter}(t + 1, p + 1))\).
\end{theorem}
\begin{proof}
The goal follows from \(\operatorname{lhcentersteponeresetshiftflipofdoublesandsinglelhbandagree}(Q, t, p, p\_0, L)\).
\end{proof}
\begin{theorem}
\label{ll:paper:lhcenter-pin-doubling-diag-true-of-doubles-one-and-lhbandagree}
For every natural number \(t\) and natural number \(p\) and natural number \(p\_0\) and natural number \(L\), if \(\operatorname{zerodiagonal}(t - (p + 2))\), then if \(\operatorname{doublesat}(1, t - (p + 2))\), then if \(\operatorname{settledby}(t - (p + 2), t)\), then if \(\operatorname{settledby}(t - (p + 2) + 1, t)\), then if \(p\_0 \le p + 2\), then if \(p + 2 \le p\_0 + L\), then if \(\operatorname{lhbandagree}(t, 1, p\_0, L)\), then \(\operatorname{lhcenter}(t, p + 1) = \operatorname{true}\).
\end{theorem}
\begin{proof}
The goal follows from \(\operatorname{lhcenterpindoublingdiagtrueofdoublesoneandlhbandagree}(t, p, p\_0, L)\).
\end{proof}
\begin{theorem}
\label{ll:paper:lhcenter-step-one-pin-reset-diag-false-of-doubles-one-and-lhbandagree}
For every natural number \(t\) and natural number \(p\) and natural number \(p\_0\) and natural number \(L\), if \(\operatorname{zerodiagonal}(t - (p + 3))\), then if \(\operatorname{doublesat}(1, t - (p + 3))\), then if \(\operatorname{settledby}(t - (p + 3), t)\), then if \(\operatorname{settledby}(t - (p + 3) + 1, t)\), then if \(\operatorname{settledby}(t - (p + 3) + 2, t)\), then if \(p\_0 \le p + 3\), then if \(p + 3 \le p\_0 + L\), then if \(\operatorname{lhbandagree}(t, 1, p\_0, L)\), then \(\operatorname{lhcenter}(t + 1, p + 1) = \operatorname{false}\).
\end{theorem}
\begin{proof}
The goal follows from \(\operatorname{lhcentersteponepinresetdiagfalseofdoublesoneandlhbandagree}(t, p, p\_0, L)\).
\end{proof}
\begin{theorem}
\label{ll:paper:not-lhbandagree-one-of-doubles}
For every natural number \(t\) and natural number \(p\) and natural number \(p\_0\) and natural number \(L\), if \(\operatorname{zerodiagonal}(t - (p + 3))\), then if \(\operatorname{doublesat}(1, t - (p + 3))\), then if \(\operatorname{settledby}(t - (p + 3), t)\), then if \(\operatorname{settledby}(t - (p + 3) + 1, t)\), then if \(\operatorname{settledby}(t - (p + 3) + 2, t)\), then if \(p\_0 \le p + 1\), then if \(p + 3 \le p\_0 + L\), then not \(\operatorname{lhbandagree}(t, 1, p\_0, L)\).
\end{theorem}
\begin{proof}
The goal follows from \(\operatorname{notlhbandagreeoneofdoubles}(t, p, p\_0, L)\).
\end{proof}
\begin{theorem}
\label{ll:paper:not-periodic-rhcolagree-one-of-doubles-one}
For every natural number \(T\), if \(\operatorname{centerperiodic}(T, 1)\), then for every natural number \(t\) and natural number \(p\) and natural number \(S\), if \(T \le t\), then if \(\operatorname{zerodiagonal}(t - (p + 3))\), then if \(\operatorname{doublesat}(1, t - (p + 3))\), then if \(\operatorname{settledby}(t - (p + 3), t)\), then if \(\operatorname{settledby}(t - (p + 3) + 1, t)\), then if \(\operatorname{settledby}(t - (p + 3) + 2, t)\), then if \(p + 2 \le S\), then not \(\operatorname{rhcolagree}(t, 1, 1, S)\).
\end{theorem}
\begin{proof}
Assume \(hper\).
The goal follows from \(\operatorname{notperiodicrhcolagreeoneofdoublesone}(T, hper)\).
\end{proof}
\begin{theorem}
\label{ll:paper:lhcenter-pin-plus-two-true-of-lhbandagree-and-zero-diag-three}
For every natural number \(Q\) and natural number \(t\) and natural number \(p\) and natural number \(p\_0\) and natural number \(L\), if \(\operatorname{zerodiagonal}(t - p)\), then if \(\operatorname{zerodiagonal}(t + Q - (p + 3))\), then if \(\operatorname{settledby}(t - p, t)\), then if \(\operatorname{settledby}(t + Q - (p + 3) + 2, t + Q)\), then if \(p\_0 \le p + 1\), then if \(p + 1 \le p\_0 + L\), then if \(\operatorname{lhbandagree}(t, Q, p\_0, L)\), then \(\operatorname{lhcenter}(t, p + 2) = \operatorname{true}\).
\end{theorem}
\begin{proof}
The goal follows from \(\operatorname{lhcenterpinplustwotrueoflhbandagreeandzerodiagthree}(Q, t, p, p\_0, L)\).
\end{proof}
\begin{theorem}
\label{ll:paper:lhcenter-step-one-pin-plus-one-false-of-lhbandagree-and-zero-diag-three}
For every natural number \(Q\) and natural number \(t\) and natural number \(p\) and natural number \(p\_0\) and natural number \(L\), if \(\operatorname{zerodiagonal}(t - p)\), then if \(\operatorname{zerodiagonal}(t + Q - (p + 3))\), then if \(\operatorname{settledby}(t - p, t)\), then if \(\operatorname{settledby}(t + Q - (p + 3) + 2, t + Q)\), then if \(p\_0 \le p + 1\), then if \(p + 1 \le p\_0 + L\), then if \(\operatorname{lhbandagree}(t, Q, p\_0, L)\), then \(\operatorname{lhcenter}(t + 1, p + 1) = \operatorname{false}\).
\end{theorem}
\begin{proof}
The goal follows from \(\operatorname{lhcentersteponepinplusonefalseoflhbandagreeandzerodiagthree}(Q, t, p, p\_0, L)\).
\end{proof}
\begin{theorem}
\label{ll:paper:lhcenter-pin-before-zero-diag-of-zero-diag-three-and-lhbandagree}
For every natural number \(t\) and natural number \(p\) and natural number \(p\_0\) and natural number \(L\), if \(\operatorname{zerodiagonal}(t - p)\), then if \(\operatorname{settledby}(t - p, t)\), then if \(\operatorname{settledby}(t - p + 2, t + 3)\), then if \(p\_0 \le p + 1\), then if \(p + 1 \le p\_0 + L\), then if \(\operatorname{lhbandagree}(t, 3, p\_0, L)\), then \(\operatorname{lhcenter}(t, p + 2) = \operatorname{true}\).
\end{theorem}
\begin{proof}
The goal follows from \(\operatorname{lhcenterpinbeforezerodiagofzerodiagthreeandlhbandagree}(t, p, p\_0, L)\).
\end{proof}
\begin{theorem}
\label{ll:paper:lhcenter-pin-doubling-diag-false-of-doubles-three-and-lhbandagree}
For every natural number \(t\) and natural number \(p\) and natural number \(p\_0\) and natural number \(L\), if \(\operatorname{zerodiagonal}(t - (p + 1))\), then if \(\operatorname{doublesat}(3, t - (p + 1))\), then if \(\operatorname{settledby}(t - (p + 1), t)\), then if \(\operatorname{settledby}(t - (p + 1) + 1, t)\), then if \(\operatorname{settledby}(t - (p + 1) + 2, t + 3)\), then if \(p\_0 \le p + 2\), then if \(p + 3 \le p\_0 + L\), then if \(\operatorname{lhbandagree}(t, 3, p\_0, L)\), then \(\operatorname{lhcenter}(t, p) = \operatorname{false}\).
\end{theorem}
\begin{proof}
The goal follows from \(\operatorname{lhcenterpindoublingdiagfalseofdoublesthreeandlhbandagree}(t, p, p\_0, L)\).
\end{proof}
\begin{theorem}
\label{ll:paper:lhcenter-step-one-pin-plus-one-true-of-doubles-three-and-lhbandagree}
For every natural number \(t\) and natural number \(p\) and natural number \(p\_0\) and natural number \(L\), if \(\operatorname{zerodiagonal}(t - (p + 1))\), then if \(\operatorname{doublesat}(3, t - (p + 1))\), then if \(\operatorname{settledby}(t - (p + 1), t)\), then if \(\operatorname{settledby}(t - (p + 1) + 1, t)\), then if \(\operatorname{settledby}(t - (p + 1) + 2, t + 3)\), then if \(p\_0 \le p + 2\), then if \(p + 3 \le p\_0 + L\), then if \(\operatorname{lhbandagree}(t, 3, p\_0, L)\), then \(\operatorname{lhcenter}(t + 1, p + 1) = \operatorname{true}\).
\end{theorem}
\begin{proof}
The goal follows from \(\operatorname{lhcentersteponepinplusonetrueofdoublesthreeandlhbandagree}(t, p, p\_0, L)\).
\end{proof}
\begin{theorem}
\label{ll:paper:lhcenter-step-two-pin-plus-two-false-of-doubles-three-and-lhbandagree}
For every natural number \(t\) and natural number \(p\) and natural number \(p\_0\) and natural number \(L\), if \(\operatorname{zerodiagonal}(t - (p + 1))\), then if \(\operatorname{doublesat}(3, t - (p + 1))\), then if \(\operatorname{settledby}(t - (p + 1), t)\), then if \(\operatorname{settledby}(t - (p + 1) + 1, t)\), then if \(\operatorname{settledby}(t - (p + 1) + 2, t + 3)\), then if \(p\_0 \le p + 2\), then if \(p + 4 \le p\_0 + L\), then if \(\operatorname{lhbandagree}(t, 3, p\_0, L)\), then \(\operatorname{lhcenter}(t + 2, p + 2) = \operatorname{false}\).
\end{theorem}
\begin{proof}
The goal follows from \(\operatorname{lhcentersteptwopinplustwofalseofdoublesthreeandlhbandagree}(t, p, p\_0, L)\).
\end{proof}
\begin{theorem}
\label{ll:paper:lhcenter-step-three-pin-plus-one-true-of-doubles-three-and-lhbandagree}
For every natural number \(t\) and natural number \(p\) and natural number \(p\_0\) and natural number \(L\), if \(\operatorname{zerodiagonal}(t - (p + 1))\), then if \(\operatorname{doublesat}(3, t - (p + 1))\), then if \(\operatorname{settledby}(t - (p + 1), t)\), then if \(\operatorname{settledby}(t - (p + 1) + 1, t)\), then if \(\operatorname{settledby}(t - (p + 1) + 2, t + 2)\), then if \(p\_0 \le p + 2\), then if \(p + 4 \le p\_0 + L\), then if \(\operatorname{lhbandagree}(t, 3, p\_0, L)\), then \(\operatorname{lhcenter}(t + 3, p + 1) = \operatorname{true}\).
\end{theorem}
\begin{proof}
The goal follows from \(\operatorname{lhcenterstepthreepinplusonetrueofdoublesthreeandlhbandagree}(t, p, p\_0, L)\).
\end{proof}
\begin{theorem}
\label{ll:paper:not-lhbandagree-three-of-doubles}
For every natural number \(t\) and natural number \(p\) and natural number \(p\_0\) and natural number \(L\), if \(\operatorname{zerodiagonal}(t - (p + 1))\), then if \(\operatorname{doublesat}(3, t - (p + 1))\), then if \(\operatorname{settledby}(t - (p + 1), t)\), then if \(\operatorname{settledby}(t - (p + 1) + 1, t)\), then if \(\operatorname{settledby}(t - (p + 1) + 2, t + 2)\), then if \(p\_0 \le p + 1\), then if \(p + 4 \le p\_0 + L\), then not \(\operatorname{lhbandagree}(t, 3, p\_0, L)\).
\end{theorem}
\begin{proof}
The goal follows from \(\operatorname{notlhbandagreethreeofdoubles}(t, p, p\_0, L)\).
\end{proof}
\begin{theorem}
\label{ll:paper:ldiag-settled-twopow-periodic}
For every natural number \(d\) and natural number \(t\), if \(\operatorname{settledby}(d, t)\), then \(\operatorname{ldiag}(d, t + \operatorname{twopow}(d)) = \operatorname{ldiag}(d, t)\).
\end{theorem}
\begin{proof}
The goal follows from \(\operatorname{ldiagsettledtwopowperiodic}(d, t)\).
\end{proof}
\begin{theorem}
\label{ll:paper:ldiag-settled-twopow-periodic-le}
For every natural number \(e\) and natural number \(d\) and natural number \(t\), if \(d \le e\), then if \(\operatorname{settledby}(d, t)\), then \(\operatorname{ldiag}(d, t + \operatorname{twopow}(e)) = \operatorname{ldiag}(d, t)\).
\end{theorem}
\begin{proof}
The goal follows from \(\operatorname{ldiagsettledtwopowperiodicle}(e, d, t)\).
\end{proof}
\begin{theorem}
\label{ll:paper:ldiag-settled-two-mul-periodic-of-doubles}
For every natural number \(Q\) and natural number \(z\) and natural number \(t\), if \(\operatorname{zerodiagonal}(z)\), then if \(\operatorname{doublesat}(Q, z)\), then if \(\operatorname{settledby}(z + 1, t)\), then \(\operatorname{ldiag}(z + 1, t + (Q + Q)) = \operatorname{ldiag}(z + 1, t)\).
\end{theorem}
\begin{proof}
The goal follows from \(\operatorname{ldiagsettledtwomulperiodicofdoubles}(Q, z, t)\).
\end{proof}
\begin{theorem}
\label{ll:paper:ldiag-settled-reset-antiperiodic-of-doubles}
For every natural number \(Q\) and natural number \(z\) and natural number \(t\), if \(\operatorname{zerodiagonal}(z)\), then if \(\operatorname{doublesat}(Q, z)\), then if \(\operatorname{settledby}(z + 3, t)\), then \(\operatorname{ldiag}(z + 3, t + Q) = \operatorname{not}(\operatorname{ldiag}(z + 3, t))\).
\end{theorem}
\begin{proof}
The goal follows from \(\operatorname{ldiagsettledresetantiperiodicofdoubles}(Q, z, t)\).
\end{proof}
\begin{theorem}
\label{ll:paper:ldiag-settled-two-mul-periodic-le-add-three-of-doubles}
For every natural number \(Q\) and natural number \(z\) and natural number \(d\) and natural number \(t\), if \(\operatorname{zerodiagonal}(z)\), then if \(\operatorname{doublesat}(Q, z)\), then if \(d \le z + 3\), then if \(\operatorname{settledby}(d, t)\), then \(\operatorname{ldiag}(d, t + (Q + Q)) = \operatorname{ldiag}(d, t)\).
\end{theorem}
\begin{proof}
The goal follows from \(\operatorname{ldiagsettledtwomulperiodicleaddthreeofdoubles}(Q, z, d, t)\).
\end{proof}
\begin{theorem}
\label{ll:paper:exists-zero-diag-doublesat-twopow}
For every natural number \(e\), there exists a natural number \(z\) such that \(\operatorname{zerodiagonal}(z)\) and \(\operatorname{doublesat}(\operatorname{twopow}(e), z)\).
\end{theorem}
\begin{proof}
The goal follows from \(\operatorname{existszerodiagdoublesattwopow}(e)\).
\end{proof}
\begin{theorem}
\label{ll:paper:exists-zero-diag-doublesat-mul}
For every natural number \(m\) and natural number \(P\), if \(0 < m\), then if \(0 < P\), then there exists a natural number \(z\) such that \(\operatorname{zerodiagonal}(z)\) and \(\operatorname{doublesat}(m \cdot P, z)\).
\end{theorem}
\begin{proof}
The goal follows from \(\operatorname{existszerodiagdoublesatmul}(m, P)\).
\end{proof}
\begin{theorem}
\label{ll:paper:doublesat-twopow-le}
For every natural number \(e\) and natural number \(z\), if \(\operatorname{doublesat}(\operatorname{twopow}(e), z)\), then \(e \le z\).
\end{theorem}
\begin{proof}
The goal follows from \(\operatorname{doublesattwopowle}(e, z)\).
\end{proof}
\begin{theorem}
\label{ll:paper:exists-zero-diag-doublesat-twopow-ge}
For every natural number \(e\), there exists a natural number \(z\) such that \(e \le z\) and \(\operatorname{zerodiagonal}(z) \land \operatorname{doublesat}(\operatorname{twopow}(e), z)\).
\end{theorem}
\begin{proof}
The goal follows from \(\operatorname{existszerodiagdoublesattwopowge}(e)\).
\end{proof}
\begin{theorem}
\label{ll:paper:lhcenter-two-shift-diag-periodic-of-doubles}
For every natural number \(Q\) and natural number \(t\) and natural number \(p\), if \(\operatorname{zerodiagonal}(t - (p + 1))\), then if \(\operatorname{doublesat}(Q, t - (p + 1))\), then if \(\operatorname{settledby}(t - (p + 1) + 1, t)\), then \(\operatorname{lhcenter}(t + Q + Q, p + Q + Q) = \operatorname{lhcenter}(t, p)\).
\end{theorem}
\begin{proof}
The goal follows from \(\operatorname{lhcentertwoshiftdiagperiodicofdoubles}(Q, t, p)\).
\end{proof}
\begin{theorem}
\label{ll:paper:lhcenter-shift-two-q-flip-of-doubles-and-lhbandagree}
For every natural number \(Q\) and natural number \(t\) and natural number \(p\) and natural number \(p\_0\) and natural number \(L\), if \(1 \le Q\), then if \(\operatorname{zerodiagonal}(t - (p + 1))\), then if \(\operatorname{doublesat}(Q, t - (p + 1))\), then if \(\operatorname{settledby}(t - (p + 1) + 1, t)\), then if \(\operatorname{settledby}(t - (p + Q), t)\), then if \(p\_0 \le p + Q\), then if \(p + Q \le p\_0 + L\), then if \(\operatorname{lhbandagree}(t, Q, p\_0, L)\), then \(\operatorname{lhcenter}(t + Q, p + Q + Q) = \operatorname{not}(\operatorname{lhcenter}(t, p))\).
\end{theorem}
\begin{proof}
The goal follows from \(\operatorname{lhcentershifttwoqflipofdoublesandlhbandagree}(Q, t, p, p\_0, L)\).
\end{proof}
\begin{theorem}
\label{ll:paper:lhcenter-same-col-two-q-flip-of-doubles-and-lhbandagree}
For every natural number \(Q\) and natural number \(t\) and natural number \(p\) and natural number \(p\_0\) and natural number \(L\), if \(1 \le Q\), then if \(\operatorname{zerodiagonal}(t - (p + 1))\), then if \(\operatorname{doublesat}(Q, t - (p + 1))\), then if \(\operatorname{settledby}(t - (p + 1) + 1, t)\), then if \(\operatorname{settledby}(t - (p + Q), t)\), then if \(p\_0 \le p + Q\), then if \(p + Q \le p\_0 + L\), then if \(\operatorname{lhbandagree}(t, Q, p\_0, L)\), then \(\operatorname{lhcenter}(t + Q + Q, p + Q + Q) = \operatorname{not}(\operatorname{lhcenter}(t + Q, p + Q + Q))\).
\end{theorem}
\begin{proof}
The goal follows from \(\operatorname{lhcentersamecoltwoqflipofdoublesandlhbandagree}(Q, t, p, p\_0, L)\).
\end{proof}
\begin{theorem}
\label{ll:paper:not-lhbandagree-step-q-of-doubles}
For every natural number \(Q\) and natural number \(t\) and natural number \(p\) and natural number \(p\_0\) and natural number \(L\) and natural number \(p\_1\) and natural number \(L\_1\), if \(1 \le Q\), then if \(\operatorname{zerodiagonal}(t - (p + 1))\), then if \(\operatorname{doublesat}(Q, t - (p + 1))\), then if \(\operatorname{settledby}(t - (p + 1) + 1, t)\), then if \(\operatorname{settledby}(t - (p + Q), t)\), then if \(p\_0 \le p + Q\), then if \(p + Q \le p\_0 + L\), then if \(p\_1 \le p + Q + Q\), then if \(p + Q + Q \le p\_1 + L\_1\), then if \(\operatorname{lhbandagree}(t, Q, p\_0, L)\), then not \(\operatorname{lhbandagree}(t + Q, Q, p\_1, L\_1)\).
\end{theorem}
\begin{proof}
The goal follows from \(\operatorname{notlhbandagreestepqofdoubles}(Q, t, p, p\_0, L, p\_1, L\_1)\).
\end{proof}
\begin{theorem}
\label{ll:paper:not-periodic-rhcolagree-one-of-doubles}
For every natural number \(T\) and natural number \(P\), if \(0 < P\), then if \(\operatorname{centerperiodic}(T, P)\), then for every natural number \(m\) and natural number \(t\) and natural number \(p\) and natural number \(S\), if \(0 < m\), then if \(T \le t\), then if \(\operatorname{zerodiagonal}(t - (p + 1))\), then if \(\operatorname{doublesat}(m \cdot P, t - (p + 1))\), then if \(\operatorname{settledby}(t - (p + 1) + 1, t)\), then if \(\operatorname{settledby}(t - (p + m \cdot P), t)\), then if \(p + 3 \cdot (m \cdot P) \le S + 1\), then not \(\operatorname{rhcolagree}(t, m \cdot P, 1, S)\).
\end{theorem}
\begin{proof}
Assume \(hP\), \(hper\).
The goal follows from \(\operatorname{notperiodicrhcolagreeoneofdoubles}(T, P, hP, hper)\).
\end{proof}
\begin{theorem}
\label{ll:paper:not-centerperiodic-of-rhcolagree-one-doubles}
For every natural number \(T\) and natural number \(P\) and natural number \(m\) and natural number \(t\) and natural number \(p\) and natural number \(S\), if \(0 < P\), then if \(0 < m\), then if \(T \le t\), then if \(\operatorname{zerodiagonal}(t - (p + 1))\), then if \(\operatorname{doublesat}(m \cdot P, t - (p + 1))\), then if \(\operatorname{settledby}(t - (p + 1) + 1, t)\), then if \(\operatorname{settledby}(t - (p + m \cdot P), t)\), then if \(p + 3 \cdot (m \cdot P) \le S + 1\), then if \(\operatorname{rhcolagree}(t, m \cdot P, 1, S)\), then not \(\operatorname{centerperiodic}(T, P)\).
\end{theorem}
\begin{proof}
The goal follows from \(\operatorname{notcenterperiodicofrhcolagreeonedoubles}(T, P, m, t, p, S)\).
\end{proof}
\begin{theorem}
\label{ll:paper:not-periodic-lhcolagree-one-of-doubles}
For every natural number \(T\) and natural number \(P\), if \(0 < P\), then if \(\operatorname{centerperiodic}(T, P)\), then for every natural number \(m\) and natural number \(t\) and natural number \(p\) and natural number \(S\), if \(0 < m\), then if \(T \le t\), then if \(\operatorname{zerodiagonal}(t - (p + 1))\), then if \(\operatorname{doublesat}(m \cdot P, t - (p + 1))\), then if \(\operatorname{settledby}(t - (p + 1) + 1, t)\), then if \(\operatorname{settledby}(t - (p + m \cdot P), t)\), then if \(p + 3 \cdot (m \cdot P) \le S + 1\), then not \(\operatorname{lhcolagree}(t, m \cdot P, 1, S)\).
\end{theorem}
\begin{proof}
Assume \(hP\), \(hper\).
The goal follows from \(\operatorname{notperiodiclhcolagreeoneofdoubles}(T, P, hP, hper)\).
\end{proof}
\begin{theorem}
\label{ll:paper:not-centerperiodic-of-lhcolagree-one-doubles}
For every natural number \(T\) and natural number \(P\) and natural number \(m\) and natural number \(t\) and natural number \(p\) and natural number \(S\), if \(0 < P\), then if \(0 < m\), then if \(T \le t\), then if \(\operatorname{zerodiagonal}(t - (p + 1))\), then if \(\operatorname{doublesat}(m \cdot P, t - (p + 1))\), then if \(\operatorname{settledby}(t - (p + 1) + 1, t)\), then if \(\operatorname{settledby}(t - (p + m \cdot P), t)\), then if \(p + 3 \cdot (m \cdot P) \le S + 1\), then if \(\operatorname{lhcolagree}(t, m \cdot P, 1, S)\), then not \(\operatorname{centerperiodic}(T, P)\).
\end{theorem}
\begin{proof}
The goal follows from \(\operatorname{notcenterperiodicoflhcolagreeonedoubles}(T, P, m, t, p, S)\).
\end{proof}
\begin{theorem}
\label{ll:paper:not-centerperiodic-of-rhbandagree-zero-doubles}
For every natural number \(T\) and natural number \(P\) and natural number \(m\) and natural number \(t\) and natural number \(p\) and natural number \(S\), if \(0 < P\), then if \(0 < m\), then if \(T \le t\), then if \(\operatorname{zerodiagonal}(t - (p + 1))\), then if \(\operatorname{doublesat}(m \cdot P, t - (p + 1))\), then if \(\operatorname{settledby}(t - (p + 1) + 1, t)\), then if \(\operatorname{settledby}(t - (p + m \cdot P), t)\), then if \(p + 3 \cdot (m \cdot P) \le S + 1\), then if \(\operatorname{rhbandagree}(t, m \cdot P, 0, S + 1)\), then not \(\operatorname{centerperiodic}(T, P)\).
\end{theorem}
\begin{proof}
The goal follows from \(\operatorname{notcenterperiodicofrhbandagreezerodoubles}(T, P, m, t, p, S)\).
\end{proof}
\begin{theorem}
\label{ll:paper:not-centerperiodic-of-rhcenter-inward-determine-doubles}
For every natural number \(T\) and natural number \(P\) and natural number \(m\) and natural number \(t\) and natural number \(p\) and natural number \(S\) and natural number \(K\), if \(0 < P\), then if \(0 < m\), then if \(T \le t\), then if \(\operatorname{zerodiagonal}(t - (p + 1))\), then if \(\operatorname{doublesat}(m \cdot P, t - (p + 1))\), then if \(\operatorname{settledby}(t - (p + 1) + 1, t)\), then if \(\operatorname{settledby}(t - (p + m \cdot P), t)\), then if \(p + 3 \cdot (m \cdot P) \le S + 1\), then if \(\operatorname{rhcolagree}(t, m \cdot P, 1 + K, S + K)\), then if \(\operatorname{rhoragree}(t, m \cdot P, 1 + K, S + K)\), then not \(\operatorname{centerperiodic}(T, P)\).
\end{theorem}
\begin{proof}
The goal follows from \(\operatorname{notcenterperiodicofrhcenterinwarddeterminedoubles}(T, P, m, t, p, S, K)\).
\end{proof}
\begin{theorem}
\label{ll:paper:not-centerperiodic-of-rhoragree-one-doubles}
For every natural number \(T\) and natural number \(P\) and natural number \(m\) and natural number \(t\) and natural number \(p\) and natural number \(S\), if \(0 < P\), then if \(0 < m\), then if \(T \le t\), then if \(\operatorname{zerodiagonal}(t - (p + 1))\), then if \(\operatorname{doublesat}(m \cdot P, t - (p + 1))\), then if \(\operatorname{settledby}(t - (p + 1) + 1, t)\), then if \(\operatorname{settledby}(t - (p + m \cdot P), t)\), then if \(p + 3 \cdot (m \cdot P) \le S + 1\), then if \(\operatorname{rhcenter}(t + m \cdot P, 1) = \operatorname{rhcenter}(t, 1)\), then if \(\operatorname{rhoragree}(t, m \cdot P, 1, S)\), then not \(\operatorname{centerperiodic}(T, P)\).
\end{theorem}
\begin{proof}
The goal follows from \(\operatorname{notcenterperiodicofrhoragreeonedoubles}(T, P, m, t, p, S)\).
\end{proof}
\begin{theorem}
\label{ll:paper:not-periodic-rhcolagree-one-at-zero-diag-doubles}
For every natural number \(T\) and natural number \(P\), if \(0 < P\), then if \(\operatorname{centerperiodic}(T, P)\), then for every natural number \(m\) and natural number \(z\) and natural number \(p\) and natural number \(S\), if \(0 < m\), then if \(T \le z + (p + 1)\), then if \(\operatorname{zerodiagonal}(z)\), then if \(\operatorname{doublesat}(m \cdot P, z)\), then if \(\operatorname{settledby}(z + 1, z + (p + 1))\), then if \(\operatorname{settledby}(z + 1 - m \cdot P, z + (p + 1))\), then if \(p + 3 \cdot (m \cdot P) \le S + 1\), then not \(\operatorname{rhcolagree}(z + (p + 1), m \cdot P, 1, S)\).
\end{theorem}
\begin{proof}
Assume \(hP\), \(hper\).
The goal follows from \(\operatorname{notperiodicrhcolagreeoneatzerodiagdoubles}(T, P, hP, hper)\).
\end{proof}
\begin{theorem}
\label{ll:paper:exists-zero-diag-not-periodic-rhcolagree-one}
For every natural number \(T\) and natural number \(P\), if \(0 < P\), then if \(\operatorname{centerperiodic}(T, P)\), then for every natural number \(m\), if \(0 < m\), then there exists a natural number \(z\) such that there exists a natural number \(p\) such that \(\operatorname{zerodiagonal}(z)\) and \(\operatorname{doublesat}(m \cdot P, z) \land \lnot\operatorname{rhcolagree}(z + (p + 1), m \cdot P, 1, p + 3 \cdot (m \cdot P))\).
\end{theorem}
\begin{proof}
Assume \(hP\), \(hper\).
The goal follows from \(\operatorname{existszerodiagnotperiodicrhcolagreeone}(T, P, hP, hper)\).
\end{proof}
\begin{theorem}
\label{ll:paper:exists-zero-diag-not-periodic-lhcolagree-one}
For every natural number \(T\) and natural number \(P\), if \(0 < P\), then if \(\operatorname{centerperiodic}(T, P)\), then for every natural number \(m\), if \(0 < m\), then there exists a natural number \(z\) such that there exists a natural number \(p\) such that \(\operatorname{zerodiagonal}(z)\) and \(\operatorname{doublesat}(m \cdot P, z) \land \lnot\operatorname{lhcolagree}(z + (p + 1), m \cdot P, 1, p + 3 \cdot (m \cdot P))\).
\end{theorem}
\begin{proof}
Assume \(hP\), \(hper\).
The goal follows from \(\operatorname{existszerodiagnotperiodiclhcolagreeone}(T, P, hP, hper)\).
\end{proof}
\begin{theorem}
\label{ll:paper:maskparity-zero-eq-not-center-and-rhcenter-one}
For every natural number \(t\), \(\operatorname{maskparity}(t, 0) = \operatorname{not}(\operatorname{center}(t)) \operatorname{and} \operatorname{rhcenter}(t, 1)\).
\end{theorem}
\begin{proof}
The goal follows from \(\operatorname{maskparityzeroeqnotcenterandrhcenterone}(t)\).
\end{proof}
\begin{theorem}
\label{ll:paper:periodic-lhcenter-one-eq-of-black-center}
For every natural number \(T\) and natural number \(P\), if \(0 < P\), then if \(\operatorname{centerperiodic}(T, P)\), then for every natural number \(m\) and natural number \(t\), if \(T \le t\), then if \(\operatorname{center}(t) = \operatorname{true}\), then \(\operatorname{lhcenter}(t + m \cdot P, 1) = \operatorname{lhcenter}(t, 1)\).
\end{theorem}
\begin{proof}
Assume \(hP\), \(hper\).
The goal follows from \(\operatorname{periodiclhcenteroneeqofblackcenter}(T, P, hP, hper)\).
\end{proof}
\begin{theorem}
\label{ll:paper:periodic-lhcenter-one-eq-of-white-center}
For every natural number \(T\) and natural number \(P\), if \(0 < P\), then if \(\operatorname{centerperiodic}(T, P)\), then for every natural number \(m\) and natural number \(t\), if \(T \le t\), then if if \(\operatorname{center}(t) = \operatorname{false}\), then \(\operatorname{rhcenter}(t + m \cdot P, 1) = \operatorname{rhcenter}(t, 1)\), then \(\operatorname{lhcenter}(t + m \cdot P, 1) = \operatorname{lhcenter}(t, 1)\).
\end{theorem}
\begin{proof}
Assume \(hP\), \(hper\).
The goal follows from \(\operatorname{periodiclhcenteroneeqofwhitecenter}(T, P, hP, hper)\).
\end{proof}
\begin{theorem}
\label{ll:paper:periodic-lhcenter-one-eq-of-maskparity-zero-eq}
For every natural number \(T\) and natural number \(P\), if \(0 < P\), then if \(\operatorname{centerperiodic}(T, P)\), then for every natural number \(m\) and natural number \(t\), if \(T \le t\), then if \(\operatorname{maskparity}(t + m \cdot P, 0) = \operatorname{maskparity}(t, 0)\), then \(\operatorname{lhcenter}(t + m \cdot P, 1) = \operatorname{lhcenter}(t, 1)\).
\end{theorem}
\begin{proof}
Assume \(hP\), \(hper\).
The goal follows from \(\operatorname{periodiclhcenteroneeqofmaskparityzeroeq}(T, P, hP, hper)\).
\end{proof}
\begin{theorem}
\label{ll:paper:periodic-lhbandagree-zero-of-white-center-rhcenter-one}
For every natural number \(T\) and natural number \(P\), if \(0 < P\), then if \(\operatorname{centerperiodic}(T, P)\), then for every natural number \(m\) and natural number \(t\) and natural number \(S\), if \(T \le t\), then if for every natural number \(s\), if \(s \le S\), then if \(\operatorname{center}(t + s) = \operatorname{false}\), then \(\operatorname{rhcenter}(t + s + m \cdot P, 1) = \operatorname{rhcenter}(t + s, 1)\), then \(\operatorname{lhbandagree}(t, m \cdot P, 0, S + 1)\).
\end{theorem}
\begin{proof}
Assume \(hP\), \(hper\).
The goal follows from \(\operatorname{periodiclhbandagreezeroofwhitecenterrhcenterone}(T, P, hP, hper)\).
\end{proof}
\begin{theorem}
\label{ll:paper:periodic-lhbandagree-zero-of-maskparity-zero-agree}
For every natural number \(T\) and natural number \(P\), if \(0 < P\), then if \(\operatorname{centerperiodic}(T, P)\), then for every natural number \(m\) and natural number \(t\) and natural number \(S\), if \(T \le t\), then if for every natural number \(s\), if \(s \le S\), then \(\operatorname{maskparity}(t + s + m \cdot P, 0) = \operatorname{maskparity}(t + s, 0)\), then \(\operatorname{lhbandagree}(t, m \cdot P, 0, S + 1)\).
\end{theorem}
\begin{proof}
Assume \(hP\), \(hper\).
The goal follows from \(\operatorname{periodiclhbandagreezeroofmaskparityzeroagree}(T, P, hP, hper)\).
\end{proof}
\begin{theorem}
\label{ll:paper:not-centerperiodic-of-white-center-rhcenter-one-doubles}
For every natural number \(T\) and natural number \(P\) and natural number \(m\) and natural number \(t\) and natural number \(p\) and natural number \(S\), if \(0 < P\), then if \(0 < m\), then if \(T \le t\), then if \(\operatorname{zerodiagonal}(t - (p + 1))\), then if \(\operatorname{doublesat}(m \cdot P, t - (p + 1))\), then if \(\operatorname{settledby}(t - (p + 1) + 1, t)\), then if \(\operatorname{settledby}(t - (p + m \cdot P), t)\), then if \(p + 3 \cdot (m \cdot P) \le S + 1\), then if for every natural number \(s\), if \(s \le S\), then if \(\operatorname{center}(t + s) = \operatorname{false}\), then \(\operatorname{rhcenter}(t + s + m \cdot P, 1) = \operatorname{rhcenter}(t + s, 1)\), then not \(\operatorname{centerperiodic}(T, P)\).
\end{theorem}
\begin{proof}
The goal follows from \(\operatorname{notcenterperiodicofwhitecenterrhcenteronedoubles}(T, P, m, t, p, S)\).
\end{proof}
\begin{theorem}
\label{ll:paper:not-centerperiodic-of-maskparity-zero-agree-doubles}
For every natural number \(T\) and natural number \(P\) and natural number \(m\) and natural number \(t\) and natural number \(p\) and natural number \(S\), if \(0 < P\), then if \(0 < m\), then if \(T \le t\), then if \(\operatorname{zerodiagonal}(t - (p + 1))\), then if \(\operatorname{doublesat}(m \cdot P, t - (p + 1))\), then if \(\operatorname{settledby}(t - (p + 1) + 1, t)\), then if \(\operatorname{settledby}(t - (p + m \cdot P), t)\), then if \(p + 3 \cdot (m \cdot P) \le S + 1\), then if for every natural number \(s\), if \(s \le S\), then \(\operatorname{maskparity}(t + s + m \cdot P, 0) = \operatorname{maskparity}(t + s, 0)\), then not \(\operatorname{centerperiodic}(T, P)\).
\end{theorem}
\begin{proof}
The goal follows from \(\operatorname{notcenterperiodicofmaskparityzeroagreedoubles}(T, P, m, t, p, S)\).
\end{proof}
\begin{theorem}
\label{ll:paper:not-centerperiodic-of-white-center-rhcenter-one-at-zero-diag-doubles}
For every natural number \(T\) and natural number \(P\) and natural number \(m\) and natural number \(z\) and natural number \(p\) and natural number \(S\), if \(0 < P\), then if \(0 < m\), then if \(T \le z + (p + 1)\), then if \(\operatorname{zerodiagonal}(z)\), then if \(\operatorname{doublesat}(m \cdot P, z)\), then if \(\operatorname{settledby}(z + 1, z + (p + 1))\), then if \(\operatorname{settledby}(z + 1 - m \cdot P, z + (p + 1))\), then if \(p + 3 \cdot (m \cdot P) \le S + 1\), then if for every natural number \(s\), if \(s \le S\), then if \(\operatorname{center}(z + (p + 1) + s) = \operatorname{false}\), then \(\operatorname{rhcenter}(z + (p + 1) + s + m \cdot P, 1) = \operatorname{rhcenter}(z + (p + 1) + s, 1)\), then not \(\operatorname{centerperiodic}(T, P)\).
\end{theorem}
\begin{proof}
The goal follows from \(\operatorname{notcenterperiodicofwhitecenterrhcenteroneatzerodiagdoubles}(T, P, m, z, p, S)\).
\end{proof}
\begin{theorem}
\label{ll:paper:not-centerperiodic-of-maskparity-zero-agree-at-zero-diag-doubles}
For every natural number \(T\) and natural number \(P\) and natural number \(m\) and natural number \(z\) and natural number \(p\) and natural number \(S\), if \(0 < P\), then if \(0 < m\), then if \(T \le z + (p + 1)\), then if \(\operatorname{zerodiagonal}(z)\), then if \(\operatorname{doublesat}(m \cdot P, z)\), then if \(\operatorname{settledby}(z + 1, z + (p + 1))\), then if \(\operatorname{settledby}(z + 1 - m \cdot P, z + (p + 1))\), then if \(p + 3 \cdot (m \cdot P) \le S + 1\), then if for every natural number \(s\), if \(s \le S\), then \(\operatorname{maskparity}(z + (p + 1) + s + m \cdot P, 0) = \operatorname{maskparity}(z + (p + 1) + s, 0)\), then not \(\operatorname{centerperiodic}(T, P)\).
\end{theorem}
\begin{proof}
The goal follows from \(\operatorname{notcenterperiodicofmaskparityzeroagreeatzerodiagdoubles}(T, P, m, z, p, S)\).
\end{proof}
\begin{theorem}
\label{ll:paper:exists-zero-diag-white-center-rhcenter-one-flip-of-periodic}
For every natural number \(T\) and natural number \(P\), if \(0 < P\), then if \(\operatorname{centerperiodic}(T, P)\), then for every natural number \(m\), if \(0 < m\), then there exists a natural number \(z\) such that there exists a natural number \(p\) such that there exists a natural number \(s\) such that \(\operatorname{zerodiagonal}(z)\) and \(\operatorname{doublesat}(m \cdot P, z) \land s \le p + 3 \cdot (m \cdot P) \land \operatorname{center}(z + (p + 1) + s) = \operatorname{false} \land \operatorname{rhcenter}(z + (p + 1) + s + m \cdot P, 1) = \operatorname{not}(\operatorname{rhcenter}(z + (p + 1) + s, 1))\).
\end{theorem}
\begin{proof}
Assume \(hP\), \(hper\).
The goal follows from \(\operatorname{existszerodiagwhitecenterrhcenteroneflipofperiodic}(T, P, hP, hper)\).
\end{proof}
\begin{theorem}
\label{ll:paper:exists-zero-diag-maskparity-zero-flip-of-periodic}
For every natural number \(T\) and natural number \(P\), if \(0 < P\), then if \(\operatorname{centerperiodic}(T, P)\), then for every natural number \(m\), if \(0 < m\), then there exists a natural number \(z\) such that there exists a natural number \(p\) such that there exists a natural number \(s\) such that \(\operatorname{zerodiagonal}(z)\) and \(\operatorname{doublesat}(m \cdot P, z) \land s \le p + 3 \cdot (m \cdot P) \land \operatorname{maskparity}(z + (p + 1) + s + m \cdot P, 0) = \operatorname{not}(\operatorname{maskparity}(z + (p + 1) + s, 0))\).
\end{theorem}
\begin{proof}
Assume \(hP\), \(hper\).
The goal follows from \(\operatorname{existszerodiagmaskparityzeroflipofperiodic}(T, P, hP, hper)\).
\end{proof}
\begin{theorem}
\label{ll:paper:exists-zero-diag-lhcenter-one-flip-of-periodic}
For every natural number \(T\) and natural number \(P\), if \(0 < P\), then if \(\operatorname{centerperiodic}(T, P)\), then for every natural number \(m\), if \(0 < m\), then there exists a natural number \(z\) such that there exists a natural number \(p\) such that there exists a natural number \(s\) such that \(\operatorname{zerodiagonal}(z)\) and \(\operatorname{doublesat}(m \cdot P, z) \land s \le p + 3 \cdot (m \cdot P) \land \operatorname{lhcenter}(z + (p + 1) + s + m \cdot P, 1) = \operatorname{not}(\operatorname{lhcenter}(z + (p + 1) + s, 1))\).
\end{theorem}
\begin{proof}
Assume \(hP\), \(hper\).
The goal follows from \(\operatorname{existszerodiaglhcenteroneflipofperiodic}(T, P, hP, hper)\).
\end{proof}
\begin{theorem}
\label{ll:paper:rhcenter-step-one-eq-center-xor-or}
For every natural number \(t\), \(\operatorname{rhcenter}(t + 1, 1) = \operatorname{bxor}(\operatorname{center}(t), \operatorname{rhcenter}(t, 1) \operatorname{or} \operatorname{rhcenter}(t, 2))\).
\end{theorem}
\begin{proof}
The goal follows from \(\operatorname{rhcentersteponeeqcenterxoror}(t)\).
\end{proof}
\begin{theorem}
\label{ll:paper:rhcenter-step-one-eq-not-center-of-black-one}
For every natural number \(t\), if \(\operatorname{rhcenter}(t, 1) = \operatorname{true}\), then \(\operatorname{rhcenter}(t + 1, 1) = \operatorname{not}(\operatorname{center}(t))\).
\end{theorem}
\begin{proof}
The goal follows from \(\operatorname{rhcentersteponeeqnotcenterofblackone}(t)\).
\end{proof}
\begin{theorem}
\label{ll:paper:rhcenter-step-one-eq-true-of-white-center-and-black-one}
For every natural number \(t\), if \(\operatorname{center}(t) = \operatorname{false}\), then if \(\operatorname{rhcenter}(t, 1) = \operatorname{true}\), then \(\operatorname{rhcenter}(t + 1, 1) = \operatorname{true}\).
\end{theorem}
\begin{proof}
The goal follows from \(\operatorname{rhcentersteponeeqtrueofwhitecenterandblackone}(t)\).
\end{proof}
\begin{theorem}
\label{ll:paper:rhcenter-step-one-eq-true-of-maskparity-zero-true}
For every natural number \(t\), if \(\operatorname{maskparity}(t, 0) = \operatorname{true}\), then \(\operatorname{rhcenter}(t + 1, 1) = \operatorname{true}\).
\end{theorem}
\begin{proof}
The goal follows from \(\operatorname{rhcentersteponeeqtrueofmaskparityzerotrue}(t)\).
\end{proof}
\begin{theorem}
\label{ll:paper:rhcenter-step-two-eq-not-center-of-maskparity-zero-true}
For every natural number \(t\), if \(\operatorname{maskparity}(t, 0) = \operatorname{true}\), then \(\operatorname{rhcenter}(t + 2, 1) = \operatorname{not}(\operatorname{center}(t + 1))\).
\end{theorem}
\begin{proof}
The goal follows from \(\operatorname{rhcentersteptwoeqnotcenterofmaskparityzerotrue}(t)\).
\end{proof}
\begin{theorem}
\label{ll:paper:periodic-rhcenter-step-one-eq-of-black-one}
For every natural number \(T\) and natural number \(P\), if \(0 < P\), then if \(\operatorname{centerperiodic}(T, P)\), then for every natural number \(m\) and natural number \(t\), if \(T \le t\), then if \(\operatorname{rhcenter}(t + m \cdot P, 1) = \operatorname{true}\), then if \(\operatorname{rhcenter}(t, 1) = \operatorname{true}\), then \(\operatorname{rhcenter}(t + 1 + m \cdot P, 1) = \operatorname{rhcenter}(t + 1, 1)\).
\end{theorem}
\begin{proof}
Assume \(hP\), \(hper\).
The goal follows from \(\operatorname{periodicrhcentersteponeeqofblackone}(T, P, hP, hper)\).
\end{proof}
\begin{theorem}
\label{ll:paper:periodic-lhcenter-step-one-eq-of-black-one}
For every natural number \(T\) and natural number \(P\), if \(0 < P\), then if \(\operatorname{centerperiodic}(T, P)\), then for every natural number \(m\) and natural number \(t\), if \(T \le t\), then if \(\operatorname{rhcenter}(t + m \cdot P, 1) = \operatorname{true}\), then if \(\operatorname{rhcenter}(t, 1) = \operatorname{true}\), then \(\operatorname{lhcenter}(t + 1 + m \cdot P, 1) = \operatorname{lhcenter}(t + 1, 1)\).
\end{theorem}
\begin{proof}
Assume \(hP\), \(hper\).
The goal follows from \(\operatorname{periodiclhcentersteponeeqofblackone}(T, P, hP, hper)\).
\end{proof}
\begin{theorem}
\label{ll:paper:periodic-lhcenter-step-one-and-two-eq-of-maskparity-zero-true}
For every natural number \(T\) and natural number \(P\), if \(0 < P\), then if \(\operatorname{centerperiodic}(T, P)\), then for every natural number \(m\) and natural number \(t\), if \(T \le t\), then if \(\operatorname{maskparity}(t + m \cdot P, 0) = \operatorname{true}\), then if \(\operatorname{maskparity}(t, 0) = \operatorname{true}\), then \(\operatorname{lhcenter}(t + 1 + m \cdot P, 1) = \operatorname{lhcenter}(t + 1, 1)\) and \(\operatorname{lhcenter}(t + 2 + m \cdot P, 1) = \operatorname{lhcenter}(t + 2, 1)\).
\end{theorem}
\begin{proof}
Assume \(hP\), \(hper\).
The goal follows from \(\operatorname{periodiclhcentersteponeandtwoeqofmaskparityzerotrue}(T, P, hP, hper)\).
\end{proof}
\begin{theorem}
\label{ll:paper:periodic-lhbandagree-zero-of-white-center-rhoragree-one}
For every natural number \(T\) and natural number \(P\), if \(0 < P\), then if \(\operatorname{centerperiodic}(T, P)\), then for every natural number \(m\) and natural number \(t\) and natural number \(S\), if \(T \le t\), then if if \(\operatorname{center}(t) = \operatorname{false}\), then \(\operatorname{rhcenter}(t + m \cdot P, 1) = \operatorname{rhcenter}(t, 1)\), then if for every natural number \(s\), if \(s < S\), then if \(\operatorname{center}(t + (s + 1)) = \operatorname{false}\), then \((\operatorname{rhcenter}(t + s + m \cdot P, 1) \operatorname{or} \operatorname{rhcenter}(t + s + m \cdot P, 2)) = (\operatorname{rhcenter}(t + s, 1) \operatorname{or} \operatorname{rhcenter}(t + s, 2))\), then \(\operatorname{lhbandagree}(t, m \cdot P, 0, S + 1)\).
\end{theorem}
\begin{proof}
Assume \(hP\), \(hper\).
The goal follows from \(\operatorname{periodiclhbandagreezeroofwhitecenterrhoragreeone}(T, P, hP, hper)\).
\end{proof}
\begin{theorem}
\label{ll:paper:not-centerperiodic-of-white-center-rhoragree-one-doubles}
For every natural number \(T\) and natural number \(P\) and natural number \(m\) and natural number \(t\) and natural number \(p\) and natural number \(S\), if \(0 < P\), then if \(0 < m\), then if \(T \le t\), then if \(\operatorname{zerodiagonal}(t - (p + 1))\), then if \(\operatorname{doublesat}(m \cdot P, t - (p + 1))\), then if \(\operatorname{settledby}(t - (p + 1) + 1, t)\), then if \(\operatorname{settledby}(t - (p + m \cdot P), t)\), then if \(p + 3 \cdot (m \cdot P) \le S + 1\), then if if \(\operatorname{center}(t) = \operatorname{false}\), then \(\operatorname{rhcenter}(t + m \cdot P, 1) = \operatorname{rhcenter}(t, 1)\), then if for every natural number \(s\), if \(s < S\), then if \(\operatorname{center}(t + (s + 1)) = \operatorname{false}\), then \((\operatorname{rhcenter}(t + s + m \cdot P, 1) \operatorname{or} \operatorname{rhcenter}(t + s + m \cdot P, 2)) = (\operatorname{rhcenter}(t + s, 1) \operatorname{or} \operatorname{rhcenter}(t + s, 2))\), then not \(\operatorname{centerperiodic}(T, P)\).
\end{theorem}
\begin{proof}
The goal follows from \(\operatorname{notcenterperiodicofwhitecenterrhoragreeonedoubles}(T, P, m, t, p, S)\).
\end{proof}
\begin{theorem}
\label{ll:paper:rhcenter-step-one-eq-center-xor-two-of-white-one}
For every natural number \(t\), if \(\operatorname{rhcenter}(t, 1) = \operatorname{false}\), then \(\operatorname{rhcenter}(t + 1, 1) = \operatorname{bxor}(\operatorname{center}(t), \operatorname{rhcenter}(t, 2))\).
\end{theorem}
\begin{proof}
The goal follows from \(\operatorname{rhcentersteponeeqcenterxortwoofwhiteone}(t)\).
\end{proof}
\begin{theorem}
\label{ll:paper:periodic-rhcenter-step-one-eq-of-white-one-and-two-eq}
For every natural number \(T\) and natural number \(P\), if \(0 < P\), then if \(\operatorname{centerperiodic}(T, P)\), then for every natural number \(m\) and natural number \(t\), if \(T \le t\), then if \(\operatorname{rhcenter}(t + m \cdot P, 1) = \operatorname{false}\), then if \(\operatorname{rhcenter}(t, 1) = \operatorname{false}\), then if \(\operatorname{rhcenter}(t + m \cdot P, 2) = \operatorname{rhcenter}(t, 2)\), then \(\operatorname{rhcenter}(t + 1 + m \cdot P, 1) = \operatorname{rhcenter}(t + 1, 1)\).
\end{theorem}
\begin{proof}
Assume \(hP\), \(hper\).
The goal follows from \(\operatorname{periodicrhcentersteponeeqofwhiteoneandtwoeq}(T, P, hP, hper)\).
\end{proof}
\begin{theorem}
\label{ll:paper:periodic-lhcenter-step-one-eq-of-white-or-black-one}
For every natural number \(T\) and natural number \(P\), if \(0 < P\), then if \(\operatorname{centerperiodic}(T, P)\), then for every natural number \(m\) and natural number \(t\), if \(T \le t\), then if \(\operatorname{rhcenter}(t + m \cdot P, 1) = \operatorname{rhcenter}(t, 1)\), then if if \(\operatorname{center}(t + 1) = \operatorname{false}\), then if \(\operatorname{rhcenter}(t, 1) = \operatorname{false}\), then \(\operatorname{rhcenter}(t + m \cdot P, 2) = \operatorname{rhcenter}(t, 2)\), then \(\operatorname{lhcenter}(t + 1 + m \cdot P, 1) = \operatorname{lhcenter}(t + 1, 1)\).
\end{theorem}
\begin{proof}
Assume \(hP\), \(hper\).
The goal follows from \(\operatorname{periodiclhcentersteponeeqofwhiteorblackone}(T, P, hP, hper)\).
\end{proof}
\begin{theorem}
\label{ll:paper:cell-negpow-eq-lhcenter-twopow}
For every natural number \(t\) and natural number \(k\), \(\operatorname{cell}(t, \operatorname{negpow}(k)) = \operatorname{lhcenter}(t, \operatorname{twopow}(k))\).
\end{theorem}
\begin{proof}
The goal follows from \(\operatorname{cellnegpoweqlhcentertwopow}(t, k)\).
\end{proof}
\begin{theorem}
\label{ll:paper:scale-center-lhcenter}
For every natural number \(t\) and natural number \(k\), \(\operatorname{center}(t + \operatorname{twopow}(k)) = \operatorname{bxor}(\operatorname{bxor}(\operatorname{center}(t), \operatorname{lhcenter}(t, \operatorname{twopow}(k))), \operatorname{response}(t, k))\).
\end{theorem}
\begin{proof}
The goal follows from \(\operatorname{scalecenterlhcenter}(t, k)\).
\end{proof}
\begin{theorem}
\label{ll:paper:periodic-lhcenter-twopow-defect-eq-response-defect}
For every natural number \(T\) and natural number \(P\), if \(0 < P\), then if \(\operatorname{centerperiodic}(T, P)\), then for every natural number \(m\) and natural number \(t\) and natural number \(k\), if \(T \le t\), then \(\operatorname{bxor}(\operatorname{lhcenter}(t + m \cdot P, \operatorname{twopow}(k)), \operatorname{lhcenter}(t, \operatorname{twopow}(k))) = \operatorname{bxor}(\operatorname{response}(t + m \cdot P, k), \operatorname{response}(t, k))\).
\end{theorem}
\begin{proof}
Assume \(hP\), \(hper\).
The goal follows from \(\operatorname{periodiclhcentertwopowdefecteqresponsedefect}(T, P, hP, hper)\).
\end{proof}
\begin{theorem}
\label{ll:paper:periodic-maskparity-defect-eq-lhcenter-twopow-xor-leftresponse-defect}
For every natural number \(T\) and natural number \(P\), if \(0 < P\), then if \(\operatorname{centerperiodic}(T, P)\), then for every natural number \(m\) and natural number \(t\) and natural number \(k\), if \(T \le t\), then \(\operatorname{bxor}(\operatorname{maskparity}(t + m \cdot P, k), \operatorname{maskparity}(t, k)) = \operatorname{bxor}(\operatorname{bxor}(\operatorname{lhcenter}(t + m \cdot P, \operatorname{twopow}(k)), \operatorname{lhcenter}(t, \operatorname{twopow}(k))), \operatorname{bxor}(\operatorname{leftresponse}(t + m \cdot P, k), \operatorname{leftresponse}(t, k)))\).
\end{theorem}
\begin{proof}
Assume \(hP\), \(hper\).
The goal follows from \(\operatorname{periodicmaskparitydefecteqlhcentertwopowxorleftresponsedefect}(T, P, hP, hper)\).
\end{proof}
\begin{theorem}
\label{ll:paper:periodic-response-eq-cone-mul}
For every natural number \(T\) and natural number \(P\), if \(0 < P\), then if \(\operatorname{centerperiodic}(T, P)\), then for every natural number \(m\) and natural number \(t\) and natural number \(k\), if \(T \le t\), then if \(t + m \cdot P < \operatorname{twopow}(k)\), then \(\operatorname{response}(t + m \cdot P, k) = \operatorname{response}(t, k)\).
\end{theorem}
\begin{proof}
Assume \(hP\), \(hper\).
The goal follows from \(\operatorname{periodicresponseeqconemul}(T, P, hP, hper)\).
\end{proof}
\begin{theorem}
\label{ll:paper:periodic-maskparity-defect-eq-leftresponse-defect-cone-mul}
For every natural number \(T\) and natural number \(P\), if \(0 < P\), then if \(\operatorname{centerperiodic}(T, P)\), then for every natural number \(m\) and natural number \(t\) and natural number \(k\), if \(T \le t\), then if \(t + m \cdot P < \operatorname{twopow}(k)\), then \(\operatorname{bxor}(\operatorname{maskparity}(t + m \cdot P, k), \operatorname{maskparity}(t, k)) = \operatorname{bxor}(\operatorname{leftresponse}(t + m \cdot P, k), \operatorname{leftresponse}(t, k))\).
\end{theorem}
\begin{proof}
Assume \(hP\), \(hper\).
The goal follows from \(\operatorname{periodicmaskparitydefecteqleftresponsedefectconemul}(T, P, hP, hper)\).
\end{proof}
\begin{theorem}
\label{ll:paper:periodic-response-flip-at-twopow-mul}
For every natural number \(T\) and natural number \(P\), if \(0 < P\), then if \(\operatorname{centerperiodic}(T, P)\), then for every natural number \(m\) and natural number \(t\) and natural number \(k\), if \(0 < m\), then if \(T \le t\), then if \(t + m \cdot P = \operatorname{twopow}(k)\), then \(\operatorname{response}(t + m \cdot P, k) = \operatorname{not}(\operatorname{response}(t, k))\).
\end{theorem}
\begin{proof}
Assume \(hP\), \(hper\).
The goal follows from \(\operatorname{periodicresponseflipattwopowmul}(T, P, hP, hper)\).
\end{proof}
\begin{theorem}
\label{ll:paper:periodic-maskparity-defect-eq-not-leftresponse-defect-at-twopow-mul}
For every natural number \(T\) and natural number \(P\), if \(0 < P\), then if \(\operatorname{centerperiodic}(T, P)\), then for every natural number \(m\) and natural number \(t\) and natural number \(k\), if \(0 < m\), then if \(T \le t\), then if \(t + m \cdot P = \operatorname{twopow}(k)\), then \(\operatorname{bxor}(\operatorname{maskparity}(t + m \cdot P, k), \operatorname{maskparity}(t, k)) = \operatorname{not}(\operatorname{bxor}(\operatorname{leftresponse}(t + m \cdot P, k), \operatorname{leftresponse}(t, k)))\).
\end{theorem}
\begin{proof}
Assume \(hP\), \(hper\).
The goal follows from \(\operatorname{periodicmaskparitydefecteqnotleftresponsedefectattwopowmul}(T, P, hP, hper)\).
\end{proof}
\begin{theorem}
\label{ll:paper:periodic-lhcenter-one-defect-eq-maskparity-zero-defect}
For every natural number \(T\) and natural number \(P\), if \(0 < P\), then if \(\operatorname{centerperiodic}(T, P)\), then for every natural number \(m\) and natural number \(t\), if \(T \le t\), then \(\operatorname{bxor}(\operatorname{lhcenter}(t + m \cdot P, 1), \operatorname{lhcenter}(t, 1)) = \operatorname{bxor}(\operatorname{maskparity}(t + m \cdot P, 0), \operatorname{maskparity}(t, 0))\).
\end{theorem}
\begin{proof}
Assume \(hP\), \(hper\).
The goal follows from \(\operatorname{periodiclhcenteronedefecteqmaskparityzerodefect}(T, P, hP, hper)\).
\end{proof}
\begin{theorem}
\label{ll:paper:periodic-maskparity-zero-eq-of-lhcenter-one-eq}
For every natural number \(T\) and natural number \(P\), if \(0 < P\), then if \(\operatorname{centerperiodic}(T, P)\), then for every natural number \(m\) and natural number \(t\), if \(T \le t\), then if \(\operatorname{lhcenter}(t + m \cdot P, 1) = \operatorname{lhcenter}(t, 1)\), then \(\operatorname{maskparity}(t + m \cdot P, 0) = \operatorname{maskparity}(t, 0)\).
\end{theorem}
\begin{proof}
Assume \(hP\), \(hper\).
The goal follows from \(\operatorname{periodicmaskparityzeroeqoflhcenteroneeq}(T, P, hP, hper)\).
\end{proof}
\begin{theorem}
\label{ll:paper:periodic-lhcenter-one-dyadic-jump-defect-eq-maskparity-step-defect}
For every natural number \(T\) and natural number \(P\), if \(0 < P\), then if \(\operatorname{centerperiodic}(T, P)\), then for every natural number \(m\) and natural number \(t\) and natural number \(k\), if \(T \le t\), then \(\operatorname{bxor}(\operatorname{bxor}(\operatorname{lhcenter}(t + \operatorname{twopow}(k) + m \cdot P, 1), \operatorname{lhcenter}(t + \operatorname{twopow}(k), 1)), \operatorname{bxor}(\operatorname{lhcenter}(t + m \cdot P, 1), \operatorname{lhcenter}(t, 1))) = \operatorname{bxor}(\operatorname{bxor}(\operatorname{maskparity}(t + 1 + m \cdot P, k), \operatorname{maskparity}(t + 1, k)), \operatorname{bxor}(\operatorname{maskparity}(t + m \cdot P, k), \operatorname{maskparity}(t, k)))\).
\end{theorem}
\begin{proof}
Assume \(hP\), \(hper\).
The goal follows from \(\operatorname{periodiclhcenteronedyadicjumpdefecteqmaskparitystepdefect}(T, P, hP, hper)\).
\end{proof}
\begin{theorem}
\label{ll:paper:periodic-lhcenter-one-dyadic-jump-defect-eq-leftresponse-step-defect-cone-mul}
For every natural number \(T\) and natural number \(P\), if \(0 < P\), then if \(\operatorname{centerperiodic}(T, P)\), then for every natural number \(m\) and natural number \(t\) and natural number \(k\), if \(T \le t\), then if \(t + 1 + m \cdot P < \operatorname{twopow}(k)\), then \(\operatorname{bxor}(\operatorname{bxor}(\operatorname{lhcenter}(t + \operatorname{twopow}(k) + m \cdot P, 1), \operatorname{lhcenter}(t + \operatorname{twopow}(k), 1)), \operatorname{bxor}(\operatorname{lhcenter}(t + m \cdot P, 1), \operatorname{lhcenter}(t, 1))) = \operatorname{bxor}(\operatorname{bxor}(\operatorname{leftresponse}(t + 1 + m \cdot P, k), \operatorname{leftresponse}(t + 1, k)), \operatorname{bxor}(\operatorname{leftresponse}(t + m \cdot P, k), \operatorname{leftresponse}(t, k)))\).
\end{theorem}
\begin{proof}
Assume \(hP\), \(hper\).
The goal follows from \(\operatorname{periodiclhcenteronedyadicjumpdefecteqleftresponsestepdefectconemul}(T, P, hP, hper)\).
\end{proof}
\begin{theorem}
\label{ll:paper:periodic-lhcenter-one-dyadic-jump-defect-eq-not-leftresponse-step-defect-at-twopow-mul}
For every natural number \(T\) and natural number \(P\), if \(0 < P\), then if \(\operatorname{centerperiodic}(T, P)\), then for every natural number \(m\) and natural number \(t\) and natural number \(k\), if \(0 < m\), then if \(T \le t\), then if \(t + 1 + m \cdot P = \operatorname{twopow}(k)\), then \(\operatorname{bxor}(\operatorname{bxor}(\operatorname{lhcenter}(t + \operatorname{twopow}(k) + m \cdot P, 1), \operatorname{lhcenter}(t + \operatorname{twopow}(k), 1)), \operatorname{bxor}(\operatorname{lhcenter}(t + m \cdot P, 1), \operatorname{lhcenter}(t, 1))) = \operatorname{not}(\operatorname{bxor}(\operatorname{bxor}(\operatorname{leftresponse}(t + 1 + m \cdot P, k), \operatorname{leftresponse}(t + 1, k)), \operatorname{bxor}(\operatorname{leftresponse}(t + m \cdot P, k), \operatorname{leftresponse}(t, k))))\).
\end{theorem}
\begin{proof}
Assume \(hP\), \(hper\).
The goal follows from \(\operatorname{periodiclhcenteronedyadicjumpdefecteqnotleftresponsestepdefectattwopowmul}(T, P, hP, hper)\).
\end{proof}
\begin{theorem}
\label{ll:paper:periodic-lhcenter-one-dyadic-jump-defect-eq-lhcenter-twopow-step-xor-leftresponse-step-defect}
For every natural number \(T\) and natural number \(P\), if \(0 < P\), then if \(\operatorname{centerperiodic}(T, P)\), then for every natural number \(m\) and natural number \(t\) and natural number \(k\), if \(T \le t\), then \(\operatorname{bxor}(\operatorname{bxor}(\operatorname{lhcenter}(t + \operatorname{twopow}(k) + m \cdot P, 1), \operatorname{lhcenter}(t + \operatorname{twopow}(k), 1)), \operatorname{bxor}(\operatorname{lhcenter}(t + m \cdot P, 1), \operatorname{lhcenter}(t, 1))) = \operatorname{bxor}(\operatorname{bxor}(\operatorname{bxor}(\operatorname{lhcenter}(t + 1 + m \cdot P, \operatorname{twopow}(k)), \operatorname{lhcenter}(t + 1, \operatorname{twopow}(k))), \operatorname{bxor}(\operatorname{lhcenter}(t + m \cdot P, \operatorname{twopow}(k)), \operatorname{lhcenter}(t, \operatorname{twopow}(k)))), \operatorname{bxor}(\operatorname{bxor}(\operatorname{leftresponse}(t + 1 + m \cdot P, k), \operatorname{leftresponse}(t + 1, k)), \operatorname{bxor}(\operatorname{leftresponse}(t + m \cdot P, k), \operatorname{leftresponse}(t, k))))\).
\end{theorem}
\begin{proof}
Assume \(hP\), \(hper\).
The goal follows from \(\operatorname{periodiclhcenteronedyadicjumpdefecteqlhcentertwopowstepxorleftresponsestepdefect}(T, P, hP, hper)\).
\end{proof}
\begin{theorem}
\label{ll:paper:lhcenter-twopow-of-zero-diag}
For every natural number \(z\) and natural number \(k\), if \(\operatorname{zerodiagonal}(z)\), then if \(\operatorname{settledby}(z, z + \operatorname{twopow}(k))\), then \(\operatorname{lhcenter}(z + \operatorname{twopow}(k), \operatorname{twopow}(k)) = \operatorname{false}\).
\end{theorem}
\begin{proof}
The goal follows from \(\operatorname{lhcentertwopowofzerodiag}(z, k)\).
\end{proof}
\begin{theorem}
\label{ll:paper:lhcenter-twopow-after-zero-diag}
For every natural number \(z\) and natural number \(k\), if \(\operatorname{zerodiagonal}(z)\), then if \(\operatorname{settledby}(z + 2, z + 2 + \operatorname{twopow}(k))\), then \(\operatorname{lhcenter}(z + 2 + \operatorname{twopow}(k), \operatorname{twopow}(k)) = \operatorname{true}\).
\end{theorem}
\begin{proof}
The goal follows from \(\operatorname{lhcentertwopowafterzerodiag}(z, k)\).
\end{proof}
\begin{theorem}
\label{ll:paper:scale-center-at-zero-diag}
For every natural number \(z\) and natural number \(k\), if \(\operatorname{zerodiagonal}(z)\), then if \(\operatorname{settledby}(z, z + \operatorname{twopow}(k))\), then \(\operatorname{center}(z + \operatorname{twopow}(k) + \operatorname{twopow}(k)) = \operatorname{bxor}(\operatorname{center}(z + \operatorname{twopow}(k)), \operatorname{response}(z + \operatorname{twopow}(k), k))\).
\end{theorem}
\begin{proof}
The goal follows from \(\operatorname{scalecenteratzerodiag}(z, k)\).
\end{proof}
\begin{theorem}
\label{ll:paper:scale-center-after-zero-diag}
For every natural number \(z\) and natural number \(k\), if \(\operatorname{zerodiagonal}(z)\), then if \(\operatorname{settledby}(z + 2, z + 2 + \operatorname{twopow}(k))\), then \(\operatorname{center}(z + 2 + \operatorname{twopow}(k) + \operatorname{twopow}(k)) = \operatorname{not}(\operatorname{bxor}(\operatorname{center}(z + 2 + \operatorname{twopow}(k)), \operatorname{response}(z + 2 + \operatorname{twopow}(k), k)))\).
\end{theorem}
\begin{proof}
The goal follows from \(\operatorname{scalecenterafterzerodiag}(z, k)\).
\end{proof}
\begin{theorem}
\label{ll:paper:periodic-response-at-zero-diag}
For every natural number \(T\) and natural number \(P\), if \(0 < P\), then if \(\operatorname{centerperiodic}(T, P)\), then for every natural number \(z\) and natural number \(k\), if \(T \le z + \operatorname{twopow}(k)\), then if \(\operatorname{zerodiagonal}(z)\), then if \(\operatorname{settledby}(z, z + \operatorname{twopow}(k))\), then \(\operatorname{response}(z + \operatorname{twopow}(k), k) = \operatorname{bxor}(\operatorname{center}(z + \operatorname{twopow}(k) + \operatorname{lag}(P, k)), \operatorname{center}(z + \operatorname{twopow}(k)))\).
\end{theorem}
\begin{proof}
Assume \(hP\), \(hper\).
The goal follows from \(\operatorname{periodicresponseatzerodiag}(T, P, hP, hper)\).
\end{proof}
\begin{theorem}
\label{ll:paper:periodic-response-after-zero-diag}
For every natural number \(T\) and natural number \(P\), if \(0 < P\), then if \(\operatorname{centerperiodic}(T, P)\), then for every natural number \(z\) and natural number \(k\), if \(T \le z + 2 + \operatorname{twopow}(k)\), then if \(\operatorname{zerodiagonal}(z)\), then if \(\operatorname{settledby}(z + 2, z + 2 + \operatorname{twopow}(k))\), then \(\operatorname{response}(z + 2 + \operatorname{twopow}(k), k) = \operatorname{not}(\operatorname{bxor}(\operatorname{center}(z + 2 + \operatorname{twopow}(k) + \operatorname{lag}(P, k)), \operatorname{center}(z + 2 + \operatorname{twopow}(k))))\).
\end{theorem}
\begin{proof}
Assume \(hP\), \(hper\).
The goal follows from \(\operatorname{periodicresponseafterzerodiag}(T, P, hP, hper)\).
\end{proof}
\begin{theorem}
\label{ll:paper:periodic-response-defect-at-zero-diag-eq-lhcenter-twopow}
For every natural number \(T\) and natural number \(P\), if \(0 < P\), then if \(\operatorname{centerperiodic}(T, P)\), then for every natural number \(m\) and natural number \(z\) and natural number \(k\), if \(T \le z + \operatorname{twopow}(k)\), then if \(\operatorname{zerodiagonal}(z)\), then if \(\operatorname{settledby}(z, z + \operatorname{twopow}(k))\), then \(\operatorname{bxor}(\operatorname{response}(z + \operatorname{twopow}(k) + m \cdot P, k), \operatorname{response}(z + \operatorname{twopow}(k), k)) = \operatorname{lhcenter}(z + \operatorname{twopow}(k) + m \cdot P, \operatorname{twopow}(k))\).
\end{theorem}
\begin{proof}
Assume \(hP\), \(hper\).
The goal follows from \(\operatorname{periodicresponsedefectatzerodiageqlhcentertwopow}(T, P, hP, hper)\).
\end{proof}
\begin{theorem}
\label{ll:paper:periodic-response-defect-after-zero-diag-eq-not-lhcenter-twopow}
For every natural number \(T\) and natural number \(P\), if \(0 < P\), then if \(\operatorname{centerperiodic}(T, P)\), then for every natural number \(m\) and natural number \(z\) and natural number \(k\), if \(T \le z + 2 + \operatorname{twopow}(k)\), then if \(\operatorname{zerodiagonal}(z)\), then if \(\operatorname{settledby}(z + 2, z + 2 + \operatorname{twopow}(k))\), then \(\operatorname{bxor}(\operatorname{response}(z + 2 + \operatorname{twopow}(k) + m \cdot P, k), \operatorname{response}(z + 2 + \operatorname{twopow}(k), k)) = \operatorname{not}(\operatorname{lhcenter}(z + 2 + \operatorname{twopow}(k) + m \cdot P, \operatorname{twopow}(k)))\).
\end{theorem}
\begin{proof}
Assume \(hP\), \(hper\).
The goal follows from \(\operatorname{periodicresponsedefectafterzerodiageqnotlhcentertwopow}(T, P, hP, hper)\).
\end{proof}
\begin{theorem}
\label{ll:paper:periodic-response-eq-of-lhbandagree-twopow}
For every natural number \(T\) and natural number \(P\), if \(0 < P\), then if \(\operatorname{centerperiodic}(T, P)\), then for every natural number \(m\) and natural number \(t\) and natural number \(k\) and natural number \(p0\) and natural number \(L\), if \(T \le t\), then if \(p0 \le \operatorname{twopow}(k)\), then if \(\operatorname{twopow}(k) \le p0 + L\), then if \(\operatorname{lhbandagree}(t, m \cdot P, p0, L)\), then \(\operatorname{response}(t + m \cdot P, k) = \operatorname{response}(t, k)\).
\end{theorem}
\begin{proof}
Assume \(hP\), \(hper\).
The goal follows from \(\operatorname{periodicresponseeqoflhbandagreetwopow}(T, P, hP, hper)\).
\end{proof}
\begin{theorem}
\label{ll:paper:periodic-maskparity-defect-eq-leftresponse-defect-of-lhbandagree-twopow}
For every natural number \(T\) and natural number \(P\), if \(0 < P\), then if \(\operatorname{centerperiodic}(T, P)\), then for every natural number \(m\) and natural number \(t\) and natural number \(k\) and natural number \(p0\) and natural number \(L\), if \(T \le t\), then if \(p0 \le \operatorname{twopow}(k)\), then if \(\operatorname{twopow}(k) \le p0 + L\), then if \(\operatorname{lhbandagree}(t, m \cdot P, p0, L)\), then \(\operatorname{bxor}(\operatorname{maskparity}(t + m \cdot P, k), \operatorname{maskparity}(t, k)) = \operatorname{bxor}(\operatorname{leftresponse}(t + m \cdot P, k), \operatorname{leftresponse}(t, k))\).
\end{theorem}
\begin{proof}
Assume \(hP\), \(hper\).
The goal follows from \(\operatorname{periodicmaskparitydefecteqleftresponsedefectoflhbandagreetwopow}(T, P, hP, hper)\).
\end{proof}
\begin{theorem}
\label{ll:paper:periodic-response-eq-of-white-center-rhcenter-one-twopow}
For every natural number \(T\) and natural number \(P\), if \(0 < P\), then if \(\operatorname{centerperiodic}(T, P)\), then for every natural number \(m\) and natural number \(t\) and natural number \(k\) and natural number \(S\), if \(T \le t\), then if \(\operatorname{twopow}(k) \le S + 1\), then if for every natural number \(s\), if \(s \le S\), then if \(\operatorname{center}(t + s) = \operatorname{false}\), then \(\operatorname{rhcenter}(t + s + m \cdot P, 1) = \operatorname{rhcenter}(t + s, 1)\), then \(\operatorname{response}(t + m \cdot P, k) = \operatorname{response}(t, k)\).
\end{theorem}
\begin{proof}
Assume \(hP\), \(hper\).
The goal follows from \(\operatorname{periodicresponseeqofwhitecenterrhcenteronetwopow}(T, P, hP, hper)\).
\end{proof}
\begin{theorem}
\label{ll:paper:periodic-maskparity-defect-eq-leftresponse-defect-of-white-center-rhcenter-one-twopow}
For every natural number \(T\) and natural number \(P\), if \(0 < P\), then if \(\operatorname{centerperiodic}(T, P)\), then for every natural number \(m\) and natural number \(t\) and natural number \(k\) and natural number \(S\), if \(T \le t\), then if \(\operatorname{twopow}(k) \le S + 1\), then if for every natural number \(s\), if \(s \le S\), then if \(\operatorname{center}(t + s) = \operatorname{false}\), then \(\operatorname{rhcenter}(t + s + m \cdot P, 1) = \operatorname{rhcenter}(t + s, 1)\), then \(\operatorname{bxor}(\operatorname{maskparity}(t + m \cdot P, k), \operatorname{maskparity}(t, k)) = \operatorname{bxor}(\operatorname{leftresponse}(t + m \cdot P, k), \operatorname{leftresponse}(t, k))\).
\end{theorem}
\begin{proof}
Assume \(hP\), \(hper\).
The goal follows from \(\operatorname{periodicmaskparitydefecteqleftresponsedefectofwhitecenterrhcenteronetwopow}(T, P, hP, hper)\).
\end{proof}
\begin{theorem}
\label{ll:paper:response-coboundary-eq-lhcenter-twopow}
For every natural number \(t\) and natural number \(k\), \(\operatorname{bxor}(\operatorname{bxor}(\operatorname{response}(t, k + 1), \operatorname{response}(t, k)), \operatorname{response}(t + \operatorname{twopow}(k), k)) = \operatorname{bxor}(\operatorname{bxor}(\operatorname{lhcenter}(t + \operatorname{twopow}(k), \operatorname{twopow}(k)), \operatorname{lhcenter}(t, \operatorname{twopow}(k))), \operatorname{lhcenter}(t, \operatorname{twopow}(k + 1)))\).
\end{theorem}
\begin{proof}
The goal follows from \(\operatorname{responsecoboundaryeqlhcentertwopow}(t, k)\).
\end{proof}
\begin{theorem}
\label{ll:paper:leftresponse-coboundary-eq-lhcenter-twopow}
For every natural number \(t\) and natural number \(k\), \(\operatorname{bxor}(\operatorname{bxor}(\operatorname{leftresponse}(t, k + 1), \operatorname{leftresponse}(t, k)), \operatorname{leftresponse}(t + \operatorname{twopow}(k), k)) = \operatorname{bxor}(\operatorname{bxor}(\operatorname{lhcenter}(t + \operatorname{twopow}(k), \operatorname{twopow}(k)), \operatorname{lhcenter}(t, \operatorname{twopow}(k))), \operatorname{lhcenter}(t, \operatorname{twopow}(k + 1)))\).
\end{theorem}
\begin{proof}
The goal follows from \(\operatorname{leftresponsecoboundaryeqlhcentertwopow}(t, k)\).
\end{proof}
\begin{theorem}
\label{ll:paper:lhcenter-twopow-eq-response-coboundary-cone}
For every natural number \(t\) and natural number \(k\), if \(t < \operatorname{twopow}(k)\), then \(\operatorname{lhcenter}(t + \operatorname{twopow}(k), \operatorname{twopow}(k)) = \operatorname{bxor}(\operatorname{bxor}(\operatorname{response}(t, k + 1), \operatorname{response}(t, k)), \operatorname{response}(t + \operatorname{twopow}(k), k))\).
\end{theorem}
\begin{proof}
The goal follows from \(\operatorname{lhcentertwopoweqresponsecoboundarycone}(t, k)\).
\end{proof}
\begin{theorem}
\label{ll:paper:lhcenter-twopow-eq-leftresponse-coboundary-cone}
For every natural number \(t\) and natural number \(k\), if \(t < \operatorname{twopow}(k)\), then \(\operatorname{lhcenter}(t + \operatorname{twopow}(k), \operatorname{twopow}(k)) = \operatorname{bxor}(\operatorname{bxor}(\operatorname{leftresponse}(t, k + 1), \operatorname{leftresponse}(t, k)), \operatorname{leftresponse}(t + \operatorname{twopow}(k), k))\).
\end{theorem}
\begin{proof}
The goal follows from \(\operatorname{lhcentertwopoweqleftresponsecoboundarycone}(t, k)\).
\end{proof}
\begin{theorem}
\label{ll:paper:leftresponse-coboundary-at-zero-diag-cone}
For every natural number \(z\) and natural number \(k\), if \(z < \operatorname{twopow}(k)\), then if \(\operatorname{zerodiagonal}(z)\), then if \(\operatorname{settledby}(z, z + \operatorname{twopow}(k))\), then \(\operatorname{bxor}(\operatorname{bxor}(\operatorname{leftresponse}(z, k + 1), \operatorname{leftresponse}(z, k)), \operatorname{leftresponse}(z + \operatorname{twopow}(k), k)) = \operatorname{false}\).
\end{theorem}
\begin{proof}
The goal follows from \(\operatorname{leftresponsecoboundaryatzerodiagcone}(z, k)\).
\end{proof}
\begin{theorem}
\label{ll:paper:leftresponse-coboundary-after-zero-diag-cone}
For every natural number \(z\) and natural number \(k\), if \(z + 2 < \operatorname{twopow}(k)\), then if \(\operatorname{zerodiagonal}(z)\), then if \(\operatorname{settledby}(z + 2, z + 2 + \operatorname{twopow}(k))\), then \(\operatorname{bxor}(\operatorname{bxor}(\operatorname{leftresponse}(z + 2, k + 1), \operatorname{leftresponse}(z + 2, k)), \operatorname{leftresponse}(z + 2 + \operatorname{twopow}(k), k)) = \operatorname{true}\).
\end{theorem}
\begin{proof}
The goal follows from \(\operatorname{leftresponsecoboundaryafterzerodiagcone}(z, k)\).
\end{proof}
\begin{theorem}
\label{ll:paper:lhcenter-twopow-succ-antiperiodic-of-doubles}
For every natural number \(z\) and natural number \(k\), if \(\operatorname{zerodiagonal}(z)\), then if \(\operatorname{doublesat}(\operatorname{twopow}(k), z)\), then if \(\operatorname{settledby}(z + 1, z + 1 + \operatorname{twopow}(k))\), then \(\operatorname{lhcenter}(z + 1 + \operatorname{twopow}(k + 1), \operatorname{twopow}(k + 1)) = \operatorname{not}(\operatorname{lhcenter}(z + 1 + \operatorname{twopow}(k), \operatorname{twopow}(k)))\).
\end{theorem}
\begin{proof}
The goal follows from \(\operatorname{lhcentertwopowsuccantiperiodicofdoubles}(z, k)\).
\end{proof}
\begin{theorem}
\label{ll:paper:periodic-maskparity-at-zero-diag}
For every natural number \(T\) and natural number \(P\), if \(0 < P\), then if \(\operatorname{centerperiodic}(T, P)\), then for every natural number \(z\) and natural number \(k\), if \(T \le z + \operatorname{twopow}(k)\), then if \(\operatorname{zerodiagonal}(z)\), then if \(\operatorname{settledby}(z, z + \operatorname{twopow}(k))\), then \(\operatorname{maskparity}(z + \operatorname{twopow}(k), k) = \operatorname{bxor}(\operatorname{leftresponse}(z + \operatorname{twopow}(k), k), \operatorname{bxor}(\operatorname{center}(z + \operatorname{twopow}(k) + \operatorname{lag}(P, k)), \operatorname{center}(z + \operatorname{twopow}(k))))\).
\end{theorem}
\begin{proof}
Assume \(hP\), \(hper\).
The goal follows from \(\operatorname{periodicmaskparityatzerodiag}(T, P, hP, hper)\).
\end{proof}
\begin{theorem}
\label{ll:paper:periodic-maskparity-after-zero-diag}
For every natural number \(T\) and natural number \(P\), if \(0 < P\), then if \(\operatorname{centerperiodic}(T, P)\), then for every natural number \(z\) and natural number \(k\), if \(T \le z + 2 + \operatorname{twopow}(k)\), then if \(\operatorname{zerodiagonal}(z)\), then if \(\operatorname{settledby}(z + 2, z + 2 + \operatorname{twopow}(k))\), then \(\operatorname{maskparity}(z + 2 + \operatorname{twopow}(k), k) = \operatorname{not}(\operatorname{bxor}(\operatorname{leftresponse}(z + 2 + \operatorname{twopow}(k), k), \operatorname{bxor}(\operatorname{center}(z + 2 + \operatorname{twopow}(k) + \operatorname{lag}(P, k)), \operatorname{center}(z + 2 + \operatorname{twopow}(k)))))\).
\end{theorem}
\begin{proof}
Assume \(hP\), \(hper\).
The goal follows from \(\operatorname{periodicmaskparityafterzerodiag}(T, P, hP, hper)\).
\end{proof}
\begin{theorem}
\label{ll:paper:periodic-maskparity-defect-at-zero-diag-eq-lhcenter-twopow-xor-leftresponse-defect}
For every natural number \(T\) and natural number \(P\), if \(0 < P\), then if \(\operatorname{centerperiodic}(T, P)\), then for every natural number \(m\) and natural number \(z\) and natural number \(k\), if \(T \le z + \operatorname{twopow}(k)\), then if \(\operatorname{zerodiagonal}(z)\), then if \(\operatorname{settledby}(z, z + \operatorname{twopow}(k))\), then \(\operatorname{bxor}(\operatorname{maskparity}(z + \operatorname{twopow}(k) + m \cdot P, k), \operatorname{maskparity}(z + \operatorname{twopow}(k), k)) = \operatorname{bxor}(\operatorname{lhcenter}(z + \operatorname{twopow}(k) + m \cdot P, \operatorname{twopow}(k)), \operatorname{bxor}(\operatorname{leftresponse}(z + \operatorname{twopow}(k) + m \cdot P, k), \operatorname{leftresponse}(z + \operatorname{twopow}(k), k)))\).
\end{theorem}
\begin{proof}
Assume \(hP\), \(hper\).
The goal follows from \(\operatorname{periodicmaskparitydefectatzerodiageqlhcentertwopowxorleftresponsedefect}(T, P, hP, hper)\).
\end{proof}
\begin{theorem}
\label{ll:paper:periodic-maskparity-defect-after-zero-diag-eq-not-lhcenter-twopow-xor-leftresponse-defect}
For every natural number \(T\) and natural number \(P\), if \(0 < P\), then if \(\operatorname{centerperiodic}(T, P)\), then for every natural number \(m\) and natural number \(z\) and natural number \(k\), if \(T \le z + 2 + \operatorname{twopow}(k)\), then if \(\operatorname{zerodiagonal}(z)\), then if \(\operatorname{settledby}(z + 2, z + 2 + \operatorname{twopow}(k))\), then \(\operatorname{bxor}(\operatorname{maskparity}(z + 2 + \operatorname{twopow}(k) + m \cdot P, k), \operatorname{maskparity}(z + 2 + \operatorname{twopow}(k), k)) = \operatorname{not}(\operatorname{bxor}(\operatorname{lhcenter}(z + 2 + \operatorname{twopow}(k) + m \cdot P, \operatorname{twopow}(k)), \operatorname{bxor}(\operatorname{leftresponse}(z + 2 + \operatorname{twopow}(k) + m \cdot P, k), \operatorname{leftresponse}(z + 2 + \operatorname{twopow}(k), k))))\).
\end{theorem}
\begin{proof}
Assume \(hP\), \(hper\).
The goal follows from \(\operatorname{periodicmaskparitydefectafterzerodiageqnotlhcentertwopowxorleftresponsedefect}(T, P, hP, hper)\).
\end{proof}
\begin{theorem}
\label{ll:paper:periodic-response-flip-at-zero-diag-step-two}
For every natural number \(T\) and natural number \(P\), if \(0 < P\), then if \(\operatorname{centerperiodic}(T, P)\), then for every natural number \(m\) and natural number \(z\) and natural number \(k\), if \(m \cdot P = 2\), then if \(T \le z + \operatorname{twopow}(k)\), then if \(\operatorname{zerodiagonal}(z)\), then if \(\operatorname{settledby}(z, z + \operatorname{twopow}(k))\), then if \(\operatorname{settledby}(z + 2, z + 2 + \operatorname{twopow}(k))\), then \(\operatorname{response}(z + \operatorname{twopow}(k) + m \cdot P, k) = \operatorname{not}(\operatorname{response}(z + \operatorname{twopow}(k), k))\).
\end{theorem}
\begin{proof}
Assume \(hP\), \(hper\).
The goal follows from \(\operatorname{periodicresponseflipatzerodiagsteptwo}(T, P, hP, hper)\).
\end{proof}
\begin{theorem}
\label{ll:paper:periodic-response-flip-of-two-zero-diags}
For every natural number \(T\) and natural number \(P\), if \(0 < P\), then if \(\operatorname{centerperiodic}(T, P)\), then for every natural number \(m\) and natural number \(z1\) and natural number \(z2\) and natural number \(k\), if \(T \le z1 + \operatorname{twopow}(k)\), then if \(z1 + m \cdot P = z2 + 2\), then if \(\operatorname{zerodiagonal}(z1)\), then if \(\operatorname{zerodiagonal}(z2)\), then if \(\operatorname{settledby}(z1, z1 + \operatorname{twopow}(k))\), then if \(\operatorname{settledby}(z2 + 2, z2 + 2 + \operatorname{twopow}(k))\), then \(\operatorname{response}(z1 + \operatorname{twopow}(k) + m \cdot P, k) = \operatorname{not}(\operatorname{response}(z1 + \operatorname{twopow}(k), k))\).
\end{theorem}
\begin{proof}
Assume \(hP\), \(hper\).
The goal follows from \(\operatorname{periodicresponseflipoftwozerodiags}(T, P, hP, hper)\).
\end{proof}
\begin{theorem}
\label{ll:paper:periodic-response-eq-of-two-zero-diags}
For every natural number \(T\) and natural number \(P\), if \(0 < P\), then if \(\operatorname{centerperiodic}(T, P)\), then for every natural number \(m\) and natural number \(z1\) and natural number \(z2\) and natural number \(k\), if \(T \le z1 + \operatorname{twopow}(k)\), then if \(z1 + m \cdot P = z2\), then if \(\operatorname{zerodiagonal}(z1)\), then if \(\operatorname{zerodiagonal}(z2)\), then if \(\operatorname{settledby}(z1, z1 + \operatorname{twopow}(k))\), then if \(\operatorname{settledby}(z2, z2 + \operatorname{twopow}(k))\), then \(\operatorname{response}(z1 + \operatorname{twopow}(k) + m \cdot P, k) = \operatorname{response}(z1 + \operatorname{twopow}(k), k)\).
\end{theorem}
\begin{proof}
Assume \(hP\), \(hper\).
The goal follows from \(\operatorname{periodicresponseeqoftwozerodiags}(T, P, hP, hper)\).
\end{proof}
\begin{theorem}
\label{ll:paper:late-leftresponse-defect-eq-lhcenter-twopow-defect-cone}
For every natural number \(Q\) and natural number \(d\) and natural number \(k\), if \(d + Q < \operatorname{twopow}(k)\), then \(\operatorname{bxor}(\operatorname{leftresponse}(d + \operatorname{twopow}(k) + Q, k), \operatorname{leftresponse}(d + \operatorname{twopow}(k), k)) = \operatorname{bxor}(\operatorname{bxor}(\operatorname{lhcenter}(d + \operatorname{twopow}(k) + Q, \operatorname{twopow}(k)), \operatorname{lhcenter}(d + \operatorname{twopow}(k), \operatorname{twopow}(k))), \operatorname{bxor}(\operatorname{bxor}(\operatorname{leftresponse}(d + Q, k + 1), \operatorname{leftresponse}(d, k + 1)), \operatorname{bxor}(\operatorname{leftresponse}(d + Q, k), \operatorname{leftresponse}(d, k))))\).
\end{theorem}
\begin{proof}
The goal follows from \(\operatorname{lateleftresponsedefecteqlhcentertwopowdefectcone}(Q, d, k)\).
\end{proof}
\begin{theorem}
\label{ll:paper:maskcount-succ-eq-xor-maskparity-zero}
For every natural number \(t\) and natural number \(n\), \(\operatorname{maskcount}(t, n + 1) = \operatorname{bxor}(\operatorname{maskcount}(t, n), \operatorname{maskparity}(t + n, 0))\).
\end{theorem}
\begin{proof}
The goal follows from \(\operatorname{maskcountsucceqxormaskparityzero}(t, n)\).
\end{proof}
\begin{theorem}
\label{ll:paper:maskcount-eq-of-maskparity-zero-agree}
For every natural number \(t\) and natural number \(Q\) and natural number \(n\), if for every natural number \(s\), if \(s < n\), then \(\operatorname{maskparity}(t + s + Q, 0) = \operatorname{maskparity}(t + s, 0)\), then \(\operatorname{maskcount}(t + Q, n) = \operatorname{maskcount}(t, n)\).
\end{theorem}
\begin{proof}
The goal follows from \(\operatorname{maskcounteqofmaskparityzeroagree}(t, Q, n)\).
\end{proof}
\begin{theorem}
\label{ll:paper:maskparity-eq-of-maskparity-zero-agree}
For every natural number \(t\) and natural number \(Q\) and natural number \(k\) and natural number \(S\), if \(\operatorname{twopow}(k) \le S + 1\), then if for every natural number \(s\), if \(s \le S\), then \(\operatorname{maskparity}(t + s + Q, 0) = \operatorname{maskparity}(t + s, 0)\), then \(\operatorname{maskparity}(t + Q, k) = \operatorname{maskparity}(t, k)\).
\end{theorem}
\begin{proof}
The goal follows from \(\operatorname{maskparityeqofmaskparityzeroagree}(t, Q, k, S)\).
\end{proof}
\begin{theorem}
\label{ll:paper:periodic-maskparity-eq-of-white-center-rhcenter-one-twopow}
For every natural number \(T\) and natural number \(P\), if \(0 < P\), then if \(\operatorname{centerperiodic}(T, P)\), then for every natural number \(m\) and natural number \(t\) and natural number \(k\) and natural number \(S\), if \(T \le t\), then if \(\operatorname{twopow}(k) \le S + 1\), then if for every natural number \(s\), if \(s \le S\), then if \(\operatorname{center}(t + s) = \operatorname{false}\), then \(\operatorname{rhcenter}(t + s + m \cdot P, 1) = \operatorname{rhcenter}(t + s, 1)\), then \(\operatorname{maskparity}(t + m \cdot P, k) = \operatorname{maskparity}(t, k)\).
\end{theorem}
\begin{proof}
Assume \(hP\), \(hper\).
The goal follows from \(\operatorname{periodicmaskparityeqofwhitecenterrhcenteronetwopow}(T, P, hP, hper)\).
\end{proof}
\begin{theorem}
\label{ll:paper:periodic-leftresponse-eq-of-white-center-rhcenter-one-twopow}
For every natural number \(T\) and natural number \(P\), if \(0 < P\), then if \(\operatorname{centerperiodic}(T, P)\), then for every natural number \(m\) and natural number \(t\) and natural number \(k\) and natural number \(S\), if \(T \le t\), then if \(\operatorname{twopow}(k) \le S + 1\), then if for every natural number \(s\), if \(s \le S\), then if \(\operatorname{center}(t + s) = \operatorname{false}\), then \(\operatorname{rhcenter}(t + s + m \cdot P, 1) = \operatorname{rhcenter}(t + s, 1)\), then \(\operatorname{leftresponse}(t + m \cdot P, k) = \operatorname{leftresponse}(t, k)\).
\end{theorem}
\begin{proof}
Assume \(hP\), \(hper\).
The goal follows from \(\operatorname{periodicleftresponseeqofwhitecenterrhcenteronetwopow}(T, P, hP, hper)\).
\end{proof}
\begin{theorem}
\label{ll:paper:periodic-leftresponse-eq-of-maskparity-zero-agree-twopow}
For every natural number \(T\) and natural number \(P\), if \(0 < P\), then if \(\operatorname{centerperiodic}(T, P)\), then for every natural number \(m\) and natural number \(t\) and natural number \(k\) and natural number \(S\), if \(T \le t\), then if \(\operatorname{twopow}(k) \le S + 1\), then if for every natural number \(s\), if \(s \le S\), then \(\operatorname{maskparity}(t + s + m \cdot P, 0) = \operatorname{maskparity}(t + s, 0)\), then \(\operatorname{leftresponse}(t + m \cdot P, k) = \operatorname{leftresponse}(t, k)\).
\end{theorem}
\begin{proof}
Assume \(hP\), \(hper\).
The goal follows from \(\operatorname{periodicleftresponseeqofmaskparityzeroagreetwopow}(T, P, hP, hper)\).
\end{proof}
\begin{theorem}
\label{ll:paper:periodic-response-eq-of-maskparity-zero-agree-twopow}
For every natural number \(T\) and natural number \(P\), if \(0 < P\), then if \(\operatorname{centerperiodic}(T, P)\), then for every natural number \(m\) and natural number \(t\) and natural number \(k\) and natural number \(S\), if \(T \le t\), then if \(\operatorname{twopow}(k) \le S + 1\), then if for every natural number \(s\), if \(s \le S\), then \(\operatorname{maskparity}(t + s + m \cdot P, 0) = \operatorname{maskparity}(t + s, 0)\), then \(\operatorname{response}(t + m \cdot P, k) = \operatorname{response}(t, k)\).
\end{theorem}
\begin{proof}
Assume \(hP\), \(hper\).
The goal follows from \(\operatorname{periodicresponseeqofmaskparityzeroagreetwopow}(T, P, hP, hper)\).
\end{proof}
\begin{theorem}
\label{ll:paper:periodic-leftresponse-eq-of-lhcolagree-one-twopow}
For every natural number \(T\) and natural number \(P\), if \(0 < P\), then if \(\operatorname{centerperiodic}(T, P)\), then for every natural number \(m\) and natural number \(t\) and natural number \(k\) and natural number \(S\), if \(T \le t\), then if \(\operatorname{twopow}(k) \le S + 1\), then if \(\operatorname{lhcolagree}(t, m \cdot P, 1, S)\), then \(\operatorname{leftresponse}(t + m \cdot P, k) = \operatorname{leftresponse}(t, k)\).
\end{theorem}
\begin{proof}
Assume \(hP\), \(hper\).
The goal follows from \(\operatorname{periodicleftresponseeqoflhcolagreeonetwopow}(T, P, hP, hper)\).
\end{proof}
\begin{theorem}
\label{ll:paper:periodic-maskparity-eq-of-lhcolagree-one-twopow}
For every natural number \(T\) and natural number \(P\), if \(0 < P\), then if \(\operatorname{centerperiodic}(T, P)\), then for every natural number \(m\) and natural number \(t\) and natural number \(k\) and natural number \(S\), if \(T \le t\), then if \(\operatorname{twopow}(k) \le S + 1\), then if \(\operatorname{lhcolagree}(t, m \cdot P, 1, S)\), then \(\operatorname{maskparity}(t + m \cdot P, k) = \operatorname{maskparity}(t, k)\).
\end{theorem}
\begin{proof}
Assume \(hP\), \(hper\).
The goal follows from \(\operatorname{periodicmaskparityeqoflhcolagreeonetwopow}(T, P, hP, hper)\).
\end{proof}
\begin{theorem}
\label{ll:paper:periodic-response-eq-of-lhcolagree-one-twopow}
For every natural number \(T\) and natural number \(P\), if \(0 < P\), then if \(\operatorname{centerperiodic}(T, P)\), then for every natural number \(m\) and natural number \(t\) and natural number \(k\) and natural number \(S\), if \(T \le t\), then if \(\operatorname{twopow}(k) \le S + 1\), then if \(\operatorname{lhcolagree}(t, m \cdot P, 1, S)\), then \(\operatorname{response}(t + m \cdot P, k) = \operatorname{response}(t, k)\).
\end{theorem}
\begin{proof}
Assume \(hP\), \(hper\).
The goal follows from \(\operatorname{periodicresponseeqoflhcolagreeonetwopow}(T, P, hP, hper)\).
\end{proof}
\begin{theorem}
\label{ll:paper:lhcenter-twopow-eq-ldiag}
For every natural number \(d\) and natural number \(k\), \(\operatorname{lhcenter}(d + \operatorname{twopow}(k), \operatorname{twopow}(k)) = \operatorname{ldiag}(d, d + \operatorname{twopow}(k))\).
\end{theorem}
\begin{proof}
The goal follows from \(\operatorname{lhcentertwopoweqldiag}(d, k)\).
\end{proof}
\begin{theorem}
\label{ll:paper:lhcenter-twopow-settled-eq-zcenter}
For every natural number \(d\) and natural number \(k\), if \(d \le k\), then if \(\operatorname{settledby}(d, d + \operatorname{twopow}(k))\), then \(\operatorname{zcentereq}(d, \operatorname{lhcenter}(d + \operatorname{twopow}(k), \operatorname{twopow}(k)))\).
\end{theorem}
\begin{proof}
The goal follows from \(\operatorname{lhcentertwopowsettledeqzcenter}(d, k)\).
\end{proof}
\begin{theorem}
\label{ll:paper:lhcenter-twopow-settled-scale-invariant}
For every natural number \(d\) and natural number \(k\), if \(d \le k\), then if \(\operatorname{settledby}(d, d + \operatorname{twopow}(k))\), then \(\operatorname{lhcenter}(d + \operatorname{twopow}(k + 1), \operatorname{twopow}(k + 1)) = \operatorname{lhcenter}(d + \operatorname{twopow}(k), \operatorname{twopow}(k))\).
\end{theorem}
\begin{proof}
The goal follows from \(\operatorname{lhcentertwopowsettledscaleinvariant}(d, k)\).
\end{proof}
\begin{theorem}
\label{ll:paper:lhcenter-twopow-succ-periodic-of-doubles-le}
For every natural number \(z\) and natural number \(d\) and natural number \(k\), if \(\operatorname{doublesat}(\operatorname{twopow}(k), z)\), then if \(d \le z\), then if \(\operatorname{settledby}(d, d + \operatorname{twopow}(k))\), then \(\operatorname{lhcenter}(d + \operatorname{twopow}(k + 1), \operatorname{twopow}(k + 1)) = \operatorname{lhcenter}(d + \operatorname{twopow}(k), \operatorname{twopow}(k))\).
\end{theorem}
\begin{proof}
The goal follows from \(\operatorname{lhcentertwopowsuccperiodicofdoublesle}(z, d, k)\).
\end{proof}
\begin{theorem}
\label{ll:paper:lhcenter-twopow-reset-succ-antiperiodic-of-doubles}
For every natural number \(z\) and natural number \(k\), if \(\operatorname{zerodiagonal}(z)\), then if \(\operatorname{doublesat}(\operatorname{twopow}(k), z)\), then if \(\operatorname{settledby}(z + 3, z + 3 + \operatorname{twopow}(k))\), then \(\operatorname{lhcenter}(z + 3 + \operatorname{twopow}(k + 1), \operatorname{twopow}(k + 1)) = \operatorname{not}(\operatorname{lhcenter}(z + 3 + \operatorname{twopow}(k), \operatorname{twopow}(k)))\).
\end{theorem}
\begin{proof}
The goal follows from \(\operatorname{lhcentertwopowresetsuccantiperiodicofdoubles}(z, k)\).
\end{proof}
\begin{theorem}
\label{ll:paper:lhcenter-twopow-eq-late-scale-center}
For every natural number \(d\) and natural number \(k\), \(\operatorname{lhcenter}(d + \operatorname{twopow}(k), \operatorname{twopow}(k)) = \operatorname{bxor}(\operatorname{bxor}(\operatorname{center}(d + \operatorname{twopow}(k + 1)), \operatorname{center}(d + \operatorname{twopow}(k))), \operatorname{response}(d + \operatorname{twopow}(k), k))\).
\end{theorem}
\begin{proof}
The goal follows from \(\operatorname{lhcentertwopoweqlatescalecenter}(d, k)\).
\end{proof}
\begin{theorem}
\label{ll:paper:late-scale-center-periodic-of-doubles-le}
For every natural number \(z\) and natural number \(d\) and natural number \(k\), if \(\operatorname{doublesat}(\operatorname{twopow}(k), z)\), then if \(d \le z\), then if \(\operatorname{settledby}(d, d + \operatorname{twopow}(k))\), then \(\operatorname{bxor}(\operatorname{bxor}(\operatorname{center}(d + \operatorname{twopow}(k + 2)), \operatorname{center}(d + \operatorname{twopow}(k + 1))), \operatorname{response}(d + \operatorname{twopow}(k + 1), k + 1)) = \operatorname{bxor}(\operatorname{bxor}(\operatorname{center}(d + \operatorname{twopow}(k + 1)), \operatorname{center}(d + \operatorname{twopow}(k))), \operatorname{response}(d + \operatorname{twopow}(k), k))\).
\end{theorem}
\begin{proof}
The goal follows from \(\operatorname{latescalecenterperiodicofdoublesle}(z, d, k)\).
\end{proof}
\begin{theorem}
\label{ll:paper:late-scale-center-antiperiodic-of-doubles}
For every natural number \(z\) and natural number \(k\), if \(\operatorname{zerodiagonal}(z)\), then if \(\operatorname{doublesat}(\operatorname{twopow}(k), z)\), then if \(\operatorname{settledby}(z + 1, z + 1 + \operatorname{twopow}(k))\), then \(\operatorname{bxor}(\operatorname{bxor}(\operatorname{center}(z + 1 + \operatorname{twopow}(k + 2)), \operatorname{center}(z + 1 + \operatorname{twopow}(k + 1))), \operatorname{response}(z + 1 + \operatorname{twopow}(k + 1), k + 1)) = \operatorname{not}(\operatorname{bxor}(\operatorname{bxor}(\operatorname{center}(z + 1 + \operatorname{twopow}(k + 1)), \operatorname{center}(z + 1 + \operatorname{twopow}(k))), \operatorname{response}(z + 1 + \operatorname{twopow}(k), k)))\).
\end{theorem}
\begin{proof}
The goal follows from \(\operatorname{latescalecenterantiperiodicofdoubles}(z, k)\).
\end{proof}
\begin{theorem}
\label{ll:paper:late-scale-center-reset-antiperiodic-of-doubles}
For every natural number \(z\) and natural number \(k\), if \(\operatorname{zerodiagonal}(z)\), then if \(\operatorname{doublesat}(\operatorname{twopow}(k), z)\), then if \(\operatorname{settledby}(z + 3, z + 3 + \operatorname{twopow}(k))\), then \(\operatorname{bxor}(\operatorname{bxor}(\operatorname{center}(z + 3 + \operatorname{twopow}(k + 2)), \operatorname{center}(z + 3 + \operatorname{twopow}(k + 1))), \operatorname{response}(z + 3 + \operatorname{twopow}(k + 1), k + 1)) = \operatorname{not}(\operatorname{bxor}(\operatorname{bxor}(\operatorname{center}(z + 3 + \operatorname{twopow}(k + 1)), \operatorname{center}(z + 3 + \operatorname{twopow}(k))), \operatorname{response}(z + 3 + \operatorname{twopow}(k), k)))\).
\end{theorem}
\begin{proof}
The goal follows from \(\operatorname{latescalecenterresetantiperiodicofdoubles}(z, k)\).
\end{proof}
\begin{theorem}
\label{ll:paper:periodic-late-response-defect-eq-of-doubles-le}
For every natural number \(T\) and natural number \(P\), if \(0 < P\), then if \(\operatorname{centerperiodic}(T, P)\), then for every natural number \(m\) and natural number \(d\) and natural number \(z\) and natural number \(k\), if \(\operatorname{doublesat}(\operatorname{twopow}(k), z)\), then if \(d + m \cdot P \le z\), then if \(T \le d + \operatorname{twopow}(k)\), then if \(\operatorname{settledby}(d, d + \operatorname{twopow}(k))\), then if \(\operatorname{settledby}(d + m \cdot P, d + m \cdot P + \operatorname{twopow}(k))\), then \(\operatorname{bxor}(\operatorname{response}(d + \operatorname{twopow}(k + 1) + m \cdot P, k + 1), \operatorname{response}(d + \operatorname{twopow}(k + 1), k + 1)) = \operatorname{bxor}(\operatorname{response}(d + \operatorname{twopow}(k) + m \cdot P, k), \operatorname{response}(d + \operatorname{twopow}(k), k))\).
\end{theorem}
\begin{proof}
Assume \(hP\), \(hper\).
The goal follows from \(\operatorname{periodiclateresponsedefecteqofdoublesle}(T, P, hP, hper)\).
\end{proof}
\begin{theorem}
\label{ll:paper:periodic-late-response-defect-flip-of-doubles}
For every natural number \(T\) and natural number \(P\), if \(0 < P\), then if \(\operatorname{centerperiodic}(T, P)\), then for every natural number \(m\) and natural number \(d\) and natural number \(z\) and natural number \(k\), if \(\operatorname{zerodiagonal}(z)\), then if \(\operatorname{doublesat}(\operatorname{twopow}(k), z)\), then if \(d \le z\), then if \(d + m \cdot P = z + 1\), then if \(T \le d + \operatorname{twopow}(k)\), then if \(\operatorname{settledby}(d, d + \operatorname{twopow}(k))\), then if \(\operatorname{settledby}(z + 1, z + 1 + \operatorname{twopow}(k))\), then \(\operatorname{bxor}(\operatorname{response}(d + \operatorname{twopow}(k + 1) + m \cdot P, k + 1), \operatorname{response}(d + \operatorname{twopow}(k + 1), k + 1)) = \operatorname{not}(\operatorname{bxor}(\operatorname{response}(d + \operatorname{twopow}(k) + m \cdot P, k), \operatorname{response}(d + \operatorname{twopow}(k), k)))\).
\end{theorem}
\begin{proof}
Assume \(hP\), \(hper\).
The goal follows from \(\operatorname{periodiclateresponsedefectflipofdoubles}(T, P, hP, hper)\).
\end{proof}
\begin{theorem}
\label{ll:paper:periodic-late-response-flip-of-doubles-and-eq}
For every natural number \(T\) and natural number \(P\), if \(0 < P\), then if \(\operatorname{centerperiodic}(T, P)\), then for every natural number \(m\) and natural number \(d\) and natural number \(z\) and natural number \(k\), if \(\operatorname{zerodiagonal}(z)\), then if \(\operatorname{doublesat}(\operatorname{twopow}(k), z)\), then if \(d \le z\), then if \(d + m \cdot P = z + 1\), then if \(T \le d + \operatorname{twopow}(k)\), then if \(\operatorname{settledby}(d, d + \operatorname{twopow}(k))\), then if \(\operatorname{settledby}(z + 1, z + 1 + \operatorname{twopow}(k))\), then if \(\operatorname{response}(d + \operatorname{twopow}(k) + m \cdot P, k) = \operatorname{response}(d + \operatorname{twopow}(k), k)\), then \(\operatorname{response}(d + \operatorname{twopow}(k + 1) + m \cdot P, k + 1) = \operatorname{not}(\operatorname{response}(d + \operatorname{twopow}(k + 1), k + 1))\).
\end{theorem}
\begin{proof}
Assume \(hP\), \(hper\).
The goal follows from \(\operatorname{periodiclateresponseflipofdoublesandeq}(T, P, hP, hper)\).
\end{proof}
\begin{theorem}
\label{ll:paper:not-centerperiodic-of-white-center-rhcenter-one-at-twopow}
For every natural number \(T\) and natural number \(P\) and natural number \(m\) and natural number \(t\) and natural number \(k\) and natural number \(S\), if \(0 < P\), then if \(0 < m\), then if \(T \le t\), then if \(t + m \cdot P = \operatorname{twopow}(k)\), then if \(\operatorname{twopow}(k) \le S + 1\), then if for every natural number \(s\), if \(s \le S\), then if \(\operatorname{center}(t + s) = \operatorname{false}\), then \(\operatorname{rhcenter}(t + s + m \cdot P, 1) = \operatorname{rhcenter}(t + s, 1)\), then not \(\operatorname{centerperiodic}(T, P)\).
\end{theorem}
\begin{proof}
The goal follows from \(\operatorname{notcenterperiodicofwhitecenterrhcenteroneattwopow}(T, P, m, t, k, S)\).
\end{proof}
\begin{theorem}
\label{ll:paper:not-centerperiodic-of-maskparity-zero-agree-at-twopow}
For every natural number \(T\) and natural number \(P\) and natural number \(m\) and natural number \(t\) and natural number \(k\) and natural number \(S\), if \(0 < P\), then if \(0 < m\), then if \(T \le t\), then if \(t + m \cdot P = \operatorname{twopow}(k)\), then if \(\operatorname{twopow}(k) \le S + 1\), then if for every natural number \(s\), if \(s \le S\), then \(\operatorname{maskparity}(t + s + m \cdot P, 0) = \operatorname{maskparity}(t + s, 0)\), then not \(\operatorname{centerperiodic}(T, P)\).
\end{theorem}
\begin{proof}
The goal follows from \(\operatorname{notcenterperiodicofmaskparityzeroagreeattwopow}(T, P, m, t, k, S)\).
\end{proof}
\begin{theorem}
\label{ll:paper:not-centerperiodic-of-white-center-rhcenter-one-at-zero-diag-step-two}
For every natural number \(T\) and natural number \(P\) and natural number \(m\) and natural number \(z\) and natural number \(k\) and natural number \(S\), if \(0 < P\), then if \(\operatorname{zerodiagonal}(z)\), then if \(m \cdot P = 2\), then if \(T \le z + \operatorname{twopow}(k)\), then if \(\operatorname{settledby}(z, z + \operatorname{twopow}(k))\), then if \(\operatorname{settledby}(z + 2, z + 2 + \operatorname{twopow}(k))\), then if \(\operatorname{twopow}(k) \le S + 1\), then if for every natural number \(s\), if \(s \le S\), then if \(\operatorname{center}(z + \operatorname{twopow}(k) + s) = \operatorname{false}\), then \(\operatorname{rhcenter}(z + \operatorname{twopow}(k) + s + m \cdot P, 1) = \operatorname{rhcenter}(z + \operatorname{twopow}(k) + s, 1)\), then not \(\operatorname{centerperiodic}(T, P)\).
\end{theorem}
\begin{proof}
The goal follows from \(\operatorname{notcenterperiodicofwhitecenterrhcenteroneatzerodiagsteptwo}(T, P, m, z, k, S)\).
\end{proof}
\begin{theorem}
\label{ll:paper:not-centerperiodic-of-maskparity-zero-agree-at-zero-diag-step-two}
For every natural number \(T\) and natural number \(P\) and natural number \(m\) and natural number \(z\) and natural number \(k\) and natural number \(S\), if \(0 < P\), then if \(\operatorname{zerodiagonal}(z)\), then if \(m \cdot P = 2\), then if \(T \le z + \operatorname{twopow}(k)\), then if \(\operatorname{settledby}(z, z + \operatorname{twopow}(k))\), then if \(\operatorname{settledby}(z + 2, z + 2 + \operatorname{twopow}(k))\), then if \(\operatorname{twopow}(k) \le S + 1\), then if for every natural number \(s\), if \(s \le S\), then \(\operatorname{maskparity}(z + \operatorname{twopow}(k) + s + m \cdot P, 0) = \operatorname{maskparity}(z + \operatorname{twopow}(k) + s, 0)\), then not \(\operatorname{centerperiodic}(T, P)\).
\end{theorem}
\begin{proof}
The goal follows from \(\operatorname{notcenterperiodicofmaskparityzeroagreeatzerodiagsteptwo}(T, P, m, z, k, S)\).
\end{proof}
\begin{theorem}
\label{ll:paper:not-centerperiodic-of-white-center-rhcenter-one-of-two-zero-diags}
For every natural number \(T\) and natural number \(P\) and natural number \(m\) and natural number \(z\) and natural number \(k\) and natural number \(S\), if \(0 < P\), then if \(\operatorname{zerodiagonal}(z)\), then if \(2 \le z + m \cdot P\), then if \(\operatorname{zerodiagonal}(z + m \cdot P - 2)\), then if \(T \le z + \operatorname{twopow}(k)\), then if \(\operatorname{settledby}(z, z + \operatorname{twopow}(k))\), then if \(\operatorname{settledby}(z + m \cdot P, z + m \cdot P + \operatorname{twopow}(k))\), then if \(\operatorname{twopow}(k) \le S + 1\), then if for every natural number \(s\), if \(s \le S\), then if \(\operatorname{center}(z + \operatorname{twopow}(k) + s) = \operatorname{false}\), then \(\operatorname{rhcenter}(z + \operatorname{twopow}(k) + s + m \cdot P, 1) = \operatorname{rhcenter}(z + \operatorname{twopow}(k) + s, 1)\), then not \(\operatorname{centerperiodic}(T, P)\).
\end{theorem}
\begin{proof}
The goal follows from \(\operatorname{notcenterperiodicofwhitecenterrhcenteroneoftwozerodiags}(T, P, m, z, k, S)\).
\end{proof}
\begin{theorem}
\label{ll:paper:not-centerperiodic-of-maskparity-zero-agree-of-two-zero-diags}
For every natural number \(T\) and natural number \(P\) and natural number \(m\) and natural number \(z\) and natural number \(k\) and natural number \(S\), if \(0 < P\), then if \(\operatorname{zerodiagonal}(z)\), then if \(2 \le z + m \cdot P\), then if \(\operatorname{zerodiagonal}(z + m \cdot P - 2)\), then if \(T \le z + \operatorname{twopow}(k)\), then if \(\operatorname{settledby}(z, z + \operatorname{twopow}(k))\), then if \(\operatorname{settledby}(z + m \cdot P, z + m \cdot P + \operatorname{twopow}(k))\), then if \(\operatorname{twopow}(k) \le S + 1\), then if for every natural number \(s\), if \(s \le S\), then \(\operatorname{maskparity}(z + \operatorname{twopow}(k) + s + m \cdot P, 0) = \operatorname{maskparity}(z + \operatorname{twopow}(k) + s, 0)\), then not \(\operatorname{centerperiodic}(T, P)\).
\end{theorem}
\begin{proof}
The goal follows from \(\operatorname{notcenterperiodicofmaskparityzeroagreeoftwozerodiags}(T, P, m, z, k, S)\).
\end{proof}
\begin{theorem}
\label{ll:paper:not-centerperiodic-of-white-center-rhcenter-one-scale-step-doubles}
For every natural number \(T\) and natural number \(P\) and natural number \(m\) and natural number \(d\) and natural number \(z\) and natural number \(k\) and natural number \(S0\) and natural number \(S1\), if \(0 < P\), then if \(\operatorname{zerodiagonal}(z)\), then if \(\operatorname{doublesat}(\operatorname{twopow}(k), z)\), then if \(d \le z\), then if \(d + m \cdot P = z + 1\), then if \(T \le d + \operatorname{twopow}(k)\), then if \(\operatorname{settledby}(d, d + \operatorname{twopow}(k))\), then if \(\operatorname{settledby}(z + 1, z + 1 + \operatorname{twopow}(k))\), then if \(\operatorname{twopow}(k) \le S0 + 1\), then if \(\operatorname{twopow}(k + 1) \le S1 + 1\), then if for every natural number \(s\), if \(s \le S0\), then if \(\operatorname{center}(d + \operatorname{twopow}(k) + s) = \operatorname{false}\), then \(\operatorname{rhcenter}(d + \operatorname{twopow}(k) + s + m \cdot P, 1) = \operatorname{rhcenter}(d + \operatorname{twopow}(k) + s, 1)\), then if for every natural number \(s\), if \(s \le S1\), then if \(\operatorname{center}(d + \operatorname{twopow}(k + 1) + s) = \operatorname{false}\), then \(\operatorname{rhcenter}(d + \operatorname{twopow}(k + 1) + s + m \cdot P, 1) = \operatorname{rhcenter}(d + \operatorname{twopow}(k + 1) + s, 1)\), then not \(\operatorname{centerperiodic}(T, P)\).
\end{theorem}
\begin{proof}
The goal follows from \(\operatorname{notcenterperiodicofwhitecenterrhcenteronescalestepdoubles}(T, P, m, d, z, k, S0, S1)\).
\end{proof}
\begin{theorem}
\label{ll:paper:not-centerperiodic-of-maskparity-zero-agree-scale-step-doubles}
For every natural number \(T\) and natural number \(P\) and natural number \(m\) and natural number \(d\) and natural number \(z\) and natural number \(k\) and natural number \(S0\) and natural number \(S1\), if \(0 < P\), then if \(\operatorname{zerodiagonal}(z)\), then if \(\operatorname{doublesat}(\operatorname{twopow}(k), z)\), then if \(d \le z\), then if \(d + m \cdot P = z + 1\), then if \(T \le d + \operatorname{twopow}(k)\), then if \(\operatorname{settledby}(d, d + \operatorname{twopow}(k))\), then if \(\operatorname{settledby}(z + 1, z + 1 + \operatorname{twopow}(k))\), then if \(\operatorname{twopow}(k) \le S0 + 1\), then if \(\operatorname{twopow}(k + 1) \le S1 + 1\), then if for every natural number \(s\), if \(s \le S0\), then \(\operatorname{maskparity}(d + \operatorname{twopow}(k) + s + m \cdot P, 0) = \operatorname{maskparity}(d + \operatorname{twopow}(k) + s, 0)\), then if for every natural number \(s\), if \(s \le S1\), then \(\operatorname{maskparity}(d + \operatorname{twopow}(k + 1) + s + m \cdot P, 0) = \operatorname{maskparity}(d + \operatorname{twopow}(k + 1) + s, 0)\), then not \(\operatorname{centerperiodic}(T, P)\).
\end{theorem}
\begin{proof}
The goal follows from \(\operatorname{notcenterperiodicofmaskparityzeroagreescalestepdoubles}(T, P, m, d, z, k, S0, S1)\).
\end{proof}
\begin{theorem}
\label{ll:paper:lt-twopow-self}
For every natural number \(n\), \(n < \operatorname{twopow}(n)\).
\end{theorem}
\begin{proof}
The goal follows from \(\operatorname{lttwopowself}(n)\).
\end{proof}
\begin{theorem}
\label{ll:paper:not-centerperiodic-of-lhcolagree-one-at-twopow}
For every natural number \(T\) and natural number \(P\) and natural number \(m\) and natural number \(t\) and natural number \(k\) and natural number \(S\), if \(0 < P\), then if \(0 < m\), then if \(T \le t\), then if \(t + m \cdot P = \operatorname{twopow}(k)\), then if \(\operatorname{twopow}(k) \le S + 1\), then if \(\operatorname{lhcolagree}(t, m \cdot P, 1, S)\), then not \(\operatorname{centerperiodic}(T, P)\).
\end{theorem}
\begin{proof}
The goal follows from \(\operatorname{notcenterperiodicoflhcolagreeoneattwopow}(T, P, m, t, k, S)\).
\end{proof}
\begin{theorem}
\label{ll:paper:not-centerperiodic-of-lhcolagree-one-at-zero-diag-step-two}
For every natural number \(T\) and natural number \(P\) and natural number \(m\) and natural number \(z\) and natural number \(k\) and natural number \(S\), if \(0 < P\), then if \(\operatorname{zerodiagonal}(z)\), then if \(m \cdot P = 2\), then if \(T \le z + \operatorname{twopow}(k)\), then if \(\operatorname{settledby}(z, z + \operatorname{twopow}(k))\), then if \(\operatorname{settledby}(z + 2, z + 2 + \operatorname{twopow}(k))\), then if \(\operatorname{twopow}(k) \le S + 1\), then if \(\operatorname{lhcolagree}(z + \operatorname{twopow}(k), m \cdot P, 1, S)\), then not \(\operatorname{centerperiodic}(T, P)\).
\end{theorem}
\begin{proof}
The goal follows from \(\operatorname{notcenterperiodicoflhcolagreeoneatzerodiagsteptwo}(T, P, m, z, k, S)\).
\end{proof}
\begin{theorem}
\label{ll:paper:not-centerperiodic-of-lhcolagree-one-of-two-zero-diags}
For every natural number \(T\) and natural number \(P\) and natural number \(m\) and natural number \(z\) and natural number \(k\) and natural number \(S\), if \(0 < P\), then if \(\operatorname{zerodiagonal}(z)\), then if \(2 \le z + m \cdot P\), then if \(\operatorname{zerodiagonal}(z + m \cdot P - 2)\), then if \(T \le z + \operatorname{twopow}(k)\), then if \(\operatorname{settledby}(z, z + \operatorname{twopow}(k))\), then if \(\operatorname{settledby}(z + m \cdot P, z + m \cdot P + \operatorname{twopow}(k))\), then if \(\operatorname{twopow}(k) \le S + 1\), then if \(\operatorname{lhcolagree}(z + \operatorname{twopow}(k), m \cdot P, 1, S)\), then not \(\operatorname{centerperiodic}(T, P)\).
\end{theorem}
\begin{proof}
The goal follows from \(\operatorname{notcenterperiodicoflhcolagreeoneoftwozerodiags}(T, P, m, z, k, S)\).
\end{proof}
\begin{theorem}
\label{ll:paper:not-centerperiodic-of-lhcolagree-one-scale-step-doubles}
For every natural number \(T\) and natural number \(P\) and natural number \(m\) and natural number \(d\) and natural number \(z\) and natural number \(k\) and natural number \(S0\) and natural number \(S1\), if \(0 < P\), then if \(\operatorname{zerodiagonal}(z)\), then if \(\operatorname{doublesat}(\operatorname{twopow}(k), z)\), then if \(d \le z\), then if \(d + m \cdot P = z + 1\), then if \(T \le d + \operatorname{twopow}(k)\), then if \(\operatorname{settledby}(d, d + \operatorname{twopow}(k))\), then if \(\operatorname{settledby}(z + 1, z + 1 + \operatorname{twopow}(k))\), then if \(\operatorname{twopow}(k) \le S0 + 1\), then if \(\operatorname{twopow}(k + 1) \le S1 + 1\), then if \(\operatorname{lhcolagree}(d + \operatorname{twopow}(k), m \cdot P, 1, S0)\), then if \(\operatorname{lhcolagree}(d + \operatorname{twopow}(k + 1), m \cdot P, 1, S1)\), then not \(\operatorname{centerperiodic}(T, P)\).
\end{theorem}
\begin{proof}
The goal follows from \(\operatorname{notcenterperiodicoflhcolagreeonescalestepdoubles}(T, P, m, d, z, k, S0, S1)\).
\end{proof}
\begin{theorem}
\label{ll:paper:not-centerperiodic-of-white-center-rhcenter-one-sub-twopow}
For every natural number \(T\) and natural number \(P\) and natural number \(m\) and natural number \(k\) and natural number \(S\), if \(0 < P\), then if \(0 < m\), then if \(T + m \cdot P \le \operatorname{twopow}(k)\), then if \(\operatorname{twopow}(k) \le S + 1\), then if for every natural number \(s\), if \(s \le S\), then if \(\operatorname{center}(\operatorname{twopow}(k) - m \cdot P + s) = \operatorname{false}\), then \(\operatorname{rhcenter}(\operatorname{twopow}(k) - m \cdot P + s + m \cdot P, 1) = \operatorname{rhcenter}(\operatorname{twopow}(k) - m \cdot P + s, 1)\), then not \(\operatorname{centerperiodic}(T, P)\).
\end{theorem}
\begin{proof}
The goal follows from \(\operatorname{notcenterperiodicofwhitecenterrhcenteronesubtwopow}(T, P, m, k, S)\).
\end{proof}
\begin{theorem}
\label{ll:paper:not-centerperiodic-of-maskparity-zero-agree-sub-twopow}
For every natural number \(T\) and natural number \(P\) and natural number \(m\) and natural number \(k\) and natural number \(S\), if \(0 < P\), then if \(0 < m\), then if \(T + m \cdot P \le \operatorname{twopow}(k)\), then if \(\operatorname{twopow}(k) \le S + 1\), then if for every natural number \(s\), if \(s \le S\), then \(\operatorname{maskparity}(\operatorname{twopow}(k) - m \cdot P + s + m \cdot P, 0) = \operatorname{maskparity}(\operatorname{twopow}(k) - m \cdot P + s, 0)\), then not \(\operatorname{centerperiodic}(T, P)\).
\end{theorem}
\begin{proof}
The goal follows from \(\operatorname{notcenterperiodicofmaskparityzeroagreesubtwopow}(T, P, m, k, S)\).
\end{proof}
\begin{theorem}
\label{ll:paper:exists-white-center-rhcenter-one-flip-at-twopow-of-periodic}
For every natural number \(T\) and natural number \(P\), if \(0 < P\), then if \(\operatorname{centerperiodic}(T, P)\), then for every natural number \(m\) and natural number \(t\) and natural number \(k\) and natural number \(S\), if \(0 < m\), then if \(T \le t\), then if \(t + m \cdot P = \operatorname{twopow}(k)\), then if \(\operatorname{twopow}(k) \le S + 1\), then there exists a natural number \(s\) such that \(s \le S\) and \(\operatorname{center}(t + s) = \operatorname{false} \land \operatorname{rhcenter}(t + s + m \cdot P, 1) = \operatorname{not}(\operatorname{rhcenter}(t + s, 1))\).
\end{theorem}
\begin{proof}
Assume \(hP\), \(hper\).
The goal follows from \(\operatorname{existswhitecenterrhcenteroneflipattwopowofperiodic}(T, P, hP, hper)\).
\end{proof}
\begin{theorem}
\label{ll:paper:exists-maskparity-zero-flip-at-twopow-of-periodic}
For every natural number \(T\) and natural number \(P\), if \(0 < P\), then if \(\operatorname{centerperiodic}(T, P)\), then for every natural number \(m\) and natural number \(t\) and natural number \(k\) and natural number \(S\), if \(0 < m\), then if \(T \le t\), then if \(t + m \cdot P = \operatorname{twopow}(k)\), then if \(\operatorname{twopow}(k) \le S + 1\), then there exists a natural number \(s\) such that \(s \le S\) and \(\operatorname{maskparity}(t + s + m \cdot P, 0) = \operatorname{not}(\operatorname{maskparity}(t + s, 0))\).
\end{theorem}
\begin{proof}
Assume \(hP\), \(hper\).
The goal follows from \(\operatorname{existsmaskparityzeroflipattwopowofperiodic}(T, P, hP, hper)\).
\end{proof}
\begin{theorem}
\label{ll:paper:exists-lhcenter-one-flip-at-twopow-of-periodic}
For every natural number \(T\) and natural number \(P\), if \(0 < P\), then if \(\operatorname{centerperiodic}(T, P)\), then for every natural number \(m\) and natural number \(t\) and natural number \(k\) and natural number \(S\), if \(0 < m\), then if \(T \le t\), then if \(t + m \cdot P = \operatorname{twopow}(k)\), then if \(\operatorname{twopow}(k) \le S + 1\), then there exists a natural number \(s\) such that \(s \le S\) and \(\operatorname{lhcenter}(t + s + m \cdot P, 1) = \operatorname{not}(\operatorname{lhcenter}(t + s, 1))\).
\end{theorem}
\begin{proof}
Assume \(hP\), \(hper\).
The goal follows from \(\operatorname{existslhcenteroneflipattwopowofperiodic}(T, P, hP, hper)\).
\end{proof}
\begin{theorem}
\label{ll:paper:exists-white-center-rhcenter-one-flip-sub-twopow-of-periodic}
For every natural number \(T\) and natural number \(P\), if \(0 < P\), then if \(\operatorname{centerperiodic}(T, P)\), then for every natural number \(m\) and natural number \(k\), if \(0 < m\), then if \(T + m \cdot P \le \operatorname{twopow}(k)\), then there exists a natural number \(s\) such that \(s < \operatorname{twopow}(k)\) and \(\operatorname{center}(\operatorname{twopow}(k) - m \cdot P + s) = \operatorname{false} \land \operatorname{rhcenter}(\operatorname{twopow}(k) - m \cdot P + s + m \cdot P, 1) = \operatorname{not}(\operatorname{rhcenter}(\operatorname{twopow}(k) - m \cdot P + s, 1))\).
\end{theorem}
\begin{proof}
Assume \(hP\), \(hper\).
The goal follows from \(\operatorname{existswhitecenterrhcenteroneflipsubtwopowofperiodic}(T, P, hP, hper)\).
\end{proof}
\begin{theorem}
\label{ll:paper:exists-maskparity-zero-flip-sub-twopow-of-periodic}
For every natural number \(T\) and natural number \(P\), if \(0 < P\), then if \(\operatorname{centerperiodic}(T, P)\), then for every natural number \(m\) and natural number \(k\), if \(0 < m\), then if \(T + m \cdot P \le \operatorname{twopow}(k)\), then there exists a natural number \(s\) such that \(s < \operatorname{twopow}(k)\) and \(\operatorname{maskparity}(\operatorname{twopow}(k) - m \cdot P + s + m \cdot P, 0) = \operatorname{not}(\operatorname{maskparity}(\operatorname{twopow}(k) - m \cdot P + s, 0))\).
\end{theorem}
\begin{proof}
Assume \(hP\), \(hper\).
The goal follows from \(\operatorname{existsmaskparityzeroflipsubtwopowofperiodic}(T, P, hP, hper)\).
\end{proof}
\begin{theorem}
\label{ll:paper:exists-lhcenter-one-flip-sub-twopow-of-periodic}
For every natural number \(T\) and natural number \(P\), if \(0 < P\), then if \(\operatorname{centerperiodic}(T, P)\), then for every natural number \(m\) and natural number \(k\), if \(0 < m\), then if \(T + m \cdot P \le \operatorname{twopow}(k)\), then there exists a natural number \(s\) such that \(s < \operatorname{twopow}(k)\) and \(\operatorname{lhcenter}(\operatorname{twopow}(k) - m \cdot P + s + m \cdot P, 1) = \operatorname{not}(\operatorname{lhcenter}(\operatorname{twopow}(k) - m \cdot P + s, 1))\).
\end{theorem}
\begin{proof}
Assume \(hP\), \(hper\).
The goal follows from \(\operatorname{existslhcenteroneflipsubtwopowofperiodic}(T, P, hP, hper)\).
\end{proof}
\begin{theorem}
\label{ll:paper:exists-white-center-rhcenter-one-flip-at-zero-diag-step-two-of-periodic}
For every natural number \(T\) and natural number \(P\), if \(0 < P\), then if \(\operatorname{centerperiodic}(T, P)\), then for every natural number \(m\) and natural number \(z\) and natural number \(k\) and natural number \(S\), if \(\operatorname{zerodiagonal}(z)\), then if \(m \cdot P = 2\), then if \(T \le z + \operatorname{twopow}(k)\), then if \(\operatorname{settledby}(z, z + \operatorname{twopow}(k))\), then if \(\operatorname{settledby}(z + 2, z + 2 + \operatorname{twopow}(k))\), then if \(\operatorname{twopow}(k) \le S + 1\), then there exists a natural number \(s\) such that \(s \le S\) and \(\operatorname{center}(z + \operatorname{twopow}(k) + s) = \operatorname{false} \land \operatorname{rhcenter}(z + \operatorname{twopow}(k) + s + m \cdot P, 1) = \operatorname{not}(\operatorname{rhcenter}(z + \operatorname{twopow}(k) + s, 1))\).
\end{theorem}
\begin{proof}
Assume \(hP\), \(hper\).
The goal follows from \(\operatorname{existswhitecenterrhcenteroneflipatzerodiagsteptwoofperiodic}(T, P, hP, hper)\).
\end{proof}
\begin{theorem}
\label{ll:paper:exists-maskparity-zero-flip-at-zero-diag-step-two-of-periodic}
For every natural number \(T\) and natural number \(P\), if \(0 < P\), then if \(\operatorname{centerperiodic}(T, P)\), then for every natural number \(m\) and natural number \(z\) and natural number \(k\) and natural number \(S\), if \(\operatorname{zerodiagonal}(z)\), then if \(m \cdot P = 2\), then if \(T \le z + \operatorname{twopow}(k)\), then if \(\operatorname{settledby}(z, z + \operatorname{twopow}(k))\), then if \(\operatorname{settledby}(z + 2, z + 2 + \operatorname{twopow}(k))\), then if \(\operatorname{twopow}(k) \le S + 1\), then there exists a natural number \(s\) such that \(s \le S\) and \(\operatorname{maskparity}(z + \operatorname{twopow}(k) + s + m \cdot P, 0) = \operatorname{not}(\operatorname{maskparity}(z + \operatorname{twopow}(k) + s, 0))\).
\end{theorem}
\begin{proof}
Assume \(hP\), \(hper\).
The goal follows from \(\operatorname{existsmaskparityzeroflipatzerodiagsteptwoofperiodic}(T, P, hP, hper)\).
\end{proof}
\begin{theorem}
\label{ll:paper:twopow-mono}
For every natural number \(a\) and natural number \(b\), if \(a \le b\), then \(\operatorname{twopow}(a) \le \operatorname{twopow}(b)\).
\end{theorem}
\begin{proof}
The goal follows from \(\operatorname{twopowmono}(a, b)\).
\end{proof}
\begin{theorem}
\label{ll:paper:lt-twopow-of-le}
For every natural number \(n\) and natural number \(k\), if \(n \le k\), then \(n < \operatorname{twopow}(k)\).
\end{theorem}
\begin{proof}
The goal follows from \(\operatorname{lttwopowofle}(n, k)\).
\end{proof}
\begin{theorem}
\label{ll:paper:not-centerperiodic-of-white-center-rhcenter-one-sub-twopow-ge}
For every natural number \(T\) and natural number \(P\) and natural number \(m\) and natural number \(k\) and natural number \(S\), if \(0 < P\), then if \(0 < m\), then if \(T + m \cdot P \le k\), then if \(\operatorname{twopow}(k) \le S + 1\), then if for every natural number \(s\), if \(s \le S\), then if \(\operatorname{center}(\operatorname{twopow}(k) - m \cdot P + s) = \operatorname{false}\), then \(\operatorname{rhcenter}(\operatorname{twopow}(k) - m \cdot P + s + m \cdot P, 1) = \operatorname{rhcenter}(\operatorname{twopow}(k) - m \cdot P + s, 1)\), then not \(\operatorname{centerperiodic}(T, P)\).
\end{theorem}
\begin{proof}
The goal follows from \(\operatorname{notcenterperiodicofwhitecenterrhcenteronesubtwopowge}(T, P, m, k, S)\).
\end{proof}
\begin{theorem}
\label{ll:paper:not-centerperiodic-of-maskparity-zero-agree-sub-twopow-ge}
For every natural number \(T\) and natural number \(P\) and natural number \(m\) and natural number \(k\) and natural number \(S\), if \(0 < P\), then if \(0 < m\), then if \(T + m \cdot P \le k\), then if \(\operatorname{twopow}(k) \le S + 1\), then if for every natural number \(s\), if \(s \le S\), then \(\operatorname{maskparity}(\operatorname{twopow}(k) - m \cdot P + s + m \cdot P, 0) = \operatorname{maskparity}(\operatorname{twopow}(k) - m \cdot P + s, 0)\), then not \(\operatorname{centerperiodic}(T, P)\).
\end{theorem}
\begin{proof}
The goal follows from \(\operatorname{notcenterperiodicofmaskparityzeroagreesubtwopowge}(T, P, m, k, S)\).
\end{proof}
\begin{theorem}
\label{ll:paper:exists-white-center-rhcenter-one-flip-sub-twopow-ge-of-periodic}
For every natural number \(T\) and natural number \(P\), if \(0 < P\), then if \(\operatorname{centerperiodic}(T, P)\), then for every natural number \(m\) and natural number \(k\), if \(0 < m\), then if \(T + m \cdot P \le k\), then there exists a natural number \(s\) such that \(s < \operatorname{twopow}(k)\) and \(\operatorname{center}(\operatorname{twopow}(k) - m \cdot P + s) = \operatorname{false} \land \operatorname{rhcenter}(\operatorname{twopow}(k) - m \cdot P + s + m \cdot P, 1) = \operatorname{not}(\operatorname{rhcenter}(\operatorname{twopow}(k) - m \cdot P + s, 1))\).
\end{theorem}
\begin{proof}
Assume \(hP\), \(hper\).
The goal follows from \(\operatorname{existswhitecenterrhcenteroneflipsubtwopowgeofperiodic}(T, P, hP, hper)\).
\end{proof}
\begin{theorem}
\label{ll:paper:exists-maskparity-zero-flip-sub-twopow-ge-of-periodic}
For every natural number \(T\) and natural number \(P\), if \(0 < P\), then if \(\operatorname{centerperiodic}(T, P)\), then for every natural number \(m\) and natural number \(k\), if \(0 < m\), then if \(T + m \cdot P \le k\), then there exists a natural number \(s\) such that \(s < \operatorname{twopow}(k)\) and \(\operatorname{maskparity}(\operatorname{twopow}(k) - m \cdot P + s + m \cdot P, 0) = \operatorname{not}(\operatorname{maskparity}(\operatorname{twopow}(k) - m \cdot P + s, 0))\).
\end{theorem}
\begin{proof}
Assume \(hP\), \(hper\).
The goal follows from \(\operatorname{existsmaskparityzeroflipsubtwopowgeofperiodic}(T, P, hP, hper)\).
\end{proof}
\begin{theorem}
\label{ll:paper:exists-lhcenter-one-flip-sub-twopow-ge-of-periodic}
For every natural number \(T\) and natural number \(P\), if \(0 < P\), then if \(\operatorname{centerperiodic}(T, P)\), then for every natural number \(m\) and natural number \(k\), if \(0 < m\), then if \(T + m \cdot P \le k\), then there exists a natural number \(s\) such that \(s < \operatorname{twopow}(k)\) and \(\operatorname{lhcenter}(\operatorname{twopow}(k) - m \cdot P + s + m \cdot P, 1) = \operatorname{not}(\operatorname{lhcenter}(\operatorname{twopow}(k) - m \cdot P + s, 1))\).
\end{theorem}
\begin{proof}
Assume \(hP\), \(hper\).
The goal follows from \(\operatorname{existslhcenteroneflipsubtwopowgeofperiodic}(T, P, hP, hper)\).
\end{proof}
\begin{theorem}
\label{ll:paper:exists-white-center-rhcenter-one-flip-sub-twopow-p-of-periodic}
For every natural number \(T\) and natural number \(P\), if \(0 < P\), then if \(\operatorname{centerperiodic}(T, P)\), then for every natural number \(k\), if \(T + P \le k\), then there exists a natural number \(s\) such that \(s < \operatorname{twopow}(k)\) and \(\operatorname{center}(\operatorname{twopow}(k) - P + s) = \operatorname{false} \land \operatorname{rhcenter}(\operatorname{twopow}(k) - P + s + P, 1) = \operatorname{not}(\operatorname{rhcenter}(\operatorname{twopow}(k) - P + s, 1))\).
\end{theorem}
\begin{proof}
Assume \(hP\), \(hper\).
The goal follows from \(\operatorname{existswhitecenterrhcenteroneflipsubtwopowpofperiodic}(T, P, hP, hper)\).
\end{proof}
\begin{theorem}
\label{ll:paper:exists-maskparity-zero-flip-sub-twopow-p-of-periodic}
For every natural number \(T\) and natural number \(P\), if \(0 < P\), then if \(\operatorname{centerperiodic}(T, P)\), then for every natural number \(k\), if \(T + P \le k\), then there exists a natural number \(s\) such that \(s < \operatorname{twopow}(k)\) and \(\operatorname{maskparity}(\operatorname{twopow}(k) - P + s + P, 0) = \operatorname{not}(\operatorname{maskparity}(\operatorname{twopow}(k) - P + s, 0))\).
\end{theorem}
\begin{proof}
Assume \(hP\), \(hper\).
The goal follows from \(\operatorname{existsmaskparityzeroflipsubtwopowpofperiodic}(T, P, hP, hper)\).
\end{proof}
\begin{theorem}
\label{ll:paper:exists-lhcenter-one-flip-sub-twopow-p-of-periodic}
For every natural number \(T\) and natural number \(P\), if \(0 < P\), then if \(\operatorname{centerperiodic}(T, P)\), then for every natural number \(k\), if \(T + P \le k\), then there exists a natural number \(s\) such that \(s < \operatorname{twopow}(k)\) and \(\operatorname{lhcenter}(\operatorname{twopow}(k) - P + s + P, 1) = \operatorname{not}(\operatorname{lhcenter}(\operatorname{twopow}(k) - P + s, 1))\).
\end{theorem}
\begin{proof}
Assume \(hP\), \(hper\).
The goal follows from \(\operatorname{existslhcenteroneflipsubtwopowpofperiodic}(T, P, hP, hper)\).
\end{proof}
\begin{theorem}
\label{ll:paper:exists-lhcenter-one-flip-at-zero-diag-step-two-of-periodic}
For every natural number \(T\) and natural number \(P\), if \(0 < P\), then if \(\operatorname{centerperiodic}(T, P)\), then for every natural number \(m\) and natural number \(z\) and natural number \(k\) and natural number \(S\), if \(\operatorname{zerodiagonal}(z)\), then if \(m \cdot P = 2\), then if \(T \le z + \operatorname{twopow}(k)\), then if \(\operatorname{settledby}(z, z + \operatorname{twopow}(k))\), then if \(\operatorname{settledby}(z + 2, z + 2 + \operatorname{twopow}(k))\), then if \(\operatorname{twopow}(k) \le S + 1\), then there exists a natural number \(s\) such that \(s \le S\) and \(\operatorname{lhcenter}(z + \operatorname{twopow}(k) + s + m \cdot P, 1) = \operatorname{not}(\operatorname{lhcenter}(z + \operatorname{twopow}(k) + s, 1))\).
\end{theorem}
\begin{proof}
Assume \(hP\), \(hper\).
The goal follows from \(\operatorname{existslhcenteroneflipatzerodiagsteptwoofperiodic}(T, P, hP, hper)\).
\end{proof}
\begin{theorem}
\label{ll:paper:exists-white-center-rhcenter-one-flip-of-two-zero-diags-and-periodic}
For every natural number \(T\) and natural number \(P\), if \(0 < P\), then if \(\operatorname{centerperiodic}(T, P)\), then for every natural number \(m\) and natural number \(z\) and natural number \(k\) and natural number \(S\), if \(\operatorname{zerodiagonal}(z)\), then if \(2 \le z + m \cdot P\), then if \(\operatorname{zerodiagonal}(z + m \cdot P - 2)\), then if \(T \le z + \operatorname{twopow}(k)\), then if \(\operatorname{settledby}(z, z + \operatorname{twopow}(k))\), then if \(\operatorname{settledby}(z + m \cdot P, z + m \cdot P + \operatorname{twopow}(k))\), then if \(\operatorname{twopow}(k) \le S + 1\), then there exists a natural number \(s\) such that \(s \le S\) and \(\operatorname{center}(z + \operatorname{twopow}(k) + s) = \operatorname{false} \land \operatorname{rhcenter}(z + \operatorname{twopow}(k) + s + m \cdot P, 1) = \operatorname{not}(\operatorname{rhcenter}(z + \operatorname{twopow}(k) + s, 1))\).
\end{theorem}
\begin{proof}
Assume \(hP\), \(hper\).
The goal follows from \(\operatorname{existswhitecenterrhcenteroneflipoftwozerodiagsandperiodic}(T, P, hP, hper)\).
\end{proof}
\begin{theorem}
\label{ll:paper:exists-maskparity-zero-flip-of-two-zero-diags-and-periodic}
For every natural number \(T\) and natural number \(P\), if \(0 < P\), then if \(\operatorname{centerperiodic}(T, P)\), then for every natural number \(m\) and natural number \(z\) and natural number \(k\) and natural number \(S\), if \(\operatorname{zerodiagonal}(z)\), then if \(2 \le z + m \cdot P\), then if \(\operatorname{zerodiagonal}(z + m \cdot P - 2)\), then if \(T \le z + \operatorname{twopow}(k)\), then if \(\operatorname{settledby}(z, z + \operatorname{twopow}(k))\), then if \(\operatorname{settledby}(z + m \cdot P, z + m \cdot P + \operatorname{twopow}(k))\), then if \(\operatorname{twopow}(k) \le S + 1\), then there exists a natural number \(s\) such that \(s \le S\) and \(\operatorname{maskparity}(z + \operatorname{twopow}(k) + s + m \cdot P, 0) = \operatorname{not}(\operatorname{maskparity}(z + \operatorname{twopow}(k) + s, 0))\).
\end{theorem}
\begin{proof}
Assume \(hP\), \(hper\).
The goal follows from \(\operatorname{existsmaskparityzeroflipoftwozerodiagsandperiodic}(T, P, hP, hper)\).
\end{proof}
\begin{theorem}
\label{ll:paper:exists-lhcenter-one-flip-of-two-zero-diags-and-periodic}
For every natural number \(T\) and natural number \(P\), if \(0 < P\), then if \(\operatorname{centerperiodic}(T, P)\), then for every natural number \(m\) and natural number \(z\) and natural number \(k\) and natural number \(S\), if \(\operatorname{zerodiagonal}(z)\), then if \(2 \le z + m \cdot P\), then if \(\operatorname{zerodiagonal}(z + m \cdot P - 2)\), then if \(T \le z + \operatorname{twopow}(k)\), then if \(\operatorname{settledby}(z, z + \operatorname{twopow}(k))\), then if \(\operatorname{settledby}(z + m \cdot P, z + m \cdot P + \operatorname{twopow}(k))\), then if \(\operatorname{twopow}(k) \le S + 1\), then there exists a natural number \(s\) such that \(s \le S\) and \(\operatorname{lhcenter}(z + \operatorname{twopow}(k) + s + m \cdot P, 1) = \operatorname{not}(\operatorname{lhcenter}(z + \operatorname{twopow}(k) + s, 1))\).
\end{theorem}
\begin{proof}
Assume \(hP\), \(hper\).
The goal follows from \(\operatorname{existslhcenteroneflipoftwozerodiagsandperiodic}(T, P, hP, hper)\).
\end{proof}
\begin{theorem}
\label{ll:paper:periodic-late-maskparity-xor-leftresponse-defect-flip-of-doubles}
For every natural number \(T\) and natural number \(P\), if \(0 < P\), then if \(\operatorname{centerperiodic}(T, P)\), then for every natural number \(m\) and natural number \(d\) and natural number \(z\) and natural number \(k\), if \(\operatorname{zerodiagonal}(z)\), then if \(\operatorname{doublesat}(\operatorname{twopow}(k), z)\), then if \(d \le z\), then if \(d + m \cdot P = z + 1\), then if \(T \le d + \operatorname{twopow}(k)\), then if \(\operatorname{settledby}(d, d + \operatorname{twopow}(k))\), then if \(\operatorname{settledby}(z + 1, z + 1 + \operatorname{twopow}(k))\), then \(\operatorname{bxor}(\operatorname{bxor}(\operatorname{maskparity}(d + \operatorname{twopow}(k + 1) + m \cdot P, k + 1), \operatorname{maskparity}(d + \operatorname{twopow}(k + 1), k + 1)), \operatorname{bxor}(\operatorname{leftresponse}(d + \operatorname{twopow}(k + 1) + m \cdot P, k + 1), \operatorname{leftresponse}(d + \operatorname{twopow}(k + 1), k + 1))) = \operatorname{not}(\operatorname{bxor}(\operatorname{bxor}(\operatorname{maskparity}(d + \operatorname{twopow}(k) + m \cdot P, k), \operatorname{maskparity}(d + \operatorname{twopow}(k), k)), \operatorname{bxor}(\operatorname{leftresponse}(d + \operatorname{twopow}(k) + m \cdot P, k), \operatorname{leftresponse}(d + \operatorname{twopow}(k), k))))\).
\end{theorem}
\begin{proof}
Assume \(hP\), \(hper\).
The goal follows from \(\operatorname{periodiclatemaskparityxorleftresponsedefectflipofdoubles}(T, P, hP, hper)\).
\end{proof}
\begin{theorem}
\label{ll:paper:periodic-late-response-reset-defect-flip-of-doubles}
For every natural number \(T\) and natural number \(P\), if \(0 < P\), then if \(\operatorname{centerperiodic}(T, P)\), then for every natural number \(m\) and natural number \(d\) and natural number \(z\) and natural number \(k\), if \(\operatorname{zerodiagonal}(z)\), then if \(\operatorname{doublesat}(\operatorname{twopow}(k), z)\), then if \(d \le z\), then if \(d + m \cdot P = z + 3\), then if \(T \le d + \operatorname{twopow}(k)\), then if \(\operatorname{settledby}(d, d + \operatorname{twopow}(k))\), then if \(\operatorname{settledby}(z + 3, z + 3 + \operatorname{twopow}(k))\), then \(\operatorname{bxor}(\operatorname{response}(d + \operatorname{twopow}(k + 1) + m \cdot P, k + 1), \operatorname{response}(d + \operatorname{twopow}(k + 1), k + 1)) = \operatorname{not}(\operatorname{bxor}(\operatorname{response}(d + \operatorname{twopow}(k) + m \cdot P, k), \operatorname{response}(d + \operatorname{twopow}(k), k)))\).
\end{theorem}
\begin{proof}
Assume \(hP\), \(hper\).
The goal follows from \(\operatorname{periodiclateresponseresetdefectflipofdoubles}(T, P, hP, hper)\).
\end{proof}
\begin{theorem}
\label{ll:paper:periodic-late-response-reset-flip-of-doubles-and-eq}
For every natural number \(T\) and natural number \(P\), if \(0 < P\), then if \(\operatorname{centerperiodic}(T, P)\), then for every natural number \(m\) and natural number \(d\) and natural number \(z\) and natural number \(k\), if \(\operatorname{zerodiagonal}(z)\), then if \(\operatorname{doublesat}(\operatorname{twopow}(k), z)\), then if \(d \le z\), then if \(d + m \cdot P = z + 3\), then if \(T \le d + \operatorname{twopow}(k)\), then if \(\operatorname{settledby}(d, d + \operatorname{twopow}(k))\), then if \(\operatorname{settledby}(z + 3, z + 3 + \operatorname{twopow}(k))\), then if \(\operatorname{response}(d + \operatorname{twopow}(k) + m \cdot P, k) = \operatorname{response}(d + \operatorname{twopow}(k), k)\), then \(\operatorname{response}(d + \operatorname{twopow}(k + 1) + m \cdot P, k + 1) = \operatorname{not}(\operatorname{response}(d + \operatorname{twopow}(k + 1), k + 1))\).
\end{theorem}
\begin{proof}
Assume \(hP\), \(hper\).
The goal follows from \(\operatorname{periodiclateresponseresetflipofdoublesandeq}(T, P, hP, hper)\).
\end{proof}
\begin{theorem}
\label{ll:paper:periodic-late-maskparity-xor-leftresponse-reset-defect-flip-of-doubles}
For every natural number \(T\) and natural number \(P\), if \(0 < P\), then if \(\operatorname{centerperiodic}(T, P)\), then for every natural number \(m\) and natural number \(d\) and natural number \(z\) and natural number \(k\), if \(\operatorname{zerodiagonal}(z)\), then if \(\operatorname{doublesat}(\operatorname{twopow}(k), z)\), then if \(d \le z\), then if \(d + m \cdot P = z + 3\), then if \(T \le d + \operatorname{twopow}(k)\), then if \(\operatorname{settledby}(d, d + \operatorname{twopow}(k))\), then if \(\operatorname{settledby}(z + 3, z + 3 + \operatorname{twopow}(k))\), then \(\operatorname{bxor}(\operatorname{bxor}(\operatorname{maskparity}(d + \operatorname{twopow}(k + 1) + m \cdot P, k + 1), \operatorname{maskparity}(d + \operatorname{twopow}(k + 1), k + 1)), \operatorname{bxor}(\operatorname{leftresponse}(d + \operatorname{twopow}(k + 1) + m \cdot P, k + 1), \operatorname{leftresponse}(d + \operatorname{twopow}(k + 1), k + 1))) = \operatorname{not}(\operatorname{bxor}(\operatorname{bxor}(\operatorname{maskparity}(d + \operatorname{twopow}(k) + m \cdot P, k), \operatorname{maskparity}(d + \operatorname{twopow}(k), k)), \operatorname{bxor}(\operatorname{leftresponse}(d + \operatorname{twopow}(k) + m \cdot P, k), \operatorname{leftresponse}(d + \operatorname{twopow}(k), k))))\).
\end{theorem}
\begin{proof}
Assume \(hP\), \(hper\).
The goal follows from \(\operatorname{periodiclatemaskparityxorleftresponseresetdefectflipofdoubles}(T, P, hP, hper)\).
\end{proof}
\begin{theorem}
\label{ll:paper:not-centerperiodic-of-white-center-rhcenter-one-three-twopow-doubles}
For every natural number \(T\) and natural number \(P\) and natural number \(m\) and natural number \(d\) and natural number \(z\) and natural number \(k\) and natural number \(S\), if \(0 < P\), then if \(\operatorname{zerodiagonal}(z)\), then if \(\operatorname{doublesat}(\operatorname{twopow}(k), z)\), then if \(d \le z\), then if \(d + m \cdot P = z + 1\), then if \(T \le d + \operatorname{twopow}(k)\), then if \(\operatorname{settledby}(d, d + \operatorname{twopow}(k))\), then if \(\operatorname{settledby}(z + 1, z + 1 + \operatorname{twopow}(k))\), then if \(\operatorname{twopow}(k) + \operatorname{twopow}(k + 1) \le S + 1\), then if for every natural number \(s\), if \(s \le S\), then if \(\operatorname{center}(d + \operatorname{twopow}(k) + s) = \operatorname{false}\), then \(\operatorname{rhcenter}(d + \operatorname{twopow}(k) + s + m \cdot P, 1) = \operatorname{rhcenter}(d + \operatorname{twopow}(k) + s, 1)\), then not \(\operatorname{centerperiodic}(T, P)\).
\end{theorem}
\begin{proof}
The goal follows from \(\operatorname{notcenterperiodicofwhitecenterrhcenteronethreetwopowdoubles}(T, P, m, d, z, k, S)\).
\end{proof}
\begin{theorem}
\label{ll:paper:not-centerperiodic-of-maskparity-zero-agree-three-twopow-doubles}
For every natural number \(T\) and natural number \(P\) and natural number \(m\) and natural number \(d\) and natural number \(z\) and natural number \(k\) and natural number \(S\), if \(0 < P\), then if \(\operatorname{zerodiagonal}(z)\), then if \(\operatorname{doublesat}(\operatorname{twopow}(k), z)\), then if \(d \le z\), then if \(d + m \cdot P = z + 1\), then if \(T \le d + \operatorname{twopow}(k)\), then if \(\operatorname{settledby}(d, d + \operatorname{twopow}(k))\), then if \(\operatorname{settledby}(z + 1, z + 1 + \operatorname{twopow}(k))\), then if \(\operatorname{twopow}(k) + \operatorname{twopow}(k + 1) \le S + 1\), then if for every natural number \(s\), if \(s \le S\), then \(\operatorname{maskparity}(d + \operatorname{twopow}(k) + s + m \cdot P, 0) = \operatorname{maskparity}(d + \operatorname{twopow}(k) + s, 0)\), then not \(\operatorname{centerperiodic}(T, P)\).
\end{theorem}
\begin{proof}
The goal follows from \(\operatorname{notcenterperiodicofmaskparityzeroagreethreetwopowdoubles}(T, P, m, d, z, k, S)\).
\end{proof}
\begin{theorem}
\label{ll:paper:not-centerperiodic-of-lhcolagree-one-three-twopow-doubles}
For every natural number \(T\) and natural number \(P\) and natural number \(m\) and natural number \(d\) and natural number \(z\) and natural number \(k\) and natural number \(S\), if \(0 < P\), then if \(\operatorname{zerodiagonal}(z)\), then if \(\operatorname{doublesat}(\operatorname{twopow}(k), z)\), then if \(d \le z\), then if \(d + m \cdot P = z + 1\), then if \(T \le d + \operatorname{twopow}(k)\), then if \(\operatorname{settledby}(d, d + \operatorname{twopow}(k))\), then if \(\operatorname{settledby}(z + 1, z + 1 + \operatorname{twopow}(k))\), then if \(\operatorname{twopow}(k) + \operatorname{twopow}(k + 1) \le S + 1\), then if \(\operatorname{lhcolagree}(d + \operatorname{twopow}(k), m \cdot P, 1, S)\), then not \(\operatorname{centerperiodic}(T, P)\).
\end{theorem}
\begin{proof}
The goal follows from \(\operatorname{notcenterperiodicoflhcolagreeonethreetwopowdoubles}(T, P, m, d, z, k, S)\).
\end{proof}
\begin{theorem}
\label{ll:paper:not-centerperiodic-of-white-center-rhcenter-one-three-twopow-reset-doubles}
For every natural number \(T\) and natural number \(P\) and natural number \(m\) and natural number \(d\) and natural number \(z\) and natural number \(k\) and natural number \(S\), if \(0 < P\), then if \(\operatorname{zerodiagonal}(z)\), then if \(\operatorname{doublesat}(\operatorname{twopow}(k), z)\), then if \(d \le z\), then if \(d + m \cdot P = z + 3\), then if \(T \le d + \operatorname{twopow}(k)\), then if \(\operatorname{settledby}(d, d + \operatorname{twopow}(k))\), then if \(\operatorname{settledby}(z + 3, z + 3 + \operatorname{twopow}(k))\), then if \(\operatorname{twopow}(k) + \operatorname{twopow}(k + 1) \le S + 1\), then if for every natural number \(s\), if \(s \le S\), then if \(\operatorname{center}(d + \operatorname{twopow}(k) + s) = \operatorname{false}\), then \(\operatorname{rhcenter}(d + \operatorname{twopow}(k) + s + m \cdot P, 1) = \operatorname{rhcenter}(d + \operatorname{twopow}(k) + s, 1)\), then not \(\operatorname{centerperiodic}(T, P)\).
\end{theorem}
\begin{proof}
The goal follows from \(\operatorname{notcenterperiodicofwhitecenterrhcenteronethreetwopowresetdoubles}(T, P, m, d, z, k, S)\).
\end{proof}
\begin{theorem}
\label{ll:paper:not-centerperiodic-of-maskparity-zero-agree-three-twopow-reset-doubles}
For every natural number \(T\) and natural number \(P\) and natural number \(m\) and natural number \(d\) and natural number \(z\) and natural number \(k\) and natural number \(S\), if \(0 < P\), then if \(\operatorname{zerodiagonal}(z)\), then if \(\operatorname{doublesat}(\operatorname{twopow}(k), z)\), then if \(d \le z\), then if \(d + m \cdot P = z + 3\), then if \(T \le d + \operatorname{twopow}(k)\), then if \(\operatorname{settledby}(d, d + \operatorname{twopow}(k))\), then if \(\operatorname{settledby}(z + 3, z + 3 + \operatorname{twopow}(k))\), then if \(\operatorname{twopow}(k) + \operatorname{twopow}(k + 1) \le S + 1\), then if for every natural number \(s\), if \(s \le S\), then \(\operatorname{maskparity}(d + \operatorname{twopow}(k) + s + m \cdot P, 0) = \operatorname{maskparity}(d + \operatorname{twopow}(k) + s, 0)\), then not \(\operatorname{centerperiodic}(T, P)\).
\end{theorem}
\begin{proof}
The goal follows from \(\operatorname{notcenterperiodicofmaskparityzeroagreethreetwopowresetdoubles}(T, P, m, d, z, k, S)\).
\end{proof}
\begin{theorem}
\label{ll:paper:not-centerperiodic-of-lhcolagree-one-three-twopow-reset-doubles}
For every natural number \(T\) and natural number \(P\) and natural number \(m\) and natural number \(d\) and natural number \(z\) and natural number \(k\) and natural number \(S\), if \(0 < P\), then if \(\operatorname{zerodiagonal}(z)\), then if \(\operatorname{doublesat}(\operatorname{twopow}(k), z)\), then if \(d \le z\), then if \(d + m \cdot P = z + 3\), then if \(T \le d + \operatorname{twopow}(k)\), then if \(\operatorname{settledby}(d, d + \operatorname{twopow}(k))\), then if \(\operatorname{settledby}(z + 3, z + 3 + \operatorname{twopow}(k))\), then if \(\operatorname{twopow}(k) + \operatorname{twopow}(k + 1) \le S + 1\), then if \(\operatorname{lhcolagree}(d + \operatorname{twopow}(k), m \cdot P, 1, S)\), then not \(\operatorname{centerperiodic}(T, P)\).
\end{theorem}
\begin{proof}
The goal follows from \(\operatorname{notcenterperiodicoflhcolagreeonethreetwopowresetdoubles}(T, P, m, d, z, k, S)\).
\end{proof}
\begin{theorem}
\label{ll:paper:rdrv-diag-eq-rhcenter}
For every natural number \(t\) and natural number \(k\) and natural number \(p\), \(\operatorname{rdrv}(t, t, k, p) = \operatorname{rhcenter}(t + k, p)\).
\end{theorem}
\begin{proof}
The goal follows from \(\operatorname{rdrvdiageqrhcenter}(t, k, p)\).
\end{proof}
\begin{theorem}
\label{ll:paper:rdrv-zero-eq-rhcenter}
For every natural number \(t\) and natural number \(p\), \(\operatorname{rdrv}(0, 0, t, p) = \operatorname{rhcenter}(t, p)\).
\end{theorem}
\begin{proof}
The goal follows from \(\operatorname{rdrvzeroeqrhcenter}(t, p)\).
\end{proof}
\begin{theorem}
\label{ll:paper:rdrv-step-succ}
For every natural number \(td\) and natural number \(tx\) and natural number \(k\) and natural number \(p\), \(\operatorname{rdrv}(td, tx, k + 1, p + 1) = \operatorname{bxor}(\operatorname{rdrv}(td, tx, k, p), \operatorname{rdrv}(td, tx, k, p + 1) \operatorname{or} \operatorname{rdrv}(td, tx, k, p + 2))\).
\end{theorem}
\begin{proof}
The goal follows from \(\operatorname{rdrvstepsucc}(td, tx, k, p)\).
\end{proof}
\begin{theorem}
\label{ll:paper:rdrv-drive-indep-of-lt}
For every natural number \(td\) and natural number \(td2\) and natural number \(tx\) and natural number \(k\) and natural number \(p\), if \(k < p\), then \(\operatorname{rdrv}(td, tx, k, p) = \operatorname{rdrv}(td2, tx, k, p)\).
\end{theorem}
\begin{proof}
The goal follows from \(\operatorname{rdrvdriveindepoflt}(td, td2, tx, k, p)\).
\end{proof}
\begin{theorem}
\label{ll:paper:periodic-rdrv-drive-shift-eq}
For every natural number \(T\) and natural number \(P\), if \(0 < P\), then if \(\operatorname{centerperiodic}(T, P)\), then for every natural number \(m\) and natural number \(td\) and natural number \(tx\) and natural number \(k\) and natural number \(p\), if \(T \le td\), then \(\operatorname{rdrv}(td + m \cdot P, tx, k, p) = \operatorname{rdrv}(td, tx, k, p)\).
\end{theorem}
\begin{proof}
Assume \(hP\), \(hper\).
The goal follows from \(\operatorname{periodicrdrvdriveshifteq}(T, P, hP, hper)\).
\end{proof}
\begin{theorem}
\label{ll:paper:periodic-rhcenter-shift-eq-rdrv-same-drive}
For every natural number \(T\) and natural number \(P\), if \(0 < P\), then if \(\operatorname{centerperiodic}(T, P)\), then for every natural number \(m\) and natural number \(t\) and natural number \(k\) and natural number \(p\), if \(T \le t\), then \(\operatorname{rhcenter}(t + m \cdot P + k, p) = \operatorname{rdrv}(t, t + m \cdot P, k, p)\).
\end{theorem}
\begin{proof}
Assume \(hP\), \(hper\).
The goal follows from \(\operatorname{periodicrhcentershifteqrdrvsamedrive}(T, P, hP, hper)\).
\end{proof}
\begin{theorem}
\label{ll:paper:rdrv-outside-cone-eq-rhcenter}
For every natural number \(td\) and natural number \(tx\) and natural number \(k\) and natural number \(p\), if \(k < p\), then \(\operatorname{rdrv}(td, tx, k, p) = \operatorname{rhcenter}(tx + k, p)\).
\end{theorem}
\begin{proof}
The goal follows from \(\operatorname{rdrvoutsideconeeqrhcenter}(td, tx, k, p)\).
\end{proof}
\begin{theorem}
\label{ll:paper:rdrv-semigroup-add}
For every natural number \(t\) and natural number \(a\) and natural number \(k\) and natural number \(p\), \(\operatorname{rdrv}(t, t, a + k, p) = \operatorname{rdrv}(t + a, t + a, k, p)\).
\end{theorem}
\begin{proof}
The goal follows from \(\operatorname{rdrvsemigroupadd}(t, a, k, p)\).
\end{proof}
\begin{theorem}
\label{ll:paper:rdrv-drive-causal-cone}
For every natural number \(td\) and natural number \(td2\) and natural number \(tx\) and natural number \(k\) and natural number \(p\), if for every natural number \(s\), if \(s + p \le k\), then \(\operatorname{center}(td + s) = \operatorname{center}(td2 + s)\), then \(\operatorname{rdrv}(td, tx, k, p) = \operatorname{rdrv}(td2, tx, k, p)\).
\end{theorem}
\begin{proof}
The goal follows from \(\operatorname{rdrvdrivecausalcone}(td, td2, tx, k, p)\).
\end{proof}
\begin{theorem}
\label{ll:paper:rdrv-init-causal-cone}
For every natural number \(td\) and natural number \(tx\) and natural number \(Q\) and natural number \(L\) and natural number \(k\) and natural number \(p\), if \(\operatorname{rhbandagree}(tx, Q, 1, L)\), then if \(p + k \le L + 1\), then \(\operatorname{rdrv}(td, tx + Q, k, p) = \operatorname{rdrv}(td, tx, k, p)\).
\end{theorem}
\begin{proof}
The goal follows from \(\operatorname{rdrvinitcausalcone}(td, tx, Q, L, k, p)\).
\end{proof}
\begin{theorem}
\label{ll:paper:rdrv-col-one-agree-of-rhbandagree-one}
For every natural number \(td\) and natural number \(tx\) and natural number \(Q\) and natural number \(S\) and natural number \(s\), if \(\operatorname{rhbandagree}(tx, Q, 1, S)\), then if \(s \le S\), then \(\operatorname{rdrv}(td, tx + Q, s, 1) = \operatorname{rdrv}(td, tx, s, 1)\).
\end{theorem}
\begin{proof}
The goal follows from \(\operatorname{rdrvcoloneagreeofrhbandagreeone}(td, tx, Q, S, s)\).
\end{proof}
\begin{theorem}
\label{ll:paper:periodic-rhcolagree-one-of-rhbandagree-one}
For every natural number \(T\) and natural number \(P\), if \(0 < P\), then if \(\operatorname{centerperiodic}(T, P)\), then for every natural number \(m\) and natural number \(t\) and natural number \(S\), if \(T \le t\), then if \(\operatorname{rhbandagree}(t, m \cdot P, 1, S)\), then \(\operatorname{rhcolagree}(t, m \cdot P, 1, S)\).
\end{theorem}
\begin{proof}
Assume \(hP\), \(hper\).
The goal follows from \(\operatorname{periodicrhcolagreeoneofrhbandagreeone}(T, P, hP, hper)\).
\end{proof}
\begin{theorem}
\label{ll:paper:periodic-maskparity-zero-agree-of-rhbandagree-one}
For every natural number \(T\) and natural number \(P\), if \(0 < P\), then if \(\operatorname{centerperiodic}(T, P)\), then for every natural number \(m\) and natural number \(t\) and natural number \(S\) and natural number \(s\), if \(T \le t\), then if \(\operatorname{rhbandagree}(t, m \cdot P, 1, S)\), then if \(s \le S\), then \(\operatorname{maskparity}(t + s + m \cdot P, 0) = \operatorname{maskparity}(t + s, 0)\).
\end{theorem}
\begin{proof}
Assume \(hP\), \(hper\).
The goal follows from \(\operatorname{periodicmaskparityzeroagreeofrhbandagreeone}(T, P, hP, hper)\).
\end{proof}
\begin{theorem}
\label{ll:paper:periodic-maskparity-zero-agree-of-rdrv-one-agree}
For every natural number \(T\) and natural number \(P\), if \(0 < P\), then if \(\operatorname{centerperiodic}(T, P)\), then for every natural number \(m\) and natural number \(t\) and natural number \(s\), if \(T \le t\), then if if \(\operatorname{center}(t + s) = \operatorname{false}\), then \(\operatorname{rdrv}(t, t + m \cdot P, s, 1) = \operatorname{rdrv}(t, t, s, 1)\), then \(\operatorname{maskparity}(t + s + m \cdot P, 0) = \operatorname{maskparity}(t + s, 0)\).
\end{theorem}
\begin{proof}
Assume \(hP\), \(hper\).
The goal follows from \(\operatorname{periodicmaskparityzeroagreeofrdrvoneagree}(T, P, hP, hper)\).
\end{proof}
\begin{theorem}
\label{ll:paper:rdrv-strip-agree-of-rhbandagree-and-col-agree}
For every natural number \(td\) and natural number \(tx\) and natural number \(Q\) and natural number \(B\) and natural number \(S\) and natural number \(s\) and natural number \(p\), if \(\operatorname{rhbandagree}(tx, Q, 1, B)\), then if for every natural number \(u\), if \(u < S\), then \(\operatorname{rdrv}(td, tx + Q, u, B + 2) = \operatorname{rdrv}(td, tx, u, B + 2)\), then if \(s \le S\), then if \(p \le B + 1\), then \(\operatorname{rdrv}(td, tx + Q, s, p) = \operatorname{rdrv}(td, tx, s, p)\).
\end{theorem}
\begin{proof}
The goal follows from \(\operatorname{rdrvstripagreeofrhbandagreeandcolagree}(td, tx, Q, B, S, s, p)\).
\end{proof}
\begin{theorem}
\label{ll:paper:periodic-maskparity-zero-agree-of-rhbandagree-and-rhcolagree}
For every natural number \(T\) and natural number \(P\), if \(0 < P\), then if \(\operatorname{centerperiodic}(T, P)\), then for every natural number \(m\) and natural number \(t\) and natural number \(B\) and natural number \(S\) and natural number \(s\), if \(T \le t\), then if \(\operatorname{rhbandagree}(t, m \cdot P, 1, B)\), then if \(\operatorname{rhcolagree}(t, m \cdot P, B + 2, S)\), then if \(s \le S\), then \(\operatorname{maskparity}(t + s + m \cdot P, 0) = \operatorname{maskparity}(t + s, 0)\).
\end{theorem}
\begin{proof}
Assume \(hP\), \(hper\).
The goal follows from \(\operatorname{periodicmaskparityzeroagreeofrhbandagreeandrhcolagree}(T, P, hP, hper)\).
\end{proof}
\begin{theorem}
\label{ll:paper:not-centerperiodic-of-rdrv-one-agree-at-zero-diag-step-two}
For every natural number \(T\) and natural number \(P\) and natural number \(m\) and natural number \(z\) and natural number \(p\) and natural number \(S\), if \(0 < P\), then if \(0 < m\), then if \(T \le z + (p + 1)\), then if \(\operatorname{zerodiagonal}(z)\), then if \(\operatorname{doublesat}(m \cdot P, z)\), then if \(\operatorname{settledby}(z + 1, z + (p + 1))\), then if \(\operatorname{settledby}(z + 1 - m \cdot P, z + (p + 1))\), then if \(p + 3 \cdot (m \cdot P) \le S + 1\), then if for every natural number \(s\), if \(s \le S\), then if \(\operatorname{center}(z + (p + 1) + s) = \operatorname{false}\), then \(\operatorname{rdrv}(z + (p + 1), z + (p + 1) + m \cdot P, s, 1) = \operatorname{rdrv}(z + (p + 1), z + (p + 1), s, 1)\), then not \(\operatorname{centerperiodic}(T, P)\).
\end{theorem}
\begin{proof}
The goal follows from \(\operatorname{notcenterperiodicofrdrvoneagreeatzerodiagsteptwo}(T, P, m, z, p, S)\).
\end{proof}
\begin{theorem}
\label{ll:paper:not-centerperiodic-of-rhbandagree-one-at-zero-diag-step-two}
For every natural number \(T\) and natural number \(P\) and natural number \(m\) and natural number \(z\) and natural number \(p\) and natural number \(S\), if \(0 < P\), then if \(0 < m\), then if \(T \le z + (p + 1)\), then if \(\operatorname{zerodiagonal}(z)\), then if \(\operatorname{doublesat}(m \cdot P, z)\), then if \(\operatorname{settledby}(z + 1, z + (p + 1))\), then if \(\operatorname{settledby}(z + 1 - m \cdot P, z + (p + 1))\), then if \(p + 3 \cdot (m \cdot P) \le S + 1\), then if \(\operatorname{rhbandagree}(z + (p + 1), m \cdot P, 1, S)\), then not \(\operatorname{centerperiodic}(T, P)\).
\end{theorem}
\begin{proof}
The goal follows from \(\operatorname{notcenterperiodicofrhbandagreeoneatzerodiagsteptwo}(T, P, m, z, p, S)\).
\end{proof}
\begin{theorem}
\label{ll:paper:not-centerperiodic-of-strip-agree-at-zero-diag-step-two}
For every natural number \(T\) and natural number \(P\) and natural number \(m\) and natural number \(z\) and natural number \(p\) and natural number \(B\) and natural number \(S\), if \(0 < P\), then if \(0 < m\), then if \(T \le z + (p + 1)\), then if \(\operatorname{zerodiagonal}(z)\), then if \(\operatorname{doublesat}(m \cdot P, z)\), then if \(\operatorname{settledby}(z + 1, z + (p + 1))\), then if \(\operatorname{settledby}(z + 1 - m \cdot P, z + (p + 1))\), then if \(p + 3 \cdot (m \cdot P) \le S + 1\), then if \(\operatorname{rhbandagree}(z + (p + 1), m \cdot P, 1, B)\), then if \(\operatorname{rhcolagree}(z + (p + 1), m \cdot P, B + 2, S)\), then not \(\operatorname{centerperiodic}(T, P)\).
\end{theorem}
\begin{proof}
The goal follows from \(\operatorname{notcenterperiodicofstripagreeatzerodiagsteptwo}(T, P, m, z, p, B, S)\).
\end{proof}
\begin{theorem}
\label{ll:paper:not-centerperiodic-of-rdrv-one-agree-sub-twopow-ge}
For every natural number \(T\) and natural number \(P\) and natural number \(m\) and natural number \(k\) and natural number \(S\), if \(0 < P\), then if \(0 < m\), then if \(m \cdot P \le \operatorname{twopow}(k)\), then if \(T + m \cdot P \le \operatorname{twopow}(k)\), then if \(\operatorname{twopow}(k) \le S + 1\), then if for every natural number \(s\), if \(s \le S\), then if \(\operatorname{center}(\operatorname{twopow}(k) - m \cdot P + s) = \operatorname{false}\), then \(\operatorname{rdrv}(\operatorname{twopow}(k) - m \cdot P, \operatorname{twopow}(k), s, 1) = \operatorname{rdrv}(\operatorname{twopow}(k) - m \cdot P, \operatorname{twopow}(k) - m \cdot P, s, 1)\), then not \(\operatorname{centerperiodic}(T, P)\).
\end{theorem}
\begin{proof}
The goal follows from \(\operatorname{notcenterperiodicofrdrvoneagreesubtwopowge}(T, P, m, k, S)\).
\end{proof}
\begin{theorem}
\label{ll:paper:not-centerperiodic-of-rhbandagree-one-sub-twopow-ge}
For every natural number \(T\) and natural number \(P\) and natural number \(m\) and natural number \(k\) and natural number \(S\), if \(0 < P\), then if \(0 < m\), then if \(m \cdot P \le \operatorname{twopow}(k)\), then if \(T + m \cdot P \le \operatorname{twopow}(k)\), then if \(\operatorname{twopow}(k) \le S + 1\), then if \(\operatorname{rhbandagree}(\operatorname{twopow}(k) - m \cdot P, m \cdot P, 1, S)\), then not \(\operatorname{centerperiodic}(T, P)\).
\end{theorem}
\begin{proof}
The goal follows from \(\operatorname{notcenterperiodicofrhbandagreeonesubtwopowge}(T, P, m, k, S)\).
\end{proof}
\begin{theorem}
\label{ll:paper:rdrv-inward-bxor}
For every natural number \(td\) and natural number \(tx\) and natural number \(k\) and natural number \(p\), \(\operatorname{rdrv}(td, tx, k, p) = \operatorname{bxor}(\operatorname{rdrv}(td, tx, k + 1, p + 1), \operatorname{rdrv}(td, tx, k, p + 1) \operatorname{or} \operatorname{rdrv}(td, tx, k, p + 2))\).
\end{theorem}
\begin{proof}
The goal follows from \(\operatorname{rdrvinwardbxor}(td, tx, k, p)\).
\end{proof}
\begin{theorem}
\label{ll:paper:rdrv-inward-of-black-left}
For every natural number \(td\) and natural number \(tx\) and natural number \(k\) and natural number \(p\), if \(\operatorname{rdrv}(td, tx, k, p + 1) = \operatorname{true}\), then \(\operatorname{rdrv}(td, tx, k, p) = \operatorname{not}(\operatorname{rdrv}(td, tx, k + 1, p + 1))\).
\end{theorem}
\begin{proof}
The goal follows from \(\operatorname{rdrvinwardofblackleft}(td, tx, k, p)\).
\end{proof}
\begin{theorem}
\label{ll:paper:rdrv-inward-of-black-right}
For every natural number \(td\) and natural number \(tx\) and natural number \(k\) and natural number \(p\), if \(\operatorname{rdrv}(td, tx, k, p + 2) = \operatorname{true}\), then \(\operatorname{rdrv}(td, tx, k, p) = \operatorname{not}(\operatorname{rdrv}(td, tx, k + 1, p + 1))\).
\end{theorem}
\begin{proof}
The goal follows from \(\operatorname{rdrvinwardofblackright}(td, tx, k, p)\).
\end{proof}
\begin{theorem}
\label{ll:paper:rdrv-step-one-eq-center-xor-or}
For every natural number \(td\) and natural number \(tx\) and natural number \(k\), \(\operatorname{rdrv}(td, tx, k + 1, 1) = \operatorname{bxor}(\operatorname{center}(td + k), \operatorname{rdrv}(td, tx, k, 1) \operatorname{or} \operatorname{rdrv}(td, tx, k, 2))\).
\end{theorem}
\begin{proof}
The goal follows from \(\operatorname{rdrvsteponeeqcenterxoror}(td, tx, k)\).
\end{proof}
\begin{theorem}
\label{ll:paper:rdrv-step-one-eq-not-center-of-black-one}
For every natural number \(td\) and natural number \(tx\) and natural number \(k\), if \(\operatorname{rdrv}(td, tx, k, 1) = \operatorname{true}\), then \(\operatorname{rdrv}(td, tx, k + 1, 1) = \operatorname{not}(\operatorname{center}(td + k))\).
\end{theorem}
\begin{proof}
The goal follows from \(\operatorname{rdrvsteponeeqnotcenterofblackone}(td, tx, k)\).
\end{proof}
\begin{theorem}
\label{ll:paper:rdrv-step-one-defect-eq-or-defect}
For every natural number \(td\) and natural number \(tx2\) and natural number \(tx\) and natural number \(k\), \(\operatorname{bxor}(\operatorname{rdrv}(td, tx2, k + 1, 1), \operatorname{rdrv}(td, tx, k + 1, 1)) = \operatorname{bxor}(\operatorname{rdrv}(td, tx2, k, 1) \operatorname{or} \operatorname{rdrv}(td, tx2, k, 2), \operatorname{rdrv}(td, tx, k, 1) \operatorname{or} \operatorname{rdrv}(td, tx, k, 2))\).
\end{theorem}
\begin{proof}
The goal follows from \(\operatorname{rdrvsteponedefecteqordefect}(td, tx2, tx, k)\).
\end{proof}
\begin{theorem}
\label{ll:paper:rdrv-step-one-eq-of-black-one}
For every natural number \(td\) and natural number \(tx2\) and natural number \(tx\) and natural number \(k\), if \(\operatorname{rdrv}(td, tx2, k, 1) = \operatorname{true}\), then if \(\operatorname{rdrv}(td, tx, k, 1) = \operatorname{true}\), then \(\operatorname{rdrv}(td, tx2, k + 1, 1) = \operatorname{rdrv}(td, tx, k + 1, 1)\).
\end{theorem}
\begin{proof}
The goal follows from \(\operatorname{rdrvsteponeeqofblackone}(td, tx2, tx, k)\).
\end{proof}
\begin{theorem}
\label{ll:paper:rdrv-gt-eq-false}
For every natural number \(td\) and natural number \(tx\) and natural number \(k\) and natural number \(p\), if \(tx + k < p\), then \(\operatorname{rdrv}(td, tx, k, p) = \operatorname{false}\).
\end{theorem}
\begin{proof}
The goal follows from \(\operatorname{rdrvgteqfalse}(td, tx, k, p)\).
\end{proof}
\begin{theorem}
\label{ll:paper:rdrv-self-eq-true}
For every natural number \(td\) and natural number \(tx\) and natural number \(k\), if \(0 < tx\), then \(\operatorname{rdrv}(td, tx, k, tx + k) = \operatorname{true}\).
\end{theorem}
\begin{proof}
The goal follows from \(\operatorname{rdrvselfeqtrue}(td, tx, k)\).
\end{proof}
\begin{theorem}
\label{ll:paper:rdrv-diag-shift-twopow-eq}
For every natural number \(td\) and natural number \(tx\) and natural number \(e\) and natural number \(k\) and natural number \(p\) and natural number \(v\), if \(k < p\), then if \(p \le tx + k\), then if \(tx + k - p \le e\), then \(\operatorname{rdrv}(td, tx + v \cdot \operatorname{twopow}(e), k, p + v \cdot \operatorname{twopow}(e)) = \operatorname{rdrv}(td, tx, k, p)\).
\end{theorem}
\begin{proof}
The goal follows from \(\operatorname{rdrvdiagshifttwopoweq}(td, tx, e, k, p, v)\).
\end{proof}
\begin{theorem}
\label{ll:paper:rdrv-inward-step-agree}
For every natural number \(td\) and natural number \(tx2\) and natural number \(tx\) and natural number \(k\) and natural number \(p\), if \(\operatorname{rdrv}(td, tx2, k + 1, p + 1) = \operatorname{rdrv}(td, tx, k + 1, p + 1)\), then if \((\operatorname{rdrv}(td, tx2, k, p + 1) \operatorname{or} \operatorname{rdrv}(td, tx2, k, p + 2)) = (\operatorname{rdrv}(td, tx, k, p + 1) \operatorname{or} \operatorname{rdrv}(td, tx, k, p + 2))\), then \(\operatorname{rdrv}(td, tx2, k, p) = \operatorname{rdrv}(td, tx, k, p)\).
\end{theorem}
\begin{proof}
The goal follows from \(\operatorname{rdrvinwardstepagree}(td, tx2, tx, k, p)\).
\end{proof}
\begin{theorem}
\label{ll:paper:rdrv-inward-two-col-step-agree}
For every natural number \(td\) and natural number \(tx2\) and natural number \(tx\) and natural number \(S\) and natural number \(p\) and natural number \(s\), if for every natural number \(u\), if \(u \le S + 1\), then \(\operatorname{rdrv}(td, tx2, u, p + 1) = \operatorname{rdrv}(td, tx, u, p + 1)\), then if for every natural number \(u\), if \(u \le S\), then \(\operatorname{rdrv}(td, tx2, u, p + 2) = \operatorname{rdrv}(td, tx, u, p + 2)\), then if \(s \le S\), then \(\operatorname{rdrv}(td, tx2, s, p) = \operatorname{rdrv}(td, tx, s, p)\).
\end{theorem}
\begin{proof}
The goal follows from \(\operatorname{rdrvinwardtwocolstepagree}(td, tx2, tx, S, p, s)\).
\end{proof}
\begin{theorem}
\label{ll:paper:rdrv-inward-two-col-propagate}
For every natural number \(td\) and natural number \(tx2\) and natural number \(tx\) and natural number \(K\) and natural number \(p0\) and natural number \(S\) and natural number \(s\), if for every natural number \(u\), if \(u \le S + K\), then \(\operatorname{rdrv}(td, tx2, u, p0 + K) = \operatorname{rdrv}(td, tx, u, p0 + K)\), then if for every natural number \(u\), if \(u \le S + K\), then \(\operatorname{rdrv}(td, tx2, u, p0 + K + 1) = \operatorname{rdrv}(td, tx, u, p0 + K + 1)\), then if \(s \le S\), then \(\operatorname{rdrv}(td, tx2, s, p0) = \operatorname{rdrv}(td, tx, s, p0)\).
\end{theorem}
\begin{proof}
The goal follows from \(\operatorname{rdrvinwardtwocolpropagate}(td, tx2, tx, K, p0, S, s)\).
\end{proof}
\begin{theorem}
\label{ll:paper:periodic-maskparity-zero-agree-of-rdrv-two-col-agree}
For every natural number \(T\) and natural number \(P\), if \(0 < P\), then if \(\operatorname{centerperiodic}(T, P)\), then for every natural number \(m\) and natural number \(t\) and natural number \(K\) and natural number \(S\) and natural number \(s\), if \(T \le t\), then if for every natural number \(u\), if \(u \le S + K\), then \(\operatorname{rdrv}(t, t + m \cdot P, u, 1 + K) = \operatorname{rdrv}(t, t, u, 1 + K)\), then if for every natural number \(u\), if \(u \le S + K\), then \(\operatorname{rdrv}(t, t + m \cdot P, u, 1 + K + 1) = \operatorname{rdrv}(t, t, u, 1 + K + 1)\), then if \(s \le S\), then \(\operatorname{maskparity}(t + s + m \cdot P, 0) = \operatorname{maskparity}(t + s, 0)\).
\end{theorem}
\begin{proof}
Assume \(hP\), \(hper\).
The goal follows from \(\operatorname{periodicmaskparityzeroagreeofrdrvtwocolagree}(T, P, hP, hper)\).
\end{proof}
\begin{theorem}
\label{ll:paper:not-centerperiodic-of-rdrv-two-col-agree-at-zero-diag}
For every natural number \(T\) and natural number \(P\) and natural number \(m\) and natural number \(z\) and natural number \(p\) and natural number \(K\) and natural number \(S\), if \(0 < P\), then if \(0 < m\), then if \(T \le z + (p + 1)\), then if \(\operatorname{zerodiagonal}(z)\), then if \(\operatorname{doublesat}(m \cdot P, z)\), then if \(\operatorname{settledby}(z + 1, z + (p + 1))\), then if \(\operatorname{settledby}(z + 1 - m \cdot P, z + (p + 1))\), then if \(p + 3 \cdot (m \cdot P) \le S + 1\), then if for every natural number \(u\), if \(u \le S + K\), then \(\operatorname{rdrv}(z + (p + 1), z + (p + 1) + m \cdot P, u, 1 + K) = \operatorname{rdrv}(z + (p + 1), z + (p + 1), u, 1 + K)\), then if for every natural number \(u\), if \(u \le S + K\), then \(\operatorname{rdrv}(z + (p + 1), z + (p + 1) + m \cdot P, u, 1 + K + 1) = \operatorname{rdrv}(z + (p + 1), z + (p + 1), u, 1 + K + 1)\), then not \(\operatorname{centerperiodic}(T, P)\).
\end{theorem}
\begin{proof}
The goal follows from \(\operatorname{notcenterperiodicofrdrvtwocolagreeatzerodiag}(T, P, m, z, p, K, S)\).
\end{proof}
\begin{theorem}
\label{ll:paper:rdrv-diag-shift-twopow-eq-of-threshold}
For every natural number \(td\) and natural number \(tx\) and natural number \(e\) and natural number \(J\) and natural number \(k\) and natural number \(p\) and natural number \(v\), if \(\operatorname{rightthreshold}(e, J)\), then if \(k < p\), then if \(p \le tx + k\), then if \(tx + k - p \le J\), then \(\operatorname{rdrv}(td, tx + v \cdot \operatorname{twopow}(e), k, p + v \cdot \operatorname{twopow}(e)) = \operatorname{rdrv}(td, tx, k, p)\).
\end{theorem}
\begin{proof}
The goal follows from \(\operatorname{rdrvdiagshifttwopoweqofthreshold}(td, tx, e, J, k, p, v)\).
\end{proof}
\begin{theorem}
\label{ll:paper:rdrv-zero-gap-of-threshold}
For every natural number \(td\) and natural number \(e\) and natural number \(J\) and natural number \(u\) and natural number \(k\) and natural number \(j\), if \(\operatorname{rightthreshold}(e, J)\), then if \(1 \le j\), then if \(j \le J\), then if \(k + j < u \cdot \operatorname{twopow}(e)\), then \(\operatorname{rdrv}(td, u \cdot \operatorname{twopow}(e) - k, k, u \cdot \operatorname{twopow}(e) - j) = \operatorname{false}\).
\end{theorem}
\begin{proof}
The goal follows from \(\operatorname{rdrvzerogapofthreshold}(td, e, J, u, k, j)\).
\end{proof}
\begin{theorem}
\label{ll:paper:rdrv-last-one-of-threshold}
For every natural number \(td\) and natural number \(e\) and natural number \(J\) and natural number \(u\) and natural number \(k\) and natural number \(j\), if \(\operatorname{rightthreshold}(e, J)\), then if \(j \le J\), then if \(k + j < (u + 1) \cdot \operatorname{twopow}(e) - 1\), then \(\operatorname{rdrv}(td, (u + 1) \cdot \operatorname{twopow}(e) - 1 - k, k, (u + 1) \cdot \operatorname{twopow}(e) - 1 - j) = \operatorname{true}\).
\end{theorem}
\begin{proof}
The goal follows from \(\operatorname{rdrvlastoneofthreshold}(td, e, J, u, k, j)\).
\end{proof}
\begin{theorem}
\label{ll:paper:rdrv-defect-step-succ}
For every natural number \(td\) and natural number \(tx2\) and natural number \(tx\) and natural number \(k\) and natural number \(p\), \(\operatorname{bxor}(\operatorname{rdrv}(td, tx2, k + 1, p + 1), \operatorname{rdrv}(td, tx, k + 1, p + 1)) = \operatorname{bxor}(\operatorname{bxor}(\operatorname{rdrv}(td, tx2, k, p), \operatorname{rdrv}(td, tx, k, p)), \operatorname{bxor}(\operatorname{rdrv}(td, tx2, k, p + 1) \operatorname{or} \operatorname{rdrv}(td, tx2, k, p + 2), \operatorname{rdrv}(td, tx, k, p + 1) \operatorname{or} \operatorname{rdrv}(td, tx, k, p + 2)))\).
\end{theorem}
\begin{proof}
The goal follows from \(\operatorname{rdrvdefectstepsucc}(td, tx2, tx, k, p)\).
\end{proof}
\begin{theorem}
\label{ll:paper:rdrv-defect-inward-step}
For every natural number \(td\) and natural number \(tx2\) and natural number \(tx\) and natural number \(k\) and natural number \(p\), \(\operatorname{bxor}(\operatorname{rdrv}(td, tx2, k, p), \operatorname{rdrv}(td, tx, k, p)) = \operatorname{bxor}(\operatorname{bxor}(\operatorname{rdrv}(td, tx2, k + 1, p + 1), \operatorname{rdrv}(td, tx, k + 1, p + 1)), \operatorname{bxor}(\operatorname{rdrv}(td, tx2, k, p + 1) \operatorname{or} \operatorname{rdrv}(td, tx2, k, p + 2), \operatorname{rdrv}(td, tx, k, p + 1) \operatorname{or} \operatorname{rdrv}(td, tx, k, p + 2)))\).
\end{theorem}
\begin{proof}
The goal follows from \(\operatorname{rdrvdefectinwardstep}(td, tx2, tx, k, p)\).
\end{proof}
\begin{theorem}
\label{ll:paper:rdrv-step-one-eq-of-or-eq}
For every natural number \(td\) and natural number \(tx2\) and natural number \(tx\) and natural number \(k\), if \((\operatorname{rdrv}(td, tx2, k, 1) \operatorname{or} \operatorname{rdrv}(td, tx2, k, 2)) = (\operatorname{rdrv}(td, tx, k, 1) \operatorname{or} \operatorname{rdrv}(td, tx, k, 2))\), then \(\operatorname{rdrv}(td, tx2, k + 1, 1) = \operatorname{rdrv}(td, tx, k + 1, 1)\).
\end{theorem}
\begin{proof}
The goal follows from \(\operatorname{rdrvsteponeeqoforeq}(td, tx2, tx, k)\).
\end{proof}
\begin{theorem}
\label{ll:paper:rdrv-or-one-eq-of-step-one-eq}
For every natural number \(td\) and natural number \(tx2\) and natural number \(tx\) and natural number \(k\), if \(\operatorname{rdrv}(td, tx2, k + 1, 1) = \operatorname{rdrv}(td, tx, k + 1, 1)\), then \((\operatorname{rdrv}(td, tx2, k, 1) \operatorname{or} \operatorname{rdrv}(td, tx2, k, 2)) = (\operatorname{rdrv}(td, tx, k, 1) \operatorname{or} \operatorname{rdrv}(td, tx, k, 2))\).
\end{theorem}
\begin{proof}
The goal follows from \(\operatorname{rdrvoroneeqofsteponeeq}(td, tx2, tx, k)\).
\end{proof}
\begin{theorem}
\label{ll:paper:rdrv-diag-eq-of-white-pair}
For every natural number \(td\) and natural number \(tx\) and natural number \(k\) and natural number \(p\), if \(\operatorname{rdrv}(td, tx, k, p + 1) = \operatorname{false}\), then if \(\operatorname{rdrv}(td, tx, k, p + 2) = \operatorname{false}\), then \(\operatorname{rdrv}(td, tx, k + 1, p + 1) = \operatorname{rdrv}(td, tx, k, p)\).
\end{theorem}
\begin{proof}
The goal follows from \(\operatorname{rdrvdiageqofwhitepair}(td, tx, k, p)\).
\end{proof}
\begin{theorem}
\label{ll:paper:rdrv-defect-diag-eq-of-or-eq}
For every natural number \(td\) and natural number \(tx2\) and natural number \(tx\) and natural number \(k\) and natural number \(p\), if \((\operatorname{rdrv}(td, tx2, k, p + 1) \operatorname{or} \operatorname{rdrv}(td, tx2, k, p + 2)) = (\operatorname{rdrv}(td, tx, k, p + 1) \operatorname{or} \operatorname{rdrv}(td, tx, k, p + 2))\), then \(\operatorname{bxor}(\operatorname{rdrv}(td, tx2, k + 1, p + 1), \operatorname{rdrv}(td, tx, k + 1, p + 1)) = \operatorname{bxor}(\operatorname{rdrv}(td, tx2, k, p), \operatorname{rdrv}(td, tx, k, p))\).
\end{theorem}
\begin{proof}
The goal follows from \(\operatorname{rdrvdefectdiageqoforeq}(td, tx2, tx, k, p)\).
\end{proof}
\begin{theorem}
\label{ll:paper:rdrv-agree-inward-of-black-left}
For every natural number \(td\) and natural number \(tx2\) and natural number \(tx\) and natural number \(k\) and natural number \(p\), if \(\operatorname{rdrv}(td, tx2, k, p + 1) = \operatorname{true}\), then if \(\operatorname{rdrv}(td, tx, k, p + 1) = \operatorname{true}\), then if \(\operatorname{rdrv}(td, tx2, k + 1, p + 1) = \operatorname{rdrv}(td, tx, k + 1, p + 1)\), then \(\operatorname{rdrv}(td, tx2, k, p) = \operatorname{rdrv}(td, tx, k, p)\).
\end{theorem}
\begin{proof}
The goal follows from \(\operatorname{rdrvagreeinwardofblackleft}(td, tx2, tx, k, p)\).
\end{proof}
\begin{theorem}
\label{ll:paper:periodic-maskparity-zero-agree-of-rdrv-or-one-agree}
For every natural number \(T\) and natural number \(P\), if \(0 < P\), then if \(\operatorname{centerperiodic}(T, P)\), then for every natural number \(m\) and natural number \(t\) and natural number \(S\) and natural number \(s\), if \(T \le t\), then if if \(\operatorname{center}(t) = \operatorname{false}\), then \(\operatorname{rdrv}(t, t + m \cdot P, 0, 1) = \operatorname{rdrv}(t, t, 0, 1)\), then if for every natural number \(u\), if \(u < S\), then if \(\operatorname{center}(t + (u + 1)) = \operatorname{false}\), then \((\operatorname{rdrv}(t, t + m \cdot P, u, 1) \operatorname{or} \operatorname{rdrv}(t, t + m \cdot P, u, 2)) = (\operatorname{rdrv}(t, t, u, 1) \operatorname{or} \operatorname{rdrv}(t, t, u, 2))\), then if \(s \le S\), then \(\operatorname{maskparity}(t + s + m \cdot P, 0) = \operatorname{maskparity}(t + s, 0)\).
\end{theorem}
\begin{proof}
Assume \(hP\), \(hper\).
The goal follows from \(\operatorname{periodicmaskparityzeroagreeofrdrvoroneagree}(T, P, hP, hper)\).
\end{proof}
\begin{theorem}
\label{ll:paper:rdrv-inward-col-or-propagate}
For every natural number \(td\) and natural number \(tx2\) and natural number \(tx\) and natural number \(K\) and natural number \(p0\) and natural number \(S\) and natural number \(s\), if for every natural number \(u\), if \(u \le S + K\), then \(\operatorname{rdrv}(td, tx2, u, p0 + K) = \operatorname{rdrv}(td, tx, u, p0 + K)\), then if for every natural number \(u\), if \(u < S + K\), then \((\operatorname{rdrv}(td, tx2, u, p0 + K) \operatorname{or} \operatorname{rdrv}(td, tx2, u, p0 + K + 1)) = (\operatorname{rdrv}(td, tx, u, p0 + K) \operatorname{or} \operatorname{rdrv}(td, tx, u, p0 + K + 1))\), then if \(s \le S\), then \(\operatorname{rdrv}(td, tx2, s, p0) = \operatorname{rdrv}(td, tx, s, p0)\).
\end{theorem}
\begin{proof}
The goal follows from \(\operatorname{rdrvinwardcolorpropagate}(td, tx2, tx, K, p0, S, s)\).
\end{proof}
\begin{theorem}
\label{ll:paper:periodic-maskparity-zero-agree-of-rdrv-col-or-agree}
For every natural number \(T\) and natural number \(P\), if \(0 < P\), then if \(\operatorname{centerperiodic}(T, P)\), then for every natural number \(m\) and natural number \(t\) and natural number \(K\) and natural number \(S\) and natural number \(s\), if \(T \le t\), then if for every natural number \(u\), if \(u \le S + K\), then \(\operatorname{rdrv}(t, t + m \cdot P, u, 1 + K) = \operatorname{rdrv}(t, t, u, 1 + K)\), then if for every natural number \(u\), if \(u < S + K\), then \((\operatorname{rdrv}(t, t + m \cdot P, u, 1 + K) \operatorname{or} \operatorname{rdrv}(t, t + m \cdot P, u, 1 + K + 1)) = (\operatorname{rdrv}(t, t, u, 1 + K) \operatorname{or} \operatorname{rdrv}(t, t, u, 1 + K + 1))\), then if \(s \le S\), then \(\operatorname{maskparity}(t + s + m \cdot P, 0) = \operatorname{maskparity}(t + s, 0)\).
\end{theorem}
\begin{proof}
Assume \(hP\), \(hper\).
The goal follows from \(\operatorname{periodicmaskparityzeroagreeofrdrvcoloragree}(T, P, hP, hper)\).
\end{proof}
\begin{theorem}
\label{ll:paper:not-centerperiodic-of-rdrv-col-or-agree-at-zero-diag}
For every natural number \(T\) and natural number \(P\) and natural number \(m\) and natural number \(z\) and natural number \(p\) and natural number \(K\) and natural number \(S\), if \(0 < P\), then if \(0 < m\), then if \(T \le z + (p + 1)\), then if \(\operatorname{zerodiagonal}(z)\), then if \(\operatorname{doublesat}(m \cdot P, z)\), then if \(\operatorname{settledby}(z + 1, z + (p + 1))\), then if \(\operatorname{settledby}(z + 1 - m \cdot P, z + (p + 1))\), then if \(p + 3 \cdot (m \cdot P) \le S + 1\), then if for every natural number \(u\), if \(u \le S + K\), then \(\operatorname{rdrv}(z + (p + 1), z + (p + 1) + m \cdot P, u, 1 + K) = \operatorname{rdrv}(z + (p + 1), z + (p + 1), u, 1 + K)\), then if for every natural number \(u\), if \(u < S + K\), then \((\operatorname{rdrv}(z + (p + 1), z + (p + 1) + m \cdot P, u, 1 + K) \operatorname{or} \operatorname{rdrv}(z + (p + 1), z + (p + 1) + m \cdot P, u, 1 + K + 1)) = (\operatorname{rdrv}(z + (p + 1), z + (p + 1), u, 1 + K) \operatorname{or} \operatorname{rdrv}(z + (p + 1), z + (p + 1), u, 1 + K + 1))\), then not \(\operatorname{centerperiodic}(T, P)\).
\end{theorem}
\begin{proof}
The goal follows from \(\operatorname{notcenterperiodicofrdrvcoloragreeatzerodiag}(T, P, m, z, p, K, S)\).
\end{proof}
\begin{theorem}
\label{ll:paper:not-centerperiodic-of-rdrv-col-or-agree-sub-twopow-ge}
For every natural number \(T\) and natural number \(P\) and natural number \(m\) and natural number \(k\) and natural number \(K\) and natural number \(S\), if \(0 < P\), then if \(0 < m\), then if \(m \cdot P \le \operatorname{twopow}(k)\), then if \(T + m \cdot P \le \operatorname{twopow}(k)\), then if \(\operatorname{twopow}(k) \le S + 1\), then if for every natural number \(u\), if \(u \le S + K\), then \(\operatorname{rdrv}(\operatorname{twopow}(k) - m \cdot P, \operatorname{twopow}(k), u, 1 + K) = \operatorname{rdrv}(\operatorname{twopow}(k) - m \cdot P, \operatorname{twopow}(k) - m \cdot P, u, 1 + K)\), then if for every natural number \(u\), if \(u < S + K\), then \((\operatorname{rdrv}(\operatorname{twopow}(k) - m \cdot P, \operatorname{twopow}(k), u, 1 + K) \operatorname{or} \operatorname{rdrv}(\operatorname{twopow}(k) - m \cdot P, \operatorname{twopow}(k), u, 1 + K + 1)) = (\operatorname{rdrv}(\operatorname{twopow}(k) - m \cdot P, \operatorname{twopow}(k) - m \cdot P, u, 1 + K) \operatorname{or} \operatorname{rdrv}(\operatorname{twopow}(k) - m \cdot P, \operatorname{twopow}(k) - m \cdot P, u, 1 + K + 1))\), then not \(\operatorname{centerperiodic}(T, P)\).
\end{theorem}
\begin{proof}
The goal follows from \(\operatorname{notcenterperiodicofrdrvcoloragreesubtwopowge}(T, P, m, k, K, S)\).
\end{proof}
\begin{theorem}
\label{ll:paper:rdrv-white-seed-init-eq-false}
For every natural number \(td\) and natural number \(p\), if \(0 < p\), then \(\operatorname{rdrv}(td, 0, 0, p) = \operatorname{false}\).
\end{theorem}
\begin{proof}
The goal follows from \(\operatorname{rdrvwhiteseediniteqfalse}(td, p)\).
\end{proof}
\begin{theorem}
\label{ll:paper:rdrv-white-seed-gt-eq-false}
For every natural number \(td\) and natural number \(k\) and natural number \(p\), if \(k < p\), then \(\operatorname{rdrv}(td, 0, k, p) = \operatorname{false}\).
\end{theorem}
\begin{proof}
The goal follows from \(\operatorname{rdrvwhiteseedgteqfalse}(td, k, p)\).
\end{proof}
\begin{theorem}
\label{ll:paper:rdrv-white-seed-diag-eq-center}
For every natural number \(td\) and natural number \(k\), \(\operatorname{rdrv}(td, 0, k, k) = \operatorname{center}(td)\).
\end{theorem}
\begin{proof}
The goal follows from \(\operatorname{rdrvwhiteseeddiageqcenter}(td, k)\).
\end{proof}
\begin{theorem}
\label{ll:paper:periodic-rdrv-white-seed-shift-eq}
For every natural number \(T\) and natural number \(P\), if \(0 < P\), then if \(\operatorname{centerperiodic}(T, P)\), then for every natural number \(m\) and natural number \(td\) and natural number \(k\) and natural number \(p\), if \(T \le td\), then \(\operatorname{rdrv}(td + m \cdot P, 0, k, p) = \operatorname{rdrv}(td, 0, k, p)\).
\end{theorem}
\begin{proof}
Assume \(hP\), \(hper\).
The goal follows from \(\operatorname{periodicrdrvwhiteseedshifteq}(T, P, hP, hper)\).
\end{proof}
\begin{theorem}
\label{ll:paper:rdrv-eq-white-seed-of-initial-white-band}
For every natural number \(td\) and natural number \(tx\) and natural number \(L\) and natural number \(k\) and natural number \(p\), if for every natural number \(i\), if \(i \le L\), then \(\operatorname{rhcenter}(tx, 1 + i) = \operatorname{false}\), then if \(p + k \le L + 1\), then \(\operatorname{rdrv}(td, tx, k, p) = \operatorname{rdrv}(td, 0, k, p)\).
\end{theorem}
\begin{proof}
The goal follows from \(\operatorname{rdrveqwhiteseedofinitialwhiteband}(td, tx, L, k, p)\).
\end{proof}
\begin{theorem}
\label{ll:paper:rdrv-col-one-eq-white-seed-of-initial-white-band}
For every natural number \(td\) and natural number \(tx\) and natural number \(S\) and natural number \(s\), if for every natural number \(i\), if \(i \le S\), then \(\operatorname{rhcenter}(tx, 1 + i) = \operatorname{false}\), then if \(s \le S\), then \(\operatorname{rdrv}(td, tx, s, 1) = \operatorname{rdrv}(td, 0, s, 1)\).
\end{theorem}
\begin{proof}
The goal follows from \(\operatorname{rdrvcoloneeqwhiteseedofinitialwhiteband}(td, tx, S, s)\).
\end{proof}
\begin{theorem}
\label{ll:paper:rdrv-step-one-eq-true-of-white-center-and-black-one}
For every natural number \(td\) and natural number \(tx\) and natural number \(k\), if \(\operatorname{center}(td + k) = \operatorname{false}\), then if \(\operatorname{rdrv}(td, tx, k, 1) = \operatorname{true}\), then \(\operatorname{rdrv}(td, tx, k + 1, 1) = \operatorname{true}\).
\end{theorem}
\begin{proof}
The goal follows from \(\operatorname{rdrvsteponeeqtrueofwhitecenterandblackone}(td, tx, k)\).
\end{proof}
\begin{theorem}
\label{ll:paper:rdrv-white-center-run-black-one-persist}
For every natural number \(td\) and natural number \(tx\) and natural number \(k\) and natural number \(W\) and natural number \(j\), if \(\operatorname{rdrv}(td, tx, k, 1) = \operatorname{true}\), then if for every natural number \(u\), if \(u < W\), then \(\operatorname{center}(td + (k + u)) = \operatorname{false}\), then if \(j \le W\), then \(\operatorname{rdrv}(td, tx, k + j, 1) = \operatorname{true}\).
\end{theorem}
\begin{proof}
The goal follows from \(\operatorname{rdrvwhitecenterrunblackonepersist}(td, tx, k, W, j)\).
\end{proof}
\begin{theorem}
\label{ll:paper:rdrv-white-center-run-black-one-agree}
For every natural number \(td\) and natural number \(tx2\) and natural number \(tx\) and natural number \(k\) and natural number \(W\) and natural number \(j\), if \(\operatorname{rdrv}(td, tx2, k, 1) = \operatorname{true}\), then if \(\operatorname{rdrv}(td, tx, k, 1) = \operatorname{true}\), then if for every natural number \(u\), if \(u < W\), then \(\operatorname{center}(td + (k + u)) = \operatorname{false}\), then if \(j \le W + 1\), then \(\operatorname{rdrv}(td, tx2, k + j, 1) = \operatorname{rdrv}(td, tx, k + j, 1)\).
\end{theorem}
\begin{proof}
The goal follows from \(\operatorname{rdrvwhitecenterrunblackoneagree}(td, tx2, tx, k, W, j)\).
\end{proof}
\begin{theorem}
\label{ll:paper:rdrv-step-one-eq-true-of-maskparity-zero-true}
For every natural number \(t\), if \(\operatorname{maskparity}(t, 0) = \operatorname{true}\), then \(\operatorname{rdrv}(t, t, 1, 1) = \operatorname{true}\).
\end{theorem}
\begin{proof}
The goal follows from \(\operatorname{rdrvsteponeeqtrueofmaskparityzerotrue}(t)\).
\end{proof}
\begin{theorem}
\label{ll:paper:periodic-rdrv-step-one-eq-true-of-maskparity-zero-true}
For every natural number \(T\) and natural number \(P\), if \(0 < P\), then if \(\operatorname{centerperiodic}(T, P)\), then for every natural number \(m\) and natural number \(t\), if \(T \le t\), then if \(\operatorname{maskparity}(t + m \cdot P, 0) = \operatorname{true}\), then \(\operatorname{rdrv}(t, t + m \cdot P, 1, 1) = \operatorname{true}\).
\end{theorem}
\begin{proof}
Assume \(hP\), \(hper\).
The goal follows from \(\operatorname{periodicrdrvsteponeeqtrueofmaskparityzerotrue}(T, P, hP, hper)\).
\end{proof}
\begin{theorem}
\label{ll:paper:periodic-maskparity-zero-agree-on-white-run-of-maskparity-zero-true}
For every natural number \(T\) and natural number \(P\), if \(0 < P\), then if \(\operatorname{centerperiodic}(T, P)\), then for every natural number \(m\) and natural number \(t\) and natural number \(W\) and natural number \(s\), if \(T \le t\), then if \(\operatorname{maskparity}(t + m \cdot P, 0) = \operatorname{true}\), then if \(\operatorname{maskparity}(t, 0) = \operatorname{true}\), then if for every natural number \(u\), if \(u < W\), then \(\operatorname{center}(t + (1 + u)) = \operatorname{false}\), then if \(s \le W + 2\), then \(\operatorname{maskparity}(t + s + m \cdot P, 0) = \operatorname{maskparity}(t + s, 0)\).
\end{theorem}
\begin{proof}
Assume \(hP\), \(hper\).
The goal follows from \(\operatorname{periodicmaskparityzeroagreeonwhiterunofmaskparityzerotrue}(T, P, hP, hper)\).
\end{proof}
\begin{theorem}
\label{ll:paper:rdrv-two-col-eq-white-seed-of-zero-gap}
For every natural number \(td\) and natural number \(e\) and natural number \(J\) and natural number \(u\) and natural number \(k\) and natural number \(j\), if \(\operatorname{rightthreshold}(e, J)\), then if \(1 \le j\), then if \(j + 1 \le J\), then if \(k + (j + 1) < u \cdot \operatorname{twopow}(e)\), then \(\operatorname{rdrv}(td, u \cdot \operatorname{twopow}(e) - k, k, u \cdot \operatorname{twopow}(e) - (j + 1)) = \operatorname{rdrv}(td, 0, k, u \cdot \operatorname{twopow}(e) - (j + 1))\) and \(\operatorname{rdrv}(td, u \cdot \operatorname{twopow}(e) - k, k, u \cdot \operatorname{twopow}(e) - j) = \operatorname{rdrv}(td, 0, k, u \cdot \operatorname{twopow}(e) - j)\).
\end{theorem}
\begin{proof}
The goal follows from \(\operatorname{rdrvtwocoleqwhiteseedofzerogap}(td, e, J, u, k, j)\).
\end{proof}
\begin{theorem}
\label{ll:paper:periodic-maskparity-zero-agree-of-both-eq-white-seed}
For every natural number \(T\) and natural number \(P\), if \(0 < P\), then if \(\operatorname{centerperiodic}(T, P)\), then for every natural number \(m\) and natural number \(t\) and natural number \(s\), if \(T \le t\), then if if \(\operatorname{center}(t + s) = \operatorname{false}\), then \(\operatorname{rdrv}(t, t + m \cdot P, s, 1) = \operatorname{rdrv}(t, 0, s, 1)\) and \(\operatorname{rdrv}(t, t, s, 1) = \operatorname{rdrv}(t, 0, s, 1)\), then \(\operatorname{maskparity}(t + s + m \cdot P, 0) = \operatorname{maskparity}(t + s, 0)\).
\end{theorem}
\begin{proof}
Assume \(hP\), \(hper\).
The goal follows from \(\operatorname{periodicmaskparityzeroagreeofbotheqwhiteseed}(T, P, hP, hper)\).
\end{proof}
\begin{theorem}
\label{ll:paper:not-centerperiodic-of-both-eq-white-seed-at-zero-diag}
For every natural number \(T\) and natural number \(P\) and natural number \(m\) and natural number \(z\) and natural number \(p\) and natural number \(S\), if \(0 < P\), then if \(0 < m\), then if \(T \le z + (p + 1)\), then if \(\operatorname{zerodiagonal}(z)\), then if \(\operatorname{doublesat}(m \cdot P, z)\), then if \(\operatorname{settledby}(z + 1, z + (p + 1))\), then if \(\operatorname{settledby}(z + 1 - m \cdot P, z + (p + 1))\), then if \(p + 3 \cdot (m \cdot P) \le S + 1\), then if for every natural number \(s\), if \(s \le S\), then if \(\operatorname{center}(z + (p + 1) + s) = \operatorname{false}\), then \(\operatorname{rdrv}(z + (p + 1), z + (p + 1) + m \cdot P, s, 1) = \operatorname{rdrv}(z + (p + 1), 0, s, 1)\) and \(\operatorname{rdrv}(z + (p + 1), z + (p + 1), s, 1) = \operatorname{rdrv}(z + (p + 1), 0, s, 1)\), then not \(\operatorname{centerperiodic}(T, P)\).
\end{theorem}
\begin{proof}
The goal follows from \(\operatorname{notcenterperiodicofbotheqwhiteseedatzerodiag}(T, P, m, z, p, S)\).
\end{proof}
\begin{theorem}
\label{ll:paper:rdrv-white-seed-subdiag-one-step}
For every natural number \(td\) and natural number \(k\), \(\operatorname{rdrv}(td, 0, k + 2, k + 1) = \operatorname{bxor}(\operatorname{rdrv}(td, 0, k + 1, k), \operatorname{center}(td))\).
\end{theorem}
\begin{proof}
The goal follows from \(\operatorname{rdrvwhiteseedsubdiagonestep}(td, k)\).
\end{proof}
\begin{theorem}
\label{ll:paper:rdrv-white-seed-subdiag-one-period-two}
For every natural number \(td\) and natural number \(k\), \(\operatorname{rdrv}(td, 0, k + 3, k + 2) = \operatorname{rdrv}(td, 0, k + 1, k)\).
\end{theorem}
\begin{proof}
The goal follows from \(\operatorname{rdrvwhiteseedsubdiagoneperiodtwo}(td, k)\).
\end{proof}
\begin{theorem}
\label{ll:paper:rdrv-white-seed-subdiag-one-even}
For every natural number \(td\) and natural number \(u\), \(\operatorname{rdrv}(td, 0, 2 \cdot u + 1, 2 \cdot u) = \operatorname{center}(td + 1)\).
\end{theorem}
\begin{proof}
The goal follows from \(\operatorname{rdrvwhiteseedsubdiagoneeven}(td, u)\).
\end{proof}
\begin{theorem}
\label{ll:paper:rdrv-white-seed-subdiag-one-odd}
For every natural number \(td\) and natural number \(u\), \(\operatorname{rdrv}(td, 0, 2 \cdot u + 2, 2 \cdot u + 1) = \operatorname{bxor}(\operatorname{center}(td + 1), \operatorname{center}(td))\).
\end{theorem}
\begin{proof}
The goal follows from \(\operatorname{rdrvwhiteseedsubdiagoneodd}(td, u)\).
\end{proof}
\begin{theorem}
\label{ll:paper:rdrv-white-seed-subdiag-period-twopow}
For every natural number \(td\) and natural number \(d\) and natural number \(k\) and natural number \(u\), \(\operatorname{rdrv}(td, 0, k + d + u \cdot \operatorname{twopow}(d), k + u \cdot \operatorname{twopow}(d)) = \operatorname{rdrv}(td, 0, k + d, k)\).
\end{theorem}
\begin{proof}
The goal follows from \(\operatorname{rdrvwhiteseedsubdiagperiodtwopow}(td, d, k, u)\).
\end{proof}
\begin{theorem}
\label{ll:paper:rdrv-white-seed-diag-shift-twopow-eq}
For every natural number \(td\) and natural number \(d\) and natural number \(k\) and natural number \(p\) and natural number \(u\), if \(p \le k\), then if \(k - p \le d\), then \(\operatorname{rdrv}(td, 0, k + u \cdot \operatorname{twopow}(d), p + u \cdot \operatorname{twopow}(d)) = \operatorname{rdrv}(td, 0, k, p)\).
\end{theorem}
\begin{proof}
The goal follows from \(\operatorname{rdrvwhiteseeddiagshifttwopoweq}(td, d, k, p, u)\).
\end{proof}
\begin{theorem}
\label{ll:paper:rdrv-white-seed-scale-mult-eq-center}
For every natural number \(td\) and natural number \(d\) and natural number \(u\), \(\operatorname{rdrv}(td, 0, d + u \cdot \operatorname{twopow}(d), u \cdot \operatorname{twopow}(d)) = \operatorname{center}(td + d)\).
\end{theorem}
\begin{proof}
The goal follows from \(\operatorname{rdrvwhiteseedscalemulteqcenter}(td, d, u)\).
\end{proof}
\begin{theorem}
\label{ll:paper:rdrv-white-seed-scale-eq-center}
For every natural number \(td\) and natural number \(d\), \(\operatorname{rdrv}(td, 0, d + \operatorname{twopow}(d), \operatorname{twopow}(d)) = \operatorname{center}(td + d)\).
\end{theorem}
\begin{proof}
The goal follows from \(\operatorname{rdrvwhiteseedscaleeqcenter}(td, d)\).
\end{proof}
\begin{theorem}
\label{ll:paper:rdrv-white-seed-scale-ge-eq-center}
For every natural number \(td\) and natural number \(d\) and natural number \(k\), if \(d \le k\), then \(\operatorname{rdrv}(td, 0, d + \operatorname{twopow}(k), \operatorname{twopow}(k)) = \operatorname{center}(td + d)\).
\end{theorem}
\begin{proof}
The goal follows from \(\operatorname{rdrvwhiteseedscalegeeqcenter}(td, d, k)\).
\end{proof}
\begin{theorem}
\label{ll:paper:rdrv-white-seed-scale-double-eq-scale}
For every natural number \(td\) and natural number \(d\) and natural number \(k\), if \(d \le k\), then \(\operatorname{rdrv}(td, 0, d + \operatorname{twopow}(k + 1), \operatorname{twopow}(k + 1)) = \operatorname{rdrv}(td, 0, d + \operatorname{twopow}(k), \operatorname{twopow}(k))\).
\end{theorem}
\begin{proof}
The goal follows from \(\operatorname{rdrvwhiteseedscaledoubleeqscale}(td, d, k)\).
\end{proof}
\begin{theorem}
\label{ll:paper:rdrv-scale-eq-center-of-initial-white-band}
For every natural number \(td\) and natural number \(tx\) and natural number \(L\) and natural number \(d\) and natural number \(k\), if for every natural number \(i\), if \(i \le L\), then \(\operatorname{rhcenter}(tx, 1 + i) = \operatorname{false}\), then if \(d \le k\), then if \(d + 2 \cdot \operatorname{twopow}(k) \le L + 1\), then \(\operatorname{rdrv}(td, tx, d + \operatorname{twopow}(k), \operatorname{twopow}(k)) = \operatorname{center}(td + d)\).
\end{theorem}
\begin{proof}
The goal follows from \(\operatorname{rdrvscaleeqcenterofinitialwhiteband}(td, tx, L, d, k)\).
\end{proof}
\begin{theorem}
\label{ll:paper:rhcenter-scale-eq-center-of-initial-white-band}
For every natural number \(t\) and natural number \(L\) and natural number \(d\) and natural number \(k\), if for every natural number \(i\), if \(i \le L\), then \(\operatorname{rhcenter}(t, 1 + i) = \operatorname{false}\), then if \(d \le k\), then if \(d + 2 \cdot \operatorname{twopow}(k) \le L + 1\), then \(\operatorname{rhcenter}(t + (d + \operatorname{twopow}(k)), \operatorname{twopow}(k)) = \operatorname{center}(t + d)\).
\end{theorem}
\begin{proof}
The goal follows from \(\operatorname{rhcenterscaleeqcenterofinitialwhiteband}(t, L, d, k)\).
\end{proof}
\begin{theorem}
\label{ll:paper:rdrv-subdiag-period-twopow-of-initial-white-band}
For every natural number \(td\) and natural number \(tx\) and natural number \(L\) and natural number \(d\) and natural number \(k\) and natural number \(u\), if for every natural number \(i\), if \(i \le L\), then \(\operatorname{rhcenter}(tx, 1 + i) = \operatorname{false}\), then if \(k + u \cdot \operatorname{twopow}(d) + (k + d + u \cdot \operatorname{twopow}(d)) \le L + 1\), then \(\operatorname{rdrv}(td, tx, k + d + u \cdot \operatorname{twopow}(d), k + u \cdot \operatorname{twopow}(d)) = \operatorname{rdrv}(td, tx, k + d, k)\).
\end{theorem}
\begin{proof}
The goal follows from \(\operatorname{rdrvsubdiagperiodtwopowofinitialwhiteband}(td, tx, L, d, k, u)\).
\end{proof}
\begin{theorem}
\label{ll:paper:rdrv-eq-of-both-initial-white-band}
For every natural number \(td\) and natural number \(tx2\) and natural number \(tx\) and natural number \(L\) and natural number \(k\) and natural number \(p\), if for every natural number \(i\), if \(i \le L\), then \(\operatorname{rhcenter}(tx2, 1 + i) = \operatorname{false}\), then if for every natural number \(i\), if \(i \le L\), then \(\operatorname{rhcenter}(tx, 1 + i) = \operatorname{false}\), then if \(p + k \le L + 1\), then \(\operatorname{rdrv}(td, tx2, k, p) = \operatorname{rdrv}(td, tx, k, p)\).
\end{theorem}
\begin{proof}
The goal follows from \(\operatorname{rdrveqofbothinitialwhiteband}(td, tx2, tx, L, k, p)\).
\end{proof}
\begin{theorem}
\label{ll:paper:periodic-rhcenter-shift-eq-of-both-initial-white-band}
For every natural number \(T\) and natural number \(P\), if \(0 < P\), then if \(\operatorname{centerperiodic}(T, P)\), then for every natural number \(m\) and natural number \(t\) and natural number \(L\) and natural number \(k\) and natural number \(p\), if \(T \le t\), then if for every natural number \(i\), if \(i \le L\), then \(\operatorname{rhcenter}(t + m \cdot P, 1 + i) = \operatorname{false}\), then if for every natural number \(i\), if \(i \le L\), then \(\operatorname{rhcenter}(t, 1 + i) = \operatorname{false}\), then if \(p + k \le L + 1\), then \(\operatorname{rhcenter}(t + m \cdot P + k, p) = \operatorname{rhcenter}(t + k, p)\).
\end{theorem}
\begin{proof}
Assume \(hP\), \(hper\).
The goal follows from \(\operatorname{periodicrhcentershifteqofbothinitialwhiteband}(T, P, hP, hper)\).
\end{proof}
\begin{theorem}
\label{ll:paper:rdrv-main-diag-step-rhcenter}
For every natural number \(td\) and natural number \(tx\) and natural number \(k\), \(\operatorname{rdrv}(td, tx, k + 1, k + 1) = \operatorname{bxor}(\operatorname{rdrv}(td, tx, k, k), \operatorname{rhcenter}(tx + k, k + 1) \operatorname{or} \operatorname{rhcenter}(tx + k, k + 2))\).
\end{theorem}
\begin{proof}
The goal follows from \(\operatorname{rdrvmaindiagsteprhcenter}(td, tx, k)\).
\end{proof}
\begin{theorem}
\label{ll:paper:rdrv-main-diag-drive-defect-const}
For every natural number \(td2\) and natural number \(td\) and natural number \(tx\) and natural number \(k\), \(\operatorname{bxor}(\operatorname{rdrv}(td2, tx, k, k), \operatorname{rdrv}(td, tx, k, k)) = \operatorname{bxor}(\operatorname{center}(td2), \operatorname{center}(td))\).
\end{theorem}
\begin{proof}
The goal follows from \(\operatorname{rdrvmaindiagdrivedefectconst}(td2, td, tx, k)\).
\end{proof}
\begin{theorem}
\label{ll:paper:rdrv-main-diag-eq-rhcenter-xor-centers}
For every natural number \(td\) and natural number \(tx\) and natural number \(k\), \(\operatorname{rdrv}(td, tx, k, k) = \operatorname{bxor}(\operatorname{rhcenter}(tx + k, k), \operatorname{bxor}(\operatorname{center}(td), \operatorname{center}(tx)))\).
\end{theorem}
\begin{proof}
The goal follows from \(\operatorname{rdrvmaindiageqrhcenterxorcenters}(td, tx, k)\).
\end{proof}
\begin{theorem}
\label{ll:paper:rdrv-main-diag-eq-rhcenter-of-center-eq}
For every natural number \(td\) and natural number \(tx\) and natural number \(k\), if \(\operatorname{center}(td) = \operatorname{center}(tx)\), then \(\operatorname{rdrv}(td, tx, k, k) = \operatorname{rhcenter}(tx + k, k)\).
\end{theorem}
\begin{proof}
The goal follows from \(\operatorname{rdrvmaindiageqrhcenterofcentereq}(td, tx, k)\).
\end{proof}
\begin{theorem}
\label{ll:paper:rdrv-drive-wavefront-defect-const}
For every natural number \(td2\) and natural number \(td\) and natural number \(tx\) and natural number \(d\) and natural number \(p\), if for every natural number \(j\), if \(j < d\), then \(\operatorname{center}(td2 + j) = \operatorname{center}(td + j)\), then \(\operatorname{bxor}(\operatorname{rdrv}(td2, tx, p + d, p), \operatorname{rdrv}(td, tx, p + d, p)) = \operatorname{bxor}(\operatorname{center}(td2 + d), \operatorname{center}(td + d))\).
\end{theorem}
\begin{proof}
The goal follows from \(\operatorname{rdrvdrivewavefrontdefectconst}(td2, td, tx, d, p)\).
\end{proof}
\begin{theorem}
\label{ll:paper:rdrv-drive-wavefront-eq-rhcenter-xor}
For every natural number \(td\) and natural number \(tx\) and natural number \(d\) and natural number \(p\), if for every natural number \(j\), if \(j < d\), then \(\operatorname{center}(td + j) = \operatorname{center}(tx + j)\), then \(\operatorname{rdrv}(td, tx, p + d, p) = \operatorname{bxor}(\operatorname{rhcenter}(tx + (p + d), p), \operatorname{bxor}(\operatorname{center}(td + d), \operatorname{center}(tx + d)))\).
\end{theorem}
\begin{proof}
The goal follows from \(\operatorname{rdrvdrivewavefronteqrhcenterxor}(td, tx, d, p)\).
\end{proof}
\begin{theorem}
\label{ll:paper:rdrv-drive-wavefront-flip-of-first-center-diff}
For every natural number \(td2\) and natural number \(td\) and natural number \(tx\) and natural number \(d\) and natural number \(p\), if for every natural number \(j\), if \(j < d\), then \(\operatorname{center}(td2 + j) = \operatorname{center}(td + j)\), then if \(\operatorname{center}(td2 + d) = \operatorname{not}(\operatorname{center}(td + d))\), then \(\operatorname{rdrv}(td2, tx, p + d, p) = \operatorname{not}(\operatorname{rdrv}(td, tx, p + d, p))\).
\end{theorem}
\begin{proof}
The goal follows from \(\operatorname{rdrvdrivewavefrontflipoffirstcenterdiff}(td2, td, tx, d, p)\).
\end{proof}
\begin{theorem}
\label{ll:paper:rdrv-inward-agree-step-of-black-shield}
For every natural number \(td\) and natural number \(tx2\) and natural number \(tx\) and natural number \(k\) and natural number \(d\), if \(\operatorname{rdrv}(td, tx2, k, d) = \operatorname{rdrv}(td, tx, k, d)\), then if \(\operatorname{rdrv}(td, tx2, k, d + 1) = \operatorname{rdrv}(td, tx, k, d + 1)\), then if \(\operatorname{rdrv}(td, tx, k, d + 1) = \operatorname{true}\), then \(\operatorname{rdrv}(td, tx2, k + 1, d + 1) = \operatorname{rdrv}(td, tx, k + 1, d + 1)\).
\end{theorem}
\begin{proof}
The goal follows from \(\operatorname{rdrvinwardagreestepofblackshield}(td, tx2, tx, k, d)\).
\end{proof}
\begin{theorem}
\label{ll:paper:rdrv-inward-defect-step-of-white-shield}
For every natural number \(td\) and natural number \(tx2\) and natural number \(tx\) and natural number \(k\) and natural number \(d\), if \(\operatorname{rdrv}(td, tx2, k, d) = \operatorname{rdrv}(td, tx, k, d)\), then if \(\operatorname{rdrv}(td, tx2, k, d + 1) = \operatorname{rdrv}(td, tx, k, d + 1)\), then if \(\operatorname{rdrv}(td, tx, k, d + 1) = \operatorname{false}\), then \(\operatorname{bxor}(\operatorname{rdrv}(td, tx2, k + 1, d + 1), \operatorname{rdrv}(td, tx, k + 1, d + 1)) = \operatorname{bxor}(\operatorname{rdrv}(td, tx2, k, d + 2), \operatorname{rdrv}(td, tx, k, d + 2))\).
\end{theorem}
\begin{proof}
The goal follows from \(\operatorname{rdrvinwarddefectstepofwhiteshield}(td, tx2, tx, k, d)\).
\end{proof}
\begin{theorem}
\label{ll:paper:rdrv-inward-agree-step-of-three-agree}
For every natural number \(td\) and natural number \(tx2\) and natural number \(tx\) and natural number \(k\) and natural number \(d\), if \(\operatorname{rdrv}(td, tx2, k, d) = \operatorname{rdrv}(td, tx, k, d)\), then if \(\operatorname{rdrv}(td, tx2, k, d + 1) = \operatorname{rdrv}(td, tx, k, d + 1)\), then if \(\operatorname{rdrv}(td, tx2, k, d + 2) = \operatorname{rdrv}(td, tx, k, d + 2)\), then \(\operatorname{rdrv}(td, tx2, k + 1, d + 1) = \operatorname{rdrv}(td, tx, k + 1, d + 1)\).
\end{theorem}
\begin{proof}
The goal follows from \(\operatorname{rdrvinwardagreestepofthreeagree}(td, tx2, tx, k, d)\).
\end{proof}
\begin{theorem}
\label{ll:paper:rdrv-init-wavefront-defect-of-white}
For every natural number \(td\) and natural number \(tx\) and natural number \(Q\) and natural number \(L\) and natural number \(d\) and natural number \(k\), if \(\operatorname{rhbandagree}(tx, Q, 1, L)\), then if \(d + k \le L\), then if \(\operatorname{rdrv}(td, tx, k, d + 1) = \operatorname{false}\), then \(\operatorname{bxor}(\operatorname{rdrv}(td, tx + Q, k + 1, d + 1), \operatorname{rdrv}(td, tx, k + 1, d + 1)) = \operatorname{bxor}(\operatorname{rdrv}(td, tx + Q, k, d + 2), \operatorname{rdrv}(td, tx, k, d + 2))\).
\end{theorem}
\begin{proof}
The goal follows from \(\operatorname{rdrvinitwavefrontdefectofwhite}(td, tx, Q, L, d, k)\).
\end{proof}
\begin{theorem}
\label{ll:paper:rdrv-init-wavefront-absorb-of-black}
For every natural number \(td\) and natural number \(tx\) and natural number \(Q\) and natural number \(L\) and natural number \(d\) and natural number \(k\), if \(\operatorname{rhbandagree}(tx, Q, 1, L)\), then if \(d + k \le L\), then if \(\operatorname{rdrv}(td, tx, k, d + 1) = \operatorname{true}\), then \(\operatorname{rdrv}(td, tx + Q, k + 1, d + 1) = \operatorname{rdrv}(td, tx, k + 1, d + 1)\).
\end{theorem}
\begin{proof}
The goal follows from \(\operatorname{rdrvinitwavefrontabsorbofblack}(td, tx, Q, L, d, k)\).
\end{proof}
\begin{theorem}
\label{ll:paper:rdrv-init-wavefront-agree-propagate}
For every natural number \(td\) and natural number \(tx\) and natural number \(Q\) and natural number \(L\) and natural number \(d\) and natural number \(k\) and natural number \(j\), if \(\operatorname{rhbandagree}(tx, Q, 1, L)\), then if \(d + (k + j) \le L + 1\), then if \(\operatorname{rdrv}(td, tx + Q, k, d + 1 + j) = \operatorname{rdrv}(td, tx, k, d + 1 + j)\), then \(\operatorname{rdrv}(td, tx + Q, k + j, d + 1) = \operatorname{rdrv}(td, tx, k + j, d + 1)\).
\end{theorem}
\begin{proof}
The goal follows from \(\operatorname{rdrvinitwavefrontagreepropagate}(td, tx, Q, L, d, k, j)\).
\end{proof}
\begin{theorem}
\label{ll:paper:rdrv-init-wavefront-col-one-agree-of-black-shield}
For every natural number \(td\) and natural number \(tx\) and natural number \(Q\) and natural number \(L\) and natural number \(k\), if \(\operatorname{rhbandagree}(tx, Q, 1, L)\), then if \(k \le L\), then if \(\operatorname{rdrv}(td, tx, k, L + 1 - k) = \operatorname{true}\), then \(\operatorname{rdrv}(td, tx + Q, L + 1, 1) = \operatorname{rdrv}(td, tx, L + 1, 1)\).
\end{theorem}
\begin{proof}
The goal follows from \(\operatorname{rdrvinitwavefrontcoloneagreeofblackshield}(td, tx, Q, L, k)\).
\end{proof}
\begin{theorem}
\label{ll:paper:periodic-maskparity-zero-agree-of-rhbandagree-and-black-shield}
For every natural number \(T\) and natural number \(P\), if \(0 < P\), then if \(\operatorname{centerperiodic}(T, P)\), then for every natural number \(m\) and natural number \(t\) and natural number \(L\) and natural number \(k\) and natural number \(s\), if \(T \le t\), then if \(\operatorname{rhbandagree}(t, m \cdot P, 1, L)\), then if \(k \le L\), then if \(\operatorname{rdrv}(t, t, k, L + 1 - k) = \operatorname{true}\), then if \(s \le L + 1\), then \(\operatorname{maskparity}(t + s + m \cdot P, 0) = \operatorname{maskparity}(t + s, 0)\).
\end{theorem}
\begin{proof}
Assume \(hP\), \(hper\).
The goal follows from \(\operatorname{periodicmaskparityzeroagreeofrhbandagreeandblackshield}(T, P, hP, hper)\).
\end{proof}
\begin{theorem}
\label{ll:paper:rdrv-strip-agree-of-white-boundary-col-agree}
For every natural number \(td\) and natural number \(tx\) and natural number \(Q\) and natural number \(B\) and natural number \(S\) and natural number \(s\) and natural number \(p\), if \(\operatorname{rhbandagree}(tx, Q, 1, B)\), then if for every natural number \(u\), if \(u < S\), then if \(\operatorname{rdrv}(td, tx, u, B + 1) = \operatorname{false}\), then \(\operatorname{rdrv}(td, tx + Q, u, B + 2) = \operatorname{rdrv}(td, tx, u, B + 2)\), then if \(s \le S\), then if \(p \le B + 1\), then \(\operatorname{rdrv}(td, tx + Q, s, p) = \operatorname{rdrv}(td, tx, s, p)\).
\end{theorem}
\begin{proof}
The goal follows from \(\operatorname{rdrvstripagreeofwhiteboundarycolagree}(td, tx, Q, B, S, s, p)\).
\end{proof}
\begin{theorem}
\label{ll:paper:rdrv-strip-agree-of-black-boundary-col}
For every natural number \(td\) and natural number \(tx\) and natural number \(Q\) and natural number \(B\) and natural number \(S\) and natural number \(s\) and natural number \(p\), if \(\operatorname{rhbandagree}(tx, Q, 1, B)\), then if for every natural number \(u\), if \(u < S\), then \(\operatorname{rdrv}(td, tx, u, B + 1) = \operatorname{true}\), then if \(s \le S\), then if \(p \le B + 1\), then \(\operatorname{rdrv}(td, tx + Q, s, p) = \operatorname{rdrv}(td, tx, s, p)\).
\end{theorem}
\begin{proof}
The goal follows from \(\operatorname{rdrvstripagreeofblackboundarycol}(td, tx, Q, B, S, s, p)\).
\end{proof}
\begin{theorem}
\label{ll:paper:periodic-maskparity-zero-agree-of-black-boundary-col}
For every natural number \(T\) and natural number \(P\), if \(0 < P\), then if \(\operatorname{centerperiodic}(T, P)\), then for every natural number \(m\) and natural number \(t\) and natural number \(B\) and natural number \(S\) and natural number \(s\), if \(T \le t\), then if \(\operatorname{rhbandagree}(t, m \cdot P, 1, B)\), then if for every natural number \(u\), if \(u < S\), then \(\operatorname{rdrv}(t, t, u, B + 1) = \operatorname{true}\), then if \(s \le S\), then \(\operatorname{maskparity}(t + s + m \cdot P, 0) = \operatorname{maskparity}(t + s, 0)\).
\end{theorem}
\begin{proof}
Assume \(hP\), \(hper\).
The goal follows from \(\operatorname{periodicmaskparityzeroagreeofblackboundarycol}(T, P, hP, hper)\).
\end{proof}
\begin{theorem}
\label{ll:paper:rdrv-subdiag-one-step-of-black-main-diag}
For every natural number \(td\) and natural number \(tx\) and natural number \(k\), if \(\operatorname{rdrv}(td, tx, k + 1, k + 1) = \operatorname{true}\), then \(\operatorname{rdrv}(td, tx, k + 2, k + 1) = \operatorname{not}(\operatorname{rdrv}(td, tx, k + 1, k))\).
\end{theorem}
\begin{proof}
The goal follows from \(\operatorname{rdrvsubdiagonestepofblackmaindiag}(td, tx, k)\).
\end{proof}
\begin{theorem}
\label{ll:paper:rdrv-subdiag-one-step-of-white-main-diag}
For every natural number \(td\) and natural number \(tx\) and natural number \(k\), if \(\operatorname{rdrv}(td, tx, k + 1, k + 1) = \operatorname{false}\), then \(\operatorname{rdrv}(td, tx, k + 2, k + 1) = \operatorname{bxor}(\operatorname{rdrv}(td, tx, k + 1, k), \operatorname{rhcenter}(tx + (k + 1), k + 2))\).
\end{theorem}
\begin{proof}
The goal follows from \(\operatorname{rdrvsubdiagonestepofwhitemaindiag}(td, tx, k)\).
\end{proof}
\begin{theorem}
\label{ll:paper:rdrv-drive-cone-agree-of-main-diag-gate}
For every natural number \(td\) and natural number \(tx2\) and natural number \(tx\) and natural number \(k\) and natural number \(p\), if for every natural number \(u\), if \(1 \le u\), then if \(u \le k\), then \(\operatorname{rdrv}(td, tx2, u, u) = \operatorname{rdrv}(td, tx, u, u)\), then if for every natural number \(u\), if \(1 \le u\), then if \(u \le k\), then if \(\operatorname{rdrv}(td, tx, u, u) = \operatorname{false}\), then \(\operatorname{rdrv}(td, tx2, u, u + 1) = \operatorname{rdrv}(td, tx, u, u + 1)\), then if \(p \le k\), then \(\operatorname{rdrv}(td, tx2, k, p) = \operatorname{rdrv}(td, tx, k, p)\).
\end{theorem}
\begin{proof}
The goal follows from \(\operatorname{rdrvdriveconeagreeofmaindiaggate}(td, tx2, tx, k, p)\).
\end{proof}
\begin{theorem}
\label{ll:paper:rdrv-drive-cone-agree-of-black-main-diag}
For every natural number \(td\) and natural number \(tx2\) and natural number \(tx\) and natural number \(k\) and natural number \(p\), if for every natural number \(u\), if \(1 \le u\), then if \(u \le k\), then \(\operatorname{rdrv}(td, tx2, u, u) = \operatorname{true}\), then if for every natural number \(u\), if \(1 \le u\), then if \(u \le k\), then \(\operatorname{rdrv}(td, tx, u, u) = \operatorname{true}\), then if \(p \le k\), then \(\operatorname{rdrv}(td, tx2, k, p) = \operatorname{rdrv}(td, tx, k, p)\).
\end{theorem}
\begin{proof}
The goal follows from \(\operatorname{rdrvdriveconeagreeofblackmaindiag}(td, tx2, tx, k, p)\).
\end{proof}
\begin{theorem}
\label{ll:paper:rdrv-drive-cone-eq-white-seed-of-black-main-diag}
For every natural number \(td\) and natural number \(tx\) and natural number \(k\) and natural number \(p\), if \(\operatorname{center}(td) = \operatorname{true}\), then if for every natural number \(u\), if \(1 \le u\), then if \(u \le k\), then \(\operatorname{rdrv}(td, tx, u, u) = \operatorname{true}\), then if \(p \le k\), then \(\operatorname{rdrv}(td, tx, k, p) = \operatorname{rdrv}(td, 0, k, p)\).
\end{theorem}
\begin{proof}
The goal follows from \(\operatorname{rdrvdriveconeeqwhiteseedofblackmaindiag}(td, tx, k, p)\).
\end{proof}
\begin{theorem}
\label{ll:paper:rdrv-drive-cone-eq-white-seed-of-black-center-and-rhcenter-diag}
For every natural number \(td\) and natural number \(tx\) and natural number \(k\) and natural number \(p\), if \(\operatorname{center}(td) = \operatorname{true}\), then if \(\operatorname{center}(tx) = \operatorname{true}\), then if for every natural number \(u\), if \(1 \le u\), then if \(u \le k\), then \(\operatorname{rhcenter}(tx + u, u) = \operatorname{true}\), then if \(p \le k\), then \(\operatorname{rdrv}(td, tx, k, p) = \operatorname{rdrv}(td, 0, k, p)\).
\end{theorem}
\begin{proof}
The goal follows from \(\operatorname{rdrvdriveconeeqwhiteseedofblackcenterandrhcenterdiag}(td, tx, k, p)\).
\end{proof}
\begin{theorem}
\label{ll:paper:rhcenter-drive-cone-eq-white-seed-of-black-diag}
For every natural number \(t\) and natural number \(k\) and natural number \(p\), if \(\operatorname{center}(t) = \operatorname{true}\), then if for every natural number \(u\), if \(1 \le u\), then if \(u \le k\), then \(\operatorname{rhcenter}(t + u, u) = \operatorname{true}\), then if \(p \le k\), then \(\operatorname{rhcenter}(t + k, p) = \operatorname{rdrv}(t, 0, k, p)\).
\end{theorem}
\begin{proof}
The goal follows from \(\operatorname{rhcenterdriveconeeqwhiteseedofblackdiag}(t, k, p)\).
\end{proof}
\begin{theorem}
\label{ll:paper:periodic-maskparity-zero-agree-of-main-diag-gate}
For every natural number \(T\) and natural number \(P\), if \(0 < P\), then if \(\operatorname{centerperiodic}(T, P)\), then for every natural number \(m\) and natural number \(t\) and natural number \(s\), if \(T \le t\), then if \(1 \le s\), then if for every natural number \(u\), if \(1 \le u\), then if \(u \le s\), then \(\operatorname{rdrv}(t, t + m \cdot P, u, u) = \operatorname{rdrv}(t, t, u, u)\), then if for every natural number \(u\), if \(1 \le u\), then if \(u \le s\), then if \(\operatorname{rdrv}(t, t, u, u) = \operatorname{false}\), then \(\operatorname{rdrv}(t, t + m \cdot P, u, u + 1) = \operatorname{rdrv}(t, t, u, u + 1)\), then \(\operatorname{maskparity}(t + s + m \cdot P, 0) = \operatorname{maskparity}(t + s, 0)\).
\end{theorem}
\begin{proof}
Assume \(hP\), \(hper\).
The goal follows from \(\operatorname{periodicmaskparityzeroagreeofmaindiaggate}(T, P, hP, hper)\).
\end{proof}
\begin{theorem}
\label{ll:paper:periodic-maskparity-zero-agree-of-black-main-diag}
For every natural number \(T\) and natural number \(P\), if \(0 < P\), then if \(\operatorname{centerperiodic}(T, P)\), then for every natural number \(m\) and natural number \(t\) and natural number \(s\), if \(T \le t\), then if \(1 \le s\), then if for every natural number \(u\), if \(1 \le u\), then if \(u \le s\), then \(\operatorname{rdrv}(t, t + m \cdot P, u, u) = \operatorname{true}\), then if for every natural number \(u\), if \(1 \le u\), then if \(u \le s\), then \(\operatorname{rdrv}(t, t, u, u) = \operatorname{true}\), then \(\operatorname{maskparity}(t + s + m \cdot P, 0) = \operatorname{maskparity}(t + s, 0)\).
\end{theorem}
\begin{proof}
Assume \(hP\), \(hper\).
The goal follows from \(\operatorname{periodicmaskparityzeroagreeofblackmaindiag}(T, P, hP, hper)\).
\end{proof}
\begin{theorem}
\label{ll:paper:not-centerperiodic-of-main-diag-gate-at-zero-diag}
For every natural number \(T\) and natural number \(P\) and natural number \(m\) and natural number \(z\) and natural number \(p\) and natural number \(S\), if \(0 < P\), then if \(0 < m\), then if \(T \le z + p\), then if \(\operatorname{zerodiagonal}(z)\), then if \(\operatorname{doublesat}(m \cdot P, z)\), then if \(\operatorname{settledby}(z + 1, z + (p + 1))\), then if \(\operatorname{settledby}(z + 1 - m \cdot P, z + (p + 1))\), then if \(p + 3 \cdot (m \cdot P) \le S + 1\), then if for every natural number \(u\), if \(1 \le u\), then if \(u \le S + 1\), then \(\operatorname{rdrv}(z + p, z + p + m \cdot P, u, u) = \operatorname{rdrv}(z + p, z + p, u, u)\), then if for every natural number \(u\), if \(1 \le u\), then if \(u \le S + 1\), then if \(\operatorname{rdrv}(z + p, z + p, u, u) = \operatorname{false}\), then \(\operatorname{rdrv}(z + p, z + p + m \cdot P, u, u + 1) = \operatorname{rdrv}(z + p, z + p, u, u + 1)\), then not \(\operatorname{centerperiodic}(T, P)\).
\end{theorem}
\begin{proof}
The goal follows from \(\operatorname{notcenterperiodicofmaindiaggateatzerodiag}(T, P, m, z, p, S)\).
\end{proof}
\begin{theorem}
\label{ll:paper:not-centerperiodic-of-black-main-diag-at-zero-diag}
For every natural number \(T\) and natural number \(P\) and natural number \(m\) and natural number \(z\) and natural number \(p\) and natural number \(S\), if \(0 < P\), then if \(0 < m\), then if \(T \le z + p\), then if \(\operatorname{zerodiagonal}(z)\), then if \(\operatorname{doublesat}(m \cdot P, z)\), then if \(\operatorname{settledby}(z + 1, z + (p + 1))\), then if \(\operatorname{settledby}(z + 1 - m \cdot P, z + (p + 1))\), then if \(p + 3 \cdot (m \cdot P) \le S + 1\), then if for every natural number \(u\), if \(1 \le u\), then if \(u \le S + 1\), then \(\operatorname{rdrv}(z + p, z + p + m \cdot P, u, u) = \operatorname{true}\), then if for every natural number \(u\), if \(1 \le u\), then if \(u \le S + 1\), then \(\operatorname{rdrv}(z + p, z + p, u, u) = \operatorname{true}\), then not \(\operatorname{centerperiodic}(T, P)\).
\end{theorem}
\begin{proof}
The goal follows from \(\operatorname{notcenterperiodicofblackmaindiagatzerodiag}(T, P, m, z, p, S)\).
\end{proof}
\begin{theorem}
\label{ll:paper:not-centerperiodic-of-black-boundary-col-at-zero-diag}
For every natural number \(T\) and natural number \(P\) and natural number \(m\) and natural number \(z\) and natural number \(p\) and natural number \(B\) and natural number \(S\), if \(0 < P\), then if \(0 < m\), then if \(T \le z + (p + 1)\), then if \(\operatorname{zerodiagonal}(z)\), then if \(\operatorname{doublesat}(m \cdot P, z)\), then if \(\operatorname{settledby}(z + 1, z + (p + 1))\), then if \(\operatorname{settledby}(z + 1 - m \cdot P, z + (p + 1))\), then if \(p + 3 \cdot (m \cdot P) \le S + 1\), then if \(\operatorname{rhbandagree}(z + (p + 1), m \cdot P, 1, B)\), then if for every natural number \(u\), if \(u < S\), then \(\operatorname{rdrv}(z + (p + 1), z + (p + 1), u, B + 1) = \operatorname{true}\), then not \(\operatorname{centerperiodic}(T, P)\).
\end{theorem}
\begin{proof}
The goal follows from \(\operatorname{notcenterperiodicofblackboundarycolatzerodiag}(T, P, m, z, p, B, S)\).
\end{proof}
\begin{theorem}
\label{ll:paper:rdrv-main-diag-agree-of-corridor-or-agree}
For every natural number \(td\) and natural number \(tx2\) and natural number \(tx\) and natural number \(k\), if for every natural number \(u\), if \(u < k\), then \((\operatorname{rhcenter}(tx2 + u, u + 1) \operatorname{or} \operatorname{rhcenter}(tx2 + u, u + 2)) = (\operatorname{rhcenter}(tx + u, u + 1) \operatorname{or} \operatorname{rhcenter}(tx + u, u + 2))\), then \(\operatorname{rdrv}(td, tx2, k, k) = \operatorname{rdrv}(td, tx, k, k)\).
\end{theorem}
\begin{proof}
The goal follows from \(\operatorname{rdrvmaindiagagreeofcorridororagree}(td, tx2, tx, k)\).
\end{proof}
\begin{theorem}
\label{ll:paper:rdrv-main-diag-eq-center-of-white-corridor}
For every natural number \(td\) and natural number \(tx\) and natural number \(k\), if for every natural number \(u\), if \(u < k\), then \(\operatorname{rhcenter}(tx + u, u + 1) = \operatorname{false}\) and \(\operatorname{rhcenter}(tx + u, u + 2) = \operatorname{false}\), then \(\operatorname{rdrv}(td, tx, k, k) = \operatorname{center}(td)\).
\end{theorem}
\begin{proof}
The goal follows from \(\operatorname{rdrvmaindiageqcenterofwhitecorridor}(td, tx, k)\).
\end{proof}
\begin{theorem}
\label{ll:paper:rdrv-drive-cone-agree-of-corridor-gate}
For every natural number \(td\) and natural number \(tx2\) and natural number \(tx\) and natural number \(k\) and natural number \(p\), if for every natural number \(u\), if \(u \le k\), then \(\operatorname{rhcenter}(tx2 + u, u + 1) = \operatorname{rhcenter}(tx + u, u + 1)\), then if for every natural number \(u\), if \(u \le k\), then if \(\operatorname{rhcenter}(tx + u, u + 1) = \operatorname{false}\), then \(\operatorname{rhcenter}(tx2 + u, u + 2) = \operatorname{rhcenter}(tx + u, u + 2)\), then if \(p \le k\), then \(\operatorname{rdrv}(td, tx2, k, p) = \operatorname{rdrv}(td, tx, k, p)\).
\end{theorem}
\begin{proof}
The goal follows from \(\operatorname{rdrvdriveconeagreeofcorridorgate}(td, tx2, tx, k, p)\).
\end{proof}
\begin{theorem}
\label{ll:paper:rdrv-drive-cone-agree-of-black-first-superdiag}
For every natural number \(td\) and natural number \(tx2\) and natural number \(tx\) and natural number \(k\) and natural number \(p\), if for every natural number \(u\), if \(u \le k\), then \(\operatorname{rhcenter}(tx2 + u, u + 1) = \operatorname{true}\), then if for every natural number \(u\), if \(u \le k\), then \(\operatorname{rhcenter}(tx + u, u + 1) = \operatorname{true}\), then if \(p \le k\), then \(\operatorname{rdrv}(td, tx2, k, p) = \operatorname{rdrv}(td, tx, k, p)\).
\end{theorem}
\begin{proof}
The goal follows from \(\operatorname{rdrvdriveconeagreeofblackfirstsuperdiag}(td, tx2, tx, k, p)\).
\end{proof}
\begin{theorem}
\label{ll:paper:rdrv-drive-cone-eq-white-seed-of-white-corridor}
For every natural number \(td\) and natural number \(tx\) and natural number \(k\) and natural number \(p\), if for every natural number \(u\), if \(u \le k\), then \(\operatorname{rhcenter}(tx + u, u + 1) = \operatorname{false}\) and \(\operatorname{rhcenter}(tx + u, u + 2) = \operatorname{false}\), then if \(p \le k\), then \(\operatorname{rdrv}(td, tx, k, p) = \operatorname{rdrv}(td, 0, k, p)\).
\end{theorem}
\begin{proof}
The goal follows from \(\operatorname{rdrvdriveconeeqwhiteseedofwhitecorridor}(td, tx, k, p)\).
\end{proof}
\begin{theorem}
\label{ll:paper:rhcenter-drive-cone-eq-white-seed-of-white-corridor}
For every natural number \(t\) and natural number \(k\) and natural number \(p\), if for every natural number \(u\), if \(u \le k\), then \(\operatorname{rhcenter}(t + u, u + 1) = \operatorname{false}\) and \(\operatorname{rhcenter}(t + u, u + 2) = \operatorname{false}\), then if \(p \le k\), then \(\operatorname{rhcenter}(t + k, p) = \operatorname{rdrv}(t, 0, k, p)\).
\end{theorem}
\begin{proof}
The goal follows from \(\operatorname{rhcenterdriveconeeqwhiteseedofwhitecorridor}(t, k, p)\).
\end{proof}
\begin{theorem}
\label{ll:paper:rdrv-scale-eq-center-of-black-main-diag}
For every natural number \(td\) and natural number \(tx\) and natural number \(d\) and natural number \(k\), if \(\operatorname{center}(td) = \operatorname{true}\), then if \(d \le k\), then if for every natural number \(u\), if \(1 \le u\), then if \(u \le d + \operatorname{twopow}(k)\), then \(\operatorname{rdrv}(td, tx, u, u) = \operatorname{true}\), then \(\operatorname{rdrv}(td, tx, d + \operatorname{twopow}(k), \operatorname{twopow}(k)) = \operatorname{center}(td + d)\).
\end{theorem}
\begin{proof}
The goal follows from \(\operatorname{rdrvscaleeqcenterofblackmaindiag}(td, tx, d, k)\).
\end{proof}
\begin{theorem}
\label{ll:paper:rhcenter-scale-eq-center-of-black-diag}
For every natural number \(t\) and natural number \(d\) and natural number \(k\), if \(\operatorname{center}(t) = \operatorname{true}\), then if \(d \le k\), then if for every natural number \(u\), if \(1 \le u\), then if \(u \le d + \operatorname{twopow}(k)\), then \(\operatorname{rhcenter}(t + u, u) = \operatorname{true}\), then \(\operatorname{rhcenter}(t + (d + \operatorname{twopow}(k)), \operatorname{twopow}(k)) = \operatorname{center}(t + d)\).
\end{theorem}
\begin{proof}
The goal follows from \(\operatorname{rhcenterscaleeqcenterofblackdiag}(t, d, k)\).
\end{proof}
\begin{theorem}
\label{ll:paper:rdrv-scale-eq-center-of-white-corridor}
For every natural number \(td\) and natural number \(tx\) and natural number \(d\) and natural number \(k\), if \(d \le k\), then if for every natural number \(u\), if \(u \le d + \operatorname{twopow}(k)\), then \(\operatorname{rhcenter}(tx + u, u + 1) = \operatorname{false}\) and \(\operatorname{rhcenter}(tx + u, u + 2) = \operatorname{false}\), then \(\operatorname{rdrv}(td, tx, d + \operatorname{twopow}(k), \operatorname{twopow}(k)) = \operatorname{center}(td + d)\).
\end{theorem}
\begin{proof}
The goal follows from \(\operatorname{rdrvscaleeqcenterofwhitecorridor}(td, tx, d, k)\).
\end{proof}
\begin{theorem}
\label{ll:paper:rhcenter-scale-eq-center-of-white-corridor}
For every natural number \(t\) and natural number \(d\) and natural number \(k\), if \(d \le k\), then if for every natural number \(u\), if \(u \le d + \operatorname{twopow}(k)\), then \(\operatorname{rhcenter}(t + u, u + 1) = \operatorname{false}\) and \(\operatorname{rhcenter}(t + u, u + 2) = \operatorname{false}\), then \(\operatorname{rhcenter}(t + (d + \operatorname{twopow}(k)), \operatorname{twopow}(k)) = \operatorname{center}(t + d)\).
\end{theorem}
\begin{proof}
The goal follows from \(\operatorname{rhcenterscaleeqcenterofwhitecorridor}(t, d, k)\).
\end{proof}
\begin{theorem}
\label{ll:paper:rdrv-subdiag-period-twopow-of-black-main-diag}
For every natural number \(td\) and natural number \(tx\) and natural number \(d\) and natural number \(k\) and natural number \(u\), if \(\operatorname{center}(td) = \operatorname{true}\), then if for every natural number \(j\), if \(1 \le j\), then if \(j \le k + d + u \cdot \operatorname{twopow}(d)\), then \(\operatorname{rdrv}(td, tx, j, j) = \operatorname{true}\), then \(\operatorname{rdrv}(td, tx, k + d + u \cdot \operatorname{twopow}(d), k + u \cdot \operatorname{twopow}(d)) = \operatorname{rdrv}(td, tx, k + d, k)\).
\end{theorem}
\begin{proof}
The goal follows from \(\operatorname{rdrvsubdiagperiodtwopowofblackmaindiag}(td, tx, d, k, u)\).
\end{proof}
\begin{theorem}
\label{ll:paper:periodic-maskparity-zero-agree-of-corridor-gate}
For every natural number \(T\) and natural number \(P\), if \(0 < P\), then if \(\operatorname{centerperiodic}(T, P)\), then for every natural number \(m\) and natural number \(t\) and natural number \(s\), if \(T \le t\), then if \(1 \le s\), then if for every natural number \(u\), if \(u \le s\), then \(\operatorname{rhcenter}(t + m \cdot P + u, u + 1) = \operatorname{rhcenter}(t + u, u + 1)\), then if for every natural number \(u\), if \(u \le s\), then if \(\operatorname{rhcenter}(t + u, u + 1) = \operatorname{false}\), then \(\operatorname{rhcenter}(t + m \cdot P + u, u + 2) = \operatorname{rhcenter}(t + u, u + 2)\), then \(\operatorname{maskparity}(t + s + m \cdot P, 0) = \operatorname{maskparity}(t + s, 0)\).
\end{theorem}
\begin{proof}
Assume \(hP\), \(hper\).
The goal follows from \(\operatorname{periodicmaskparityzeroagreeofcorridorgate}(T, P, hP, hper)\).
\end{proof}
\begin{theorem}
\label{ll:paper:periodic-maskparity-zero-agree-of-black-first-superdiag}
For every natural number \(T\) and natural number \(P\), if \(0 < P\), then if \(\operatorname{centerperiodic}(T, P)\), then for every natural number \(m\) and natural number \(t\) and natural number \(s\), if \(T \le t\), then if \(1 \le s\), then if for every natural number \(u\), if \(u \le s\), then \(\operatorname{rhcenter}(t + m \cdot P + u, u + 1) = \operatorname{true}\), then if for every natural number \(u\), if \(u \le s\), then \(\operatorname{rhcenter}(t + u, u + 1) = \operatorname{true}\), then \(\operatorname{maskparity}(t + s + m \cdot P, 0) = \operatorname{maskparity}(t + s, 0)\).
\end{theorem}
\begin{proof}
Assume \(hP\), \(hper\).
The goal follows from \(\operatorname{periodicmaskparityzeroagreeofblackfirstsuperdiag}(T, P, hP, hper)\).
\end{proof}
\begin{theorem}
\label{ll:paper:not-centerperiodic-of-corridor-gate-at-zero-diag}
For every natural number \(T\) and natural number \(P\) and natural number \(m\) and natural number \(z\) and natural number \(p\) and natural number \(S\), if \(0 < P\), then if \(0 < m\), then if \(T \le z + p\), then if \(\operatorname{zerodiagonal}(z)\), then if \(\operatorname{doublesat}(m \cdot P, z)\), then if \(\operatorname{settledby}(z + 1, z + (p + 1))\), then if \(\operatorname{settledby}(z + 1 - m \cdot P, z + (p + 1))\), then if \(p + 3 \cdot (m \cdot P) \le S + 1\), then if for every natural number \(u\), if \(u \le S + 1\), then \(\operatorname{rhcenter}(z + p + m \cdot P + u, u + 1) = \operatorname{rhcenter}(z + p + u, u + 1)\), then if for every natural number \(u\), if \(u \le S + 1\), then if \(\operatorname{rhcenter}(z + p + u, u + 1) = \operatorname{false}\), then \(\operatorname{rhcenter}(z + p + m \cdot P + u, u + 2) = \operatorname{rhcenter}(z + p + u, u + 2)\), then not \(\operatorname{centerperiodic}(T, P)\).
\end{theorem}
\begin{proof}
The goal follows from \(\operatorname{notcenterperiodicofcorridorgateatzerodiag}(T, P, m, z, p, S)\).
\end{proof}
\begin{theorem}
\label{ll:paper:not-centerperiodic-of-black-first-superdiag-at-zero-diag}
For every natural number \(T\) and natural number \(P\) and natural number \(m\) and natural number \(z\) and natural number \(p\) and natural number \(S\), if \(0 < P\), then if \(0 < m\), then if \(T \le z + p\), then if \(\operatorname{zerodiagonal}(z)\), then if \(\operatorname{doublesat}(m \cdot P, z)\), then if \(\operatorname{settledby}(z + 1, z + (p + 1))\), then if \(\operatorname{settledby}(z + 1 - m \cdot P, z + (p + 1))\), then if \(p + 3 \cdot (m \cdot P) \le S + 1\), then if for every natural number \(u\), if \(u \le S + 1\), then \(\operatorname{rhcenter}(z + p + m \cdot P + u, u + 1) = \operatorname{true}\), then if for every natural number \(u\), if \(u \le S + 1\), then \(\operatorname{rhcenter}(z + p + u, u + 1) = \operatorname{true}\), then not \(\operatorname{centerperiodic}(T, P)\).
\end{theorem}
\begin{proof}
The goal follows from \(\operatorname{notcenterperiodicofblackfirstsuperdiagatzerodiag}(T, P, m, z, p, S)\).
\end{proof}
\begin{theorem}
\label{ll:paper:rdrv-subdiag-one-eq-true-of-false-true}
For every natural number \(td\) and natural number \(u\), if \(\operatorname{center}(td) = \operatorname{false}\), then if \(\operatorname{center}(td + 1) = \operatorname{true}\), then \(\operatorname{rdrv}(td, 0, u + 1, u) = \operatorname{true}\).
\end{theorem}
\begin{proof}
The goal follows from \(\operatorname{rdrvsubdiagoneeqtrueoffalsetrue}(td, u)\).
\end{proof}
\begin{theorem}
\label{ll:paper:rdrv-white-seed-eq-false-of-white-drive-prefix}
For every natural number \(td\) and natural number \(d\) and natural number \(k\) and natural number \(p\), if for every natural number \(j\), if \(j < d\), then \(\operatorname{center}(td + j) = \operatorname{false}\), then if \(k < p + d\), then \(\operatorname{rdrv}(td, 0, k, p) = \operatorname{false}\).
\end{theorem}
\begin{proof}
The goal follows from \(\operatorname{rdrvwhiteseedeqfalseofwhitedriveprefix}(td, d, k, p)\).
\end{proof}
\begin{theorem}
\label{ll:paper:rdrv-white-seed-shift-of-white-drive-prefix}
For every natural number \(td\) and natural number \(d\) and natural number \(k\) and natural number \(p\), if for every natural number \(j\), if \(j < d\), then \(\operatorname{center}(td + j) = \operatorname{false}\), then \(\operatorname{rdrv}(td, 0, d + k, p) = \operatorname{rdrv}(td + d, 0, k, p)\).
\end{theorem}
\begin{proof}
The goal follows from \(\operatorname{rdrvwhiteseedshiftofwhitedriveprefix}(td, d, k, p)\).
\end{proof}
\begin{theorem}
\label{ll:paper:rdrv-white-seed-subdiag-eq-center-of-white-drive-prefix}
For every natural number \(td\) and natural number \(d\) and natural number \(u\), if for every natural number \(j\), if \(j < d\), then \(\operatorname{center}(td + j) = \operatorname{false}\), then \(\operatorname{rdrv}(td, 0, u + d, u) = \operatorname{center}(td + d)\).
\end{theorem}
\begin{proof}
The goal follows from \(\operatorname{rdrvwhiteseedsubdiageqcenterofwhitedriveprefix}(td, d, u)\).
\end{proof}
\begin{theorem}
\label{ll:paper:rdrv-subcone-agree-of-subdiag-gate}
For every natural number \(td\) and natural number \(tx2\) and natural number \(tx\) and natural number \(d\) and natural number \(k\) and natural number \(p\), if for every natural number \(u\), if \(1 \le u\), then if \(u + d \le k\), then \(\operatorname{rdrv}(td, tx2, u + d, u) = \operatorname{rdrv}(td, tx, u + d, u)\), then if for every natural number \(u\), if \(1 \le u\), then if \(u + d \le k\), then if \(\operatorname{rdrv}(td, tx, u + d, u) = \operatorname{false}\), then \(\operatorname{rdrv}(td, tx2, u + d, u + 1) = \operatorname{rdrv}(td, tx, u + d, u + 1)\), then if \(p + d \le k\), then \(\operatorname{rdrv}(td, tx2, k, p) = \operatorname{rdrv}(td, tx, k, p)\).
\end{theorem}
\begin{proof}
The goal follows from \(\operatorname{rdrvsubconeagreeofsubdiaggate}(td, tx2, tx, d, k, p)\).
\end{proof}
\begin{theorem}
\label{ll:paper:rdrv-subcone-agree-of-black-subdiag}
For every natural number \(td\) and natural number \(tx2\) and natural number \(tx\) and natural number \(d\) and natural number \(k\) and natural number \(p\), if for every natural number \(u\), if \(1 \le u\), then if \(u + d \le k\), then \(\operatorname{rdrv}(td, tx2, u + d, u) = \operatorname{true}\), then if for every natural number \(u\), if \(1 \le u\), then if \(u + d \le k\), then \(\operatorname{rdrv}(td, tx, u + d, u) = \operatorname{true}\), then if \(p + d \le k\), then \(\operatorname{rdrv}(td, tx2, k, p) = \operatorname{rdrv}(td, tx, k, p)\).
\end{theorem}
\begin{proof}
The goal follows from \(\operatorname{rdrvsubconeagreeofblacksubdiag}(td, tx2, tx, d, k, p)\).
\end{proof}
\begin{theorem}
\label{ll:paper:rdrv-subcone-eq-white-seed-of-white-prefix-and-black-subdiag}
For every natural number \(td\) and natural number \(tx\) and natural number \(d\) and natural number \(k\) and natural number \(p\), if for every natural number \(j\), if \(j < d\), then \(\operatorname{center}(td + j) = \operatorname{false}\), then if \(\operatorname{center}(td + d) = \operatorname{true}\), then if for every natural number \(u\), if \(1 \le u\), then if \(u + d \le k\), then \(\operatorname{rdrv}(td, tx, u + d, u) = \operatorname{true}\), then if \(p + d \le k\), then \(\operatorname{rdrv}(td, tx, k, p) = \operatorname{rdrv}(td, 0, k, p)\).
\end{theorem}
\begin{proof}
The goal follows from \(\operatorname{rdrvsubconeeqwhiteseedofwhiteprefixandblacksubdiag}(td, tx, d, k, p)\).
\end{proof}
\begin{theorem}
\label{ll:paper:rdrv-subcone-one-eq-white-seed-of-false-true-and-black-subdiag-one}
For every natural number \(td\) and natural number \(tx\) and natural number \(k\) and natural number \(p\), if \(\operatorname{center}(td) = \operatorname{false}\), then if \(\operatorname{center}(td + 1) = \operatorname{true}\), then if for every natural number \(u\), if \(1 \le u\), then if \(u + 1 \le k\), then \(\operatorname{rdrv}(td, tx, u + 1, u) = \operatorname{true}\), then if \(p + 1 \le k\), then \(\operatorname{rdrv}(td, tx, k, p) = \operatorname{rdrv}(td, 0, k, p)\).
\end{theorem}
\begin{proof}
The goal follows from \(\operatorname{rdrvsubconeoneeqwhiteseedoffalsetrueandblacksubdiagone}(td, tx, k, p)\).
\end{proof}
\begin{theorem}
\label{ll:paper:periodic-maskparity-zero-agree-of-subdiag-gate}
For every natural number \(T\) and natural number \(P\), if \(0 < P\), then if \(\operatorname{centerperiodic}(T, P)\), then for every natural number \(m\) and natural number \(t\) and natural number \(d\) and natural number \(s\), if \(T \le t\), then if \(1 + d \le s\), then if for every natural number \(u\), if \(1 \le u\), then if \(u + d \le s\), then \(\operatorname{rdrv}(t, t + m \cdot P, u + d, u) = \operatorname{rdrv}(t, t, u + d, u)\), then if for every natural number \(u\), if \(1 \le u\), then if \(u + d \le s\), then if \(\operatorname{rdrv}(t, t, u + d, u) = \operatorname{false}\), then \(\operatorname{rdrv}(t, t + m \cdot P, u + d, u + 1) = \operatorname{rdrv}(t, t, u + d, u + 1)\), then \(\operatorname{maskparity}(t + s + m \cdot P, 0) = \operatorname{maskparity}(t + s, 0)\).
\end{theorem}
\begin{proof}
Assume \(hP\), \(hper\).
The goal follows from \(\operatorname{periodicmaskparityzeroagreeofsubdiaggate}(T, P, hP, hper)\).
\end{proof}
\begin{theorem}
\label{ll:paper:periodic-maskparity-zero-agree-of-black-subdiag}
For every natural number \(T\) and natural number \(P\), if \(0 < P\), then if \(\operatorname{centerperiodic}(T, P)\), then for every natural number \(m\) and natural number \(t\) and natural number \(d\) and natural number \(s\), if \(T \le t\), then if \(1 + d \le s\), then if for every natural number \(u\), if \(1 \le u\), then if \(u + d \le s\), then \(\operatorname{rdrv}(t, t + m \cdot P, u + d, u) = \operatorname{true}\), then if for every natural number \(u\), if \(1 \le u\), then if \(u + d \le s\), then \(\operatorname{rdrv}(t, t, u + d, u) = \operatorname{true}\), then \(\operatorname{maskparity}(t + s + m \cdot P, 0) = \operatorname{maskparity}(t + s, 0)\).
\end{theorem}
\begin{proof}
Assume \(hP\), \(hper\).
The goal follows from \(\operatorname{periodicmaskparityzeroagreeofblacksubdiag}(T, P, hP, hper)\).
\end{proof}
\begin{theorem}
\label{ll:paper:not-centerperiodic-of-subdiag-gate-at-zero-diag}
For every natural number \(T\) and natural number \(P\) and natural number \(m\) and natural number \(z\) and natural number \(p\) and natural number \(t\) and natural number \(d\) and natural number \(S\), if \(0 < P\), then if \(0 < m\), then if \(T \le t\), then if \(t + (d + 1) = z + (p + 1)\), then if \(\operatorname{zerodiagonal}(z)\), then if \(\operatorname{doublesat}(m \cdot P, z)\), then if \(\operatorname{settledby}(z + 1, z + (p + 1))\), then if \(\operatorname{settledby}(z + 1 - m \cdot P, z + (p + 1))\), then if \(p + 3 \cdot (m \cdot P) \le S + 1\), then if for every natural number \(u\), if \(1 \le u\), then if \(u + d \le S + (d + 1)\), then \(\operatorname{rdrv}(t, t + m \cdot P, u + d, u) = \operatorname{rdrv}(t, t, u + d, u)\), then if for every natural number \(u\), if \(1 \le u\), then if \(u + d \le S + (d + 1)\), then if \(\operatorname{rdrv}(t, t, u + d, u) = \operatorname{false}\), then \(\operatorname{rdrv}(t, t + m \cdot P, u + d, u + 1) = \operatorname{rdrv}(t, t, u + d, u + 1)\), then not \(\operatorname{centerperiodic}(T, P)\).
\end{theorem}
\begin{proof}
The goal follows from \(\operatorname{notcenterperiodicofsubdiaggateatzerodiag}(T, P, m, z, p, t, d, S)\).
\end{proof}
\begin{theorem}
\label{ll:paper:not-centerperiodic-of-black-subdiag-at-zero-diag}
For every natural number \(T\) and natural number \(P\) and natural number \(m\) and natural number \(z\) and natural number \(p\) and natural number \(t\) and natural number \(d\) and natural number \(S\), if \(0 < P\), then if \(0 < m\), then if \(T \le t\), then if \(t + (d + 1) = z + (p + 1)\), then if \(\operatorname{zerodiagonal}(z)\), then if \(\operatorname{doublesat}(m \cdot P, z)\), then if \(\operatorname{settledby}(z + 1, z + (p + 1))\), then if \(\operatorname{settledby}(z + 1 - m \cdot P, z + (p + 1))\), then if \(p + 3 \cdot (m \cdot P) \le S + 1\), then if for every natural number \(u\), if \(1 \le u\), then if \(u + d \le S + (d + 1)\), then \(\operatorname{rdrv}(t, t + m \cdot P, u + d, u) = \operatorname{true}\), then if for every natural number \(u\), if \(1 \le u\), then if \(u + d \le S + (d + 1)\), then \(\operatorname{rdrv}(t, t, u + d, u) = \operatorname{true}\), then not \(\operatorname{centerperiodic}(T, P)\).
\end{theorem}
\begin{proof}
The goal follows from \(\operatorname{notcenterperiodicofblacksubdiagatzerodiag}(T, P, m, z, p, t, d, S)\).
\end{proof}
\begin{theorem}
\label{ll:paper:rdrv-drive-wavefront-eq-true-of-first-diff-and-white-rhcenter}
For every natural number \(td\) and natural number \(tx\) and natural number \(d\) and natural number \(p\), if for every natural number \(j\), if \(j < d\), then \(\operatorname{center}(td + j) = \operatorname{center}(tx + j)\), then if \(\operatorname{center}(td + d) = \operatorname{not}(\operatorname{center}(tx + d))\), then if \(\operatorname{rhcenter}(tx + (p + d), p) = \operatorname{false}\), then \(\operatorname{rdrv}(td, tx, p + d, p) = \operatorname{true}\).
\end{theorem}
\begin{proof}
The goal follows from \(\operatorname{rdrvdrivewavefronteqtrueoffirstdiffandwhiterhcenter}(td, tx, d, p)\).
\end{proof}
\begin{theorem}
\label{ll:paper:rdrv-drive-wavefront-eq-rhcenter-of-center-prefix-eq}
For every natural number \(td\) and natural number \(tx\) and natural number \(d\) and natural number \(p\), if for every natural number \(j\), if \(j \le d\), then \(\operatorname{center}(td + j) = \operatorname{center}(tx + j)\), then \(\operatorname{rdrv}(td, tx, p + d, p) = \operatorname{rhcenter}(tx + (p + d), p)\).
\end{theorem}
\begin{proof}
The goal follows from \(\operatorname{rdrvdrivewavefronteqrhcenterofcenterprefixeq}(td, tx, d, p)\).
\end{proof}
\begin{theorem}
\label{ll:paper:periodic-rdrv-subdiag-eq-rhcenter}
For every natural number \(T\) and natural number \(P\), if \(0 < P\), then if \(\operatorname{centerperiodic}(T, P)\), then for every natural number \(m\) and natural number \(t\) and natural number \(d\) and natural number \(p\), if \(T \le t\), then \(\operatorname{rdrv}(t, t + m \cdot P, p + d, p) = \operatorname{rhcenter}(t + m \cdot P + (p + d), p)\).
\end{theorem}
\begin{proof}
Assume \(hP\), \(hper\).
The goal follows from \(\operatorname{periodicrdrvsubdiageqrhcenter}(T, P, hP, hper)\).
\end{proof}
\begin{theorem}
\label{ll:paper:rdrv-trapezoid-agree-of-superdiag-gate}
For every natural number \(td\) and natural number \(tx\) and natural number \(Q\) and natural number \(B\) and natural number \(k\) and natural number \(p\), if \(\operatorname{rhbandagree}(tx, Q, 1, B)\), then if for every natural number \(u\), if \(u \le k\), then \(\operatorname{rhcenter}(tx + Q + u, u + (B + 1)) = \operatorname{rhcenter}(tx + u, u + (B + 1))\), then if for every natural number \(u\), if \(u \le k\), then if \(\operatorname{rhcenter}(tx + u, u + (B + 1)) = \operatorname{false}\), then \(\operatorname{rhcenter}(tx + Q + u, u + (B + 2)) = \operatorname{rhcenter}(tx + u, u + (B + 2))\), then if \(p \le k + (B + 1)\), then \(\operatorname{rdrv}(td, tx + Q, k, p) = \operatorname{rdrv}(td, tx, k, p)\).
\end{theorem}
\begin{proof}
The goal follows from \(\operatorname{rdrvtrapezoidagreeofsuperdiaggate}(td, tx, Q, B, k, p)\).
\end{proof}
\begin{theorem}
\label{ll:paper:rdrv-trapezoid-agree-of-black-superdiag}
For every natural number \(td\) and natural number \(tx\) and natural number \(Q\) and natural number \(B\) and natural number \(k\) and natural number \(p\), if \(\operatorname{rhbandagree}(tx, Q, 1, B)\), then if for every natural number \(u\), if \(u \le k\), then \(\operatorname{rhcenter}(tx + Q + u, u + (B + 1)) = \operatorname{true}\), then if for every natural number \(u\), if \(u \le k\), then \(\operatorname{rhcenter}(tx + u, u + (B + 1)) = \operatorname{true}\), then if \(p \le k + (B + 1)\), then \(\operatorname{rdrv}(td, tx + Q, k, p) = \operatorname{rdrv}(td, tx, k, p)\).
\end{theorem}
\begin{proof}
The goal follows from \(\operatorname{rdrvtrapezoidagreeofblacksuperdiag}(td, tx, Q, B, k, p)\).
\end{proof}
\begin{theorem}
\label{ll:paper:rdrv-trapezoid-eq-white-seed-of-white-band-and-corridor}
For every natural number \(td\) and natural number \(tx\) and natural number \(B\) and natural number \(k\) and natural number \(p\), if for every natural number \(i\), if \(i \le B\), then \(\operatorname{rhcenter}(tx, 1 + i) = \operatorname{false}\), then if for every natural number \(u\), if \(u \le k\), then \(\operatorname{rhcenter}(tx + u, u + (B + 1)) = \operatorname{false}\) and \(\operatorname{rhcenter}(tx + u, u + (B + 2)) = \operatorname{false}\), then if \(p \le k + (B + 1)\), then \(\operatorname{rdrv}(td, tx, k, p) = \operatorname{rdrv}(td, 0, k, p)\).
\end{theorem}
\begin{proof}
The goal follows from \(\operatorname{rdrvtrapezoideqwhiteseedofwhitebandandcorridor}(td, tx, B, k, p)\).
\end{proof}
\begin{theorem}
\label{ll:paper:rhcenter-trapezoid-eq-white-seed-of-white-band-and-corridor}
For every natural number \(t\) and natural number \(B\) and natural number \(k\) and natural number \(p\), if for every natural number \(i\), if \(i \le B\), then \(\operatorname{rhcenter}(t, 1 + i) = \operatorname{false}\), then if for every natural number \(u\), if \(u \le k\), then \(\operatorname{rhcenter}(t + u, u + (B + 1)) = \operatorname{false}\) and \(\operatorname{rhcenter}(t + u, u + (B + 2)) = \operatorname{false}\), then if \(p \le k + (B + 1)\), then \(\operatorname{rhcenter}(t + k, p) = \operatorname{rdrv}(t, 0, k, p)\).
\end{theorem}
\begin{proof}
The goal follows from \(\operatorname{rhcentertrapezoideqwhiteseedofwhitebandandcorridor}(t, B, k, p)\).
\end{proof}
\begin{theorem}
\label{ll:paper:periodic-maskparity-zero-agree-of-superdiag-gate}
For every natural number \(T\) and natural number \(P\), if \(0 < P\), then if \(\operatorname{centerperiodic}(T, P)\), then for every natural number \(m\) and natural number \(t\) and natural number \(B\) and natural number \(s\), if \(T \le t\), then if \(\operatorname{rhbandagree}(t, m \cdot P, 1, B)\), then if for every natural number \(u\), if \(u \le s\), then \(\operatorname{rhcenter}(t + m \cdot P + u, u + (B + 1)) = \operatorname{rhcenter}(t + u, u + (B + 1))\), then if for every natural number \(u\), if \(u \le s\), then if \(\operatorname{rhcenter}(t + u, u + (B + 1)) = \operatorname{false}\), then \(\operatorname{rhcenter}(t + m \cdot P + u, u + (B + 2)) = \operatorname{rhcenter}(t + u, u + (B + 2))\), then \(\operatorname{maskparity}(t + s + m \cdot P, 0) = \operatorname{maskparity}(t + s, 0)\).
\end{theorem}
\begin{proof}
Assume \(hP\), \(hper\).
The goal follows from \(\operatorname{periodicmaskparityzeroagreeofsuperdiaggate}(T, P, hP, hper)\).
\end{proof}
\begin{theorem}
\label{ll:paper:periodic-maskparity-zero-agree-of-black-superdiag}
For every natural number \(T\) and natural number \(P\), if \(0 < P\), then if \(\operatorname{centerperiodic}(T, P)\), then for every natural number \(m\) and natural number \(t\) and natural number \(B\) and natural number \(s\), if \(T \le t\), then if \(\operatorname{rhbandagree}(t, m \cdot P, 1, B)\), then if for every natural number \(u\), if \(u \le s\), then \(\operatorname{rhcenter}(t + m \cdot P + u, u + (B + 1)) = \operatorname{true}\), then if for every natural number \(u\), if \(u \le s\), then \(\operatorname{rhcenter}(t + u, u + (B + 1)) = \operatorname{true}\), then \(\operatorname{maskparity}(t + s + m \cdot P, 0) = \operatorname{maskparity}(t + s, 0)\).
\end{theorem}
\begin{proof}
Assume \(hP\), \(hper\).
The goal follows from \(\operatorname{periodicmaskparityzeroagreeofblacksuperdiag}(T, P, hP, hper)\).
\end{proof}
\begin{theorem}
\label{ll:paper:not-centerperiodic-of-superdiag-gate-at-zero-diag}
For every natural number \(T\) and natural number \(P\) and natural number \(m\) and natural number \(z\) and natural number \(p\) and natural number \(B\) and natural number \(S\), if \(0 < P\), then if \(0 < m\), then if \(T \le z + (p + 1)\), then if \(\operatorname{zerodiagonal}(z)\), then if \(\operatorname{doublesat}(m \cdot P, z)\), then if \(\operatorname{settledby}(z + 1, z + (p + 1))\), then if \(\operatorname{settledby}(z + 1 - m \cdot P, z + (p + 1))\), then if \(p + 3 \cdot (m \cdot P) \le S + 1\), then if \(\operatorname{rhbandagree}(z + (p + 1), m \cdot P, 1, B)\), then if for every natural number \(u\), if \(u \le S\), then \(\operatorname{rhcenter}(z + (p + 1) + m \cdot P + u, u + (B + 1)) = \operatorname{rhcenter}(z + (p + 1) + u, u + (B + 1))\), then if for every natural number \(u\), if \(u \le S\), then if \(\operatorname{rhcenter}(z + (p + 1) + u, u + (B + 1)) = \operatorname{false}\), then \(\operatorname{rhcenter}(z + (p + 1) + m \cdot P + u, u + (B + 2)) = \operatorname{rhcenter}(z + (p + 1) + u, u + (B + 2))\), then not \(\operatorname{centerperiodic}(T, P)\).
\end{theorem}
\begin{proof}
The goal follows from \(\operatorname{notcenterperiodicofsuperdiaggateatzerodiag}(T, P, m, z, p, B, S)\).
\end{proof}
\begin{theorem}
\label{ll:paper:not-centerperiodic-of-black-superdiag-at-zero-diag}
For every natural number \(T\) and natural number \(P\) and natural number \(m\) and natural number \(z\) and natural number \(p\) and natural number \(B\) and natural number \(S\), if \(0 < P\), then if \(0 < m\), then if \(T \le z + (p + 1)\), then if \(\operatorname{zerodiagonal}(z)\), then if \(\operatorname{doublesat}(m \cdot P, z)\), then if \(\operatorname{settledby}(z + 1, z + (p + 1))\), then if \(\operatorname{settledby}(z + 1 - m \cdot P, z + (p + 1))\), then if \(p + 3 \cdot (m \cdot P) \le S + 1\), then if \(\operatorname{rhbandagree}(z + (p + 1), m \cdot P, 1, B)\), then if for every natural number \(u\), if \(u \le S\), then \(\operatorname{rhcenter}(z + (p + 1) + m \cdot P + u, u + (B + 1)) = \operatorname{true}\), then if for every natural number \(u\), if \(u \le S\), then \(\operatorname{rhcenter}(z + (p + 1) + u, u + (B + 1)) = \operatorname{true}\), then not \(\operatorname{centerperiodic}(T, P)\).
\end{theorem}
\begin{proof}
The goal follows from \(\operatorname{notcenterperiodicofblacksuperdiagatzerodiag}(T, P, m, z, p, B, S)\).
\end{proof}
\begin{theorem}
\label{ll:paper:periodic-maskparity-zero-agree-of-corridor-gate-all}
For every natural number \(T\) and natural number \(P\), if \(0 < P\), then if \(\operatorname{centerperiodic}(T, P)\), then for every natural number \(m\) and natural number \(t\) and natural number \(s\), if \(T \le t\), then if for every natural number \(u\), if \(u \le s\), then \(\operatorname{rhcenter}(t + m \cdot P + u, u + 1) = \operatorname{rhcenter}(t + u, u + 1)\), then if for every natural number \(u\), if \(u \le s\), then if \(\operatorname{rhcenter}(t + u, u + 1) = \operatorname{false}\), then \(\operatorname{rhcenter}(t + m \cdot P + u, u + 2) = \operatorname{rhcenter}(t + u, u + 2)\), then \(\operatorname{maskparity}(t + s + m \cdot P, 0) = \operatorname{maskparity}(t + s, 0)\).
\end{theorem}
\begin{proof}
Assume \(hP\), \(hper\).
The goal follows from \(\operatorname{periodicmaskparityzeroagreeofcorridorgateall}(T, P, hP, hper)\).
\end{proof}
\begin{theorem}
\label{ll:paper:periodic-maskparity-zero-agree-of-black-first-superdiag-all}
For every natural number \(T\) and natural number \(P\), if \(0 < P\), then if \(\operatorname{centerperiodic}(T, P)\), then for every natural number \(m\) and natural number \(t\) and natural number \(s\), if \(T \le t\), then if for every natural number \(u\), if \(u \le s\), then \(\operatorname{rhcenter}(t + m \cdot P + u, u + 1) = \operatorname{true}\), then if for every natural number \(u\), if \(u \le s\), then \(\operatorname{rhcenter}(t + u, u + 1) = \operatorname{true}\), then \(\operatorname{maskparity}(t + s + m \cdot P, 0) = \operatorname{maskparity}(t + s, 0)\).
\end{theorem}
\begin{proof}
Assume \(hP\), \(hper\).
The goal follows from \(\operatorname{periodicmaskparityzeroagreeofblackfirstsuperdiagall}(T, P, hP, hper)\).
\end{proof}
\begin{theorem}
\label{ll:paper:not-centerperiodic-of-corridor-gate-at-zero-diag-direct}
For every natural number \(T\) and natural number \(P\) and natural number \(m\) and natural number \(z\) and natural number \(p\) and natural number \(S\), if \(0 < P\), then if \(0 < m\), then if \(T \le z + (p + 1)\), then if \(\operatorname{zerodiagonal}(z)\), then if \(\operatorname{doublesat}(m \cdot P, z)\), then if \(\operatorname{settledby}(z + 1, z + (p + 1))\), then if \(\operatorname{settledby}(z + 1 - m \cdot P, z + (p + 1))\), then if \(p + 3 \cdot (m \cdot P) \le S + 1\), then if for every natural number \(u\), if \(u \le S\), then \(\operatorname{rhcenter}(z + (p + 1) + m \cdot P + u, u + 1) = \operatorname{rhcenter}(z + (p + 1) + u, u + 1)\), then if for every natural number \(u\), if \(u \le S\), then if \(\operatorname{rhcenter}(z + (p + 1) + u, u + 1) = \operatorname{false}\), then \(\operatorname{rhcenter}(z + (p + 1) + m \cdot P + u, u + 2) = \operatorname{rhcenter}(z + (p + 1) + u, u + 2)\), then not \(\operatorname{centerperiodic}(T, P)\).
\end{theorem}
\begin{proof}
The goal follows from \(\operatorname{notcenterperiodicofcorridorgateatzerodiagdirect}(T, P, m, z, p, S)\).
\end{proof}
\begin{theorem}
\label{ll:paper:not-centerperiodic-of-black-first-superdiag-at-zero-diag-direct}
For every natural number \(T\) and natural number \(P\) and natural number \(m\) and natural number \(z\) and natural number \(p\) and natural number \(S\), if \(0 < P\), then if \(0 < m\), then if \(T \le z + (p + 1)\), then if \(\operatorname{zerodiagonal}(z)\), then if \(\operatorname{doublesat}(m \cdot P, z)\), then if \(\operatorname{settledby}(z + 1, z + (p + 1))\), then if \(\operatorname{settledby}(z + 1 - m \cdot P, z + (p + 1))\), then if \(p + 3 \cdot (m \cdot P) \le S + 1\), then if for every natural number \(u\), if \(u \le S\), then \(\operatorname{rhcenter}(z + (p + 1) + m \cdot P + u, u + 1) = \operatorname{true}\), then if for every natural number \(u\), if \(u \le S\), then \(\operatorname{rhcenter}(z + (p + 1) + u, u + 1) = \operatorname{true}\), then not \(\operatorname{centerperiodic}(T, P)\).
\end{theorem}
\begin{proof}
The goal follows from \(\operatorname{notcenterperiodicofblackfirstsuperdiagatzerodiagdirect}(T, P, m, z, p, S)\).
\end{proof}
\begin{theorem}
\label{ll:paper:not-centerperiodic-of-corridor-gate-sub-twopow}
For every natural number \(T\) and natural number \(P\) and natural number \(m\) and natural number \(k\) and natural number \(S\), if \(0 < P\), then if \(0 < m\), then if \(m \cdot P \le \operatorname{twopow}(k)\), then if \(T + m \cdot P \le \operatorname{twopow}(k)\), then if \(\operatorname{twopow}(k) \le S + 1\), then if for every natural number \(u\), if \(u \le S\), then \(\operatorname{rhcenter}(\operatorname{twopow}(k) - m \cdot P + m \cdot P + u, u + 1) = \operatorname{rhcenter}(\operatorname{twopow}(k) - m \cdot P + u, u + 1)\), then if for every natural number \(u\), if \(u \le S\), then if \(\operatorname{rhcenter}(\operatorname{twopow}(k) - m \cdot P + u, u + 1) = \operatorname{false}\), then \(\operatorname{rhcenter}(\operatorname{twopow}(k) - m \cdot P + m \cdot P + u, u + 2) = \operatorname{rhcenter}(\operatorname{twopow}(k) - m \cdot P + u, u + 2)\), then not \(\operatorname{centerperiodic}(T, P)\).
\end{theorem}
\begin{proof}
The goal follows from \(\operatorname{notcenterperiodicofcorridorgatesubtwopow}(T, P, m, k, S)\).
\end{proof}
\begin{theorem}
\label{ll:paper:not-centerperiodic-of-black-first-superdiag-sub-twopow}
For every natural number \(T\) and natural number \(P\) and natural number \(m\) and natural number \(k\) and natural number \(S\), if \(0 < P\), then if \(0 < m\), then if \(m \cdot P \le \operatorname{twopow}(k)\), then if \(T + m \cdot P \le \operatorname{twopow}(k)\), then if \(\operatorname{twopow}(k) \le S + 1\), then if for every natural number \(u\), if \(u \le S\), then \(\operatorname{rhcenter}(\operatorname{twopow}(k) - m \cdot P + m \cdot P + u, u + 1) = \operatorname{true}\), then if for every natural number \(u\), if \(u \le S\), then \(\operatorname{rhcenter}(\operatorname{twopow}(k) - m \cdot P + u, u + 1) = \operatorname{true}\), then not \(\operatorname{centerperiodic}(T, P)\).
\end{theorem}
\begin{proof}
The goal follows from \(\operatorname{notcenterperiodicofblackfirstsuperdiagsubtwopow}(T, P, m, k, S)\).
\end{proof}
\begin{theorem}
\label{ll:paper:not-centerperiodic-of-superdiag-gate-sub-twopow}
For every natural number \(T\) and natural number \(P\) and natural number \(m\) and natural number \(k\) and natural number \(B\) and natural number \(S\), if \(0 < P\), then if \(0 < m\), then if \(m \cdot P \le \operatorname{twopow}(k)\), then if \(T + m \cdot P \le \operatorname{twopow}(k)\), then if \(\operatorname{twopow}(k) \le S + 1\), then if \(\operatorname{rhbandagree}(\operatorname{twopow}(k) - m \cdot P, m \cdot P, 1, B)\), then if for every natural number \(u\), if \(u \le S\), then \(\operatorname{rhcenter}(\operatorname{twopow}(k) - m \cdot P + m \cdot P + u, u + (B + 1)) = \operatorname{rhcenter}(\operatorname{twopow}(k) - m \cdot P + u, u + (B + 1))\), then if for every natural number \(u\), if \(u \le S\), then if \(\operatorname{rhcenter}(\operatorname{twopow}(k) - m \cdot P + u, u + (B + 1)) = \operatorname{false}\), then \(\operatorname{rhcenter}(\operatorname{twopow}(k) - m \cdot P + m \cdot P + u, u + (B + 2)) = \operatorname{rhcenter}(\operatorname{twopow}(k) - m \cdot P + u, u + (B + 2))\), then not \(\operatorname{centerperiodic}(T, P)\).
\end{theorem}
\begin{proof}
The goal follows from \(\operatorname{notcenterperiodicofsuperdiaggatesubtwopow}(T, P, m, k, B, S)\).
\end{proof}
\begin{theorem}
\label{ll:paper:rdrv-black-stays-step}
For every natural number \(td\) and natural number \(tx\) and natural number \(k\) and natural number \(p\), if \(\operatorname{rdrv}(td, tx, k, p) = \operatorname{false}\), then if \(\operatorname{rdrv}(td, tx, k, p + 1) = \operatorname{true}\), then \(\operatorname{rdrv}(td, tx, k + 1, p + 1) = \operatorname{true}\).
\end{theorem}
\begin{proof}
The goal follows from \(\operatorname{rdrvblackstaysstep}(td, tx, k, p)\).
\end{proof}
\begin{theorem}
\label{ll:paper:rdrv-black-stays}
For every natural number \(td\) and natural number \(tx\) and natural number \(a\) and natural number \(p\) and natural number \(n\) and natural number \(k\), if for every natural number \(s\), if \(s < n\), then \(\operatorname{rdrv}(td, tx, a + s, p) = \operatorname{false}\), then if \(\operatorname{rdrv}(td, tx, a, p + 1) = \operatorname{true}\), then if \(k \le n\), then \(\operatorname{rdrv}(td, tx, a + k, p + 1) = \operatorname{true}\).
\end{theorem}
\begin{proof}
The goal follows from \(\operatorname{rdrvblackstays}(td, tx, a, p, n, k)\).
\end{proof}
\begin{theorem}
\label{ll:paper:rdrv-col-one-black-stays-of-white-drive}
For every natural number \(td\) and natural number \(tx\) and natural number \(a\) and natural number \(r\) and natural number \(k\), if for every natural number \(s\), if \(s < r\), then \(\operatorname{center}(td + (a + s)) = \operatorname{false}\), then if \(\operatorname{rdrv}(td, tx, a, 1) = \operatorname{true}\), then if \(k \le r\), then \(\operatorname{rdrv}(td, tx, a + k, 1) = \operatorname{true}\).
\end{theorem}
\begin{proof}
The goal follows from \(\operatorname{rdrvcoloneblackstaysofwhitedrive}(td, tx, a, r, k)\).
\end{proof}
\begin{theorem}
\label{ll:paper:rdrv-col-one-flip-after-white-drive}
For every natural number \(td\) and natural number \(tx\) and natural number \(a\) and natural number \(r\), if for every natural number \(s\), if \(s < r\), then \(\operatorname{center}(td + (a + s)) = \operatorname{false}\), then if \(\operatorname{center}(td + (a + r)) = \operatorname{true}\), then if \(\operatorname{rdrv}(td, tx, a, 1) = \operatorname{true}\), then \(\operatorname{rdrv}(td, tx, a + (r + 1), 1) = \operatorname{false}\).
\end{theorem}
\begin{proof}
The goal follows from \(\operatorname{rdrvcoloneflipafterwhitedrive}(td, tx, a, r)\).
\end{proof}
\begin{theorem}
\label{ll:paper:rdrv-black-reaches-boundary}
For every natural number \(td\) and natural number \(tx\) and natural number \(a\) and natural number \(p\) and natural number \(n\), if for every natural number \(s\), if \(s < n\), then \(\operatorname{rdrv}(td, tx, a + s, p) = \operatorname{false}\), then if \(\operatorname{rdrv}(td, tx, a, p + (n + 1)) = \operatorname{true}\), then \(\operatorname{rdrv}(td, tx, a + n, p + 1) = \operatorname{true}\).
\end{theorem}
\begin{proof}
The goal follows from \(\operatorname{rdrvblackreachesboundary}(td, tx, a, p, n)\).
\end{proof}
\begin{theorem}
\label{ll:paper:rdrv-col-one-black-of-white-drive-and-black-cell}
For every natural number \(td\) and natural number \(tx\) and natural number \(a\) and natural number \(n\) and natural number \(r\) and natural number \(k\), if for every natural number \(s\), if \(s < r\), then \(\operatorname{center}(td + (a + s)) = \operatorname{false}\), then if \(\operatorname{rdrv}(td, tx, a, n + 1) = \operatorname{true}\), then if \(n \le k\), then if \(k \le r\), then \(\operatorname{rdrv}(td, tx, a + k, 1) = \operatorname{true}\).
\end{theorem}
\begin{proof}
The goal follows from \(\operatorname{rdrvcoloneblackofwhitedriveandblackcell}(td, tx, a, n, r, k)\).
\end{proof}
\begin{theorem}
\label{ll:paper:rdrv-col-one-black-from-init-of-white-drive}
For every natural number \(td\) and natural number \(tx\) and natural number \(n\) and natural number \(r\) and natural number \(k\), if for every natural number \(s\), if \(s < r\), then \(\operatorname{center}(td + s) = \operatorname{false}\), then if \(\operatorname{rhcenter}(tx, n + 1) = \operatorname{true}\), then if \(n \le k\), then if \(k \le r\), then \(\operatorname{rdrv}(td, tx, k, 1) = \operatorname{true}\).
\end{theorem}
\begin{proof}
The goal follows from \(\operatorname{rdrvcoloneblackfrominitofwhitedrive}(td, tx, n, r, k)\).
\end{proof}
\begin{theorem}
\label{ll:paper:rdrv-col-one-black-from-right-edge-of-white-drive}
For every natural number \(td\) and natural number \(tx\) and natural number \(r\) and natural number \(k\), if \(1 \le tx\), then if for every natural number \(s\), if \(s < r\), then \(\operatorname{center}(td + s) = \operatorname{false}\), then if \(tx \le k + 1\), then if \(k \le r\), then \(\operatorname{rdrv}(td, tx, k, 1) = \operatorname{true}\).
\end{theorem}
\begin{proof}
The goal follows from \(\operatorname{rdrvcoloneblackfromrightedgeofwhitedrive}(td, tx, r, k)\).
\end{proof}
\begin{theorem}
\label{ll:paper:rdrv-col-one-agree-of-white-drive-and-black-cells}
For every natural number \(td\) and natural number \(tx2\) and natural number \(tx\) and natural number \(a\) and natural number \(n2\) and natural number \(n1\) and natural number \(r\) and natural number \(k\), if for every natural number \(s\), if \(s < r\), then \(\operatorname{center}(td + (a + s)) = \operatorname{false}\), then if \(\operatorname{rdrv}(td, tx2, a, n2 + 1) = \operatorname{true}\), then if \(\operatorname{rdrv}(td, tx, a, n1 + 1) = \operatorname{true}\), then if \(n2 \le k\), then if \(n1 \le k\), then if \(k \le r\), then \(\operatorname{rdrv}(td, tx2, a + k, 1) = \operatorname{rdrv}(td, tx, a + k, 1)\).
\end{theorem}
\begin{proof}
The goal follows from \(\operatorname{rdrvcoloneagreeofwhitedriveandblackcells}(td, tx2, tx, a, n2, n1, r, k)\).
\end{proof}
\begin{theorem}
\label{ll:paper:rdrv-col-one-agree-after-white-drive-and-black-cells}
For every natural number \(td\) and natural number \(tx2\) and natural number \(tx\) and natural number \(a\) and natural number \(n2\) and natural number \(n1\) and natural number \(r\), if for every natural number \(s\), if \(s < r\), then \(\operatorname{center}(td + (a + s)) = \operatorname{false}\), then if \(\operatorname{rdrv}(td, tx2, a, n2 + 1) = \operatorname{true}\), then if \(\operatorname{rdrv}(td, tx, a, n1 + 1) = \operatorname{true}\), then if \(n2 \le r\), then if \(n1 \le r\), then \(\operatorname{rdrv}(td, tx2, a + (r + 1), 1) = \operatorname{rdrv}(td, tx, a + (r + 1), 1)\).
\end{theorem}
\begin{proof}
The goal follows from \(\operatorname{rdrvcoloneagreeafterwhitedriveandblackcells}(td, tx2, tx, a, n2, n1, r)\).
\end{proof}
\begin{theorem}
\label{ll:paper:periodic-maskparity-zero-agree-of-white-drive-and-black-cells}
For every natural number \(T\) and natural number \(P\), if \(0 < P\), then if \(\operatorname{centerperiodic}(T, P)\), then for every natural number \(m\) and natural number \(t\) and natural number \(a\) and natural number \(n2\) and natural number \(n1\) and natural number \(r\) and natural number \(k\), if \(T \le t\), then if for every natural number \(s\), if \(s < r\), then \(\operatorname{center}(t + (a + s)) = \operatorname{false}\), then if \(\operatorname{rdrv}(t, t + m \cdot P, a, n2 + 1) = \operatorname{true}\), then if \(\operatorname{rdrv}(t, t, a, n1 + 1) = \operatorname{true}\), then if \(n2 \le k\), then if \(n1 \le k\), then if \(k \le r\), then \(\operatorname{maskparity}(t + (a + k) + m \cdot P, 0) = \operatorname{maskparity}(t + (a + k), 0)\).
\end{theorem}
\begin{proof}
Assume \(hP\), \(hper\).
The goal follows from \(\operatorname{periodicmaskparityzeroagreeofwhitedriveandblackcells}(T, P, hP, hper)\).
\end{proof}
\begin{theorem}
\label{ll:paper:periodic-maskparity-zero-agree-from-init-of-white-drive}
For every natural number \(T\) and natural number \(P\), if \(0 < P\), then if \(\operatorname{centerperiodic}(T, P)\), then for every natural number \(m\) and natural number \(t\) and natural number \(n2\) and natural number \(n1\) and natural number \(r\) and natural number \(k\), if \(T \le t\), then if for every natural number \(s\), if \(s < r\), then \(\operatorname{center}(t + s) = \operatorname{false}\), then if \(\operatorname{rhcenter}(t + m \cdot P, n2 + 1) = \operatorname{true}\), then if \(\operatorname{rhcenter}(t, n1 + 1) = \operatorname{true}\), then if \(n2 \le k\), then if \(n1 \le k\), then if \(k \le r\), then \(\operatorname{maskparity}(t + k + m \cdot P, 0) = \operatorname{maskparity}(t + k, 0)\).
\end{theorem}
\begin{proof}
Assume \(hP\), \(hper\).
The goal follows from \(\operatorname{periodicmaskparityzeroagreefrominitofwhitedrive}(T, P, hP, hper)\).
\end{proof}
\begin{theorem}
\label{ll:paper:rdrv-white-prefix-forced-of-col-one-white}
For every natural number \(td\) and natural number \(tx\) and natural number \(a\) and natural number \(r\) and natural number \(q\), if for every natural number \(s\), if \(s < r\), then \(\operatorname{center}(td + (a + s)) = \operatorname{false}\), then if \(\operatorname{rdrv}(td, tx, a + (q - 1), 1) = \operatorname{false}\), then if \(1 \le q\), then if \(q \le r + 1\), then \(\operatorname{rdrv}(td, tx, a, q) = \operatorname{false}\).
\end{theorem}
\begin{proof}
The goal follows from \(\operatorname{rdrvwhiteprefixforcedofcolonewhite}(td, tx, a, r, q)\).
\end{proof}
\begin{theorem}
\label{ll:paper:rhcenter-white-prefix-forced-of-col-one-white}
For every natural number \(t\) and natural number \(r\) and natural number \(q\), if for every natural number \(s\), if \(s < r\), then \(\operatorname{center}(t + s) = \operatorname{false}\), then if \(\operatorname{rhcenter}(t + (q - 1), 1) = \operatorname{false}\), then if \(1 \le q\), then if \(q \le r + 1\), then \(\operatorname{rhcenter}(t, q) = \operatorname{false}\).
\end{theorem}
\begin{proof}
The goal follows from \(\operatorname{rhcenterwhiteprefixforcedofcolonewhite}(t, r, q)\).
\end{proof}
\begin{theorem}
\label{ll:paper:rhcenter-white-band-forced-of-col-one-white-run}
For every natural number \(t\) and natural number \(r\) and natural number \(q\), if for every natural number \(s\), if \(s < r\), then \(\operatorname{center}(t + s) = \operatorname{false}\), then if for every natural number \(s\), if \(s \le r\), then \(\operatorname{rhcenter}(t + s, 1) = \operatorname{false}\), then if \(1 \le q\), then if \(q \le r + 1\), then \(\operatorname{rhcenter}(t, q) = \operatorname{false}\).
\end{theorem}
\begin{proof}
The goal follows from \(\operatorname{rhcenterwhitebandforcedofcolonewhiterun}(t, r, q)\).
\end{proof}
\begin{theorem}
\label{ll:paper:rdrv-subdiag-succ-two-step}
For every natural number \(td\) and natural number \(tx\) and natural number \(c\) and natural number \(u\), \(\operatorname{rdrv}(td, tx, u + 1 + (c + 2), u + 1) = \operatorname{bxor}(\operatorname{rdrv}(td, tx, u + (c + 2), u), \operatorname{rdrv}(td, tx, u + 1 + (c + 1), u + 1) \operatorname{or} \operatorname{rdrv}(td, tx, u + 2 + c, u + 2))\).
\end{theorem}
\begin{proof}
The goal follows from \(\operatorname{rdrvsubdiagsucctwostep}(td, tx, c, u)\).
\end{proof}
\begin{theorem}
\label{ll:paper:rdrv-subdiag-agree-of-prev-or-agree}
For every natural number \(td\) and natural number \(tx2\) and natural number \(tx\) and natural number \(c\) and natural number \(k\), if for every natural number \(u\), if \(u < k\), then \((\operatorname{rdrv}(td, tx2, u + 1 + (c + 1), u + 1) \operatorname{or} \operatorname{rdrv}(td, tx2, u + 2 + c, u + 2)) = (\operatorname{rdrv}(td, tx, u + 1 + (c + 1), u + 1) \operatorname{or} \operatorname{rdrv}(td, tx, u + 2 + c, u + 2))\), then \(\operatorname{rdrv}(td, tx2, k + (c + 2), k) = \operatorname{rdrv}(td, tx, k + (c + 2), k)\).
\end{theorem}
\begin{proof}
The goal follows from \(\operatorname{rdrvsubdiagagreeofprevoragree}(td, tx2, tx, c, k)\).
\end{proof}
\begin{theorem}
\label{ll:paper:rdrv-subdiag-agree-of-two-prev-subdiags}
For every natural number \(td\) and natural number \(tx2\) and natural number \(tx\) and natural number \(c\) and natural number \(k\), if for every natural number \(u\), if \(u \le k\), then \(\operatorname{rdrv}(td, tx2, u + (c + 1), u) = \operatorname{rdrv}(td, tx, u + (c + 1), u)\), then if for every natural number \(u\), if \(u \le k + 1\), then \(\operatorname{rdrv}(td, tx2, u + c, u) = \operatorname{rdrv}(td, tx, u + c, u)\), then \(\operatorname{rdrv}(td, tx2, k + (c + 2), k) = \operatorname{rdrv}(td, tx, k + (c + 2), k)\).
\end{theorem}
\begin{proof}
The goal follows from \(\operatorname{rdrvsubdiagagreeoftwoprevsubdiags}(td, tx2, tx, c, k)\).
\end{proof}
\begin{theorem}
\label{ll:paper:rdrv-subcone-agree-of-two-subdiags}
For every natural number \(td\) and natural number \(tx2\) and natural number \(tx\) and natural number \(c\) and natural number \(k\) and natural number \(p\), if for every natural number \(u\), if \(u \le k\), then \(\operatorname{rdrv}(td, tx2, u + (c + 1), u) = \operatorname{rdrv}(td, tx, u + (c + 1), u)\), then if for every natural number \(u\), if \(u \le k\), then \(\operatorname{rdrv}(td, tx2, u + 1 + c, u + 1) = \operatorname{rdrv}(td, tx, u + 1 + c, u + 1)\), then if \(p + (c + 1) \le k\), then \(\operatorname{rdrv}(td, tx2, k, p) = \operatorname{rdrv}(td, tx, k, p)\).
\end{theorem}
\begin{proof}
The goal follows from \(\operatorname{rdrvsubconeagreeoftwosubdiags}(td, tx2, tx, c, k, p)\).
\end{proof}
\begin{theorem}
\label{ll:paper:rdrv-subdiag-eq-white-seed-of-prev-or-eq}
For every natural number \(td\) and natural number \(tx\) and natural number \(c\) and natural number \(k\), if for every natural number \(u\), if \(u < k\), then \((\operatorname{rdrv}(td, tx, u + 1 + (c + 1), u + 1) \operatorname{or} \operatorname{rdrv}(td, tx, u + 2 + c, u + 2)) = (\operatorname{rdrv}(td, 0, u + 1 + (c + 1), u + 1) \operatorname{or} \operatorname{rdrv}(td, 0, u + 2 + c, u + 2))\), then \(\operatorname{rdrv}(td, tx, k + (c + 2), k) = \operatorname{rdrv}(td, 0, k + (c + 2), k)\).
\end{theorem}
\begin{proof}
The goal follows from \(\operatorname{rdrvsubdiageqwhiteseedofprevoreq}(td, tx, c, k)\).
\end{proof}
\begin{theorem}
\label{ll:paper:rdrv-subdiag-accum-shift-xor}
For every natural number \(td\) and natural number \(tx\) and natural number \(c\) and natural number \(a\) and natural number \(H\) and natural number \(k\), if for every natural number \(u\), if \(u < k\), then \((\operatorname{rdrv}(td, tx, a + H + (u + 1) + (c + 1), a + H + (u + 1)) \operatorname{or} \operatorname{rdrv}(td, tx, a + H + (u + 2) + c, a + H + (u + 2))) = (\operatorname{rdrv}(td, tx, a + (u + 1) + (c + 1), a + (u + 1)) \operatorname{or} \operatorname{rdrv}(td, tx, a + (u + 2) + c, a + (u + 2)))\), then \(\operatorname{bxor}(\operatorname{rdrv}(td, tx, a + H + k + (c + 2), a + H + k), \operatorname{rdrv}(td, tx, a + k + (c + 2), a + k)) = \operatorname{bxor}(\operatorname{rdrv}(td, tx, a + H + (c + 2), a + H), \operatorname{rdrv}(td, tx, a + (c + 2), a))\).
\end{theorem}
\begin{proof}
The goal follows from \(\operatorname{rdrvsubdiagaccumshiftxor}(td, tx, c, a, H, k)\).
\end{proof}
\begin{theorem}
\label{ll:paper:rdrv-subdiag-period-double-of-prev-or-period}
For every natural number \(td\) and natural number \(tx\) and natural number \(c\) and natural number \(a\) and natural number \(H\) and natural number \(k\), if for every natural number \(u\), if \(u < H + k\), then \((\operatorname{rdrv}(td, tx, a + H + (u + 1) + (c + 1), a + H + (u + 1)) \operatorname{or} \operatorname{rdrv}(td, tx, a + H + (u + 2) + c, a + H + (u + 2))) = (\operatorname{rdrv}(td, tx, a + (u + 1) + (c + 1), a + (u + 1)) \operatorname{or} \operatorname{rdrv}(td, tx, a + (u + 2) + c, a + (u + 2)))\), then \(\operatorname{rdrv}(td, tx, a + 2 \cdot H + k + (c + 2), a + 2 \cdot H + k) = \operatorname{rdrv}(td, tx, a + k + (c + 2), a + k)\).
\end{theorem}
\begin{proof}
The goal follows from \(\operatorname{rdrvsubdiagperioddoubleofprevorperiod}(td, tx, c, a, H, k)\).
\end{proof}
\begin{theorem}
\label{ll:paper:rdrv-subdiag-period-double-of-two-prev-periods}
For every natural number \(td\) and natural number \(tx\) and natural number \(c\) and natural number \(a\) and natural number \(H\) and natural number \(k\), if for every natural number \(u\), if \(u \le H + k\), then \(\operatorname{rdrv}(td, tx, a + H + u + (c + 1), a + H + u) = \operatorname{rdrv}(td, tx, a + u + (c + 1), a + u)\), then if for every natural number \(u\), if \(u \le H + k + 1\), then \(\operatorname{rdrv}(td, tx, a + H + u + c, a + H + u) = \operatorname{rdrv}(td, tx, a + u + c, a + u)\), then \(\operatorname{rdrv}(td, tx, a + 2 \cdot H + k + (c + 2), a + 2 \cdot H + k) = \operatorname{rdrv}(td, tx, a + k + (c + 2), a + k)\).
\end{theorem}
\begin{proof}
The goal follows from \(\operatorname{rdrvsubdiagperioddoubleoftwoprevperiods}(td, tx, c, a, H, k)\).
\end{proof}
\begin{theorem}
\label{ll:paper:rdrv-subdiag-succ-two-step-of-black-prev}
For every natural number \(td\) and natural number \(tx\) and natural number \(c\) and natural number \(u\), if \(\operatorname{rdrv}(td, tx, u + 1 + (c + 1), u + 1) = \operatorname{true}\), then \(\operatorname{rdrv}(td, tx, u + 1 + (c + 2), u + 1) = \operatorname{not}(\operatorname{rdrv}(td, tx, u + (c + 2), u))\).
\end{theorem}
\begin{proof}
The goal follows from \(\operatorname{rdrvsubdiagsucctwostepofblackprev}(td, tx, c, u)\).
\end{proof}
\begin{theorem}
\label{ll:paper:rdrv-subdiag-one-period-two-of-black-main-diag}
For every natural number \(td\) and natural number \(tx\) and natural number \(u\), if \(\operatorname{rdrv}(td, tx, u + 1, u + 1) = \operatorname{true}\), then if \(\operatorname{rdrv}(td, tx, u + 2, u + 2) = \operatorname{true}\), then \(\operatorname{rdrv}(td, tx, u + 2 + 1, u + 2) = \operatorname{rdrv}(td, tx, u + 1, u)\).
\end{theorem}
\begin{proof}
The goal follows from \(\operatorname{rdrvsubdiagoneperiodtwoofblackmaindiag}(td, tx, u)\).
\end{proof}
\begin{theorem}
\label{ll:paper:rdrv-subdiag-one-agree-of-black-main-diag}
For every natural number \(td\) and natural number \(tx2\) and natural number \(tx\) and natural number \(k\), if for every natural number \(u\), if \(u < k\), then \(\operatorname{rdrv}(td, tx2, u + 1, u + 1) = \operatorname{true}\) and \(\operatorname{rdrv}(td, tx, u + 1, u + 1) = \operatorname{true}\), then \(\operatorname{rdrv}(td, tx2, k + 1, k) = \operatorname{rdrv}(td, tx, k + 1, k)\).
\end{theorem}
\begin{proof}
The goal follows from \(\operatorname{rdrvsubdiagoneagreeofblackmaindiag}(td, tx2, tx, k)\).
\end{proof}
\begin{theorem}
\label{ll:paper:rdrv-subdiag-two-step-of-black-main-diag}
For every natural number \(td\) and natural number \(tx\) and natural number \(u\), if \(\operatorname{rdrv}(td, tx, u + 2, u + 2) = \operatorname{true}\), then \(\operatorname{rdrv}(td, tx, u + 1 + 2, u + 1) = \operatorname{not}(\operatorname{rdrv}(td, tx, u + 2, u))\).
\end{theorem}
\begin{proof}
The goal follows from \(\operatorname{rdrvsubdiagtwostepofblackmaindiag}(td, tx, u)\).
\end{proof}
\begin{theorem}
\label{ll:paper:rdrv-subdiag-two-period-two-of-black-main-diag}
For every natural number \(td\) and natural number \(tx\) and natural number \(u\), if \(\operatorname{rdrv}(td, tx, u + 2, u + 2) = \operatorname{true}\), then if \(\operatorname{rdrv}(td, tx, u + 3, u + 3) = \operatorname{true}\), then \(\operatorname{rdrv}(td, tx, u + 2 + 2, u + 2) = \operatorname{rdrv}(td, tx, u + 2, u)\).
\end{theorem}
\begin{proof}
The goal follows from \(\operatorname{rdrvsubdiagtwoperiodtwoofblackmaindiag}(td, tx, u)\).
\end{proof}
\begin{theorem}
\label{ll:paper:rdrv-subdiag-two-agree-of-black-main-diag}
For every natural number \(td\) and natural number \(tx2\) and natural number \(tx\) and natural number \(k\), if for every natural number \(u\), if \(u < k\), then \(\operatorname{rdrv}(td, tx2, u + 2, u + 2) = \operatorname{true}\) and \(\operatorname{rdrv}(td, tx, u + 2, u + 2) = \operatorname{true}\), then \(\operatorname{rdrv}(td, tx2, k + 2, k) = \operatorname{rdrv}(td, tx, k + 2, k)\).
\end{theorem}
\begin{proof}
The goal follows from \(\operatorname{rdrvsubdiagtwoagreeofblackmaindiag}(td, tx2, tx, k)\).
\end{proof}
\begin{theorem}
\label{ll:paper:rdrv-subcone-two-agree-of-black-main-diag}
For every natural number \(td\) and natural number \(tx2\) and natural number \(tx\) and natural number \(k\) and natural number \(p\), if for every natural number \(u\), if \(u \le k\), then \(\operatorname{rdrv}(td, tx2, u + 1, u + 1) = \operatorname{true}\) and \(\operatorname{rdrv}(td, tx, u + 1, u + 1) = \operatorname{true}\), then if \(p + 2 \le k\), then \(\operatorname{rdrv}(td, tx2, k, p) = \operatorname{rdrv}(td, tx, k, p)\).
\end{theorem}
\begin{proof}
The goal follows from \(\operatorname{rdrvsubconetwoagreeofblackmaindiag}(td, tx2, tx, k, p)\).
\end{proof}
\begin{theorem}
\label{ll:paper:rhcenter-superdiag-eq-rdiag}
For every natural number \(t\) and natural number \(c\) and natural number \(u\), if \(c \le t\), then \(\operatorname{rhcenter}(t + u, u + c) = \operatorname{rdiag}(t - c, t + u)\).
\end{theorem}
\begin{proof}
The goal follows from \(\operatorname{rhcentersuperdiageqrdiag}(t, c, u)\).
\end{proof}
\begin{theorem}
\label{ll:paper:rdrv-superdiag-eq-rdiag}
For every natural number \(td\) and natural number \(tx\) and natural number \(c\) and natural number \(u\), if \(1 \le c\), then if \(c \le tx\), then \(\operatorname{rdrv}(td, tx, u, u + c) = \operatorname{rdiag}(tx - c, tx + u)\).
\end{theorem}
\begin{proof}
The goal follows from \(\operatorname{rdrvsuperdiageqrdiag}(td, tx, c, u)\).
\end{proof}
\begin{theorem}
\label{ll:paper:rhcenter-superdiag-period-twopow}
For every natural number \(t\) and natural number \(c\) and natural number \(k\) and natural number \(u\) and natural number \(v\), if \(t - c \le k\), then \(\operatorname{rhcenter}(t + (u + v \cdot \operatorname{twopow}(k)), u + v \cdot \operatorname{twopow}(k) + c) = \operatorname{rhcenter}(t + u, u + c)\).
\end{theorem}
\begin{proof}
The goal follows from \(\operatorname{rhcentersuperdiagperiodtwopow}(t, c, k, u, v)\).
\end{proof}
\begin{theorem}
\label{ll:paper:rhcenter-superdiag-period-twopow-of-threshold}
For every natural number \(t\) and natural number \(c\) and natural number \(k\) and natural number \(J\) and natural number \(u\) and natural number \(v\), if \(\operatorname{rightthreshold}(k, J)\), then if \(t - c \le J\), then \(\operatorname{rhcenter}(t + (u + v \cdot \operatorname{twopow}(k)), u + v \cdot \operatorname{twopow}(k) + c) = \operatorname{rhcenter}(t + u, u + c)\).
\end{theorem}
\begin{proof}
The goal follows from \(\operatorname{rhcentersuperdiagperiodtwopowofthreshold}(t, c, k, J, u, v)\).
\end{proof}
\begin{theorem}
\label{ll:paper:rdrv-superdiag-period-twopow}
For every natural number \(td\) and natural number \(tx\) and natural number \(c\) and natural number \(k\) and natural number \(u\) and natural number \(v\), if \(1 \le c\), then if \(tx - c \le k\), then \(\operatorname{rdrv}(td, tx, u + v \cdot \operatorname{twopow}(k), u + v \cdot \operatorname{twopow}(k) + c) = \operatorname{rdrv}(td, tx, u, u + c)\).
\end{theorem}
\begin{proof}
The goal follows from \(\operatorname{rdrvsuperdiagperiodtwopow}(td, tx, c, k, u, v)\).
\end{proof}
\begin{theorem}
\label{ll:paper:rdrv-superdiag-period-twopow-of-threshold}
For every natural number \(td\) and natural number \(tx\) and natural number \(c\) and natural number \(k\) and natural number \(J\) and natural number \(u\) and natural number \(v\), if \(1 \le c\), then if \(\operatorname{rightthreshold}(k, J)\), then if \(tx - c \le J\), then \(\operatorname{rdrv}(td, tx, u + v \cdot \operatorname{twopow}(k), u + v \cdot \operatorname{twopow}(k) + c) = \operatorname{rdrv}(td, tx, u, u + c)\).
\end{theorem}
\begin{proof}
The goal follows from \(\operatorname{rdrvsuperdiagperiodtwopowofthreshold}(td, tx, c, k, J, u, v)\).
\end{proof}
\begin{theorem}
\label{ll:paper:rhcenter-superdiag-agree-all-of-lt-twopow}
For every natural number \(t\) and natural number \(Q\) and natural number \(c\) and natural number \(k\) and natural number \(u\), if \(t + Q - c \le k\), then if for every natural number \(r\), if \(r < \operatorname{twopow}(k)\), then \(\operatorname{rhcenter}(t + Q + r, r + c) = \operatorname{rhcenter}(t + r, r + c)\), then \(\operatorname{rhcenter}(t + Q + u, u + c) = \operatorname{rhcenter}(t + u, u + c)\).
\end{theorem}
\begin{proof}
The goal follows from \(\operatorname{rhcentersuperdiagagreealloflttwopow}(t, Q, c, k, u)\).
\end{proof}
\begin{theorem}
\label{ll:paper:rhcenter-superdiag-agree-all-of-threshold}
For every natural number \(t\) and natural number \(Q\) and natural number \(c\) and natural number \(k\) and natural number \(J\) and natural number \(u\), if \(\operatorname{rightthreshold}(k, J)\), then if \(t + Q - c \le J\), then if for every natural number \(r\), if \(r < \operatorname{twopow}(k)\), then \(\operatorname{rhcenter}(t + Q + r, r + c) = \operatorname{rhcenter}(t + r, r + c)\), then \(\operatorname{rhcenter}(t + Q + u, u + c) = \operatorname{rhcenter}(t + u, u + c)\).
\end{theorem}
\begin{proof}
The goal follows from \(\operatorname{rhcentersuperdiagagreeallofthreshold}(t, Q, c, k, J, u)\).
\end{proof}
\begin{theorem}
\label{ll:paper:rhcenter-superdiag-gate-all-of-lt-twopow}
For every natural number \(t\) and natural number \(Q\) and natural number \(c\) and natural number \(k\) and natural number \(u\), if \(t + Q - c \le k\), then if for every natural number \(r\), if \(r < \operatorname{twopow}(k)\), then if \(\operatorname{rhcenter}(t + r, r + c) = \operatorname{false}\), then \(\operatorname{rhcenter}(t + Q + r, r + (c + 1)) = \operatorname{rhcenter}(t + r, r + (c + 1))\), then if \(\operatorname{rhcenter}(t + u, u + c) = \operatorname{false}\), then \(\operatorname{rhcenter}(t + Q + u, u + (c + 1)) = \operatorname{rhcenter}(t + u, u + (c + 1))\).
\end{theorem}
\begin{proof}
The goal follows from \(\operatorname{rhcentersuperdiaggatealloflttwopow}(t, Q, c, k, u)\).
\end{proof}
\begin{theorem}
\label{ll:paper:rhcenter-superdiag-gate-all-of-threshold}
For every natural number \(t\) and natural number \(Q\) and natural number \(c\) and natural number \(k\) and natural number \(J\) and natural number \(u\), if \(\operatorname{rightthreshold}(k, J)\), then if \(t + Q - c \le J\), then if for every natural number \(r\), if \(r < \operatorname{twopow}(k)\), then if \(\operatorname{rhcenter}(t + r, r + c) = \operatorname{false}\), then \(\operatorname{rhcenter}(t + Q + r, r + (c + 1)) = \operatorname{rhcenter}(t + r, r + (c + 1))\), then if \(\operatorname{rhcenter}(t + u, u + c) = \operatorname{false}\), then \(\operatorname{rhcenter}(t + Q + u, u + (c + 1)) = \operatorname{rhcenter}(t + u, u + (c + 1))\).
\end{theorem}
\begin{proof}
The goal follows from \(\operatorname{rhcentersuperdiaggateallofthreshold}(t, Q, c, k, J, u)\).
\end{proof}
\begin{theorem}
\label{ll:paper:rdrv-trapezoid-agree-of-superdiag-gate-twopow}
For every natural number \(td\) and natural number \(tx\) and natural number \(Q\) and natural number \(B\) and natural number \(k\) and natural number \(s\) and natural number \(p\), if \(tx + Q - (B + 1) \le k\), then if \(\operatorname{rhbandagree}(tx, Q, 1, B)\), then if for every natural number \(r\), if \(r < \operatorname{twopow}(k)\), then \(\operatorname{rhcenter}(tx + Q + r, r + (B + 1)) = \operatorname{rhcenter}(tx + r, r + (B + 1))\), then if for every natural number \(r\), if \(r < \operatorname{twopow}(k)\), then if \(\operatorname{rhcenter}(tx + r, r + (B + 1)) = \operatorname{false}\), then \(\operatorname{rhcenter}(tx + Q + r, r + (B + 2)) = \operatorname{rhcenter}(tx + r, r + (B + 2))\), then if \(p \le s + (B + 1)\), then \(\operatorname{rdrv}(td, tx + Q, s, p) = \operatorname{rdrv}(td, tx, s, p)\).
\end{theorem}
\begin{proof}
The goal follows from \(\operatorname{rdrvtrapezoidagreeofsuperdiaggatetwopow}(td, tx, Q, B, k, s, p)\).
\end{proof}
\begin{theorem}
\label{ll:paper:rdrv-trapezoid-agree-of-superdiag-gate-threshold}
For every natural number \(td\) and natural number \(tx\) and natural number \(Q\) and natural number \(B\) and natural number \(k\) and natural number \(J\) and natural number \(s\) and natural number \(p\), if \(\operatorname{rightthreshold}(k, J)\), then if \(tx + Q - (B + 1) \le J\), then if \(\operatorname{rhbandagree}(tx, Q, 1, B)\), then if for every natural number \(r\), if \(r < \operatorname{twopow}(k)\), then \(\operatorname{rhcenter}(tx + Q + r, r + (B + 1)) = \operatorname{rhcenter}(tx + r, r + (B + 1))\), then if for every natural number \(r\), if \(r < \operatorname{twopow}(k)\), then if \(\operatorname{rhcenter}(tx + r, r + (B + 1)) = \operatorname{false}\), then \(\operatorname{rhcenter}(tx + Q + r, r + (B + 2)) = \operatorname{rhcenter}(tx + r, r + (B + 2))\), then if \(p \le s + (B + 1)\), then \(\operatorname{rdrv}(td, tx + Q, s, p) = \operatorname{rdrv}(td, tx, s, p)\).
\end{theorem}
\begin{proof}
The goal follows from \(\operatorname{rdrvtrapezoidagreeofsuperdiaggatethreshold}(td, tx, Q, B, k, J, s, p)\).
\end{proof}
\begin{theorem}
\label{ll:paper:periodic-maskparity-zero-agree-of-superdiag-gate-twopow}
For every natural number \(T\) and natural number \(P\), if \(0 < P\), then if \(\operatorname{centerperiodic}(T, P)\), then for every natural number \(m\) and natural number \(t\) and natural number \(B\) and natural number \(k\) and natural number \(s\), if \(T \le t\), then if \(t + m \cdot P - (B + 1) \le k\), then if \(\operatorname{rhbandagree}(t, m \cdot P, 1, B)\), then if for every natural number \(r\), if \(r < \operatorname{twopow}(k)\), then \(\operatorname{rhcenter}(t + m \cdot P + r, r + (B + 1)) = \operatorname{rhcenter}(t + r, r + (B + 1))\), then if for every natural number \(r\), if \(r < \operatorname{twopow}(k)\), then if \(\operatorname{rhcenter}(t + r, r + (B + 1)) = \operatorname{false}\), then \(\operatorname{rhcenter}(t + m \cdot P + r, r + (B + 2)) = \operatorname{rhcenter}(t + r, r + (B + 2))\), then \(\operatorname{maskparity}(t + s + m \cdot P, 0) = \operatorname{maskparity}(t + s, 0)\).
\end{theorem}
\begin{proof}
Assume \(hP\), \(hper\).
The goal follows from \(\operatorname{periodicmaskparityzeroagreeofsuperdiaggatetwopow}(T, P, hP, hper)\).
\end{proof}
\begin{theorem}
\label{ll:paper:not-centerperiodic-of-superdiag-gate-twopow-at-zero-diag}
For every natural number \(T\) and natural number \(P\) and natural number \(m\) and natural number \(z\) and natural number \(p\) and natural number \(B\) and natural number \(k\), if \(0 < P\), then if \(0 < m\), then if \(T \le z + (p + 1)\), then if \(\operatorname{zerodiagonal}(z)\), then if \(\operatorname{doublesat}(m \cdot P, z)\), then if \(\operatorname{settledby}(z + 1, z + (p + 1))\), then if \(\operatorname{settledby}(z + 1 - m \cdot P, z + (p + 1))\), then if \(z + (p + 1) + m \cdot P - (B + 1) \le k\), then if \(\operatorname{rhbandagree}(z + (p + 1), m \cdot P, 1, B)\), then if for every natural number \(r\), if \(r < \operatorname{twopow}(k)\), then \(\operatorname{rhcenter}(z + (p + 1) + m \cdot P + r, r + (B + 1)) = \operatorname{rhcenter}(z + (p + 1) + r, r + (B + 1))\), then if for every natural number \(r\), if \(r < \operatorname{twopow}(k)\), then if \(\operatorname{rhcenter}(z + (p + 1) + r, r + (B + 1)) = \operatorname{false}\), then \(\operatorname{rhcenter}(z + (p + 1) + m \cdot P + r, r + (B + 2)) = \operatorname{rhcenter}(z + (p + 1) + r, r + (B + 2))\), then not \(\operatorname{centerperiodic}(T, P)\).
\end{theorem}
\begin{proof}
The goal follows from \(\operatorname{notcenterperiodicofsuperdiaggatetwopowatzerodiag}(T, P, m, z, p, B, k)\).
\end{proof}
\begin{theorem}
\label{ll:paper:not-centerperiodic-of-superdiag-gate-threshold-at-zero-diag}
For every natural number \(T\) and natural number \(P\) and natural number \(m\) and natural number \(z\) and natural number \(p\) and natural number \(B\) and natural number \(k\) and natural number \(J\), if \(0 < P\), then if \(0 < m\), then if \(T \le z + (p + 1)\), then if \(\operatorname{zerodiagonal}(z)\), then if \(\operatorname{doublesat}(m \cdot P, z)\), then if \(\operatorname{settledby}(z + 1, z + (p + 1))\), then if \(\operatorname{settledby}(z + 1 - m \cdot P, z + (p + 1))\), then if \(\operatorname{rightthreshold}(k, J)\), then if \(z + (p + 1) + m \cdot P - (B + 1) \le J\), then if \(\operatorname{rhbandagree}(z + (p + 1), m \cdot P, 1, B)\), then if for every natural number \(r\), if \(r < \operatorname{twopow}(k)\), then \(\operatorname{rhcenter}(z + (p + 1) + m \cdot P + r, r + (B + 1)) = \operatorname{rhcenter}(z + (p + 1) + r, r + (B + 1))\), then if for every natural number \(r\), if \(r < \operatorname{twopow}(k)\), then if \(\operatorname{rhcenter}(z + (p + 1) + r, r + (B + 1)) = \operatorname{false}\), then \(\operatorname{rhcenter}(z + (p + 1) + m \cdot P + r, r + (B + 2)) = \operatorname{rhcenter}(z + (p + 1) + r, r + (B + 2))\), then not \(\operatorname{centerperiodic}(T, P)\).
\end{theorem}
\begin{proof}
The goal follows from \(\operatorname{notcenterperiodicofsuperdiaggatethresholdatzerodiag}(T, P, m, z, p, B, k, J)\).
\end{proof}
\begin{theorem}
\label{ll:paper:rdrv-inward-cone-agree-from-step}
For every natural number \(td\) and natural number \(tx2\) and natural number \(tx\) and natural number \(a\) and natural number \(B\) and natural number \(k\) and natural number \(p\), if for every natural number \(q\), if \(q \le B + 1\), then \(\operatorname{rdrv}(td, tx2, a, q) = \operatorname{rdrv}(td, tx, a, q)\), then if \(p + k \le B + 1\), then \(\operatorname{rdrv}(td, tx2, a + k, p) = \operatorname{rdrv}(td, tx, a + k, p)\).
\end{theorem}
\begin{proof}
The goal follows from \(\operatorname{rdrvinwardconeagreefromstep}(td, tx2, tx, a, B, k, p)\).
\end{proof}
\begin{theorem}
\label{ll:paper:rdrv-inward-cone-agree-of-init-agree}
For every natural number \(td\) and natural number \(tx2\) and natural number \(tx\) and natural number \(B\) and natural number \(k\) and natural number \(p\), if for every natural number \(q\), if \(q \le B + 1\), then \(\operatorname{rdrv}(td, tx2, 0, q) = \operatorname{rdrv}(td, tx, 0, q)\), then if \(p + k \le B + 1\), then \(\operatorname{rdrv}(td, tx2, k, p) = \operatorname{rdrv}(td, tx, k, p)\).
\end{theorem}
\begin{proof}
The goal follows from \(\operatorname{rdrvinwardconeagreeofinitagree}(td, tx2, tx, B, k, p)\).
\end{proof}
\begin{theorem}
\label{ll:paper:rdrv-inward-cone-agree-of-rhbandagree}
For every natural number \(td\) and natural number \(tx\) and natural number \(Q\) and natural number \(B\) and natural number \(k\) and natural number \(p\), if \(\operatorname{rhbandagree}(tx, Q, 1, B)\), then if \(p + k \le B + 1\), then \(\operatorname{rdrv}(td, tx + Q, k, p) = \operatorname{rdrv}(td, tx, k, p)\).
\end{theorem}
\begin{proof}
The goal follows from \(\operatorname{rdrvinwardconeagreeofrhbandagree}(td, tx, Q, B, k, p)\).
\end{proof}
\begin{theorem}
\label{ll:paper:rdrv-col-one-agree-of-rhbandagree}
For every natural number \(td\) and natural number \(tx\) and natural number \(Q\) and natural number \(B\) and natural number \(k\), if \(\operatorname{rhbandagree}(tx, Q, 1, B)\), then if \(k \le B\), then \(\operatorname{rdrv}(td, tx + Q, k, 1) = \operatorname{rdrv}(td, tx, k, 1)\).
\end{theorem}
\begin{proof}
The goal follows from \(\operatorname{rdrvcoloneagreeofrhbandagree}(td, tx, Q, B, k)\).
\end{proof}
\begin{theorem}
\label{ll:paper:rdrv-inward-cone-eq-white-seed-of-initial-white-band}
For every natural number \(td\) and natural number \(tx\) and natural number \(B\) and natural number \(k\) and natural number \(p\), if for every natural number \(i\), if \(i \le B\), then \(\operatorname{rhcenter}(tx, 1 + i) = \operatorname{false}\), then if \(p + k \le B + 1\), then \(\operatorname{rdrv}(td, tx, k, p) = \operatorname{rdrv}(td, 0, k, p)\).
\end{theorem}
\begin{proof}
The goal follows from \(\operatorname{rdrvinwardconeeqwhiteseedofinitialwhiteband}(td, tx, B, k, p)\).
\end{proof}
\begin{theorem}
\label{ll:paper:rhcenter-inward-cone-eq-white-seed-of-initial-white-band}
For every natural number \(t\) and natural number \(B\) and natural number \(k\) and natural number \(p\), if for every natural number \(i\), if \(i \le B\), then \(\operatorname{rhcenter}(t, 1 + i) = \operatorname{false}\), then if \(p + k \le B + 1\), then \(\operatorname{rhcenter}(t + k, p) = \operatorname{rdrv}(t, 0, k, p)\).
\end{theorem}
\begin{proof}
The goal follows from \(\operatorname{rhcenterinwardconeeqwhiteseedofinitialwhiteband}(t, B, k, p)\).
\end{proof}
\begin{theorem}
\label{ll:paper:periodic-rhcenter-inward-cone-agree-of-rhbandagree}
For every natural number \(T\) and natural number \(P\), if \(0 < P\), then if \(\operatorname{centerperiodic}(T, P)\), then for every natural number \(m\) and natural number \(t\) and natural number \(B\) and natural number \(s\) and natural number \(p\), if \(T \le t\), then if \(\operatorname{rhbandagree}(t, m \cdot P, 1, B)\), then if \(p + s \le B + 1\), then \(\operatorname{rhcenter}(t + s + m \cdot P, p) = \operatorname{rhcenter}(t + s, p)\).
\end{theorem}
\begin{proof}
Assume \(hP\), \(hper\).
The goal follows from \(\operatorname{periodicrhcenterinwardconeagreeofrhbandagree}(T, P, hP, hper)\).
\end{proof}
\begin{theorem}
\label{ll:paper:periodic-rhbandagree-step-of-rhbandagree}
For every natural number \(T\) and natural number \(P\), if \(0 < P\), then if \(\operatorname{centerperiodic}(T, P)\), then for every natural number \(m\) and natural number \(t\) and natural number \(B\) and natural number \(s\), if \(T \le t\), then if \(\operatorname{rhbandagree}(t, m \cdot P, 1, B)\), then if \(s \le B\), then \(\operatorname{rhbandagree}(t + s, m \cdot P, 1, B - s)\).
\end{theorem}
\begin{proof}
Assume \(hP\), \(hper\).
The goal follows from \(\operatorname{periodicrhbandagreestepofrhbandagree}(T, P, hP, hper)\).
\end{proof}
\begin{theorem}
\label{ll:paper:periodic-rhcenter-one-agree-of-rhbandagree}
For every natural number \(T\) and natural number \(P\), if \(0 < P\), then if \(\operatorname{centerperiodic}(T, P)\), then for every natural number \(m\) and natural number \(t\) and natural number \(B\) and natural number \(s\), if \(T \le t\), then if \(\operatorname{rhbandagree}(t, m \cdot P, 1, B)\), then if \(s \le B\), then \(\operatorname{rhcenter}(t + s + m \cdot P, 1) = \operatorname{rhcenter}(t + s, 1)\).
\end{theorem}
\begin{proof}
Assume \(hP\), \(hper\).
The goal follows from \(\operatorname{periodicrhcenteroneagreeofrhbandagree}(T, P, hP, hper)\).
\end{proof}
\begin{theorem}
\label{ll:paper:periodic-maskparity-zero-agree-of-rhbandagree}
For every natural number \(T\) and natural number \(P\), if \(0 < P\), then if \(\operatorname{centerperiodic}(T, P)\), then for every natural number \(m\) and natural number \(t\) and natural number \(B\) and natural number \(s\), if \(T \le t\), then if \(\operatorname{rhbandagree}(t, m \cdot P, 1, B)\), then if \(s \le B\), then \(\operatorname{maskparity}(t + s + m \cdot P, 0) = \operatorname{maskparity}(t + s, 0)\).
\end{theorem}
\begin{proof}
Assume \(hP\), \(hper\).
The goal follows from \(\operatorname{periodicmaskparityzeroagreeofrhbandagree}(T, P, hP, hper)\).
\end{proof}
\begin{theorem}
\label{ll:paper:not-centerperiodic-of-rhbandagree-at-zero-diag-doubles}
For every natural number \(T\) and natural number \(P\) and natural number \(m\) and natural number \(z\) and natural number \(p\) and natural number \(B\), if \(0 < P\), then if \(0 < m\), then if \(T \le z + (p + 1)\), then if \(\operatorname{zerodiagonal}(z)\), then if \(\operatorname{doublesat}(m \cdot P, z)\), then if \(\operatorname{settledby}(z + 1, z + (p + 1))\), then if \(\operatorname{settledby}(z + 1 - m \cdot P, z + (p + 1))\), then if \(p + 3 \cdot (m \cdot P) \le B + 1\), then if \(\operatorname{rhbandagree}(z + (p + 1), m \cdot P, 1, B)\), then not \(\operatorname{centerperiodic}(T, P)\).
\end{theorem}
\begin{proof}
The goal follows from \(\operatorname{notcenterperiodicofrhbandagreeatzerodiagdoubles}(T, P, m, z, p, B)\).
\end{proof}
\begin{theorem}
\label{ll:paper:not-centerperiodic-of-rhbandagree-twopow}
For every natural number \(T\) and natural number \(P\) and natural number \(m\) and natural number \(t\) and natural number \(k\) and natural number \(B\), if \(0 < P\), then if \(0 < m\), then if \(T \le t\), then if \(t + m \cdot P = \operatorname{twopow}(k)\), then if \(\operatorname{twopow}(k) \le B + 1\), then if \(\operatorname{rhbandagree}(t, m \cdot P, 1, B)\), then not \(\operatorname{centerperiodic}(T, P)\).
\end{theorem}
\begin{proof}
The goal follows from \(\operatorname{notcenterperiodicofrhbandagreetwopow}(T, P, m, t, k, B)\).
\end{proof}
\begin{theorem}
\label{ll:paper:not-centerperiodic-of-rhbandagree-at-zero-diag-pair}
For every natural number \(T\) and natural number \(P\) and natural number \(m\) and natural number \(z\) and natural number \(k\) and natural number \(B\), if \(0 < P\), then if \(\operatorname{zerodiagonal}(z)\), then if \(2 \le z + m \cdot P\), then if \(\operatorname{zerodiagonal}(z + m \cdot P - 2)\), then if \(T \le z + \operatorname{twopow}(k)\), then if \(\operatorname{settledby}(z, z + \operatorname{twopow}(k))\), then if \(\operatorname{settledby}(z + m \cdot P, z + m \cdot P + \operatorname{twopow}(k))\), then if \(\operatorname{twopow}(k) \le B + 1\), then if \(\operatorname{rhbandagree}(z + \operatorname{twopow}(k), m \cdot P, 1, B)\), then not \(\operatorname{centerperiodic}(T, P)\).
\end{theorem}
\begin{proof}
The goal follows from \(\operatorname{notcenterperiodicofrhbandagreeatzerodiagpair}(T, P, m, z, k, B)\).
\end{proof}
\begin{theorem}
\label{ll:paper:rdrv-inward-cone-agree-of-short-boundary-gate}
For every natural number \(td\) and natural number \(tx\) and natural number \(Q\) and natural number \(B\) and natural number \(a\) and natural number \(k\) and natural number \(p\), if \(\operatorname{rhbandagree}(tx, Q, 1, B)\), then if for every natural number \(u\), if \(u < a\), then if \(\operatorname{rdrv}(td, tx, u, B + 1) = \operatorname{false}\), then \(\operatorname{rdrv}(td, tx + Q, u, B + 2) = \operatorname{rdrv}(td, tx, u, B + 2)\), then if \(p + k \le B + 1\), then \(\operatorname{rdrv}(td, tx + Q, a + k, p) = \operatorname{rdrv}(td, tx, a + k, p)\).
\end{theorem}
\begin{proof}
The goal follows from \(\operatorname{rdrvinwardconeagreeofshortboundarygate}(td, tx, Q, B, a, k, p)\).
\end{proof}
\begin{theorem}
\label{ll:paper:periodic-maskparity-zero-agree-of-short-boundary-gate}
For every natural number \(T\) and natural number \(P\), if \(0 < P\), then if \(\operatorname{centerperiodic}(T, P)\), then for every natural number \(m\) and natural number \(t\) and natural number \(B\) and natural number \(a\) and natural number \(k\), if \(T \le t\), then if \(\operatorname{rhbandagree}(t, m \cdot P, 1, B)\), then if for every natural number \(u\), if \(u < a\), then if \(\operatorname{rdrv}(t, t, u, B + 1) = \operatorname{false}\), then \(\operatorname{rdrv}(t, t + m \cdot P, u, B + 2) = \operatorname{rdrv}(t, t, u, B + 2)\), then if \(k \le B\), then \(\operatorname{maskparity}(t + (a + k) + m \cdot P, 0) = \operatorname{maskparity}(t + (a + k), 0)\).
\end{theorem}
\begin{proof}
Assume \(hP\), \(hper\).
The goal follows from \(\operatorname{periodicmaskparityzeroagreeofshortboundarygate}(T, P, hP, hper)\).
\end{proof}
\begin{theorem}
\label{ll:paper:rdrv-black-reaches-and-stays-boundary}
For every natural number \(td\) and natural number \(tx\) and natural number \(a\) and natural number \(p\) and natural number \(n\) and natural number \(r\) and natural number \(k\), if for every natural number \(s\), if \(s < r\), then \(\operatorname{rdrv}(td, tx, a + s, p) = \operatorname{false}\), then if \(\operatorname{rdrv}(td, tx, a, p + (n + 1)) = \operatorname{true}\), then if \(n \le k\), then if \(k \le r\), then \(\operatorname{rdrv}(td, tx, a + k, p + 1) = \operatorname{true}\).
\end{theorem}
\begin{proof}
The goal follows from \(\operatorname{rdrvblackreachesandstaysboundary}(td, tx, a, p, n, r, k)\).
\end{proof}
\begin{theorem}
\label{ll:paper:rdrv-col-succ-flip-after-white-col-and-black-cell}
For every natural number \(td\) and natural number \(tx\) and natural number \(a\) and natural number \(p\) and natural number \(n\) and natural number \(r\), if for every natural number \(s\), if \(s < r\), then \(\operatorname{rdrv}(td, tx, a + s, p) = \operatorname{false}\), then if \(\operatorname{rdrv}(td, tx, a + r, p) = \operatorname{true}\), then if \(\operatorname{rdrv}(td, tx, a, p + (n + 1)) = \operatorname{true}\), then if \(n \le r\), then \(\operatorname{rdrv}(td, tx, a + (r + 1), p + 1) = \operatorname{false}\).
\end{theorem}
\begin{proof}
The goal follows from \(\operatorname{rdrvcolsuccflipafterwhitecolandblackcell}(td, tx, a, p, n, r)\).
\end{proof}
\begin{theorem}
\label{ll:paper:rdrv-col-succ-agree-of-white-col-and-black-cells}
For every natural number \(td\) and natural number \(tx2\) and natural number \(tx\) and natural number \(a\) and natural number \(p\) and natural number \(n2\) and natural number \(n1\) and natural number \(r\) and natural number \(k\), if for every natural number \(s\), if \(s < r\), then \(\operatorname{rdrv}(td, tx2, a + s, p) = \operatorname{false}\) and \(\operatorname{rdrv}(td, tx, a + s, p) = \operatorname{false}\), then if \(\operatorname{rdrv}(td, tx2, a, p + (n2 + 1)) = \operatorname{true}\), then if \(\operatorname{rdrv}(td, tx, a, p + (n1 + 1)) = \operatorname{true}\), then if \(n2 \le k\), then if \(n1 \le k\), then if \(k \le r\), then \(\operatorname{rdrv}(td, tx2, a + k, p + 1) = \operatorname{rdrv}(td, tx, a + k, p + 1)\).
\end{theorem}
\begin{proof}
The goal follows from \(\operatorname{rdrvcolsuccagreeofwhitecolandblackcells}(td, tx2, tx, a, p, n2, n1, r, k)\).
\end{proof}
\begin{theorem}
\label{ll:paper:rdrv-col-succ-agree-after-white-col-and-black-cells}
For every natural number \(td\) and natural number \(tx2\) and natural number \(tx\) and natural number \(a\) and natural number \(p\) and natural number \(n2\) and natural number \(n1\) and natural number \(r\), if for every natural number \(s\), if \(s < r\), then \(\operatorname{rdrv}(td, tx2, a + s, p) = \operatorname{false}\) and \(\operatorname{rdrv}(td, tx, a + s, p) = \operatorname{false}\), then if \(\operatorname{rdrv}(td, tx2, a + r, p) = \operatorname{rdrv}(td, tx, a + r, p)\), then if \(\operatorname{rdrv}(td, tx2, a, p + (n2 + 1)) = \operatorname{true}\), then if \(\operatorname{rdrv}(td, tx, a, p + (n1 + 1)) = \operatorname{true}\), then if \(n2 \le r\), then if \(n1 \le r\), then \(\operatorname{rdrv}(td, tx2, a + (r + 1), p + 1) = \operatorname{rdrv}(td, tx, a + (r + 1), p + 1)\).
\end{theorem}
\begin{proof}
The goal follows from \(\operatorname{rdrvcolsuccagreeafterwhitecolandblackcells}(td, tx2, tx, a, p, n2, n1, r)\).
\end{proof}
\begin{theorem}
\label{ll:paper:rdrv-white-offset-forced-of-white-col}
For every natural number \(td\) and natural number \(tx\) and natural number \(a\) and natural number \(p\) and natural number \(r\) and natural number \(q\), if for every natural number \(s\), if \(s < r\), then \(\operatorname{rdrv}(td, tx, a + s, p) = \operatorname{false}\), then if \(\operatorname{rdrv}(td, tx, a + (q - 1), p + 1) = \operatorname{false}\), then if \(1 \le q\), then if \(q \le r + 1\), then \(\operatorname{rdrv}(td, tx, a, p + q) = \operatorname{false}\).
\end{theorem}
\begin{proof}
The goal follows from \(\operatorname{rdrvwhiteoffsetforcedofwhitecol}(td, tx, a, p, r, q)\).
\end{proof}
\begin{theorem}
\label{ll:paper:rdrv-white-triangle-forced-of-two-white-cols}
For every natural number \(td\) and natural number \(tx\) and natural number \(a\) and natural number \(p\) and natural number \(r\) and natural number \(u\) and natural number \(q\), if for every natural number \(s\), if \(s < r\), then \(\operatorname{rdrv}(td, tx, a + s, p) = \operatorname{false}\), then if for every natural number \(s\), if \(s \le r\), then \(\operatorname{rdrv}(td, tx, a + s, p + 1) = \operatorname{false}\), then if \(1 \le q\), then if \(u + q \le r + 1\), then \(\operatorname{rdrv}(td, tx, a + u, p + q) = \operatorname{false}\).
\end{theorem}
\begin{proof}
The goal follows from \(\operatorname{rdrvwhitetriangleforcedoftwowhitecols}(td, tx, a, p, r, u, q)\).
\end{proof}
\begin{theorem}
\label{ll:paper:rdrv-white-triangle-forced-of-white-drive-and-col-one}
For every natural number \(td\) and natural number \(tx\) and natural number \(a\) and natural number \(r\) and natural number \(u\) and natural number \(q\), if for every natural number \(s\), if \(s < r\), then \(\operatorname{center}(td + (a + s)) = \operatorname{false}\), then if for every natural number \(s\), if \(s \le r\), then \(\operatorname{rdrv}(td, tx, a + s, 1) = \operatorname{false}\), then if \(1 \le q\), then if \(u + q \le r + 1\), then \(\operatorname{rdrv}(td, tx, a + u, q) = \operatorname{false}\).
\end{theorem}
\begin{proof}
The goal follows from \(\operatorname{rdrvwhitetriangleforcedofwhitedriveandcolone}(td, tx, a, r, u, q)\).
\end{proof}
\begin{theorem}
\label{ll:paper:rhcenter-white-triangle-forced-of-white-drive-and-col-one}
For every natural number \(t\) and natural number \(r\) and natural number \(u\) and natural number \(q\), if for every natural number \(s\), if \(s < r\), then \(\operatorname{center}(t + s) = \operatorname{false}\), then if for every natural number \(s\), if \(s \le r\), then \(\operatorname{rhcenter}(t + s, 1) = \operatorname{false}\), then if \(1 \le q\), then if \(u + q \le r + 1\), then \(\operatorname{rhcenter}(t + u, q) = \operatorname{false}\).
\end{theorem}
\begin{proof}
The goal follows from \(\operatorname{rhcenterwhitetriangleforcedofwhitedriveandcolone}(t, r, u, q)\).
\end{proof}
\begin{theorem}
\label{ll:paper:rdrv-white-step-of-three-white}
For every natural number \(td\) and natural number \(tx\) and natural number \(k\) and natural number \(p\), if \(\operatorname{rdrv}(td, tx, k, p) = \operatorname{false}\), then if \(\operatorname{rdrv}(td, tx, k, p + 1) = \operatorname{false}\), then if \(\operatorname{rdrv}(td, tx, k, p + 2) = \operatorname{false}\), then \(\operatorname{rdrv}(td, tx, k + 1, p + 1) = \operatorname{false}\).
\end{theorem}
\begin{proof}
The goal follows from \(\operatorname{rdrvwhitestepofthreewhite}(td, tx, k, p)\).
\end{proof}
\begin{theorem}
\label{ll:paper:rdrv-white-triangle-of-white-col-and-spatial-interval}
For every natural number \(td\) and natural number \(tx\) and natural number \(a\) and natural number \(p\) and natural number \(r\) and natural number \(u\) and natural number \(q\), if for every natural number \(s\), if \(s < r\), then \(\operatorname{rdrv}(td, tx, a + s, p) = \operatorname{false}\), then if for every natural number \(j\), if \(1 \le j\), then if \(j \le r + 1\), then \(\operatorname{rdrv}(td, tx, a, p + j) = \operatorname{false}\), then if \(1 \le q\), then if \(u + q \le r + 1\), then \(\operatorname{rdrv}(td, tx, a + u, p + q) = \operatorname{false}\).
\end{theorem}
\begin{proof}
The goal follows from \(\operatorname{rdrvwhitetriangleofwhitecolandspatialinterval}(td, tx, a, p, r, u, q)\).
\end{proof}
\begin{theorem}
\label{ll:paper:rdrv-white-diamond-of-spatial-white-interval}
For every natural number \(td\) and natural number \(tx\) and natural number \(a\) and natural number \(p\) and natural number \(L\) and natural number \(u\) and natural number \(q\), if for every natural number \(j\), if \(1 \le j\), then if \(j \le L\), then \(\operatorname{rdrv}(td, tx, a, p + j) = \operatorname{false}\), then if \(u + 1 \le q\), then if \(u + q \le L\), then \(\operatorname{rdrv}(td, tx, a + u, p + q) = \operatorname{false}\).
\end{theorem}
\begin{proof}
The goal follows from \(\operatorname{rdrvwhitediamondofspatialwhiteinterval}(td, tx, a, p, L, u, q)\).
\end{proof}
\begin{theorem}
\label{ll:paper:rdrv-superdiag-white-corridor-of-spatial-white-interval}
For every natural number \(td\) and natural number \(tx\) and natural number \(a\) and natural number \(p\) and natural number \(L\) and natural number \(u\), if for every natural number \(j\), if \(1 \le j\), then if \(j \le L\), then \(\operatorname{rdrv}(td, tx, a, p + j) = \operatorname{false}\), then if \(2 \cdot u + 2 \le L\), then \(\operatorname{rdrv}(td, tx, a + u, p + (u + 1)) = \operatorname{false}\) and \(\operatorname{rdrv}(td, tx, a + u, p + (u + 2)) = \operatorname{false}\).
\end{theorem}
\begin{proof}
The goal follows from \(\operatorname{rdrvsuperdiagwhitecorridorofspatialwhiteinterval}(td, tx, a, p, L, u)\).
\end{proof}
\begin{theorem}
\label{ll:paper:rhcenter-white-diamond-of-spatial-white-interval}
For every natural number \(t\) and natural number \(p\) and natural number \(L\) and natural number \(u\) and natural number \(q\), if for every natural number \(j\), if \(1 \le j\), then if \(j \le L\), then \(\operatorname{rhcenter}(t, p + j) = \operatorname{false}\), then if \(u + 1 \le q\), then if \(u + q \le L\), then \(\operatorname{rhcenter}(t + u, p + q) = \operatorname{false}\).
\end{theorem}
\begin{proof}
The goal follows from \(\operatorname{rhcenterwhitediamondofspatialwhiteinterval}(t, p, L, u, q)\).
\end{proof}
\begin{theorem}
\label{ll:paper:rhcenter-superdiag-white-corridor-of-spatial-white-interval}
For every natural number \(t\) and natural number \(B\) and natural number \(L\) and natural number \(u\), if for every natural number \(j\), if \(1 \le j\), then if \(j \le L\), then \(\operatorname{rhcenter}(t, B + j) = \operatorname{false}\), then if \(2 \cdot u + 2 \le L\), then \(\operatorname{rhcenter}(t + u, u + (B + 1)) = \operatorname{false}\) and \(\operatorname{rhcenter}(t + u, u + (B + 2)) = \operatorname{false}\).
\end{theorem}
\begin{proof}
The goal follows from \(\operatorname{rhcentersuperdiagwhitecorridorofspatialwhiteinterval}(t, B, L, u)\).
\end{proof}
\begin{theorem}
\label{ll:paper:rdrv-trapezoid-eq-white-seed-of-wide-white-band}
For every natural number \(td\) and natural number \(tx\) and natural number \(B\) and natural number \(k\) and natural number \(p\), if for every natural number \(i\), if \(i \le B + (2 \cdot k + 2)\), then \(\operatorname{rhcenter}(tx, 1 + i) = \operatorname{false}\), then if \(p \le k + (B + 1)\), then \(\operatorname{rdrv}(td, tx, k, p) = \operatorname{rdrv}(td, 0, k, p)\).
\end{theorem}
\begin{proof}
The goal follows from \(\operatorname{rdrvtrapezoideqwhiteseedofwidewhiteband}(td, tx, B, k, p)\).
\end{proof}
\begin{theorem}
\label{ll:paper:rdrv-main-diag-step-eq-rhcenter}
For every natural number \(td\) and natural number \(tx\) and natural number \(u\), \(\operatorname{rdrv}(td, tx, u + 1, u + 1) = \operatorname{bxor}(\operatorname{rdrv}(td, tx, u, u), \operatorname{rhcenter}(tx + u, u + 1) \operatorname{or} \operatorname{rhcenter}(tx + u, u + 2))\).
\end{theorem}
\begin{proof}
The goal follows from \(\operatorname{rdrvmaindiagstepeqrhcenter}(td, tx, u)\).
\end{proof}
\begin{theorem}
\label{ll:paper:rdrv-main-diag-accum-shift-xor}
For every natural number \(td\) and natural number \(tx\) and natural number \(a\) and natural number \(H\) and natural number \(n\), if for every natural number \(u\), if \(u < n\), then \((\operatorname{rhcenter}(tx + (a + H + u), a + H + u + 1) \operatorname{or} \operatorname{rhcenter}(tx + (a + H + u), a + H + u + 2)) = (\operatorname{rhcenter}(tx + (a + u), a + u + 1) \operatorname{or} \operatorname{rhcenter}(tx + (a + u), a + u + 2))\), then \(\operatorname{bxor}(\operatorname{rdrv}(td, tx, a + H + n, a + H + n), \operatorname{rdrv}(td, tx, a + n, a + n)) = \operatorname{bxor}(\operatorname{rdrv}(td, tx, a + H, a + H), \operatorname{rdrv}(td, tx, a, a))\).
\end{theorem}
\begin{proof}
The goal follows from \(\operatorname{rdrvmaindiagaccumshiftxor}(td, tx, a, H, n)\).
\end{proof}
\begin{theorem}
\label{ll:paper:rdrv-main-diag-period-double-of-superdiag-or-period}
For every natural number \(td\) and natural number \(tx\) and natural number \(a\) and natural number \(H\) and natural number \(n\), if for every natural number \(u\), if \(u < H + n\), then \((\operatorname{rhcenter}(tx + (a + H + u), a + H + u + 1) \operatorname{or} \operatorname{rhcenter}(tx + (a + H + u), a + H + u + 2)) = (\operatorname{rhcenter}(tx + (a + u), a + u + 1) \operatorname{or} \operatorname{rhcenter}(tx + (a + u), a + u + 2))\), then \(\operatorname{rdrv}(td, tx, a + 2 \cdot H + n, a + 2 \cdot H + n) = \operatorname{rdrv}(td, tx, a + n, a + n)\).
\end{theorem}
\begin{proof}
The goal follows from \(\operatorname{rdrvmaindiagperioddoubleofsuperdiagorperiod}(td, tx, a, H, n)\).
\end{proof}
\begin{theorem}
\label{ll:paper:rdrv-main-diag-period-twopow-succ}
For every natural number \(td\) and natural number \(tx\) and natural number \(k\) and natural number \(u\), if \(tx - 1 \le k\), then \(\operatorname{rdrv}(td, tx, u + \operatorname{twopow}(k + 1), u + \operatorname{twopow}(k + 1)) = \operatorname{rdrv}(td, tx, u, u)\).
\end{theorem}
\begin{proof}
The goal follows from \(\operatorname{rdrvmaindiagperiodtwopowsucc}(td, tx, k, u)\).
\end{proof}
\begin{theorem}
\label{ll:paper:rdrv-main-diag-period-twopow-succ-of-threshold}
For every natural number \(td\) and natural number \(tx\) and natural number \(k\) and natural number \(J\) and natural number \(u\), if \(\operatorname{rightthreshold}(k, J)\), then if \(tx - 1 \le J\), then \(\operatorname{rdrv}(td, tx, u + \operatorname{twopow}(k + 1), u + \operatorname{twopow}(k + 1)) = \operatorname{rdrv}(td, tx, u, u)\).
\end{theorem}
\begin{proof}
The goal follows from \(\operatorname{rdrvmaindiagperiodtwopowsuccofthreshold}(td, tx, k, J, u)\).
\end{proof}
\begin{theorem}
\label{ll:paper:rdrv-main-diag-period-twopow-succ-mul}
For every natural number \(td\) and natural number \(tx\) and natural number \(k\) and natural number \(u\) and natural number \(v\), if \(tx - 1 \le k\), then \(\operatorname{rdrv}(td, tx, u + v \cdot \operatorname{twopow}(k + 1), u + v \cdot \operatorname{twopow}(k + 1)) = \operatorname{rdrv}(td, tx, u, u)\).
\end{theorem}
\begin{proof}
The goal follows from \(\operatorname{rdrvmaindiagperiodtwopowsuccmul}(td, tx, k, u, v)\).
\end{proof}
\begin{theorem}
\label{ll:paper:rdrv-main-diag-agree-all-of-lt-twopow-succ}
For every natural number \(td\) and natural number \(tx\) and natural number \(Q\) and natural number \(k\) and natural number \(u\), if \(tx + Q - 1 \le k\), then if for every natural number \(r\), if \(r < \operatorname{twopow}(k + 1)\), then \(\operatorname{rdrv}(td, tx + Q, r, r) = \operatorname{rdrv}(td, tx, r, r)\), then \(\operatorname{rdrv}(td, tx + Q, u, u) = \operatorname{rdrv}(td, tx, u, u)\).
\end{theorem}
\begin{proof}
The goal follows from \(\operatorname{rdrvmaindiagagreealloflttwopowsucc}(td, tx, Q, k, u)\).
\end{proof}
\begin{theorem}
\label{ll:paper:rdrv-subdiag-one-step-eq-rhcenter}
For every natural number \(td\) and natural number \(tx\) and natural number \(u\), \(\operatorname{rdrv}(td, tx, u + 1 + 1, u + 1) = \operatorname{bxor}(\operatorname{rdrv}(td, tx, u + 1, u), \operatorname{rdrv}(td, tx, u + 1, u + 1) \operatorname{or} \operatorname{rhcenter}(tx + (u + 1), u + 1 + 1))\).
\end{theorem}
\begin{proof}
The goal follows from \(\operatorname{rdrvsubdiagonestepeqrhcenter}(td, tx, u)\).
\end{proof}
\begin{theorem}
\label{ll:paper:rdrv-subdiag-one-accum-shift-xor}
For every natural number \(td\) and natural number \(tx\) and natural number \(a\) and natural number \(H\) and natural number \(n\), if for every natural number \(u\), if \(u < n\), then \((\operatorname{rdrv}(td, tx, a + H + (u + 1), a + H + (u + 1)) \operatorname{or} \operatorname{rhcenter}(tx + (a + H + (u + 1)), a + H + (u + 1) + 1)) = (\operatorname{rdrv}(td, tx, a + (u + 1), a + (u + 1)) \operatorname{or} \operatorname{rhcenter}(tx + (a + (u + 1)), a + (u + 1) + 1))\), then \(\operatorname{bxor}(\operatorname{rdrv}(td, tx, a + H + n + 1, a + H + n), \operatorname{rdrv}(td, tx, a + n + 1, a + n)) = \operatorname{bxor}(\operatorname{rdrv}(td, tx, a + H + 1, a + H), \operatorname{rdrv}(td, tx, a + 1, a))\).
\end{theorem}
\begin{proof}
The goal follows from \(\operatorname{rdrvsubdiagoneaccumshiftxor}(td, tx, a, H, n)\).
\end{proof}
\begin{theorem}
\label{ll:paper:rdrv-subdiag-one-period-double-of-or-period}
For every natural number \(td\) and natural number \(tx\) and natural number \(a\) and natural number \(H\) and natural number \(n\), if for every natural number \(u\), if \(u < H + n\), then \((\operatorname{rdrv}(td, tx, a + H + (u + 1), a + H + (u + 1)) \operatorname{or} \operatorname{rhcenter}(tx + (a + H + (u + 1)), a + H + (u + 1) + 1)) = (\operatorname{rdrv}(td, tx, a + (u + 1), a + (u + 1)) \operatorname{or} \operatorname{rhcenter}(tx + (a + (u + 1)), a + (u + 1) + 1))\), then \(\operatorname{rdrv}(td, tx, a + 2 \cdot H + n + 1, a + 2 \cdot H + n) = \operatorname{rdrv}(td, tx, a + n + 1, a + n)\).
\end{theorem}
\begin{proof}
The goal follows from \(\operatorname{rdrvsubdiagoneperioddoubleoforperiod}(td, tx, a, H, n)\).
\end{proof}
\begin{theorem}
\label{ll:paper:rdrv-subdiag-one-period-twopow-succ-two}
For every natural number \(td\) and natural number \(tx\) and natural number \(k\) and natural number \(u\), if \(tx - 1 \le k\), then \(\operatorname{rdrv}(td, tx, u + \operatorname{twopow}(k + 2) + 1, u + \operatorname{twopow}(k + 2)) = \operatorname{rdrv}(td, tx, u + 1, u)\).
\end{theorem}
\begin{proof}
The goal follows from \(\operatorname{rdrvsubdiagoneperiodtwopowsucctwo}(td, tx, k, u)\).
\end{proof}
\begin{theorem}
\label{ll:paper:rdrv-subdiag-one-period-twopow-succ-two-of-threshold}
For every natural number \(td\) and natural number \(tx\) and natural number \(k\) and natural number \(J\) and natural number \(u\), if \(\operatorname{rightthreshold}(k, J)\), then if \(tx - 1 \le J\), then \(\operatorname{rdrv}(td, tx, u + \operatorname{twopow}(k + 2) + 1, u + \operatorname{twopow}(k + 2)) = \operatorname{rdrv}(td, tx, u + 1, u)\).
\end{theorem}
\begin{proof}
The goal follows from \(\operatorname{rdrvsubdiagoneperiodtwopowsucctwoofthreshold}(td, tx, k, J, u)\).
\end{proof}
\begin{theorem}
\label{ll:paper:rdrv-subdiag-period-twopow-all}
For every natural number \(td\) and natural number \(tx\) and natural number \(k\) and natural number \(c\) and natural number \(u\), if \(tx - 1 \le k\), then \(\operatorname{rdrv}(td, tx, u + \operatorname{twopow}(k + 1 + c) + c, u + \operatorname{twopow}(k + 1 + c)) = \operatorname{rdrv}(td, tx, u + c, u)\).
\end{theorem}
\begin{proof}
The goal follows from \(\operatorname{rdrvsubdiagperiodtwopowall}(td, tx, k, c, u)\).
\end{proof}
\begin{theorem}
\label{ll:paper:rdrv-subdiag-period-twopow-all-mul}
For every natural number \(td\) and natural number \(tx\) and natural number \(k\) and natural number \(c\) and natural number \(u\) and natural number \(v\), if \(tx - 1 \le k\), then \(\operatorname{rdrv}(td, tx, u + v \cdot \operatorname{twopow}(k + 1 + c) + c, u + v \cdot \operatorname{twopow}(k + 1 + c)) = \operatorname{rdrv}(td, tx, u + c, u)\).
\end{theorem}
\begin{proof}
The goal follows from \(\operatorname{rdrvsubdiagperiodtwopowallmul}(td, tx, k, c, u, v)\).
\end{proof}
\begin{theorem}
\label{ll:paper:rdrv-subdiag-period-twopow-all-of-threshold}
For every natural number \(td\) and natural number \(tx\) and natural number \(k\) and natural number \(J\) and natural number \(c\) and natural number \(u\), if \(\operatorname{rightthreshold}(k, J)\), then if \(tx - 1 \le J\), then \(\operatorname{rdrv}(td, tx, u + \operatorname{twopow}(k + 1 + c) + c, u + \operatorname{twopow}(k + 1 + c)) = \operatorname{rdrv}(td, tx, u + c, u)\).
\end{theorem}
\begin{proof}
The goal follows from \(\operatorname{rdrvsubdiagperiodtwopowallofthreshold}(td, tx, k, J, c, u)\).
\end{proof}
\begin{theorem}
\label{ll:paper:rdrv-main-diag-at-twopow-succ-eq-center}
For every natural number \(td\) and natural number \(tx\) and natural number \(k\) and natural number \(v\), if \(tx - 1 \le k\), then \(\operatorname{rdrv}(td, tx, v \cdot \operatorname{twopow}(k + 1), v \cdot \operatorname{twopow}(k + 1)) = \operatorname{center}(td)\).
\end{theorem}
\begin{proof}
The goal follows from \(\operatorname{rdrvmaindiagattwopowsucceqcenter}(td, tx, k, v)\).
\end{proof}
\begin{theorem}
\label{ll:paper:rdrv-subdiag-at-twopow-all-eq-center}
For every natural number \(td\) and natural number \(tx\) and natural number \(k\) and natural number \(c\) and natural number \(v\), if \(tx - 1 \le k\), then \(\operatorname{rdrv}(td, tx, v \cdot \operatorname{twopow}(k + 1 + c) + c, v \cdot \operatorname{twopow}(k + 1 + c)) = \operatorname{center}(td + c)\).
\end{theorem}
\begin{proof}
The goal follows from \(\operatorname{rdrvsubdiagattwopowalleqcenter}(td, tx, k, c, v)\).
\end{proof}
\begin{theorem}
\label{ll:paper:rdrv-subdiag-at-twopow-all-eq-center-of-threshold}
For every natural number \(td\) and natural number \(tx\) and natural number \(k\) and natural number \(J\) and natural number \(c\), if \(\operatorname{rightthreshold}(k, J)\), then if \(tx - 1 \le J\), then \(\operatorname{rdrv}(td, tx, \operatorname{twopow}(k + 1 + c) + c, \operatorname{twopow}(k + 1 + c)) = \operatorname{center}(td + c)\).
\end{theorem}
\begin{proof}
The goal follows from \(\operatorname{rdrvsubdiagattwopowalleqcenterofthreshold}(td, tx, k, J, c)\).
\end{proof}
\begin{theorem}
\label{ll:paper:rhcenter-subdiag-at-twopow-all-eq-center}
For every natural number \(t\) and natural number \(k\) and natural number \(c\) and natural number \(v\), if \(t - 1 \le k\), then \(\operatorname{rhcenter}(t + (v \cdot \operatorname{twopow}(k + 1 + c) + c), v \cdot \operatorname{twopow}(k + 1 + c)) = \operatorname{center}(t + c)\).
\end{theorem}
\begin{proof}
The goal follows from \(\operatorname{rhcentersubdiagattwopowalleqcenter}(t, k, c, v)\).
\end{proof}
\begin{theorem}
\label{ll:paper:rhcenter-subdiag-at-twopow-all-eq-center-of-threshold}
For every natural number \(t\) and natural number \(k\) and natural number \(J\) and natural number \(c\), if \(\operatorname{rightthreshold}(k, J)\), then if \(t - 1 \le J\), then \(\operatorname{rhcenter}(t + (\operatorname{twopow}(k + 1 + c) + c), \operatorname{twopow}(k + 1 + c)) = \operatorname{center}(t + c)\).
\end{theorem}
\begin{proof}
The goal follows from \(\operatorname{rhcentersubdiagattwopowalleqcenterofthreshold}(t, k, J, c)\).
\end{proof}
\begin{theorem}
\label{ll:paper:rdrv-subdiag-at-twopow-all-agree}
For every natural number \(td\) and natural number \(tx2\) and natural number \(tx\) and natural number \(k\) and natural number \(c\) and natural number \(v\), if \(tx2 - 1 \le k\), then if \(tx - 1 \le k\), then \(\operatorname{rdrv}(td, tx2, v \cdot \operatorname{twopow}(k + 1 + c) + c, v \cdot \operatorname{twopow}(k + 1 + c)) = \operatorname{rdrv}(td, tx, v \cdot \operatorname{twopow}(k + 1 + c) + c, v \cdot \operatorname{twopow}(k + 1 + c))\).
\end{theorem}
\begin{proof}
The goal follows from \(\operatorname{rdrvsubdiagattwopowallagree}(td, tx2, tx, k, c, v)\).
\end{proof}
\begin{theorem}
\label{ll:paper:rdrv-col-one-eq-subdiag-at-twopow-all}
For every natural number \(td\) and natural number \(tx\) and natural number \(k\) and natural number \(c\) and natural number \(v\), if \(tx - 1 \le k\), then \(\operatorname{rdrv}(td, tx, 1 + c, 1) = \operatorname{rdrv}(td, tx, 1 + v \cdot \operatorname{twopow}(k + 1 + c) + c, 1 + v \cdot \operatorname{twopow}(k + 1 + c))\).
\end{theorem}
\begin{proof}
The goal follows from \(\operatorname{rdrvcoloneeqsubdiagattwopowall}(td, tx, k, c, v)\).
\end{proof}
\begin{theorem}
\label{ll:paper:rdrv-col-one-eq-subdiag-at-twopow-all-of-threshold}
For every natural number \(td\) and natural number \(tx\) and natural number \(k\) and natural number \(J\) and natural number \(c\), if \(\operatorname{rightthreshold}(k, J)\), then if \(tx - 1 \le J\), then \(\operatorname{rdrv}(td, tx, 1 + c, 1) = \operatorname{rdrv}(td, tx, 1 + \operatorname{twopow}(k + 1 + c) + c, 1 + \operatorname{twopow}(k + 1 + c))\).
\end{theorem}
\begin{proof}
The goal follows from \(\operatorname{rdrvcoloneeqsubdiagattwopowallofthreshold}(td, tx, k, J, c)\).
\end{proof}
\begin{theorem}
\label{ll:paper:rhcenter-col-one-eq-subdiag-at-twopow-all}
For every natural number \(t\) and natural number \(k\) and natural number \(c\) and natural number \(v\), if \(t - 1 \le k\), then \(\operatorname{rhcenter}(t + (1 + c), 1) = \operatorname{rhcenter}(t + (1 + v \cdot \operatorname{twopow}(k + 1 + c) + c), 1 + v \cdot \operatorname{twopow}(k + 1 + c))\).
\end{theorem}
\begin{proof}
The goal follows from \(\operatorname{rhcentercoloneeqsubdiagattwopowall}(t, k, c, v)\).
\end{proof}
\begin{theorem}
\label{ll:paper:rhcenter-col-one-eq-subdiag-at-twopow-all-of-threshold}
For every natural number \(t\) and natural number \(k\) and natural number \(J\) and natural number \(c\), if \(\operatorname{rightthreshold}(k, J)\), then if \(t - 1 \le J\), then \(\operatorname{rhcenter}(t + (1 + c), 1) = \operatorname{rhcenter}(t + (1 + \operatorname{twopow}(k + 1 + c) + c), 1 + \operatorname{twopow}(k + 1 + c))\).
\end{theorem}
\begin{proof}
The goal follows from \(\operatorname{rhcentercoloneeqsubdiagattwopowallofthreshold}(t, k, J, c)\).
\end{proof}
\begin{theorem}
\label{ll:paper:rdrv-col-one-agree-of-subdiag-twopow-agree}
For every natural number \(td\) and natural number \(tx2\) and natural number \(tx\) and natural number \(k\) and natural number \(c\) and natural number \(v\), if \(tx2 - 1 \le k\), then if \(tx - 1 \le k\), then if \(\operatorname{rdrv}(td, tx2, 1 + v \cdot \operatorname{twopow}(k + 1 + c) + c, 1 + v \cdot \operatorname{twopow}(k + 1 + c)) = \operatorname{rdrv}(td, tx, 1 + v \cdot \operatorname{twopow}(k + 1 + c) + c, 1 + v \cdot \operatorname{twopow}(k + 1 + c))\), then \(\operatorname{rdrv}(td, tx2, 1 + c, 1) = \operatorname{rdrv}(td, tx, 1 + c, 1)\).
\end{theorem}
\begin{proof}
The goal follows from \(\operatorname{rdrvcoloneagreeofsubdiagtwopowagree}(td, tx2, tx, k, c, v)\).
\end{proof}
\begin{theorem}
\label{ll:paper:rdrv-subdiag-agree-all-of-lt-twopow-all}
For every natural number \(td\) and natural number \(tx\) and natural number \(Q\) and natural number \(k\) and natural number \(c\) and natural number \(u\), if \(tx + Q - 1 \le k\), then if for every natural number \(r\), if \(r < \operatorname{twopow}(k + 1 + c)\), then \(\operatorname{rdrv}(td, tx + Q, r + c, r) = \operatorname{rdrv}(td, tx, r + c, r)\), then \(\operatorname{rdrv}(td, tx + Q, u + c, u) = \operatorname{rdrv}(td, tx, u + c, u)\).
\end{theorem}
\begin{proof}
The goal follows from \(\operatorname{rdrvsubdiagagreealloflttwopowall}(td, tx, Q, k, c, u)\).
\end{proof}
\begin{theorem}
\label{ll:paper:rdrv-subdiag-period-twopow-all-mul-of-threshold}
For every natural number \(td\) and natural number \(tx\) and natural number \(k\) and natural number \(J\) and natural number \(c\) and natural number \(u\) and natural number \(v\), if \(\operatorname{rightthreshold}(k, J)\), then if \(tx - 1 \le J\), then \(\operatorname{rdrv}(td, tx, u + v \cdot \operatorname{twopow}(k + 1 + c) + c, u + v \cdot \operatorname{twopow}(k + 1 + c)) = \operatorname{rdrv}(td, tx, u + c, u)\).
\end{theorem}
\begin{proof}
The goal follows from \(\operatorname{rdrvsubdiagperiodtwopowallmulofthreshold}(td, tx, k, J, c, u, v)\).
\end{proof}
\begin{theorem}
\label{ll:paper:rdrv-subdiag-agree-all-of-threshold}
For every natural number \(td\) and natural number \(tx\) and natural number \(Q\) and natural number \(k\) and natural number \(J\) and natural number \(c\) and natural number \(u\), if \(\operatorname{rightthreshold}(k, J)\), then if \(tx + Q - 1 \le J\), then if for every natural number \(r\), if \(r < \operatorname{twopow}(k + 1 + c)\), then \(\operatorname{rdrv}(td, tx + Q, r + c, r) = \operatorname{rdrv}(td, tx, r + c, r)\), then \(\operatorname{rdrv}(td, tx + Q, u + c, u) = \operatorname{rdrv}(td, tx, u + c, u)\).
\end{theorem}
\begin{proof}
The goal follows from \(\operatorname{rdrvsubdiagagreeallofthreshold}(td, tx, Q, k, J, c, u)\).
\end{proof}
\begin{theorem}
\label{ll:paper:periodic-maskparity-zero-agree-of-subdiag-twopow-agree}
For every natural number \(T\) and natural number \(P\), if \(0 < P\), then if \(\operatorname{centerperiodic}(T, P)\), then for every natural number \(m\) and natural number \(t\) and natural number \(k\) and natural number \(c\) and natural number \(v\), if \(T \le t\), then if \(t + m \cdot P - 1 \le k\), then if \(\operatorname{rdrv}(t, t + m \cdot P, 1 + v \cdot \operatorname{twopow}(k + 1 + c) + c, 1 + v \cdot \operatorname{twopow}(k + 1 + c)) = \operatorname{rdrv}(t, t, 1 + v \cdot \operatorname{twopow}(k + 1 + c) + c, 1 + v \cdot \operatorname{twopow}(k + 1 + c))\), then \(\operatorname{maskparity}(t + (1 + c) + m \cdot P, 0) = \operatorname{maskparity}(t + (1 + c), 0)\).
\end{theorem}
\begin{proof}
Assume \(hP\), \(hper\).
The goal follows from \(\operatorname{periodicmaskparityzeroagreeofsubdiagtwopowagree}(T, P, hP, hper)\).
\end{proof}
\begin{theorem}
\label{ll:paper:twopow-add-eq-mul}
For every natural number \(a\) and natural number \(b\), \(\operatorname{twopow}(a + b) = \operatorname{twopow}(b) \cdot \operatorname{twopow}(a)\).
\end{theorem}
\begin{proof}
The goal follows from \(\operatorname{twopowaddeqmul}(a, b)\).
\end{proof}
\begin{theorem}
\label{ll:paper:twopow-subdiag-scale-split}
For every natural number \(k\) and natural number \(c\) and natural number \(C\) and natural number \(v\), if \(c \le C\), then \(v \cdot \operatorname{twopow}(k + 1 + C) = v \cdot \operatorname{twopow}(C - c) \cdot \operatorname{twopow}(k + 1 + c)\).
\end{theorem}
\begin{proof}
The goal follows from \(\operatorname{twopowsubdiagscalesplit}(k, c, C, v)\).
\end{proof}
\begin{theorem}
\label{ll:paper:rdrv-dyadic-quadrant-eq}
For every natural number \(td\) and natural number \(tx\) and natural number \(k\) and natural number \(C\) and natural number \(v\) and natural number \(j\) and natural number \(q\), if \(tx - 1 \le k\), then if \(j \le q + C\), then \(\operatorname{rdrv}(td, tx, v \cdot \operatorname{twopow}(k + 1 + C) + j, v \cdot \operatorname{twopow}(k + 1 + C) + q) = \operatorname{rdrv}(td, tx, j, q)\).
\end{theorem}
\begin{proof}
The goal follows from \(\operatorname{rdrvdyadicquadranteq}(td, tx, k, C, v, j, q)\).
\end{proof}
\begin{theorem}
\label{ll:paper:rdrv-dyadic-quadrant-eq-of-threshold}
For every natural number \(td\) and natural number \(tx\) and natural number \(k\) and natural number \(J\) and natural number \(C\) and natural number \(v\) and natural number \(j\) and natural number \(q\), if \(\operatorname{rightthreshold}(k, J)\), then if \(tx - 1 \le J\), then if \(j \le q + C\), then \(\operatorname{rdrv}(td, tx, v \cdot \operatorname{twopow}(k + 1 + C) + j, v \cdot \operatorname{twopow}(k + 1 + C) + q) = \operatorname{rdrv}(td, tx, j, q)\).
\end{theorem}
\begin{proof}
The goal follows from \(\operatorname{rdrvdyadicquadranteqofthreshold}(td, tx, k, J, C, v, j, q)\).
\end{proof}
\begin{theorem}
\label{ll:paper:rdrv-dyadic-slab-eq}
For every natural number \(td\) and natural number \(tx\) and natural number \(k\) and natural number \(C\) and natural number \(v\) and natural number \(j\) and natural number \(q\), if \(tx - 1 \le k\), then if \(j \le C\), then \(\operatorname{rdrv}(td, tx, v \cdot \operatorname{twopow}(k + 1 + C) + j, v \cdot \operatorname{twopow}(k + 1 + C) + q) = \operatorname{rdrv}(td, tx, j, q)\).
\end{theorem}
\begin{proof}
The goal follows from \(\operatorname{rdrvdyadicslabeq}(td, tx, k, C, v, j, q)\).
\end{proof}
\begin{theorem}
\label{ll:paper:rdrv-dyadic-row-eq-init}
For every natural number \(td\) and natural number \(tx\) and natural number \(k\) and natural number \(v\) and natural number \(q\), if \(tx - 1 \le k\), then \(\operatorname{rdrv}(td, tx, v \cdot \operatorname{twopow}(k + 1), v \cdot \operatorname{twopow}(k + 1) + q) = \operatorname{rdrv}(td, tx, 0, q)\).
\end{theorem}
\begin{proof}
The goal follows from \(\operatorname{rdrvdyadicroweqinit}(td, tx, k, v, q)\).
\end{proof}
\begin{theorem}
\label{ll:paper:rhcenter-dyadic-quadrant-eq}
For every natural number \(t\) and natural number \(k\) and natural number \(C\) and natural number \(v\) and natural number \(j\) and natural number \(q\), if \(t - 1 \le k\), then if \(j \le q + C\), then \(\operatorname{rhcenter}(t + (v \cdot \operatorname{twopow}(k + 1 + C) + j), v \cdot \operatorname{twopow}(k + 1 + C) + q) = \operatorname{rhcenter}(t + j, q)\).
\end{theorem}
\begin{proof}
The goal follows from \(\operatorname{rhcenterdyadicquadranteq}(t, k, C, v, j, q)\).
\end{proof}
\begin{theorem}
\label{ll:paper:rhcenter-dyadic-quadrant-eq-of-threshold}
For every natural number \(t\) and natural number \(k\) and natural number \(J\) and natural number \(C\) and natural number \(v\) and natural number \(j\) and natural number \(q\), if \(\operatorname{rightthreshold}(k, J)\), then if \(t - 1 \le J\), then if \(j \le q + C\), then \(\operatorname{rhcenter}(t + (v \cdot \operatorname{twopow}(k + 1 + C) + j), v \cdot \operatorname{twopow}(k + 1 + C) + q) = \operatorname{rhcenter}(t + j, q)\).
\end{theorem}
\begin{proof}
The goal follows from \(\operatorname{rhcenterdyadicquadranteqofthreshold}(t, k, J, C, v, j, q)\).
\end{proof}
\begin{theorem}
\label{ll:paper:rhcenter-dyadic-slab-eq}
For every natural number \(t\) and natural number \(k\) and natural number \(C\) and natural number \(v\) and natural number \(j\) and natural number \(q\), if \(t - 1 \le k\), then if \(j \le C\), then \(\operatorname{rhcenter}(t + (v \cdot \operatorname{twopow}(k + 1 + C) + j), v \cdot \operatorname{twopow}(k + 1 + C) + q) = \operatorname{rhcenter}(t + j, q)\).
\end{theorem}
\begin{proof}
The goal follows from \(\operatorname{rhcenterdyadicslabeq}(t, k, C, v, j, q)\).
\end{proof}
\begin{theorem}
\label{ll:paper:rhcenter-dyadic-row-eq}
For every natural number \(t\) and natural number \(k\) and natural number \(v\) and natural number \(q\), if \(t - 1 \le k\), then \(\operatorname{rhcenter}(t + v \cdot \operatorname{twopow}(k + 1), v \cdot \operatorname{twopow}(k + 1) + q) = \operatorname{rhcenter}(t, q)\).
\end{theorem}
\begin{proof}
The goal follows from \(\operatorname{rhcenterdyadicroweq}(t, k, v, q)\).
\end{proof}
\begin{theorem}
\label{ll:paper:rhcenter-dyadic-row-eq-of-threshold}
For every natural number \(t\) and natural number \(k\) and natural number \(J\) and natural number \(v\) and natural number \(q\), if \(\operatorname{rightthreshold}(k, J)\), then if \(t - 1 \le J\), then \(\operatorname{rhcenter}(t + v \cdot \operatorname{twopow}(k + 1), v \cdot \operatorname{twopow}(k + 1) + q) = \operatorname{rhcenter}(t, q)\).
\end{theorem}
\begin{proof}
The goal follows from \(\operatorname{rhcenterdyadicroweqofthreshold}(t, k, J, v, q)\).
\end{proof}
\begin{theorem}
\label{ll:paper:rhbandagree-of-dyadic-row-shift-agree}
For every natural number \(t\) and natural number \(Q\) and natural number \(k\) and natural number \(v\) and natural number \(a\) and natural number \(B\), if \(t + Q - 1 \le k\), then if \(\operatorname{rhbandagree}(t + v \cdot \operatorname{twopow}(k + 1), Q, v \cdot \operatorname{twopow}(k + 1) + a, B)\), then \(\operatorname{rhbandagree}(t, Q, a, B)\).
\end{theorem}
\begin{proof}
The goal follows from \(\operatorname{rhbandagreeofdyadicrowshiftagree}(t, Q, k, v, a, B)\).
\end{proof}
\begin{theorem}
\label{ll:paper:rhbandagree-dyadic-row-shift-of-rhbandagree}
For every natural number \(t\) and natural number \(Q\) and natural number \(k\) and natural number \(v\) and natural number \(a\) and natural number \(B\), if \(t + Q - 1 \le k\), then if \(\operatorname{rhbandagree}(t, Q, a, B)\), then \(\operatorname{rhbandagree}(t + v \cdot \operatorname{twopow}(k + 1), Q, v \cdot \operatorname{twopow}(k + 1) + a, B)\).
\end{theorem}
\begin{proof}
The goal follows from \(\operatorname{rhbandagreedyadicrowshiftofrhbandagree}(t, Q, k, v, a, B)\).
\end{proof}
\begin{theorem}
\label{ll:paper:periodic-maskparity-zero-agree-of-dyadic-row-shift-agree}
For every natural number \(T\) and natural number \(P\), if \(0 < P\), then if \(\operatorname{centerperiodic}(T, P)\), then for every natural number \(m\) and natural number \(t\) and natural number \(k\) and natural number \(v\) and natural number \(B\) and natural number \(s\), if \(T \le t\), then if \(t + m \cdot P - 1 \le k\), then if \(\operatorname{rhbandagree}(t + v \cdot \operatorname{twopow}(k + 1), m \cdot P, v \cdot \operatorname{twopow}(k + 1) + 1, B)\), then if \(s \le B\), then \(\operatorname{maskparity}(t + s + m \cdot P, 0) = \operatorname{maskparity}(t + s, 0)\).
\end{theorem}
\begin{proof}
Assume \(hP\), \(hper\).
The goal follows from \(\operatorname{periodicmaskparityzeroagreeofdyadicrowshiftagree}(T, P, hP, hper)\).
\end{proof}
\begin{theorem}
\label{ll:paper:not-centerperiodic-of-dyadic-row-shift-agree-at-zero-diag-doubles}
For every natural number \(T\) and natural number \(P\) and natural number \(m\) and natural number \(z\) and natural number \(p\) and natural number \(k\) and natural number \(v\) and natural number \(B\), if \(0 < P\), then if \(0 < m\), then if \(T \le z + (p + 1)\), then if \(\operatorname{zerodiagonal}(z)\), then if \(\operatorname{doublesat}(m \cdot P, z)\), then if \(\operatorname{settledby}(z + 1, z + (p + 1))\), then if \(\operatorname{settledby}(z + 1 - m \cdot P, z + (p + 1))\), then if \(z + (p + 1) + m \cdot P - 1 \le k\), then if \(p + 3 \cdot (m \cdot P) \le B\), then if \(\operatorname{rhbandagree}(z + (p + 1) + v \cdot \operatorname{twopow}(k + 1), m \cdot P, v \cdot \operatorname{twopow}(k + 1) + 1, B)\), then not \(\operatorname{centerperiodic}(T, P)\).
\end{theorem}
\begin{proof}
The goal follows from \(\operatorname{notcenterperiodicofdyadicrowshiftagreeatzerodiagdoubles}(T, P, m, z, p, k, v, B)\).
\end{proof}
\begin{theorem}
\label{ll:paper:rdrv-white-seed-subdiag-one-step-eq-center}
For every natural number \(td\) and natural number \(u\), \(\operatorname{rdrv}(td, 0, u + 1 + 1, u + 1) = \operatorname{bxor}(\operatorname{rdrv}(td, 0, u + 1, u), \operatorname{center}(td))\).
\end{theorem}
\begin{proof}
The goal follows from \(\operatorname{rdrvwhiteseedsubdiagonestepeqcenter}(td, u)\).
\end{proof}
\begin{theorem}
\label{ll:paper:rdrv-white-seed-subdiag-one-period-two-all}
For every natural number \(td\) and natural number \(u\), \(\operatorname{rdrv}(td, 0, u + 2 + 1, u + 2) = \operatorname{rdrv}(td, 0, u + 1, u)\).
\end{theorem}
\begin{proof}
The goal follows from \(\operatorname{rdrvwhiteseedsubdiagoneperiodtwoall}(td, u)\).
\end{proof}
\begin{theorem}
\label{ll:paper:rdrv-white-seed-subdiag-period-twopow-all}
For every natural number \(td\) and natural number \(c\) and natural number \(u\), \(\operatorname{rdrv}(td, 0, u + \operatorname{twopow}(c) + c, u + \operatorname{twopow}(c)) = \operatorname{rdrv}(td, 0, u + c, u)\).
\end{theorem}
\begin{proof}
The goal follows from \(\operatorname{rdrvwhiteseedsubdiagperiodtwopowall}(td, c, u)\).
\end{proof}
\begin{theorem}
\label{ll:paper:rdrv-white-seed-subdiag-period-twopow-all-mul}
For every natural number \(td\) and natural number \(c\) and natural number \(u\) and natural number \(v\), \(\operatorname{rdrv}(td, 0, u + v \cdot \operatorname{twopow}(c) + c, u + v \cdot \operatorname{twopow}(c)) = \operatorname{rdrv}(td, 0, u + c, u)\).
\end{theorem}
\begin{proof}
The goal follows from \(\operatorname{rdrvwhiteseedsubdiagperiodtwopowallmul}(td, c, u, v)\).
\end{proof}
\begin{theorem}
\label{ll:paper:rdrv-white-seed-subdiag-at-twopow-eq-center}
For every natural number \(td\) and natural number \(c\) and natural number \(v\), \(\operatorname{rdrv}(td, 0, v \cdot \operatorname{twopow}(c) + c, v \cdot \operatorname{twopow}(c)) = \operatorname{center}(td + c)\).
\end{theorem}
\begin{proof}
The goal follows from \(\operatorname{rdrvwhiteseedsubdiagattwopoweqcenter}(td, c, v)\).
\end{proof}
\begin{theorem}
\label{ll:paper:rdrv-white-seed-col-one-eq-subdiag-at-twopow}
For every natural number \(td\) and natural number \(c\) and natural number \(v\), \(\operatorname{rdrv}(td, 0, 1 + c, 1) = \operatorname{rdrv}(td, 0, 1 + v \cdot \operatorname{twopow}(c) + c, 1 + v \cdot \operatorname{twopow}(c))\).
\end{theorem}
\begin{proof}
The goal follows from \(\operatorname{rdrvwhiteseedcoloneeqsubdiagattwopow}(td, c, v)\).
\end{proof}
\begin{theorem}
\label{ll:paper:rdrv-white-seed-dyadic-quadrant-eq}
For every natural number \(td\) and natural number \(C\) and natural number \(v\) and natural number \(j\) and natural number \(q\), if \(j \le q + C\), then \(\operatorname{rdrv}(td, 0, v \cdot \operatorname{twopow}(C) + j, v \cdot \operatorname{twopow}(C) + q) = \operatorname{rdrv}(td, 0, j, q)\).
\end{theorem}
\begin{proof}
The goal follows from \(\operatorname{rdrvwhiteseeddyadicquadranteq}(td, C, v, j, q)\).
\end{proof}
\begin{theorem}
\label{ll:paper:rdrv-white-seed-dyadic-slab-eq}
For every natural number \(td\) and natural number \(C\) and natural number \(v\) and natural number \(j\) and natural number \(q\), if \(j \le C\), then \(\operatorname{rdrv}(td, 0, v \cdot \operatorname{twopow}(C) + j, v \cdot \operatorname{twopow}(C) + q) = \operatorname{rdrv}(td, 0, j, q)\).
\end{theorem}
\begin{proof}
The goal follows from \(\operatorname{rdrvwhiteseeddyadicslabeq}(td, C, v, j, q)\).
\end{proof}
\begin{theorem}
\label{ll:paper:rhcenter-subdiag-period-twopow-zero-delay}
For every natural number \(c\) and natural number \(u\), \(\operatorname{rhcenter}(u + \operatorname{twopow}(c) + c, u + \operatorname{twopow}(c)) = \operatorname{rhcenter}(u + c, u)\).
\end{theorem}
\begin{proof}
The goal follows from \(\operatorname{rhcentersubdiagperiodtwopowzerodelay}(c, u)\).
\end{proof}
\begin{theorem}
\label{ll:paper:rhcenter-subdiag-period-twopow-zero-delay-mul}
For every natural number \(c\) and natural number \(u\) and natural number \(v\), \(\operatorname{rhcenter}(u + v \cdot \operatorname{twopow}(c) + c, u + v \cdot \operatorname{twopow}(c)) = \operatorname{rhcenter}(u + c, u)\).
\end{theorem}
\begin{proof}
The goal follows from \(\operatorname{rhcentersubdiagperiodtwopowzerodelaymul}(c, u, v)\).
\end{proof}
\begin{theorem}
\label{ll:paper:rhcenter-subdiag-at-twopow-zero-delay-eq-center}
For every natural number \(c\) and natural number \(v\), \(\operatorname{rhcenter}(v \cdot \operatorname{twopow}(c) + c, v \cdot \operatorname{twopow}(c)) = \operatorname{center}(c)\).
\end{theorem}
\begin{proof}
The goal follows from \(\operatorname{rhcentersubdiagattwopowzerodelayeqcenter}(c, v)\).
\end{proof}
\begin{theorem}
\label{ll:paper:rhcenter-col-one-eq-subdiag-at-twopow-zero-delay}
For every natural number \(c\) and natural number \(v\), \(\operatorname{rhcenter}(1 + c, 1) = \operatorname{rhcenter}(1 + v \cdot \operatorname{twopow}(c) + c, 1 + v \cdot \operatorname{twopow}(c))\).
\end{theorem}
\begin{proof}
The goal follows from \(\operatorname{rhcentercoloneeqsubdiagattwopowzerodelay}(c, v)\).
\end{proof}
\begin{theorem}
\label{ll:paper:rhcenter-dyadic-quadrant-zero-delay-eq}
For every natural number \(C\) and natural number \(v\) and natural number \(j\) and natural number \(q\), if \(j \le q + C\), then \(\operatorname{rhcenter}(v \cdot \operatorname{twopow}(C) + j, v \cdot \operatorname{twopow}(C) + q) = \operatorname{rhcenter}(j, q)\).
\end{theorem}
\begin{proof}
The goal follows from \(\operatorname{rhcenterdyadicquadrantzerodelayeq}(C, v, j, q)\).
\end{proof}
\begin{theorem}
\label{ll:paper:rhcenter-dyadic-slab-zero-delay-eq}
For every natural number \(C\) and natural number \(v\) and natural number \(j\) and natural number \(q\), if \(j \le C\), then \(\operatorname{rhcenter}(v \cdot \operatorname{twopow}(C) + j, v \cdot \operatorname{twopow}(C) + q) = \operatorname{rhcenter}(j, q)\).
\end{theorem}
\begin{proof}
The goal follows from \(\operatorname{rhcenterdyadicslabzerodelayeq}(C, v, j, q)\).
\end{proof}
\begin{theorem}
\label{ll:paper:rhcenter-dyadic-triangle-white-corridor}
For every natural number \(C\) and natural number \(v\) and natural number \(j\) and natural number \(q\), if \(j < q\), then \(\operatorname{rhcenter}(v \cdot \operatorname{twopow}(C) + j, v \cdot \operatorname{twopow}(C) + q) = \operatorname{false}\).
\end{theorem}
\begin{proof}
The goal follows from \(\operatorname{rhcenterdyadictrianglewhitecorridor}(C, v, j, q)\).
\end{proof}
\begin{theorem}
\label{ll:paper:rdrv-subdiag-period-twopow-tx-add}
For every natural number \(td\) and natural number \(tx\) and natural number \(c\) and natural number \(u\), \(\operatorname{rdrv}(td, tx, u + \operatorname{twopow}(tx + c) + c, u + \operatorname{twopow}(tx + c)) = \operatorname{rdrv}(td, tx, u + c, u)\).
\end{theorem}
\begin{proof}
The goal follows from \(\operatorname{rdrvsubdiagperiodtwopowtxadd}(td, tx, c, u)\).
\end{proof}
\begin{theorem}
\label{ll:paper:rdrv-subdiag-period-twopow-tx-add-mul}
For every natural number \(td\) and natural number \(tx\) and natural number \(c\) and natural number \(u\) and natural number \(v\), \(\operatorname{rdrv}(td, tx, u + v \cdot \operatorname{twopow}(tx + c) + c, u + v \cdot \operatorname{twopow}(tx + c)) = \operatorname{rdrv}(td, tx, u + c, u)\).
\end{theorem}
\begin{proof}
The goal follows from \(\operatorname{rdrvsubdiagperiodtwopowtxaddmul}(td, tx, c, u, v)\).
\end{proof}
\begin{theorem}
\label{ll:paper:rdrv-main-diag-period-twopow-tx}
For every natural number \(td\) and natural number \(tx\) and natural number \(u\) and natural number \(v\), \(\operatorname{rdrv}(td, tx, u + v \cdot \operatorname{twopow}(tx), u + v \cdot \operatorname{twopow}(tx)) = \operatorname{rdrv}(td, tx, u, u)\).
\end{theorem}
\begin{proof}
The goal follows from \(\operatorname{rdrvmaindiagperiodtwopowtx}(td, tx, u, v)\).
\end{proof}
\begin{theorem}
\label{ll:paper:rdrv-main-diag-at-twopow-tx-eq-center}
For every natural number \(td\) and natural number \(tx\) and natural number \(v\), \(\operatorname{rdrv}(td, tx, v \cdot \operatorname{twopow}(tx), v \cdot \operatorname{twopow}(tx)) = \operatorname{center}(td)\).
\end{theorem}
\begin{proof}
The goal follows from \(\operatorname{rdrvmaindiagattwopowtxeqcenter}(td, tx, v)\).
\end{proof}
\begin{theorem}
\label{ll:paper:rdrv-subdiag-at-twopow-tx-add-eq-center}
For every natural number \(td\) and natural number \(tx\) and natural number \(c\) and natural number \(v\), \(\operatorname{rdrv}(td, tx, v \cdot \operatorname{twopow}(tx + c) + c, v \cdot \operatorname{twopow}(tx + c)) = \operatorname{center}(td + c)\).
\end{theorem}
\begin{proof}
The goal follows from \(\operatorname{rdrvsubdiagattwopowtxaddeqcenter}(td, tx, c, v)\).
\end{proof}
\begin{theorem}
\label{ll:paper:rdrv-col-one-eq-subdiag-at-twopow-tx-add}
For every natural number \(td\) and natural number \(tx\) and natural number \(c\) and natural number \(v\), \(\operatorname{rdrv}(td, tx, 1 + c, 1) = \operatorname{rdrv}(td, tx, 1 + v \cdot \operatorname{twopow}(tx + c) + c, 1 + v \cdot \operatorname{twopow}(tx + c))\).
\end{theorem}
\begin{proof}
The goal follows from \(\operatorname{rdrvcoloneeqsubdiagattwopowtxadd}(td, tx, c, v)\).
\end{proof}
\begin{theorem}
\label{ll:paper:rdrv-dyadic-quadrant-tx-add-eq}
For every natural number \(td\) and natural number \(tx\) and natural number \(C\) and natural number \(v\) and natural number \(j\) and natural number \(q\), if \(j \le q + C\), then \(\operatorname{rdrv}(td, tx, v \cdot \operatorname{twopow}(tx + C) + j, v \cdot \operatorname{twopow}(tx + C) + q) = \operatorname{rdrv}(td, tx, j, q)\).
\end{theorem}
\begin{proof}
The goal follows from \(\operatorname{rdrvdyadicquadranttxaddeq}(td, tx, C, v, j, q)\).
\end{proof}
\begin{theorem}
\label{ll:paper:rdrv-dyadic-slab-tx-add-eq}
For every natural number \(td\) and natural number \(tx\) and natural number \(C\) and natural number \(v\) and natural number \(j\) and natural number \(q\), if \(j \le C\), then \(\operatorname{rdrv}(td, tx, v \cdot \operatorname{twopow}(tx + C) + j, v \cdot \operatorname{twopow}(tx + C) + q) = \operatorname{rdrv}(td, tx, j, q)\).
\end{theorem}
\begin{proof}
The goal follows from \(\operatorname{rdrvdyadicslabtxaddeq}(td, tx, C, v, j, q)\).
\end{proof}
\begin{theorem}
\label{ll:paper:rdrv-dyadic-row-tx-eq-init}
For every natural number \(td\) and natural number \(tx\) and natural number \(v\) and natural number \(q\), \(\operatorname{rdrv}(td, tx, v \cdot \operatorname{twopow}(tx), v \cdot \operatorname{twopow}(tx) + q) = \operatorname{rdrv}(td, tx, 0, q)\).
\end{theorem}
\begin{proof}
The goal follows from \(\operatorname{rdrvdyadicrowtxeqinit}(td, tx, v, q)\).
\end{proof}
\begin{theorem}
\label{ll:paper:rhcenter-dyadic-quadrant-t-add-eq}
For every natural number \(t\) and natural number \(C\) and natural number \(v\) and natural number \(j\) and natural number \(q\), if \(j \le q + C\), then \(\operatorname{rhcenter}(t + (v \cdot \operatorname{twopow}(t + C) + j), v \cdot \operatorname{twopow}(t + C) + q) = \operatorname{rhcenter}(t + j, q)\).
\end{theorem}
\begin{proof}
The goal follows from \(\operatorname{rhcenterdyadicquadranttaddeq}(t, C, v, j, q)\).
\end{proof}
\begin{theorem}
\label{ll:paper:rhcenter-dyadic-slab-t-add-eq}
For every natural number \(t\) and natural number \(C\) and natural number \(v\) and natural number \(j\) and natural number \(q\), if \(j \le C\), then \(\operatorname{rhcenter}(t + (v \cdot \operatorname{twopow}(t + C) + j), v \cdot \operatorname{twopow}(t + C) + q) = \operatorname{rhcenter}(t + j, q)\).
\end{theorem}
\begin{proof}
The goal follows from \(\operatorname{rhcenterdyadicslabtaddeq}(t, C, v, j, q)\).
\end{proof}
\begin{theorem}
\label{ll:paper:rhcenter-dyadic-row-t-eq}
For every natural number \(t\) and natural number \(v\) and natural number \(q\), \(\operatorname{rhcenter}(t + v \cdot \operatorname{twopow}(t), v \cdot \operatorname{twopow}(t) + q) = \operatorname{rhcenter}(t, q)\).
\end{theorem}
\begin{proof}
The goal follows from \(\operatorname{rhcenterdyadicrowteq}(t, v, q)\).
\end{proof}
\begin{theorem}
\label{ll:paper:rhcenter-dyadic-row-scale-eq}
For every natural number \(t\) and natural number \(K\) and natural number \(v\) and natural number \(q\), if \(t \le K\), then \(\operatorname{rhcenter}(t + v \cdot \operatorname{twopow}(K), v \cdot \operatorname{twopow}(K) + q) = \operatorname{rhcenter}(t, q)\).
\end{theorem}
\begin{proof}
The goal follows from \(\operatorname{rhcenterdyadicrowscaleeq}(t, K, v, q)\).
\end{proof}
\begin{theorem}
\label{ll:paper:rhbandagree-of-dyadic-row-scale-agree}
For every natural number \(t\) and natural number \(Q\) and natural number \(K\) and natural number \(v\) and natural number \(a\) and natural number \(B\), if \(t + Q \le K\), then if \(\operatorname{rhbandagree}(t + v \cdot \operatorname{twopow}(K), Q, v \cdot \operatorname{twopow}(K) + a, B)\), then \(\operatorname{rhbandagree}(t, Q, a, B)\).
\end{theorem}
\begin{proof}
The goal follows from \(\operatorname{rhbandagreeofdyadicrowscaleagree}(t, Q, K, v, a, B)\).
\end{proof}
\begin{theorem}
\label{ll:paper:periodic-maskparity-zero-agree-of-dyadic-row-scale-agree}
For every natural number \(T\) and natural number \(P\), if \(0 < P\), then if \(\operatorname{centerperiodic}(T, P)\), then for every natural number \(m\) and natural number \(t\) and natural number \(K\) and natural number \(v\) and natural number \(B\) and natural number \(s\), if \(T \le t\), then if \(t + m \cdot P \le K\), then if \(\operatorname{rhbandagree}(t + v \cdot \operatorname{twopow}(K), m \cdot P, v \cdot \operatorname{twopow}(K) + 1, B)\), then if \(s \le B\), then \(\operatorname{maskparity}(t + s + m \cdot P, 0) = \operatorname{maskparity}(t + s, 0)\).
\end{theorem}
\begin{proof}
Assume \(hP\), \(hper\).
The goal follows from \(\operatorname{periodicmaskparityzeroagreeofdyadicrowscaleagree}(T, P, hP, hper)\).
\end{proof}
\begin{theorem}
\label{ll:paper:rdrv-init-gt-eq-false}
For every natural number \(td\) and natural number \(tx\) and natural number \(q\), if \(tx < q\), then \(\operatorname{rdrv}(td, tx, 0, q) = \operatorname{false}\).
\end{theorem}
\begin{proof}
The goal follows from \(\operatorname{rdrvinitgteqfalse}(td, tx, q)\).
\end{proof}
\begin{theorem}
\label{ll:paper:rdrv-dyadic-slab-tx-gt-eq-false}
For every natural number \(td\) and natural number \(tx\) and natural number \(C\) and natural number \(v\) and natural number \(j\) and natural number \(q\), if \(j \le C\), then if \(tx + j < q\), then \(\operatorname{rdrv}(td, tx, v \cdot \operatorname{twopow}(tx + C) + j, v \cdot \operatorname{twopow}(tx + C) + q) = \operatorname{false}\).
\end{theorem}
\begin{proof}
The goal follows from \(\operatorname{rdrvdyadicslabtxgteqfalse}(td, tx, C, v, j, q)\).
\end{proof}
\begin{theorem}
\label{ll:paper:rdrv-init-agree-of-gt}
For every natural number \(td\) and natural number \(tx2\) and natural number \(tx\) and natural number \(q\), if \(tx2 < q\), then if \(tx < q\), then \(\operatorname{rdrv}(td, tx2, 0, q) = \operatorname{rdrv}(td, tx, 0, q)\).
\end{theorem}
\begin{proof}
The goal follows from \(\operatorname{rdrvinitagreeofgt}(td, tx2, tx, q)\).
\end{proof}
\begin{theorem}
\label{ll:paper:rdrv-agree-of-gt}
For every natural number \(td\) and natural number \(tx2\) and natural number \(tx\) and natural number \(k\) and natural number \(p\), if \(tx2 + k < p\), then if \(tx + k < p\), then \(\operatorname{rdrv}(td, tx2, k, p) = \operatorname{rdrv}(td, tx, k, p)\).
\end{theorem}
\begin{proof}
The goal follows from \(\operatorname{rdrvagreeofgt}(td, tx2, tx, k, p)\).
\end{proof}
\begin{theorem}
\label{ll:paper:rhbandagree-all-of-support-agree}
For every natural number \(tx\) and natural number \(Q\) and natural number \(B\), if \(\operatorname{rhbandagree}(tx, Q, 1, tx + Q)\), then \(\operatorname{rhbandagree}(tx, Q, 1, B)\).
\end{theorem}
\begin{proof}
The goal follows from \(\operatorname{rhbandagreeallofsupportagree}(tx, Q, B)\).
\end{proof}
\begin{theorem}
\label{ll:paper:rdrv-agree-all-of-support-rhbandagree}
For every natural number \(td\) and natural number \(tx\) and natural number \(Q\) and natural number \(k\) and natural number \(p\), if \(\operatorname{rhbandagree}(tx, Q, 1, tx + Q)\), then \(\operatorname{rdrv}(td, tx + Q, k, p) = \operatorname{rdrv}(td, tx, k, p)\).
\end{theorem}
\begin{proof}
The goal follows from \(\operatorname{rdrvagreeallofsupportrhbandagree}(td, tx, Q, k, p)\).
\end{proof}
\begin{theorem}
\label{ll:paper:periodic-rhcenter-agree-all-of-support-rhbandagree}
For every natural number \(T\) and natural number \(P\), if \(0 < P\), then if \(\operatorname{centerperiodic}(T, P)\), then for every natural number \(m\) and natural number \(t\) and natural number \(s\) and natural number \(p\), if \(T \le t\), then if \(\operatorname{rhbandagree}(t, m \cdot P, 1, t + m \cdot P)\), then \(\operatorname{rhcenter}(t + s + m \cdot P, p) = \operatorname{rhcenter}(t + s, p)\).
\end{theorem}
\begin{proof}
Assume \(hP\), \(hper\).
The goal follows from \(\operatorname{periodicrhcenteragreeallofsupportrhbandagree}(T, P, hP, hper)\).
\end{proof}
\begin{theorem}
\label{ll:paper:periodic-maskparity-zero-agree-all-of-support-rhbandagree}
For every natural number \(T\) and natural number \(P\), if \(0 < P\), then if \(\operatorname{centerperiodic}(T, P)\), then for every natural number \(m\) and natural number \(t\) and natural number \(s\), if \(T \le t\), then if \(\operatorname{rhbandagree}(t, m \cdot P, 1, t + m \cdot P)\), then \(\operatorname{maskparity}(t + s + m \cdot P, 0) = \operatorname{maskparity}(t + s, 0)\).
\end{theorem}
\begin{proof}
Assume \(hP\), \(hper\).
The goal follows from \(\operatorname{periodicmaskparityzeroagreeallofsupportrhbandagree}(T, P, hP, hper)\).
\end{proof}
\begin{theorem}
\label{ll:paper:not-centerperiodic-of-support-rhbandagree-at-zero-diag-doubles}
For every natural number \(T\) and natural number \(P\) and natural number \(m\) and natural number \(z\) and natural number \(p\), if \(0 < P\), then if \(0 < m\), then if \(T \le z + (p + 1)\), then if \(\operatorname{zerodiagonal}(z)\), then if \(\operatorname{doublesat}(m \cdot P, z)\), then if \(\operatorname{settledby}(z + 1, z + (p + 1))\), then if \(\operatorname{settledby}(z + 1 - m \cdot P, z + (p + 1))\), then if \(\operatorname{rhbandagree}(z + (p + 1), m \cdot P, 1, z + (p + 1) + m \cdot P)\), then not \(\operatorname{centerperiodic}(T, P)\).
\end{theorem}
\begin{proof}
The goal follows from \(\operatorname{notcenterperiodicofsupportrhbandagreeatzerodiagdoubles}(T, P, m, z, p)\).
\end{proof}
\begin{theorem}
\label{ll:paper:not-centerperiodic-of-support-rhbandagree-twopow}
For every natural number \(T\) and natural number \(P\) and natural number \(m\) and natural number \(t\) and natural number \(k\), if \(0 < P\), then if \(0 < m\), then if \(T \le t\), then if \(t + m \cdot P = \operatorname{twopow}(k)\), then if \(\operatorname{rhbandagree}(t, m \cdot P, 1, t + m \cdot P)\), then not \(\operatorname{centerperiodic}(T, P)\).
\end{theorem}
\begin{proof}
The goal follows from \(\operatorname{notcenterperiodicofsupportrhbandagreetwopow}(T, P, m, t, k)\).
\end{proof}
\begin{theorem}
\label{ll:paper:not-centerperiodic-of-support-rhbandagree-at-zero-diag-pair}
For every natural number \(T\) and natural number \(P\) and natural number \(m\) and natural number \(z\) and natural number \(k\), if \(0 < P\), then if \(\operatorname{zerodiagonal}(z)\), then if \(2 \le z + m \cdot P\), then if \(\operatorname{zerodiagonal}(z + m \cdot P - 2)\), then if \(T \le z + \operatorname{twopow}(k)\), then if \(\operatorname{settledby}(z, z + \operatorname{twopow}(k))\), then if \(\operatorname{settledby}(z + m \cdot P, z + m \cdot P + \operatorname{twopow}(k))\), then if \(\operatorname{rhbandagree}(z + \operatorname{twopow}(k), m \cdot P, 1, z + \operatorname{twopow}(k) + m \cdot P)\), then not \(\operatorname{centerperiodic}(T, P)\).
\end{theorem}
\begin{proof}
The goal follows from \(\operatorname{notcenterperiodicofsupportrhbandagreeatzerodiagpair}(T, P, m, z, k)\).
\end{proof}
\begin{theorem}
\label{ll:paper:rdrv-dyadic-row-tx-gt-eq-false}
For every natural number \(td\) and natural number \(tx\) and natural number \(v\) and natural number \(q\), if \(tx < q\), then \(\operatorname{rdrv}(td, tx, v \cdot \operatorname{twopow}(tx), v \cdot \operatorname{twopow}(tx) + q) = \operatorname{false}\).
\end{theorem}
\begin{proof}
The goal follows from \(\operatorname{rdrvdyadicrowtxgteqfalse}(td, tx, v, q)\).
\end{proof}
\begin{theorem}
\label{ll:paper:rhcenter-dyadic-row-t-gt-eq-false}
For every natural number \(t\) and natural number \(v\) and natural number \(q\), if \(t < q\), then \(\operatorname{rhcenter}(t + v \cdot \operatorname{twopow}(t), v \cdot \operatorname{twopow}(t) + q) = \operatorname{false}\).
\end{theorem}
\begin{proof}
The goal follows from \(\operatorname{rhcenterdyadicrowtgteqfalse}(t, v, q)\).
\end{proof}
\begin{theorem}
\label{ll:paper:rhbandagree-all-of-dyadic-row-support-agree}
For every natural number \(t\) and natural number \(Q\) and natural number \(K\) and natural number \(v\) and natural number \(B\), if \(t + Q \le K\), then if \(\operatorname{rhbandagree}(t + v \cdot \operatorname{twopow}(K), Q, v \cdot \operatorname{twopow}(K) + 1, t + Q)\), then \(\operatorname{rhbandagree}(t, Q, 1, B)\).
\end{theorem}
\begin{proof}
The goal follows from \(\operatorname{rhbandagreeallofdyadicrowsupportagree}(t, Q, K, v, B)\).
\end{proof}
\begin{theorem}
\label{ll:paper:periodic-maskparity-zero-agree-all-of-dyadic-row-support-agree}
For every natural number \(T\) and natural number \(P\), if \(0 < P\), then if \(\operatorname{centerperiodic}(T, P)\), then for every natural number \(m\) and natural number \(t\) and natural number \(K\) and natural number \(v\) and natural number \(s\), if \(T \le t\), then if \(t + m \cdot P \le K\), then if \(\operatorname{rhbandagree}(t + v \cdot \operatorname{twopow}(K), m \cdot P, v \cdot \operatorname{twopow}(K) + 1, t + m \cdot P)\), then \(\operatorname{maskparity}(t + s + m \cdot P, 0) = \operatorname{maskparity}(t + s, 0)\).
\end{theorem}
\begin{proof}
Assume \(hP\), \(hper\).
The goal follows from \(\operatorname{periodicmaskparityzeroagreeallofdyadicrowsupportagree}(T, P, hP, hper)\).
\end{proof}
\begin{theorem}
\label{ll:paper:rdrv-left-step-inv}
For every natural number \(td\) and natural number \(tx\) and natural number \(k\) and natural number \(p\), \(\operatorname{rdrv}(td, tx, k, p) = \operatorname{bxor}(\operatorname{rdrv}(td, tx, k + 1, p + 1), \operatorname{rdrv}(td, tx, k, p + 1) \operatorname{or} \operatorname{rdrv}(td, tx, k, p + 2))\).
\end{theorem}
\begin{proof}
The goal follows from \(\operatorname{rdrvleftstepinv}(td, tx, k, p)\).
\end{proof}
\begin{theorem}
\label{ll:paper:rdrv-left-step-inv-sub}
For every natural number \(td\) and natural number \(tx\) and natural number \(k\) and natural number \(U\) and natural number \(d\), if \(d + 1 \le U\), then \(\operatorname{rdrv}(td, tx, k, U - (d + 1)) = \operatorname{bxor}(\operatorname{rdrv}(td, tx, k + 1, U - d), \operatorname{rdrv}(td, tx, k, U - d) \operatorname{or} \operatorname{rdrv}(td, tx, k, U - d + 1))\).
\end{theorem}
\begin{proof}
The goal follows from \(\operatorname{rdrvleftstepinvsub}(td, tx, k, U, d)\).
\end{proof}
\begin{theorem}
\label{ll:paper:rhcenter-left-step-inv}
For every natural number \(t\) and natural number \(p\), \(\operatorname{rhcenter}(t, p) = \operatorname{bxor}(\operatorname{rhcenter}(t + 1, p + 1), \operatorname{rhcenter}(t, p + 1) \operatorname{or} \operatorname{rhcenter}(t, p + 2))\).
\end{theorem}
\begin{proof}
The goal follows from \(\operatorname{rhcenterleftstepinv}(t, p)\).
\end{proof}
\begin{theorem}
\label{ll:paper:rhcenter-left-step-inv-sub}
For every natural number \(t\) and natural number \(U\) and natural number \(d\), if \(d + 1 \le U\), then \(\operatorname{rhcenter}(t, U - (d + 1)) = \operatorname{bxor}(\operatorname{rhcenter}(t + 1, U - d), \operatorname{rhcenter}(t, U - d) \operatorname{or} \operatorname{rhcenter}(t, U - d + 1))\).
\end{theorem}
\begin{proof}
The goal follows from \(\operatorname{rhcenterleftstepinvsub}(t, U, d)\).
\end{proof}
\begin{theorem}
\label{ll:paper:rdrv-dyadic-pred-col-eq-center-xor}
For every natural number \(td\) and natural number \(tx\) and natural number \(C\) and natural number \(v\) and natural number \(j\), if \(1 \le v\), then if \(j + 1 \le C\), then \(\operatorname{rdrv}(td, tx, v \cdot \operatorname{twopow}(tx + C) + j, v \cdot \operatorname{twopow}(tx + C) - 1) = \operatorname{bxor}(\operatorname{center}(td + (j + 1)), \operatorname{center}(td + j) \operatorname{or} \operatorname{rdrv}(td, tx, j, 1))\).
\end{theorem}
\begin{proof}
The goal follows from \(\operatorname{rdrvdyadicpredcoleqcenterxor}(td, tx, C, v, j)\).
\end{proof}
\begin{theorem}
\label{ll:paper:rhcenter-dyadic-pred-col-t-add-eq-lhcenter-one}
For every natural number \(t\) and natural number \(C\) and natural number \(v\) and natural number \(j\), if \(1 \le v\), then if \(j + 1 \le C\), then \(\operatorname{rhcenter}(t + (v \cdot \operatorname{twopow}(t + C) + j), v \cdot \operatorname{twopow}(t + C) - 1) = \operatorname{lhcenter}(t + j, 1)\).
\end{theorem}
\begin{proof}
The goal follows from \(\operatorname{rhcenterdyadicpredcoltaddeqlhcenterone}(t, C, v, j)\).
\end{proof}
\begin{theorem}
\label{ll:paper:rhcenter-dyadic-pred-col-zero-delay-eq-lhcenter-one}
For every natural number \(C\) and natural number \(v\) and natural number \(j\), if \(1 \le v\), then if \(j + 1 \le C\), then \(\operatorname{rhcenter}(v \cdot \operatorname{twopow}(C) + j, v \cdot \operatorname{twopow}(C) - 1) = \operatorname{lhcenter}(j, 1)\).
\end{theorem}
\begin{proof}
The goal follows from \(\operatorname{rhcenterdyadicpredcolzerodelayeqlhcenterone}(C, v, j)\).
\end{proof}
\begin{theorem}
\label{ll:paper:rhcenter-sub-eq-lhcenter-of-col-eq-center-and-one}
For every natural number \(t\) and natural number \(T0\) and natural number \(U\) and natural number \(C\) and natural number \(d\) and natural number \(j\), if \(d \le U\), then if \(j + d \le C\), then if for every natural number \(s\), if \(s \le C\), then \(\operatorname{rhcenter}(T0 + s, U) = \operatorname{center}(t + s)\), then if for every natural number \(s\), if \(s + 1 \le C\), then \(\operatorname{rhcenter}(T0 + s, U - 1) = \operatorname{lhcenter}(t + s, 1)\), then \(\operatorname{rhcenter}(T0 + j, U - d) = \operatorname{lhcenter}(t + j, d)\).
\end{theorem}
\begin{proof}
The goal follows from \(\operatorname{rhcentersubeqlhcenterofcoleqcenterandone}(t, T0, U, C, d, j)\).
\end{proof}
\begin{theorem}
\label{ll:paper:rhcenter-dyadic-sub-eq-lhcenter-t-add}
For every natural number \(t\) and natural number \(C\) and natural number \(v\) and natural number \(d\) and natural number \(j\), if \(1 \le v\), then if \(d \le v \cdot \operatorname{twopow}(t + C)\), then if \(j + d \le C\), then \(\operatorname{rhcenter}(t + (v \cdot \operatorname{twopow}(t + C) + j), v \cdot \operatorname{twopow}(t + C) - d) = \operatorname{lhcenter}(t + j, d)\).
\end{theorem}
\begin{proof}
The goal follows from \(\operatorname{rhcenterdyadicsubeqlhcentertadd}(t, C, v, d, j)\).
\end{proof}
\begin{theorem}
\label{ll:paper:rhcenter-dyadic-sub-eq-lhcenter-zero-delay}
For every natural number \(C\) and natural number \(v\) and natural number \(d\) and natural number \(j\), if \(1 \le v\), then if \(d \le v \cdot \operatorname{twopow}(C)\), then if \(j + d \le C\), then \(\operatorname{rhcenter}(v \cdot \operatorname{twopow}(C) + j, v \cdot \operatorname{twopow}(C) - d) = \operatorname{lhcenter}(j, d)\).
\end{theorem}
\begin{proof}
The goal follows from \(\operatorname{rhcenterdyadicsubeqlhcenterzerodelay}(C, v, d, j)\).
\end{proof}
\begin{theorem}
\label{ll:paper:rhcenter-dyadic-left-white-corridor-t-add}
For every natural number \(t\) and natural number \(C\) and natural number \(v\) and natural number \(d\) and natural number \(j\), if \(1 \le v\), then if \(d \le v \cdot \operatorname{twopow}(t + C)\), then if \(j + d \le C\), then if \(t + j < d\), then \(\operatorname{rhcenter}(t + (v \cdot \operatorname{twopow}(t + C) + j), v \cdot \operatorname{twopow}(t + C) - d) = \operatorname{false}\).
\end{theorem}
\begin{proof}
The goal follows from \(\operatorname{rhcenterdyadicleftwhitecorridortadd}(t, C, v, d, j)\).
\end{proof}
\begin{theorem}
\label{ll:paper:rhcenter-dyadic-left-white-corridor-zero-delay}
For every natural number \(C\) and natural number \(v\) and natural number \(d\) and natural number \(j\), if \(1 \le v\), then if \(d \le v \cdot \operatorname{twopow}(C)\), then if \(j + d \le C\), then if \(j < d\), then \(\operatorname{rhcenter}(v \cdot \operatorname{twopow}(C) + j, v \cdot \operatorname{twopow}(C) - d) = \operatorname{false}\).
\end{theorem}
\begin{proof}
The goal follows from \(\operatorname{rhcenterdyadicleftwhitecorridorzerodelay}(C, v, d, j)\).
\end{proof}
\begin{theorem}
\label{ll:paper:rhcenter-dyadic-row-left-white-corridor-t-add}
For every natural number \(t\) and natural number \(C\) and natural number \(v\) and natural number \(d\), if \(1 \le v\), then if \(d \le C\), then if \(t < d\), then \(\operatorname{rhcenter}(t + v \cdot \operatorname{twopow}(t + C), v \cdot \operatorname{twopow}(t + C) - d) = \operatorname{false}\).
\end{theorem}
\begin{proof}
The goal follows from \(\operatorname{rhcenterdyadicrowleftwhitecorridortadd}(t, C, v, d)\).
\end{proof}
\begin{theorem}
\label{ll:paper:rhcenter-dyadic-row-left-white-corridor-zero-delay}
For every natural number \(C\) and natural number \(v\) and natural number \(d\), if \(1 \le v\), then if \(1 \le d\), then if \(d \le C\), then \(\operatorname{rhcenter}(v \cdot \operatorname{twopow}(C), v \cdot \operatorname{twopow}(C) - d) = \operatorname{false}\).
\end{theorem}
\begin{proof}
The goal follows from \(\operatorname{rhcenterdyadicrowleftwhitecorridorzerodelay}(C, v, d)\).
\end{proof}
\begin{theorem}
\label{ll:paper:rhcenter-dyadic-sub-eq-ldiag-t-add}
For every natural number \(t\) and natural number \(C\) and natural number \(v\) and natural number \(d\) and natural number \(j\), if \(1 \le v\), then if \(d \le v \cdot \operatorname{twopow}(t + C)\), then if \(j + d \le C\), then if \(d \le t + j\), then \(\operatorname{rhcenter}(t + (v \cdot \operatorname{twopow}(t + C) + j), v \cdot \operatorname{twopow}(t + C) - d) = \operatorname{ldiag}(t + j - d, t + j)\).
\end{theorem}
\begin{proof}
The goal follows from \(\operatorname{rhcenterdyadicsubeqldiagtadd}(t, C, v, d, j)\).
\end{proof}
\begin{theorem}
\label{ll:paper:rhcenter-dyadic-sub-eq-lhcenter-of-le}
For every natural number \(K\) and natural number \(v\) and natural number \(d\) and natural number \(j\), if \(1 \le v\), then if \(j + d \le K\), then \(\operatorname{rhcenter}(v \cdot \operatorname{twopow}(K) + j, v \cdot \operatorname{twopow}(K) - d) = \operatorname{lhcenter}(j, d)\).
\end{theorem}
\begin{proof}
The goal follows from \(\operatorname{rhcenterdyadicsubeqlhcenterofle}(K, v, d, j)\).
\end{proof}
\begin{theorem}
\label{ll:paper:rhcenter-dyadic-sub-shift-eq-lhcenter}
For every natural number \(t\) and natural number \(K\) and natural number \(v\) and natural number \(d\) and natural number \(s\), if \(1 \le v\), then if \(t + s + d \le K\), then \(\operatorname{rhcenter}(t + v \cdot \operatorname{twopow}(K) + s, v \cdot \operatorname{twopow}(K) - d) = \operatorname{lhcenter}(t + s, d)\).
\end{theorem}
\begin{proof}
The goal follows from \(\operatorname{rhcenterdyadicsubshifteqlhcenter}(t, K, v, d, s)\).
\end{proof}
\begin{theorem}
\label{ll:paper:rhcenter-dyadic-sub-period-shift-eq-lhcenter}
For every natural number \(t\) and natural number \(Q\) and natural number \(K\) and natural number \(v\) and natural number \(d\) and natural number \(s\), if \(1 \le v\), then if \(t + s + Q + d \le K\), then \(\operatorname{rhcenter}(t + v \cdot \operatorname{twopow}(K) + s + Q, v \cdot \operatorname{twopow}(K) - d) = \operatorname{lhcenter}(t + s + Q, d)\).
\end{theorem}
\begin{proof}
The goal follows from \(\operatorname{rhcenterdyadicsubperiodshifteqlhcenter}(t, Q, K, v, d, s)\).
\end{proof}
\begin{theorem}
\label{ll:paper:rhcenter-dyadic-pred-col-shift-eq-lhcenter-one}
For every natural number \(t\) and natural number \(K\) and natural number \(v\) and natural number \(s\), if \(1 \le v\), then if \(t + s + 1 \le K\), then \(\operatorname{rhcenter}(t + v \cdot \operatorname{twopow}(K) + s, v \cdot \operatorname{twopow}(K) - 1) = \operatorname{lhcenter}(t + s, 1)\).
\end{theorem}
\begin{proof}
The goal follows from \(\operatorname{rhcenterdyadicpredcolshifteqlhcenterone}(t, K, v, s)\).
\end{proof}
\begin{theorem}
\label{ll:paper:lhcenter-shift-agree-of-rhcenter-dyadic-sub-agree}
For every natural number \(t\) and natural number \(Q\) and natural number \(K\) and natural number \(v\) and natural number \(d\) and natural number \(s\), if \(1 \le v\), then if \(t + s + Q + d \le K\), then if \(\operatorname{rhcenter}(t + v \cdot \operatorname{twopow}(K) + s + Q, v \cdot \operatorname{twopow}(K) - d) = \operatorname{rhcenter}(t + v \cdot \operatorname{twopow}(K) + s, v \cdot \operatorname{twopow}(K) - d)\), then \(\operatorname{lhcenter}(t + s + Q, d) = \operatorname{lhcenter}(t + s, d)\).
\end{theorem}
\begin{proof}
The goal follows from \(\operatorname{lhcentershiftagreeofrhcenterdyadicsubagree}(t, Q, K, v, d, s)\).
\end{proof}
\begin{theorem}
\label{ll:paper:lhcolagree-of-rhcolagree-dyadic-sub}
For every natural number \(t\) and natural number \(Q\) and natural number \(K\) and natural number \(v\) and natural number \(d\) and natural number \(S\), if \(1 \le v\), then if \(t + S + Q + d \le K\), then if \(\operatorname{rhcolagree}(t + v \cdot \operatorname{twopow}(K), Q, v \cdot \operatorname{twopow}(K) - d, S)\), then \(\operatorname{lhcolagree}(t, Q, d, S)\).
\end{theorem}
\begin{proof}
The goal follows from \(\operatorname{lhcolagreeofrhcolagreedyadicsub}(t, Q, K, v, d, S)\).
\end{proof}
\begin{theorem}
\label{ll:paper:lhcolagree-one-of-rhcolagree-dyadic-pred-col}
For every natural number \(t\) and natural number \(Q\) and natural number \(K\) and natural number \(v\) and natural number \(S\), if \(1 \le v\), then if \(t + S + Q + 1 \le K\), then if \(\operatorname{rhcolagree}(t + v \cdot \operatorname{twopow}(K), Q, v \cdot \operatorname{twopow}(K) - 1, S)\), then \(\operatorname{lhcolagree}(t, Q, 1, S)\).
\end{theorem}
\begin{proof}
The goal follows from \(\operatorname{lhcolagreeoneofrhcolagreedyadicpredcol}(t, Q, K, v, S)\).
\end{proof}
\begin{theorem}
\label{ll:paper:periodic-maskparity-zero-agree-of-rhcolagree-dyadic-pred-col}
For every natural number \(T\) and natural number \(P\), if \(0 < P\), then if \(\operatorname{centerperiodic}(T, P)\), then for every natural number \(m\) and natural number \(t\) and natural number \(K\) and natural number \(v\) and natural number \(S\) and natural number \(s\), if \(T \le t\), then if \(1 \le v\), then if \(t + S + m \cdot P + 1 \le K\), then if \(\operatorname{rhcolagree}(t + v \cdot \operatorname{twopow}(K), m \cdot P, v \cdot \operatorname{twopow}(K) - 1, S)\), then if \(s \le S\), then \(\operatorname{maskparity}(t + s + m \cdot P, 0) = \operatorname{maskparity}(t + s, 0)\).
\end{theorem}
\begin{proof}
Assume \(hP\), \(hper\), \(m\), \(t\), \(K\), \(v\), \(S\), \(s\).
The goal follows from \(\operatorname{periodicmaskparityzeroagreeofrhcolagreedyadicpredcol}(T, P, hP, hper, m, t, K, v, S, s)\).
\end{proof}
\begin{theorem}
\label{ll:paper:not-centerperiodic-of-rhcolagree-dyadic-pred-col-at-zero-diag-doubles}
For every natural number \(T\) and natural number \(P\) and natural number \(m\) and natural number \(z\) and natural number \(p\) and natural number \(K\) and natural number \(v\) and natural number \(S\), if \(0 < P\), then if \(0 < m\), then if \(T \le z + (p + 1)\), then if \(\operatorname{zerodiagonal}(z)\), then if \(\operatorname{doublesat}(m \cdot P, z)\), then if \(\operatorname{settledby}(z + 1, z + (p + 1))\), then if \(\operatorname{settledby}(z + 1 - m \cdot P, z + (p + 1))\), then if \(p + 3 \cdot (m \cdot P) \le S + 1\), then if \(1 \le v\), then if \(z + (p + 1) + S + m \cdot P + 1 \le K\), then if \(\operatorname{rhcolagree}(z + (p + 1) + v \cdot \operatorname{twopow}(K), m \cdot P, v \cdot \operatorname{twopow}(K) - 1, S)\), then not \(\operatorname{centerperiodic}(T, P)\).
\end{theorem}
\begin{proof}
The goal follows from \(\operatorname{notcenterperiodicofrhcolagreedyadicpredcolatzerodiagdoubles}(T, P, m, z, p, K, v, S)\).
\end{proof}
\begin{theorem}
\label{ll:paper:not-centerperiodic-of-rhcolagree-dyadic-pred-col-twopow}
For every natural number \(T\) and natural number \(P\) and natural number \(m\) and natural number \(t\) and natural number \(k\) and natural number \(K\) and natural number \(v\) and natural number \(S\), if \(0 < P\), then if \(0 < m\), then if \(T \le t\), then if \(t + m \cdot P = \operatorname{twopow}(k)\), then if \(\operatorname{twopow}(k) \le S + 1\), then if \(1 \le v\), then if \(t + S + m \cdot P + 1 \le K\), then if \(\operatorname{rhcolagree}(t + v \cdot \operatorname{twopow}(K), m \cdot P, v \cdot \operatorname{twopow}(K) - 1, S)\), then not \(\operatorname{centerperiodic}(T, P)\).
\end{theorem}
\begin{proof}
The goal follows from \(\operatorname{notcenterperiodicofrhcolagreedyadicpredcoltwopow}(T, P, m, t, k, K, v, S)\).
\end{proof}
\begin{theorem}
\label{ll:paper:not-centerperiodic-of-rhcolagree-dyadic-pred-col-at-zero-diag-pair}
For every natural number \(T\) and natural number \(P\) and natural number \(m\) and natural number \(z\) and natural number \(k\) and natural number \(K\) and natural number \(v\) and natural number \(S\), if \(0 < P\), then if \(\operatorname{zerodiagonal}(z)\), then if \(2 \le z + m \cdot P\), then if \(\operatorname{zerodiagonal}(z + m \cdot P - 2)\), then if \(T \le z + \operatorname{twopow}(k)\), then if \(\operatorname{settledby}(z, z + \operatorname{twopow}(k))\), then if \(\operatorname{settledby}(z + m \cdot P, z + m \cdot P + \operatorname{twopow}(k))\), then if \(\operatorname{twopow}(k) \le S + 1\), then if \(1 \le v\), then if \(z + \operatorname{twopow}(k) + S + m \cdot P + 1 \le K\), then if \(\operatorname{rhcolagree}(z + \operatorname{twopow}(k) + v \cdot \operatorname{twopow}(K), m \cdot P, v \cdot \operatorname{twopow}(K) - 1, S)\), then not \(\operatorname{centerperiodic}(T, P)\).
\end{theorem}
\begin{proof}
The goal follows from \(\operatorname{notcenterperiodicofrhcolagreedyadicpredcolatzerodiagpair}(T, P, m, z, k, K, v, S)\).
\end{proof}
\begin{theorem}
\label{ll:paper:rhcenter-dyadic-sub-eq-ldiag-of-le}
For every natural number \(K\) and natural number \(v\) and natural number \(d\) and natural number \(j\), if \(1 \le v\), then if \(j + d \le K\), then if \(d \le j\), then \(\operatorname{rhcenter}(v \cdot \operatorname{twopow}(K) + j, v \cdot \operatorname{twopow}(K) - d) = \operatorname{ldiag}(j - d, j)\).
\end{theorem}
\begin{proof}
The goal follows from \(\operatorname{rhcenterdyadicsubeqldiagofle}(K, v, d, j)\).
\end{proof}
\begin{theorem}
\label{ll:paper:rhcenter-dyadic-sub-zero-diag-eq-false}
For every natural number \(K\) and natural number \(v\) and natural number \(z\) and natural number \(d\), if \(1 \le v\), then if \(\operatorname{zerodiagonal}(z)\), then if \(\operatorname{settledby}(z, z + d)\), then if \(z + d + d \le K\), then \(\operatorname{rhcenter}(v \cdot \operatorname{twopow}(K) + (z + d), v \cdot \operatorname{twopow}(K) - d) = \operatorname{false}\).
\end{theorem}
\begin{proof}
The goal follows from \(\operatorname{rhcenterdyadicsubzerodiageqfalse}(K, v, z, d)\).
\end{proof}
\begin{theorem}
\label{ll:paper:rhcenter-dyadic-sub-after-zero-diag-eq-true}
For every natural number \(K\) and natural number \(v\) and natural number \(z\) and natural number \(d\), if \(1 \le v\), then if \(\operatorname{zerodiagonal}(z)\), then if \(\operatorname{settledby}(z + 2, z + 2 + d)\), then if \(z + 2 + d + d \le K\), then \(\operatorname{rhcenter}(v \cdot \operatorname{twopow}(K) + (z + 2 + d), v \cdot \operatorname{twopow}(K) - d) = \operatorname{true}\).
\end{theorem}
\begin{proof}
The goal follows from \(\operatorname{rhcenterdyadicsubafterzerodiageqtrue}(K, v, z, d)\).
\end{proof}
\begin{theorem}
\label{ll:paper:rhcenter-dyadic-sub-doubling-diag-antiperiodic}
For every natural number \(K\) and natural number \(v\) and natural number \(Q\) and natural number \(z\) and natural number \(d\), if \(1 \le v\), then if \(\operatorname{zerodiagonal}(z)\), then if \(\operatorname{doublesat}(Q, z)\), then if \(\operatorname{settledby}(z + 1, z + 1 + d)\), then if \(z + 1 + d + Q + (d + Q) \le K\), then \(\operatorname{rhcenter}(v \cdot \operatorname{twopow}(K) + (z + 1 + d + Q), v \cdot \operatorname{twopow}(K) - (d + Q)) = \operatorname{not}(\operatorname{rhcenter}(v \cdot \operatorname{twopow}(K) + (z + 1 + d), v \cdot \operatorname{twopow}(K) - d))\).
\end{theorem}
\begin{proof}
The goal follows from \(\operatorname{rhcenterdyadicsubdoublingdiagantiperiodic}(K, v, Q, z, d)\).
\end{proof}
\begin{theorem}
\label{ll:paper:rhcenter-step-succ-interior}
For every natural number \(t\) and natural number \(p\), \(\operatorname{rhcenter}(t + 1, p + 1) = \operatorname{bxor}(\operatorname{rhcenter}(t, p), \operatorname{rhcenter}(t, p + 1) \operatorname{or} \operatorname{rhcenter}(t, p + 2))\).
\end{theorem}
\begin{proof}
The goal follows from \(\operatorname{rhcenterstepsuccinterior}(t, p)\).
\end{proof}
\begin{theorem}
\label{ll:paper:rhcenter-bilateral-cone-agree-step}
For every natural number \(t2\) and natural number \(t1\) and natural number \(a\) and natural number \(L\) and natural number \(i\), if for every natural number \(q\), if \(q \le L + 2\), then \(\operatorname{rhcenter}(t2, a + q) = \operatorname{rhcenter}(t1, a + q)\), then if \(i \le L\), then \(\operatorname{rhcenter}(t2 + 1, a + 1 + i) = \operatorname{rhcenter}(t1 + 1, a + 1 + i)\).
\end{theorem}
\begin{proof}
The goal follows from \(\operatorname{rhcenterbilateralconeagreestep}(t2, t1, a, L, i)\).
\end{proof}
\begin{theorem}
\label{ll:paper:rhcenter-bilateral-cone-agree}
For every natural number \(t2\) and natural number \(t1\) and natural number \(a\) and natural number \(L\) and natural number \(s\) and natural number \(i\), if for every natural number \(q\), if \(q \le L + 2 \cdot s\), then \(\operatorname{rhcenter}(t2, a + q) = \operatorname{rhcenter}(t1, a + q)\), then if \(i \le L\), then \(\operatorname{rhcenter}(t2 + s, a + s + i) = \operatorname{rhcenter}(t1 + s, a + s + i)\).
\end{theorem}
\begin{proof}
The goal follows from \(\operatorname{rhcenterbilateralconeagree}(t2, t1, a, L, s, i)\).
\end{proof}
\begin{theorem}
\label{ll:paper:rhcenter-bilateral-cone-point-agree}
For every natural number \(t2\) and natural number \(t1\) and natural number \(P0\) and natural number \(s\), if \(s \le P0\), then if for every natural number \(q\), if \(q \le 2 \cdot s\), then \(\operatorname{rhcenter}(t2, P0 - s + q) = \operatorname{rhcenter}(t1, P0 - s + q)\), then \(\operatorname{rhcenter}(t2 + s, P0) = \operatorname{rhcenter}(t1 + s, P0)\).
\end{theorem}
\begin{proof}
The goal follows from \(\operatorname{rhcenterbilateralconepointagree}(t2, t1, P0, s)\).
\end{proof}
\begin{theorem}
\label{ll:paper:rhcolagree-of-rhbandagree-bilateral-cone}
For every natural number \(T0\) and natural number \(Q\) and natural number \(P0\) and natural number \(S\), if \(S \le P0\), then if \(\operatorname{rhbandagree}(T0, Q, P0 - S, 2 \cdot S)\), then \(\operatorname{rhcolagree}(T0, Q, P0, S)\).
\end{theorem}
\begin{proof}
The goal follows from \(\operatorname{rhcolagreeofrhbandagreebilateralcone}(T0, Q, P0, S)\).
\end{proof}
\begin{theorem}
\label{ll:paper:lhcolagree-of-rhbandagree-dyadic-sub-cone}
For every natural number \(t\) and natural number \(Q\) and natural number \(K\) and natural number \(v\) and natural number \(d\) and natural number \(S\), if \(1 \le v\), then if \(t + S + Q + d \le K\), then if \(\operatorname{rhbandagree}(t + v \cdot \operatorname{twopow}(K), Q, v \cdot \operatorname{twopow}(K) - d - S, 2 \cdot S)\), then \(\operatorname{lhcolagree}(t, Q, d, S)\).
\end{theorem}
\begin{proof}
The goal follows from \(\operatorname{lhcolagreeofrhbandagreedyadicsubcone}(t, Q, K, v, d, S)\).
\end{proof}
\begin{theorem}
\label{ll:paper:lhcolagree-one-of-rhbandagree-dyadic-pred-col-cone}
For every natural number \(t\) and natural number \(Q\) and natural number \(K\) and natural number \(v\) and natural number \(S\), if \(1 \le v\), then if \(t + S + Q + 1 \le K\), then if \(\operatorname{rhbandagree}(t + v \cdot \operatorname{twopow}(K), Q, v \cdot \operatorname{twopow}(K) - 1 - S, 2 \cdot S)\), then \(\operatorname{lhcolagree}(t, Q, 1, S)\).
\end{theorem}
\begin{proof}
The goal follows from \(\operatorname{lhcolagreeoneofrhbandagreedyadicpredcolcone}(t, Q, K, v, S)\).
\end{proof}
\begin{theorem}
\label{ll:paper:periodic-maskparity-zero-agree-of-rhbandagree-dyadic-pred-col-cone}
For every natural number \(T\) and natural number \(P\), if \(0 < P\), then if \(\operatorname{centerperiodic}(T, P)\), then for every natural number \(m\) and natural number \(t\) and natural number \(K\) and natural number \(v\) and natural number \(S\) and natural number \(s\), if \(T \le t\), then if \(1 \le v\), then if \(t + S + m \cdot P + 1 \le K\), then if \(\operatorname{rhbandagree}(t + v \cdot \operatorname{twopow}(K), m \cdot P, v \cdot \operatorname{twopow}(K) - 1 - S, 2 \cdot S)\), then if \(s \le S\), then \(\operatorname{maskparity}(t + s + m \cdot P, 0) = \operatorname{maskparity}(t + s, 0)\).
\end{theorem}
\begin{proof}
Assume \(hP\), \(hper\), \(m\), \(t\), \(K\), \(v\), \(S\), \(s\).
The goal follows from \(\operatorname{periodicmaskparityzeroagreeofrhbandagreedyadicpredcolcone}(T, P, hP, hper, m, t, K, v, S, s)\).
\end{proof}
\begin{theorem}
\label{ll:paper:not-centerperiodic-of-rhbandagree-dyadic-pred-col-cone-at-zero-diag-doubles}
For every natural number \(T\) and natural number \(P\) and natural number \(m\) and natural number \(z\) and natural number \(p\) and natural number \(K\) and natural number \(v\) and natural number \(S\), if \(0 < P\), then if \(0 < m\), then if \(T \le z + (p + 1)\), then if \(\operatorname{zerodiagonal}(z)\), then if \(\operatorname{doublesat}(m \cdot P, z)\), then if \(\operatorname{settledby}(z + 1, z + (p + 1))\), then if \(\operatorname{settledby}(z + 1 - m \cdot P, z + (p + 1))\), then if \(p + 3 \cdot (m \cdot P) \le S + 1\), then if \(1 \le v\), then if \(z + (p + 1) + S + m \cdot P + 1 \le K\), then if \(\operatorname{rhbandagree}(z + (p + 1) + v \cdot \operatorname{twopow}(K), m \cdot P, v \cdot \operatorname{twopow}(K) - 1 - S, 2 \cdot S)\), then not \(\operatorname{centerperiodic}(T, P)\).
\end{theorem}
\begin{proof}
The goal follows from \(\operatorname{notcenterperiodicofrhbandagreedyadicpredcolconeatzerodiagdoubles}(T, P, m, z, p, K, v, S)\).
\end{proof}
\begin{theorem}
\label{ll:paper:not-centerperiodic-of-rhbandagree-dyadic-pred-col-cone-twopow}
For every natural number \(T\) and natural number \(P\) and natural number \(m\) and natural number \(t\) and natural number \(k\) and natural number \(K\) and natural number \(v\) and natural number \(S\), if \(0 < P\), then if \(0 < m\), then if \(T \le t\), then if \(t + m \cdot P = \operatorname{twopow}(k)\), then if \(\operatorname{twopow}(k) \le S + 1\), then if \(1 \le v\), then if \(t + S + m \cdot P + 1 \le K\), then if \(\operatorname{rhbandagree}(t + v \cdot \operatorname{twopow}(K), m \cdot P, v \cdot \operatorname{twopow}(K) - 1 - S, 2 \cdot S)\), then not \(\operatorname{centerperiodic}(T, P)\).
\end{theorem}
\begin{proof}
The goal follows from \(\operatorname{notcenterperiodicofrhbandagreedyadicpredcolconetwopow}(T, P, m, t, k, K, v, S)\).
\end{proof}
\begin{theorem}
\label{ll:paper:not-centerperiodic-of-rhbandagree-dyadic-pred-col-cone-at-zero-diag-pair}
For every natural number \(T\) and natural number \(P\) and natural number \(m\) and natural number \(z\) and natural number \(k\) and natural number \(K\) and natural number \(v\) and natural number \(S\), if \(0 < P\), then if \(\operatorname{zerodiagonal}(z)\), then if \(2 \le z + m \cdot P\), then if \(\operatorname{zerodiagonal}(z + m \cdot P - 2)\), then if \(T \le z + \operatorname{twopow}(k)\), then if \(\operatorname{settledby}(z, z + \operatorname{twopow}(k))\), then if \(\operatorname{settledby}(z + m \cdot P, z + m \cdot P + \operatorname{twopow}(k))\), then if \(\operatorname{twopow}(k) \le S + 1\), then if \(1 \le v\), then if \(z + \operatorname{twopow}(k) + S + m \cdot P + 1 \le K\), then if \(\operatorname{rhbandagree}(z + \operatorname{twopow}(k) + v \cdot \operatorname{twopow}(K), m \cdot P, v \cdot \operatorname{twopow}(K) - 1 - S, 2 \cdot S)\), then not \(\operatorname{centerperiodic}(T, P)\).
\end{theorem}
\begin{proof}
The goal follows from \(\operatorname{notcenterperiodicofrhbandagreedyadicpredcolconeatzerodiagpair}(T, P, m, z, k, K, v, S)\).
\end{proof}
\begin{theorem}
\label{ll:paper:rhcenter-dyadic-row-left-white-corridor-of-le}
For every natural number \(t\) and natural number \(K\) and natural number \(v\) and natural number \(d\), if \(1 \le v\), then if \(t + d \le K\), then if \(t < d\), then \(\operatorname{rhcenter}(t + v \cdot \operatorname{twopow}(K), v \cdot \operatorname{twopow}(K) - d) = \operatorname{false}\).
\end{theorem}
\begin{proof}
The goal follows from \(\operatorname{rhcenterdyadicrowleftwhitecorridorofle}(t, K, v, d)\).
\end{proof}
\begin{theorem}
\label{ll:paper:rhcenter-dyadic-row-left-white-agree}
For every natural number \(t\) and natural number \(Q\) and natural number \(K\) and natural number \(v\) and natural number \(d\), if \(1 \le v\), then if \(t + Q + d \le K\), then if \(t + Q < d\), then \(\operatorname{rhcenter}(t + v \cdot \operatorname{twopow}(K) + Q, v \cdot \operatorname{twopow}(K) - d) = \operatorname{rhcenter}(t + v \cdot \operatorname{twopow}(K), v \cdot \operatorname{twopow}(K) - d)\).
\end{theorem}
\begin{proof}
The goal follows from \(\operatorname{rhcenterdyadicrowleftwhiteagree}(t, Q, K, v, d)\).
\end{proof}
\begin{theorem}
\label{ll:paper:rhcenter-dyadic-row-right-white-of-le}
For every natural number \(t\) and natural number \(K\) and natural number \(v\) and natural number \(r\), if \(t \le K\), then if \(t < r\), then \(\operatorname{rhcenter}(t + v \cdot \operatorname{twopow}(K), v \cdot \operatorname{twopow}(K) + r) = \operatorname{false}\).
\end{theorem}
\begin{proof}
The goal follows from \(\operatorname{rhcenterdyadicrowrightwhiteofle}(t, K, v, r)\).
\end{proof}
\begin{theorem}
\label{ll:paper:rhcenter-dyadic-row-right-white-agree}
For every natural number \(t\) and natural number \(Q\) and natural number \(K\) and natural number \(v\) and natural number \(r\), if \(t + Q \le K\), then if \(t + Q < r\), then \(\operatorname{rhcenter}(t + v \cdot \operatorname{twopow}(K) + Q, v \cdot \operatorname{twopow}(K) + r) = \operatorname{rhcenter}(t + v \cdot \operatorname{twopow}(K), v \cdot \operatorname{twopow}(K) + r)\).
\end{theorem}
\begin{proof}
The goal follows from \(\operatorname{rhcenterdyadicrowrightwhiteagree}(t, Q, K, v, r)\).
\end{proof}
\begin{theorem}
\label{ll:paper:rhbandagree-dyadic-core-expand-to-cone}
For every natural number \(t\) and natural number \(Q\) and natural number \(K\) and natural number \(v\) and natural number \(S\), if \(1 \le v\), then if \(t + S + Q + (t + Q) + 1 \le K\), then if \(\operatorname{rhbandagree}(t + v \cdot \operatorname{twopow}(K), Q, v \cdot \operatorname{twopow}(K) - (t + Q), 2 \cdot (t + Q))\), then \(\operatorname{rhbandagree}(t + v \cdot \operatorname{twopow}(K), Q, v \cdot \operatorname{twopow}(K) - 1 - S, 2 \cdot S)\).
\end{theorem}
\begin{proof}
The goal follows from \(\operatorname{rhbandagreedyadiccoreexpandtocone}(t, Q, K, v, S)\).
\end{proof}
\begin{theorem}
\label{ll:paper:lhcolagree-one-of-rhbandagree-dyadic-core}
For every natural number \(t\) and natural number \(Q\) and natural number \(K\) and natural number \(v\) and natural number \(S\), if \(1 \le v\), then if \(t + S + Q + (t + Q) + 1 \le K\), then if \(\operatorname{rhbandagree}(t + v \cdot \operatorname{twopow}(K), Q, v \cdot \operatorname{twopow}(K) - (t + Q), 2 \cdot (t + Q))\), then \(\operatorname{lhcolagree}(t, Q, 1, S)\).
\end{theorem}
\begin{proof}
The goal follows from \(\operatorname{lhcolagreeoneofrhbandagreedyadiccore}(t, Q, K, v, S)\).
\end{proof}
\begin{theorem}
\label{ll:paper:periodic-maskparity-zero-agree-of-rhbandagree-dyadic-core}
For every natural number \(T\) and natural number \(P\), if \(0 < P\), then if \(\operatorname{centerperiodic}(T, P)\), then for every natural number \(m\) and natural number \(t\) and natural number \(K\) and natural number \(v\) and natural number \(S\) and natural number \(s\), if \(T \le t\), then if \(1 \le v\), then if \(t + S + m \cdot P + (t + m \cdot P) + 1 \le K\), then if \(\operatorname{rhbandagree}(t + v \cdot \operatorname{twopow}(K), m \cdot P, v \cdot \operatorname{twopow}(K) - (t + m \cdot P), 2 \cdot (t + m \cdot P))\), then if \(s \le S\), then \(\operatorname{maskparity}(t + s + m \cdot P, 0) = \operatorname{maskparity}(t + s, 0)\).
\end{theorem}
\begin{proof}
Assume \(hP\), \(hper\), \(m\), \(t\), \(K\), \(v\), \(S\), \(s\).
The goal follows from \(\operatorname{periodicmaskparityzeroagreeofrhbandagreedyadiccore}(T, P, hP, hper, m, t, K, v, S, s)\).
\end{proof}
\begin{theorem}
\label{ll:paper:not-centerperiodic-of-rhbandagree-dyadic-core-at-zero-diag-doubles}
For every natural number \(T\) and natural number \(P\) and natural number \(m\) and natural number \(z\) and natural number \(p\) and natural number \(K\) and natural number \(v\), if \(0 < P\), then if \(0 < m\), then if \(T \le z + (p + 1)\), then if \(\operatorname{zerodiagonal}(z)\), then if \(\operatorname{doublesat}(m \cdot P, z)\), then if \(\operatorname{settledby}(z + 1, z + (p + 1))\), then if \(\operatorname{settledby}(z + 1 - m \cdot P, z + (p + 1))\), then if \(1 \le v\), then if \(z + (p + 1) + (p + 3 \cdot (m \cdot P)) + m \cdot P + (z + (p + 1) + m \cdot P) + 1 \le K\), then if \(\operatorname{rhbandagree}(z + (p + 1) + v \cdot \operatorname{twopow}(K), m \cdot P, v \cdot \operatorname{twopow}(K) - (z + (p + 1) + m \cdot P), 2 \cdot (z + (p + 1) + m \cdot P))\), then not \(\operatorname{centerperiodic}(T, P)\).
\end{theorem}
\begin{proof}
The goal follows from \(\operatorname{notcenterperiodicofrhbandagreedyadiccoreatzerodiagdoubles}(T, P, m, z, p, K, v)\).
\end{proof}
\begin{theorem}
\label{ll:paper:not-centerperiodic-of-rhbandagree-dyadic-core-twopow}
For every natural number \(T\) and natural number \(P\) and natural number \(m\) and natural number \(t\) and natural number \(k\) and natural number \(K\) and natural number \(v\), if \(0 < P\), then if \(0 < m\), then if \(T \le t\), then if \(t + m \cdot P = \operatorname{twopow}(k)\), then if \(1 \le v\), then if \(t + \operatorname{twopow}(k) + m \cdot P + (t + m \cdot P) + 1 \le K\), then if \(\operatorname{rhbandagree}(t + v \cdot \operatorname{twopow}(K), m \cdot P, v \cdot \operatorname{twopow}(K) - (t + m \cdot P), 2 \cdot (t + m \cdot P))\), then not \(\operatorname{centerperiodic}(T, P)\).
\end{theorem}
\begin{proof}
The goal follows from \(\operatorname{notcenterperiodicofrhbandagreedyadiccoretwopow}(T, P, m, t, k, K, v)\).
\end{proof}
\begin{theorem}
\label{ll:paper:not-centerperiodic-of-rhbandagree-dyadic-core-at-zero-diag-pair}
For every natural number \(T\) and natural number \(P\) and natural number \(m\) and natural number \(z\) and natural number \(k\) and natural number \(K\) and natural number \(v\), if \(0 < P\), then if \(\operatorname{zerodiagonal}(z)\), then if \(2 \le z + m \cdot P\), then if \(\operatorname{zerodiagonal}(z + m \cdot P - 2)\), then if \(T \le z + \operatorname{twopow}(k)\), then if \(\operatorname{settledby}(z, z + \operatorname{twopow}(k))\), then if \(\operatorname{settledby}(z + m \cdot P, z + m \cdot P + \operatorname{twopow}(k))\), then if \(1 \le v\), then if \(z + \operatorname{twopow}(k) + \operatorname{twopow}(k) + m \cdot P + (z + \operatorname{twopow}(k) + m \cdot P) + 1 \le K\), then if \(\operatorname{rhbandagree}(z + \operatorname{twopow}(k) + v \cdot \operatorname{twopow}(K), m \cdot P, v \cdot \operatorname{twopow}(K) - (z + \operatorname{twopow}(k) + m \cdot P), 2 \cdot (z + \operatorname{twopow}(k) + m \cdot P))\), then not \(\operatorname{centerperiodic}(T, P)\).
\end{theorem}
\begin{proof}
The goal follows from \(\operatorname{notcenterperiodicofrhbandagreedyadiccoreatzerodiagpair}(T, P, m, z, k, K, v)\).
\end{proof}
\begin{theorem}
\label{ll:paper:rhbandagree-dyadic-left-half-of-lhbandagree}
For every natural number \(t\) and natural number \(Q\) and natural number \(K\) and natural number \(v\), if \(1 \le v\), then if \(2 \cdot (t + Q) \le K\), then if \(\operatorname{lhbandagree}(t, Q, 0, t + Q)\), then \(\operatorname{rhbandagree}(t + v \cdot \operatorname{twopow}(K), Q, v \cdot \operatorname{twopow}(K) - (t + Q), t + Q)\).
\end{theorem}
\begin{proof}
The goal follows from \(\operatorname{rhbandagreedyadiclefthalfoflhbandagree}(t, Q, K, v)\).
\end{proof}
\begin{theorem}
\label{ll:paper:rhbandagree-dyadic-right-half-of-rhbandagree}
For every natural number \(t\) and natural number \(Q\) and natural number \(K\) and natural number \(v\), if \(t + Q \le K\), then if \(\operatorname{rhbandagree}(t, Q, 0, t + Q)\), then \(\operatorname{rhbandagree}(t + v \cdot \operatorname{twopow}(K), Q, v \cdot \operatorname{twopow}(K), t + Q)\).
\end{theorem}
\begin{proof}
The goal follows from \(\operatorname{rhbandagreedyadicrighthalfofrhbandagree}(t, Q, K, v)\).
\end{proof}
\begin{theorem}
\label{ll:paper:rhbandagree-dyadic-core-of-lhbandagree-and-rhbandagree}
For every natural number \(t\) and natural number \(Q\) and natural number \(K\) and natural number \(v\), if \(1 \le v\), then if \(2 \cdot (t + Q) \le K\), then if \(\operatorname{lhbandagree}(t, Q, 0, t + Q)\), then if \(\operatorname{rhbandagree}(t, Q, 0, t + Q)\), then \(\operatorname{rhbandagree}(t + v \cdot \operatorname{twopow}(K), Q, v \cdot \operatorname{twopow}(K) - (t + Q), 2 \cdot (t + Q))\).
\end{theorem}
\begin{proof}
The goal follows from \(\operatorname{rhbandagreedyadiccoreoflhbandagreeandrhbandagree}(t, Q, K, v)\).
\end{proof}
\begin{theorem}
\label{ll:paper:lhcolagree-one-of-full-row-bandagree}
For every natural number \(t\) and natural number \(Q\) and natural number \(S\), if \(\operatorname{lhbandagree}(t, Q, 0, t + Q)\), then if \(\operatorname{rhbandagree}(t, Q, 0, t + Q)\), then \(\operatorname{lhcolagree}(t, Q, 1, S)\).
\end{theorem}
\begin{proof}
The goal follows from \(\operatorname{lhcolagreeoneoffullrowbandagree}(t, Q, S)\).
\end{proof}
\begin{theorem}
\label{ll:paper:lhcenter-zero-eq-rhcenter-zero}
For every natural number \(t\), \(\operatorname{lhcenter}(t, 0) = \operatorname{rhcenter}(t, 0)\).
\end{theorem}
\begin{proof}
The goal follows from \(\operatorname{lhcenterzeroeqrhcenterzero}(t)\).
\end{proof}
\begin{theorem}
\label{ll:paper:lhbandagree-zero-of-lhbandagree-one-and-rhbandagree-zero}
For every natural number \(t\) and natural number \(Q\) and natural number \(B\), if \(\operatorname{lhbandagree}(t, Q, 1, B)\), then if \(\operatorname{rhbandagree}(t, Q, 0, B)\), then \(\operatorname{lhbandagree}(t, Q, 0, B)\).
\end{theorem}
\begin{proof}
The goal follows from \(\operatorname{lhbandagreezerooflhbandagreeoneandrhbandagreezero}(t, Q, B)\).
\end{proof}
\begin{theorem}
\label{ll:paper:lhcolagree-one-of-lhbandagree-one-and-rhbandagree-zero}
For every natural number \(t\) and natural number \(Q\) and natural number \(S\), if \(\operatorname{lhbandagree}(t, Q, 1, t + Q)\), then if \(\operatorname{rhbandagree}(t, Q, 0, t + Q)\), then \(\operatorname{lhcolagree}(t, Q, 1, S)\).
\end{theorem}
\begin{proof}
The goal follows from \(\operatorname{lhcolagreeoneoflhbandagreeoneandrhbandagreezero}(t, Q, S)\).
\end{proof}
\begin{theorem}
\label{ll:paper:not-centerperiodic-of-lhbandagree-one-and-rhbandagree-zero-at-zero-diag-doubles}
For every natural number \(T\) and natural number \(P\) and natural number \(m\) and natural number \(z\) and natural number \(p\), if \(0 < P\), then if \(0 < m\), then if \(T \le z + (p + 1)\), then if \(\operatorname{zerodiagonal}(z)\), then if \(\operatorname{doublesat}(m \cdot P, z)\), then if \(\operatorname{settledby}(z + 1, z + (p + 1))\), then if \(\operatorname{settledby}(z + 1 - m \cdot P, z + (p + 1))\), then if \(\operatorname{lhbandagree}(z + (p + 1), m \cdot P, 1, z + (p + 1) + m \cdot P)\), then if \(\operatorname{rhbandagree}(z + (p + 1), m \cdot P, 0, z + (p + 1) + m \cdot P)\), then not \(\operatorname{centerperiodic}(T, P)\).
\end{theorem}
\begin{proof}
The goal follows from \(\operatorname{notcenterperiodicoflhbandagreeoneandrhbandagreezeroatzerodiagdoubles}(T, P, m, z, p)\).
\end{proof}
\begin{theorem}
\label{ll:paper:not-centerperiodic-of-lhbandagree-one-and-rhbandagree-zero-twopow}
For every natural number \(T\) and natural number \(P\) and natural number \(m\) and natural number \(t\) and natural number \(k\), if \(0 < P\), then if \(0 < m\), then if \(T \le t\), then if \(t + m \cdot P = \operatorname{twopow}(k)\), then if \(\operatorname{lhbandagree}(t, m \cdot P, 1, t + m \cdot P)\), then if \(\operatorname{rhbandagree}(t, m \cdot P, 0, t + m \cdot P)\), then not \(\operatorname{centerperiodic}(T, P)\).
\end{theorem}
\begin{proof}
The goal follows from \(\operatorname{notcenterperiodicoflhbandagreeoneandrhbandagreezerotwopow}(T, P, m, t, k)\).
\end{proof}
\begin{theorem}
\label{ll:paper:rhbandagree-dyadic-left-cone-of-lhbandagree}
For every natural number \(t\) and natural number \(Q\) and natural number \(K\) and natural number \(v\) and natural number \(S\), if \(1 \le v\), then if \(t + S + Q + 1 \le K\), then if \(\operatorname{lhbandagree}(t, Q, 1, S)\), then \(\operatorname{rhbandagree}(t + v \cdot \operatorname{twopow}(K), Q, v \cdot \operatorname{twopow}(K) - 1 - S, S)\).
\end{theorem}
\begin{proof}
The goal follows from \(\operatorname{rhbandagreedyadicleftconeoflhbandagree}(t, Q, K, v, S)\).
\end{proof}
\begin{theorem}
\label{ll:paper:rhbandagree-dyadic-right-cone-of-rhbandagree}
For every natural number \(t\) and natural number \(Q\) and natural number \(K\) and natural number \(v\) and natural number \(B\), if \(t + Q \le K\), then if \(\operatorname{rhbandagree}(t, Q, 0, B)\), then \(\operatorname{rhbandagree}(t + v \cdot \operatorname{twopow}(K), Q, v \cdot \operatorname{twopow}(K), B)\).
\end{theorem}
\begin{proof}
The goal follows from \(\operatorname{rhbandagreedyadicrightconeofrhbandagree}(t, Q, K, v, B)\).
\end{proof}
\begin{theorem}
\label{ll:paper:rhbandagree-dyadic-pred-col-cone-of-lhbandagree-and-rhbandagree}
For every natural number \(t\) and natural number \(Q\) and natural number \(K\) and natural number \(v\) and natural number \(S\), if \(1 \le v\), then if \(t + S + Q + 1 \le K\), then if \(\operatorname{lhbandagree}(t, Q, 1, S)\), then if \(\operatorname{rhbandagree}(t, Q, 0, S)\), then \(\operatorname{rhbandagree}(t + v \cdot \operatorname{twopow}(K), Q, v \cdot \operatorname{twopow}(K) - 1 - S, 2 \cdot S)\).
\end{theorem}
\begin{proof}
The goal follows from \(\operatorname{rhbandagreedyadicpredcolconeoflhbandagreeandrhbandagree}(t, Q, K, v, S)\).
\end{proof}
\begin{theorem}
\label{ll:paper:lhcolagree-one-of-lhbandagree-one-and-rhbandagree-zero-cone}
For every natural number \(t\) and natural number \(Q\) and natural number \(S\), if \(\operatorname{lhbandagree}(t, Q, 1, S)\), then if \(\operatorname{rhbandagree}(t, Q, 0, S)\), then \(\operatorname{lhcolagree}(t, Q, 1, S)\).
\end{theorem}
\begin{proof}
The goal follows from \(\operatorname{lhcolagreeoneoflhbandagreeoneandrhbandagreezerocone}(t, Q, S)\).
\end{proof}
\begin{theorem}
\label{ll:paper:periodic-maskparity-zero-agree-of-lhbandagree-one-and-rhbandagree-zero-cone}
For every natural number \(T\) and natural number \(P\), if \(0 < P\), then if \(\operatorname{centerperiodic}(T, P)\), then for every natural number \(m\) and natural number \(t\) and natural number \(S\) and natural number \(s\), if \(T \le t\), then if \(\operatorname{lhbandagree}(t, m \cdot P, 1, S)\), then if \(\operatorname{rhbandagree}(t, m \cdot P, 0, S)\), then if \(s \le S\), then \(\operatorname{maskparity}(t + s + m \cdot P, 0) = \operatorname{maskparity}(t + s, 0)\).
\end{theorem}
\begin{proof}
Assume \(hP\), \(hper\), \(m\), \(t\), \(S\), \(s\).
The goal follows from \(\operatorname{periodicmaskparityzeroagreeoflhbandagreeoneandrhbandagreezerocone}(T, P, hP, hper, m, t, S, s)\).
\end{proof}
\begin{theorem}
\label{ll:paper:not-centerperiodic-of-lhbandagree-one-and-rhbandagree-zero-cone-at-zero-diag-doubles}
For every natural number \(T\) and natural number \(P\) and natural number \(m\) and natural number \(z\) and natural number \(p\) and natural number \(S\), if \(0 < P\), then if \(0 < m\), then if \(T \le z + (p + 1)\), then if \(\operatorname{zerodiagonal}(z)\), then if \(\operatorname{doublesat}(m \cdot P, z)\), then if \(\operatorname{settledby}(z + 1, z + (p + 1))\), then if \(\operatorname{settledby}(z + 1 - m \cdot P, z + (p + 1))\), then if \(p + 3 \cdot (m \cdot P) \le S + 1\), then if \(\operatorname{lhbandagree}(z + (p + 1), m \cdot P, 1, S)\), then if \(\operatorname{rhbandagree}(z + (p + 1), m \cdot P, 0, S)\), then not \(\operatorname{centerperiodic}(T, P)\).
\end{theorem}
\begin{proof}
The goal follows from \(\operatorname{notcenterperiodicoflhbandagreeoneandrhbandagreezeroconeatzerodiagdoubles}(T, P, m, z, p, S)\).
\end{proof}
\begin{theorem}
\label{ll:paper:not-centerperiodic-of-lhbandagree-one-and-rhbandagree-zero-cone-twopow}
For every natural number \(T\) and natural number \(P\) and natural number \(m\) and natural number \(t\) and natural number \(k\) and natural number \(S\), if \(0 < P\), then if \(0 < m\), then if \(T \le t\), then if \(t + m \cdot P = \operatorname{twopow}(k)\), then if \(\operatorname{twopow}(k) \le S + 1\), then if \(\operatorname{lhbandagree}(t, m \cdot P, 1, S)\), then if \(\operatorname{rhbandagree}(t, m \cdot P, 0, S)\), then not \(\operatorname{centerperiodic}(T, P)\).
\end{theorem}
\begin{proof}
The goal follows from \(\operatorname{notcenterperiodicoflhbandagreeoneandrhbandagreezeroconetwopow}(T, P, m, t, k, S)\).
\end{proof}
\begin{theorem}
\label{ll:paper:not-centerperiodic-of-lhbandagree-one-and-rhbandagree-zero-cone-at-zero-diag-pair}
For every natural number \(T\) and natural number \(P\) and natural number \(m\) and natural number \(z\) and natural number \(k\) and natural number \(S\), if \(0 < P\), then if \(\operatorname{zerodiagonal}(z)\), then if \(2 \le z + m \cdot P\), then if \(\operatorname{zerodiagonal}(z + m \cdot P - 2)\), then if \(T \le z + \operatorname{twopow}(k)\), then if \(\operatorname{settledby}(z, z + \operatorname{twopow}(k))\), then if \(\operatorname{settledby}(z + m \cdot P, z + m \cdot P + \operatorname{twopow}(k))\), then if \(\operatorname{twopow}(k) \le S + 1\), then if \(\operatorname{lhbandagree}(z + \operatorname{twopow}(k), m \cdot P, 1, S)\), then if \(\operatorname{rhbandagree}(z + \operatorname{twopow}(k), m \cdot P, 0, S)\), then not \(\operatorname{centerperiodic}(T, P)\).
\end{theorem}
\begin{proof}
The goal follows from \(\operatorname{notcenterperiodicoflhbandagreeoneandrhbandagreezeroconeatzerodiagpair}(T, P, m, z, k, S)\).
\end{proof}
\begin{theorem}
\label{ll:paper:rhcenter-dyadic-diag-shift-left-eq}
For every natural number \(t\) and natural number \(K\) and natural number \(v\) and natural number \(w\) and natural number \(d\), if \(1 \le v\), then if \(t + d \le K\), then \(\operatorname{rhcenter}(t + v \cdot \operatorname{twopow}(K) + w \cdot \operatorname{twopow}(K), v \cdot \operatorname{twopow}(K) + w \cdot \operatorname{twopow}(K) - d) = \operatorname{rhcenter}(t + v \cdot \operatorname{twopow}(K), v \cdot \operatorname{twopow}(K) - d)\).
\end{theorem}
\begin{proof}
The goal follows from \(\operatorname{rhcenterdyadicdiagshiftlefteq}(t, K, v, w, d)\).
\end{proof}
\begin{theorem}
\label{ll:paper:rhcenter-dyadic-diag-shift-right-eq}
For every natural number \(t\) and natural number \(K\) and natural number \(v\) and natural number \(w\) and natural number \(q\), if \(t \le K\), then \(\operatorname{rhcenter}(t + v \cdot \operatorname{twopow}(K) + w \cdot \operatorname{twopow}(K), v \cdot \operatorname{twopow}(K) + w \cdot \operatorname{twopow}(K) + q) = \operatorname{rhcenter}(t + v \cdot \operatorname{twopow}(K), v \cdot \operatorname{twopow}(K) + q)\).
\end{theorem}
\begin{proof}
The goal follows from \(\operatorname{rhcenterdyadicdiagshiftrighteq}(t, K, v, w, q)\).
\end{proof}
\begin{theorem}
\label{ll:paper:rhcenter-dyadic-diag-shift-band-eq}
For every natural number \(t\) and natural number \(K\) and natural number \(v\) and natural number \(w\) and natural number \(W\) and natural number \(i\), if \(1 \le v\), then if \(t + W \le K\), then \(\operatorname{rhcenter}(t + v \cdot \operatorname{twopow}(K) + w \cdot \operatorname{twopow}(K), v \cdot \operatorname{twopow}(K) + w \cdot \operatorname{twopow}(K) - W + i) = \operatorname{rhcenter}(t + v \cdot \operatorname{twopow}(K), v \cdot \operatorname{twopow}(K) - W + i)\).
\end{theorem}
\begin{proof}
The goal follows from \(\operatorname{rhcenterdyadicdiagshiftbandeq}(t, K, v, w, W, i)\).
\end{proof}
\begin{theorem}
\label{ll:paper:rhbandagree-dyadic-of-spatial-period-shift-agree}
For every natural number \(t\) and natural number \(K\) and natural number \(v\) and natural number \(w\) and natural number \(W\) and natural number \(B\), if \(1 \le v\), then if \(t + W \le K\), then if for every natural number \(i\), if \(i \le B\), then \(\operatorname{rhcenter}(t + v \cdot \operatorname{twopow}(K) + w \cdot \operatorname{twopow}(K), v \cdot \operatorname{twopow}(K) - W + i) = \operatorname{rhcenter}(t + v \cdot \operatorname{twopow}(K) + w \cdot \operatorname{twopow}(K), v \cdot \operatorname{twopow}(K) + w \cdot \operatorname{twopow}(K) - W + i)\), then \(\operatorname{rhbandagree}(t + v \cdot \operatorname{twopow}(K), w \cdot \operatorname{twopow}(K), v \cdot \operatorname{twopow}(K) - W, B)\).
\end{theorem}
\begin{proof}
The goal follows from \(\operatorname{rhbandagreedyadicofspatialperiodshiftagree}(t, K, v, w, W, B)\).
\end{proof}
\begin{theorem}
\label{ll:paper:rhcenter-two-scale-row-left-eq-lhcenter}
For every natural number \(t\) and natural number \(K\) and natural number \(L\) and natural number \(u\) and natural number \(w\) and natural number \(d\), if \(1 \le u\), then if \(K \le L\), then if \(t + d \le K\), then \(\operatorname{rhcenter}(t + u \cdot \operatorname{twopow}(L) + w \cdot \operatorname{twopow}(K), u \cdot \operatorname{twopow}(L) + w \cdot \operatorname{twopow}(K) - d) = \operatorname{lhcenter}(t, d)\).
\end{theorem}
\begin{proof}
The goal follows from \(\operatorname{rhcentertwoscalerowlefteqlhcenter}(t, K, L, u, w, d)\).
\end{proof}
\begin{theorem}
\label{ll:paper:rhcenter-two-scale-row-right-eq-rhcenter}
For every natural number \(t\) and natural number \(K\) and natural number \(L\) and natural number \(u\) and natural number \(w\) and natural number \(q\), if \(K \le L\), then if \(t \le K\), then \(\operatorname{rhcenter}(t + u \cdot \operatorname{twopow}(L) + w \cdot \operatorname{twopow}(K), u \cdot \operatorname{twopow}(L) + w \cdot \operatorname{twopow}(K) + q) = \operatorname{rhcenter}(t, q)\).
\end{theorem}
\begin{proof}
The goal follows from \(\operatorname{rhcentertwoscalerowrighteqrhcenter}(t, K, L, u, w, q)\).
\end{proof}
\begin{theorem}
\label{ll:paper:lhcolagree-one-of-two-scale-single-row-spatial-agree}
For every natural number \(t\) and natural number \(K\) and natural number \(L\) and natural number \(u\) and natural number \(w\) and natural number \(S\), if \(1 \le u\), then if \(t + S + 1 \le K\), then if \(t + S + w \cdot \operatorname{twopow}(K) + 1 \le L\), then if \(K \le L\), then if for every natural number \(i\), if \(i \le 2 \cdot S\), then \(\operatorname{rhcenter}(t + u \cdot \operatorname{twopow}(L) + w \cdot \operatorname{twopow}(K), u \cdot \operatorname{twopow}(L) - 1 - S + i) = \operatorname{rhcenter}(t + u \cdot \operatorname{twopow}(L) + w \cdot \operatorname{twopow}(K), u \cdot \operatorname{twopow}(L) + w \cdot \operatorname{twopow}(K) - 1 - S + i)\), then \(\operatorname{lhcolagree}(t, w \cdot \operatorname{twopow}(K), 1, S)\).
\end{theorem}
\begin{proof}
The goal follows from \(\operatorname{lhcolagreeoneoftwoscalesinglerowspatialagree}(t, K, L, u, w, S)\).
\end{proof}
\begin{theorem}
\label{ll:paper:lhcolagree-one-of-two-scale-single-row-mul-period-agree}
For every natural number \(t\) and natural number \(P\) and natural number \(K\) and natural number \(L\) and natural number \(u\) and natural number \(c\) and natural number \(S\), if \(1 \le u\), then if \(t + S + 1 \le K\), then if \(t + S + c \cdot \operatorname{twopow}(K) \cdot P + 1 \le L\), then if \(K \le L\), then if for every natural number \(i\), if \(i \le 2 \cdot S\), then \(\operatorname{rhcenter}(t + u \cdot \operatorname{twopow}(L) + c \cdot \operatorname{twopow}(K) \cdot P, u \cdot \operatorname{twopow}(L) - 1 - S + i) = \operatorname{rhcenter}(t + u \cdot \operatorname{twopow}(L) + c \cdot \operatorname{twopow}(K) \cdot P, u \cdot \operatorname{twopow}(L) + c \cdot \operatorname{twopow}(K) \cdot P - 1 - S + i)\), then \(\operatorname{lhcolagree}(t, c \cdot \operatorname{twopow}(K) \cdot P, 1, S)\).
\end{theorem}
\begin{proof}
The goal follows from \(\operatorname{lhcolagreeoneoftwoscalesinglerowmulperiodagree}(t, P, K, L, u, c, S)\).
\end{proof}
\begin{theorem}
\label{ll:paper:periodic-maskparity-zero-agree-of-two-scale-single-row-mul-period-agree}
For every natural number \(T\) and natural number \(P\), if \(0 < P\), then if \(\operatorname{centerperiodic}(T, P)\), then for every natural number \(t\) and natural number \(K\) and natural number \(L\) and natural number \(u\) and natural number \(c\) and natural number \(S\) and natural number \(s\), if \(T \le t\), then if \(1 \le u\), then if \(t + S + 1 \le K\), then if \(t + S + c \cdot \operatorname{twopow}(K) \cdot P + 1 \le L\), then if \(K \le L\), then if for every natural number \(i\), if \(i \le 2 \cdot S\), then \(\operatorname{rhcenter}(t + u \cdot \operatorname{twopow}(L) + c \cdot \operatorname{twopow}(K) \cdot P, u \cdot \operatorname{twopow}(L) - 1 - S + i) = \operatorname{rhcenter}(t + u \cdot \operatorname{twopow}(L) + c \cdot \operatorname{twopow}(K) \cdot P, u \cdot \operatorname{twopow}(L) + c \cdot \operatorname{twopow}(K) \cdot P - 1 - S + i)\), then if \(s \le S\), then \(\operatorname{maskparity}(t + s + c \cdot \operatorname{twopow}(K) \cdot P, 0) = \operatorname{maskparity}(t + s, 0)\).
\end{theorem}
\begin{proof}
Assume \(hP\), \(hper\), \(t\), \(K\), \(L\), \(u\), \(c\), \(S\), \(s\).
The goal follows from \(\operatorname{periodicmaskparityzeroagreeoftwoscalesinglerowmulperiodagree}(T, P, hP, hper, t, K, L, u, c, S, s)\).
\end{proof}
\begin{theorem}
\label{ll:paper:not-centerperiodic-of-two-scale-single-row-at-zero-diag-doubles}
For every natural number \(T\) and natural number \(P\) and natural number \(c\) and natural number \(z\) and natural number \(p\) and natural number \(K\) and natural number \(L\) and natural number \(u\) and natural number \(S\), if \(0 < P\), then if \(0 < c\), then if \(T \le z + (p + 1)\), then if \(\operatorname{zerodiagonal}(z)\), then if \(\operatorname{doublesat}(c \cdot \operatorname{twopow}(K) \cdot P, z)\), then if \(\operatorname{settledby}(z + 1, z + (p + 1))\), then if \(\operatorname{settledby}(z + 1 - c \cdot \operatorname{twopow}(K) \cdot P, z + (p + 1))\), then if \(p + 3 \cdot (c \cdot \operatorname{twopow}(K) \cdot P) \le S + 1\), then if \(1 \le u\), then if \(z + (p + 1) + S + 1 \le K\), then if \(z + (p + 1) + S + c \cdot \operatorname{twopow}(K) \cdot P + 1 \le L\), then if \(K \le L\), then if for every natural number \(i\), if \(i \le 2 \cdot S\), then \(\operatorname{rhcenter}(z + (p + 1) + u \cdot \operatorname{twopow}(L) + c \cdot \operatorname{twopow}(K) \cdot P, u \cdot \operatorname{twopow}(L) - 1 - S + i) = \operatorname{rhcenter}(z + (p + 1) + u \cdot \operatorname{twopow}(L) + c \cdot \operatorname{twopow}(K) \cdot P, u \cdot \operatorname{twopow}(L) + c \cdot \operatorname{twopow}(K) \cdot P - 1 - S + i)\), then not \(\operatorname{centerperiodic}(T, P)\).
\end{theorem}
\begin{proof}
The goal follows from \(\operatorname{notcenterperiodicoftwoscalesinglerowatzerodiagdoubles}(T, P, c, z, p, K, L, u, S)\).
\end{proof}
\begin{theorem}
\label{ll:paper:not-centerperiodic-of-two-scale-single-row-twopow}
For every natural number \(T\) and natural number \(P\) and natural number \(c\) and natural number \(t\) and natural number \(k\) and natural number \(K\) and natural number \(L\) and natural number \(u\) and natural number \(S\), if \(0 < P\), then if \(0 < c\), then if \(T \le t\), then if \(t + c \cdot \operatorname{twopow}(K) \cdot P = \operatorname{twopow}(k)\), then if \(\operatorname{twopow}(k) \le S + 1\), then if \(1 \le u\), then if \(t + S + 1 \le K\), then if \(t + S + c \cdot \operatorname{twopow}(K) \cdot P + 1 \le L\), then if \(K \le L\), then if for every natural number \(i\), if \(i \le 2 \cdot S\), then \(\operatorname{rhcenter}(t + u \cdot \operatorname{twopow}(L) + c \cdot \operatorname{twopow}(K) \cdot P, u \cdot \operatorname{twopow}(L) - 1 - S + i) = \operatorname{rhcenter}(t + u \cdot \operatorname{twopow}(L) + c \cdot \operatorname{twopow}(K) \cdot P, u \cdot \operatorname{twopow}(L) + c \cdot \operatorname{twopow}(K) \cdot P - 1 - S + i)\), then not \(\operatorname{centerperiodic}(T, P)\).
\end{theorem}
\begin{proof}
The goal follows from \(\operatorname{notcenterperiodicoftwoscalesinglerowtwopow}(T, P, c, t, k, K, L, u, S)\).
\end{proof}
\begin{theorem}
\label{ll:paper:not-centerperiodic-of-two-scale-single-row-at-zero-diag-pair}
For every natural number \(T\) and natural number \(P\) and natural number \(c\) and natural number \(z\) and natural number \(k\) and natural number \(K\) and natural number \(L\) and natural number \(u\) and natural number \(S\), if \(0 < P\), then if \(\operatorname{zerodiagonal}(z)\), then if \(2 \le z + c \cdot \operatorname{twopow}(K) \cdot P\), then if \(\operatorname{zerodiagonal}(z + c \cdot \operatorname{twopow}(K) \cdot P - 2)\), then if \(T \le z + \operatorname{twopow}(k)\), then if \(\operatorname{settledby}(z, z + \operatorname{twopow}(k))\), then if \(\operatorname{settledby}(z + c \cdot \operatorname{twopow}(K) \cdot P, z + c \cdot \operatorname{twopow}(K) \cdot P + \operatorname{twopow}(k))\), then if \(\operatorname{twopow}(k) \le S + 1\), then if \(1 \le u\), then if \(z + \operatorname{twopow}(k) + S + 1 \le K\), then if \(z + \operatorname{twopow}(k) + S + c \cdot \operatorname{twopow}(K) \cdot P + 1 \le L\), then if \(K \le L\), then if for every natural number \(i\), if \(i \le 2 \cdot S\), then \(\operatorname{rhcenter}(z + \operatorname{twopow}(k) + u \cdot \operatorname{twopow}(L) + c \cdot \operatorname{twopow}(K) \cdot P, u \cdot \operatorname{twopow}(L) - 1 - S + i) = \operatorname{rhcenter}(z + \operatorname{twopow}(k) + u \cdot \operatorname{twopow}(L) + c \cdot \operatorname{twopow}(K) \cdot P, u \cdot \operatorname{twopow}(L) + c \cdot \operatorname{twopow}(K) \cdot P - 1 - S + i)\), then not \(\operatorname{centerperiodic}(T, P)\).
\end{theorem}
\begin{proof}
The goal follows from \(\operatorname{notcenterperiodicoftwoscalesinglerowatzerodiagpair}(T, P, c, z, k, K, L, u, S)\).
\end{proof}
\begin{theorem}
\label{ll:paper:rhcenter-dyadic-exact-row-left-eq-false}
For every natural number \(L\) and natural number \(u\) and natural number \(d\), if \(1 \le u\), then if \(d \le L\), then if \(0 < d\), then \(\operatorname{rhcenter}(u \cdot \operatorname{twopow}(L), u \cdot \operatorname{twopow}(L) - d) = \operatorname{false}\).
\end{theorem}
\begin{proof}
The goal follows from \(\operatorname{rhcenterdyadicexactrowlefteqfalse}(L, u, d)\).
\end{proof}
\begin{theorem}
\label{ll:paper:rhcenter-dyadic-exact-row-self-eq-true}
For every natural number \(L\) and natural number \(u\), \(\operatorname{rhcenter}(u \cdot \operatorname{twopow}(L), u \cdot \operatorname{twopow}(L)) = \operatorname{true}\).
\end{theorem}
\begin{proof}
The goal follows from \(\operatorname{rhcenterdyadicexactrowselfeqtrue}(L, u)\).
\end{proof}
\begin{theorem}
\label{ll:paper:rhcenter-dyadic-exact-row-right-eq-false}
For every natural number \(L\) and natural number \(u\) and natural number \(r\), if \(0 < r\), then \(\operatorname{rhcenter}(u \cdot \operatorname{twopow}(L), u \cdot \operatorname{twopow}(L) + r) = \operatorname{false}\).
\end{theorem}
\begin{proof}
The goal follows from \(\operatorname{rhcenterdyadicexactrowrighteqfalse}(L, u, r)\).
\end{proof}
\begin{theorem}
\label{ll:paper:rhcenter-dyadic-exact-row-translate-agree}
For every natural number \(L\) and natural number \(u\) and natural number \(v\) and natural number \(W\) and natural number \(i\), if \(1 \le u\), then if \(1 \le v\), then if \(W \le L\), then \(\operatorname{rhcenter}(v \cdot \operatorname{twopow}(L), v \cdot \operatorname{twopow}(L) - W + i) = \operatorname{rhcenter}(u \cdot \operatorname{twopow}(L), u \cdot \operatorname{twopow}(L) - W + i)\).
\end{theorem}
\begin{proof}
The goal follows from \(\operatorname{rhcenterdyadicexactrowtranslateagree}(L, u, v, W, i)\).
\end{proof}
\begin{theorem}
\label{ll:paper:rhcenter-dyadic-exact-row-evolve-translate-agree}
For every natural number \(L\) and natural number \(u\) and natural number \(v\) and natural number \(W\) and natural number \(s\) and natural number \(i\), if \(1 \le u\), then if \(1 \le v\), then if \(s + W \le L\), then \(\operatorname{rhcenter}(v \cdot \operatorname{twopow}(L) + s, v \cdot \operatorname{twopow}(L) - W + i) = \operatorname{rhcenter}(u \cdot \operatorname{twopow}(L) + s, u \cdot \operatorname{twopow}(L) - W + i)\).
\end{theorem}
\begin{proof}
The goal follows from \(\operatorname{rhcenterdyadicexactrowevolvetranslateagree}(L, u, v, W, s, i)\).
\end{proof}
\begin{theorem}
\label{ll:paper:rdrv-dyadic-exact-row-left-eq-false}
For every natural number \(td\) and natural number \(L\) and natural number \(u\) and natural number \(d\), if \(1 \le u\), then if \(d \le L\), then if \(0 < d\), then \(\operatorname{rdrv}(td, u \cdot \operatorname{twopow}(L), 0, u \cdot \operatorname{twopow}(L) - d) = \operatorname{false}\).
\end{theorem}
\begin{proof}
The goal follows from \(\operatorname{rdrvdyadicexactrowlefteqfalse}(td, L, u, d)\).
\end{proof}
\begin{theorem}
\label{ll:paper:rdrv-dyadic-exact-row-self-eq-true}
For every natural number \(td\) and natural number \(L\) and natural number \(u\), if \(1 \le u\), then \(\operatorname{rdrv}(td, u \cdot \operatorname{twopow}(L), 0, u \cdot \operatorname{twopow}(L)) = \operatorname{true}\).
\end{theorem}
\begin{proof}
The goal follows from \(\operatorname{rdrvdyadicexactrowselfeqtrue}(td, L, u)\).
\end{proof}
\begin{theorem}
\label{ll:paper:rdrv-dyadic-exact-row-right-eq-false}
For every natural number \(td\) and natural number \(L\) and natural number \(u\) and natural number \(r\), if \(0 < r\), then \(\operatorname{rdrv}(td, u \cdot \operatorname{twopow}(L), 0, u \cdot \operatorname{twopow}(L) + r) = \operatorname{false}\).
\end{theorem}
\begin{proof}
The goal follows from \(\operatorname{rdrvdyadicexactrowrighteqfalse}(td, L, u, r)\).
\end{proof}
\begin{theorem}
\label{ll:paper:rdrv-dyadic-exact-row-shift-left-eq}
For every natural number \(td\) and natural number \(L\) and natural number \(u\) and natural number \(d\) and natural number \(s\), if \(1 \le u\), then if \(s + d \le L\), then \(\operatorname{rdrv}(td, u \cdot \operatorname{twopow}(L), s, u \cdot \operatorname{twopow}(L) - d) = \operatorname{lhcenter}(s, d)\).
\end{theorem}
\begin{proof}
The goal follows from \(\operatorname{rdrvdyadicexactrowshiftlefteq}(td, L, u, d, s)\).
\end{proof}
\begin{theorem}
\label{ll:paper:rdrv-dyadic-exact-row-shift-right-eq}
For every natural number \(td\) and natural number \(L\) and natural number \(u\) and natural number \(q\) and natural number \(s\), if \(1 \le u\), then if \(s \le L\), then \(\operatorname{rdrv}(td, u \cdot \operatorname{twopow}(L), s, u \cdot \operatorname{twopow}(L) + q) = \operatorname{rhcenter}(s, q)\).
\end{theorem}
\begin{proof}
The goal follows from \(\operatorname{rdrvdyadicexactrowshiftrighteq}(td, L, u, q, s)\).
\end{proof}
\begin{theorem}
\label{ll:paper:rdrv-dyadic-exact-row-time-shift-left-eq}
For every natural number \(td\) and natural number \(t\) and natural number \(L\) and natural number \(u\) and natural number \(d\) and natural number \(s\), if \(1 \le u\), then if \(t + s + d \le L\), then \(\operatorname{rdrv}(td, u \cdot \operatorname{twopow}(L), t + s, u \cdot \operatorname{twopow}(L) - d) = \operatorname{lhcenter}(t + s, d)\).
\end{theorem}
\begin{proof}
The goal follows from \(\operatorname{rdrvdyadicexactrowtimeshiftlefteq}(td, t, L, u, d, s)\).
\end{proof}
\begin{theorem}
\label{ll:paper:rdrv-dyadic-exact-row-time-shift-right-eq}
For every natural number \(td\) and natural number \(t\) and natural number \(L\) and natural number \(u\) and natural number \(q\) and natural number \(s\), if \(1 \le u\), then if \(t + s \le L\), then \(\operatorname{rdrv}(td, u \cdot \operatorname{twopow}(L), t + s, u \cdot \operatorname{twopow}(L) + q) = \operatorname{rhcenter}(t + s, q)\).
\end{theorem}
\begin{proof}
The goal follows from \(\operatorname{rdrvdyadicexactrowtimeshiftrighteq}(td, t, L, u, q, s)\).
\end{proof}
\begin{theorem}
\label{ll:paper:lhcolagree-of-rdrv-dyadic-exact-row-step-agree}
For every natural number \(td\) and natural number \(t\) and natural number \(Q\) and natural number \(L\) and natural number \(u\) and natural number \(d\) and natural number \(S\), if \(1 \le u\), then if \(t + S + Q + d \le L\), then if for every natural number \(s\), if \(s \le S\), then \(\operatorname{rdrv}(td, u \cdot \operatorname{twopow}(L), t + s + Q, u \cdot \operatorname{twopow}(L) - d) = \operatorname{rdrv}(td, u \cdot \operatorname{twopow}(L), t + s, u \cdot \operatorname{twopow}(L) - d)\), then \(\operatorname{lhcolagree}(t, Q, d, S)\).
\end{theorem}
\begin{proof}
The goal follows from \(\operatorname{lhcolagreeofrdrvdyadicexactrowstepagree}(td, t, Q, L, u, d, S)\).
\end{proof}
\begin{theorem}
\label{ll:paper:lhcolagree-one-of-rdrv-dyadic-exact-row-step-agree}
For every natural number \(td\) and natural number \(t\) and natural number \(Q\) and natural number \(L\) and natural number \(u\) and natural number \(S\), if \(1 \le u\), then if \(t + S + Q + 1 \le L\), then if for every natural number \(s\), if \(s \le S\), then \(\operatorname{rdrv}(td, u \cdot \operatorname{twopow}(L), t + s + Q, u \cdot \operatorname{twopow}(L) - 1) = \operatorname{rdrv}(td, u \cdot \operatorname{twopow}(L), t + s, u \cdot \operatorname{twopow}(L) - 1)\), then \(\operatorname{lhcolagree}(t, Q, 1, S)\).
\end{theorem}
\begin{proof}
The goal follows from \(\operatorname{lhcolagreeoneofrdrvdyadicexactrowstepagree}(td, t, Q, L, u, S)\).
\end{proof}
\begin{theorem}
\label{ll:paper:periodic-maskparity-zero-agree-of-rdrv-dyadic-exact-row-step-agree}
For every natural number \(T\) and natural number \(P\), if \(0 < P\), then if \(\operatorname{centerperiodic}(T, P)\), then for every natural number \(td\) and natural number \(m\) and natural number \(t\) and natural number \(L\) and natural number \(u\) and natural number \(S\) and natural number \(s\), if \(T \le t\), then if \(1 \le u\), then if \(t + S + m \cdot P + 1 \le L\), then if for every natural number \(j\), if \(j \le S\), then \(\operatorname{rdrv}(td, u \cdot \operatorname{twopow}(L), t + j + m \cdot P, u \cdot \operatorname{twopow}(L) - 1) = \operatorname{rdrv}(td, u \cdot \operatorname{twopow}(L), t + j, u \cdot \operatorname{twopow}(L) - 1)\), then if \(s \le S\), then \(\operatorname{maskparity}(t + s + m \cdot P, 0) = \operatorname{maskparity}(t + s, 0)\).
\end{theorem}
\begin{proof}
Assume \(hP\), \(hper\), \(td\), \(m\), \(t\), \(L\), \(u\), \(S\), \(s\).
The goal follows from \(\operatorname{periodicmaskparityzeroagreeofrdrvdyadicexactrowstepagree}(T, P, hP, hper, td, m, t, L, u, S, s)\).
\end{proof}
\begin{theorem}
\label{ll:paper:not-centerperiodic-of-rdrv-dyadic-exact-row-step-agree-at-zero-diag-doubles}
For every natural number \(T\) and natural number \(P\) and natural number \(td\) and natural number \(m\) and natural number \(z\) and natural number \(p\) and natural number \(L\) and natural number \(u\) and natural number \(S\), if \(0 < P\), then if \(0 < m\), then if \(T \le z + (p + 1)\), then if \(\operatorname{zerodiagonal}(z)\), then if \(\operatorname{doublesat}(m \cdot P, z)\), then if \(\operatorname{settledby}(z + 1, z + (p + 1))\), then if \(\operatorname{settledby}(z + 1 - m \cdot P, z + (p + 1))\), then if \(p + 3 \cdot (m \cdot P) \le S + 1\), then if \(1 \le u\), then if \(z + (p + 1) + S + m \cdot P + 1 \le L\), then if for every natural number \(s\), if \(s \le S\), then \(\operatorname{rdrv}(td, u \cdot \operatorname{twopow}(L), z + (p + 1) + s + m \cdot P, u \cdot \operatorname{twopow}(L) - 1) = \operatorname{rdrv}(td, u \cdot \operatorname{twopow}(L), z + (p + 1) + s, u \cdot \operatorname{twopow}(L) - 1)\), then not \(\operatorname{centerperiodic}(T, P)\).
\end{theorem}
\begin{proof}
The goal follows from \(\operatorname{notcenterperiodicofrdrvdyadicexactrowstepagreeatzerodiagdoubles}(T, P, td, m, z, p, L, u, S)\).
\end{proof}
\begin{theorem}
\label{ll:paper:not-centerperiodic-of-rdrv-dyadic-exact-row-step-agree-twopow}
For every natural number \(T\) and natural number \(P\) and natural number \(td\) and natural number \(m\) and natural number \(t\) and natural number \(k\) and natural number \(L\) and natural number \(u\) and natural number \(S\), if \(0 < P\), then if \(0 < m\), then if \(T \le t\), then if \(t + m \cdot P = \operatorname{twopow}(k)\), then if \(\operatorname{twopow}(k) \le S + 1\), then if \(1 \le u\), then if \(t + S + m \cdot P + 1 \le L\), then if for every natural number \(s\), if \(s \le S\), then \(\operatorname{rdrv}(td, u \cdot \operatorname{twopow}(L), t + s + m \cdot P, u \cdot \operatorname{twopow}(L) - 1) = \operatorname{rdrv}(td, u \cdot \operatorname{twopow}(L), t + s, u \cdot \operatorname{twopow}(L) - 1)\), then not \(\operatorname{centerperiodic}(T, P)\).
\end{theorem}
\begin{proof}
The goal follows from \(\operatorname{notcenterperiodicofrdrvdyadicexactrowstepagreetwopow}(T, P, td, m, t, k, L, u, S)\).
\end{proof}
\begin{theorem}
\label{ll:paper:not-centerperiodic-of-rdrv-dyadic-exact-row-step-agree-at-zero-diag-pair}
For every natural number \(T\) and natural number \(P\) and natural number \(td\) and natural number \(m\) and natural number \(z\) and natural number \(k\) and natural number \(L\) and natural number \(u\) and natural number \(S\), if \(0 < P\), then if \(\operatorname{zerodiagonal}(z)\), then if \(2 \le z + m \cdot P\), then if \(\operatorname{zerodiagonal}(z + m \cdot P - 2)\), then if \(T \le z + \operatorname{twopow}(k)\), then if \(\operatorname{settledby}(z, z + \operatorname{twopow}(k))\), then if \(\operatorname{settledby}(z + m \cdot P, z + m \cdot P + \operatorname{twopow}(k))\), then if \(\operatorname{twopow}(k) \le S + 1\), then if \(1 \le u\), then if \(z + \operatorname{twopow}(k) + S + m \cdot P + 1 \le L\), then if for every natural number \(s\), if \(s \le S\), then \(\operatorname{rdrv}(td, u \cdot \operatorname{twopow}(L), z + \operatorname{twopow}(k) + s + m \cdot P, u \cdot \operatorname{twopow}(L) - 1) = \operatorname{rdrv}(td, u \cdot \operatorname{twopow}(L), z + \operatorname{twopow}(k) + s, u \cdot \operatorname{twopow}(L) - 1)\), then not \(\operatorname{centerperiodic}(T, P)\).
\end{theorem}
\begin{proof}
The goal follows from \(\operatorname{notcenterperiodicofrdrvdyadicexactrowstepagreeatzerodiagpair}(T, P, td, m, z, k, L, u, S)\).
\end{proof}
\begin{theorem}
\label{ll:paper:rdrv-dyadic-shifted-init-left-eq-lhcenter}
For every natural number \(td\) and natural number \(t\) and natural number \(L\) and natural number \(u\) and natural number \(d\) and natural number \(s\), if \(1 \le u\), then if \(t + s + d \le L\), then \(\operatorname{rdrv}(td, t + u \cdot \operatorname{twopow}(L), s, u \cdot \operatorname{twopow}(L) - d) = \operatorname{lhcenter}(t + s, d)\).
\end{theorem}
\begin{proof}
The goal follows from \(\operatorname{rdrvdyadicshiftedinitlefteqlhcenter}(td, t, L, u, d, s)\).
\end{proof}
\begin{theorem}
\label{ll:paper:rdrv-dyadic-shifted-init-right-eq-rhcenter}
For every natural number \(td\) and natural number \(t\) and natural number \(L\) and natural number \(u\) and natural number \(q\) and natural number \(s\), if \(1 \le u\), then if \(t + s \le L\), then \(\operatorname{rdrv}(td, t + u \cdot \operatorname{twopow}(L), s, u \cdot \operatorname{twopow}(L) + q) = \operatorname{rhcenter}(t + s, q)\).
\end{theorem}
\begin{proof}
The goal follows from \(\operatorname{rdrvdyadicshiftedinitrighteqrhcenter}(td, t, L, u, q, s)\).
\end{proof}
\begin{theorem}
\label{ll:paper:lhcolagree-of-rdrv-dyadic-shifted-init-agree}
For every natural number \(td\) and natural number \(t\) and natural number \(Q\) and natural number \(L\) and natural number \(u\) and natural number \(d\) and natural number \(S\), if \(1 \le u\), then if \(t + S + Q + d \le L\), then if for every natural number \(s\), if \(s \le S\), then \(\operatorname{rdrv}(td, t + u \cdot \operatorname{twopow}(L) + Q, s, u \cdot \operatorname{twopow}(L) - d) = \operatorname{rdrv}(td, t + u \cdot \operatorname{twopow}(L), s, u \cdot \operatorname{twopow}(L) - d)\), then \(\operatorname{lhcolagree}(t, Q, d, S)\).
\end{theorem}
\begin{proof}
The goal follows from \(\operatorname{lhcolagreeofrdrvdyadicshiftedinitagree}(td, t, Q, L, u, d, S)\).
\end{proof}
\begin{theorem}
\label{ll:paper:lhcolagree-one-of-rdrv-dyadic-shifted-init-agree}
For every natural number \(td\) and natural number \(t\) and natural number \(Q\) and natural number \(L\) and natural number \(u\) and natural number \(S\), if \(1 \le u\), then if \(t + S + Q + 1 \le L\), then if for every natural number \(s\), if \(s \le S\), then \(\operatorname{rdrv}(td, t + u \cdot \operatorname{twopow}(L) + Q, s, u \cdot \operatorname{twopow}(L) - 1) = \operatorname{rdrv}(td, t + u \cdot \operatorname{twopow}(L), s, u \cdot \operatorname{twopow}(L) - 1)\), then \(\operatorname{lhcolagree}(t, Q, 1, S)\).
\end{theorem}
\begin{proof}
The goal follows from \(\operatorname{lhcolagreeoneofrdrvdyadicshiftedinitagree}(td, t, Q, L, u, S)\).
\end{proof}
\begin{theorem}
\label{ll:paper:not-centerperiodic-of-rdrv-dyadic-shifted-init-agree-at-zero-diag-doubles}
For every natural number \(T\) and natural number \(P\) and natural number \(td\) and natural number \(m\) and natural number \(z\) and natural number \(p\) and natural number \(L\) and natural number \(u\) and natural number \(S\), if \(0 < P\), then if \(0 < m\), then if \(T \le z + (p + 1)\), then if \(\operatorname{zerodiagonal}(z)\), then if \(\operatorname{doublesat}(m \cdot P, z)\), then if \(\operatorname{settledby}(z + 1, z + (p + 1))\), then if \(\operatorname{settledby}(z + 1 - m \cdot P, z + (p + 1))\), then if \(p + 3 \cdot (m \cdot P) \le S + 1\), then if \(1 \le u\), then if \(z + (p + 1) + S + m \cdot P + 1 \le L\), then if for every natural number \(s\), if \(s \le S\), then \(\operatorname{rdrv}(td, z + (p + 1) + u \cdot \operatorname{twopow}(L) + m \cdot P, s, u \cdot \operatorname{twopow}(L) - 1) = \operatorname{rdrv}(td, z + (p + 1) + u \cdot \operatorname{twopow}(L), s, u \cdot \operatorname{twopow}(L) - 1)\), then not \(\operatorname{centerperiodic}(T, P)\).
\end{theorem}
\begin{proof}
The goal follows from \(\operatorname{notcenterperiodicofrdrvdyadicshiftedinitagreeatzerodiagdoubles}(T, P, td, m, z, p, L, u, S)\).
\end{proof}
\begin{theorem}
\label{ll:paper:not-centerperiodic-of-rdrv-dyadic-shifted-init-agree-twopow}
For every natural number \(T\) and natural number \(P\) and natural number \(td\) and natural number \(m\) and natural number \(t\) and natural number \(k\) and natural number \(L\) and natural number \(u\) and natural number \(S\), if \(0 < P\), then if \(0 < m\), then if \(T \le t\), then if \(t + m \cdot P = \operatorname{twopow}(k)\), then if \(\operatorname{twopow}(k) \le S + 1\), then if \(1 \le u\), then if \(t + S + m \cdot P + 1 \le L\), then if for every natural number \(s\), if \(s \le S\), then \(\operatorname{rdrv}(td, t + u \cdot \operatorname{twopow}(L) + m \cdot P, s, u \cdot \operatorname{twopow}(L) - 1) = \operatorname{rdrv}(td, t + u \cdot \operatorname{twopow}(L), s, u \cdot \operatorname{twopow}(L) - 1)\), then not \(\operatorname{centerperiodic}(T, P)\).
\end{theorem}
\begin{proof}
The goal follows from \(\operatorname{notcenterperiodicofrdrvdyadicshiftedinitagreetwopow}(T, P, td, m, t, k, L, u, S)\).
\end{proof}
\begin{theorem}
\label{ll:paper:not-centerperiodic-of-rdrv-dyadic-shifted-init-agree-at-zero-diag-pair}
For every natural number \(T\) and natural number \(P\) and natural number \(td\) and natural number \(m\) and natural number \(z\) and natural number \(k\) and natural number \(L\) and natural number \(u\) and natural number \(S\), if \(0 < P\), then if \(\operatorname{zerodiagonal}(z)\), then if \(2 \le z + m \cdot P\), then if \(\operatorname{zerodiagonal}(z + m \cdot P - 2)\), then if \(T \le z + \operatorname{twopow}(k)\), then if \(\operatorname{settledby}(z, z + \operatorname{twopow}(k))\), then if \(\operatorname{settledby}(z + m \cdot P, z + m \cdot P + \operatorname{twopow}(k))\), then if \(\operatorname{twopow}(k) \le S + 1\), then if \(1 \le u\), then if \(z + \operatorname{twopow}(k) + S + m \cdot P + 1 \le L\), then if for every natural number \(s\), if \(s \le S\), then \(\operatorname{rdrv}(td, z + \operatorname{twopow}(k) + u \cdot \operatorname{twopow}(L) + m \cdot P, s, u \cdot \operatorname{twopow}(L) - 1) = \operatorname{rdrv}(td, z + \operatorname{twopow}(k) + u \cdot \operatorname{twopow}(L), s, u \cdot \operatorname{twopow}(L) - 1)\), then not \(\operatorname{centerperiodic}(T, P)\).
\end{theorem}
\begin{proof}
The goal follows from \(\operatorname{notcenterperiodicofrdrvdyadicshiftedinitagreeatzerodiagpair}(T, P, td, m, z, k, L, u, S)\).
\end{proof}
\begin{theorem}
\label{ll:paper:periodic-lhbandagree-one-step}
For every natural number \(T\) and natural number \(P\), if \(0 < P\), then if \(\operatorname{centerperiodic}(T, P)\), then for every natural number \(m\) and natural number \(t\) and natural number \(L\), if \(T \le t\), then if \(\operatorname{lhbandagree}(t, m \cdot P, 1, L + 1)\), then \(\operatorname{lhbandagree}(t + 1, m \cdot P, 1, L)\).
\end{theorem}
\begin{proof}
Assume \(hP\), \(hper\), \(m\), \(t\), \(L\).
The goal follows from \(\operatorname{periodiclhbandagreeonestep}(T, P, hP, hper, m, t, L)\).
\end{proof}
\begin{theorem}
\label{ll:paper:periodic-lhbandagree-one-steps}
For every natural number \(T\) and natural number \(P\), if \(0 < P\), then if \(\operatorname{centerperiodic}(T, P)\), then for every natural number \(m\) and natural number \(t\) and natural number \(L\) and natural number \(s\), if \(T \le t\), then if \(\operatorname{lhbandagree}(t, m \cdot P, 1, L + s)\), then \(\operatorname{lhbandagree}(t + s, m \cdot P, 1, L)\).
\end{theorem}
\begin{proof}
Assume \(hP\), \(hper\), \(m\), \(t\), \(L\), \(s\).
The goal follows from \(\operatorname{periodiclhbandagreeonesteps}(T, P, hP, hper, m, t, L, s)\).
\end{proof}
\begin{theorem}
\label{ll:paper:periodic-lhcenter-one-agree-of-lhbandagree-one}
For every natural number \(T\) and natural number \(P\), if \(0 < P\), then if \(\operatorname{centerperiodic}(T, P)\), then for every natural number \(m\) and natural number \(t\) and natural number \(S\) and natural number \(s\), if \(T \le t\), then if \(\operatorname{lhbandagree}(t, m \cdot P, 1, S)\), then if \(s \le S\), then \(\operatorname{lhcenter}(t + s + m \cdot P, 1) = \operatorname{lhcenter}(t + s, 1)\).
\end{theorem}
\begin{proof}
Assume \(hP\), \(hper\), \(m\), \(t\), \(S\), \(s\).
The goal follows from \(\operatorname{periodiclhcenteroneagreeoflhbandagreeone}(T, P, hP, hper, m, t, S, s)\).
\end{proof}
\begin{theorem}
\label{ll:paper:periodic-lhcolagree-one-of-lhbandagree-one}
For every natural number \(T\) and natural number \(P\), if \(0 < P\), then if \(\operatorname{centerperiodic}(T, P)\), then for every natural number \(m\) and natural number \(t\) and natural number \(S\), if \(T \le t\), then if \(\operatorname{lhbandagree}(t, m \cdot P, 1, S)\), then \(\operatorname{lhcolagree}(t, m \cdot P, 1, S)\).
\end{theorem}
\begin{proof}
Assume \(hP\), \(hper\), \(m\), \(t\), \(S\).
The goal follows from \(\operatorname{periodiclhcolagreeoneoflhbandagreeone}(T, P, hP, hper, m, t, S)\).
\end{proof}
\begin{theorem}
\label{ll:paper:periodic-maskparity-zero-agree-of-lhbandagree-one}
For every natural number \(T\) and natural number \(P\), if \(0 < P\), then if \(\operatorname{centerperiodic}(T, P)\), then for every natural number \(m\) and natural number \(t\) and natural number \(S\) and natural number \(s\), if \(T \le t\), then if \(\operatorname{lhbandagree}(t, m \cdot P, 1, S)\), then if \(s \le S\), then \(\operatorname{maskparity}(t + s + m \cdot P, 0) = \operatorname{maskparity}(t + s, 0)\).
\end{theorem}
\begin{proof}
Assume \(hP\), \(hper\), \(m\), \(t\), \(S\), \(s\).
The goal follows from \(\operatorname{periodicmaskparityzeroagreeoflhbandagreeone}(T, P, hP, hper, m, t, S, s)\).
\end{proof}
\begin{theorem}
\label{ll:paper:not-centerperiodic-of-lhbandagree-one-at-zero-diag-doubles}
For every natural number \(T\) and natural number \(P\) and natural number \(m\) and natural number \(z\) and natural number \(p\) and natural number \(S\), if \(0 < P\), then if \(0 < m\), then if \(T \le z + (p + 1)\), then if \(\operatorname{zerodiagonal}(z)\), then if \(\operatorname{doublesat}(m \cdot P, z)\), then if \(\operatorname{settledby}(z + 1, z + (p + 1))\), then if \(\operatorname{settledby}(z + 1 - m \cdot P, z + (p + 1))\), then if \(p + 3 \cdot (m \cdot P) \le S + 1\), then if \(\operatorname{lhbandagree}(z + (p + 1), m \cdot P, 1, S)\), then not \(\operatorname{centerperiodic}(T, P)\).
\end{theorem}
\begin{proof}
The goal follows from \(\operatorname{notcenterperiodicoflhbandagreeoneatzerodiagdoubles}(T, P, m, z, p, S)\).
\end{proof}
\begin{theorem}
\label{ll:paper:not-centerperiodic-of-lhbandagree-one-twopow}
For every natural number \(T\) and natural number \(P\) and natural number \(m\) and natural number \(t\) and natural number \(k\) and natural number \(S\), if \(0 < P\), then if \(0 < m\), then if \(T \le t\), then if \(t + m \cdot P = \operatorname{twopow}(k)\), then if \(\operatorname{twopow}(k) \le S + 1\), then if \(\operatorname{lhbandagree}(t, m \cdot P, 1, S)\), then not \(\operatorname{centerperiodic}(T, P)\).
\end{theorem}
\begin{proof}
The goal follows from \(\operatorname{notcenterperiodicoflhbandagreeonetwopow}(T, P, m, t, k, S)\).
\end{proof}
\begin{theorem}
\label{ll:paper:not-centerperiodic-of-lhbandagree-one-at-zero-diag-pair}
For every natural number \(T\) and natural number \(P\) and natural number \(m\) and natural number \(z\) and natural number \(k\) and natural number \(S\), if \(0 < P\), then if \(\operatorname{zerodiagonal}(z)\), then if \(2 \le z + m \cdot P\), then if \(\operatorname{zerodiagonal}(z + m \cdot P - 2)\), then if \(T \le z + \operatorname{twopow}(k)\), then if \(\operatorname{settledby}(z, z + \operatorname{twopow}(k))\), then if \(\operatorname{settledby}(z + m \cdot P, z + m \cdot P + \operatorname{twopow}(k))\), then if \(\operatorname{twopow}(k) \le S + 1\), then if \(\operatorname{lhbandagree}(z + \operatorname{twopow}(k), m \cdot P, 1, S)\), then not \(\operatorname{centerperiodic}(T, P)\).
\end{theorem}
\begin{proof}
The goal follows from \(\operatorname{notcenterperiodicoflhbandagreeoneatzerodiagpair}(T, P, m, z, k, S)\).
\end{proof}
\begin{theorem}
\label{ll:paper:lhbandagree-one-of-rhcenter-dyadic-left-corridor-agree}
For every natural number \(t\) and natural number \(Q\) and natural number \(K\) and natural number \(v\) and natural number \(S\), if \(1 \le v\), then if \(t + Q + (S + 1) \le K\), then if for every natural number \(i\), if \(i \le S\), then \(\operatorname{rhcenter}(t + v \cdot \operatorname{twopow}(K) + Q, v \cdot \operatorname{twopow}(K) - (1 + i)) = \operatorname{rhcenter}(t + v \cdot \operatorname{twopow}(K), v \cdot \operatorname{twopow}(K) - (1 + i))\), then \(\operatorname{lhbandagree}(t, Q, 1, S)\).
\end{theorem}
\begin{proof}
The goal follows from \(\operatorname{lhbandagreeoneofrhcenterdyadicleftcorridoragree}(t, Q, K, v, S)\).
\end{proof}
\begin{theorem}
\label{ll:paper:lhbandagree-one-of-rhbandagree-dyadic-left-corridor}
For every natural number \(t\) and natural number \(Q\) and natural number \(K\) and natural number \(v\) and natural number \(S\), if \(1 \le v\), then if \(t + Q + (S + 1) \le K\), then if \(\operatorname{rhbandagree}(t + v \cdot \operatorname{twopow}(K), Q, v \cdot \operatorname{twopow}(K) - (S + 1), S)\), then \(\operatorname{lhbandagree}(t, Q, 1, S)\).
\end{theorem}
\begin{proof}
The goal follows from \(\operatorname{lhbandagreeoneofrhbandagreedyadicleftcorridor}(t, Q, K, v, S)\).
\end{proof}
\begin{theorem}
\label{ll:paper:periodic-lhcolagree-one-of-rhbandagree-dyadic-left-corridor}
For every natural number \(T\) and natural number \(P\), if \(0 < P\), then if \(\operatorname{centerperiodic}(T, P)\), then for every natural number \(m\) and natural number \(t\) and natural number \(K\) and natural number \(v\) and natural number \(S\), if \(T \le t\), then if \(1 \le v\), then if \(t + m \cdot P + (S + 1) \le K\), then if \(\operatorname{rhbandagree}(t + v \cdot \operatorname{twopow}(K), m \cdot P, v \cdot \operatorname{twopow}(K) - (S + 1), S)\), then \(\operatorname{lhcolagree}(t, m \cdot P, 1, S)\).
\end{theorem}
\begin{proof}
Assume \(hP\), \(hper\), \(m\), \(t\), \(K\), \(v\), \(S\).
The goal follows from \(\operatorname{periodiclhcolagreeoneofrhbandagreedyadicleftcorridor}(T, P, hP, hper, m, t, K, v, S)\).
\end{proof}
\begin{theorem}
\label{ll:paper:not-centerperiodic-of-rhbandagree-dyadic-left-corridor-at-zero-diag-doubles}
For every natural number \(T\) and natural number \(P\) and natural number \(m\) and natural number \(z\) and natural number \(p\) and natural number \(K\) and natural number \(v\) and natural number \(S\), if \(0 < P\), then if \(0 < m\), then if \(T \le z + (p + 1)\), then if \(\operatorname{zerodiagonal}(z)\), then if \(\operatorname{doublesat}(m \cdot P, z)\), then if \(\operatorname{settledby}(z + 1, z + (p + 1))\), then if \(\operatorname{settledby}(z + 1 - m \cdot P, z + (p + 1))\), then if \(p + 3 \cdot (m \cdot P) \le S + 1\), then if \(1 \le v\), then if \(z + (p + 1) + m \cdot P + (S + 1) \le K\), then if \(\operatorname{rhbandagree}(z + (p + 1) + v \cdot \operatorname{twopow}(K), m \cdot P, v \cdot \operatorname{twopow}(K) - (S + 1), S)\), then not \(\operatorname{centerperiodic}(T, P)\).
\end{theorem}
\begin{proof}
The goal follows from \(\operatorname{notcenterperiodicofrhbandagreedyadicleftcorridoratzerodiagdoubles}(T, P, m, z, p, K, v, S)\).
\end{proof}
\begin{theorem}
\label{ll:paper:not-centerperiodic-of-rhbandagree-dyadic-left-corridor-twopow}
For every natural number \(T\) and natural number \(P\) and natural number \(m\) and natural number \(t\) and natural number \(k\) and natural number \(K\) and natural number \(v\) and natural number \(S\), if \(0 < P\), then if \(0 < m\), then if \(T \le t\), then if \(t + m \cdot P = \operatorname{twopow}(k)\), then if \(\operatorname{twopow}(k) \le S + 1\), then if \(1 \le v\), then if \(t + m \cdot P + (S + 1) \le K\), then if \(\operatorname{rhbandagree}(t + v \cdot \operatorname{twopow}(K), m \cdot P, v \cdot \operatorname{twopow}(K) - (S + 1), S)\), then not \(\operatorname{centerperiodic}(T, P)\).
\end{theorem}
\begin{proof}
The goal follows from \(\operatorname{notcenterperiodicofrhbandagreedyadicleftcorridortwopow}(T, P, m, t, k, K, v, S)\).
\end{proof}
\begin{theorem}
\label{ll:paper:not-centerperiodic-of-rhbandagree-dyadic-left-corridor-at-zero-diag-pair}
For every natural number \(T\) and natural number \(P\) and natural number \(m\) and natural number \(z\) and natural number \(k\) and natural number \(K\) and natural number \(v\) and natural number \(S\), if \(0 < P\), then if \(\operatorname{zerodiagonal}(z)\), then if \(2 \le z + m \cdot P\), then if \(\operatorname{zerodiagonal}(z + m \cdot P - 2)\), then if \(T \le z + \operatorname{twopow}(k)\), then if \(\operatorname{settledby}(z, z + \operatorname{twopow}(k))\), then if \(\operatorname{settledby}(z + m \cdot P, z + m \cdot P + \operatorname{twopow}(k))\), then if \(\operatorname{twopow}(k) \le S + 1\), then if \(1 \le v\), then if \(z + \operatorname{twopow}(k) + m \cdot P + (S + 1) \le K\), then if \(\operatorname{rhbandagree}(z + \operatorname{twopow}(k) + v \cdot \operatorname{twopow}(K), m \cdot P, v \cdot \operatorname{twopow}(K) - (S + 1), S)\), then not \(\operatorname{centerperiodic}(T, P)\).
\end{theorem}
\begin{proof}
The goal follows from \(\operatorname{notcenterperiodicofrhbandagreedyadicleftcorridoratzerodiagpair}(T, P, m, z, k, K, v, S)\).
\end{proof}
\begin{theorem}
\label{ll:paper:lhbandagree-one-of-rdrv-dyadic-exact-row-left-corridor-agree}
For every natural number \(td\) and natural number \(t\) and natural number \(Q\) and natural number \(K\) and natural number \(v\) and natural number \(S\), if \(1 \le v\), then if \(t + Q + (S + 1) \le K\), then if for every natural number \(i\), if \(i \le S\), then \(\operatorname{rdrv}(td, v \cdot \operatorname{twopow}(K), t + Q, v \cdot \operatorname{twopow}(K) - (1 + i)) = \operatorname{rdrv}(td, v \cdot \operatorname{twopow}(K), t, v \cdot \operatorname{twopow}(K) - (1 + i))\), then \(\operatorname{lhbandagree}(t, Q, 1, S)\).
\end{theorem}
\begin{proof}
The goal follows from \(\operatorname{lhbandagreeoneofrdrvdyadicexactrowleftcorridoragree}(td, t, Q, K, v, S)\).
\end{proof}
\begin{theorem}
\label{ll:paper:lhbandagree-one-of-two-scale-single-row-left-corridor-spatial-agree}
For every natural number \(t\) and natural number \(K\) and natural number \(L\) and natural number \(u\) and natural number \(w\) and natural number \(S\), if \(1 \le u\), then if \(t + S + 1 \le K\), then if \(t + S + w \cdot \operatorname{twopow}(K) + 1 \le L\), then if \(K \le L\), then if for every natural number \(i\), if \(i \le S\), then \(\operatorname{rhcenter}(t + u \cdot \operatorname{twopow}(L) + w \cdot \operatorname{twopow}(K), u \cdot \operatorname{twopow}(L) - (1 + i)) = \operatorname{rhcenter}(t + u \cdot \operatorname{twopow}(L) + w \cdot \operatorname{twopow}(K), u \cdot \operatorname{twopow}(L) + w \cdot \operatorname{twopow}(K) - (1 + i))\), then \(\operatorname{lhbandagree}(t, w \cdot \operatorname{twopow}(K), 1, S)\).
\end{theorem}
\begin{proof}
The goal follows from \(\operatorname{lhbandagreeoneoftwoscalesinglerowleftcorridorspatialagree}(t, K, L, u, w, S)\).
\end{proof}
\begin{theorem}
\label{ll:paper:lhbandagree-one-of-dyadic-mul-period-single-row-left-corridor-spatial-agree}
For every natural number \(t\) and natural number \(P\) and natural number \(c\) and natural number \(K\) and natural number \(L\) and natural number \(u\) and natural number \(S\), if \(1 \le u\), then if \(t + S + 1 \le K\), then if \(t + S + c \cdot \operatorname{twopow}(K) \cdot P + 1 \le L\), then if \(K \le L\), then if for every natural number \(i\), if \(i \le S\), then \(\operatorname{rhcenter}(t + u \cdot \operatorname{twopow}(L) + c \cdot \operatorname{twopow}(K) \cdot P, u \cdot \operatorname{twopow}(L) - (1 + i)) = \operatorname{rhcenter}(t + u \cdot \operatorname{twopow}(L) + c \cdot \operatorname{twopow}(K) \cdot P, u \cdot \operatorname{twopow}(L) + c \cdot \operatorname{twopow}(K) \cdot P - (1 + i))\), then \(\operatorname{lhbandagree}(t, c \cdot \operatorname{twopow}(K) \cdot P, 1, S)\).
\end{theorem}
\begin{proof}
The goal follows from \(\operatorname{lhbandagreeoneofdyadicmulperiodsinglerowleftcorridorspatialagree}(t, P, c, K, L, u, S)\).
\end{proof}
\begin{theorem}
\label{ll:paper:periodic-lhcolagree-one-of-dyadic-mul-period-single-row-left-corridor-spatial-agree}
For every natural number \(T\) and natural number \(P\), if \(0 < P\), then if \(\operatorname{centerperiodic}(T, P)\), then for every natural number \(c\) and natural number \(t\) and natural number \(K\) and natural number \(L\) and natural number \(u\) and natural number \(S\), if \(T \le t\), then if \(1 \le u\), then if \(t + S + 1 \le K\), then if \(t + S + c \cdot \operatorname{twopow}(K) \cdot P + 1 \le L\), then if \(K \le L\), then if for every natural number \(i\), if \(i \le S\), then \(\operatorname{rhcenter}(t + u \cdot \operatorname{twopow}(L) + c \cdot \operatorname{twopow}(K) \cdot P, u \cdot \operatorname{twopow}(L) - (1 + i)) = \operatorname{rhcenter}(t + u \cdot \operatorname{twopow}(L) + c \cdot \operatorname{twopow}(K) \cdot P, u \cdot \operatorname{twopow}(L) + c \cdot \operatorname{twopow}(K) \cdot P - (1 + i))\), then \(\operatorname{lhcolagree}(t, c \cdot \operatorname{twopow}(K) \cdot P, 1, S)\).
\end{theorem}
\begin{proof}
Assume \(hP\), \(hper\), \(c\), \(t\), \(K\), \(L\), \(u\), \(S\).
The goal follows from \(\operatorname{periodiclhcolagreeoneofdyadicmulperiodsinglerowleftcorridorspatialagree}(T, P, hP, hper, c, t, K, L, u, S)\).
\end{proof}
\begin{theorem}
\label{ll:paper:periodic-maskparity-zero-agree-of-dyadic-mul-period-single-row-left-corridor-spatial-agree}
For every natural number \(T\) and natural number \(P\), if \(0 < P\), then if \(\operatorname{centerperiodic}(T, P)\), then for every natural number \(c\) and natural number \(t\) and natural number \(K\) and natural number \(L\) and natural number \(u\) and natural number \(S\) and natural number \(s\), if \(T \le t\), then if \(1 \le u\), then if \(t + S + 1 \le K\), then if \(t + S + c \cdot \operatorname{twopow}(K) \cdot P + 1 \le L\), then if \(K \le L\), then if for every natural number \(i\), if \(i \le S\), then \(\operatorname{rhcenter}(t + u \cdot \operatorname{twopow}(L) + c \cdot \operatorname{twopow}(K) \cdot P, u \cdot \operatorname{twopow}(L) - (1 + i)) = \operatorname{rhcenter}(t + u \cdot \operatorname{twopow}(L) + c \cdot \operatorname{twopow}(K) \cdot P, u \cdot \operatorname{twopow}(L) + c \cdot \operatorname{twopow}(K) \cdot P - (1 + i))\), then if \(s \le S\), then \(\operatorname{maskparity}(t + s + c \cdot \operatorname{twopow}(K) \cdot P, 0) = \operatorname{maskparity}(t + s, 0)\).
\end{theorem}
\begin{proof}
Assume \(hP\), \(hper\), \(c\), \(t\), \(K\), \(L\), \(u\), \(S\), \(s\).
The goal follows from \(\operatorname{periodicmaskparityzeroagreeofdyadicmulperiodsinglerowleftcorridorspatialagree}(T, P, hP, hper, c, t, K, L, u, S, s)\).
\end{proof}
\begin{theorem}
\label{ll:paper:not-centerperiodic-of-dyadic-mul-period-single-row-left-corridor-spatial-agree-at-zero-diag-doubles}
For every natural number \(T\) and natural number \(P\) and natural number \(c\) and natural number \(z\) and natural number \(p\) and natural number \(K\) and natural number \(L\) and natural number \(u\) and natural number \(S\), if \(0 < P\), then if \(0 < c\), then if \(T \le z + (p + 1)\), then if \(\operatorname{zerodiagonal}(z)\), then if \(\operatorname{doublesat}(c \cdot \operatorname{twopow}(K) \cdot P, z)\), then if \(\operatorname{settledby}(z + 1, z + (p + 1))\), then if \(\operatorname{settledby}(z + 1 - c \cdot \operatorname{twopow}(K) \cdot P, z + (p + 1))\), then if \(p + 3 \cdot (c \cdot \operatorname{twopow}(K) \cdot P) \le S + 1\), then if \(1 \le u\), then if \(z + (p + 1) + S + 1 \le K\), then if \(z + (p + 1) + S + c \cdot \operatorname{twopow}(K) \cdot P + 1 \le L\), then if \(K \le L\), then if for every natural number \(i\), if \(i \le S\), then \(\operatorname{rhcenter}(z + (p + 1) + u \cdot \operatorname{twopow}(L) + c \cdot \operatorname{twopow}(K) \cdot P, u \cdot \operatorname{twopow}(L) - (1 + i)) = \operatorname{rhcenter}(z + (p + 1) + u \cdot \operatorname{twopow}(L) + c \cdot \operatorname{twopow}(K) \cdot P, u \cdot \operatorname{twopow}(L) + c \cdot \operatorname{twopow}(K) \cdot P - (1 + i))\), then not \(\operatorname{centerperiodic}(T, P)\).
\end{theorem}
\begin{proof}
The goal follows from \(\operatorname{notcenterperiodicofdyadicmulperiodsinglerowleftcorridorspatialagreeatzerodiagdoubles}(T, P, c, z, p, K, L, u, S)\).
\end{proof}
\begin{theorem}
\label{ll:paper:not-centerperiodic-of-dyadic-mul-period-single-row-left-corridor-spatial-agree-twopow}
For every natural number \(T\) and natural number \(P\) and natural number \(c\) and natural number \(t\) and natural number \(k\) and natural number \(K\) and natural number \(L\) and natural number \(u\) and natural number \(S\), if \(0 < P\), then if \(0 < c\), then if \(T \le t\), then if \(t + c \cdot \operatorname{twopow}(K) \cdot P = \operatorname{twopow}(k)\), then if \(\operatorname{twopow}(k) \le S + 1\), then if \(1 \le u\), then if \(t + S + 1 \le K\), then if \(t + S + c \cdot \operatorname{twopow}(K) \cdot P + 1 \le L\), then if \(K \le L\), then if for every natural number \(i\), if \(i \le S\), then \(\operatorname{rhcenter}(t + u \cdot \operatorname{twopow}(L) + c \cdot \operatorname{twopow}(K) \cdot P, u \cdot \operatorname{twopow}(L) - (1 + i)) = \operatorname{rhcenter}(t + u \cdot \operatorname{twopow}(L) + c \cdot \operatorname{twopow}(K) \cdot P, u \cdot \operatorname{twopow}(L) + c \cdot \operatorname{twopow}(K) \cdot P - (1 + i))\), then not \(\operatorname{centerperiodic}(T, P)\).
\end{theorem}
\begin{proof}
The goal follows from \(\operatorname{notcenterperiodicofdyadicmulperiodsinglerowleftcorridorspatialagreetwopow}(T, P, c, t, k, K, L, u, S)\).
\end{proof}
\begin{theorem}
\label{ll:paper:not-centerperiodic-of-dyadic-mul-period-single-row-left-corridor-spatial-agree-at-zero-diag-pair}
For every natural number \(T\) and natural number \(P\) and natural number \(c\) and natural number \(z\) and natural number \(k\) and natural number \(K\) and natural number \(L\) and natural number \(u\) and natural number \(S\), if \(0 < P\), then if \(\operatorname{zerodiagonal}(z)\), then if \(2 \le z + c \cdot \operatorname{twopow}(K) \cdot P\), then if \(\operatorname{zerodiagonal}(z + c \cdot \operatorname{twopow}(K) \cdot P - 2)\), then if \(T \le z + \operatorname{twopow}(k)\), then if \(\operatorname{settledby}(z, z + \operatorname{twopow}(k))\), then if \(\operatorname{settledby}(z + c \cdot \operatorname{twopow}(K) \cdot P, z + c \cdot \operatorname{twopow}(K) \cdot P + \operatorname{twopow}(k))\), then if \(\operatorname{twopow}(k) \le S + 1\), then if \(1 \le u\), then if \(z + \operatorname{twopow}(k) + S + 1 \le K\), then if \(z + \operatorname{twopow}(k) + S + c \cdot \operatorname{twopow}(K) \cdot P + 1 \le L\), then if \(K \le L\), then if for every natural number \(i\), if \(i \le S\), then \(\operatorname{rhcenter}(z + \operatorname{twopow}(k) + u \cdot \operatorname{twopow}(L) + c \cdot \operatorname{twopow}(K) \cdot P, u \cdot \operatorname{twopow}(L) - (1 + i)) = \operatorname{rhcenter}(z + \operatorname{twopow}(k) + u \cdot \operatorname{twopow}(L) + c \cdot \operatorname{twopow}(K) \cdot P, u \cdot \operatorname{twopow}(L) + c \cdot \operatorname{twopow}(K) \cdot P - (1 + i))\), then not \(\operatorname{centerperiodic}(T, P)\).
\end{theorem}
\begin{proof}
The goal follows from \(\operatorname{notcenterperiodicofdyadicmulperiodsinglerowleftcorridorspatialagreeatzerodiagpair}(T, P, c, z, k, K, L, u, S)\).
\end{proof}
\begin{theorem}
\label{ll:paper:rdrv-two-scale-row-left-eq-lhcenter}
For every natural number \(td\) and natural number \(t\) and natural number \(K\) and natural number \(L\) and natural number \(u\) and natural number \(w\) and natural number \(d\), if \(1 \le u\), then if \(K \le L\), then if \(t + d \le K\), then if \(t + w \cdot \operatorname{twopow}(K) + d \le L\), then \(\operatorname{rdrv}(td, u \cdot \operatorname{twopow}(L), t + w \cdot \operatorname{twopow}(K), u \cdot \operatorname{twopow}(L) + w \cdot \operatorname{twopow}(K) - d) = \operatorname{lhcenter}(t, d)\).
\end{theorem}
\begin{proof}
The goal follows from \(\operatorname{rdrvtwoscalerowlefteqlhcenter}(td, t, K, L, u, w, d)\).
\end{proof}
\begin{theorem}
\label{ll:paper:rdrv-two-scale-row-right-eq-rhcenter}
For every natural number \(td\) and natural number \(t\) and natural number \(K\) and natural number \(L\) and natural number \(u\) and natural number \(w\) and natural number \(q\), if \(1 \le u\), then if \(K \le L\), then if \(t \le K\), then if \(t + w \cdot \operatorname{twopow}(K) \le L\), then \(\operatorname{rdrv}(td, u \cdot \operatorname{twopow}(L), t + w \cdot \operatorname{twopow}(K), u \cdot \operatorname{twopow}(L) + w \cdot \operatorname{twopow}(K) + q) = \operatorname{rhcenter}(t, q)\).
\end{theorem}
\begin{proof}
The goal follows from \(\operatorname{rdrvtwoscalerowrighteqrhcenter}(td, t, K, L, u, w, q)\).
\end{proof}
\begin{theorem}
\label{ll:paper:lhbandagree-one-of-rdrv-two-scale-single-row-left-corridor-spatial-agree}
For every natural number \(td\) and natural number \(t\) and natural number \(K\) and natural number \(L\) and natural number \(u\) and natural number \(w\) and natural number \(S\), if \(1 \le u\), then if \(t + S + 1 \le K\), then if \(t + S + w \cdot \operatorname{twopow}(K) + 1 \le L\), then if \(K \le L\), then if for every natural number \(i\), if \(i \le S\), then \(\operatorname{rdrv}(td, u \cdot \operatorname{twopow}(L), t + w \cdot \operatorname{twopow}(K), u \cdot \operatorname{twopow}(L) - (1 + i)) = \operatorname{rdrv}(td, u \cdot \operatorname{twopow}(L), t + w \cdot \operatorname{twopow}(K), u \cdot \operatorname{twopow}(L) + w \cdot \operatorname{twopow}(K) - (1 + i))\), then \(\operatorname{lhbandagree}(t, w \cdot \operatorname{twopow}(K), 1, S)\).
\end{theorem}
\begin{proof}
The goal follows from \(\operatorname{lhbandagreeoneofrdrvtwoscalesinglerowleftcorridorspatialagree}(td, t, K, L, u, w, S)\).
\end{proof}
\begin{theorem}
\label{ll:paper:periodic-lhcolagree-one-of-rdrv-dyadic-mul-period-single-row-left-corridor-spatial-agree}
For every natural number \(T\) and natural number \(P\), if \(0 < P\), then if \(\operatorname{centerperiodic}(T, P)\), then for every natural number \(td\) and natural number \(c\) and natural number \(t\) and natural number \(K\) and natural number \(L\) and natural number \(u\) and natural number \(S\), if \(T \le t\), then if \(1 \le u\), then if \(t + S + 1 \le K\), then if \(t + S + c \cdot \operatorname{twopow}(K) \cdot P + 1 \le L\), then if \(K \le L\), then if for every natural number \(i\), if \(i \le S\), then \(\operatorname{rdrv}(td, u \cdot \operatorname{twopow}(L), t + c \cdot \operatorname{twopow}(K) \cdot P, u \cdot \operatorname{twopow}(L) - (1 + i)) = \operatorname{rdrv}(td, u \cdot \operatorname{twopow}(L), t + c \cdot \operatorname{twopow}(K) \cdot P, u \cdot \operatorname{twopow}(L) + c \cdot \operatorname{twopow}(K) \cdot P - (1 + i))\), then \(\operatorname{lhcolagree}(t, c \cdot \operatorname{twopow}(K) \cdot P, 1, S)\).
\end{theorem}
\begin{proof}
Assume \(hP\), \(hper\), \(td\), \(c\), \(t\), \(K\), \(L\), \(u\), \(S\).
The goal follows from \(\operatorname{periodiclhcolagreeoneofrdrvdyadicmulperiodsinglerowleftcorridorspatialagree}(T, P, hP, hper, td, c, t, K, L, u, S)\).
\end{proof}
\begin{theorem}
\label{ll:paper:periodic-maskparity-zero-agree-of-rdrv-dyadic-mul-period-single-row-left-corridor-spatial-agree}
For every natural number \(T\) and natural number \(P\), if \(0 < P\), then if \(\operatorname{centerperiodic}(T, P)\), then for every natural number \(td\) and natural number \(c\) and natural number \(t\) and natural number \(K\) and natural number \(L\) and natural number \(u\) and natural number \(S\) and natural number \(s\), if \(T \le t\), then if \(1 \le u\), then if \(t + S + 1 \le K\), then if \(t + S + c \cdot \operatorname{twopow}(K) \cdot P + 1 \le L\), then if \(K \le L\), then if for every natural number \(i\), if \(i \le S\), then \(\operatorname{rdrv}(td, u \cdot \operatorname{twopow}(L), t + c \cdot \operatorname{twopow}(K) \cdot P, u \cdot \operatorname{twopow}(L) - (1 + i)) = \operatorname{rdrv}(td, u \cdot \operatorname{twopow}(L), t + c \cdot \operatorname{twopow}(K) \cdot P, u \cdot \operatorname{twopow}(L) + c \cdot \operatorname{twopow}(K) \cdot P - (1 + i))\), then if \(s \le S\), then \(\operatorname{maskparity}(t + s + c \cdot \operatorname{twopow}(K) \cdot P, 0) = \operatorname{maskparity}(t + s, 0)\).
\end{theorem}
\begin{proof}
Assume \(hP\), \(hper\), \(td\), \(c\), \(t\), \(K\), \(L\), \(u\), \(S\), \(s\).
The goal follows from \(\operatorname{periodicmaskparityzeroagreeofrdrvdyadicmulperiodsinglerowleftcorridorspatialagree}(T, P, hP, hper, td, c, t, K, L, u, S, s)\).
\end{proof}
\begin{theorem}
\label{ll:paper:not-centerperiodic-of-rdrv-dyadic-mul-period-single-row-left-corridor-spatial-agree-at-zero-diag-doubles}
For every natural number \(T\) and natural number \(P\) and natural number \(td\) and natural number \(c\) and natural number \(z\) and natural number \(p\) and natural number \(K\) and natural number \(L\) and natural number \(u\) and natural number \(S\), if \(0 < P\), then if \(0 < c\), then if \(T \le z + (p + 1)\), then if \(\operatorname{zerodiagonal}(z)\), then if \(\operatorname{doublesat}(c \cdot \operatorname{twopow}(K) \cdot P, z)\), then if \(\operatorname{settledby}(z + 1, z + (p + 1))\), then if \(\operatorname{settledby}(z + 1 - c \cdot \operatorname{twopow}(K) \cdot P, z + (p + 1))\), then if \(p + 3 \cdot (c \cdot \operatorname{twopow}(K) \cdot P) \le S + 1\), then if \(1 \le u\), then if \(z + (p + 1) + S + 1 \le K\), then if \(z + (p + 1) + S + c \cdot \operatorname{twopow}(K) \cdot P + 1 \le L\), then if \(K \le L\), then if for every natural number \(i\), if \(i \le S\), then \(\operatorname{rdrv}(td, u \cdot \operatorname{twopow}(L), z + (p + 1) + c \cdot \operatorname{twopow}(K) \cdot P, u \cdot \operatorname{twopow}(L) - (1 + i)) = \operatorname{rdrv}(td, u \cdot \operatorname{twopow}(L), z + (p + 1) + c \cdot \operatorname{twopow}(K) \cdot P, u \cdot \operatorname{twopow}(L) + c \cdot \operatorname{twopow}(K) \cdot P - (1 + i))\), then not \(\operatorname{centerperiodic}(T, P)\).
\end{theorem}
\begin{proof}
The goal follows from \(\operatorname{notcenterperiodicofrdrvdyadicmulperiodsinglerowleftcorridorspatialagreeatzerodiagdoubles}(T, P, td, c, z, p, K, L, u, S)\).
\end{proof}
\begin{theorem}
\label{ll:paper:not-centerperiodic-of-rdrv-dyadic-mul-period-single-row-left-corridor-spatial-agree-twopow}
For every natural number \(T\) and natural number \(P\) and natural number \(td\) and natural number \(c\) and natural number \(t\) and natural number \(k\) and natural number \(K\) and natural number \(L\) and natural number \(u\) and natural number \(S\), if \(0 < P\), then if \(0 < c\), then if \(T \le t\), then if \(t + c \cdot \operatorname{twopow}(K) \cdot P = \operatorname{twopow}(k)\), then if \(\operatorname{twopow}(k) \le S + 1\), then if \(1 \le u\), then if \(t + S + 1 \le K\), then if \(t + S + c \cdot \operatorname{twopow}(K) \cdot P + 1 \le L\), then if \(K \le L\), then if for every natural number \(i\), if \(i \le S\), then \(\operatorname{rdrv}(td, u \cdot \operatorname{twopow}(L), t + c \cdot \operatorname{twopow}(K) \cdot P, u \cdot \operatorname{twopow}(L) - (1 + i)) = \operatorname{rdrv}(td, u \cdot \operatorname{twopow}(L), t + c \cdot \operatorname{twopow}(K) \cdot P, u \cdot \operatorname{twopow}(L) + c \cdot \operatorname{twopow}(K) \cdot P - (1 + i))\), then not \(\operatorname{centerperiodic}(T, P)\).
\end{theorem}
\begin{proof}
The goal follows from \(\operatorname{notcenterperiodicofrdrvdyadicmulperiodsinglerowleftcorridorspatialagreetwopow}(T, P, td, c, t, k, K, L, u, S)\).
\end{proof}
\begin{theorem}
\label{ll:paper:not-centerperiodic-of-rdrv-dyadic-mul-period-single-row-left-corridor-spatial-agree-at-zero-diag-pair}
For every natural number \(T\) and natural number \(P\) and natural number \(td\) and natural number \(c\) and natural number \(z\) and natural number \(k\) and natural number \(K\) and natural number \(L\) and natural number \(u\) and natural number \(S\), if \(0 < P\), then if \(\operatorname{zerodiagonal}(z)\), then if \(2 \le z + c \cdot \operatorname{twopow}(K) \cdot P\), then if \(\operatorname{zerodiagonal}(z + c \cdot \operatorname{twopow}(K) \cdot P - 2)\), then if \(T \le z + \operatorname{twopow}(k)\), then if \(\operatorname{settledby}(z, z + \operatorname{twopow}(k))\), then if \(\operatorname{settledby}(z + c \cdot \operatorname{twopow}(K) \cdot P, z + c \cdot \operatorname{twopow}(K) \cdot P + \operatorname{twopow}(k))\), then if \(\operatorname{twopow}(k) \le S + 1\), then if \(1 \le u\), then if \(z + \operatorname{twopow}(k) + S + 1 \le K\), then if \(z + \operatorname{twopow}(k) + S + c \cdot \operatorname{twopow}(K) \cdot P + 1 \le L\), then if \(K \le L\), then if for every natural number \(i\), if \(i \le S\), then \(\operatorname{rdrv}(td, u \cdot \operatorname{twopow}(L), z + \operatorname{twopow}(k) + c \cdot \operatorname{twopow}(K) \cdot P, u \cdot \operatorname{twopow}(L) - (1 + i)) = \operatorname{rdrv}(td, u \cdot \operatorname{twopow}(L), z + \operatorname{twopow}(k) + c \cdot \operatorname{twopow}(K) \cdot P, u \cdot \operatorname{twopow}(L) + c \cdot \operatorname{twopow}(K) \cdot P - (1 + i))\), then not \(\operatorname{centerperiodic}(T, P)\).
\end{theorem}
\begin{proof}
The goal follows from \(\operatorname{notcenterperiodicofrdrvdyadicmulperiodsinglerowleftcorridorspatialagreeatzerodiagpair}(T, P, td, c, z, k, K, L, u, S)\).
\end{proof}
\begin{theorem}
\label{ll:paper:ldiag-settled-mul-twopow-periodic-le}
For every natural number \(e\) and natural number \(w\) and natural number \(d\) and natural number \(t\), if \(d \le e\), then if \(\operatorname{settledby}(d, t)\), then \(\operatorname{ldiag}(d, t + w \cdot \operatorname{twopow}(e)) = \operatorname{ldiag}(d, t)\).
\end{theorem}
\begin{proof}
The goal follows from \(\operatorname{ldiagsettledmultwopowperiodicle}(e, w, d, t)\).
\end{proof}
\begin{theorem}
\label{ll:paper:lhcenter-diag-shift-mul-twopow-periodic-le}
For every natural number \(e\) and natural number \(w\) and natural number \(t\) and natural number \(p\), if \(t - p \le e\), then if \(\operatorname{settledby}(t - p, t)\), then \(\operatorname{lhcenter}(t + w \cdot \operatorname{twopow}(e), p + w \cdot \operatorname{twopow}(e)) = \operatorname{lhcenter}(t, p)\).
\end{theorem}
\begin{proof}
The goal follows from \(\operatorname{lhcenterdiagshiftmultwopowperiodicle}(e, w, t, p)\).
\end{proof}
\begin{theorem}
\label{ll:paper:lhcenter-eq-shifted-row-of-mul-twopow-and-spatial-agree}
For every natural number \(e\) and natural number \(w\) and natural number \(t\) and natural number \(p\), if \(t - p \le e\), then if \(\operatorname{settledby}(t - p, t)\), then if \(\operatorname{lhcenter}(t + w \cdot \operatorname{twopow}(e), p) = \operatorname{lhcenter}(t + w \cdot \operatorname{twopow}(e), p + w \cdot \operatorname{twopow}(e))\), then \(\operatorname{lhcenter}(t + w \cdot \operatorname{twopow}(e), p) = \operatorname{lhcenter}(t, p)\).
\end{theorem}
\begin{proof}
The goal follows from \(\operatorname{lhcentereqshiftedrowofmultwopowandspatialagree}(e, w, t, p)\).
\end{proof}
\begin{theorem}
\label{ll:paper:lhbandagree-of-single-lh-row-mul-twopow-spatial-agree}
For every natural number \(e\) and natural number \(w\) and natural number \(t\) and natural number \(p0\) and natural number \(L\), if \(t - p0 \le e\), then if for every natural number \(i\), if \(i \le L\), then \(\operatorname{settledby}(t - (p0 + i), t)\), then if for every natural number \(i\), if \(i \le L\), then \(\operatorname{lhcenter}(t + w \cdot \operatorname{twopow}(e), p0 + i) = \operatorname{lhcenter}(t + w \cdot \operatorname{twopow}(e), p0 + i + w \cdot \operatorname{twopow}(e))\), then \(\operatorname{lhbandagree}(t, w \cdot \operatorname{twopow}(e), p0, L)\).
\end{theorem}
\begin{proof}
The goal follows from \(\operatorname{lhbandagreeofsinglelhrowmultwopowspatialagree}(e, w, t, p0, L)\).
\end{proof}
\begin{theorem}
\label{ll:paper:lhbandagree-one-of-single-lh-row-mul-twopow-spatial-agree}
For every natural number \(e\) and natural number \(w\) and natural number \(t\) and natural number \(S\), if \(t - 1 \le e\), then if for every natural number \(i\), if \(i \le S\), then \(\operatorname{settledby}(t - (1 + i), t)\), then if for every natural number \(i\), if \(i \le S\), then \(\operatorname{lhcenter}(t + w \cdot \operatorname{twopow}(e), 1 + i) = \operatorname{lhcenter}(t + w \cdot \operatorname{twopow}(e), 1 + i + w \cdot \operatorname{twopow}(e))\), then \(\operatorname{lhbandagree}(t, w \cdot \operatorname{twopow}(e), 1, S)\).
\end{theorem}
\begin{proof}
The goal follows from \(\operatorname{lhbandagreeoneofsinglelhrowmultwopowspatialagree}(e, w, t, S)\).
\end{proof}
\begin{theorem}
\label{ll:paper:lhbandagree-one-of-single-lh-row-dyadic-mul-period-spatial-agree}
For every natural number \(e\) and natural number \(c\) and natural number \(P\) and natural number \(t\) and natural number \(S\), if \(t - 1 \le e\), then if for every natural number \(i\), if \(i \le S\), then \(\operatorname{settledby}(t - (1 + i), t)\), then if for every natural number \(i\), if \(i \le S\), then \(\operatorname{lhcenter}(t + c \cdot \operatorname{twopow}(e) \cdot P, 1 + i) = \operatorname{lhcenter}(t + c \cdot \operatorname{twopow}(e) \cdot P, 1 + i + c \cdot \operatorname{twopow}(e) \cdot P)\), then \(\operatorname{lhbandagree}(t, c \cdot \operatorname{twopow}(e) \cdot P, 1, S)\).
\end{theorem}
\begin{proof}
The goal follows from \(\operatorname{lhbandagreeoneofsinglelhrowdyadicmulperiodspatialagree}(e, c, P, t, S)\).
\end{proof}
\begin{theorem}
\label{ll:paper:periodic-lhcolagree-one-of-single-lh-row-dyadic-mul-period-spatial-agree}
For every natural number \(T\) and natural number \(P\), if \(0 < P\), then if \(\operatorname{centerperiodic}(T, P)\), then for every natural number \(c\) and natural number \(e\) and natural number \(t\) and natural number \(S\), if \(T \le t\), then if \(t - 1 \le e\), then if for every natural number \(i\), if \(i \le S\), then \(\operatorname{settledby}(t - (1 + i), t)\), then if for every natural number \(i\), if \(i \le S\), then \(\operatorname{lhcenter}(t + c \cdot \operatorname{twopow}(e) \cdot P, 1 + i) = \operatorname{lhcenter}(t + c \cdot \operatorname{twopow}(e) \cdot P, 1 + i + c \cdot \operatorname{twopow}(e) \cdot P)\), then \(\operatorname{lhcolagree}(t, c \cdot \operatorname{twopow}(e) \cdot P, 1, S)\).
\end{theorem}
\begin{proof}
Assume \(hP\), \(hper\), \(c\), \(e\), \(t\), \(S\).
The goal follows from \(\operatorname{periodiclhcolagreeoneofsinglelhrowdyadicmulperiodspatialagree}(T, P, hP, hper, c, e, t, S)\).
\end{proof}
\begin{theorem}
\label{ll:paper:periodic-maskparity-zero-agree-of-single-lh-row-dyadic-mul-period-spatial-agree}
For every natural number \(T\) and natural number \(P\), if \(0 < P\), then if \(\operatorname{centerperiodic}(T, P)\), then for every natural number \(c\) and natural number \(e\) and natural number \(t\) and natural number \(S\) and natural number \(s\), if \(T \le t\), then if \(t - 1 \le e\), then if for every natural number \(i\), if \(i \le S\), then \(\operatorname{settledby}(t - (1 + i), t)\), then if for every natural number \(i\), if \(i \le S\), then \(\operatorname{lhcenter}(t + c \cdot \operatorname{twopow}(e) \cdot P, 1 + i) = \operatorname{lhcenter}(t + c \cdot \operatorname{twopow}(e) \cdot P, 1 + i + c \cdot \operatorname{twopow}(e) \cdot P)\), then if \(s \le S\), then \(\operatorname{maskparity}(t + s + c \cdot \operatorname{twopow}(e) \cdot P, 0) = \operatorname{maskparity}(t + s, 0)\).
\end{theorem}
\begin{proof}
Assume \(hP\), \(hper\), \(c\), \(e\), \(t\), \(S\), \(s\).
The goal follows from \(\operatorname{periodicmaskparityzeroagreeofsinglelhrowdyadicmulperiodspatialagree}(T, P, hP, hper, c, e, t, S, s)\).
\end{proof}
\begin{theorem}
\label{ll:paper:not-centerperiodic-of-single-lh-row-dyadic-mul-period-spatial-agree-at-zero-diag-doubles}
For every natural number \(T\) and natural number \(P\) and natural number \(c\) and natural number \(z\) and natural number \(p\) and natural number \(e\) and natural number \(S\), if \(0 < P\), then if \(0 < c\), then if \(T \le z + (p + 1)\), then if \(\operatorname{zerodiagonal}(z)\), then if \(\operatorname{doublesat}(c \cdot \operatorname{twopow}(e) \cdot P, z)\), then if \(\operatorname{settledby}(z + 1, z + (p + 1))\), then if \(\operatorname{settledby}(z + 1 - c \cdot \operatorname{twopow}(e) \cdot P, z + (p + 1))\), then if \(p + 3 \cdot (c \cdot \operatorname{twopow}(e) \cdot P) \le S + 1\), then if \(z + (p + 1) - 1 \le e\), then if for every natural number \(i\), if \(i \le S\), then \(\operatorname{settledby}(z + (p + 1) - (1 + i), z + (p + 1))\), then if for every natural number \(i\), if \(i \le S\), then \(\operatorname{lhcenter}(z + (p + 1) + c \cdot \operatorname{twopow}(e) \cdot P, 1 + i) = \operatorname{lhcenter}(z + (p + 1) + c \cdot \operatorname{twopow}(e) \cdot P, 1 + i + c \cdot \operatorname{twopow}(e) \cdot P)\), then not \(\operatorname{centerperiodic}(T, P)\).
\end{theorem}
\begin{proof}
The goal follows from \(\operatorname{notcenterperiodicofsinglelhrowdyadicmulperiodspatialagreeatzerodiagdoubles}(T, P, c, z, p, e, S)\).
\end{proof}
\begin{theorem}
\label{ll:paper:not-centerperiodic-of-single-lh-row-dyadic-mul-period-spatial-agree-twopow}
For every natural number \(T\) and natural number \(P\) and natural number \(c\) and natural number \(t\) and natural number \(k\) and natural number \(e\) and natural number \(S\), if \(0 < P\), then if \(0 < c\), then if \(T \le t\), then if \(t + c \cdot \operatorname{twopow}(e) \cdot P = \operatorname{twopow}(k)\), then if \(\operatorname{twopow}(k) \le S + 1\), then if \(t - 1 \le e\), then if for every natural number \(i\), if \(i \le S\), then \(\operatorname{settledby}(t - (1 + i), t)\), then if for every natural number \(i\), if \(i \le S\), then \(\operatorname{lhcenter}(t + c \cdot \operatorname{twopow}(e) \cdot P, 1 + i) = \operatorname{lhcenter}(t + c \cdot \operatorname{twopow}(e) \cdot P, 1 + i + c \cdot \operatorname{twopow}(e) \cdot P)\), then not \(\operatorname{centerperiodic}(T, P)\).
\end{theorem}
\begin{proof}
The goal follows from \(\operatorname{notcenterperiodicofsinglelhrowdyadicmulperiodspatialagreetwopow}(T, P, c, t, k, e, S)\).
\end{proof}
\begin{theorem}
\label{ll:paper:lhbandagree-of-doubles-le-and-single-lh-row-spatial-agree}
For every natural number \(Q\) and natural number \(z\) and natural number \(t\) and natural number \(p0\) and natural number \(L\), if \(\operatorname{doublesat}(Q, z)\), then if \(t - p0 \le z\), then if for every natural number \(i\), if \(i \le L\), then \(\operatorname{settledby}(t - (p0 + i), t)\), then if for every natural number \(i\), if \(i \le L\), then \(\operatorname{lhcenter}(t + Q, p0 + i) = \operatorname{lhcenter}(t + Q, p0 + i + Q)\), then \(\operatorname{lhbandagree}(t, Q, p0, L)\).
\end{theorem}
\begin{proof}
The goal follows from \(\operatorname{lhbandagreeofdoublesleandsinglelhrowspatialagree}(Q, z, t, p0, L)\).
\end{proof}
\begin{theorem}
\label{ll:paper:lhbandagree-one-at-zero-diag-succ-of-doubles-and-single-lh-row-spatial-agree}
For every natural number \(Q\) and natural number \(z\) and natural number \(S\), if \(\operatorname{doublesat}(Q, z)\), then if for every natural number \(i\), if \(i \le S\), then \(\operatorname{settledby}(z + 1 - (1 + i), z + 1)\), then if for every natural number \(i\), if \(i \le S\), then \(\operatorname{lhcenter}(z + 1 + Q, 1 + i) = \operatorname{lhcenter}(z + 1 + Q, 1 + i + Q)\), then \(\operatorname{lhbandagree}(z + 1, Q, 1, S)\).
\end{theorem}
\begin{proof}
The goal follows from \(\operatorname{lhbandagreeoneatzerodiagsuccofdoublesandsinglelhrowspatialagree}(Q, z, S)\).
\end{proof}
\begin{theorem}
\label{ll:paper:periodic-lhcolagree-one-at-zero-diag-succ-of-doubles-and-single-lh-row-spatial-agree}
For every natural number \(T\) and natural number \(P\), if \(0 < P\), then if \(\operatorname{centerperiodic}(T, P)\), then for every natural number \(m\) and natural number \(z\) and natural number \(S\), if \(T \le z + 1\), then if \(\operatorname{doublesat}(m \cdot P, z)\), then if for every natural number \(i\), if \(i \le S\), then \(\operatorname{settledby}(z + 1 - (1 + i), z + 1)\), then if for every natural number \(i\), if \(i \le S\), then \(\operatorname{lhcenter}(z + 1 + m \cdot P, 1 + i) = \operatorname{lhcenter}(z + 1 + m \cdot P, 1 + i + m \cdot P)\), then \(\operatorname{lhcolagree}(z + 1, m \cdot P, 1, S)\).
\end{theorem}
\begin{proof}
Assume \(hP\), \(hper\), \(m\), \(z\), \(S\).
The goal follows from \(\operatorname{periodiclhcolagreeoneatzerodiagsuccofdoublesandsinglelhrowspatialagree}(T, P, hP, hper, m, z, S)\).
\end{proof}
\begin{theorem}
\label{ll:paper:periodic-maskparity-zero-agree-at-zero-diag-succ-of-doubles-and-single-lh-row-spatial-agree}
For every natural number \(T\) and natural number \(P\), if \(0 < P\), then if \(\operatorname{centerperiodic}(T, P)\), then for every natural number \(m\) and natural number \(z\) and natural number \(S\) and natural number \(s\), if \(T \le z + 1\), then if \(\operatorname{doublesat}(m \cdot P, z)\), then if for every natural number \(i\), if \(i \le S\), then \(\operatorname{settledby}(z + 1 - (1 + i), z + 1)\), then if for every natural number \(i\), if \(i \le S\), then \(\operatorname{lhcenter}(z + 1 + m \cdot P, 1 + i) = \operatorname{lhcenter}(z + 1 + m \cdot P, 1 + i + m \cdot P)\), then if \(s \le S\), then \(\operatorname{maskparity}(z + 1 + s + m \cdot P, 0) = \operatorname{maskparity}(z + 1 + s, 0)\).
\end{theorem}
\begin{proof}
Assume \(hP\), \(hper\), \(m\), \(z\), \(S\), \(s\).
The goal follows from \(\operatorname{periodicmaskparityzeroagreeatzerodiagsuccofdoublesandsinglelhrowspatialagree}(T, P, hP, hper, m, z, S, s)\).
\end{proof}
\begin{theorem}
\label{ll:paper:not-centerperiodic-of-single-lh-row-spatial-agree-at-zero-diag-succ-doubles}
For every natural number \(T\) and natural number \(P\) and natural number \(m\) and natural number \(z\) and natural number \(S\), if \(0 < P\), then if \(0 < m\), then if \(T \le z + 1\), then if \(\operatorname{zerodiagonal}(z)\), then if \(\operatorname{doublesat}(m \cdot P, z)\), then if \(\operatorname{settledby}(z + 1, z + 1)\), then if \(\operatorname{settledby}(z + 1 - m \cdot P, z + 1)\), then if \(3 \cdot (m \cdot P) \le S + 1\), then if for every natural number \(i\), if \(i \le S\), then \(\operatorname{settledby}(z + 1 - (1 + i), z + 1)\), then if for every natural number \(i\), if \(i \le S\), then \(\operatorname{lhcenter}(z + 1 + m \cdot P, 1 + i) = \operatorname{lhcenter}(z + 1 + m \cdot P, 1 + i + m \cdot P)\), then not \(\operatorname{centerperiodic}(T, P)\).
\end{theorem}
\begin{proof}
The goal follows from \(\operatorname{notcenterperiodicofsinglelhrowspatialagreeatzerodiagsuccdoubles}(T, P, m, z, S)\).
\end{proof}
\begin{theorem}
\label{ll:paper:lhbandagree-one-at-zero-diag-of-doubles-prefix-and-single-lh-row-tail}
For every natural number \(Q\) and natural number \(z\) and natural number \(p\) and natural number \(S\), if \(\operatorname{doublesat}(Q, z)\), then if for every natural number \(i\), if \(i < p\), then \(\operatorname{lhcenter}(z + (p + 1) + Q, 1 + i) = \operatorname{lhcenter}(z + (p + 1), 1 + i)\), then if for every natural number \(i\), if \(p \le i\), then if \(i \le S\), then \(\operatorname{settledby}(z + (p + 1) - (1 + i), z + (p + 1))\), then if for every natural number \(i\), if \(p \le i\), then if \(i \le S\), then \(\operatorname{lhcenter}(z + (p + 1) + Q, 1 + i) = \operatorname{lhcenter}(z + (p + 1) + Q, 1 + i + Q)\), then \(\operatorname{lhbandagree}(z + (p + 1), Q, 1, S)\).
\end{theorem}
\begin{proof}
The goal follows from \(\operatorname{lhbandagreeoneatzerodiagofdoublesprefixandsinglelhrowtail}(Q, z, p, S)\).
\end{proof}
\begin{theorem}
\label{ll:paper:periodic-lhcolagree-one-at-zero-diag-of-doubles-prefix-and-single-lh-row-tail}
For every natural number \(T\) and natural number \(P\), if \(0 < P\), then if \(\operatorname{centerperiodic}(T, P)\), then for every natural number \(m\) and natural number \(z\) and natural number \(p\) and natural number \(S\), if \(T \le z + (p + 1)\), then if \(\operatorname{doublesat}(m \cdot P, z)\), then if for every natural number \(i\), if \(i < p\), then \(\operatorname{lhcenter}(z + (p + 1) + m \cdot P, 1 + i) = \operatorname{lhcenter}(z + (p + 1), 1 + i)\), then if for every natural number \(i\), if \(p \le i\), then if \(i \le S\), then \(\operatorname{settledby}(z + (p + 1) - (1 + i), z + (p + 1))\), then if for every natural number \(i\), if \(p \le i\), then if \(i \le S\), then \(\operatorname{lhcenter}(z + (p + 1) + m \cdot P, 1 + i) = \operatorname{lhcenter}(z + (p + 1) + m \cdot P, 1 + i + m \cdot P)\), then \(\operatorname{lhcolagree}(z + (p + 1), m \cdot P, 1, S)\).
\end{theorem}
\begin{proof}
Assume \(hP\), \(hper\), \(m\), \(z\), \(p\), \(S\).
The goal follows from \(\operatorname{periodiclhcolagreeoneatzerodiagofdoublesprefixandsinglelhrowtail}(T, P, hP, hper, m, z, p, S)\).
\end{proof}
\begin{theorem}
\label{ll:paper:not-centerperiodic-of-doubles-prefix-and-single-lh-row-tail-at-zero-diag}
For every natural number \(T\) and natural number \(P\) and natural number \(m\) and natural number \(z\) and natural number \(p\) and natural number \(S\), if \(0 < P\), then if \(0 < m\), then if \(T \le z + (p + 1)\), then if \(\operatorname{zerodiagonal}(z)\), then if \(\operatorname{doublesat}(m \cdot P, z)\), then if \(\operatorname{settledby}(z + 1, z + (p + 1))\), then if \(\operatorname{settledby}(z + 1 - m \cdot P, z + (p + 1))\), then if \(p + 3 \cdot (m \cdot P) \le S + 1\), then if for every natural number \(i\), if \(i < p\), then \(\operatorname{lhcenter}(z + (p + 1) + m \cdot P, 1 + i) = \operatorname{lhcenter}(z + (p + 1), 1 + i)\), then if for every natural number \(i\), if \(p \le i\), then if \(i \le S\), then \(\operatorname{settledby}(z + (p + 1) - (1 + i), z + (p + 1))\), then if for every natural number \(i\), if \(p \le i\), then if \(i \le S\), then \(\operatorname{lhcenter}(z + (p + 1) + m \cdot P, 1 + i) = \operatorname{lhcenter}(z + (p + 1) + m \cdot P, 1 + i + m \cdot P)\), then not \(\operatorname{centerperiodic}(T, P)\).
\end{theorem}
\begin{proof}
The goal follows from \(\operatorname{notcenterperiodicofdoublesprefixandsinglelhrowtailatzerodiag}(T, P, m, z, p, S)\).
\end{proof}
\begin{theorem}
\label{ll:paper:rhcenter-dyadic-left-shift-pair-eq-lhcenter-pair}
For every natural number \(t\) and natural number \(Q\) and natural number \(K\) and natural number \(v\) and natural number \(d\), if \(1 \le v\), then if \(t + Q + (d + Q) \le K\), then if \(\operatorname{rhcenter}(t + Q + v \cdot \operatorname{twopow}(K), v \cdot \operatorname{twopow}(K) - d) = \operatorname{rhcenter}(t + Q + v \cdot \operatorname{twopow}(K), v \cdot \operatorname{twopow}(K) - (d + Q))\), then \(\operatorname{lhcenter}(t + Q, d) = \operatorname{lhcenter}(t + Q, d + Q)\).
\end{theorem}
\begin{proof}
The goal follows from \(\operatorname{rhcenterdyadicleftshiftpaireqlhcenterpair}(t, Q, K, v, d)\).
\end{proof}
\begin{theorem}
\label{ll:paper:lhbandagree-of-doubles-le-and-single-rh-row-dyadic-left-spatial-agree}
For every natural number \(Q\) and natural number \(z\) and natural number \(t\) and natural number \(p0\) and natural number \(L\) and natural number \(K\) and natural number \(v\), if \(\operatorname{doublesat}(Q, z)\), then if \(t - p0 \le z\), then if \(1 \le v\), then if \(t + Q + (p0 + L + Q) \le K\), then if for every natural number \(i\), if \(i \le L\), then \(\operatorname{settledby}(t - (p0 + i), t)\), then if for every natural number \(i\), if \(i \le L\), then \(\operatorname{rhcenter}(t + Q + v \cdot \operatorname{twopow}(K), v \cdot \operatorname{twopow}(K) - (p0 + i)) = \operatorname{rhcenter}(t + Q + v \cdot \operatorname{twopow}(K), v \cdot \operatorname{twopow}(K) - (p0 + i + Q))\), then \(\operatorname{lhbandagree}(t, Q, p0, L)\).
\end{theorem}
\begin{proof}
The goal follows from \(\operatorname{lhbandagreeofdoublesleandsinglerhrowdyadicleftspatialagree}(Q, z, t, p0, L, K, v)\).
\end{proof}
\begin{theorem}
\label{ll:paper:lhbandagree-one-at-zero-diag-succ-of-doubles-and-single-rh-row-dyadic-left-spatial-agree}
For every natural number \(Q\) and natural number \(z\) and natural number \(S\) and natural number \(K\) and natural number \(v\), if \(\operatorname{doublesat}(Q, z)\), then if \(1 \le v\), then if \(z + 1 + Q + (1 + S + Q) \le K\), then if for every natural number \(i\), if \(i \le S\), then \(\operatorname{settledby}(z + 1 - (1 + i), z + 1)\), then if for every natural number \(i\), if \(i \le S\), then \(\operatorname{rhcenter}(z + 1 + Q + v \cdot \operatorname{twopow}(K), v \cdot \operatorname{twopow}(K) - (1 + i)) = \operatorname{rhcenter}(z + 1 + Q + v \cdot \operatorname{twopow}(K), v \cdot \operatorname{twopow}(K) - (1 + i + Q))\), then \(\operatorname{lhbandagree}(z + 1, Q, 1, S)\).
\end{theorem}
\begin{proof}
The goal follows from \(\operatorname{lhbandagreeoneatzerodiagsuccofdoublesandsinglerhrowdyadicleftspatialagree}(Q, z, S, K, v)\).
\end{proof}
\begin{theorem}
\label{ll:paper:periodic-lhcolagree-one-at-zero-diag-succ-of-doubles-and-single-rh-row-dyadic-left-spatial-agree}
For every natural number \(T\) and natural number \(P\), if \(0 < P\), then if \(\operatorname{centerperiodic}(T, P)\), then for every natural number \(m\) and natural number \(z\) and natural number \(S\) and natural number \(K\) and natural number \(v\), if \(T \le z + 1\), then if \(\operatorname{doublesat}(m \cdot P, z)\), then if \(1 \le v\), then if \(z + 1 + m \cdot P + (1 + S + m \cdot P) \le K\), then if for every natural number \(i\), if \(i \le S\), then \(\operatorname{settledby}(z + 1 - (1 + i), z + 1)\), then if for every natural number \(i\), if \(i \le S\), then \(\operatorname{rhcenter}(z + 1 + m \cdot P + v \cdot \operatorname{twopow}(K), v \cdot \operatorname{twopow}(K) - (1 + i)) = \operatorname{rhcenter}(z + 1 + m \cdot P + v \cdot \operatorname{twopow}(K), v \cdot \operatorname{twopow}(K) - (1 + i + m \cdot P))\), then \(\operatorname{lhcolagree}(z + 1, m \cdot P, 1, S)\).
\end{theorem}
\begin{proof}
Assume \(hP\), \(hper\), \(m\), \(z\), \(S\), \(K\), \(v\).
The goal follows from \(\operatorname{periodiclhcolagreeoneatzerodiagsuccofdoublesandsinglerhrowdyadicleftspatialagree}(T, P, hP, hper, m, z, S, K, v)\).
\end{proof}
\begin{theorem}
\label{ll:paper:periodic-maskparity-zero-agree-at-zero-diag-succ-of-doubles-and-single-rh-row-dyadic-left-spatial-agree}
For every natural number \(T\) and natural number \(P\), if \(0 < P\), then if \(\operatorname{centerperiodic}(T, P)\), then for every natural number \(m\) and natural number \(z\) and natural number \(S\) and natural number \(K\) and natural number \(v\) and natural number \(s\), if \(T \le z + 1\), then if \(\operatorname{doublesat}(m \cdot P, z)\), then if \(1 \le v\), then if \(z + 1 + m \cdot P + (1 + S + m \cdot P) \le K\), then if for every natural number \(i\), if \(i \le S\), then \(\operatorname{settledby}(z + 1 - (1 + i), z + 1)\), then if for every natural number \(i\), if \(i \le S\), then \(\operatorname{rhcenter}(z + 1 + m \cdot P + v \cdot \operatorname{twopow}(K), v \cdot \operatorname{twopow}(K) - (1 + i)) = \operatorname{rhcenter}(z + 1 + m \cdot P + v \cdot \operatorname{twopow}(K), v \cdot \operatorname{twopow}(K) - (1 + i + m \cdot P))\), then if \(s \le S\), then \(\operatorname{maskparity}(z + 1 + s + m \cdot P, 0) = \operatorname{maskparity}(z + 1 + s, 0)\).
\end{theorem}
\begin{proof}
Assume \(hP\), \(hper\), \(m\), \(z\), \(S\), \(K\), \(v\), \(s\).
The goal follows from \(\operatorname{periodicmaskparityzeroagreeatzerodiagsuccofdoublesandsinglerhrowdyadicleftspatialagree}(T, P, hP, hper, m, z, S, K, v, s)\).
\end{proof}
\begin{theorem}
\label{ll:paper:not-centerperiodic-of-single-rh-row-dyadic-left-spatial-agree-at-zero-diag-succ-doubles}
For every natural number \(T\) and natural number \(P\) and natural number \(m\) and natural number \(z\) and natural number \(S\) and natural number \(K\) and natural number \(v\), if \(0 < P\), then if \(0 < m\), then if \(T \le z + 1\), then if \(\operatorname{zerodiagonal}(z)\), then if \(\operatorname{doublesat}(m \cdot P, z)\), then if \(\operatorname{settledby}(z + 1, z + 1)\), then if \(\operatorname{settledby}(z + 1 - m \cdot P, z + 1)\), then if \(3 \cdot (m \cdot P) \le S + 1\), then if \(1 \le v\), then if \(z + 1 + m \cdot P + (1 + S + m \cdot P) \le K\), then if for every natural number \(i\), if \(i \le S\), then \(\operatorname{settledby}(z + 1 - (1 + i), z + 1)\), then if for every natural number \(i\), if \(i \le S\), then \(\operatorname{rhcenter}(z + 1 + m \cdot P + v \cdot \operatorname{twopow}(K), v \cdot \operatorname{twopow}(K) - (1 + i)) = \operatorname{rhcenter}(z + 1 + m \cdot P + v \cdot \operatorname{twopow}(K), v \cdot \operatorname{twopow}(K) - (1 + i + m \cdot P))\), then not \(\operatorname{centerperiodic}(T, P)\).
\end{theorem}
\begin{proof}
The goal follows from \(\operatorname{notcenterperiodicofsinglerhrowdyadicleftspatialagreeatzerodiagsuccdoubles}(T, P, m, z, S, K, v)\).
\end{proof}
\begin{theorem}
\label{ll:paper:rdrv-dyadic-left-shift-pair-eq-lhcenter-pair}
For every natural number \(td\) and natural number \(t\) and natural number \(Q\) and natural number \(K\) and natural number \(v\) and natural number \(d\), if \(1 \le v\), then if \(t + Q + (d + Q) \le K\), then if \(\operatorname{rdrv}(td, v \cdot \operatorname{twopow}(K), t + Q, v \cdot \operatorname{twopow}(K) - d) = \operatorname{rdrv}(td, v \cdot \operatorname{twopow}(K), t + Q, v \cdot \operatorname{twopow}(K) - (d + Q))\), then \(\operatorname{lhcenter}(t + Q, d) = \operatorname{lhcenter}(t + Q, d + Q)\).
\end{theorem}
\begin{proof}
The goal follows from \(\operatorname{rdrvdyadicleftshiftpaireqlhcenterpair}(td, t, Q, K, v, d)\).
\end{proof}
\begin{theorem}
\label{ll:paper:lhbandagree-of-doubles-le-and-single-rdrv-row-dyadic-left-spatial-agree}
For every natural number \(td\) and natural number \(Q\) and natural number \(z\) and natural number \(t\) and natural number \(p0\) and natural number \(L\) and natural number \(K\) and natural number \(v\), if \(\operatorname{doublesat}(Q, z)\), then if \(t - p0 \le z\), then if \(1 \le v\), then if \(t + Q + (p0 + L + Q) \le K\), then if for every natural number \(i\), if \(i \le L\), then \(\operatorname{settledby}(t - (p0 + i), t)\), then if for every natural number \(i\), if \(i \le L\), then \(\operatorname{rdrv}(td, v \cdot \operatorname{twopow}(K), t + Q, v \cdot \operatorname{twopow}(K) - (p0 + i)) = \operatorname{rdrv}(td, v \cdot \operatorname{twopow}(K), t + Q, v \cdot \operatorname{twopow}(K) - (p0 + i + Q))\), then \(\operatorname{lhbandagree}(t, Q, p0, L)\).
\end{theorem}
\begin{proof}
The goal follows from \(\operatorname{lhbandagreeofdoublesleandsinglerdrvrowdyadicleftspatialagree}(td, Q, z, t, p0, L, K, v)\).
\end{proof}
\begin{theorem}
\label{ll:paper:lhbandagree-one-at-zero-diag-succ-of-doubles-and-single-rdrv-row-dyadic-left-spatial-agree}
For every natural number \(td\) and natural number \(Q\) and natural number \(z\) and natural number \(S\) and natural number \(K\) and natural number \(v\), if \(\operatorname{doublesat}(Q, z)\), then if \(1 \le v\), then if \(z + 1 + Q + (1 + S + Q) \le K\), then if for every natural number \(i\), if \(i \le S\), then \(\operatorname{settledby}(z + 1 - (1 + i), z + 1)\), then if for every natural number \(i\), if \(i \le S\), then \(\operatorname{rdrv}(td, v \cdot \operatorname{twopow}(K), z + 1 + Q, v \cdot \operatorname{twopow}(K) - (1 + i)) = \operatorname{rdrv}(td, v \cdot \operatorname{twopow}(K), z + 1 + Q, v \cdot \operatorname{twopow}(K) - (1 + i + Q))\), then \(\operatorname{lhbandagree}(z + 1, Q, 1, S)\).
\end{theorem}
\begin{proof}
The goal follows from \(\operatorname{lhbandagreeoneatzerodiagsuccofdoublesandsinglerdrvrowdyadicleftspatialagree}(td, Q, z, S, K, v)\).
\end{proof}
\begin{theorem}
\label{ll:paper:periodic-lhcolagree-one-at-zero-diag-succ-of-doubles-and-single-rdrv-row-dyadic-left-spatial-agree}
For every natural number \(T\) and natural number \(P\), if \(0 < P\), then if \(\operatorname{centerperiodic}(T, P)\), then for every natural number \(td\) and natural number \(m\) and natural number \(z\) and natural number \(S\) and natural number \(K\) and natural number \(v\), if \(T \le z + 1\), then if \(\operatorname{doublesat}(m \cdot P, z)\), then if \(1 \le v\), then if \(z + 1 + m \cdot P + (1 + S + m \cdot P) \le K\), then if for every natural number \(i\), if \(i \le S\), then \(\operatorname{settledby}(z + 1 - (1 + i), z + 1)\), then if for every natural number \(i\), if \(i \le S\), then \(\operatorname{rdrv}(td, v \cdot \operatorname{twopow}(K), z + 1 + m \cdot P, v \cdot \operatorname{twopow}(K) - (1 + i)) = \operatorname{rdrv}(td, v \cdot \operatorname{twopow}(K), z + 1 + m \cdot P, v \cdot \operatorname{twopow}(K) - (1 + i + m \cdot P))\), then \(\operatorname{lhcolagree}(z + 1, m \cdot P, 1, S)\).
\end{theorem}
\begin{proof}
Assume \(hP\), \(hper\), \(td\), \(m\), \(z\), \(S\), \(K\), \(v\).
The goal follows from \(\operatorname{periodiclhcolagreeoneatzerodiagsuccofdoublesandsinglerdrvrowdyadicleftspatialagree}(T, P, hP, hper, td, m, z, S, K, v)\).
\end{proof}
\begin{theorem}
\label{ll:paper:periodic-maskparity-zero-agree-at-zero-diag-succ-of-doubles-and-single-rdrv-row-dyadic-left-spatial-agree}
For every natural number \(T\) and natural number \(P\), if \(0 < P\), then if \(\operatorname{centerperiodic}(T, P)\), then for every natural number \(td\) and natural number \(m\) and natural number \(z\) and natural number \(S\) and natural number \(K\) and natural number \(v\) and natural number \(s\), if \(T \le z + 1\), then if \(\operatorname{doublesat}(m \cdot P, z)\), then if \(1 \le v\), then if \(z + 1 + m \cdot P + (1 + S + m \cdot P) \le K\), then if for every natural number \(i\), if \(i \le S\), then \(\operatorname{settledby}(z + 1 - (1 + i), z + 1)\), then if for every natural number \(i\), if \(i \le S\), then \(\operatorname{rdrv}(td, v \cdot \operatorname{twopow}(K), z + 1 + m \cdot P, v \cdot \operatorname{twopow}(K) - (1 + i)) = \operatorname{rdrv}(td, v \cdot \operatorname{twopow}(K), z + 1 + m \cdot P, v \cdot \operatorname{twopow}(K) - (1 + i + m \cdot P))\), then if \(s \le S\), then \(\operatorname{maskparity}(z + 1 + s + m \cdot P, 0) = \operatorname{maskparity}(z + 1 + s, 0)\).
\end{theorem}
\begin{proof}
Assume \(hP\), \(hper\), \(td\), \(m\), \(z\), \(S\), \(K\), \(v\), \(s\).
The goal follows from \(\operatorname{periodicmaskparityzeroagreeatzerodiagsuccofdoublesandsinglerdrvrowdyadicleftspatialagree}(T, P, hP, hper, td, m, z, S, K, v, s)\).
\end{proof}
\begin{theorem}
\label{ll:paper:not-centerperiodic-of-single-rdrv-row-dyadic-left-spatial-agree-at-zero-diag-succ-doubles}
For every natural number \(T\) and natural number \(P\) and natural number \(td\) and natural number \(m\) and natural number \(z\) and natural number \(S\) and natural number \(K\) and natural number \(v\), if \(0 < P\), then if \(0 < m\), then if \(T \le z + 1\), then if \(\operatorname{zerodiagonal}(z)\), then if \(\operatorname{doublesat}(m \cdot P, z)\), then if \(\operatorname{settledby}(z + 1, z + 1)\), then if \(\operatorname{settledby}(z + 1 - m \cdot P, z + 1)\), then if \(3 \cdot (m \cdot P) \le S + 1\), then if \(1 \le v\), then if \(z + 1 + m \cdot P + (1 + S + m \cdot P) \le K\), then if for every natural number \(i\), if \(i \le S\), then \(\operatorname{settledby}(z + 1 - (1 + i), z + 1)\), then if for every natural number \(i\), if \(i \le S\), then \(\operatorname{rdrv}(td, v \cdot \operatorname{twopow}(K), z + 1 + m \cdot P, v \cdot \operatorname{twopow}(K) - (1 + i)) = \operatorname{rdrv}(td, v \cdot \operatorname{twopow}(K), z + 1 + m \cdot P, v \cdot \operatorname{twopow}(K) - (1 + i + m \cdot P))\), then not \(\operatorname{centerperiodic}(T, P)\).
\end{theorem}
\begin{proof}
The goal follows from \(\operatorname{notcenterperiodicofsinglerdrvrowdyadicleftspatialagreeatzerodiagsuccdoubles}(T, P, td, m, z, S, K, v)\).
\end{proof}
\begin{theorem}
\label{ll:paper:lhcenter-late-scale-diag-shift-periodic-of-doubles-le}
For every natural number \(Q\) and natural number \(z\) and natural number \(t\) and natural number \(p\) and natural number \(k\), if \(\operatorname{doublesat}(Q, z)\), then if \(t - p \le z\), then if \(\operatorname{settledby}(t - p, t + \operatorname{twopow}(k))\), then \(\operatorname{lhcenter}(t + Q + \operatorname{twopow}(k), p + Q + \operatorname{twopow}(k)) = \operatorname{lhcenter}(t + \operatorname{twopow}(k), p + \operatorname{twopow}(k))\).
\end{theorem}
\begin{proof}
The goal follows from \(\operatorname{lhcenterlatescalediagshiftperiodicofdoublesle}(Q, z, t, p, k)\).
\end{proof}
\begin{theorem}
\label{ll:paper:lhcenter-late-scale-diag-shift-mul-twopow-periodic-le}
For every natural number \(e\) and natural number \(w\) and natural number \(t\) and natural number \(p\) and natural number \(k\), if \(t - p \le e\), then if \(\operatorname{settledby}(t - p, t + \operatorname{twopow}(k))\), then \(\operatorname{lhcenter}(t + w \cdot \operatorname{twopow}(e) + \operatorname{twopow}(k), p + w \cdot \operatorname{twopow}(e) + \operatorname{twopow}(k)) = \operatorname{lhcenter}(t + \operatorname{twopow}(k), p + \operatorname{twopow}(k))\).
\end{theorem}
\begin{proof}
The goal follows from \(\operatorname{lhcenterlatescalediagshiftmultwopowperiodicle}(e, w, t, p, k)\).
\end{proof}
\begin{theorem}
\label{ll:paper:lhcenter-late-scale-eq-of-doubles-le-and-spatial-agree}
For every natural number \(Q\) and natural number \(z\) and natural number \(t\) and natural number \(p\) and natural number \(k\), if \(\operatorname{doublesat}(Q, z)\), then if \(t - p \le z\), then if \(\operatorname{settledby}(t - p, t + \operatorname{twopow}(k))\), then if \(\operatorname{lhcenter}(t + Q + \operatorname{twopow}(k), p + \operatorname{twopow}(k)) = \operatorname{lhcenter}(t + Q + \operatorname{twopow}(k), p + Q + \operatorname{twopow}(k))\), then \(\operatorname{lhcenter}(t + Q + \operatorname{twopow}(k), p + \operatorname{twopow}(k)) = \operatorname{lhcenter}(t + \operatorname{twopow}(k), p + \operatorname{twopow}(k))\).
\end{theorem}
\begin{proof}
The goal follows from \(\operatorname{lhcenterlatescaleeqofdoublesleandspatialagree}(Q, z, t, p, k)\).
\end{proof}
\begin{theorem}
\label{ll:paper:lhbandagree-late-scale-of-doubles-le-and-single-lh-row-spatial-agree}
For every natural number \(Q\) and natural number \(z\) and natural number \(t\) and natural number \(p0\) and natural number \(L\) and natural number \(k\), if \(\operatorname{doublesat}(Q, z)\), then if \(t - p0 \le z\), then if for every natural number \(i\), if \(i \le L\), then \(\operatorname{settledby}(t - (p0 + i), t + \operatorname{twopow}(k))\), then if for every natural number \(i\), if \(i \le L\), then \(\operatorname{lhcenter}(t + Q + \operatorname{twopow}(k), p0 + \operatorname{twopow}(k) + i) = \operatorname{lhcenter}(t + Q + \operatorname{twopow}(k), p0 + \operatorname{twopow}(k) + i + Q)\), then \(\operatorname{lhbandagree}(t + \operatorname{twopow}(k), Q, p0 + \operatorname{twopow}(k), L)\).
\end{theorem}
\begin{proof}
The goal follows from \(\operatorname{lhbandagreelatescaleofdoublesleandsinglelhrowspatialagree}(Q, z, t, p0, L, k)\).
\end{proof}
\begin{theorem}
\label{ll:paper:lhbandagree-late-scale-at-zero-diag-succ-of-doubles-and-single-lh-row-spatial-agree}
For every natural number \(Q\) and natural number \(z\) and natural number \(S\) and natural number \(k\), if \(\operatorname{doublesat}(Q, z)\), then if for every natural number \(i\), if \(i \le S\), then \(\operatorname{settledby}(z + 1 - (1 + i), z + 1 + \operatorname{twopow}(k))\), then if for every natural number \(i\), if \(i \le S\), then \(\operatorname{lhcenter}(z + 1 + Q + \operatorname{twopow}(k), 1 + \operatorname{twopow}(k) + i) = \operatorname{lhcenter}(z + 1 + Q + \operatorname{twopow}(k), 1 + \operatorname{twopow}(k) + i + Q)\), then \(\operatorname{lhbandagree}(z + 1 + \operatorname{twopow}(k), Q, 1 + \operatorname{twopow}(k), S)\).
\end{theorem}
\begin{proof}
The goal follows from \(\operatorname{lhbandagreelatescaleatzerodiagsuccofdoublesandsinglelhrowspatialagree}(Q, z, S, k)\).
\end{proof}
\begin{theorem}
\label{ll:paper:lhbandagree-append-of-split}
For every natural number \(t\) and natural number \(Q\) and natural number \(p0\) and natural number \(A\) and natural number \(B\), if \(\operatorname{lhbandagree}(t, Q, p0, A)\), then if \(\operatorname{lhbandagree}(t, Q, p0 + A, B)\), then \(\operatorname{lhbandagree}(t, Q, p0, A + B)\).
\end{theorem}
\begin{proof}
The goal follows from \(\operatorname{lhbandagreeappendofsplit}(t, Q, p0, A, B)\).
\end{proof}
\begin{theorem}
\label{ll:paper:lhbandagree-one-of-prefix-and-shifted-band}
For every natural number \(t\) and natural number \(Q\) and natural number \(P0\) and natural number \(L\), if for every natural number \(i\), if \(i < P0\), then \(\operatorname{lhcenter}(t + Q, 1 + i) = \operatorname{lhcenter}(t, 1 + i)\), then if \(\operatorname{lhbandagree}(t, Q, 1 + P0, L)\), then \(\operatorname{lhbandagree}(t, Q, 1, P0 + L)\).
\end{theorem}
\begin{proof}
The goal follows from \(\operatorname{lhbandagreeoneofprefixandshiftedband}(t, Q, P0, L)\).
\end{proof}
\begin{theorem}
\label{ll:paper:lhbandagree-one-late-scale-at-zero-diag-succ-of-doubles-prefix-and-single-lh-row-tail}
For every natural number \(Q\) and natural number \(z\) and natural number \(L\) and natural number \(k\), if \(\operatorname{doublesat}(Q, z)\), then if for every natural number \(i\), if \(i < \operatorname{twopow}(k)\), then \(\operatorname{lhcenter}(z + 1 + \operatorname{twopow}(k) + Q, 1 + i) = \operatorname{lhcenter}(z + 1 + \operatorname{twopow}(k), 1 + i)\), then if for every natural number \(i\), if \(i \le L\), then \(\operatorname{settledby}(z + 1 - (1 + i), z + 1 + \operatorname{twopow}(k))\), then if for every natural number \(i\), if \(i \le L\), then \(\operatorname{lhcenter}(z + 1 + Q + \operatorname{twopow}(k), 1 + \operatorname{twopow}(k) + i) = \operatorname{lhcenter}(z + 1 + Q + \operatorname{twopow}(k), 1 + \operatorname{twopow}(k) + i + Q)\), then \(\operatorname{lhbandagree}(z + 1 + \operatorname{twopow}(k), Q, 1, \operatorname{twopow}(k) + L)\).
\end{theorem}
\begin{proof}
The goal follows from \(\operatorname{lhbandagreeonelatescaleatzerodiagsuccofdoublesprefixandsinglelhrowtail}(Q, z, L, k)\).
\end{proof}
\begin{theorem}
\label{ll:paper:periodic-lhcolagree-one-late-scale-at-zero-diag-succ-of-doubles-prefix-and-single-lh-row-tail}
For every natural number \(T\) and natural number \(P\), if \(0 < P\), then if \(\operatorname{centerperiodic}(T, P)\), then for every natural number \(m\) and natural number \(z\) and natural number \(L\) and natural number \(k\), if \(T \le z + 1 + \operatorname{twopow}(k)\), then if \(\operatorname{doublesat}(m \cdot P, z)\), then if for every natural number \(i\), if \(i < \operatorname{twopow}(k)\), then \(\operatorname{lhcenter}(z + 1 + \operatorname{twopow}(k) + m \cdot P, 1 + i) = \operatorname{lhcenter}(z + 1 + \operatorname{twopow}(k), 1 + i)\), then if for every natural number \(i\), if \(i \le L\), then \(\operatorname{settledby}(z + 1 - (1 + i), z + 1 + \operatorname{twopow}(k))\), then if for every natural number \(i\), if \(i \le L\), then \(\operatorname{lhcenter}(z + 1 + m \cdot P + \operatorname{twopow}(k), 1 + \operatorname{twopow}(k) + i) = \operatorname{lhcenter}(z + 1 + m \cdot P + \operatorname{twopow}(k), 1 + \operatorname{twopow}(k) + i + m \cdot P)\), then \(\operatorname{lhcolagree}(z + 1 + \operatorname{twopow}(k), m \cdot P, 1, \operatorname{twopow}(k) + L)\).
\end{theorem}
\begin{proof}
Assume \(hP\), \(hper\), \(m\), \(z\), \(L\), \(k\).
The goal follows from \(\operatorname{periodiclhcolagreeonelatescaleatzerodiagsuccofdoublesprefixandsinglelhrowtail}(T, P, hP, hper, m, z, L, k)\).
\end{proof}
\begin{theorem}
\label{ll:paper:not-centerperiodic-of-late-scale-doubles-prefix-and-single-lh-row-tail}
For every natural number \(T\) and natural number \(P\) and natural number \(m\) and natural number \(z\) and natural number \(L\) and natural number \(k\), if \(0 < P\), then if \(0 < m\), then if \(T \le z + 1 + \operatorname{twopow}(k)\), then if \(\operatorname{zerodiagonal}(z)\), then if \(\operatorname{doublesat}(m \cdot P, z)\), then if \(\operatorname{settledby}(z + 1, z + 1 + \operatorname{twopow}(k))\), then if \(\operatorname{settledby}(z + 1 - m \cdot P, z + 1 + \operatorname{twopow}(k))\), then if \(3 \cdot (m \cdot P) \le L + 1\), then if for every natural number \(i\), if \(i < \operatorname{twopow}(k)\), then \(\operatorname{lhcenter}(z + 1 + \operatorname{twopow}(k) + m \cdot P, 1 + i) = \operatorname{lhcenter}(z + 1 + \operatorname{twopow}(k), 1 + i)\), then if for every natural number \(i\), if \(i \le L\), then \(\operatorname{settledby}(z + 1 - (1 + i), z + 1 + \operatorname{twopow}(k))\), then if for every natural number \(i\), if \(i \le L\), then \(\operatorname{lhcenter}(z + 1 + m \cdot P + \operatorname{twopow}(k), 1 + \operatorname{twopow}(k) + i) = \operatorname{lhcenter}(z + 1 + m \cdot P + \operatorname{twopow}(k), 1 + \operatorname{twopow}(k) + i + m \cdot P)\), then not \(\operatorname{centerperiodic}(T, P)\).
\end{theorem}
\begin{proof}
The goal follows from \(\operatorname{notcenterperiodicoflatescaledoublesprefixandsinglelhrowtail}(T, P, m, z, L, k)\).
\end{proof}
\begin{theorem}
\label{ll:paper:lhbandagree-late-scale-at-zero-diag-succ-of-doubles-and-single-rh-row-dyadic-left-spatial-agree}
For every natural number \(Q\) and natural number \(z\) and natural number \(L\) and natural number \(k\) and natural number \(K\) and natural number \(v\), if \(\operatorname{doublesat}(Q, z)\), then if \(1 \le v\), then if \(z + 1 + Q + \operatorname{twopow}(k) + (1 + \operatorname{twopow}(k) + L + Q) \le K\), then if for every natural number \(i\), if \(i \le L\), then \(\operatorname{settledby}(z + 1 - (1 + i), z + 1 + \operatorname{twopow}(k))\), then if for every natural number \(i\), if \(i \le L\), then \(\operatorname{rhcenter}(z + 1 + Q + \operatorname{twopow}(k) + v \cdot \operatorname{twopow}(K), v \cdot \operatorname{twopow}(K) - (1 + \operatorname{twopow}(k) + i)) = \operatorname{rhcenter}(z + 1 + Q + \operatorname{twopow}(k) + v \cdot \operatorname{twopow}(K), v \cdot \operatorname{twopow}(K) - (1 + \operatorname{twopow}(k) + i + Q))\), then \(\operatorname{lhbandagree}(z + 1 + \operatorname{twopow}(k), Q, 1 + \operatorname{twopow}(k), L)\).
\end{theorem}
\begin{proof}
The goal follows from \(\operatorname{lhbandagreelatescaleatzerodiagsuccofdoublesandsinglerhrowdyadicleftspatialagree}(Q, z, L, k, K, v)\).
\end{proof}
\begin{theorem}
\label{ll:paper:lhbandagree-one-late-scale-at-zero-diag-succ-of-doubles-prefix-and-single-rh-row-tail}
For every natural number \(Q\) and natural number \(z\) and natural number \(L\) and natural number \(k\) and natural number \(K\) and natural number \(v\), if \(\operatorname{doublesat}(Q, z)\), then if \(1 \le v\), then if \(z + 1 + Q + \operatorname{twopow}(k) + (1 + \operatorname{twopow}(k) + L + Q) \le K\), then if for every natural number \(i\), if \(i < \operatorname{twopow}(k)\), then \(\operatorname{lhcenter}(z + 1 + \operatorname{twopow}(k) + Q, 1 + i) = \operatorname{lhcenter}(z + 1 + \operatorname{twopow}(k), 1 + i)\), then if for every natural number \(i\), if \(i \le L\), then \(\operatorname{settledby}(z + 1 - (1 + i), z + 1 + \operatorname{twopow}(k))\), then if for every natural number \(i\), if \(i \le L\), then \(\operatorname{rhcenter}(z + 1 + Q + \operatorname{twopow}(k) + v \cdot \operatorname{twopow}(K), v \cdot \operatorname{twopow}(K) - (1 + \operatorname{twopow}(k) + i)) = \operatorname{rhcenter}(z + 1 + Q + \operatorname{twopow}(k) + v \cdot \operatorname{twopow}(K), v \cdot \operatorname{twopow}(K) - (1 + \operatorname{twopow}(k) + i + Q))\), then \(\operatorname{lhbandagree}(z + 1 + \operatorname{twopow}(k), Q, 1, \operatorname{twopow}(k) + L)\).
\end{theorem}
\begin{proof}
The goal follows from \(\operatorname{lhbandagreeonelatescaleatzerodiagsuccofdoublesprefixandsinglerhrowtail}(Q, z, L, k, K, v)\).
\end{proof}
\begin{theorem}
\label{ll:paper:not-centerperiodic-of-late-scale-doubles-prefix-and-single-rh-row-tail}
For every natural number \(T\) and natural number \(P\) and natural number \(m\) and natural number \(z\) and natural number \(L\) and natural number \(k\) and natural number \(K\) and natural number \(v\), if \(0 < P\), then if \(0 < m\), then if \(T \le z + 1 + \operatorname{twopow}(k)\), then if \(\operatorname{zerodiagonal}(z)\), then if \(\operatorname{doublesat}(m \cdot P, z)\), then if \(\operatorname{settledby}(z + 1, z + 1 + \operatorname{twopow}(k))\), then if \(\operatorname{settledby}(z + 1 - m \cdot P, z + 1 + \operatorname{twopow}(k))\), then if \(3 \cdot (m \cdot P) \le L + 1\), then if \(1 \le v\), then if \(z + 1 + m \cdot P + \operatorname{twopow}(k) + (1 + \operatorname{twopow}(k) + L + m \cdot P) \le K\), then if for every natural number \(i\), if \(i < \operatorname{twopow}(k)\), then \(\operatorname{lhcenter}(z + 1 + \operatorname{twopow}(k) + m \cdot P, 1 + i) = \operatorname{lhcenter}(z + 1 + \operatorname{twopow}(k), 1 + i)\), then if for every natural number \(i\), if \(i \le L\), then \(\operatorname{settledby}(z + 1 - (1 + i), z + 1 + \operatorname{twopow}(k))\), then if for every natural number \(i\), if \(i \le L\), then \(\operatorname{rhcenter}(z + 1 + m \cdot P + \operatorname{twopow}(k) + v \cdot \operatorname{twopow}(K), v \cdot \operatorname{twopow}(K) - (1 + \operatorname{twopow}(k) + i)) = \operatorname{rhcenter}(z + 1 + m \cdot P + \operatorname{twopow}(k) + v \cdot \operatorname{twopow}(K), v \cdot \operatorname{twopow}(K) - (1 + \operatorname{twopow}(k) + i + m \cdot P))\), then not \(\operatorname{centerperiodic}(T, P)\).
\end{theorem}
\begin{proof}
The goal follows from \(\operatorname{notcenterperiodicoflatescaledoublesprefixandsinglerhrowtail}(T, P, m, z, L, k, K, v)\).
\end{proof}
\begin{theorem}
\label{ll:paper:lhbandagree-one-late-scale-at-zero-diag-succ-of-doubles-prefix-and-single-rdrv-row-tail}
For every natural number \(td\) and natural number \(Q\) and natural number \(z\) and natural number \(L\) and natural number \(k\) and natural number \(K\) and natural number \(v\), if \(\operatorname{doublesat}(Q, z)\), then if \(1 \le v\), then if \(z + 1 + Q + \operatorname{twopow}(k) + (1 + \operatorname{twopow}(k) + L + Q) \le K\), then if for every natural number \(i\), if \(i < \operatorname{twopow}(k)\), then \(\operatorname{lhcenter}(z + 1 + \operatorname{twopow}(k) + Q, 1 + i) = \operatorname{lhcenter}(z + 1 + \operatorname{twopow}(k), 1 + i)\), then if for every natural number \(i\), if \(i \le L\), then \(\operatorname{settledby}(z + 1 - (1 + i), z + 1 + \operatorname{twopow}(k))\), then if for every natural number \(i\), if \(i \le L\), then \(\operatorname{rdrv}(td, v \cdot \operatorname{twopow}(K), z + 1 + Q + \operatorname{twopow}(k), v \cdot \operatorname{twopow}(K) - (1 + \operatorname{twopow}(k) + i)) = \operatorname{rdrv}(td, v \cdot \operatorname{twopow}(K), z + 1 + Q + \operatorname{twopow}(k), v \cdot \operatorname{twopow}(K) - (1 + \operatorname{twopow}(k) + i + Q))\), then \(\operatorname{lhbandagree}(z + 1 + \operatorname{twopow}(k), Q, 1, \operatorname{twopow}(k) + L)\).
\end{theorem}
\begin{proof}
The goal follows from \(\operatorname{lhbandagreeonelatescaleatzerodiagsuccofdoublesprefixandsinglerdrvrowtail}(td, Q, z, L, k, K, v)\).
\end{proof}
\begin{theorem}
\label{ll:paper:not-centerperiodic-of-late-scale-doubles-prefix-and-single-rdrv-row-tail}
For every natural number \(T\) and natural number \(P\) and natural number \(td\) and natural number \(m\) and natural number \(z\) and natural number \(L\) and natural number \(k\) and natural number \(K\) and natural number \(v\), if \(0 < P\), then if \(0 < m\), then if \(T \le z + 1 + \operatorname{twopow}(k)\), then if \(\operatorname{zerodiagonal}(z)\), then if \(\operatorname{doublesat}(m \cdot P, z)\), then if \(\operatorname{settledby}(z + 1, z + 1 + \operatorname{twopow}(k))\), then if \(\operatorname{settledby}(z + 1 - m \cdot P, z + 1 + \operatorname{twopow}(k))\), then if \(3 \cdot (m \cdot P) \le L + 1\), then if \(1 \le v\), then if \(z + 1 + m \cdot P + \operatorname{twopow}(k) + (1 + \operatorname{twopow}(k) + L + m \cdot P) \le K\), then if for every natural number \(i\), if \(i < \operatorname{twopow}(k)\), then \(\operatorname{lhcenter}(z + 1 + \operatorname{twopow}(k) + m \cdot P, 1 + i) = \operatorname{lhcenter}(z + 1 + \operatorname{twopow}(k), 1 + i)\), then if for every natural number \(i\), if \(i \le L\), then \(\operatorname{settledby}(z + 1 - (1 + i), z + 1 + \operatorname{twopow}(k))\), then if for every natural number \(i\), if \(i \le L\), then \(\operatorname{rdrv}(td, v \cdot \operatorname{twopow}(K), z + 1 + m \cdot P + \operatorname{twopow}(k), v \cdot \operatorname{twopow}(K) - (1 + \operatorname{twopow}(k) + i)) = \operatorname{rdrv}(td, v \cdot \operatorname{twopow}(K), z + 1 + m \cdot P + \operatorname{twopow}(k), v \cdot \operatorname{twopow}(K) - (1 + \operatorname{twopow}(k) + i + m \cdot P))\), then not \(\operatorname{centerperiodic}(T, P)\).
\end{theorem}
\begin{proof}
The goal follows from \(\operatorname{notcenterperiodicoflatescaledoublesprefixandsinglerdrvrowtail}(T, P, td, m, z, L, k, K, v)\).
\end{proof}
\begin{theorem}
\label{ll:paper:settledby-of-le-and-settlebound-le}
For every natural number \(D\) and natural number \(d\) and natural number \(t\), if \(d \le D\), then if \(\operatorname{settledupto}(D, t)\), then \(\operatorname{settledby}(d, t)\).
\end{theorem}
\begin{proof}
The goal follows from \(\operatorname{settledbyofleandsettleboundle}(D, d, t)\).
\end{proof}
\begin{theorem}
\label{ll:paper:ldiag-settled-periodic-of-doubles-le-settlebound}
For every natural number \(Q\) and natural number \(z\) and natural number \(d\) and natural number \(t\), if \(\operatorname{doublesat}(Q, z)\), then if \(d \le z\), then if \(\operatorname{settledupto}(z, t)\), then \(\operatorname{ldiag}(d, t + Q) = \operatorname{ldiag}(d, t)\).
\end{theorem}
\begin{proof}
The goal follows from \(\operatorname{ldiagsettledperiodicofdoubleslesettlebound}(Q, z, d, t)\).
\end{proof}
\begin{theorem}
\label{ll:paper:lhcenter-diag-shift-periodic-of-doubles-le-settlebound}
For every natural number \(Q\) and natural number \(z\) and natural number \(t\) and natural number \(p\), if \(\operatorname{doublesat}(Q, z)\), then if \(t - p \le z\), then if \(\operatorname{settledupto}(z, t)\), then \(\operatorname{lhcenter}(t + Q, p + Q) = \operatorname{lhcenter}(t, p)\).
\end{theorem}
\begin{proof}
The goal follows from \(\operatorname{lhcenterdiagshiftperiodicofdoubleslesettlebound}(Q, z, t, p)\).
\end{proof}
\begin{theorem}
\label{ll:paper:lhcenter-diag-shift-mul-twopow-periodic-le-settlebound}
For every natural number \(e\) and natural number \(w\) and natural number \(t\) and natural number \(p\), if \(t - p \le e\), then if \(\operatorname{settledupto}(e, t)\), then \(\operatorname{lhcenter}(t + w \cdot \operatorname{twopow}(e), p + w \cdot \operatorname{twopow}(e)) = \operatorname{lhcenter}(t, p)\).
\end{theorem}
\begin{proof}
The goal follows from \(\operatorname{lhcenterdiagshiftmultwopowperiodiclesettlebound}(e, w, t, p)\).
\end{proof}
\begin{theorem}
\label{ll:paper:lhbandagree-of-doubles-le-settlebound-and-single-lh-row-spatial-agree}
For every natural number \(Q\) and natural number \(z\) and natural number \(t\) and natural number \(p0\) and natural number \(L\), if \(\operatorname{doublesat}(Q, z)\), then if \(t - p0 \le z\), then if \(\operatorname{settledupto}(z, t)\), then if for every natural number \(i\), if \(i \le L\), then \(\operatorname{lhcenter}(t + Q, p0 + i) = \operatorname{lhcenter}(t + Q, p0 + i + Q)\), then \(\operatorname{lhbandagree}(t, Q, p0, L)\).
\end{theorem}
\begin{proof}
The goal follows from \(\operatorname{lhbandagreeofdoubleslesettleboundandsinglelhrowspatialagree}(Q, z, t, p0, L)\).
\end{proof}
\begin{theorem}
\label{ll:paper:lhbandagree-one-at-zero-diag-of-doubles-settlebound-prefix-and-single-lh-row-tail}
For every natural number \(Q\) and natural number \(z\) and natural number \(p\) and natural number \(S\), if \(\operatorname{doublesat}(Q, z)\), then if \(\operatorname{settledupto}(z + 1, z + (p + 1))\), then if for every natural number \(i\), if \(i < p\), then \(\operatorname{lhcenter}(z + (p + 1) + Q, 1 + i) = \operatorname{lhcenter}(z + (p + 1), 1 + i)\), then if for every natural number \(i\), if \(p \le i\), then if \(i \le S\), then \(\operatorname{lhcenter}(z + (p + 1) + Q, 1 + i) = \operatorname{lhcenter}(z + (p + 1) + Q, 1 + i + Q)\), then \(\operatorname{lhbandagree}(z + (p + 1), Q, 1, S)\).
\end{theorem}
\begin{proof}
The goal follows from \(\operatorname{lhbandagreeoneatzerodiagofdoublessettleboundprefixandsinglelhrowtail}(Q, z, p, S)\).
\end{proof}
\begin{theorem}
\label{ll:paper:not-centerperiodic-of-lhbandagree-one-at-zero-diag-doubles-settlebound}
For every natural number \(T\) and natural number \(P\) and natural number \(m\) and natural number \(z\) and natural number \(p\) and natural number \(S\), if \(0 < P\), then if \(0 < m\), then if \(T \le z + (p + 1)\), then if \(\operatorname{zerodiagonal}(z)\), then if \(\operatorname{doublesat}(m \cdot P, z)\), then if \(\operatorname{settledupto}(z + 1, z + (p + 1))\), then if \(p + 3 \cdot (m \cdot P) \le S + 1\), then if \(\operatorname{lhbandagree}(z + (p + 1), m \cdot P, 1, S)\), then not \(\operatorname{centerperiodic}(T, P)\).
\end{theorem}
\begin{proof}
The goal follows from \(\operatorname{notcenterperiodicoflhbandagreeoneatzerodiagdoublessettlebound}(T, P, m, z, p, S)\).
\end{proof}
\begin{theorem}
\label{ll:paper:not-centerperiodic-of-doubles-settlebound-prefix-and-single-lh-row-tail-at-zero-diag}
For every natural number \(T\) and natural number \(P\) and natural number \(m\) and natural number \(z\) and natural number \(p\) and natural number \(S\), if \(0 < P\), then if \(0 < m\), then if \(T \le z + (p + 1)\), then if \(\operatorname{zerodiagonal}(z)\), then if \(\operatorname{doublesat}(m \cdot P, z)\), then if \(\operatorname{settledupto}(z + 1, z + (p + 1))\), then if \(p + 3 \cdot (m \cdot P) \le S + 1\), then if for every natural number \(i\), if \(i < p\), then \(\operatorname{lhcenter}(z + (p + 1) + m \cdot P, 1 + i) = \operatorname{lhcenter}(z + (p + 1), 1 + i)\), then if for every natural number \(i\), if \(p \le i\), then if \(i \le S\), then \(\operatorname{lhcenter}(z + (p + 1) + m \cdot P, 1 + i) = \operatorname{lhcenter}(z + (p + 1) + m \cdot P, 1 + i + m \cdot P)\), then not \(\operatorname{centerperiodic}(T, P)\).
\end{theorem}
\begin{proof}
The goal follows from \(\operatorname{notcenterperiodicofdoublessettleboundprefixandsinglelhrowtailatzerodiag}(T, P, m, z, p, S)\).
\end{proof}
\begin{theorem}
\label{ll:paper:not-centerperiodic-of-doubles-settlebound-prefix-and-single-rh-row-tail-at-zero-diag}
For every natural number \(T\) and natural number \(P\) and natural number \(m\) and natural number \(z\) and natural number \(p\) and natural number \(S\) and natural number \(K\) and natural number \(v\), if \(0 < P\), then if \(0 < m\), then if \(T \le z + (p + 1)\), then if \(\operatorname{zerodiagonal}(z)\), then if \(\operatorname{doublesat}(m \cdot P, z)\), then if \(\operatorname{settledupto}(z + 1, z + (p + 1))\), then if \(p + 3 \cdot (m \cdot P) \le S + 1\), then if \(1 \le v\), then if \(z + (p + 1) + m \cdot P + (1 + S + m \cdot P) \le K\), then if for every natural number \(i\), if \(i < p\), then \(\operatorname{lhcenter}(z + (p + 1) + m \cdot P, 1 + i) = \operatorname{lhcenter}(z + (p + 1), 1 + i)\), then if for every natural number \(i\), if \(p \le i\), then if \(i \le S\), then \(\operatorname{rhcenter}(z + (p + 1) + m \cdot P + v \cdot \operatorname{twopow}(K), v \cdot \operatorname{twopow}(K) - (1 + i)) = \operatorname{rhcenter}(z + (p + 1) + m \cdot P + v \cdot \operatorname{twopow}(K), v \cdot \operatorname{twopow}(K) - (1 + i + m \cdot P))\), then not \(\operatorname{centerperiodic}(T, P)\).
\end{theorem}
\begin{proof}
The goal follows from \(\operatorname{notcenterperiodicofdoublessettleboundprefixandsinglerhrowtailatzerodiag}(T, P, m, z, p, S, K, v)\).
\end{proof}
\begin{theorem}
\label{ll:paper:not-centerperiodic-of-doubles-settlebound-prefix-and-single-rdrv-row-tail-at-zero-diag}
For every natural number \(T\) and natural number \(P\) and natural number \(td\) and natural number \(m\) and natural number \(z\) and natural number \(p\) and natural number \(S\) and natural number \(K\) and natural number \(v\), if \(0 < P\), then if \(0 < m\), then if \(T \le z + (p + 1)\), then if \(\operatorname{zerodiagonal}(z)\), then if \(\operatorname{doublesat}(m \cdot P, z)\), then if \(\operatorname{settledupto}(z + 1, z + (p + 1))\), then if \(p + 3 \cdot (m \cdot P) \le S + 1\), then if \(1 \le v\), then if \(z + (p + 1) + m \cdot P + (1 + S + m \cdot P) \le K\), then if for every natural number \(i\), if \(i < p\), then \(\operatorname{lhcenter}(z + (p + 1) + m \cdot P, 1 + i) = \operatorname{lhcenter}(z + (p + 1), 1 + i)\), then if for every natural number \(i\), if \(p \le i\), then if \(i \le S\), then \(\operatorname{rdrv}(td, v \cdot \operatorname{twopow}(K), z + (p + 1) + m \cdot P, v \cdot \operatorname{twopow}(K) - (1 + i)) = \operatorname{rdrv}(td, v \cdot \operatorname{twopow}(K), z + (p + 1) + m \cdot P, v \cdot \operatorname{twopow}(K) - (1 + i + m \cdot P))\), then not \(\operatorname{centerperiodic}(T, P)\).
\end{theorem}
\begin{proof}
The goal follows from \(\operatorname{notcenterperiodicofdoublessettleboundprefixandsinglerdrvrowtailatzerodiag}(T, P, td, m, z, p, S, K, v)\).
\end{proof}
\begin{theorem}
\label{ll:paper:lhbandagree-one-at-zero-diag-of-two-scale-single-rh-row-prefix-and-tail-settlebound}
For every natural number \(z\) and natural number \(p\) and natural number \(S\) and natural number \(K\) and natural number \(L\) and natural number \(u\) and natural number \(w\), if \(\operatorname{doublesat}(w \cdot \operatorname{twopow}(K), z)\), then if \(\operatorname{settledupto}(z + 1, z + (p + 1))\), then if \(1 \le u\), then if \(z + (p + 1) + p \le K\), then if \(z + (p + 1) + w \cdot \operatorname{twopow}(K) + (1 + S + w \cdot \operatorname{twopow}(K)) \le L\), then if \(K \le L\), then if for every natural number \(i\), if \(i < p\), then \(\operatorname{rhcenter}(z + (p + 1) + u \cdot \operatorname{twopow}(L) + w \cdot \operatorname{twopow}(K), u \cdot \operatorname{twopow}(L) - (1 + i)) = \operatorname{rhcenter}(z + (p + 1) + u \cdot \operatorname{twopow}(L) + w \cdot \operatorname{twopow}(K), u \cdot \operatorname{twopow}(L) + w \cdot \operatorname{twopow}(K) - (1 + i))\), then if for every natural number \(i\), if \(p \le i\), then if \(i \le S\), then \(\operatorname{rhcenter}(z + (p + 1) + u \cdot \operatorname{twopow}(L) + w \cdot \operatorname{twopow}(K), u \cdot \operatorname{twopow}(L) - (1 + i)) = \operatorname{rhcenter}(z + (p + 1) + u \cdot \operatorname{twopow}(L) + w \cdot \operatorname{twopow}(K), u \cdot \operatorname{twopow}(L) - (1 + i + w \cdot \operatorname{twopow}(K)))\), then \(\operatorname{lhbandagree}(z + (p + 1), w \cdot \operatorname{twopow}(K), 1, S)\).
\end{theorem}
\begin{proof}
The goal follows from \(\operatorname{lhbandagreeoneatzerodiagoftwoscalesinglerhrowprefixandtailsettlebound}(z, p, S, K, L, u, w)\).
\end{proof}
\begin{theorem}
\label{ll:paper:not-centerperiodic-of-two-scale-single-rh-row-prefix-and-tail-settlebound}
For every natural number \(T\) and natural number \(P\) and natural number \(c\) and natural number \(z\) and natural number \(p\) and natural number \(S\) and natural number \(K\) and natural number \(L\) and natural number \(u\), if \(0 < P\), then if \(0 < c\), then if \(T \le z + (p + 1)\), then if \(\operatorname{zerodiagonal}(z)\), then if \(\operatorname{doublesat}(c \cdot \operatorname{twopow}(K) \cdot P, z)\), then if \(\operatorname{settledupto}(z + 1, z + (p + 1))\), then if \(p + 3 \cdot (c \cdot \operatorname{twopow}(K) \cdot P) \le S + 1\), then if \(1 \le u\), then if \(z + (p + 1) + p \le K\), then if \(z + (p + 1) + c \cdot \operatorname{twopow}(K) \cdot P + (1 + S + c \cdot \operatorname{twopow}(K) \cdot P) \le L\), then if \(K \le L\), then if for every natural number \(i\), if \(i < p\), then \(\operatorname{rhcenter}(z + (p + 1) + u \cdot \operatorname{twopow}(L) + c \cdot \operatorname{twopow}(K) \cdot P, u \cdot \operatorname{twopow}(L) - (1 + i)) = \operatorname{rhcenter}(z + (p + 1) + u \cdot \operatorname{twopow}(L) + c \cdot \operatorname{twopow}(K) \cdot P, u \cdot \operatorname{twopow}(L) + c \cdot \operatorname{twopow}(K) \cdot P - (1 + i))\), then if for every natural number \(i\), if \(p \le i\), then if \(i \le S\), then \(\operatorname{rhcenter}(z + (p + 1) + u \cdot \operatorname{twopow}(L) + c \cdot \operatorname{twopow}(K) \cdot P, u \cdot \operatorname{twopow}(L) - (1 + i)) = \operatorname{rhcenter}(z + (p + 1) + u \cdot \operatorname{twopow}(L) + c \cdot \operatorname{twopow}(K) \cdot P, u \cdot \operatorname{twopow}(L) - (1 + i + c \cdot \operatorname{twopow}(K) \cdot P))\), then not \(\operatorname{centerperiodic}(T, P)\).
\end{theorem}
\begin{proof}
The goal follows from \(\operatorname{notcenterperiodicoftwoscalesinglerhrowprefixandtailsettlebound}(T, P, c, z, p, S, K, L, u)\).
\end{proof}
\begin{theorem}
\label{ll:paper:lhbandagree-one-at-zero-diag-of-two-scale-single-rdrv-row-prefix-and-tail-settlebound}
For every natural number \(td\) and natural number \(z\) and natural number \(p\) and natural number \(S\) and natural number \(K\) and natural number \(L\) and natural number \(u\) and natural number \(w\), if \(\operatorname{doublesat}(w \cdot \operatorname{twopow}(K), z)\), then if \(\operatorname{settledupto}(z + 1, z + (p + 1))\), then if \(1 \le u\), then if \(z + (p + 1) + p \le K\), then if \(z + (p + 1) + w \cdot \operatorname{twopow}(K) + (1 + S + w \cdot \operatorname{twopow}(K)) \le L\), then if \(K \le L\), then if for every natural number \(i\), if \(i < p\), then \(\operatorname{rdrv}(td, u \cdot \operatorname{twopow}(L), z + (p + 1) + w \cdot \operatorname{twopow}(K), u \cdot \operatorname{twopow}(L) - (1 + i)) = \operatorname{rdrv}(td, u \cdot \operatorname{twopow}(L), z + (p + 1) + w \cdot \operatorname{twopow}(K), u \cdot \operatorname{twopow}(L) + w \cdot \operatorname{twopow}(K) - (1 + i))\), then if for every natural number \(i\), if \(p \le i\), then if \(i \le S\), then \(\operatorname{rdrv}(td, u \cdot \operatorname{twopow}(L), z + (p + 1) + w \cdot \operatorname{twopow}(K), u \cdot \operatorname{twopow}(L) - (1 + i)) = \operatorname{rdrv}(td, u \cdot \operatorname{twopow}(L), z + (p + 1) + w \cdot \operatorname{twopow}(K), u \cdot \operatorname{twopow}(L) - (1 + i + w \cdot \operatorname{twopow}(K)))\), then \(\operatorname{lhbandagree}(z + (p + 1), w \cdot \operatorname{twopow}(K), 1, S)\).
\end{theorem}
\begin{proof}
The goal follows from \(\operatorname{lhbandagreeoneatzerodiagoftwoscalesinglerdrvrowprefixandtailsettlebound}(td, z, p, S, K, L, u, w)\).
\end{proof}
\begin{theorem}
\label{ll:paper:not-centerperiodic-of-two-scale-single-rdrv-row-prefix-and-tail-settlebound}
For every natural number \(T\) and natural number \(P\) and natural number \(td\) and natural number \(c\) and natural number \(z\) and natural number \(p\) and natural number \(S\) and natural number \(K\) and natural number \(L\) and natural number \(u\), if \(0 < P\), then if \(0 < c\), then if \(T \le z + (p + 1)\), then if \(\operatorname{zerodiagonal}(z)\), then if \(\operatorname{doublesat}(c \cdot \operatorname{twopow}(K) \cdot P, z)\), then if \(\operatorname{settledupto}(z + 1, z + (p + 1))\), then if \(p + 3 \cdot (c \cdot \operatorname{twopow}(K) \cdot P) \le S + 1\), then if \(1 \le u\), then if \(z + (p + 1) + p \le K\), then if \(z + (p + 1) + c \cdot \operatorname{twopow}(K) \cdot P + (1 + S + c \cdot \operatorname{twopow}(K) \cdot P) \le L\), then if \(K \le L\), then if for every natural number \(i\), if \(i < p\), then \(\operatorname{rdrv}(td, u \cdot \operatorname{twopow}(L), z + (p + 1) + c \cdot \operatorname{twopow}(K) \cdot P, u \cdot \operatorname{twopow}(L) - (1 + i)) = \operatorname{rdrv}(td, u \cdot \operatorname{twopow}(L), z + (p + 1) + c \cdot \operatorname{twopow}(K) \cdot P, u \cdot \operatorname{twopow}(L) + c \cdot \operatorname{twopow}(K) \cdot P - (1 + i))\), then if for every natural number \(i\), if \(p \le i\), then if \(i \le S\), then \(\operatorname{rdrv}(td, u \cdot \operatorname{twopow}(L), z + (p + 1) + c \cdot \operatorname{twopow}(K) \cdot P, u \cdot \operatorname{twopow}(L) - (1 + i)) = \operatorname{rdrv}(td, u \cdot \operatorname{twopow}(L), z + (p + 1) + c \cdot \operatorname{twopow}(K) \cdot P, u \cdot \operatorname{twopow}(L) - (1 + i + c \cdot \operatorname{twopow}(K) \cdot P))\), then not \(\operatorname{centerperiodic}(T, P)\).
\end{theorem}
\begin{proof}
The goal follows from \(\operatorname{notcenterperiodicoftwoscalesinglerdrvrowprefixandtailsettlebound}(T, P, td, c, z, p, S, K, L, u)\).
\end{proof}
\begin{theorem}
\label{ll:paper:not-centerperiodic-of-late-scale-doubles-settlebound-prefix-and-single-lh-row-tail}
For every natural number \(T\) and natural number \(P\) and natural number \(m\) and natural number \(z\) and natural number \(L\) and natural number \(k\), if \(0 < P\), then if \(0 < m\), then if \(T \le z + 1 + \operatorname{twopow}(k)\), then if \(\operatorname{zerodiagonal}(z)\), then if \(\operatorname{doublesat}(m \cdot P, z)\), then if \(\operatorname{settledupto}(z + 1, z + 1 + \operatorname{twopow}(k))\), then if \(3 \cdot (m \cdot P) \le L + 1\), then if for every natural number \(i\), if \(i < \operatorname{twopow}(k)\), then \(\operatorname{lhcenter}(z + 1 + \operatorname{twopow}(k) + m \cdot P, 1 + i) = \operatorname{lhcenter}(z + 1 + \operatorname{twopow}(k), 1 + i)\), then if for every natural number \(i\), if \(i \le L\), then \(\operatorname{lhcenter}(z + 1 + m \cdot P + \operatorname{twopow}(k), 1 + \operatorname{twopow}(k) + i) = \operatorname{lhcenter}(z + 1 + m \cdot P + \operatorname{twopow}(k), 1 + \operatorname{twopow}(k) + i + m \cdot P)\), then not \(\operatorname{centerperiodic}(T, P)\).
\end{theorem}
\begin{proof}
The goal follows from \(\operatorname{notcenterperiodicoflatescaledoublessettleboundprefixandsinglelhrowtail}(T, P, m, z, L, k)\).
\end{proof}
\begin{theorem}
\label{ll:paper:settledupto-mono}
For every natural number \(D\) and natural number \(Dp\) and natural number \(t\), if \(D \le Dp\), then if \(\operatorname{settledupto}(Dp, t)\), then \(\operatorname{settledupto}(D, t)\).
\end{theorem}
\begin{proof}
The goal follows from \(\operatorname{settleduptomono}(D, Dp, t)\).
\end{proof}
\begin{theorem}
\label{ll:paper:exists-settledupto}
For every natural number \(D\), there exists a natural number \(t\) such that \(\operatorname{settledupto}(D, t)\).
\end{theorem}
\begin{proof}
The goal follows from \(\operatorname{existssettledupto}(D)\).
\end{proof}
\begin{theorem}
\label{ll:paper:not-lhbandagree-of-doubles-single}
For every natural number \(Q\) and natural number \(t\) and natural number \(p\) and natural number \(p0\) and natural number \(L\), if \(1 \le Q\), then if \(\operatorname{zerodiagonal}(t - (p + 1))\), then if \(\operatorname{doublesat}(Q, t - (p + 1))\), then if \(\operatorname{settledby}(t - (p + 1) + 1, t)\), then if \(\operatorname{settledby}(t - (p + Q), t)\), then if \(p0 \le p + Q\), then if \(p + Q + 2 \cdot Q \le p0 + L\), then not \(\operatorname{lhbandagree}(t, Q, p0, L)\).
\end{theorem}
\begin{proof}
The goal follows from \(\operatorname{notlhbandagreeofdoublessingle}(Q, t, p, p0, L)\).
\end{proof}
\begin{theorem}
\label{ll:paper:not-lhbandagree-of-doubles-settledupto}
For every natural number \(Q\) and natural number \(t\) and natural number \(p\) and natural number \(p0\) and natural number \(L\), if \(1 \le Q\), then if \(\operatorname{zerodiagonal}(t - (p + 1))\), then if \(\operatorname{doublesat}(Q, t - (p + 1))\), then if \(\operatorname{settledupto}(t - (p + 1) + 1, t)\), then if \(p0 \le p + Q\), then if \(p + Q + 2 \cdot Q \le p0 + L\), then not \(\operatorname{lhbandagree}(t, Q, p0, L)\).
\end{theorem}
\begin{proof}
The goal follows from \(\operatorname{notlhbandagreeofdoublessettledupto}(Q, t, p, p0, L)\).
\end{proof}
\begin{theorem}
\label{ll:paper:not-lhbandagree-one-at-zero-diag-of-doubles-settledupto}
For every natural number \(Q\) and natural number \(z\) and natural number \(p\) and natural number \(S\), if \(1 \le Q\), then if \(\operatorname{zerodiagonal}(z)\), then if \(\operatorname{doublesat}(Q, z)\), then if \(\operatorname{settledupto}(z + 1, z + (p + 1))\), then if \(p + 3 \cdot Q \le S + 1\), then not \(\operatorname{lhbandagree}(z + (p + 1), Q, 1, S)\).
\end{theorem}
\begin{proof}
The goal follows from \(\operatorname{notlhbandagreeoneatzerodiagofdoublessettledupto}(Q, z, p, S)\).
\end{proof}
\begin{theorem}
\label{ll:paper:exists-lh-row-spatial-disagree-of-doubles-settledupto}
For every natural number \(Q\) and natural number \(z\) and natural number \(p\), if \(1 \le Q\), then if \(\operatorname{zerodiagonal}(z)\), then if \(\operatorname{doublesat}(Q, z)\), then if \(\operatorname{settledupto}(z + 1, z + (p + 1))\), then there exists a natural number \(i\) such that \(i \le 2 \cdot Q\) and \(\operatorname{lhcenter}(z + (p + 1) + Q, p + Q + i) = \operatorname{not}(\operatorname{lhcenter}(z + (p + 1) + Q, p + Q + i + Q))\).
\end{theorem}
\begin{proof}
The goal follows from \(\operatorname{existslhrowspatialdisagreeofdoublessettledupto}(Q, z, p)\).
\end{proof}
\begin{theorem}
\label{ll:paper:exists-lh-row-spatial-disagree-ge}
For every natural number \(Q\) and natural number \(N\), if \(1 \le Q\), then there exists a natural number \(r\) such that there exists a natural number \(c\) such that there exists a natural number \(i\) such that \(N \le r\) and \(i \le 2 \cdot Q \land \operatorname{lhcenter}(r, c + i) = \operatorname{not}(\operatorname{lhcenter}(r, c + i + Q))\).
\end{theorem}
\begin{proof}
The goal follows from \(\operatorname{existslhrowspatialdisagreege}(Q, N)\).
\end{proof}
\begin{theorem}
\label{ll:paper:lhcenter-shift-agree-iff-spatial-antiagree-of-diag-flip}
For every natural number \(Q\) and natural number \(t\) and natural number \(x\), if \(\operatorname{lhcenter}(t + Q, x + Q) = \operatorname{not}(\operatorname{lhcenter}(t, x))\), then \(\operatorname{lhcenter}(t + Q, x) = \operatorname{lhcenter}(t, x)\) if and only if \(\operatorname{lhcenter}(t + Q, x) = \operatorname{not}(\operatorname{lhcenter}(t + Q, x + Q))\).
\end{theorem}
\begin{proof}
The goal follows from \(\operatorname{lhcentershiftagreeiffspatialantiagreeofdiagflip}(Q, t, x)\).
\end{proof}
\begin{theorem}
\label{ll:paper:lhcenter-shift-agree-iff-spatial-agree-of-diag-eq}
For every natural number \(Q\) and natural number \(t\) and natural number \(x\), if \(\operatorname{lhcenter}(t + Q, x + Q) = \operatorname{lhcenter}(t, x)\), then \(\operatorname{lhcenter}(t + Q, x) = \operatorname{lhcenter}(t, x)\) if and only if \(\operatorname{lhcenter}(t + Q, x) = \operatorname{lhcenter}(t + Q, x + Q)\).
\end{theorem}
\begin{proof}
The goal follows from \(\operatorname{lhcentershiftagreeiffspatialagreeofdiageq}(Q, t, x)\).
\end{proof}
\begin{theorem}
\label{ll:paper:lhcenter-shift-agree-iff-spatial-antiagree-on-doubling-diag}
For every natural number \(Q\) and natural number \(z\) and natural number \(c\), if \(\operatorname{zerodiagonal}(z)\), then if \(\operatorname{doublesat}(Q, z)\), then if \(\operatorname{settledupto}(z + 1, z + 1 + c)\), then \(\operatorname{lhcenter}(z + 1 + c + Q, c) = \operatorname{lhcenter}(z + 1 + c, c)\) if and only if \(\operatorname{lhcenter}(z + 1 + c + Q, c) = \operatorname{not}(\operatorname{lhcenter}(z + 1 + c + Q, c + Q))\).
\end{theorem}
\begin{proof}
The goal follows from \(\operatorname{lhcentershiftagreeiffspatialantiagreeondoublingdiag}(Q, z, c)\).
\end{proof}
\begin{theorem}
\label{ll:paper:lhcenter-shift-agree-iff-spatial-agree-on-black-diag}
For every natural number \(Q\) and natural number \(z\) and natural number \(c\), if \(\operatorname{zerodiagonal}(z)\), then if \(\operatorname{settledupto}(z + 2, z + 2 + c)\), then \(\operatorname{lhcenter}(z + 2 + c + Q, c) = \operatorname{lhcenter}(z + 2 + c, c)\) if and only if \(\operatorname{lhcenter}(z + 2 + c + Q, c) = \operatorname{lhcenter}(z + 2 + c + Q, c + Q)\).
\end{theorem}
\begin{proof}
The goal follows from \(\operatorname{lhcentershiftagreeiffspatialagreeonblackdiag}(Q, z, c)\).
\end{proof}
\begin{theorem}
\label{ll:paper:lhcenter-shift-agree-iff-spatial-antiagree-on-reset-diag}
For every natural number \(Q\) and natural number \(z\) and natural number \(c\), if \(\operatorname{zerodiagonal}(z)\), then if \(\operatorname{doublesat}(Q, z)\), then if \(\operatorname{settledupto}(z + 3, z + 3 + c)\), then \(\operatorname{lhcenter}(z + 3 + c + Q, c) = \operatorname{lhcenter}(z + 3 + c, c)\) if and only if \(\operatorname{lhcenter}(z + 3 + c + Q, c) = \operatorname{not}(\operatorname{lhcenter}(z + 3 + c + Q, c + Q))\).
\end{theorem}
\begin{proof}
The goal follows from \(\operatorname{lhcentershiftagreeiffspatialantiagreeonresetdiag}(Q, z, c)\).
\end{proof}
\begin{theorem}
\label{ll:paper:lhcenter-shift-agree-iff-spatial-agree-below-doubling}
For every natural number \(Q\) and natural number \(z\) and natural number \(t\) and natural number \(c\), if \(\operatorname{doublesat}(Q, z)\), then if \(t - c \le z\), then if \(\operatorname{settledupto}(z, t)\), then \(\operatorname{lhcenter}(t + Q, c) = \operatorname{lhcenter}(t, c)\) if and only if \(\operatorname{lhcenter}(t + Q, c) = \operatorname{lhcenter}(t + Q, c + Q)\).
\end{theorem}
\begin{proof}
The goal follows from \(\operatorname{lhcentershiftagreeiffspatialagreebelowdoubling}(Q, z, t, c)\).
\end{proof}
\begin{theorem}
\label{ll:paper:lhcenter-diag-shift-two-mul-periodic-le-add-three-of-doubles}
For every natural number \(Q\) and natural number \(z\) and natural number \(t\) and natural number \(c\), if \(\operatorname{zerodiagonal}(z)\), then if \(\operatorname{doublesat}(Q, z)\), then if \(c \le t\), then if \(t - c \le z + 3\), then if \(\operatorname{settledupto}(z + 3, t)\), then \(\operatorname{lhcenter}(t + (Q + Q), c + (Q + Q)) = \operatorname{lhcenter}(t, c)\).
\end{theorem}
\begin{proof}
The goal follows from \(\operatorname{lhcenterdiagshifttwomulperiodicleaddthreeofdoubles}(Q, z, t, c)\).
\end{proof}
\begin{theorem}
\label{ll:paper:lhbandagree-two-mul-of-doubles-and-single-lh-row-spatial-agree}
For every natural number \(Q\) and natural number \(z\) and natural number \(t\) and natural number \(p0\) and natural number \(L\), if \(\operatorname{zerodiagonal}(z)\), then if \(\operatorname{doublesat}(Q, z)\), then if \(t - p0 \le z + 3\), then if \(p0 + L \le t\), then if \(\operatorname{settledupto}(z + 3, t)\), then if for every natural number \(i\), if \(i \le L\), then \(\operatorname{lhcenter}(t + (Q + Q), p0 + i) = \operatorname{lhcenter}(t + (Q + Q), p0 + i + (Q + Q))\), then \(\operatorname{lhbandagree}(t, Q + Q, p0, L)\).
\end{theorem}
\begin{proof}
The goal follows from \(\operatorname{lhbandagreetwomulofdoublesandsinglelhrowspatialagree}(Q, z, t, p0, L)\).
\end{proof}
\end{document}
